\documentclass[11pt,a4paper,twoside]{book}
% ===================================================================
%  1. ΒΑΣΙΚΑ ΠΑΚΕΤΑ & ΓΛΩΣΣΑ
% ===================================================================
\usepackage[english,greek]{babel}
% ΣΗΜΕΙΩΣΗ: Χρήση του parskip για αυτόματο κενό μεταξύ παραγράφων,
% αντί για χειροκίνητο μηδενισμό της εσοχής (\parindent).
\usepackage{parskip}
% ===================================================================
%  2. ΓΡΑΜΜΑΤΟΣΕΙΡΕΣ (Απαιτείται XeLaTeX ή LuaLaTeX)
% ===================================================================
% ΣΗΜΕΙΩΣΗ: Τα inputenc και fontenc είναι περιττά με το XeLaTeX.
% Το xunicode είναι παρωχημένο και η λειτουργικότητά του
% είναι ενσωματωμένη στο fontspec.
\usepackage[no-math]{fontspec}
\usepackage{xgreek}
\let\hbar\relax % Για συμβατότητα με το πακέτο mtpro2
\defaultfontfeatures{Mapping=tex-text,Scale=MatchLowercase}
% --- Κύριες γραμματοσειρές ---
\setmainfont[Mapping=tex-text,Numbers=Lining,Scale=1.0,BoldFont={Minion Pro Bold}]{Minion Pro}
\newfontfamily\titlefont{Century Gothic Bold}
\newfontfamily\scfont{Nimbus Roman}
\newfontfamily\notoserif{Noto Serif}
\newfontfamily\urwgothic{URW Gothic}
% --- Προσαρμοσμένες εντολές για γραμματοσειρές ---
% ΣΗΜΕΙΩΣΗ: Αντικατάσταση της παλιάς εντολής \font με σύγχρονες
% εντολές του LaTeX για καλύτερη διαχείριση μεγέθους.
\newcommand{\kefalaio}[1]{{\notoserif\bfseries\fontsize{36}{42}\selectfont #1}}
\newcommand{\ArKef}[1]{{\notoserif\bfseries\itshape\fontsize{72}{80}\selectfont #1}}
\newcommand{\OnKef}[1]{{\notoserif\fontsize{40}{46}\selectfont #1}}
\newcommand{\OnPar}[1]{{\notoserif\bfseries\fontsize{18}{22}\selectfont #1}}
\newcommand{\onoma}[1]{{\urwgothic\bfseries\normalsize #1}}
\newcommand{\Askhseis}[1]{{\urwgothic\bfseries\fontsize{16}{19}\selectfont #1}}
% ===================================================================
%  3. ΓΕΩΜΕΤΡΙΑ ΣΕΛΙΔΑΣ
% ===================================================================
\usepackage[
    a4paper,
    twoside,
    inner=1.5cm,
    top=3cm,
    bottom=2cm,
    textwidth=18cm,
%    marginparsep=5mm,
%    marginparwidth=5.5cm,
    ]{geometry}
% ===================================================================
%  4. ΧΡΩΜΑΤΑ
% ===================================================================
\usepackage[dvipsnames,table,x11names]{xcolor}
\definecolor{maincolor}{RGB}{35,150,210}
\definecolor{secondarycolor}{RGB}{210,20,210}
% ===================================================================
%  5. ΜΑΘΗΜΑΤΙΚΑ
% ===================================================================
\usepackage{mathtools} % Το mathtools φορτώνει αυτόματα το amsmath
\let\myBbbk\Bbbk
\let\Bbbk\relax
\usepackage[amsbb,subscriptcorrection,zswash,mtpcal,mtphrb,mtpfrak]{mtpro2}
\usepackage{mathimatika,venndiagram,diffcoeff}
\allowdisplaybreaks
\providecommand\PARENS[1]{\left(#1\right)}
\makeatletter
\newsavebox{\mtp@matrixbox}
\renewenvironment{pmatrix}{%
  \matrix@check\pmatrix
  \setbox\mtp@matrixbox=\hbox\bgroup$\begin{matrix}%
}{%
  \end{matrix}$\egroup
  \PARENS{\copy\mtp@matrixbox}%
}
\renewenvironment{pmatrix*}[1][c]{%
  \matrix@check\pmatrix
  \setbox\mtp@matrixbox=\hbox\bgroup$\begin{matrix*}[#1]%
}{%
  \end{matrix*}$\egroup
  \PARENS{\copy\mtp@matrixbox}%
}
\makeatother
\DeclareMathSizes{10.95}{10.95}{7}{5}
\DeclareMathSizes{6}{6}{3.8}{2.7}
\DeclareMathSizes{8}{8}{5.1}{3.6}
\DeclareMathSizes{9}{9}{5.8}{4.1}
\DeclareMathSizes{10}{10}{6.4}{4.5}
\DeclareMathSizes{12}{12}{7.7}{5.5}
\DeclareMathSizes{14.4}{14.4}{9.2}{6.5}
\DeclareMathSizes{17.28}{17.28}{11}{7.9}
\DeclareMathSizes{20.74}{20.74}{13.3}{9.4}
\DeclareMathSizes{24.88}{24.88}{16}{11.3}
\makeatletter
\AtBeginDocument{
\check@mathfonts
\fontdimen16\textfont2=2.5pt
\fontdimen17\textfont2=2.5pt
\fontdimen14\textfont2=4.5pt
\fontdimen13\textfont2=4.5pt}
\makeatother
\DeclareMathOperator{\tr}{tr}
\DeclareMathOperator{\adj}{adj}
\DeclareMathOperator{\diag}{diag}
% ===================================================================
%  6. ΓΡΑΦΙΚΑ & ΣΧΗΜΑΤΑ
% ===================================================================
\usepackage{graphicx,tikz,pgfplots,tkz-euclide,tkz-tab}
\usetikzlibrary{shadows,calc,fadings,patterns}
\pgfplotsset{compat=newest}
\tkzSetUpPoint[size=2.8,fill=white]
\makeatletter
\xpatchcmd{\tkzTabLine}
{\node at (Z\thetkz@cnt@impair\thetkz@cnt@lg){0};} % search
{\node[fill=white,inner sep=.5mm] at (Z\thetkz@cnt@impair\thetkz@cnt@lg){0};} % replace
{}{}
\makeatother
% ΣΗΜΕΙΩΣΗ: Χρήση της σύγχρονης εντολής \tikzset αντί για \tikzstyle.
\tikzset{
    >=latex,
    l3/.style={line width=0.3mm},
    l4/.style={line width=0.4mm},
    labelbox/.style args={#1}{align=center,draw=#1,fill=#1!30!white,rectangle,rounded corners=1,font=\tiny\linespread{0.8}\selectfont}
}
\usetikzlibrary{positioning, shapes.geometric, arrows.meta}
\usepackage{tikz-3dplot}
\newcommand{\DrawSphereInCylinder}[3]{\begin{tikzpicture}[declare function={R=#2;% radius of sphere
        r=#3;% radius of cylinder
        theta=#1;% view angle
    }]
    \pgfmathsetmacro{\myalpha}{asin(r/R)}
    \pgfmathsetmacro{\myh}{R*cos(\myalpha)}
    \pgfmathtruncatemacro{\iflag}{ifthenelse(theta>90-\myalpha,1,0)}
    \tdplotsetmaincoords{theta}{0}
    \ifnum\iflag=0\relax
        \path[ball color=gray] (0,0) circle[radius=R];          
    \else
        \pgfmathsetmacro{\myhprime}{\myh*sin(theta)+r*cos(theta)}
        \pgfmathsetmacro{\mytcrit}{-asin((r*cos(theta)*\myhprime)/(r*r-\myhprime*\myhprime))}
        \pgfmathsetmacro{\mybeta}{90-atan2(\myhprime+r*cos(theta)*sin(\mytcrit),r*cos(\mytcrit))}
        \pgfmathsetmacro{\mytcrit}{asin((cos(theta)*\myhprime)/(r*sin(theta)*sin(theta)))}
        \path[ball color=secondarycolor]  (90-\mybeta:R) arc[start angle=90-\mybeta,end angle=-90+\mybeta,radius=R]
        [tdplot_main_coords] 
        -- ({r*cos(-\mytcrit)},{-r*sin(\mytcrit)},-\myh)
        arc[start angle=-\mytcrit,end angle=-180+\mytcrit,radius=r]
        [tdplot_screen_coords]
        -- (-90-\mybeta:R) arc[start angle=-90-\mybeta,end angle=-270+\mybeta,radius=R];
    \fi
    \begin{scope}[tdplot_main_coords]
        \ifnum\iflag=0\relax
            \draw[dashed] (0,0,-\myh) circle[radius=r] (-r,0,\myh) edge (-r,0,-\myh) (r,0,\myh) edge (r,0,-\myh);
        \else 
            \draw[dashed] ({r*cos(-\mytcrit)},{-r*sin(\mytcrit)},-\myh)
            arc[start angle=-\mytcrit,end angle=180+\mytcrit,radius=r]
            (-r,0,\myh) edge (-r,0,-\myh) (r,0,\myh) edge (r,0,-\myh);
        \fi    
        \clip (0,0,\myh) circle[radius=r];
        \path[left color=secondarycolor!50!black,right color=secondarycolor,middle color=secondarycolor!20] (r/2,0,\myh) circle[radius=1.5*r];
        \fill[white] (0,0,-\myh) circle[radius=r];  
    \end{scope}
    \path (current bounding box.north) node[above]{};
\end{tikzpicture}}
% ===================================================================
%  7. ΔΙΑΦΟΡΑ ΠΑΚΕΤΑ
% ===================================================================
\usepackage{gensymb,fontawesome5,eurosym,titletoc,sidenotes,multicol,adjmulticol,siunitx,fullwidth,svg,wrap-rl}
\usepackage[framemethod=tikz]{mdframed}
\usepackage[explicit]{titlesec}
% ===================================================================
%  8. ΚΕΦΑΛΙΔΑ & ΥΠΟΣΕΛΙΔΟ (fancyhdr)
% ===================================================================
\usepackage{fancyhdr}
\pagestyle{fancy}
%\fancyheadoffset[$LE$]{\marginparwidth+\marginparsep}
%\fancyheadoffset[$RO$]{\marginparwidth+\marginparsep}
\renewcommand{\headrulewidth}{\iftopfloat{0pt}{.5pt}}
\renewcommand{\chaptermark}[1]{\markboth{#1}{}}
\renewcommand{\sectionmark}[1]{\markright{\it\thesection\ #1}}
\fancyhf{}
\fancyhead[LE]{\thepage\ $\cdot$\ \nouppercase{\leftmark}}
\fancyhead[RO]{\nouppercase{\rightmark} $\cdot$\ \thepage}
\fancypagestyle{plain}{%
\fancyhead{}
\renewcommand{\headrulewidth}{0pt}}
% ===================================================================
%  9. ΟΡΙΣΜΟΙ - ΘΕΩΡΗΜΑΤΑ - ΔΙΑΦΟΡΑ ΠΛΑΙΣΙΑ (tcolorbox)
% ===================================================================
% ΣΗΜΕΙΩΣΗ: Ολόκληρη η ενότητα έχει αναδομηθεί με τη βιβλιοθήκη 'theorems'
% του tcolorbox. Αυτό κάνει τον κώδικα δραματικά πιο καθαρό,
% ευέλικτο και εύκολο στη συντήρηση.
\usepackage[most]{tcolorbox}
\tcbuselibrary{skins,theorems,breakable}
% $\varSigmaTY\varLambda $ ΛΥΣΗΣ - ΑΠΟΔΕΙΞΗΣ
\newcounter{lysh}[chapter]
\renewcommand{\thelysh}{
	\bf{\titlefont{\arabic{lysh}}}
}
\newcommand{\lysh}[2]{\refstepcounter{lysh}{\vspace{4mm}\titlefont{\faCheck\ {\large Λύση άσκησης  1}.\thelysh\ - \textcolor{maincolor!80!black}{\titlefont{Θέμα #2ο (#1)}}}}}
\newcommand{\apodeiksi}{\textcolor{secondarycolor!80!black}{\titlefont{\faPen\ \textbf{ΑΠΟΔΕΙΞΗ}\\}}}
% Κάρτα πληροφοριών με φωτογραφία
\newcommand{\Info}[5]{
\begin{tcolorbox}[
    enhanced,
    width=5.5cm,
    arc=2mm,
    colback=white,
    colframe=black,
    boxrule=0pt,
    bottom=0pt, top=0pt, left=0pt, right=0pt,
	enhanced standard,
	fuzzy shadow={2mm}{-2mm}{0mm}{0.4mm}{black!50!white},
	boxsep=0mm
]
% --- Εικόνα χωρίς καθόλου περιθώρια ---
\begin{tikzpicture}
\node[
    inner sep=0pt,
    outer sep=0pt,
    rounded corners=2mm,
    clip
] {\includegraphics[width=\linewidth,height=4cm]{#5}};
\end{tikzpicture}
% --- Κείμενο ---
\begin{tcolorbox}[
    colback=white,
    boxrule=0pt,
    top=2mm,bottom=0mm,left=3mm,right=3mm,
    sharp corners=all
]
    \begin{center}
    {\bfseries\LARGE \titlefont{#1}} \\[4pt]
	{\color{secondarycolor!80!black} \titlefont{#2}}
	\end{center}
    {\small #3}
\end{tcolorbox}
% --- Kάτω μπάρα ---
\begin{tcolorbox}[
    colback=maincolor!90!black,
    colupper=white,
    boxrule=0pt,
    bottom=3pt,top=3pt,left=2mm,right=2mm,
    rounded corners=south,
	sharp corners=north,
    arc=2mm
]
\titlefont{#4}
\end{tcolorbox}
\end{tcolorbox}}
% ===================================================================
%  10. ΛIΣTEΣ (enumitem)
% ===================================================================
\usepackage{enumitem,moreenum}
\newlist{rlist}{enumerate}{3} \setlist[rlist]{itemsep=0mm,label=\roman*.}
\newlist{alist}{enumerate}{3} \setlist[alist]{itemsep=0mm,label=\alph*.}
\newlist{balist}{enumerate}{3} \setlist[balist]{itemsep=0mm,label=\bf\alph*.}
\newlist{Alist}{enumerate}{3} \setlist[Alist]{itemsep=0mm,label=\Alph*.}
\newlist{bAlist}{enumerate}{3} \setlist[bAlist]{itemsep=0mm,label=\bf\Alph*.}
\newlist{askhseis}{enumerate}{3} \setlist[askhseis]{label={\Large\thesection}.\arabic*.}
\renewcommand{\textstigma}{\textsigma\texttau}
%\renewcommand{\textdexiakeraia}{}
\AtBeginDocument{\renewcommand{\textdexiakeraia}{}}
\setlist[itemize]{itemsep=0mm}
% --- Στυλ άσκησης ---
\newcounter{askhsh}[chapter]
\renewcommand{\theaskhsh}{
	\bf{\titlefont{\Large{\thechapter}}.\arabic{askhsh}}
}
\newcommand{\Askhsh}[2]{\refstepcounter{askhsh}{\vspace{4mm}\titlefont{{\large Άσκηση}\theaskhsh} - \textcolor{maincolor}{ΘΕΜΑ #2}} ({\small{\textbf{Κωδικός:} #1}})}{}
% --- Περιβάλλον Ύλης ---
\newenvironment{ylh}[1]{\raggedright
\begin{tcolorbox}[
breakable,
halign=flush left,
enhanced standard,
titlerule=0pt,
toptitle=1mm,
toprule=0pt, 
rightrule=0pt, 
bottomrule=0pt,
title={\titlefont{Ύλη άσκησης}},
colback=white,
opacityback=0,
left=1mm,
top=0mm,
bottom=2mm,
boxrule=0pt,
colframe=white,
%borderline west={0mm}{0pt}{maincolor},
borderline north={0.2mm}{0pt}{maincolor},
borderline south={0.2mm}{0pt}{maincolor},
leftrule=2mm,
sharp corners,
coltitle=black
]
\begin{multicols}{#1}
\begin{itemize}[leftmargin=5mm, itemsep=0mm, label=\textcolor{maincolor}{{\small\faBook}}]
}{
\end{itemize}
\end{multicols}
\end{tcolorbox}
}
% ===================================================================
%  11. BΙBΛΙOΓΡAΦΙA (biblatex)
% ===================================================================
\usepackage{csquotes} % Για σωστή διαχείριση εισαγωγικών, συνιστάται από το biblatex
\usepackage[backend=biber,style=alphabetic,sorting=ynt]{biblatex}
% ===================================================================
%  12. ΠΙNAKEΣ (tabularray, longtable, caption)
% ===================================================================
\usepackage{tabularx,longtable,tabularray,capt-of,caption,multirow}
\DeclareTblrTemplate{caption}{nocaptemplate}{}
\DeclareTblrTemplate{capcont}{nocaptemplate}{}
\DeclareTblrTemplate{contfoot}{nocaptemplate}{}
\NewTblrTheme{mytabletheme}{
\SetTblrTemplate{caption}{nocaptemplate}{}
\SetTblrTemplate{capcont}{nocaptemplate}{}
\SetTblrTemplate{contfoot}{nocaptemplate}{}}

\NewTblrEnviron{myhtblr}
\SetTblrOuter[myhtblr]{theme=mytabletheme}
\SetTblrInner[myhtblr]{
    hlines = {1pt,white},
	vlines = {0.5pt,white},
	column{odd} = {maincolor!40,c},
	column{even} = {maincolor!10,c},
	column{1} = {maincolor,c,fg={white}}
}

\NewTblrEnviron{mytblr}
\SetTblrStyle{firsthead}{font=\bfseries}
\SetTblrStyle{firstfoot}{fg=secondarycolor}
\SetTblrOuter[mytblr]{theme=mytabletheme}
\SetTblrInner[mytblr]{
    rowspec={t{7mm}},
    columns = {c},
    width = 0.85\linewidth,
    row{odd} = {bg=maincolor!20!white,fg=black,ht=8mm},
    row{even} = {bg=gray!10!white,fg=black,ht=8mm},
    hlines={white},
    vlines={white},
    row{1} = {bg=maincolor, fg=white, font=\large\bfseries\titlefont},
    rowhead = 1,
    hline{2} = {.7mm},
}
\captionsetup[figure]{format=hang,labelsep=period,name={\titlefont{\textbf{Σχήμα}}}}
\renewcommand\thefigure{{\bf\titlefont{\thechapter.\arabic{figure}}}}
% ===================================================================
%  13. ΠEΡΙEΧOMENA (etoc & titletoc)
% ===================================================================
% 2. $NEO$: Μορφοποίηση για άτιτλα κεφάλαια (Περιεχόμενα, Βιβλιογραφία)
\titleformat{name=\chapter,numberless}[block]
  {\normalfont\sffamily}
  {}
  {0pt}
  {%
    \begin{tikzpicture}[remember picture, overlay]
      \fill[maincolor] (current page.north west) rectangle ($(current page.north east) + (0,-4.5cm)$);%
      \fill[maincolor!50!black] (current page.north west) rectangle ($(current page.north west) + (2cm,-4.5cm)$);%
      \fill[maincolor!50] ($(current page.north west) + (0,-4.5cm)$) rectangle ($(current page.north east) + (0,-4.65cm)$);%
      \node[anchor=north west, text=white, font=\titlefont\Huge, text width=\paperwidth-5cm] 
        at ($(current page.north west) + (2.4cm,-2.5cm)$) {#1};
    \end{tikzpicture}%
  }
\titlespacing*{\chapter}{0pt}{-1cm}{3.5cm}
\titlespacing*{name=\chapter,numberless}{0pt}{-1cm}{3.5cm}
\titleformat{\section}
  {\normalfont\titlefont\Large\color{maincolor}}{\thesection}{1em}{}
\titlecontents{chapter}
  [0em] % Αριστερό περιθώριο
  {\addvspace{1.5ex}\titlefont\large\bfseries\color{maincolor}} % Κενό και γραμματοσειρά
  {\color{maincolor!50}\thecontentslabel\hspace*{1em}\color{maincolor}} % Αριθμός κεφαλαίου σε maincolor
  {} % Για τα άτιτλα (δεν χρειάζεται κάτι)
  {\hfill\color{maincolor!50}\bfseries\contentspage} % Αριθμός σελίδας δεξιά
% Εμφάνιση των ενοτήτων στον πίνακα περιεχομένων
\titlecontents{section}
  [2.5em] % Εσοχή
  {\titlefont\color{maincolor!50!black}}
  {\thecontentslabel\hspace*{1em}}
  {}
  {\titlerule*[0.5pc]{\color{gray}.}\color{maincolor!50!black}\contentspage} % Τελίτσες μέχρι τον αριθμό σελίδας
\usepackage{tikzpagenodes,eso-pic}
\newcommand{\chaptertoc}{%
\AddToShipoutPictureFG*{%
\begin{tikzpicture}[overlay]
  \coordinate ($NE$) at ($(current page text area.north east)+(-59mm,-14.5cm)$);
  \coordinate ($SE$) at ($(current page text area.south east)+(-59mm,0)$);
  \draw[secondarycolor!80!black] ([xshift=0.5\marginparsep]$NE$) -- ([xshift=0.5\marginparsep]$SE$);
\end{tikzpicture}%
}%
\marginpar{%
\vspace*{8mm}
\begin{minipage}{\dimexpr\marginparwidth+4mm}
\begin{flushleft}
\localtableofcontents
\end{flushleft}
\end{minipage}%
}%
}
% ===================================================================
%  14. $\varSigmaTY\varLambda $ KEΦAΛAΙOY & ENOTHTΩN
% ===================================================================
\newcommand*\chapterlabel{}
\newcommand{\fancychapter}{%
\titleformat{\chapter}
{\normalfont\huge }
{\gdef\chapterlabel{\thechapter\ }}{-2cm}
{    \begin{tikzpicture}[remember picture, overlay]
      \fill[maincolor] (current page.north west) rectangle ($(current page.north east) + (0,-4.5cm)$);%
      \fill[maincolor!50!black] (current page.north west) rectangle ($(current page.north west) + (2cm,-4.5cm)$);%
      \fill[maincolor!50] ($(current page.north west) + (0,-4.5cm)$) rectangle ($(current page.north east) + (0,-4.65cm)$);%
      \node[anchor=south west, text=maincolor!50, font=\titlefont\Large] 
        at ($(current page.north west) + (2.5cm,-2.4cm)$) {ΚΕΦΑΛΑΙΟ \thechapter};%
      \node[anchor=north west, text=white, font=\titlefont\Huge, text width=\paperwidth-5cm] 
        at ($(current page.north west) + (2.4cm,-2.5cm)$) {##1};
    \end{tikzpicture}%
}
}
\titleformat{\section}[hang]
  {\normalfont\titlefont\Large\color{maincolor}}
  {\textcolor{maincolor!50}{\rule[-0.15em]{3pt}{1.1em}}\hspace{0.5em}\textcolor{maincolor!50!black}{\thesection}}
  {0.5em}
  {#1}
% --- 4. Υπο-ενότητες (Subsections) ---
% Πιο διακριτικές, με τον αριθμό σε maincolor και τον τίτλο σε σκούρο μπλε
\titleformat{\subsection}[hang]
  {\normalfont\titlefont\large\color{maincolor!50!black}}
  {\textcolor{maincolor!50}{\thesubsection}}
  {0.5em}
  {#1}
\apptocmd{\mainmatter}{\fancychapter}{}{}
\titlespacing{\chapter}{0pt}{0cm}{1cm}
\titleformat{\paragraph}{\large}{}{0em}{\textcolor{maincolor}{\faSquare\ \ \titlefont{\bmath{#1}}}}
\titlespacing{\paragraph}{0mm}{2mm}{2mm}
% ===================================================================
%  15. BOHΘHTΙKEΣ ENTOΛEΣ
% ===================================================================
\newcommand{\eng}[1]{\selectlanguage{english}#1\selectlanguage{greek}}
\newcommand{\tss}[1]{\textsuperscript{#1}}
\newcommand{\tssL}[1]{\MakeLowercase{\textsuperscript{#1}}}
\newcommand{\full}[1]{\begin{mdframed}[outermargin=\dimexpr-\marginparwidth-\marginparsep\relax,innerleftmargin=0mm,innerrightmargin=0mm,hidealllines=true] #1 \end{mdframed}}
\newcommand{\fulltwoc}[1]{\begin{adjmulticols}{2}{5mm}{\dimexpr-\marginparwidth-\marginparsep+5mm\relax} #1 \end{adjmulticols}}
\newcommand{\monades}[1]{
\hspace*{\fill}
\textbf{\textit{\textcolor{maincolor}{Μονάδες #1}}}}
% ===================================================================
%  TEΛOΣ PREAMBLE
% ===================================================================
\begin{document}
\pagestyle{empty}
\frontmatter
\pagenumbering{gobble}
\newpage
\newgeometry{inner=2.00cm, top=3.00cm, bottom=2.00cm,outer=1.50cm,headheight=14pt}
\tableofcontents
\restoregeometry
\mainmatter
\pagestyle{fancy}

\chapter{Τράπεζα Θεμάτων - Γεωμετρία}

% --- Άσκηση 1706 (1706.tex) ---
\Askhsh{1706}{4}
\begin{ylh}{1}
    \item 3.1. Είδη και στοιχεία τριγώνων
    \item 3.2. 1ο Κριτήριο ισότητας τριγώνων
    \item 5.7. Βαρύκεντρο τριγώνου.
\end{ylh}
Έστω τρίγωνο $AB\varGamma $ και $\mu_\beta, \mu_\gamma$ οι διάμεσοι του τριγώνου που αντιστοιχούν στις πλευρές $\beta$ και $\gamma$ αντίστοιχα. Δίνεται η ακόλουθη πρόταση:\\[0.2cm]
$\Pi$: Αν το τρίγωνο $AB\varGamma $ είναι ισοσκελές με $\beta = \gamma$, τότε οι διάμεσοι $\mu_\beta, \mu_\gamma$ είναι ίσες.
\begin{alist}
    \item Να εξετάσετε αν ισχύει η πρόταση $\Pi$, αιτιολογώντας την απάντησή σας. \monades{10}
    \item Να διατυπώσετε την αντίστροφη πρόταση της $\Pi$ και να εξετάσετε αν ισχύει αιτιολογώντας την απάντησή σας. \monades{10}
    \item Στην περίπτωση που οι δυο προτάσεις, η $\Pi$ και η αντίστροφή της ισχύουν, να τις διατυπώσετε ως ενιαία πρόταση. \monades{5}
\end{alist}

\vspace{4mm}

% --- Άσκηση 1707 (1707.tex) ---
\Askhsh{1707}{4}
\begin{ylh}{1}
    \item 3.6. Κριτήρια ισότητας ορθογώνιων τριγώνων.
\end{ylh}
Στο τρίγωνο $AB\varGamma $ ($AB < A\varGamma $) του παρακάτω σχήματος, η κάθετη στο μέσο $ M$ της $B\varGamma $ τέμνει την προέκταση της διχοτόμου $ A\varDelta  $ στο σημείο $ E$. Αν $\varTheta ,Z$ είναι οι προβολές του $E$ στις  $AB,  A\varGamma $, να αποδείξετε ότι:
\begin{alist}
    \item Το τρίγωνο $EB\varGamma $ είναι ισοσκελές. \monades{5}
    \item Τα τρίγωνα $\varTheta  B E$ και $ Z\varGamma  E$ είναι ίσα. \monades{8}
    \item $A\hat{\varGamma } E + A\hat{B} E = 180^\circ$. \monades{12}
\end{alist}

\begin{center}
\begin{tikzpicture}[scale=1]
% Triangle $AB\varGamma $ with $AB$ < $A\varGamma $
\tkzDefPoint(0,0){B}
\tkzDefPoint(6,0){Gamma}
\tkzDefTriangle[two angles = 65 and 40](B,Gamma) \tkzGetPoint{A}

% Angle bisector $A\varDelta  $: Δ is on $B\varGamma $
\tkzDefLine[bisector](B,A,Gamma) \tkzGetPoint{bisect_dir}
\tkzInterLL(A,bisect_dir)(B,Gamma) \tkzGetPoint{Delta}

% Midpoint M of $B\varGamma $
\tkzDefMidPoint(B,Gamma) \tkzGetPoint{M}

% Perpendicular to $B\varGamma $ at M
\tkzDefLine[orthogonal=through M](B,Gamma) \tkzGetPoint{perp_dir}

% E is the intersection of the bisector extension $A\varDelta  $ with the perpendicular at M
\tkzInterLL(A,Delta)(M,perp_dir) \tkzGetPoint{E}

% H is the projection of E onto line $AB$
\tkzDefPointBy[projection=onto A--B](E) \tkzGetPoint{H}

% Z is the projection of E onto line $A\varGamma $
\tkzDefPointBy[projection=onto A--Gamma](E) \tkzGetPoint{Z}

% Draw triangle
\draw[l3] (A)--(B)--(Gamma)--cycle;

% Draw bisector and extension to E
\draw[l3] (A)--(Delta);
\draw[l3] (Delta)--(E);

% Draw perpendicular at M
\draw[l4,maincolor] (M)--(E);

% Draw projections
\draw[l3,maincolor] (E)--(H);
\draw[l3] (E)--(Z);

% Draw segments $BE$ and $\varGamma E$
\draw[l3] (B)--(E);
\draw[l3] (Gamma)--(E);
\draw[l3] (B)--(H);

% Draw points
\tkzDrawPoints(A,B,Gamma,Delta,M,E,H,Z)
\tkzLabelPoint[above](A){$A$}
\tkzLabelPoint[left](B){$B$}
\tkzLabelPoint[below right](Gamma){$\varGamma $}
\tkzLabelPoint[above left](Delta){$\varDelta  $}
\tkzLabelPoint[above](M){$ M$}
\tkzLabelPoint[below](E){$ E$}
\tkzLabelPoint[left](H){$\varTheta $}
\tkzLabelPoint[above right](Z){$ Z$}

% Right angle marks
\tkzMarkRightAngle[size=0.25](B,M,E)
\tkzMarkRightAngle[size=0.25](A,H,E)
\tkzMarkRightAngle[size=0.25](A,Z,E)
\end{tikzpicture}
\end{center}

\vspace{4mm}

% --- Άσκηση 1708 (1708.tex) ---
\Askhsh{1708}{4}
\begin{ylh}{1}
    \item 3.2. 1ο Κριτήριο ισότητας τριγώνων
    \item 3.6. Κριτήρια ισότητας ορθογώνιων τριγώνων
    \item 4.6. Άθροισμα γωνιών τριγώνου.
\end{ylh}
\begin{wrapfigure}[5]{r}{4cm}
\begin{tikzpicture}[scale=1]
\tkzDefPoints{0/0/A, 3/0/B}
\tkzDefTriangle[two angles=90 and 50](A,B)\tkzGetPoint{C}
\tkzMarkRightAngle[size=.3](B,A,C)
\tkzFillAngle[maincolor,size=0.4](C,B,A)
\tkzMarkAngle[size=.4](C,B,A)
\tkzLabelPoint[left](A){$A$}
\tkzLabelPoint[right](B){$B$}
\tkzLabelPoint[above](C){$\varGamma $}
\tkzLabelAngle[pos=.7,font=\scriptsize](C,B,A){$50^\circ$}
\tkzDefLine[altitude](B,A,C) \tkzGetPoint{D}
\tkzDefPointBy[symmetry=center D](B)\tkzGetPoint{E}
\tkzDefPointBy[projection = onto A--E](C)\tkzGetPoint{Z}
\tkzDrawSegments[l3,maincolor](A,D C,Z)
\tkzDrawSegment[dashed](A,Z)
\tkzLabelPoint[above right](D){$\varDelta  $}
\tkzLabelPoint[right](E){$E$}
\tkzLabelPoint[right](Z){$Z$}
\tkzDrawPolygon[l3](A,B,C)
\tkzDrawPoints(A,B,C,D,E,Z)
\end{tikzpicture}
\end{wrapfigure}
Δίνεται ορθογώνιο τρίγωνο $AB\varGamma $ με $(\hat{A} = 90^\circ)$ και $\hat{B} = 50^\circ$, το ύψος του $ A\varDelta  $ και σημείο $ E$ στην $\varDelta  \varGamma $ ώστε $\varDelta  E = B\varDelta  $. Το σημείο $ Z$ είναι η προβολή του $\varGamma $ στην $AE$.
\begin{alist}
    \item Να αποδείξετε ότι:
    \begin{rlist}
        \item Το τρίγωνο $AB E$ είναι ισοσκελές. \monades{6}
        \item $\varGamma \hat{A} E = 10^\circ$. \monades{10}
    \end{rlist}
    \item Να υπολογίσετε τις γωνίες του τριγώνου $ Z\varGamma  E$. \monades{9}
\end{alist}
\vspace{4mm}

% --- Άσκηση 1709 (1709.tex) ---
\Askhsh{1709}{4}
\begin{ylh}{1}
    \item 3.11. Ανισοτικές σχέσεις πλευρών και γωνιών
    \item 4.2. Τέμνουσα δύο ευθειών - Ευκλείδειο αίτημα
    \item 4.6. Άθροισμα γωνιών τριγώνου
    \item 5.2. Παραλληλόγραμμα.
\end{ylh}
Δίνεται τρίγωνο $AB\varGamma $, στο οποίο η εξωτερική του γωνία $\hat{\varGamma }_{\varepsilon\xi}$ είναι διπλάσια της εσωτερικής του γωνίας $\hat{ A}$. Από την κορυφή $ A$ διέρχεται ημιευθεία $Ax \parallel B\varGamma $ στο ημιεπίπεδο $(AB, \varGamma )$. Στην ημιευθεία $Ax$ θεωρούμε σημείο $\varDelta  $ τέτοιο ώστε $A\varDelta = B\varGamma $. Να αποδείξετε ότι:
\begin{alist}
    \item Η $B\varDelta  $ διέρχεται από το μέσο του τμήματος $A\varGamma $. \monades{7}
    \item Η $\varGamma \varDelta  $ είναι διχοτόμος της $\hat{\varGamma }_{\varepsilon\xi}$. \monades{9}
    \item Το τρίγωνο $AB\varGamma $ είναι ισοσκελές. \monades{9}
\end{alist}

\vspace{4mm}

% --- Άσκηση 1723 (1723.tex) ---
\Askhsh{1723}{4}
\begin{ylh}{1}
    \item 3.1. Είδη και στοιχεία τριγώνων
    \item 3.6. Κριτήρια ισότητας ορθογώνιων τριγώνων
    \item 5.6. Εφαρμογές στα τρίγωνα.
\end{ylh}
\begin{wrapfigure}[5]{r}{6cm}
\begin{tikzpicture}[scale=1]
\tkzDefPoints{0/0/B, 5/0/C, 1.3/3/A}
\tkzDrawPolygon[l3](A,B,C)
\tkzDefMidPoint(B,C)\tkzGetPoint{M}
\tkzDefLine[bisector](B,A,C) \tkzGetPoint{d}
\tkzInterLL(A,d)(B,C) \tkzGetPoint{D}
\tkzDefPointBy[projection=onto A--D](B) \tkzGetPoint{E}
\tkzInterLL(A,C)(B,E)\tkzGetPoint{H}
\tkzDrawSegments[l3,maincolor](A,D B,H E,M)
\tkzMarkRightAngle[size=.2](A,E,B)
\tkzLabelPoint[above](A){$A$}
\tkzLabelPoint[left](B){$B$}
\tkzLabelPoint[right](C){$\varGamma $}
\tkzLabelPoint[below](D){$\varDelta  $}
\tkzLabelPoint[above right](E){$E$}
\tkzLabelPoint[right](H){$H$}
\tkzLabelPoint[below](M){$M$}
\tkzMarkAngle[size=.7,arc=ll](B,A,D)
\tkzMarkAngle[arc=ll,size=.5](D,A,C)

\tkzDrawPoints(A,B,C,D,M,E,H)
\end{tikzpicture}
\end{wrapfigure}
Δίνεται τρίγωνο $AB\varGamma $ ($AB < A\varGamma $) και η διχοτόμος του $A\varDelta  $. Φέρουμε από το $B$ κάθετη στην $ A\varDelta  $ που τέμνει την $ A\varDelta  $ στο $ E$ και την πλευρά $A\varGamma $ στο $ H$. Αν $ M$ είναι το μέσο της πλευράς $B\varGamma $, να αποδείξετε ότι:
\begin{alist}
    \item Το τρίγωνο $AB H$ είναι ισοσκελές. \monades{9}
    \item $ E M \parallel  H\varGamma $ \monades{8}
    \item $ E M = \frac{A\varGamma  - AB}{2}$ \monades{8}
\end{alist}

\vspace{4mm}

% --- Άσκηση 1727 (1727.tex) ---
\Askhsh{1727}{4}
\begin{ylh}{1}
    \item 3.1. Είδη και στοιχεία τριγώνων
    \item 3.11. Ανισοτικές σχέσεις πλευρών και γωνιών
    \item 4.2. Τέμνουσα δύο ευθειών - Ευκλείδειο αίτημα
    \item 4.6. Άθροισμα γωνιών τριγώνου
    \item 5.2. Παραλληλόγραμμα
    \item 5.3. Ορθογώνιο
    \item 5.6. Εφαρμογές στα τρίγωνα.
\end{ylh}
\begin{wrapfigure}[5]{r}{7cm}
\begin{tikzpicture}[scale=1]
\tkzDefPoints{0/0/D, 6/0/C, 0/3/A, 3/3/B, 3/0/E}
\tkzDrawPolygon[l3](A,B,C,D)
\tkzDrawSegments[l3](B,E A,C)
\tkzDrawSegments[l4,maincolor](A,E D,B)
\tkzInterLL(A,C)(B,E)\tkzGetPoint{K}
\tkzInterLL(A,E)(B,D)\tkzGetPoint{L}
\tkzDrawSegment[dashed,maincolor](L,K)
\tkzLabelPoint[above](A){$A$}
\tkzLabelPoint[above](B){$B$}
\tkzLabelPoint[right](C){$\varGamma $}
\tkzLabelPoint[left](D){$\varDelta  $}
\tkzLabelPoint[below](E){$E$}
\tkzLabelPoint[above right](K){$K$}
\tkzLabelPoint[below](L){$\varLambda $}
\tkzDrawPolygon[l3](A,B,C,D)
\tkzDrawSegments[l3](B,E A,C)
\tkzDrawSegments[l4,maincolor](A,E D,B)
\tkzMarkRightAngle[size=.3](D,A,B)
\tkzMarkRightAngle[size=.3](E,D,A)
\tkzMarkRightAngle[size=.3](C,E,B)
\tkzDrawPoints(A,B,C,D,E,K,L)
\end{tikzpicture}
\end{wrapfigure}
Δίνεται τραπέζιο $AB\varGamma \varDelta  $ με $\hat{ A} = \hat{\varDelta  } = 90^\circ$, $\varDelta  \varGamma  = 2AB$ και $\hat{B} = 3\hat{\varGamma }$. Από το $B$ φέρνουμε κάθετη στη $\varGamma \varDelta  $ που τέμνει την $A\varGamma $ στο σημείο $ K$ και την $\varGamma \varDelta  $ στο $E$. Επίσης φέρνουμε την $AE$ που τέμνει τη $B\varDelta  $ στο σημείο $\Lambda$.
Να αποδείξετε ότι:
\begin{alist}
    \item $\hat{\varGamma } = 45^\circ$ \monades{8}
    \item $B\varDelta   = A E$ \monades{9}
    \item $ K\Lambda = \frac{1}{4}\varDelta  \varGamma $ \monades{8}
\end{alist}

\vspace{4mm}

% --- Άσκηση 1728 (1728.tex) ---
\Askhsh{1728}{4}
\begin{ylh}{1}
    \item 5.2. Παραλληλόγραμμα
    \item 5.6. Εφαρμογές στα τρίγωνα.
\end{ylh}
Έστω παραλληλόγραμμο $AB\varGamma \varDelta  $. Αν τα σημεία $ E$ και $ Z$ είναι τα μέσα των πλευρών του $ AB$ και $\varGamma \varDelta  $ αντίστοιχα, να αποδείξετε ότι:
\begin{alist}
    \item Το τετράπλευρο $\varDelta   $EB$ Z$ είναι παραλληλόγραμμο. \monades{8}
    \item $A\hat{ E}\varDelta   = B\hat{ Z}\varGamma $ \monades{8}
    \item Οι $\varDelta   E$ και $B Z$ τριχοτομούν τη διαγώνιο $A\varGamma $ του παραλληλογράμμου $ $AB$\varGamma \varDelta  $. \monades{9}
\end{alist}

\vspace{4mm}

% --- Άσκηση 1729 (1729.tex) ---
\Askhsh{1729}{4}
\begin{ylh}{1}
    \item 3.1. Είδη και στοιχεία τριγώνων
    \item 3.6. Κριτήρια ισότητας ορθογώνιων τριγώνων
    \item 4.6. Άθροισμα γωνιών τριγώνου
    \item 5.3. Ορθογώνιο.
\end{ylh}
Στο ορθογώνιο παραλληλόγραμμο $AB\varGamma \varDelta  $ είναι $\varDelta  \hat{\varGamma }A = 30^\circ$ και $ O$ το κέντρο του. Φέρνουμε $\varDelta   E \perp A\varGamma $.
\begin{alist}
    \item Να αποδείξετε ότι η γωνία $A\hat{\varDelta  }\varGamma $ χωρίζεται από τη $\varDelta   E$ και τη διαγώνιο $\varDelta   B$ σε τρεις ίσες γωνίες. \monades{13}
    \item Φέρνουμε κάθετη στην $A\varGamma $ στο σημείο $ O$ η οποία τέμνει την προέκταση της $ A\varDelta  $ στο $ Z$. Να δείξετε ότι τα τρίγωνα $ A Z O$ και $AB\varGamma $ είναι ίσα. \monades{12}
\end{alist}
\begin{center}
\begin{tikzpicture}[scale=.8]
\tkzDefPoints{0/0/D, 6/0/C, 0/3/A, 6/3/B}
\tkzDrawPolygon[l3](A,B,C,D)
\tkzDrawSegments[l3](A,C B,D)
\tkzDefPointBy[projection=onto A--C](D) \tkzGetPoint{E}
\tkzDrawSegments[l3,maincolor](D,E)
\tkzInterLL(A,C)(B,D)\tkzGetPoint{O}
\tkzDefLine[perpendicular=through O](C,A)\tkzGetPoint{c}
\tkzInterLL(O,c)(A,D)\tkzGetPoint{Z}
\tkzDrawSegments[l3,maincolor](D,Z O,Z)
\tkzMarkRightAngles[size=.3](A,E,D A,O,Z)
\tkzMarkAngle[size=.7](A,C,D)
\tkzLabelAngle[pos=1,font=\scriptsize](A,C,D){$30\degree$}
\tkzLabelPoint[above left](A){$A$}
\tkzLabelPoint[above right](B){$B$}
\tkzLabelPoint[right](C){$\varGamma $}
\tkzLabelPoint[left](D){$\varDelta  $}
\tkzLabelPoint[above right](E){$E$}
\tkzLabelPoint[above](O){$O$}
\tkzLabelPoint[left](Z){$Z$}
\tkzDrawPoints(A,B,C,D,E,O,Z)
\end{tikzpicture}
\end{center}

\vspace{4mm}

% --- Άσκηση 1730 (1730.tex) ---
\Askhsh{1730}{4}
\begin{ylh}{1}
    \item 3.2. 1ο Κριτήριο ισότητας τριγώνων
    \item 4.2. Τέμνουσα δύο ευθειών - Ευκλείδειο αίτημα
    \item 5.2. Παραλληλόγραμμα.
\end{ylh}
Έστω ότι $E$ και $\mathrm{Z}$ είναι τα μέσα των πλευρών $AB$ και $\varGamma \Delta$ παραλληλογράμμου $AB\varGamma \Delta$ αντίστοιχα. Αν για το παραλληλόγραμμο $ AB\varGamma \Delta$ επιπλέον ισχύει $AB > A\Delta$, να εξετάσετε αν είναι αληθείς ή όχι οι ακόλουθοι ισχυρισμοί:\\[0.2cm]
Ισχυρισμός 1: Το τετράπλευρο $\varDelta EBZ$ είναι παραλληλόγραμμο. \\
Ισχυρισμός 2: $A\hat{ E}\varDelta   = B\hat{ Z}\varGamma $. \\
Ισχυρισμός 3: Οι $\Delta E$ και $B\mathrm{Z}$ είναι διχοτόμοι των απέναντι γωνιών $\hat{\Delta}$ και $\hat{B}$.
\begin{alist}
    \item Στην περίπτωση που θεωρείτε ότι κάποιος ισχυρισμός είναι αληθής να τον αποδείξετε. \monades{16}
    \item Στην περίπτωση που κάποιος ισχυρισμός δεν είναι αληθής, να βρείτε τη σχέση των διαδοχικών πλευρών του παραλληλογράμμου ώστε να είναι αληθής. Να αιτιολογήσετε την απάντησή σας. \monades{9}
\end{alist}

\vspace{4mm}

% --- Άσκηση 1731 (1731.tex) ---
\Askhsh{1731}{4}
\begin{ylh}{1}
    \item 3.2. 1ο Κριτήριο ισότητας τριγώνων
    \item 5.2. Παραλληλόγραμμα.
\end{ylh}
Έστω ότι $E$ και $\mathrm{Z}$ είναι τα μέσα των πλευρών $AB$ και $\varGamma \Delta$ παραλληλογράμμου $ AB\varGamma \Delta$ αντίστοιχα. Αν για το παραλληλόγραμμο $ AB\varGamma \Delta$ επιπλέον ισχύουν $AB > A\Delta$ και γωνία $\hat{A}$ αμβλεία, να εξετάσετε αν είναι αληθείς οι ακόλουθοι ισχυρισμοί:\\[0.2cm]
Ισχυρισμός 1: Το τετράπλευρο $\varDelta EBZ$ είναι παραλληλόγραμμο. \\
Ισχυρισμός 2: Τα τρίγωνα $A\varDelta   E$ και $B Z\varGamma $ είναι ίσα. \\
Ισχυρισμός 3: Τα τρίγωνα $A\varDelta   E$ και $B Z\varGamma $ είναι ισοσκελή.
\begin{alist}
    \item Στην περίπτωση που θεωρείτε ότι κάποιος ισχυρισμός είναι αληθής να τον αποδείξετε. \monades{16}
    \item Στην περίπτωση που κάποιος ισχυρισμός δεν είναι αληθής, να βρείτε τη σχέση των διαδοχικών πλευρών του παραλληλογράμμου ώστε να είναι αληθής. Να αιτιολογήσετε την απάντησή σας. \monades{9}
\end{alist}

\vspace{4mm}

% --- Άσκηση 1733 (1733.tex) ---
\Askhsh{1733}{4}
\begin{ylh}{1}
    \item 3.4. 3ο Κριτήριο ισότητας τριγώνων
    \item 3.7. Κύκλος - Μεσοκάθετος - Διχοτόμος
    \item 5.3. Ορθογώνιο.
\end{ylh}
Έστω $\varepsilon_1, \varepsilon_2$ δυο κάθετες ευθείες που τέμνονται στο $ O$ και τυχαίο σημείο $ M$ του επιπέδου που δεν ανήκει στις ευθείες.
\begin{alist}
    \item Αν $ M_1$ είναι το συμμετρικό του $ M$ ως προς την $\varepsilon_1$ και $ M_2$ το συμμετρικό του $ M_1$ ως προς την $\varepsilon_2$, να αποδείξετε ότι:
    \begin{rlist}
        \item $ O M =  O M_1 =  O M_2$ \monades{6}
        \item Τα σημεία $ M,  O$ και $ M_2$ είναι συνευθειακά. \monades{8}
        \item Το τρίγωνο $ M M_1 M_2$ είναι ορθογώνιο. \monades{6}
    \end{rlist}
    \item Αν $ M_3$ είναι το συμμετρικό σημείο του $ M$ ως προς την $\varepsilon_2$, τι είδους παραλληλόγραμμο είναι το $ M M_1 M_2 M_3$; Να αιτιολογήσετε την απάντησή σας. \monades{5}
\end{alist}

\vspace{4mm}

% --- Άσκηση 1734 (1734.tex) ---
\Askhsh{1734}{4}
\begin{ylh}{1}
    \item 3.2. 1ο Κριτήριο ισότητας τριγώνων
    \item 3.4. 3ο Κριτήριο ισότητας τριγώνων
    \item 5.4. Ρόμβος
    \item 5.5. Τετράγωνο.
\end{ylh}
\begin{wrapfigure}[7]{r}{5cm}
\begin{tikzpicture}[scale=.35]
\tkzDefPoints{0/0/A, 6/0/B, 0/3/D, 6/3/C}
\tkzDrawPolygon[l3](A,B,C,D)
\tkzDefTriangle[equilateral](B,A)\tkzGetPoint{E}
\tkzDefTriangle[equilateral](C,B)\tkzGetPoint{Z}
\tkzDefTriangle[equilateral](D,C)\tkzGetPoint{H}
\tkzDefTriangle[equilateral](A,D)\tkzGetPoint{T}
\tkzMarkRightAngles[size=.3](B,A,D C,B,A D,C,B A,D,C)
\tkzDrawPolygons[l3](A,B,E C,B,Z D,C,H A,D,T)
\tkzDrawPolygon[l3](E,Z,H,T)
\tkzLabelPoint[below left,fill=white,inner sep=0.2mm](A){$A$}
\tkzLabelPoint[below right,fill=white,inner sep=0.2mm](B){$B$}
\tkzLabelPoint[above right,fill=white,inner sep=0.2mm](C){$\varGamma $}
\tkzLabelPoint[above left,fill=white,inner sep=0.2mm](D){$\varDelta  $}
\tkzLabelPoint[below](E){$E$}
\tkzLabelPoint[right](Z){$Z$}
\tkzLabelPoint[above](H){$H$}
\tkzLabelPoint[left](T){$\varTheta $}
\tkzDrawPoints(A,B,C,D,E)
\end{tikzpicture}
\end{wrapfigure}
Δίνεται ορθογώνιο $AB\varGamma \varDelta  $ και έξω από αυτό, κατασκευάζουμε τέσσερα ισόπλευρα τρίγωνα $ $AB$ E, B\varGamma  Z, \varGamma \varDelta   H, \varDelta   A\varTheta $.
\begin{alist}
    \item Να αποδείξετε ότι το τετράπλευρο $ E Z H\varTheta $ είναι ρόμβος. \monades{15}
    \item Αν το αρχικό τετράπλευρο $AB\varGamma \varDelta  $ είναι τετράγωνο, τότε το $ E Z H\varTheta $ τι είδους παραλληλόγραμμο είναι; Δικαιολογήστε την απάντησή σας. \monades{10}
\end{alist}

\vspace{4mm}

% --- Άσκηση 1735 (1735.tex) ---
\Askhsh{1735}{4}
\begin{ylh}{1}
    \item 3.7. Κύκλος - Μεσοκάθετος - Διχοτόμος
    \item 5.3. Ορθογώνιο
    \item 5.10. Τραπέζιο.
\end{ylh}
Θεωρούμε ευθεία $(\varepsilon)$ και δυο σημεία $A$ και $B$ εκτός αυτής, τα οποία βρίσκονται στο ίδιο ημιεπίπεδο σε σχέση με την $(\varepsilon)$ έτσι ώστε, η ευθεία $ AB$ να μην είναι κάθετη στην $(\varepsilon)$. Έστω $A'$ και $B'$ τα συμμετρικά σημεία των $ A$ και $B$ αντίστοιχα ως προς την ευθεία $(\varepsilon)$.
\begin{alist}
    \item Να αποδείξετε ότι $AA' \parallel BB'$. \monades{6}
    \item Αν η μεσοκάθετος του $AB$ τέμνει την ευθεία $(\varepsilon)$ στο σημείο $ K$, να αποδείξετε ότι το $ K$ ανήκει και στη μεσοκάθετο του $ A'B'$. \monades{10}
    \item Να βρείτε τη σχέση της ευθείας $AB$ με την ευθεία $(\varepsilon)$ ώστε το τετράπλευρο $ ABB'A'$ να είναι ορθογώνιο. Να αιτιολογήσετε την απάντησή σας. \monades{9}
\end{alist}

\vspace{4mm}

% --- Άσκηση 1736 (1736.tex) ---
\Askhsh{1736}{4}
\begin{ylh}{1}
    \item 4.2. Τέμνουσα δύο ευθειών - Ευκλείδειο αίτημα
    \item 5.9. Μια ιδιότητα του ορθογώνιου τριγώνου
    \item 5.11. Ισοσκελές τραπέζιο.
\end{ylh}
\begin{wrapfigure}[5]{r}{7cm}
\begin{tikzpicture}[scale=0.75]
\tkzDefPoint(0,0){D} % Σημείο \Delta
\tkzDefPoint(8,0){C} % Σημείο \varGamma 
\tkzDefPoint(2, {2*sqrt(3)}){E}
\tkzDefPointBy[homothety=center E ratio 0.45](D) \tkzGetPoint{A}
\tkzDefPointBy[homothety=center E ratio 0.45](C) \tkzGetPoint{B}
\tkzDefMidPoint(B,D) \tkzGetPoint{K}
\tkzDefMidPoint(A,C) \tkzGetPoint{L}
\tkzDrawSegments[l3](D,C C,B B,A A,D)
\tkzDrawSegments[l3,maincolor](A,C B,D K,L)
\tkzDrawSegments[dashed](A,E B,E)
\tkzMarkRightAngle[size=0.3](D,E,C)
\tkzMarkAngle[size=1.2,mark=none](E,C,D)
\tkzLabelAngle[pos=1.7,fill=white,inner sep=0.3mm](E,C,D){\small $30^\circ$}
\tkzDrawPoints(A,B,C,D,E,K,L)
\tkzLabelPoint[left](A){A}
\tkzLabelPoint[right](B){B}
\tkzLabelPoint[below right](C){$\varGamma $}
\tkzLabelPoint[below left](D){$\varDelta  $}
\tkzLabelPoint[above](E){E}
\tkzLabelPoint[below](K){K}
\tkzLabelPoint[below](L){$\Lambda$}
\end{tikzpicture}
\end{wrapfigure}
Δίνεται τραπέζιο $AB\varGamma \varDelta (AB \parallel \varGamma \varDelta  $) με τη γωνία $\hat{\varGamma }$ ίση με $30^\circ$ και έστω $ K, \Lambda$ τα μέσα των διαγωνίων του. Οι μη παράλληλες πλευρές του $\varDelta   A$ και $\varGamma  B$ προεκτεινόμενες τέμνονται κάθετα στο σημείο $ E$.
Να αποδείξετε ότι:
\begin{alist}
    \item $AB = 2A E$ \monades{10}
    \item $ K\Lambda = A\varDelta  $ \monades{10}
    \item Σε ποια περίπτωση το $AB\Lambda K$ είναι παραλληλόγραμμο; Να αιτιολογήσετε την απάντησή σας. 

\monades{5}
\end{alist}

\vspace{4mm}

% --- Άσκηση 1737 (1737.tex) ---
\Askhsh{1737}{4}
\begin{ylh}{1}
    \item 3.7. Κύκλος - Μεσοκάθετος - Διχοτόμος
    \item 5.3. Ορθογώνιο
    \item 5.5. Τετράγωνο
    \item 5.9. Μια ιδιότητα του ορθογώνιου τριγώνου.
\end{ylh}
% ===================================================================
%  Αρχείο: 1737.tex (Fragment)
%  Αυτόματη μετατροπή μέσω doc2latex.py
% ===================================================================



Θεωρούμε ορθογώνιο τρίγωνο $AB\varGamma $ ($\hat{A}$ = 90\textsuperscript{ο}) και το ύψος του $AH$. Ονομάζουμε $\varDelta  $ και $E$ τα συμμετρικά σημεία του $H$ ως προς τις ευθείες $AB$ και $A\varGamma $ αντίστοιχα. Αν $M$ είναι το σημείο τομής του τμήματος $H\varDelta  $ με την πλευρά $AB$ και $N$ είναι το σημείο τομής του $HE$ με την πλευρά $A\varGamma $, να αποδείξετε ότι:

\begin{alist}
\item $AH$=$A\varDelta  $=$AE$. \monades{10}

\item Η γωνία $EH\varDelta  $ είναι ορθή. \monades{8}

\item Τα σημεία $A$, Ε και $\varDelta  $ είναι συνευθειακά και $MN$ = $\frac{\varDelta   E}{2}$ . \monades{7}
\end{alist}

\vspace{4mm}

% --- Άσκηση 1738 (1738.tex) ---
\Askhsh{1738}{4}
\begin{ylh}{1}
    \item 3.11. Ανισοτικές σχέσεις πλευρών και γωνιών
    \item 4.2. Τέμνουσα δύο ευθειών - Ευκλείδειο αίτημα
    \item 5.9. Μια ιδιότητα του ορθογώνιου τριγώνου.
\end{ylh}
\begin{wrapfigure}[5]{r}{7cm}
\begin{tikzpicture}[scale=0.75]
\tkzDefPoint(0,0){B}
\tkzDefPoint(8,0){C}
\tkzDefTriangle[two angles = 60 and 30](B,C) \tkzGetPoint{A}
\tkzDefLine[bisector](C,B,A) \tkzGetPoint{b_bis}
\tkzInterLL(B,b_bis)(A,C) \tkzGetPoint{D}
\tkzDefMidPoint(A,C) \tkzGetPoint{M}
\tkzDefLine[parallel=through M](B,D) \tkzGetPoint{m_par}
\tkzInterLL(M,m_par)(B,C) \tkzGetPoint{N}
\tkzDrawPolygon[l3](A,B,C)
\tkzDrawSegments[l4,maincolor](B,D A,N)
\tkzDrawSegments[dashed,l3](M,N)
\tkzDrawPoints(A,B,C,D,M,N)
\tkzLabelPoint[above=2pt](A){A}
\tkzLabelPoint[below left](B){B}
\tkzLabelPoint[below right](C){$\varGamma $}
\tkzLabelPoint[above right](D){$\varDelta  $}
\tkzLabelPoint[above right](M){M}
\tkzLabelPoint[below](N){N}
\tkzMarkAngle[size=0.6,mark=none](C,B,D)
\tkzMarkAngle[size=0.7,mark=none](C,B,D)
\tkzMarkAngle[size=0.6,mark=none](D,B,A)
\tkzMarkAngle[size=0.7,mark=none](D,B,A)
\end{tikzpicture}
\end{wrapfigure}
Δίνεται τρίγωνο $AB\varGamma $ με $\hat{B} = 2\hat{\varGamma }$, και η διχοτόμος $B\varDelta  $ της γωνίας $\hat{B}$. Από το μέσο $ M$ της $A\varGamma $ φέρνουμε παράλληλη στη διχοτόμο $B\varDelta  $ που τέμνει την πλευρά $B\varGamma $ στο $ N$.
Να αποδείξετε ότι:
\begin{alist}
    \item Το τρίγωνο $B\varDelta  \varGamma $ είναι ισοσκελές. \monades{5}
    \item Το τρίγωνο $ M N\varGamma $ είναι ισοσκελές. \monades{10}
    \item $A N \perp B\varGamma $ \monades{10}
\end{alist}

\vspace{4mm}

% --- Άσκηση 1740 (1740.tex) ---
\Askhsh{1740}{4}
\begin{ylh}{1}
    \item 3.6. Κριτήρια ισότητας ορθογώνιων τριγώνων
    \item 5.4. Ρόμβος.
\end{ylh}
Δίνονται οι ακόλουθες προτάσεις $\Pi_1$ και $\Pi_2$:

$\Pi_1$: Αν ένα παραλληλόγραμμο είναι ρόμβος, τότε οι αποστάσεις των απέναντι πλευρών του είναι ίσες.\\[0.2cm]
$\Pi_2$: Αν οι αποστάσεις των απέναντι πλευρών ενός παραλληλογράμμου είναι ίσες, τότε το παραλληλόγραμμο είναι ρόμβος.

\begin{alist}
    \item Να εξετάσετε αν ισχύουν οι προτάσεις $\Pi_1$ και $\Pi_2$ αιτιολογώντας πλήρως την απάντησή σας. \monades{20}
    \item Στην περίπτωση που και οι δύο προτάσεις ισχύουν, να τις διατυπώσετε ως μια ενιαία πρόταση. \monades{5}
\end{alist}

\vspace{4mm}

% --- Άσκηση 1741 (1741.tex) ---
\Askhsh{1741}{4}
\begin{ylh}{1}
    \item 5.2. Παραλληλόγραμμα
    \item 5.6. Εφαρμογές στα τρίγωνα.
\end{ylh}
Δίνεται τρίγωνο $AB\varGamma $ και έστω $ K, \Lambda$ τα μέσα των πλευρών του $ AB$ και $A\varGamma $ αντίστοιχα.
\begin{alist}
    \item Θεωρούμε τυχαίο σημείο $ M$ στο εσωτερικό του τριγώνου και $\varDelta ,E$ τα συμμετρικά του $ M$ ως προς $ K$ και $\Lambda$ αντίστοιχα. Να αποδείξετε ότι $\varDelta  E \parallel B\varGamma $. \monades{15}
    \item Στην περίπτωση που το σημείο $ M$ είναι το μέσο της πλευράς $B\varGamma $, και $\varDelta,E$ τα συμμετρικά του $ M$ ως προς $ K$ και $\Lambda$ αντίστοιχα. Να αποδείξετε ότι τα σημεία $\varDelta  , A$ και $ E$ είναι συνευθειακά. \monades{10}
\end{alist}

\vspace{4mm}

% --- Άσκηση 1742 (1742.tex) ---
\Askhsh{1742}{4}
\begin{ylh}{1}
    \item 3.4. 3ο Κριτήριο ισότητας τριγώνων
    \item 3.6. Κριτήρια ισότητας ορθογώνιων τριγώνων
    \item 3.7. Κύκλος - Μεσοκάθετος - Διχοτόμος
    \item 4.2. Τέμνουσα δύο ευθειών - Ευκλείδειο αίτημα
    \item 5.4. Ρόμβος
    \item 5.6. Εφαρμογές στα τρίγωνα
    \item 5.11. Ισοσκελές τραπέζιο.
\end{ylh}
Το τετράπλευρο $AB\varGamma \varDelta  $ του παρακάτω σχήματος είναι ρόμβος. Θεωρούμε $ A Z \perp \varGamma \varDelta  $ και $A E \perp \varGamma  B$.
Να αποδείξετε ότι:
\begin{alist}
    \item Το τρίγωνο $ZAE$ είναι ισοσκελές. \monades{6}
    \item Η ευθεία $A\varGamma $ είναι μεσοκάθετος του τμήματος $ Z E$. \monades{9}
    \item Αν $ M$ και $ N$ τα μέσα των πλευρών $A\varDelta  $ και $AB$ αντίστοιχα, να αποδείξετε ότι $ M N \parallel  Z E$ και $ Z M = E N$. \monades{10}

\begin{center}
\begin{tikzpicture}[scale=0.9]
% Rhombus $AB\varGamma \varDelta  $: A bottom, B right, Γ top, Δ left
\tkzDefPoint(0,-2.5){A}
\tkzDefPoint(3,0){B}
\tkzDefPoint(0,2.5){Gamma}
\tkzDefPoint(-3,0){Delta}

% Z = projection of A onto $\varGamma \varDelta  $
\tkzDefPointBy[projection=onto Gamma--Delta](A) \tkzGetPoint{Z}

% E = projection of A onto $\varGamma B$
\tkzDefPointBy[projection=onto Gamma--B](A) \tkzGetPoint{E}

% Draw rhombus
\draw[l3] (A)--(B)--(Gamma)--(Delta)--cycle;

% Draw diagonals
\draw[l3, dashed] (A)--(Gamma);

% Draw perpendiculars
\draw[l3] (A)--(Z);
\draw[l3] (A)--(E);

% Draw segment $ZE$
\draw[l3] (Z)--(E);

% Points and labels
\tkzDrawPoints(A,B,Gamma,Delta,Z,E)
\tkzLabelPoint[below](A){$A$}
\tkzLabelPoint[right](B){$B$}
\tkzLabelPoint[above](Gamma){$\varGamma $}
\tkzLabelPoint[left](Delta){$\varDelta  $}
\tkzLabelPoint[left](Z){$ Z$}
\tkzLabelPoint[right](E){$ E$}

% Right angle marks
\tkzMarkRightAngle[size=0.25](A,Z,Gamma)
\tkzMarkRightAngle[size=0.25](A,E,Gamma)
\end{tikzpicture}
\end{center}
\end{alist}

\vspace{4mm}

% --- Άσκηση 1743 (1743.tex) ---
\Askhsh{1743}{4}
\begin{ylh}{1}
    \item 5.2. Παραλληλόγραμμα
    \item 5.3. Ορθογώνιο
    \item 5.4. Ρόμβος
    \item 5.6. Εφαρμογές στα τρίγωνα.
\end{ylh}
Δίνεται ρόμβος $AB\varGamma \varDelta  $ με $\hat{\varGamma } = 120^\circ$. Έστω ότι $ A E$ και $A H$ είναι οι αποστάσεις του σημείου $ A$ στις πλευρές $\varGamma \varDelta  $ και $\varGamma  B$ αντίστοιχα.
\begin{alist}
    \item Να αποδείξετε ότι:
    \begin{rlist}
        \item Τα σημεία $ E$ και $ H$ είναι τα μέσα των πλευρών $\varGamma \varDelta  $ και $\varGamma  B$ αντίστοιχα. \monades{8}
        \item $A\varGamma  \perp  E H$. \monades{8}
    \end{rlist}
    \item Αν $ M$ και $\mathrm{N}$ τα μέσα των πλευρών $A\Delta$ και $AB$ αντίστοιχα, να αποδείξετε ότι το τετράπλευρο $E M\mathrm{N}\mathrm{H}$ είναι ορθογώνιο παραλληλόγραμμο. \monades{9}
\end{alist}

\vspace{4mm}

% --- Άσκηση 1744 (1744.tex) ---
\Askhsh{1744}{4}
\begin{ylh}{1}
    \item 3.2. 1ο Κριτήριο ισότητας τριγώνων
    \item 3.11. Ανισοτικές σχέσεις πλευρών και γωνιών
    \item 4.2. Τέμνουσα δύο ευθειών - Ευκλείδειο αίτημα.
\end{ylh}
Στο ισοσκελές τρίγωνο $AB\varGamma $ ($AB = A\varGamma $) φέρνουμε τις διαμέσους $B\varDelta  $ και $\varGamma  E$. Μία ευθεία $\varepsilon$ παράλληλη στη βάση $B\varGamma $ τέμνει τις πλευρές $AB$ και $ A\varGamma $ στα $ Z$ και $ H$ αντίστοιχα και τις διαμέσους $B\varDelta  $ και $\varGamma  E$ στα σημεία $\varTheta $ και $ K$ αντίστοιχα.
Να αποδείξετε ότι:
\begin{alist}
    \item $B Z = \varGamma  H$. \monades{8}
    \item τα τρίγωνα $ $ZB$\varTheta $ και $ H K\varGamma $ είναι ίσα. \monades{9}
    \item $ Z K =  H\varTheta $. \monades{8}
\end{alist}

\vspace{4mm}

% --- Άσκηση 1745 (1745.tex) ---
\Askhsh{1745}{4}
\begin{ylh}{1}
    \item 3.2. 1ο Κριτήριο ισότητας τριγώνων
    \item 3.4. 3ο Κριτήριο ισότητας τριγώνων
    \item 3.11. Ανισοτικές σχέσεις πλευρών και γωνιών
    \item 5.2. Παραλληλόγραμμα
    \item 5.3. Ορθογώνιο
    \item 5.6. Εφαρμογές στα τρίγωνα.
\end{ylh}
Δίνεται κυρτό τετράπλευρο $AB\varGamma \Delta$ με $BA = B\varGamma $ και $\hat{A} = \hat{\varGamma }$.
Να αποδείξετε ότι:
\begin{alist}
    \item Το τρίγωνο $A\varDelta  \varGamma $ είναι ισοσκελές. \monades{9}
    \item Οι διαγώνιοι του τετραπλεύρου $AB\varGamma \varDelta  $ τέμνονται κάθετα. \monades{6}
    \item Το τετράπλευρο που έχει για κορυφές τα μέσα των πλευρών του $AB\varGamma \varDelta  $ είναι ορθογώνιο. \monades{10}
\end{alist}

\vspace{4mm}

% --- Άσκηση 1746 (1746.tex) ---
\Askhsh{1746}{4}
\begin{ylh}{1}
    \item 4.2. Τέμνουσα δύο ευθειών - Ευκλείδειο αίτημα
    \item 4.6. Άθροισμα γωνιών τριγώνου
    \item 4.8. Άθροισμα γωνιών κυρτού ν-γώνου
    \item 5.2. Παραλληλόγραμμα.
\end{ylh}
Στο κυρτό εξάγωνο $AB\varGamma \varDelta   E Z$ ισχύουν τα εξής: $\hat{\alpha} = \hat{\beta}, \hat{\varGamma } = \hat{\delta}$ και $\hat{\varepsilon} = \hat{\zeta}$.
\begin{alist}
    \item Να υπολογίσετε το άθροισμα $\hat{\alpha} + \hat{\varGamma } + \hat{\varepsilon}$. \monades{8}
    \item Αν οι πλευρές $A Z$ και $\varDelta   E$ προεκτεινόμενες τέμνονται στο $ H$ και οι πλευρές $ AB$ και $\varDelta  \varGamma $ προεκτεινόμενες τέμνονται στο $\varTheta $, να αποδείξετε ότι:
    \begin{rlist}
        \item Οι γωνίες $\hat{A}$ και $\hat{ H}$ είναι παραπληρωματικές. \monades{10}
        \item Το τετράπλευρο $A\varTheta \Delta H$ είναι παραλληλόγραμμο. \monades{7}
    \end{rlist}
\end{alist}

\vspace{4mm}

% --- Άσκηση 1747 (1747.tex) ---
\Askhsh{1747}{4}
\begin{ylh}{1}
    \item 3.14. Σχετικές θέσεις ευθείας και κύκλου
    \item 3.15. Εφαπτόμενα τμήματα
    \item 4.2. Τέμνουσα δύο ευθειών - Ευκλείδειο αίτημα
    \item 5.2. Παραλληλόγραμμα
    \item 5.3. Ορθογώνιο
    \item 5.5. Τετράγωνο
    \item 5.10. Τραπέζιο.
\end{ylh}
Δίνεται κύκλος $( O, R)$ με διάμετρο $AB$ και δυο ευθείες $\varepsilon_1, \varepsilon_2$ εφαπτόμενες του κύκλου στα άκρα της διαμέτρου $ AB$. Έστω ότι, μια τρίτη ευθεία $\varepsilon$ εφάπτεται του κύκλου σ' ένα σημείο του $ E$ και τέμνει τις $\varepsilon_1$ και $\varepsilon_2$ στα $\varDelta  $ και $\varGamma $ αντίστοιχα.
\begin{alist}
    \item Αν το σημείο $ E$ δεν είναι το μέσο του τόξου $AB$, να αποδείξετε ότι:
    \begin{rlist}
        \item Το τετράπλευρο $AB\varGamma \varDelta  $ είναι τραπέζιο. \monades{8}
        \item $\varGamma \varDelta   = A\varDelta   + B\varGamma $. \monades{8}
    \end{rlist}
    \item Αν το σημείο $ E$ βρίσκεται στο μέσον του τόξου $AB$, να αποδείξετε ότι το τετράπλευρο $ A\varDelta  \varGamma  B$ είναι ορθογώνιο. Στην περίπτωση αυτή να εκφράσετε την περίμετρο του ορθογωνίου $ A\varDelta  \varGamma  B$ ως συνάρτηση της ακτίνας $R$ του κύκλου. \monades{9}
\end{alist}

\vspace{4mm}

% --- Άσκηση 1748 (1748.tex) ---
\Askhsh{1748}{4}
\begin{ylh}{1}
    \item 3.3. 2ο Κριτήριο ισότητας τριγώνων
    \item 4.6. Άθροισμα γωνιών τριγώνου
    \item 5.5. Τετράγωνο
    \item 5.8. Το ορθόκεντρο τριγώνου.
\end{ylh}
\begin{wrapfigure}[5]{l}{5cm}
\vspace{-7mm}
\begin{tikzpicture}[scale=1]
% Square $AB\varGamma \varDelta  $: A bottom-left, B bottom-right, Γ top-right, Δ top-left
\tkzDefPoint(0,0){A}
\tkzDefPoint(4,0){B}
\tkzDefPoint(4,4){Gamma}
\tkzDefPoint(0,4){Delta}

% Center O at intersection of diagonals
\tkzDefMidPoint(A,Gamma) \tkzGetPoint{O}

% E is on segment $O\varDelta  $ (arbitrary position, e.g. 40% from O to Δ)
\tkzDefPointOnLine[pos=0.4](O,Delta) \tkzGetPoint{E}

% Z is foot of perpendicular from B to $AE$, on segment $AO$
\tkzDefPointBy[projection=onto A--E](B) \tkzGetPoint{F} % Foot of perp on $AE$
\tkzInterLL(B,F)(A,Gamma) \tkzGetPoint{Z} % Z on diagonal $A\varGamma $

% Draw square
\draw[l3] (A)--(B)--(Gamma)--(Delta)--cycle;

% Draw diagonals
\draw[l3] (A)--(Gamma);
\draw[l3] (B)--(Delta);

% Draw $AE$
\draw[l3,maincolor] (A)--(E);

% Draw $BZ$ (perpendicular to $AE$)
\draw[l3,maincolor] (B)--(F);
\draw[l3] (Z)--(F);

% Draw $EZ$
\draw[l3, dashed] (E)--(Z);

% Draw $\varGamma Z$
\draw[l3, dashed] (Gamma)--(Z);

% Points and labels
\tkzDrawPoints(A,B,Gamma,Delta,O,E,Z)
\tkzLabelPoint[below left](A){$A$}
\tkzLabelPoint[below right](B){$B$}
\tkzLabelPoint[above right](Gamma){$\varGamma $}
\tkzLabelPoint[above left](Delta){$\varDelta  $}
\tkzLabelPoint[right](O){$ O$}
\tkzLabelPoint[above right](E){$ E$}
\tkzLabelPoint[below](Z){$ Z$}

% Right angle at intersection of $BZ$ and $AE$
\tkzMarkRightAngle[size=0.2](B,F,A)

% Angle ω at vertex A (between $A\varDelta  $ and $AE$)
\tkzMarkAngle[size=0.7](O,A,E)
\tkzLabelAngle[pos=0.9](O,A,E){$\omega$}

% Angle φ at vertex B (between $AB$ and $BZ$)
\tkzMarkAngle[size=0.7](O,B,Z)
\tkzLabelAngle[pos=0.9](O,B,Z){$\varphi$}
\end{tikzpicture}
\end{wrapfigure}
Στο τετράγωνο $AB\varGamma \varDelta  $ ονομάζουμε $ O$ το κέντρο του και θεωρούμε τυχαίο σημείο $ E$ του τμήματος $\mathrm{O}\varDelta  $. Φέρνουμε την κάθετη από το $B$ στην $AE$, που τέμνει το τμήμα $ A\mathrm{O}$ στο $ Z$.
Να αποδείξετε ότι:
\begin{alist}
    \item Οι γωνίες $\omega$ και $\varphi$ του παρακάτω σχήματος είναι ίσες. \monades{6}
    \item $B Z = A E$ και $\varGamma  Z = B E$. \monades{12}
    \item Το τμήμα $ E Z$ είναι κάθετο στο $AB$. \monades{7}
\end{alist}
\vspace{4mm}

% --- Άσκηση 1749 (1749.tex) ---
\Askhsh{1749}{4}
\begin{ylh}{1}
    \item 3.7. Κύκλος - Μεσοκάθετος - Διχοτόμος
    \item 3.10. Σχέση εξωτερικής και απέναντι γωνίας
    \item 3.12. Τριγωνική ανισότητα.
\end{ylh}
% ===================================================================
%  Αρχείο: 1749.tex (Fragment)
%  Αυτόματη μετατροπή μέσω doc2latex.py
% ===================================================================
Θεωρούμε δυο σημεία $A$ και Β τα οποία βρίσκονται στο ίδιο μέρος ως προς μια ευθεία (ε), τέτοια ώστε η ευθεία $AB$ δεν είναι κάθετη στην (ε). Έστω Α΄ το συμμετρικό του Α ως προς την ευθεία (ε), δηλαδή η (ε) είναι μεσοκάθετος του $AA$΄. $\ $

\begin{alist}
\item Αν η Α΄Β τέμνει την ευθεία (ε) στο σημείο $O$, να αποδείξετε ότι:
\end{alist}
\begin{alist}
\item
  Η ευθεία (ε) διχοτομεί τη γωνία $\hat{AOA΄}$. \monades{6}
\item
  Οι ημιευθείες $OA$ και $OB$ σχηματίζουν ίσες οξείες γωνίες με την ευθεία (ε).

\monades{6}

\item Αν Κ είναι ένα άλλο σημείο πάνω στην ευθεία (ε), να αποδείξετε ότι:

\end{alist}
\begin{alist}
\item
  $KA$=$KA$΄ \monades{6}
\item
  $KA$+$KB$\textgreater $AO$+$OB$ \monades{7}
\end{alist}

\vspace{4mm}

% --- Άσκηση 1750 (1750.tex) ---
\Askhsh{1750}{4}
\begin{ylh}{1}
    \item 3.6. Κριτήρια ισότητας ορθογώνιων τριγώνων
    \item 5.5. Τετράγωνο.
\end{ylh}
Στο τετράγωνο $AB\varGamma \varDelta  $ προεκτείνουμε την πλευρά $ AB$ κατά τμήμα $B N$ και την πλευρά $B\varGamma $ κατά τμήμα $\varGamma  M = A\mathrm{N}$.
Να αποδείξετε ότι:
\begin{alist}
    \item $\varDelta   N = \varDelta   M$. \monades{12}
    \item $\varDelta   N \perp \varDelta  \mathrm{M}$. \monades{13}
\end{alist}

\vspace{4mm}

% --- Άσκηση 1751 (1751.tex) ---
\Askhsh{1751}{4}
\begin{ylh}{1}
    \item 3.4. 3ο Κριτήριο ισότητας τριγώνων
    \item 3.6. Κριτήρια ισότητας ορθογώνιων τριγώνων
    \item 3.15. Εφαπτόμενα τμήματα.
\end{ylh}
Έστω ότι ο κύκλος $( O, \rho)$ εφάπτεται των πλευρών του τριγώνου $\mathrm{P}\varGamma  E$ στα σημεία $A, \varDelta  $ και $B$.
\begin{alist}
    \item Να αποδείξετε ότι:
    \begin{rlist}
        \item $\mathrm{P}\varGamma  = \varGamma \varDelta   + A\mathrm{P}$ \monades{6}
        \item $\mathrm{P}\varGamma  - \varGamma \varDelta   = \mathrm{P} E - \varDelta   E$ \monades{8}
    \end{rlist}
    \item Αν $A\varGamma  = B E$, να αποδείξετε ότι
    \begin{rlist}
        \item Το τρίγωνο $\mathrm{P}\varGamma  E$ είναι ισοσκελές. \monades{6}
        \item Τα σημεία $\mathrm{P},  O$ και $\varDelta  $ είναι συνευθειακά. \monades{5}
    \end{rlist}
\end{alist}

\vspace{4mm}

% --- Άσκηση 1752 (1752.tex) ---
\Askhsh{1752}{4}
\begin{ylh}{1}
    \item 3.15. Εφαπτόμενα τμήματα.
\end{ylh}
\begin{wrapfigure}[5]{r}{4.5cm}
\vspace{-7mm}
\begin{tikzpicture}[scale=0.6]

% 1. Ορισμός των κορυφών του τριγώνου 
% (Επιλέγουμε συντεταγμένες που δημιουργούν ένα ισοσκελές τρίγωνο)
\tkzDefPoint(0,0){G} % Σημείο \varGamma 
\tkzDefPoint(6,0){E} % Σημείο E
\tkzDefPoint(3,6.5){P} % Σημείο P

% 2. Εύρεση του έγκεντρου Ο του τριγώνου P\varGamma  E
\tkzDefTriangleCenter[in](P,G,E) \tkzGetPoint{O}

% 3. Εύρεση των σημείων επαφής A, B, \Delta ως προβολές του Ο στις πλευρές
\tkzDefPointBy[projection=onto P--G](O) \tkzGetPoint{A}
\tkzDefPointBy[projection=onto P--E](O) \tkzGetPoint{B}
\tkzDefPointBy[projection=onto G--E](O) \tkzGetPoint{D} % Σημείο \Delta

% 4. Σχεδίαση του τριγώνου P\varGamma  E
\tkzDrawPolygon[l3](P,G,E)

% 5. Σχεδίαση του εγγεγραμμένου κύκλου (Ο, ρ)
\tkzDrawCircle[l4,maincolor](O,D)

% 6. Εμφάνιση των κουκκίδων στα σημεία
\tkzDrawPoints(P,G,E,A,B,D,O)

% 7. Τοποθέτηση των ετικετών (γραμμάτων) στις σωστές θέσεις
\tkzLabelPoint[above=2pt](P){P}
\tkzLabelPoint[below left=2pt](G){$\varGamma $}
\tkzLabelPoint[below right=2pt](E){E}
\tkzLabelPoint[left=3pt](A){A}
\tkzLabelPoint[right=3pt](B){B}
\tkzLabelPoint[below=3pt](D){$\varDelta  $}
\tkzLabelPoint[above=4pt](O){O}

\end{tikzpicture}
\end{wrapfigure}
Θεωρούμε κύκλο κέντρου $ O$ και εξωτερικό σημείο του $\mathrm{P}$. Από το $\mathrm{P}$ φέρνουμε τα εφαπτόμενα τμήματα $\mathrm{P}A$ και $\mathrm{P}B$. Η διακεντρική ευθεία $\mathrm{P} O$ τέμνει τον κύκλο στο σημείο $\Lambda$. Η εφαπτόμενη του κύκλου στο $\Lambda$ τέμνει τα $\mathrm{P}A$ και $\mathrm{P}B$ στα σημεία $\varGamma $ και $\varDelta  $ αντίστοιχα.
Να αποδείξετε ότι:
\begin{alist}
    \item το τρίγωνο $\mathrm{P}\varGamma \varDelta  $ είναι ισοσκελές. \monades{10}
    \item $\varGamma  A = \varDelta   B$. \monades{8}
    \item η περίμετρος του τριγώνου $\mathrm{P}\varGamma \varDelta  $ είναι ίση με $\mathrm{P}A + \mathrm{P}B$. \monades{7}
\end{alist}

\vspace{4mm}

% --- Άσκηση 1754 (1754.tex) ---
\Askhsh{1754}{4}
\begin{ylh}{1}
    \item 3.4. 3ο Κριτήριο ισότητας τριγώνων
    \item 3.6. Κριτήρια ισότητας ορθογώνιων τριγώνων
    \item 5.8. Το ορθόκεντρο τριγώνου.
\end{ylh}
\begin{wrapfigure}[5]{r}{4.5cm}
\begin{tikzpicture}[scale=0.5]
\tkzDefPoint(-4,0){B}
\tkzDefPoint(4,0){C} % Σημείο \varGamma 
\tkzDefPoint(0,0){D} % Σημείο \Delta
\tkzDefPoint(0,6){A}
\tkzDefLine[mediator](A,B) \tkzGetPoints{med1}{med2}
\tkzInterLL(A,D)(med1,med2) \tkzGetPoint{H}
\tkzInterLL(B,H)(A,C) \tkzGetPoint{E}
\tkzDefPointBy[projection=onto B--E](A) \tkzGetPoint{Z}
\tkzInterLL(A,Z)(B,C) \tkzGetPoint{Th}
\tkzDrawSegments[l3](A,B B,C A,C)
\tkzDrawSegments[l4,maincolor](A,D B,E A,Th)
\tkzMarkSegments[mark=||, size=4pt](A,H B,H)
\tkzMarkRightAngle[fill=maincolor!20, size=0.35](A,D,B)
\tkzMarkRightAngle[fill=maincolor!20, size=0.35](A,Z,H)
\tkzDrawPoints(A,B,C,D,E,Z,H,Th)
\tkzLabelPoint[above](A){A}
\tkzLabelPoint[left](B){B}
\tkzLabelPoint[right](C){$\varGamma $}
\tkzLabelPoint[below](D){$\varDelta  $}
\tkzLabelPoint[above right](E){E}
\tkzLabelPoint[below](Th){$\varTheta $}
\tkzLabelPoint[left](H){H}
\tkzLabelPoint[right=2pt](Z){Z}
\end{tikzpicture}
\end{wrapfigure}
Δίνεται ισοσκελές τρίγωνο $AB\varGamma $ και το ύψος του $ A\Delta$. Στο $A\Delta$ θεωρούμε σημείο $\Theta$ τέτοιο ώστε $\Theta A = \Theta B$. Έστω ότι $E$ είναι το σημείο τομής της $B\Theta$ με την $A\varGamma $. Φέρνουμε την $A\mathrm{H}$ κάθετη στην $BE$, η οποία τέμνει την πλευρά $B\varGamma $ στο $\mathrm{Z}$.
\begin{alist}
    \item Να αποδείξετε ότι:
    \begin{rlist}
        \item Τα τρίγωνα $\varTheta\varDelta   B$ και $\varTheta  HA$ είναι ίσα. \monades{6}
        \item $\varDelta  \varTheta  = \varTheta  H$. \monades{6}
        \item Η ευθεία $\varDelta  \varTheta $ είναι μεσοκάθετος του τμήματος $AB$. \monades{6}
    \end{rlist}
    \item Ποιο από τα σημεία του σχήματος είναι το ορθόκεντρο του τριγώνου $A\varTheta  B$; Να δικαιολογήσετε την απάντησή σας. \monades{7}
\end{alist}

\vspace{4mm}

% --- Άσκηση 1755 (1755.tex) ---
\Askhsh{1755}{4}
\begin{ylh}{1}
    \item 4.2. Τέμνουσα δύο ευθειών - Ευκλείδειο αίτημα
    \item 4.6. Άθροισμα γωνιών τριγώνου
    \item 5.4. Ρόμβος
    \item 5.11. Ισοσκελές τραπέζιο.
\end{ylh}
Σε ισοσκελές τραπέζιο $AB\varGamma \varDelta  $ ($AB \parallel \varGamma \varDelta  $) είναι $AB =  A\varDelta  $.
\begin{alist}
    \item Να αποδείξετε ότι η $B\varDelta  $ είναι διχοτόμος της γωνίας $\hat{\varDelta  }$. \monades{7}
    \item Να προσδιορίσετε τη θέση ενός σημείου $ E$, ώστε το τετράπλευρο $ABE\varDelta  $ να είναι ρόμβος. \monades{10}
    \item Αν επιπλέον είναι γωνία $B\hat{A}\varDelta   = 120^\circ$ και οι διαγώνιοι του ρόμβου τέμνονται στο σημείο $ O$, να υπολογίσετε τις γωνίες του τετραπλεύρου $ E\mathrm{O}B\varGamma $. \monades{8}
\end{alist}

\vspace{4mm}

% --- Άσκηση 1757 (1757.tex) ---
\Askhsh{1757}{4}
\begin{ylh}{1}
    \item 5.2. Παραλληλόγραμμα
    \item 5.3. Ορθογώνιο
    \item 5.10. Τραπέζιο.
\end{ylh}
Θεωρούμε τραπέζιο $AB\varGamma \varDelta  $, τέτοιο ώστε $\hat{ A} = \hat{\varDelta  } = 90^\circ$, $AB = \frac{1}{4}\varDelta  \varGamma $ και $AB = \frac{1}{3}A\varDelta  $. Επιπλέον, φέρνουμε $B E \perp \varDelta  \varGamma $.
\begin{alist}
    \item Να αποδείξετε ότι το τετράπλευρο $AB E\varDelta  $ είναι ορθογώνιο. \monades{6}
    \item Να αποδείξετε ότι το τρίγωνο $B E\varGamma $ είναι ορθογώνιο και ισοσκελές. \monades{10}
    \item Αν $ K, \Lambda$ είναι τα μέσα των $B\varDelta  $ και $A\varGamma $ αντίστοιχα, να αποδείξετε ότι η $ A\varGamma $ διέρχεται από το μέσο του ευθυγράμμου τμήματος $B K$. \monades{9}
\end{alist}

\vspace{4mm}

% --- Άσκηση 1758 (1758.tex) ---
\Askhsh{1758}{4}
\begin{ylh}{1}
    \item 3.6. Κριτήρια ισότητας ορθογώνιων τριγώνων
    \item 3.15. Εφαπτόμενα τμήματα
    \item 4.2. Τέμνουσα δύο ευθειών - Ευκλείδειο αίτημα
    \item 5.3. Ορθογώνιο
    \item 5.10. Τραπέζιο.
\end{ylh}
\begin{wrapfigure}[5]{r}{5cm}
\begin{tikzpicture}[scale=0.6]

% 1. Ορισμός κέντρου και άκρων διαμέτρου
\tkzDefPoint(0,0){O}
\tkzDefPoint(-3,0){A}
\tkzDefPoint(3,0){B}

% 2. Σχεδίαση κύκλου (O, R) και διαμέτρου $AB$
\tkzDrawCircle[l4,maincolor](O,A)
\tkzDrawSegment(A,B)

% 3. Ορισμός του σημείου E στο ημικύκλιο (επιλογή 55 μοιρών για οπτική ταύτιση)
\tkzDefPoint(65:3){E}

% 4. Εύρεση της εφαπτομένης στο Ε
\tkzDefLine[tangent at=E](O) \tkzGetPoint{e_tan}

% 5. Βοηθητικά σημεία για τις κατακόρυφες εφαπτόμενες ε1 και ε2
\tkzDefPoint(-3,1){A_up}
\tkzDefPoint(3,1){B_up}

% 6. Εύρεση των σημείων τομής \Delta και \varGamma 
\tkzInterLL(A,A_up)(E,e_tan) \tkzGetPoint{D} % Σημείο \Delta
\tkzInterLL(B,B_up)(E,e_tan) \tkzGetPoint{C} % Σημείο \varGamma 

% 7. Καθορισμός των άκρων για τις ευθείες ώστε να προεξέχουν σωστά
\tkzDefPoint(-3,-2){eps1_bot}
\tkzDefPoint(-3,5.5){eps1_top}
\tkzDefPoint(3,-2){eps2_bot}
\tkzDefPoint(3,4){eps2_top}

% Επέκταση της ευθείας ε
\tkzDefPointBy[homothety=center C ratio -0.2](D) \tkzGetPoint{eps_top}
\tkzDefPointBy[homothety=center D ratio -0.1](C) \tkzGetPoint{eps_bot}

% 8. Σχεδίαση των τριών ευθειών (ε1, ε2, ε)
\tkzDrawSegment(eps1_bot,eps1_top)
\tkzDrawSegment(eps2_bot,eps2_top)
\tkzDrawSegment(eps_top,eps_bot)

% 9. Σχεδίαση των κουκκίδων στα σημεία ενδιαφέροντος
\tkzDrawPoints[size=3](O,A,B,E,C,D)

% 10. Τοποθέτηση των ετικετών (Γραμμάτων)
\tkzLabelPoint[left](A){A}
\tkzLabelPoint[right](B){B}
\tkzLabelPoint[below=3pt](O){O}
\tkzLabelPoint[above right](E){E}
\tkzLabelPoint[above right](C){$\varGamma $}
\tkzLabelPoint[above right](D){$\varDelta  $}

% 11. Τοποθέτηση των ονομάτων των ευθειών
\tkzLabelPoint[right](eps1_bot){$\varepsilon_1$}
\tkzLabelPoint[right](eps2_bot){$\varepsilon_2$}
\tkzLabelPoint[above right](eps_bot){$\varepsilon$}

\end{tikzpicture}
\end{wrapfigure}
Δίνεται κύκλος $( O, R)$ με διάμετρο $AB$ και ευθείες $\varepsilon_1, \varepsilon_2$ εφαπτόμενες του κύκλου στα άκρα της διαμέτρου $ AB$. Θεωρούμε ευθεία $\varepsilon$ εφαπτομένη του κύκλου σε σημείο του $ E$, η οποία τέμνει τις $\varepsilon_1$ και $\varepsilon_2$ στα $\varDelta  $ και $\varGamma $ αντίστοιχα.
\begin{alist}
    \item Να αποδείξετε ότι $\varGamma \varDelta   = A\varDelta   + B\varGamma $. \monades{9}
    \item Να αποδείξετε ότι το τρίγωνο $\varGamma  O\varDelta  $ είναι ορθογώνιο. \monades{9}
    \item Να διερευνήσετε το είδος του τετραπλεύρου $AB\varGamma \varDelta  $ ανάλο$\varGamma  A$ με τη θέση του σημείου $ E$ στο ημικύκλιο $ AB$. \monades{7}
\end{alist}

\vspace{4mm}

% --- Άσκηση 1759 (1759.tex) ---
\Askhsh{1759}{4}
\begin{ylh}{1}
    \item 3.2. 1ο Κριτήριο ισότητας τριγώνων
    \item 4.2. Τέμνουσα δύο ευθειών - Ευκλείδειο αίτημα
    \item 5.2. Παραλληλόγραμμα
    \item 5.4. Ρόμβος
    \item 5.9. Μια ιδιότητα του ορθογώνιου τριγώνου.
\end{ylh}
Σε παραλληλόγραμμο $AB\varGamma \varDelta  $ με γωνία $\hat{ A}$ αμβλεία, ισχύει ότι $ $AB$ = 2A\varDelta  $. Τα σημεία $ E$ και $ Z$, είναι μέσα των πλευρών του $AB$ και $\varGamma \varDelta  $ αντίστοιχα. Από το $\varDelta  $ φέρνουμε τη $\varDelta   H$ κάθετη στην προέκταση της $B\varGamma $.
Να αποδείξετε ότι:
\begin{alist}
    \item Το τετράπλευρο $A E Z\varDelta  $ είναι ρόμβος. \monades{8}
    \item Το τρίγωνο $ E Z H$ είναι ισοσκελές. \monades{9}
    \item Το τμήμα $ H E$, είναι διχοτόμος της γωνίας $ Z\hat{\mathrm{H}}\varGamma $. \monades{8}
\end{alist}

\vspace{4mm}

% --- Άσκηση 1760 (1760.tex) ---
\Askhsh{1760}{4}
\begin{ylh}{1}
    \item 3.6. Κριτήρια ισότητας ορθογώνιων τριγώνων
    \item 5.4. Ρόμβος
    \item 5.7. Βαρύκεντρο τριγώνου.
\end{ylh}
\begin{wrapfigure}[5]{r}{5cm}
\begin{tikzpicture}[scale=0.65]
\tkzDefPoint(-1.5,0){B}
\tkzDefPoint(1.5,0){C} % Σημείο \varGamma 
\tkzDefPoint(0,0){M}
\tkzDefPoint(0,2.5){A}
\tkzDefPointBy[symmetry=center M](A) \tkzGetPoint{N}
\tkzDefPointBy[translation=from B to C](C) \tkzGetPoint{D}
\tkzFillPolygon[color=maincolor!10](A,B,C)
\tkzDrawSegments[l3](A,B B,D A,C A,D C,N D,N A,N)
\tkzDrawSegment[dotted, thick](B,N)
\tkzDrawPoints(A,B,C,D,M,N)
\tkzLabelPoint[above](A){A}
\tkzLabelPoint[left](B){B}
\tkzLabelPoint[below left](M){M}
\tkzLabelPoint[below right](C){$\varGamma $}
\tkzLabelPoint[right](D){$\varDelta  $}
\tkzLabelPoint[below](N){N}
\end{tikzpicture}
\end{wrapfigure}
Δίνεται ισοσκελές τρίγωνο $AB\varGamma (AB = A\varGamma )$ και $A\mathrm{M}$ το ύψος του στην πλευρά $B\varGamma $. Στην προέκταση του $A\mathrm{M}$ θεωρούμε τμήμα $\mathrm{M}\mathrm{N} = A\mathrm{M}$. Στην προέκταση του $B\varGamma $ προς το μέρος του $\varGamma $ θεωρούμε τμήμα $\varGamma \Delta = B\varGamma $.
Να αποδείξετε ότι:
\begin{alist}
    \item Το τετράπλευρο $AB\mathrm{N}\varGamma $ είναι ρόμβος. \monades{8}
    \item Το τρίγωνο $A\varDelta  \mathrm{N}$ είναι ισοσκελές. \monades{8}
    \item Το σημείο $\varGamma $ είναι το βαρύκεντρο του τριγώνου $A\varDelta  \mathrm{N}$. \monades{9}
\end{alist}

\vspace{4mm}

% --- Άσκηση 1761 (1761.tex) ---
\Askhsh{1761}{4}
\begin{ylh}{1}
    \item 3.4. 3ο Κριτήριο ισότητας τριγώνων
    \item 4.6. Άθροισμα γωνιών τριγώνου
    \item 5.9. Μια ιδιότητα του ορθογώνιου τριγώνου.
\end{ylh}
\begin{wrapfigure}[5]{r}{5cm}
\begin{tikzpicture}[scale=0.7]
\tkzDefPoint(0,0){B}
\tkzDefPoint(8,0){G}
\tkzDefTriangle[two angles = 45 and 15](B,G) \tkzGetPoint{A}
\tkzDefPointBy[homothety=center B ratio 3](A) \tkzGetPoint{D}
\tkzDefPointBy[homothety=center B ratio 1.15](D) \tkzGetPoint{D_ext}
\tkzDefPointBy[projection=onto A--G](D) \tkzGetPoint{K}
\tkzDrawSegments[l3](B,D_ext B,G A,G D,K B,K)
\tkzMarkAngle[size=0.7,mark=none](G,B,A)
\tkzLabelAngle[pos=1.1](G,B,A){\small $45^\circ$}
\tkzMarkAngle[size=0.5,mark=none](B,A,G)
\tkzLabelAngle[pos=0.8](B,A,G){\small $120^\circ$}
\tkzMarkRightAngle[fill=maincolor!20, size=0.3](D,K,G)
\tkzDrawPoints(A,B,G,D,K)
\tkzLabelPoint[left=3pt](A){A}
\tkzLabelPoint[below left](B){B}
\tkzLabelPoint[above right](G){$\varGamma $}
\tkzLabelPoint[left=3pt](D){$\varDelta  $}
\tkzLabelPoint[below=3pt](K){K}
\end{tikzpicture}
\end{wrapfigure}
Δίνεται τρίγωνο $AB\varGamma $ με γωνία $\hat{ A}$ ίση με $120^\circ$ και γωνία $\hat{B}$ είναι ίση με $45^\circ$. Στην προέκταση της $BA$ προς το $A$, παίρνουμε τμήμα $A\Delta = 2 AB$. Από το $\Delta$ φέρνουμε την κάθετη στην $ A\varGamma $ που την τέμνει στο σημείο $ K$.
Να αποδείξετε ότι:
\begin{alist}
    \item Η γωνία $A\hat{\varDelta  } K$ είναι ίση με $30^\circ$. \monades{6}
    \item Το τρίγωνο $ KAB$ είναι ισοσκελές. \monades{6}
    \item Αν $ Z$ το μέσο της $\varDelta   A$, τότε $ Z\hat{ K}B = 90^\circ$. \monades{6}
    \item Το σημείο $ K$ ανήκει στη μεσοκάθετο του τμήματος $B\varDelta  $. \monades{7}
\end{alist}

\vspace{4mm}

% --- Άσκηση 1764 (1764.tex) ---
\Askhsh{1764}{4}
\begin{ylh}{1}
    \item 5.2. Παραλληλόγραμμα
    \item 5.8. Το ορθόκεντρο τριγώνου.
\end{ylh}
Δίνεται ορθογώνιο $AB\varGamma \varDelta  $ με κέντρο $ O$ και $A\varGamma  = 2B\varGamma $. Στην προέκταση της πλευράς $\varDelta   A$, προς το $A$, παίρνουμε σημείο $ E$ ώστε $\varDelta   A = AE$.
Να αποδείξετε ότι:
\begin{alist}
    \item Το τετράπλευρο $A $EB$\varGamma $ είναι παραλληλόγραμμο. \monades{8}
    \item Το τρίγωνο $ $EB$\varDelta  $ είναι ισόπλευρο. \monades{9}
    \item Αν η $ E O$ τέμνει την πλευρά $AB$ στο σημείο $ Z$, να αποδείξετε ότι $\Delta Z \perp EB$. \monades{8}
\end{alist}

\vspace{4mm}

% --- Άσκηση 1766 (1766.tex) ---
\Askhsh{1766}{4}
\begin{ylh}{1}
    \item 3.6. Κριτήρια ισότητας ορθογώνιων τριγώνων
    \item 5.6. Εφαρμογές στα τρίγωνα.
\end{ylh}
\begin{wrapfigure}[5]{r}{7cm}
\begin{tikzpicture}[scale=0.7]
\tkzDefPoint(0,4){A}
\tkzDefPoint(4,4){B}
\tkzDefPoint(4,0){C}
\tkzDefPoint(0,0){D}
\tkzDefPointBy[symmetry=center D](B) \tkzGetPoint{E}
\tkzDefMidPoint(A,D) \tkzGetPoint{Z}
\tkzInterLL(A,E)(C,D) \tkzGetPoint{H}
\tkzDrawPolygon[l3](A,B,C,D)
\tkzDrawSegments[l4,maincolor](B,E A,E Z,C)
\tkzDrawSegment[dashed](D,H)
\tkzDrawPoints(A,B,C,D,E,Z,H)
\tkzLabelPoint[above left](A){A}
\tkzLabelPoint[above right](B){B}
\tkzLabelPoint[below right](C){$\varGamma $}
\tkzLabelPoint[below right](D){$\varDelta  $}
\tkzLabelPoint[below](E){E}
\tkzLabelPoint[left](H){H}
\tkzLabelPoint[above right=1pt](Z){Z}
\end{tikzpicture}
\end{wrapfigure}
Δίνεται τετράγωνο $AB\varGamma \varDelta  $. Έστω $ E$ το συμμετρικό σημείο του $B$ ως προς το $\varDelta  $ και $ H$ είναι το μέσο της $ A\varDelta  $. Η προέκταση της $\varGamma \varDelta  $ τέμνει την $ A E$ στο $\varTheta $.
Να αποδείξετε ότι:
\begin{alist}
    \item $\varDelta  \varTheta  = \frac{AB}{2}$. \monades{8}
    \item Τα τρίγωνα $A\varDelta  \varTheta $ και $ H\varDelta  \varGamma $ είναι ίσα. \monades{9}
    \item Η $\varGamma  H$ είναι κάθετη στην $A E$. \monades{8}
\end{alist}

\vspace{4mm}

% --- Άσκηση 1767 (1767.tex) ---
\Askhsh{1767}{4}
\begin{ylh}{1}
    \item 4.2. Τέμνουσα δύο ευθειών - Ευκλείδειο αίτημα
    \item 5.2. Παραλληλόγραμμα
    \item 5.3. Ορθογώνιο
    \item 5.4. Ρόμβος.
\end{ylh}
\begin{wrapfigure}[5]{r}{7cm}
\begin{tikzpicture}[scale=0.7]
    \tkzDefPoint(0,0){D}
    \tkzDefPoint(8,0){C}
    \tkzDefPoint(0,4){A}
    \tkzDefPoint(4,4){B}
    \tkzDefPointBy[projection=onto D--C](B) \tkzGetPoint{E}
    \tkzInterLL(B,E)(A,C) \tkzGetPoint{M}
    \tkzInterLL(A,E)(B,D) \tkzGetPoint{N}
    \tkzDrawSegments[l3](A,B B,C C,D D,A B,E A,C)
    \tkzDrawSegments[l4,maincolor,dashed](A,E B,D)
    \tkzMarkRightAngle[fill=maincolor!15, size=0.4](D,A,B)
    \tkzMarkRightAngle[fill=maincolor!15, size=0.4](C,D,A)
    \tkzMarkRightAngle[fill=maincolor!15, size=0.4](C,E,B)
    \tkzDrawPoints(A,B,C,D,E,M,N)
    \tkzLabelPoint[above left](A){A}
    \tkzLabelPoint[above right](B){B}
    \tkzLabelPoint[right](C){$\varGamma $}
    \tkzLabelPoint[left](D){$\varDelta  $}
    \tkzLabelPoint[below](E){E}
    \tkzLabelPoint[right=2pt](M){M}
    \tkzLabelPoint[below=4pt](N){N}
\end{tikzpicture}
\end{wrapfigure}
Δίνεται τραπέζιο $AB\varGamma \Delta$ ($AB \parallel \varGamma \Delta$) με $\hat{A} = \hat{\Delta} = 90^\circ$, $\Delta\varGamma  = 2AB$ και $\hat{B} = 3\hat{\varGamma }$. Φέρνουμε $BE \perp \Delta\varGamma $ που τέμνει τη διαγώνιο $A\varGamma $ στο $\mathrm{M}$. Φέρνουμε την $AE$ που τέμνει τη διαγώνιο $B\Delta$ στο $\mathrm{N}$.
Να αποδείξετε ότι:
\begin{alist}
    \item $\hat{\varGamma } = 45^\circ$. \monades{7}
    \item Το τετράπλευρο $AB\varGamma  E$ είναι παραλληλόγραμμο. \monades{9}
    \item $A E \perp B\varDelta  $. \monades{9}
\end{alist}

\vspace{4mm}

% --- Άσκηση 1770 (1770.tex) ---
\Askhsh{1770}{4}
\begin{ylh}{1}
    \item 3.6. Κριτήρια ισότητας ορθογώνιων τριγώνων
    \item 5.2. Παραλληλόγραμμα
    \item 5.9. Μια ιδιότητα του ορθογώνιου τριγώνου.
\end{ylh}
\begin{wrapfigure}[5]{r}{7cm}
\begin{tikzpicture}[scale=0.7]
    \tkzDefPoint(0,0){D}
    \tkzDefPoint(7,0){C}
    \tkzDefPoint(1,3.5){A}
    \tkzDefPointBy[translation=from D to C](A) \tkzGetPoint{B}
    \tkzDefMidPoint(A,C) \tkzGetPoint{O}
    \tkzDefPointBy[projection=onto A--C](D) \tkzGetPoint{K}
    \tkzDefPointBy[symmetry=center K](D) \tkzGetPoint{E}
    \tkzDrawPolygon[l3](A,B,C,D)
    \tkzDrawSegments[l3](A,C B,D D,E)
    \tkzDrawSegment[dashed](O,E)
    \tkzMarkRightAngle[fill=maincolor!15, size=0.3](D,K,A)
    \tkzDrawPoints(A,B,C,D,O,K,E)
    \tkzLabelPoint[above left](A){A}
    \tkzLabelPoint[above right](B){B}
    \tkzLabelPoint[below right](C){$\varGamma $}
    \tkzLabelPoint[below left](D){$\varDelta  $}
    \tkzLabelPoint[below](O){O}
    \tkzLabelPoint[above=2pt](K){K}
    \tkzLabelPoint[above](E){E}
\end{tikzpicture}
\end{wrapfigure}
Δίνεται παραλληλόγραμμο $AB\varGamma \varDelta  $ με $ O$ το κέντρο του. Από την κορυφή $\varDelta  $ φέρνουμε το τμήμα $\varDelta   K$ κάθετο στην $ A\varGamma $ και στην προέκτασή του προς το $ K$ θεωρούμε σημείο $ E$, ώστε $\mathrm{K} E = \varDelta   K$.
Να αποδείξετε ότι:
\begin{alist}
    \item $ E O = \frac{B\varDelta  }{2}$. \monades{8}
    \item Η γωνία $\varDelta  \hat{ E}B$ είναι ορθή. \monades{8}
    \item Το τετράπλευρο $A $EB$\varGamma $ είναι ισοσκελές τραπέζιο. \monades{9}
\end{alist}

\vspace{4mm}

% --- Άσκηση 1771 (1771.tex) ---
\Askhsh{1771}{4}
\begin{ylh}{1}
    \item 3.15. Εφαπτόμενα τμήματα
    \item 5.9. Μια ιδιότητα του ορθογώνιου τριγώνου.
\end{ylh}
\begin{wrapfigure}[5]{r}{7cm}
\begin{tikzpicture}[scale=0.7]
    \tkzDefPoint(0,0){O}
    \tkzDefPoint(5,0){K}
    \tkzDefPoint(3,0){N}
    \tkzDrawCircle[l3](O,N)
    \tkzDrawCircle[l3](K,N)
    \pgfmathsetmacro{\ang}{90 + asin(1/5)}
    \tkzDefPoint({3*cos(\ang)},{3*sin(\ang)}){A}
    \tkzDefPoint({5 + 2*cos(\ang)},{2*sin(\ang)}){B}
    \tkzDefPoint(3,4){delta_top}
    \tkzDefPoint(3,-3){delta_bot}
    \tkzInterLL(A,B)(delta_top,delta_bot) \tkzGetPoint{M}
    \tkzDefPointBy[homothety=center A ratio -0.5](B) \tkzGetPoint{eps_left}
    \tkzDefPointBy[homothety=center B ratio -0.1](A) \tkzGetPoint{eps_right}
    \tkzDrawSegment[l3](eps_left,eps_right)
    \tkzDrawSegment[l3](delta_top,delta_bot)
    \tkzDrawSegments[l4,maincolor](O,K O,A K,B)
    \tkzDrawPoints(O,K,A,B,N,M)
    \tkzLabelPoint[below left](O){O}
    \tkzLabelPoint[below right](K){K}
    \tkzLabelPoint[above left=2pt](A){{A}}
    \tkzLabelPoint[above right=2pt](B){{B}}
    \tkzLabelPoint[below right](N){{N}}
    \tkzLabelPoint[above right](M){{M}}
    \tkzLabelPoint[above](eps_left){$(\epsilon)$}
    \tkzLabelPoint[right](delta_top){$(\delta)$}
\end{tikzpicture}
\end{wrapfigure}
Δύο κύκλοι $( O, \rho_1), ( K, \rho_2)$ εφάπτονται εξωτερικά στο $ N$. Μια ευθεία $(\varepsilon)$ εφάπτεται στους δύο κύκλους στα σημεία $A, B$ αντίστοιχα. Η κοινή εφαπτομένη των κύκλων στο $\mathrm{N}$ τέμνει την $(\varepsilon)$ στο $ M$.
Να αποδείξετε ότι:
\begin{alist}
    \item Το $\mathrm{M}$ είναι μέσον του $AB$. \monades{7}
    \item $ O\hat{ M} K = 90^\circ$. \monades{9}
    \item $A\hat{\mathrm{N}}B = 90^\circ$. \monades{9}
\end{alist}

\vspace{4mm}

% --- Άσκηση 1773 (1773.tex) ---
\Askhsh{1773}{4}
\begin{ylh}{1}
    \item 5.2. Παραλληλόγραμμα
    \item 5.4. Ρόμβος
    \item 5.6. Εφαρμογές στα τρίγωνα.
\end{ylh}
\begin{wrapfigure}[5]{r}{7cm}
\begin{tikzpicture}[scale=0.7]
    \tkzDefPoint(0,0){D}
    \tkzDefPoint(10,-0.5){C}
    \tkzDefPoint(2,3.5){A}
    \tkzDefPoint(8.2, 2.7){B} 
    \tkzDefMidPoint(A,B) \tkzGetPoint{E}
    \tkzDefMidPoint(B,C) \tkzGetPoint{L}
    \tkzDefMidPoint(C,D) \tkzGetPoint{Z}
    \tkzDefMidPoint(D,A) \tkzGetPoint{K}
    \tkzDefMidPoint(D,B) \tkzGetPoint{N}
    \tkzDefMidPoint(A,C) \tkzGetPoint{M}
    \tkzDrawSegments[l4,maincolor](A,D B,C)
    \tkzDrawSegments[l3](A,B C,D)
    \tkzDrawSegments[l3](A,C B,D)
    \tkzDrawPoints(A,B,C,D,E,L,Z,K,N,M)
    \tkzLabelPoint[above](A){A}
    \tkzLabelPoint[above right](B){B}
    \tkzLabelPoint[below right](C){$\varGamma $}
    \tkzLabelPoint[below left](D){$\varDelta  $}
    \tkzLabelPoint[above](E){E}
    \tkzLabelPoint[right=2pt](L){$\Lambda$}
    \tkzLabelPoint[below](Z){Z}
    \tkzLabelPoint[left=2pt](K){K}
    \tkzLabelPoint[below right](N){N}
    \tkzLabelPoint[below](M){M}
\end{tikzpicture}
\end{wrapfigure}
Δίνεται τετράπλευρο $AB\varGamma \varDelta  $ με $A\varDelta   = B\varGamma $. Αν $ E, \Lambda,  Z,  K,  N,  M$ είναι τα μέσα των $ AB, B\varGamma , \varGamma \varDelta  , \varDelta   A, \varDelta   B$ και $A\varGamma $ αντίστοιχα, να αποδείξετε ότι:
\begin{alist}
    \item Το τετράπλευρο $ E M Z N$ ρόμβος. \monades{8}
    \item Η $E\mathrm{Z}$ είναι μεσοκάθετος του ευθύγραμμου τμήματος $\mathrm{M}\mathrm{N}$. \monades{7}
    \item $ K E =  Z\Lambda$. \monades{5}
    \item Τα ευθύγραμμα τμήματα $ K\Lambda,  M N,  E Z$ διέρχονται από το ίδιο σημείο. \monades{5}
\end{alist}

\vspace{4mm}

% --- Άσκηση 1775 (1775.tex) ---
\Askhsh{1775}{4}
\begin{ylh}{1}
    \item 3.6. Κριτήρια ισότητας ορθογώνιων τριγώνων
    \item 5.2. Παραλληλόγραμμα
    \item 5.6. Εφαρμογές στα τρίγωνα.
\end{ylh}
\begin{wrapfigure}[5]{r}{5cm}
\begin{tikzpicture}[scale=0.7]
    \tkzDefPoint(0,0){D} % \Delta
    \tkzDefPoint(5,0){C} % \varGamma 
    \tkzDefPoint(3,3){B}
    \tkzDefPoint(-2,3){A}
    \tkzDefMidPoint(A,D) \tkzGetPoint{M}
    \tkzDefLine[orthogonal=through C](M,B) \tkzGetPoint{c_perp}
    \tkzInterLL(C,c_perp)(M,B) \tkzGetPoint{E}
    \tkzDefLine[parallel=through D](M,B) \tkzGetPoint{d_par}
    \tkzInterLL(D,d_par)(B,C) \tkzGetPoint{N}
    \tkzInterLL(D,d_par)(C,E) \tkzGetPoint{Z}
    \tkzDefPointBy[homothety=center D ratio 1.25](Z) \tkzGetPoint{X_pt}
    \tkzDrawPolygon[l3](A,B,C,D)
    \tkzDrawSegments[l4,maincolor](M,E C,E D,E)
    \tkzDrawSegment[dashed, thick](D,X_pt)
    \tkzMarkRightAngle[size=0.3](C,E,M)
    \tkzDrawPoints(A,B,C,D,M,E,N,Z)
    \tkzLabelPoint[above left](A){A}
    \tkzLabelPoint[above](B){B}
    \tkzLabelPoint[below right](C){$\varGamma $}
    \tkzLabelPoint[below left](D){$\varDelta  $}
    \tkzLabelPoint[left](M){M}
    \tkzLabelPoint[above right](E){E}
    \tkzLabelPoint[below right](N){N}
    \tkzLabelPoint[below right](Z){Z}
    \tkzLabelPoint[above](X_pt){$x$}
\end{tikzpicture}
\end{wrapfigure}
Δίνεται παραλληλόγραμμο $AB\varGamma \Delta$. Θεωρούμε το μέσο $\mathrm{M}$ της πλευράς $ A\Delta$ και $\varGamma E$ κάθετος από τη κορυφή $\varGamma $ στην ευθεία $\mathrm{M}B$ ($\varGamma E \perp \mathrm{M}B$). Η παράλληλη από την κορυφή $\Delta$ στην ευθεία $\mathrm{M}B$ ($\Delta x \parallel \mathrm{M}B$) τέμνει τις $B\varGamma $ και $\varGamma E$ στα σημεία $\mathrm{N}, \mathrm{Z}$ αντίστοιχα.
Να αποδείξετε ότι:
\begin{alist}
    \item Το τετράπλευρο $ $MB$ N\varDelta  $ είναι παραλληλόγραμμο. \monades{7}
    \item Το σημείο $ Z$ είναι μέσον του ευθυγράμμου τμήματος $\varGamma  E$. \monades{9}
    \item $\varDelta   E = \varDelta  \varGamma $. \monades{9}
\end{alist}

\vspace{4mm}

% --- Άσκηση 1777 (1777.tex) ---
\Askhsh{1777}{4}
\begin{ylh}{1}
    \item 5.2. Παραλληλόγραμμα
    \item 5.3. Ορθογώνιο
    \item 5.6. Εφαρμογές στα τρίγωνα
    \item 5.8. Το ορθόκεντρο τριγώνου.
\end{ylh}
\begin{wrapfigure}[10]{r}{7cm}
\begin{tikzpicture}[scale=0.7]
    % Ορισμός κορυφών τριγώνου
    \tkzDefPoint(0,0){B}
    \tkzDefPoint(8,0){C} % \varGamma 
    \tkzDefPoint(3,5){A}
    \tkzDefPointBy[projection=onto A--C](B) \tkzGetPoint{E}
    \tkzDefPointBy[projection=onto A--B](C) \tkzGetPoint{Z}
    \tkzInterLL(B,E)(C,Z) \tkzGetPoint{H}
    \tkzDefMidPoint(A,B) \tkzGetPoint{M}
    \tkzDefMidPoint(A,C) \tkzGetPoint{N}
    \tkzDefMidPoint(C,H) \tkzGetPoint{K}
    \tkzDefMidPoint(B,H) \tkzGetPoint{L} % \Lambda
    \tkzDefMidPoint(B,C) \tkzGetPoint{O}
    \tkzDrawPolygon[l3](A,B,C)
    \tkzDrawSegments[l4,maincolor](B,E C,Z)
    \tkzDrawPolygon[dashed, thick](M,N,K,L)
    \tkzMarkRightAngle[size=0.3](A,E,B)
    \tkzMarkRightAngle[size=0.3](A,Z,C)
    \tkzDrawPoints(A,B,C,E,Z,H,M,N,K,L,O)
    \tkzLabelPoint[above](A){A}
    \tkzLabelPoint[below left](B){B}
    \tkzLabelPoint[below right](C){$\varGamma $}
    \tkzLabelPoint[above right](E){E}
    \tkzLabelPoint[above left](Z){Z}
    \tkzLabelPoint[above=3pt](H){H}
    \tkzLabelPoint[left=2pt](M){M}
    \tkzLabelPoint[right=2pt](N){N}
    \tkzLabelPoint[below right](K){K}
    \tkzLabelPoint[below left](L){$\Lambda$}
    \tkzLabelPoint[below](O){O}
\end{tikzpicture}
\end{wrapfigure}
Δίνονται οξυγώνιο τρίγωνο $AB\varGamma $, $B E, \varGamma  Z$, τα ύψη από τις κορυφές $B, \varGamma $ αντίστοιχα και $ H$ το ορθόκεντρο του τριγώνου. Επίσης δίνονται τα $ M,  N,  K, \Lambda$ μέσα των ευθυγράμμων τμημάτων $ AB, A\varGamma , \varGamma  H, B H$ αντίστοιχα.
\begin{alist}
    \item Να αποδείξετε ότι:
    \begin{rlist}
        \item $ M N = \Lambda K$ \monades{6}
        \item $ N K =  M\Lambda = \frac{A H}{2}$ \monades{6}
        \item Το τετράπλευρο $\mathrm{M}\mathrm{N} K\Lambda$ είναι ορθογώνιο. \monades{6}
    \end{rlist}
    \item Αν το $ O$ είναι το μέσο της $B\varGamma $, να αποδείξετε ότι $ M\hat{\mathrm{O}} K = 90^\circ$. \monades{7}
\end{alist}

\vspace{4mm}

% --- Άσκηση 1778 (1778.tex) ---
\Askhsh{1778}{4}
\begin{ylh}{1}
    \item 3.6. Κριτήρια ισότητας ορθογώνιων τριγώνων
    \item 5.9. Μια ιδιότητα του ορθογώνιου τριγώνου
    \item 5.10. Τραπέζιο.
\end{ylh}
\begin{wrapfigure}[9]{r}{5cm}
\begin{tikzpicture}[scale=0.7]
    \tkzDefPoint(0,0){O}
    \tkzDefPoint(0,3.5){A}
    \tkzDefPoint(3.5,0){B}
    \tkzDefPoint(4.5,0){X}
    \tkzDefPoint(0,4.5){Y}
    \tkzDefPoint({5*cos(108 deg)},{5*sin(108 deg)}){Eps_top}
    \tkzDefPoint({-2.5*cos(108 deg)},{-2.5*sin(108 deg)}){Eps_bot}
    \tkzDefPointBy[projection=onto Eps_top--Eps_bot](A) \tkzGetPoint{D}
    \tkzDefPointBy[projection=onto Eps_top--Eps_bot](B) \tkzGetPoint{E}
    \tkzDefMidPoint(A,B) \tkzGetPoint{M}
    \tkzDefMidPoint(D,E) \tkzGetPoint{N}
    \tkzDrawSegment(O,X)
    \tkzDrawSegment(O,Y)
    \tkzDrawSegment[thick](Eps_top,Eps_bot)
    \tkzDrawSegments[l3](A,B A,D B,E)
    \tkzDrawSegments[l4, color=maincolor](M,D M,E M,N M,O)
    \tkzMarkRightAngle[size=0.3](Y,O,X)
    \tkzMarkRightAngle[size=0.3](A,D,Eps_top)
    \tkzMarkRightAngle[size=0.3](B,E,Eps_bot)
    \tkzDrawPoints(O,A,B,D,E,M,N)
    \tkzLabelPoint[below left](O){O}
    \tkzLabelPoint[right=2pt](A){A}
    \tkzLabelPoint[above right](B){B}
    \tkzLabelPoint[left=2pt](D){$\varDelta  $}
    \tkzLabelPoint[below left](E){E}
    \tkzLabelPoint[right=2pt](M){M}
    \tkzLabelPoint[left=2pt](N){N}
    \tkzLabelPoint[below](X){$x$}
    \tkzLabelPoint[left](Y){$y$}
    \tkzLabelPoint[right](Eps_top){$\varepsilon$}
\end{tikzpicture}
\end{wrapfigure}
Δίνεται ορθή γωνία $x\hat{\mathrm{O}}y$ και τα σημεία $A$ και $B$ των ημιευθειών $\mathrm{O}y$ και $\mathrm{O}x$ αντίστοιχα με $\mathrm{O}A = \mathrm{O}B$. Μία ευθεία $(\varepsilon)$ η οποία δεν είναι παράλληλη στην $ AB$ διέρχεται από το $\mathrm{O}$ ώστε τα σημεία $A$ και $B$ να είναι στο ίδιο ημιεπίπεδο. Η κάθετη από το $ A$ στην $(\varepsilon)$ την τέμνει στο $\Delta$ και η κάθετη από το $B$ στην $(\varepsilon)$ την τέμνει στο $E$.
Να αποδείξετε ότι:
\begin{alist}
    \item Τα τρίγωνα $ $OA$\varDelta  $ και $ O EB$ είναι ίσα. \monades{7}
    \item $A\varDelta   + B E = \varDelta   E$. \monades{7}
    \item $\mathrm{M}\mathrm{N} = \frac{\Delta E}{2}$, όπου $\mathrm{M}$ και $\mathrm{N}$ τα μέσα των $AB$ και $\Delta E$ αντίστοιχα. \monades{7}
    \item Το τρίγωνο $\varDelta   M E$ είναι ορθογώνιο και ισοσκελές. \monades{4}
\end{alist}

\vspace{4mm}

% --- Άσκηση 1780 (1780.tex) ---
\Askhsh{1780}{4}
\begin{ylh}{1}
    \item 5.5. Τετράγωνο
    \item 5.6. Εφαρμογές στα τρίγωνα
    \item 5.8. Το ορθόκεντρο τριγώνου.
\end{ylh}
\begin{wrapfigure}[5]{r}{7cm}
\begin{tikzpicture}[scale=0.7]
    \tkzDefPoint(0,0){D} % \Delta
    \tkzDefPoint(3,-1){C} % \varGamma 
    \tkzDefPoint(1,3){A}
    \tkzDefPointBy[translation=from D to C](A) \tkzGetPoint{B} 
    \tkzDefPointBy[symmetry=center D](B) \tkzGetPoint{E}
    \tkzDefMidPoint(A,D) \tkzGetPoint{M}
    \tkzInterLL(A,E)(C,D) \tkzGetPoint{N}
    \tkzDefPointBy[homothety=center D ratio 1.4](N) \tkzGetPoint{N_ext}
    \tkzDrawPolygon[l3](A,B,C,D)
    \tkzDrawSegments[l4,maincolor](D,B E,D A,E)
    \tkzDrawSegment[dashed, l3](D,N_ext)
    \tkzDrawPoints(A,B,C,D,E,M,N)
    \tkzLabelPoint[above right](A){{A}}
    \tkzLabelPoint[above right](B){{B}}
    \tkzLabelPoint[below](C){{$\varGamma $}}
    \tkzLabelPoint[below right](D){{$\varDelta  $}}
    \tkzLabelPoint[below left](E){{E}}
    \tkzLabelPoint[right=2pt](M){{M}}
    \tkzLabelPoint[above left](N){{N}}
\end{tikzpicture}
\end{wrapfigure}
Σε τετράγωνο $AB\varGamma \varDelta  $ προεκτείνουμε τη διαγώνιο $B\varDelta  $ (προς το $\varDelta  $) κατά τμήμα $\varDelta   E = \varDelta   B$. Έστω $ M$ το μέσο της $ A\varDelta  $ και $ N$ το σημείο τομής των ευθειών $ AE$ και $\varGamma \varDelta  $.
\begin{alist}
    \item Να αποδείξετε ότι $\varDelta   N = \varDelta   M$. \monades{6}
    \item Να υπολογίσετε τις γωνίες του τριγώνου $\mathrm{N} M\varDelta  $. \monades{5}
    \item Να αποδείξετε ότι:
    \begin{rlist}
        \item $ M N \perp A\varGamma $ \monades{7}
        \item $\varGamma  M \perp A N$ \monades{7}
    \end{rlist}
\end{alist}

\vspace{4mm}

% --- Άσκηση 1781 (1781.tex) ---
\Askhsh{1781}{4}
\begin{ylh}{1}
    \item 5.2. Παραλληλόγραμμα
    \item 5.3. Ορθογώνιο
    \item 5.5. Τετράγωνο
    \item 5.6. Εφαρμογές στα τρίγωνα.
\end{ylh}
\begin{wrapfigure}[5]{r}{5cm}
\begin{tikzpicture}[scale=0.7]
    % Ορισμός κορυφών τετραγώνου
    \tkzDefPoint(0,0){D} % \Delta
    \tkzDefPoint(4,0){C} % \varGamma 
    \tkzDefPoint(4,4){B}
    \tkzDefPoint(0,4){A}
    
    % Ορισμός του κέντρου O
    \tkzDefPoint(2,2){O}
    
    % Ορισμός των σημείων I και H στη διαγώνιο (1/4 και 3/4 της απόστασης)
    \tkzDefMidPoint(A,O) \tkzGetPoint{I}
    \tkzDefMidPoint(O,C) \tkzGetPoint{H}
    
    % Ορισμός των μέσων των πλευρών
    \tkzDefMidPoint(D,C) \tkzGetPoint{E}
    \tkzDefMidPoint(B,C) \tkzGetPoint{Z}
    \tkzDefMidPoint(A,B) \tkzGetPoint{Th} % \Theta
    
    % Σχεδίαση του περιγράμματος
    \tkzDrawPolygon[l3](A,B,C,D)
    
    % Σχεδίαση της διαγωνίου
    \tkzDrawSegment[l3](A,C)
    
    % Σχεδίαση των εσωτερικών τμημάτων
    \tkzDrawSegments[l4,maincolor](Th,I Th,Z Z,H H,E E,O O,Z)
    
    % Σημεία (κόκκινες κουκκίδες όπως στην εικόνα)
    \tkzDrawPoints(A,B,C,D,O,I,H,E,Z,Th)
    
    % Ετικέτες
    \tkzLabelPoint[above left](A){A}
    \tkzLabelPoint[above right](B){B}
    \tkzLabelPoint[below right](C){$\varGamma $}
    \tkzLabelPoint[below left](D){$\varDelta  $}
    \tkzLabelPoint[below](E){E}
    \tkzLabelPoint[right](Z){Z}
    \tkzLabelPoint[above](Th){$\varTheta $}
    \tkzLabelPoint[below left](I){I}
    \tkzLabelPoint[left=2pt](O){O}
    \tkzLabelPoint[below left](H){H}
\end{tikzpicture}
\end{wrapfigure}
Δίνεται το τετράγωνο $AB\varGamma \varDelta  $. Στη διαγώνιο $ A\varGamma $ θεωρούμε σημεία $ I,  O,  H$ ώστε $A\mathrm{I} =  I O =  O H =  H\varGamma $. Αν $ E, \varTheta $ και $ Z$ τα μέσα των πλευρών $\varDelta  \varGamma , AB$ και $B\varGamma $ αντίστοιχα να αποδείξετε ότι:
\begin{alist}
    \item Το τετράπλευρο $ O Z\varGamma  E$ είναι τετράγωνο. \monades{7}
    \item $ Z H = \frac{A\varGamma }{4}$. \monades{8}
    \item Το τετράπλευρο $ I\varTheta  Z H$ είναι ορθογώνιο παραλληλόγραμμο, με $\varTheta  Z = 2\varTheta  I$. \monades{10}
\end{alist}

\vspace{4mm}

% --- Άσκηση 1783 (1783.tex) ---
\Askhsh{1783}{4}
\begin{ylh}{1}
    \item 3.3. 2ο Κριτήριο ισότητας τριγώνων
    \item 4.2. Τέμνουσα δύο ευθειών - Ευκλείδειο αίτημα.
\end{ylh}
\begin{wrapfigure}[5]{r}{5cm}
\begin{tikzpicture}[scale=0.7]
    \tkzDefPoint(0,0){D} % \Delta
    \tkzDefPoint(3.5,0){C} % \varGamma 
    \tkzDefPoint(5,0){Z}
    \tkzDefPoint(1.4,4.8){A}
    \tkzDefPoint(2.9,4.8){B}
    \tkzInterLL(A,Z)(B,C) \tkzGetPoint{E}
    
    % Σχεδίαση περιγράμματος τραπεζίου
    \tkzDrawSegments[l3](A,B B,C C,D D,A)
    \tkzDrawSegments[l4,maincolor](A,Z D,E C,Z)
    
    % Σημεία και Ετικέτες
    \tkzDrawPoints(A,B,C,D,Z,E)
    \tkzLabelPoint[above](A){A}
    \tkzLabelPoint[above](B){B}
    \tkzLabelPoint[below](C){$\varGamma $}
    \tkzLabelPoint[below](D){$\varDelta  $}
    \tkzLabelPoint[below](Z){Z}
    \tkzLabelPoint[right=2pt](E){E}
\end{tikzpicture}
\end{wrapfigure}
Σε τραπέζιο $AB\varGamma \Delta$ ($AB \parallel \varGamma \Delta$) ισχύει $AB + \varGamma \Delta = A\Delta$. Αν η διχοτόμος της γωνίας $\hat{A}$ τέμνει την $B\varGamma $ στο $E$ και την προέκταση της $\Delta\varGamma $ στο $\mathrm{Z}$, να αποδείξετε ότι:
\begin{alist}
    \item Το τρίγωνο $\varDelta   A Z$ είναι ισοσκελές. \monades{7}
    \item Το $ E$ είναι το μέσο της $B\varGamma $. \monades{10}
    \item Η $\varDelta   E$ είναι διχοτόμος της γωνίας $\hat{\varDelta  }$ του τραπεζίου. \monades{8}
\end{alist}

\vspace{4mm}

% --- Άσκηση 1784 (1784.tex) ---
\Askhsh{1784}{4}
\begin{ylh}{1}
    \item 4.2. Τέμνουσα δύο ευθειών - Ευκλείδειο αίτημα
    \item 4.6. Άθροισμα γωνιών τριγώνου.
\end{ylh}
\begin{wrapfigure}[5]{r}{7cm}
\begin{tikzpicture}[scale=0.7]
    % Ορισμός των σημείων στη βάση $AB$
    \tkzDefPoint(0,0){A}
    \tkzDefPoint(3.5,0){Z}
    \tkzDefPoint(6,0){B}
    
    % Ορισμός του $\varDelta$ ώστε $A\varDelta  $ = $AZ$ (μήκος 3.5, γωνία 70 μοίρες)
    \tkzDefPoint(70:3.5){D} % \Delta
    
    % Ορισμός του $E$ ώστε $BE$ παράλληλο στο $A\varDelta  $ και $BE$ = $BZ$ (μήκος 2.5)
    \tkzDefShiftPoint[B](70:2.5){E}
    
    % Ορισμός του Ο ως μέσο του $\varDelta  E$
    \tkzDefMidPoint(D,E) \tkzGetPoint{O}
    
    % Σχεδίαση περιγράμματος
    \tkzDrawSegments[l3](A,D D,E E,B B,A)
    
    % Σχεδίαση εσωτερικών τμημάτων
    \tkzDrawSegments[l4,maincolor](D,Z E,Z O,A O,B)
    
    % Προσθήκη συμβόλων ισότητας (// και /)
    \tkzMarkSegments[mark=||, size=3pt](A,D A,Z)
    \tkzMarkSegments[mark=|, size=3pt](B,E B,Z)
    
    % Σημεία και Ετικέτες
    \tkzDrawPoints(A,Z,B,D,E,O)
    \tkzLabelPoint[below left](A){A}
    \tkzLabelPoint[below](Z){Z}
    \tkzLabelPoint[below right](B){B}
    \tkzLabelPoint[above left](D){$\varDelta  $}
    \tkzLabelPoint[right](E){E}
    \tkzLabelPoint[above](O){O}
\end{tikzpicture}
\end{wrapfigure}
Δίνεται τραπέζιο $A\varDelta   EB$, με $A\varDelta   \parallel B E$, στο οποίο ισχύει ότι $ AB = A\varDelta   + BE$, και $ O$ το μέσον της $\varDelta  E$. Θεωρούμε σημείο $ Z$ στην $ AB$ τέτοιο ώστε $A Z =  A\varDelta  $ και $B Z = B E$.
Αν γωνία $\varDelta  \hat{A} Z = \varphi$,
\begin{alist}
    \item να εκφράσετε τη γωνία $A\hat{ Z}\varDelta  $ σε συνάρτηση με τη $\varphi$. \monades{8}
    \item να εκφράσετε τη γωνία $ E\hat{ Z}B$ σε συνάρτηση με τη $\varphi$. \monades{8}
    \item να αποδείξετε ότι οι $ OA$ και $ OB$ είναι μεσοκάθετοι των τμημάτων $\varDelta   Z$ και $ Z E$ αντίστοιχα. \monades{9}
\end{alist}

\vspace{4mm}

% --- Άσκηση 1785 (1785.tex) ---
\Askhsh{1785}{4}
\begin{ylh}{1}
    \item 3.11. Ανισοτικές σχέσεις πλευρών και γωνιών
    \item 4.2. Τέμνουσα δύο ευθειών - Ευκλείδειο αίτημα
    \item 4.6. Άθροισμα γωνιών τριγώνου
    \item 5.2. Παραλληλόγραμμα.
\end{ylh}
\begin{wrapfigure}[5]{r}{7cm}
\begin{tikzpicture}[scale=0.7]
    % Ορισμός κορυφών
    \tkzDefPoint(0,0){D} % \Delta
    \tkzDefPoint(6,0){C} % \varGamma 
    \tkzDefPoint(2,4){A}
    \tkzDefPoint(8,4){B}
    \tkzDefPoint(1.2, 2.4){K} 
    \tkzDefPoint(3.7888, 4){L}
    \tkzDefMidPoint(K,L) \tkzGetPoint{M}
    \tkzInterLL(A,M)(D,C) \tkzGetPoint{E}
    \tkzDrawPolygon[l3](A,B,C,D)
    \tkzDrawSegments[l4,maincolor](K,L A,E)
    \tkzDrawPoints(A,B,C,D,K,L,M,E)
    \tkzLabelPoint[above](A){A}
    \tkzLabelPoint[above](B){B}
    \tkzLabelPoint[right](C){$\varGamma $}
    \tkzLabelPoint[left](D){$\varDelta  $}
    \tkzLabelPoint[left=2pt](K){K}
    \tkzLabelPoint[above](L){$\Lambda$}
    \tkzLabelPoint[right=2pt](M){M}
    \tkzLabelPoint[above right](E){E}
\end{tikzpicture}
\end{wrapfigure}
Δίνεται παραλληλόγραμμο $AB\varGamma \Delta$, με $AB > A\Delta$. Θεωρούμε σημεία $ K, \Lambda$, των $A\Delta$ και $AB$ αντίστοιχα ώστε $ A\mathrm{K} = A\Lambda$. Έστω $\mathrm{M}$ το μέσο του $\mathrm{K}\Lambda$ και η προέκταση της $ A\mathrm{M}$ (προς το $\mathrm{M}$) τέμνει τη $\Delta\varGamma $ στο σημείο $E$.
Να αποδείξετε ότι:
\begin{alist}
    \item $A\varDelta   = \varDelta   E$. \monades{8}
    \item $B\varGamma  + \varGamma  E = AB$. \monades{10}
    \item $\hat{B} = 2 \cdot A\hat{\Lambda} K$. \monades{7}
\end{alist}

\vspace{4mm}

% --- Άσκηση 1786 (1786.tex) ---
\Askhsh{1786}{4}
\begin{ylh}{1}
    \item 3.11. Ανισοτικές σχέσεις πλευρών και γωνιών
    \item 5.2. Παραλληλόγραμμα
    \item 5.4. Ρόμβος
    \item 5.9. Μια ιδιότητα του ορθογώνιου τριγώνου.
\end{ylh}
Δίνεται παραλληλόγραμμο $AB\varGamma \varDelta  $ με $AB = 2B\varGamma $ και τη γωνία $\hat{B}$ αμβλεία. Από την κορυφή $ A$ φέρνουμε την $A E$ κάθετη στην ευθεία $B\varGamma $ και έστω $ M,  N$ τα μέσα των $AB, \varDelta  \varGamma $ αντίστοιχα.
Να αποδείξετε ότι:
\begin{alist}
    \item Το τετράπλευρο $\mathrm{M}B\varGamma \mathrm{N}$ είναι ρόμβος. \monades{8}
    \item Το τετράπλευρο $\mathrm{M}E\varGamma \mathrm{N}$ είναι ισοσκελές τραπέζιο. \monades{9}
    \item Η $ E N$ είναι διχοτόμος της γωνίας $ M\hat{ E}\varGamma $. \monades{8}
\end{alist}

\vspace{4mm}

% --- Άσκηση 1787 (1787.tex) ---
\Askhsh{1787}{4}
\begin{ylh}{1}
    \item 3.3. 2ο Κριτήριο ισότητας τριγώνων
    \item 4.2. Τέμνουσα δύο ευθειών - Ευκλείδειο αίτημα
    \item 5.6. Εφαρμογές στα τρίγωνα.
\end{ylh}
\begin{wrapfigure}[5]{r}{7cm}
\begin{tikzpicture}[scale=0.7]
    \tkzDefPoint(2,5){A}
    \tkzDefPoint(8,4){B}
    \tkzDefPointBy[rotation=center A angle -105](B) \tkzGetPoint{dirD}
    \tkzDefPointBy[homothety=center A ratio 0.5](dirD) \tkzGetPoint{D}
    \tkzDefPointBy[translation=from A to D](B) \tkzGetPoint{C} 
    \tkzDefMidPoint(D,C) \tkzGetPoint{M}
    \tkzDefLine[orthogonal=through A](A,D) \tkzGetPoint{a_perp}
    \tkzInterLL(A,a_perp)(B,C) \tkzGetPoint{H}
    \tkzInterLL(H,M)(A,D) \tkzGetPoint{E}
    \tkzDrawPolygon[l3](A,B,C,D)
    \tkzDrawSegments[l4,maincolor](A,H E,H D,E)
    \tkzDrawSegment[l3](A,M)
    \tkzMarkRightAngle[size=0.3](D,A,H)
    \tkzDrawPoints(A,B,C,D,M,H,E)
    \tkzLabelPoint[above](A){A}
    \tkzLabelPoint[right](B){B}
    \tkzLabelPoint[below right](C){$\varGamma $}
    \tkzLabelPoint[above left](D){$\varDelta  $}
    \tkzLabelPoint[below](M){M}
    \tkzLabelPoint[right=2pt](H){H}
    \tkzLabelPoint[left](E){E}
\end{tikzpicture}
\end{wrapfigure}
Δίνεται παραλληλόγραμμο $AB\varGamma \Delta$ με $AB = 2 B\varGamma $, τη γωνία $\hat{A}$ αμβλεία και $ M$ το μέσο της $\varGamma \Delta$. Φέρουμε κάθετη στην $A\Delta$ στο σημείο $ A$, η οποία τέμνει την $B\varGamma $ στο $\mathrm{H}$. Αν η προέκταση της $\mathrm{H} M$ τέμνει την προέκταση της $A\Delta$ στο $E$, να αποδείξετε ότι:
\begin{alist}
    \item Η $A M$ είναι διχοτόμος της γωνίας $\varDelta  \hat{A}B$. \monades{9}
    \item Τα τμήματα $ E H, \varDelta  \varGamma $ διχοτομούνται. \monades{8}
    \item $\hat{ E} = \varDelta  \hat{ M}A$. \monades{8}
\end{alist}

\vspace{4mm}

% --- Άσκηση 1788 (1788.tex) ---
\Askhsh{1788}{4}
\begin{ylh}{1}
    \item 3.2. 1ο Κριτήριο ισότητας τριγώνων
    \item 4.6. Άθροισμα γωνιών τριγώνου
    \item 5.5. Τετράγωνο.
\end{ylh}
Δίνεται οξυγώνιο τρίγωνο $AB\varGamma $ και στο εξωτερικό του σχηματίζονται τα τετράγωνα $ $AB$\varDelta   E$ και $A\varGamma  Z H$.
Να αποδείξετε ότι:
\begin{alist}
    \item $ E\hat{A} H = A\hat{B}\varGamma  + A\hat{\varGamma }B$ \monades{8}
    \item $ E\varGamma  = B H$ \monades{9}
    \item Η $ E\varGamma $ είναι κάθετη στη $B H$. \monades{8}
\end{alist}
\begin{center}
\begin{tikzpicture}[scale=0.7]
    \tkzDefPoint(3,4){A}
    \tkzDefPoint(0.5,0){B}
    \tkzDefPoint(5,0){C}
    \tkzDefPointBy[rotation=center A angle -90](B) \tkzGetPoint{E}
    \tkzDefPointBy[rotation=center B angle 90](A) \tkzGetPoint{D}
    \tkzDefPointBy[rotation=center A angle 90](C) \tkzGetPoint{H}
    \tkzDefPointBy[rotation=center C angle -90](A) \tkzGetPoint{Z}
    \tkzDrawPolygon[l3](A,B,C)
    \tkzDrawPolygon[l4](A,B,D,E)
    \tkzDrawPolygon[l4](A,C,Z,H)
    \tkzDrawSegments[l3,maincolor](E,H E,C B,H)
    \tkzDrawPoints(A,B,C,D,E,Z,H)
    \tkzLabelPoint[above](A){A}
    \tkzLabelPoint[below left](B){B}
    \tkzLabelPoint[below right](C){$\varGamma $}
    \tkzLabelPoint[left](D){$\varDelta  $}
    \tkzLabelPoint[above left](E){E}
    \tkzLabelPoint[right](Z){Z}
    \tkzLabelPoint[above right](H){H}
\end{tikzpicture}
\end{center}
\vspace{4mm}

% --- Άσκηση 1789 (1789.tex) ---
\Askhsh{1789}{4}
\begin{ylh}{1}
    \item 5.2. Παραλληλόγραμμα
    \item 5.6. Εφαρμογές στα τρίγωνα
    \item 5.9. Μια ιδιότητα του ορθογώνιου τριγώνου
    \item 5.11. Ισοσκελές τραπέζιο.
\end{ylh}
\begin{wrapfigure}[5]{r}{5.5cm}
\begin{tikzpicture}[scale=0.7]
    \tkzDefPoint(0,0){A}
    \tkzDefPoint(6,0){B}
    \tkzDefPoint(0,4){C}
    \tkzDefPoint(2.5,0){D}
    \tkzDefMidPoint(C,D) \tkzGetPoint{K}
    \tkzDefMidPoint(B,C) \tkzGetPoint{M}
    \tkzDefMidPoint(B,D) \tkzGetPoint{N}
    \tkzDrawPolygon[l3](A,B,C)
    \tkzDrawSegments[l4,maincolor](C,D M,N)
    \tkzDrawSegments[l3](K,M K,N)
    \tkzMarkRightAngle[size=0.3](B,A,C)
    \tkzDrawPoints(A,B,C,D,K,M,N)
    \tkzLabelPoint[below left](A){A}
    \tkzLabelPoint[below right](B){B}
    \tkzLabelPoint[above](C){$\varGamma $}
    \tkzLabelPoint[below](D){$\varDelta  $} 
    \tkzLabelPoint[left](K){K}
    \tkzLabelPoint[above right](M){M}
    \tkzLabelPoint[above right](N){N}
\end{tikzpicture}
\end{wrapfigure}
Δίνεται ορθογώνιο τρίγωνο $AB\varGamma $ με τη γωνία $\hat{ A}$ ορθή, και τυχαίο σημείο $\varDelta  $ της πλευράς $ AB$. Έστω $ K,  M,  N$ τα μέσα των $\varGamma \varDelta  , B\varGamma , B\varDelta  $ αντίστοιχα.
Να αποδείξετε ότι:
\begin{alist}
    \item Το τετράπλευρο $ K\mathrm{M}\mathrm{N}\Delta$ είναι παραλληλόγραμμο. \monades{8}
    \item Το τετράπλευρο $A K M N$ είναι ισοσκελές τραπέζιο. \monades{9}
    \item Η διάμεσος του τραπεζίου $A K M N$ είναι ίση με $\frac{AB}{2}$. \monades{8}
\end{alist}

\vspace{4mm}

% --- Άσκηση 1790 (1790.tex) ---
\Askhsh{1790}{4}
\begin{ylh}{1}
    \item 3.2. 1ο Κριτήριο ισότητας τριγώνων
    \item 4.2. Τέμνουσα δύο ευθειών - Ευκλείδειο αίτημα
    \item 5.2. Παραλληλόγραμμα
    \item 5.6. Εφαρμογές στα τρίγωνα
    \item 5.10. Τραπέζιο
    \item 5.11. Ισοσκελές τραπέζιο.
\end{ylh}
\begin{wrapfigure}[5]{r}{7cm}
\begin{tikzpicture}[scale=0.7]
    \tkzDefPoint(0,0){B}
    \tkzDefPoint(7,0){C}
    \tkzDefPoint(70:4){A}
    \tkzDefPointBy[translation=from B to C](A) \tkzGetPoint{D}
    \tkzDefPointBy[projection=onto B--C](A) \tkzGetPoint{E}
    \tkzDefPointBy[symmetry=center E](B) \tkzGetPoint{Z}
    \tkzDefMidPoint(B,D) \tkzGetPoint{M}
    \tkzDrawPolygon[l3](A,B,C,D)
    \tkzDrawSegments[l3](A,E A,Z B,D)
	\tkzDrawSegments[l4,maincolor](E,M Z,D)
    \tkzMarkRightAngle[size=0.3](A,E,B)
    \tkzDrawPoints(A,B,C,D,E,Z,M)
    \tkzLabelPoint[above](A){A}
    \tkzLabelPoint[below left](B){B}
    \tkzLabelPoint[below right](C){$\varGamma $}
    \tkzLabelPoint[above](D){$\varDelta  $}
    \tkzLabelPoint[below](E){E}
    \tkzLabelPoint[below](Z){Z}
    \tkzLabelPoint[above](M){M}
\end{tikzpicture}
\end{wrapfigure}
Δίνεται παραλληλόγραμμο $AB\varGamma \Delta$ με τη γωνία του $\hat{B}$ να είναι ίση με $70^\circ$ και το ύψος του $ AE$. Έστω $\mathrm{H}$ σημείο της $B\varGamma $ ώστε $BE = E\mathrm{H}$.
\begin{alist}
    \item Να αποδείξετε ότι το τετράπλευρο $A H\varGamma \varDelta  $ είναι ισοσκελές τραπέζιο. \monades{8}
    \item Να υπολογίσετε τις γωνίες του τραπεζίου $A H\varGamma \varDelta  $. \monades{9}
    \item Αν $ M$ το μέσο της $B\varDelta  $, να αποδείξετε ότι $ E M = \frac{A\varGamma }{2}$. \monades{8}
\end{alist}

\vspace{4mm}

% --- Άσκηση 1791 (1791.tex) ---
\Askhsh{1791}{4}
\begin{ylh}{1}
    \item 3.2. 1ο Κριτήριο ισότητας τριγώνων
    \item 4.2. Τέμνουσα δύο ευθειών - Ευκλείδειο αίτημα
    \item 5.2. Παραλληλόγραμμα
    \item 5.6. Εφαρμογές στα τρίγωνα
    \item 5.11. Ισοσκελές τραπέζιο.
\end{ylh}
\begin{wrapfigure}[5]{r}{3cm}
\begin{tikzpicture}[scale=0.65]
    \tkzDefPoint(0,0){A}
    \tkzDefPoint(3,0){B}
    \tkzDefPoint(0,{3*sqrt(3)}){C}
    \tkzDefMidPoint(B,C) \tkzGetPoint{M}
    \tkzDefPointBy[projection=onto B--C](A) \tkzGetPoint{D}
    \tkzDefLine[orthogonal=through C](A,M) \tkzGetPoint{c_perp}
    \tkzInterLL(C,c_perp)(A,M) \tkzGetPoint{E}
    \tkzDrawPolygon[l3](A,B,C)
    \tkzDrawSegments[l4,maincolor](A,D A,E C,E D,E)
    \tkzMarkRightAngle[size=0.3](B,A,C)
    \tkzDrawPoints(A,B,C,M,D,E)
    \tkzLabelPoint[below left](A){A}
    \tkzLabelPoint[below right](B){B}
    \tkzLabelPoint[above left](C){$\varGamma $}
    \tkzLabelPoint[right=2pt](M){M}
    \tkzLabelPoint[right](D){$\varDelta  $}
    \tkzLabelPoint[above right](E){E}
\end{tikzpicture}
\end{wrapfigure}
Δίνεται ορθογώνιο τρίγωνο $AB\varGamma $ με $\hat{ A} = 90^\circ$ και $\hat{\varGamma } = 30^\circ$. Φέρνουμε το ύψος του $ A\Delta$ και τη διάμεσό του $A\mathrm{M}$. Από το $\varGamma $ φέρνουμε κάθετη στην ευθεία $A\mathrm{M}$, η οποία την τέμνει στο $E$.
Να αποδείξετε ότι:
\begin{alist}
    \item Το τρίγωνο $A\mathrm{M}B$ είναι ισόπλευρο. \monades{8}
    \item $ M E =  M\varDelta   = \frac{B\varGamma }{4}$. \monades{9}
    \item Το $A\Delta E\varGamma $ είναι ισοσκελές τραπέζιο. \monades{8}
\end{alist}

\vspace{4mm}

% --- Άσκηση 1792 (1792.tex) ---
\Askhsh{1792}{4}
\begin{ylh}{1}
    \item 3.2. 1ο Κριτήριο ισότητας τριγώνων
    \item 4.6. Άθροισμα γωνιών τριγώνου.
\end{ylh}
Δίνεται τρίγωνο $AB\varGamma $ με $AB < A\varGamma $. Φέρνουμε τη διχοτόμο του $ A K$ και σε τυχαίο σημείο της $E$ φέρνουμε ευθεία κάθετη στη διχοτόμο $ A\mathrm{K}$, η οποία τέμνει τις $AB$ και $ A\varGamma $ στα σημεία $\mathrm{Z}$ και $\Delta$ αντίστοιχα και την προέκταση της $\varGamma B$ στο σημείο $\mathrm{H}$.
Να αποδείξετε ότι:
\begin{alist}
    \item $ Z\hat{\varDelta  }\varGamma  = 90^\circ + \frac{\hat{A}}{2}$. \monades{7}
    \item $ Z K =  K\varDelta  $. \monades{8}
    \item $ Z\hat{ H}\varGamma  = \frac{\hat{B} - \hat{\varGamma }}{2}$. \monades{10}
\end{alist}
\begin{center}
\begin{tikzpicture}[scale=0.8]
    \tkzDefPoint(2,0){B}
    \tkzDefPoint(10,0){C}
    \tkzDefPoint(3,4){A}
    \tkzDefLine[bisector](B,A,C) \tkzGetPoint{k_dir}
    \tkzInterLL(A,k_dir)(B,C) \tkzGetPoint{K}
    \tkzDefPointBy[homothety=center K ratio 0.45](A) \tkzGetPoint{E}
    \tkzDefLine[orthogonal=through E](A,K) \tkzGetPoint{e_perp}
    \tkzInterLL(E,e_perp)(A,B) \tkzGetPoint{Z}
    \tkzInterLL(E,e_perp)(A,C) \tkzGetPoint{D}
    \tkzInterLL(E,e_perp)(B,C) \tkzGetPoint{H}
    \tkzDrawPolygon[l3](A,B,C)
    \tkzDrawSegments[l4,maincolor](A,K H,D H,B)
    \tkzMarkRightAngle[size=0.3](A,E,Z)
    \tkzDrawPoints(A,B,C,K,E,Z,D,H)
    \tkzLabelPoint[above](A){A}
    \tkzLabelPoint[below](B){B}
    \tkzLabelPoint[right](C){$\varGamma $}
    \tkzLabelPoint[below](K){K}
    \tkzLabelPoint[below right=3pt](E){E}
    \tkzLabelPoint[above left=2pt](Z){Z}
    \tkzLabelPoint[above right](D){$\varDelta  $}
    \tkzLabelPoint[left](H){H}
\end{tikzpicture}
\end{center}
\vspace{4mm}

% --- Άσκηση 1794 (1794.tex) ---
\Askhsh{1794}{4}
\begin{ylh}{1}
    \item 5.3. Ορθογώνιο
    \item 5.4. Ρόμβος
    \item 5.6. Εφαρμογές στα τρίγωνα.
\end{ylh}
\begin{alist}
    \item Σε ορθογώνιο $AB\varGamma \Delta$ θεωρούμε $ K, \Lambda, \mathrm{M}, \mathrm{N}$ τα μέσα των πλευρών του $ $AB$, B\varGamma , \varGamma \Delta, \Delta A$ αντίστοιχα. Να αποδείξετε ότι το τετράπλευρο $\mathrm{K}\Lambda\mathrm{M}\mathrm{N}$ είναι ρόμβος. \monades{15}
    \item Σε ένα τετράπλευρο $AB\varGamma \varDelta  $ τα μέσα $ K, \Lambda,  M,  N$ των πλευρών του $ $AB$, B\varGamma , \varGamma \varDelta  , \varDelta   A$ αντίστοιχα είναι κορυφές ρόμβου. Το τετράπλευρο $ AB\varGamma \varDelta  $, πρέπει να είναι απαραίτητα ορθογώνιο; Να τεκμηριώσετε τη θετική ή αρνητική σας απάντηση. \monades{10}
\end{alist}

\vspace{4mm}

% --- Άσκηση 1795 (1795.tex) ---
\Askhsh{1795}{4}
\begin{ylh}{1}
    \item 3.2. 1ο Κριτήριο ισότητας τριγώνων
    \item 4.2. Τέμνουσα δύο ευθειών - Ευκλείδειο αίτημα
    \item 4.6. Άθροισμα γωνιών τριγώνου
    \item 5.2. Παραλληλόγραμμα
    \item 5.5. Τετράγωνο.
\end{ylh}
Εκτός τριγώνου $AB\varGamma $ κατασκευάζουμε τετράγωνα $ $AB$\Delta E$ και $A\varGamma \mathrm{Z}\mathrm{H}$. Αν $\mathrm{M}$ το μέσο του $B\varGamma $ και $\Lambda$ σημείο στην προέκταση της $ A\mathrm{M}$ τέτοιο ώστε $A\mathrm{M} = \mathrm{M}\Lambda$, να αποδείξετε ότι:
\begin{alist}
    \item $\varGamma \Lambda = A E$. \monades{10}
    \item Οι γωνίες $A\hat{\varGamma }\Lambda$ και $ E\hat{A} H$ είναι ίσες. \monades{10}
    \item Η προέκταση της $ MA$ (προς το $ A$) τέμνει κάθετα την $ E H$. \monades{5}
\end{alist}
\begin{center}
\begin{tikzpicture}[scale=0.85]
    \tkzDefPoint(2,0){B}
    \tkzDefPoint(7,0){C}
    \tkzDefPoint(3.5,3.5){A}
    \tkzDefPointBy[rotation=center A angle -90](B) \tkzGetPoint{E}
    \tkzDefPointBy[rotation=center B angle 90](A) \tkzGetPoint{D}
    \tkzDefPointBy[rotation=center A angle 90](C) \tkzGetPoint{H}
    \tkzDefPointBy[rotation=center C angle -90](A) \tkzGetPoint{Z}
    \tkzDefMidPoint(B,C) \tkzGetPoint{M}
    \tkzDefPointBy[symmetry=center M](A) \tkzGetPoint{L}
    \tkzInterLL(A,M)(E,H) \tkzGetPoint{P}
    \tkzDrawPolygon[l3](A,B,C)
    \tkzDrawPolygon[l3](A,B,D,E)
    \tkzDrawPolygon[l3](A,C,Z,H)
    \tkzDrawSegments[l4,maincolor](P,L B,L C,L E,H)
    \tkzDrawPoints(A,B,C,D,E,Z,H,M,L,P)
    \tkzLabelPoint[right=3pt](A){A}
    \tkzLabelPoint[below left](B){B}
    \tkzLabelPoint[below right](C){$\varGamma $}
    \tkzLabelPoint[left](D){$\varDelta  $}
    \tkzLabelPoint[above left](E){E}
    \tkzLabelPoint[right](Z){Z}
    \tkzLabelPoint[above right](H){H}
    \tkzLabelPoint[above right=2pt](M){M}
    \tkzLabelPoint[below](L){$\Lambda$}
    \tkzLabelPoint[above left](P){P}
\end{tikzpicture}
\end{center}
\vspace{4mm}

% --- Άσκηση 1796 (1796.tex) ---
\Askhsh{1796}{4}
\begin{ylh}{1}
    \item 3.2. 1ο Κριτήριο ισότητας τριγώνων
    \item 3.14. Σχετικές θέσεις ευθείας και κύκλου
    \item 3.15. Εφαπτόμενα τμήματα
    \item 5.4. Ρόμβος
    \item 5.9. Μια ιδιότητα του ορθογώνιου τριγώνου.
\end{ylh}
\begin{wrapfigure}[5]{r}{7cm}
\begin{tikzpicture}[scale=0.7]
    \tkzDefPoint(0,0){O}
    \tkzDefPoint(4,0){K}
    \tkzDefPoint(2,0){E}
    \tkzDrawCircle[l3](O,E)
    \tkzDrawCircle[l3](K,E)
    \tkzInterCC(E,K)(K,E) \tkzGetPoints{A}{B}
    \tkzDefPoint(-2.5,0){L_left}
    \tkzDefPoint(6.5,0){L_right}
    \tkzDrawSegment[l3](L_left,L_right)
    \tkzDrawSegments[l3](O,A O,B K,A K,B A,E B,E)
    \tkzDrawPoints(O,K,A,B)
    \tkzDrawPoint(E)
    \tkzLabelPoint[above left](O){O}
    \tkzLabelPoint[above right](K){K}
    \tkzLabelPoint[above left](E){E}
    \tkzLabelPoint[above](A){A}
    \tkzLabelPoint[below](B){B}
\end{tikzpicture}
\end{wrapfigure}
Δυο ίσοι κύκλοι $( O, \rho)$ και $( K, \rho)$ εφάπτονται εξωτερικά στο σημείο $ E$. Αν $ OA$ και $ OB$ είναι τα εφαπτόμενα τμήματα από το σημείο $\mathrm{O}$ στον κύκλο $( K, \rho)$ να αποδείξετε ότι:
\begin{alist}
    \item $A E = B E$. \monades{9}
    \item $A\hat{ O} K = 30^\circ$. \monades{8}
    \item Το τετράπλευρο $A $KB$ E$ είναι ρόμβος. \monades{8}
\end{alist}

\vspace{4mm}

% --- Άσκηση 1797 (1797.tex) ---
\Askhsh{1797}{4}
\begin{ylh}{1}
    \item 5.4. Ρόμβος
    \item 5.6. Εφαρμογές στα τρίγωνα
    \item 5.11. Ισοσκελές τραπέζιο.
\end{ylh}
\begin{alist}
    \item Σε ισοσκελές τραπέζιο $AB\varGamma \Delta$ θεωρούμε $ K, \Lambda, \mathrm{M}, \mathrm{N}$ τα μέσα των πλευρών του $ $AB$, B\varGamma , \varGamma \Delta, \Delta A$ αντίστοιχα. Να αποδείξετε ότι το τετράπλευρο $\mathrm{K}\Lambda\mathrm{M}\mathrm{N}$ είναι ρόμβος. \monades{13}
    \item Σε ένα τετράπλευρο $AB\varGamma \varDelta  $ τα μέσα $ K, \Lambda,  M,  N$ των πλευρών του $ $AB$, B\varGamma , \varGamma \varDelta  , \varDelta   A$ αντίστοιχα είναι κορυφές ρόμβου. Για να σχηματίζεται ρόμβος το $ AB\varGamma \varDelta  $ πρέπει να είναι ισοσκελές τραπέζιο; Να αιτιολογήσετε πλήρως τη θετική ή αρνητική απάντησή σας. \monades{12}
\end{alist}

\vspace{4mm}

% --- Άσκηση 1798 (1798.tex) ---
\Askhsh{1798}{4}
\begin{ylh}{1}
    \item 5.2. Παραλληλόγραμμα
    \item 5.3. Ορθογώνιο
    \item 5.4. Ρόμβος
    \item 5.6. Εφαρμογές στα τρίγωνα.
\end{ylh}
\begin{alist}
    \item Σε ρόμβο $AB\varGamma \Delta$ θεωρούμε $ K, \Lambda, \mathrm{M}, \mathrm{N}$ τα μέσα των πλευρών του $ $AB$, B\varGamma , \varGamma \Delta, \Delta A$ αντίστοιχα. Να αποδείξετε ότι το τετράπλευρο $\mathrm{K}\Lambda\mathrm{M}\mathrm{N}$ είναι ορθογώνιο. \monades{13}
    \item Να αποδείξετε ότι τα μέσα των πλευρών ενός ορθογωνίου είναι κορυφές ρόμβου. \monades{12}
\end{alist}

\vspace{4mm}

% --- Άσκηση 1800 (1800.tex) ---
\Askhsh{1800}{4}
\begin{ylh}{1}
    \item 3.2. 1ο Κριτήριο ισότητας τριγώνων
    \item 3.6. Κριτήρια ισότητας ορθογώνιων τριγώνων
    \item 4.2. Τέμνουσα δύο ευθειών - Ευκλείδειο αίτημα
    \item 5.3. Ορθογώνιο.
\end{ylh}
\begin{wrapfigure}[5]{r}{5cm}
\begin{tikzpicture}[scale=0.7]
    \tkzDefPoint(0,0){B}
    \tkzDefPoint(6,0){C} % \varGamma 
    \tkzDefPoint(3,5.5){A}
    \tkzDefPoint(2.5,0){M}
    \tkzDefPointBy[projection=onto A--C](B) \tkzGetPoint{H}
    \tkzDefPointBy[projection=onto A--B](M) \tkzGetPoint{D} % \Delta
    \tkzDefPointBy[projection=onto A--C](M) \tkzGetPoint{E}
    \tkzDefPointBy[projection=onto B--H](M) \tkzGetPoint{Th}
    \tkzDrawPolygon[l3](A,B,C)
    \tkzDrawSegments[l4,maincolor](B,H M,D M,E M,Th)
    \tkzMarkRightAngle[size=0.3](B,H,A)
    \tkzMarkRightAngle[size=0.3](A,D,M)
    \tkzMarkRightAngle[size=0.3](M,E,C)
    \tkzMarkRightAngle[size=0.3](M,Th,H)
    \tkzDrawPoints(A,B,C,M,H,D,E,Th)
    \tkzLabelPoint[above](A){A}
    \tkzLabelPoint[below left](B){B}
    \tkzLabelPoint[below right](C){$\varGamma $}
    \tkzLabelPoint[below](M){M}
    \tkzLabelPoint[above right](H){H}
    \tkzLabelPoint[above left](D){$\varDelta  $}
    \tkzLabelPoint[right](E){E}
    \tkzLabelPoint[above left=1pt](Th){$\varTheta $}
\end{tikzpicture}
\end{wrapfigure}
Δίνεται ισοσκελές τρίγωνο $AB\varGamma $ με $AB = A\varGamma $, τυχαίο σημείο $\mathrm{M}$ της βάσης του $B\varGamma $ και το ύψος του $B\mathrm{H}$. Από το $\mathrm{M}$ φέρνουμε κάθετες $\mathrm{M}\Delta, \mathrm{M}E$ και $\mathrm{M}\Theta$ στις $ AB,  A\varGamma $ και $B\mathrm{H}$ αντίστοιχα.
Να αποδείξετε ότι:
\begin{alist}
    \item Το τετράπλευρο $ M E H\varTheta $ είναι ορθογώνιο. \monades{9}
    \item $B\varTheta  = \varDelta   M$. \monades{9}
    \item Το άθροισμα $ M\varDelta   +  M E = B H$. \monades{7}
\end{alist}

\vspace{4mm}

% --- Άσκηση 1801 (1801.tex) ---
\Askhsh{1801}{4}
\begin{ylh}{1}
    \item 4.2. Τέμνουσα δύο ευθειών - Ευκλείδειο αίτημα
    \item 5.2. Παραλληλόγραμμα
    \item 5.6. Εφαρμογές στα τρίγωνα.
\end{ylh}
\begin{wrapfigure}[5]{r}{5cm}
\begin{tikzpicture}[scale=0.7]
    \tkzDefPoint(0,0){B}
    \tkzDefPoint(8,0){C} % \varGamma 
    \tkzDefPoint(2,4.5){A}
    \tkzDefMidPoint(B,C) \tkzGetPoint{D} % \Delta
    \tkzDefMidPoint(A,C) \tkzGetPoint{E}
    \tkzDefMidPoint(A,B) \tkzGetPoint{Z}
    \tkzDefLine[bisector](C,B,A) \tkzGetPoint{b_bis}
    \tkzInterLL(Z,E)(B,b_bis) \tkzGetPoint{M}
    \tkzInterLL(D,E)(B,b_bis) \tkzGetPoint{N}
    \tkzDrawPolygon[l3](A,B,C)
    \tkzDrawSegments[l4,maincolor](Z,E D,N B,N)
    \tkzDrawPoints(A,B,C,D,E,Z,M,N)
    \tkzLabelPoint[above](A){A}
    \tkzLabelPoint[below left](B){B}
    \tkzLabelPoint[below right](C){$\varGamma $}
    \tkzLabelPoint[below](D){$\varDelta  $}
    \tkzLabelPoint[right=2pt](E){E}
    \tkzLabelPoint[left](Z){Z}
    \tkzLabelPoint[above left](M){M}
    \tkzLabelPoint[above](N){N}
\end{tikzpicture}
\end{wrapfigure}
Δίνεται τρίγωνο $AB\varGamma $ με $A\varGamma  > AB$ και $\varDelta  ,  E,  Z$ τα μέσα των πλευρών του $B\varGamma ,  A\varGamma , AB$ αντίστοιχα. Αν η διχοτόμος της γωνίας $\hat{B}$ τέμνει την $ Z E$ στο σημείο $ M$ και την προέκταση της $\varDelta  E$ στο σημείο $ N$, να αποδείξετε ότι:
\begin{alist}
    \item Το τετράπλευρο $ Z E\varDelta   B$ είναι παραλληλόγραμμο. \monades{7}
    \item Τα τρίγωνα $B Z M$ και $ M E N$ είναι ισοσκελή. \monades{10}
    \item $B Z + \mathrm{N} E = \varDelta  \varGamma $. \monades{8}
\end{alist}

\vspace{4mm}

% --- Άσκηση 1802 (1802.tex) ---
\Askhsh{1802}{4}
\begin{ylh}{1}
    \item 4.6. Άθροισμα γωνιών τριγώνου
    \item 5.6. Εφαρμογές στα τρίγωνα.
\end{ylh}
\begin{wrapfigure}[5]{r}{5cm}
\begin{tikzpicture}[scale=0.7]
    \tkzDefPoint(0,0){B}
    \tkzDefPoint(7,0){C} % \varGamma 
    \tkzDefPoint(2.5,4){A}
    \tkzDefMidPoint(B,C) \tkzGetPoint{M}
    \tkzDefMidPoint(A,M) \tkzGetPoint{K}
    \tkzInterLL(B,K)(A,C) \tkzGetPoint{N}
    \tkzDefMidPoint(C,N) \tkzGetPoint{L}
    \tkzDrawPolygon[l3](A,B,C)
    \tkzDrawSegments[l4,maincolor](A,M B,N)
    \tkzDrawPoints(A,B,C,M,K,N,L)
    \tkzLabelPoint[above](A){A}
    \tkzLabelPoint[below left](B){B}
    \tkzLabelPoint[below right](C){$\varGamma $}
    \tkzLabelPoint[below](M){M}
    \tkzLabelPoint[left=2pt](K){K}
    \tkzLabelPoint[right=2pt](N){N}
    \tkzLabelPoint[right=2pt](L){$\Lambda$}
\end{tikzpicture}
\end{wrapfigure}
Δίνεται τρίγωνο $AB\varGamma $, $A M$ διάμεσός του και $ K$ το μέσο του $A M$. Αν η προέκταση της $B K$ τέμνει την $A\varGamma $ στο σημείο $ N$, και $\Lambda$ είναι το μέσο του $\varGamma  N$, να αποδείξετε ότι:
\begin{alist}
    \item Το σημείο $\mathrm{N}$ είναι μέσο του $A\Lambda$. \monades{9}
    \item $ K\hat{ M}\varGamma  =  M\hat{B} K + A\hat{\mathrm{K}} N$. \monades{9}
    \item $B K = 3 K\mathrm{N}$. \monades{7}
\end{alist}

\vspace{4mm}

% --- Άσκηση 1803 (1803.tex) ---
\Askhsh{1803}{4}
\begin{ylh}{1}
    \item 3.2. 1ο Κριτήριο ισότητας τριγώνων
    \item 4.2. Τέμνουσα δύο ευθειών - Ευκλείδειο αίτημα
    \item 5.6. Εφαρμογές στα τρίγωνα.
\end{ylh}
\begin{wrapfigure}[5]{r}{4cm}
\begin{tikzpicture}[scale=0.8]
\tkzDefPoint(0,0){C}
\tkzDefPoint(6,0){B}
\tkzDefPoint(1.5,2.5){A} % Just an estimate to make it look decent.

% $BC$ = 2 $AC$. $AC$ = 3 if scale allows.
% Let's construct it more carefully:
\tkzDefPoint(0,0){C}
\tkzDefPoint(3,0){A}
% B is at distance 6 from C. Angle at C can be anything.
\tkzDefPointBy[rotation=center C angle 60](A) \tkzGetPoint{B_dir}
\tkzDefPointOnLine[pos=2](C,B_dir) \tkzGetPoint{B}

\tkzDefMidPoint(B,C) \tkzGetPoint{M}
\tkzDefMidPoint(M,C) \tkzGetPoint{K}
\tkzDefMidPoint(A,B) \tkzGetPoint{L}

\draw[l3] (A)--(B)--(C)--cycle;
\draw[l3] (A)--(M);
\draw[l3] (M)--(L);
\draw[l3] (A)--(K); % Not explicitly asked but often in these diagrams

\tkzDrawPoints(A,B,C,M,K,L)
\tkzLabelPoint[right](A){$A$}
\tkzLabelPoint[above](B){$B$}
\tkzLabelPoint[left](C){$\varGamma $}
\tkzLabelPoint[left](M){$M$}
\tkzLabelPoint[left](K){$K$}
\tkzLabelPoint[right](L){$\Lambda$}

\tkzMarkSegments[mark=s|](C,K K,M)
\tkzMarkSegments[mark=s|](M,B) % $BC$=2AC => $CM$=$MB$=$AC$
\tkzMarkSegments[mark=s|](A,C)
\tkzMarkSegments[mark=s||](A,L L,B)
\end{tikzpicture}
\end{wrapfigure}
Δίνεται τρίγωνο $AB\varGamma $ με $B\varGamma $ = 2ΑΓ. Έστω $AM$ διάμεσος του $AB\varGamma $ και Κ, Λ τα μέσα των $M\varGamma $ και $AB$ αντίστοιχα.
Να αποδείξετε~ ότι:
\begin{alist}
\item Μ$\hat{A}$Γ = A$\hat{Μ}\varGamma $ \monades{7}
\item $M\varLambda $ = $MK$. \monades{9}
\item Η $AM$ είναι~διχοτόμος της γωνίας $\varLambda  MK$. \monades{9}
\end{alist}

\vspace{4mm}

% --- Άσκηση 1804 (1804.tex) ---
\Askhsh{1804}{4}
\begin{ylh}{1}
    \item 4.2. Τέμνουσα δύο ευθειών - Ευκλείδειο αίτημα
    \item 5.6. Εφαρμογές στα τρίγωνα.
\end{ylh}
Δίνεται τετράπλευρο $AB\varGamma \varDelta  $ με $AB = \varGamma \varDelta  $ και $ M,  N,  K$ τα μέσα των $A\varDelta  , B\varGamma , B\varDelta  $ αντίστοιχα. Αν οι προεκτάσεις των $ AB$ και $\varDelta  \varGamma $ τέμνουν την προέκταση της $ M\mathrm{N}$ στα σημεία $ E$ και $ Z$ αντίστοιχα να αποδείξετε ότι:
\begin{alist}
    \item $\mathrm{M} K =  K\mathrm{N}$. \monades{13}
    \item $\mathrm{M}\hat{E}A = \mathrm{M}\hat{\mathrm{Z}}\Delta$. \monades{12}
\end{alist}

\vspace{4mm}

% --- Άσκηση 1805 (1805.tex) ---
\Askhsh{1805}{4}
\begin{ylh}{1}
    \item 3.2. 1ο Κριτήριο ισότητας τριγώνων
    \item 4.6. Άθροισμα γωνιών τριγώνου
    \item 5.2. Παραλληλόγραμμα.
\end{ylh}
Δίνεται παραλληλόγραμμο $AB\varGamma \varDelta  $ και στην προέκταση της $ A\varDelta  $ θεωρούμε σημείο $ E$ τέτοιο ώστε $\varDelta  E = \varDelta  \varGamma $ ενώ στην προέκταση της $ AB$ θεωρούμε σημείο $ H$ τέτοιο ώστε $B\mathrm{H} = B\varGamma $.
\begin{alist}
    \item Να αποδείξετε ότι:
    \begin{rlist}
        \item $B\hat{\varGamma } H = \varDelta  \hat{\varGamma } E$. \monades{10}
        \item τα σημεία $ H, \varGamma ,  E$ είναι συνευθειακά. \monades{10}
    \end{rlist}
    \item Ένας μαθητής για να αποδείξει ότι τα σημεία $ H, \varGamma ,  E$ είναι συνευθειακά ανέπτυξε τον παρακάτω συλλογισμό. « Έχουμε:

$B\hat{\varGamma } H = \varDelta  \hat{ E}\varGamma $ ($\omega\varsigma$ εντός εκτός και επί τα αυτά μέρη των παραλλήλων $\varDelta  E$ και $B\varGamma $ που τέμνονται από τη $\mathrm{H} E$) και

$B\hat{\varGamma }\varDelta   = \varGamma \hat{\varDelta  } E$ ($\omega\varsigma$ εντός εναλλάξ των παραλλήλων $\varDelta  E$ και $B\varGamma $ που τέμνονται από την $\varDelta  \varGamma $).

Όμως $\varDelta  \hat{\varGamma } E + \varGamma \hat{\varDelta  } E + \varDelta  \hat{ E}\varGamma  = 180^\circ$ ($\omega\varsigma$ άθροισμα των γωνιών του τριγώνου $\varDelta  E\varGamma $). Άρα σύμφωνα με τα προηγούμενα: $\varDelta  \hat{\varGamma }E + B\hat{\varGamma }\varDelta   + B\hat{\varGamma } H = 180^\circ$. Οπότε τα σημεία $\mathrm{H}, \varGamma ,  E$ είναι συνευθειακά.»

Όμως ο καθηγητής υπέδειξε ένα λάθος στο συλλογισμό αυτό. Να βρείτε το λάθος στο συγκεκριμένο συλλογισμό. \monades{5}
\end{alist}

\vspace{4mm}

% --- Άσκηση 1806 (1806.tex) ---
\Askhsh{1806}{4}
\begin{ylh}{1}
    \item 3.2. 1ο Κριτήριο ισότητας τριγώνων
    \item 4.6. Άθροισμα γωνιών τριγώνου
    \item 5.9. Μια ιδιότητα του ορθογώνιου τριγώνου.
\end{ylh}
Δίνεται ορθογώνιο τρίγωνο $AB\varGamma $ με τη γωνία $\hat{ A}$ ορθή. Φέρνουμε τη διάμεσό του $ A M$ και σε τυχαίο σημείο $ K$ αυτής φέρνουμε κάθετη στην $A M$ η οποία τέμνει τις $AB$ και $A\varGamma $ στα σημεία $\Delta$ και $E$ αντίστοιχα. Αν $\mathrm{H}$ είναι το μέσο του $\Delta E$ να αποδείξετε ότι:
\begin{alist}
    \item $\hat{B} = B\hat{A} M$. \monades{8}
    \item $A\hat{\varDelta  } H = \varDelta  \hat{A} H$. \monades{9}
    \item Η ευθεία $A H$ τέμνει κάθετα τη $B\varGamma $. \monades{8}
\end{alist}

\vspace{4mm}

% --- Άσκηση 1808 (1808.tex) ---
\Askhsh{1808}{4}
\begin{ylh}{1}
    \item 3.2. 1ο Κριτήριο ισότητας τριγώνων
    \item 3.6. Κριτήρια ισότητας ορθογώνιων τριγώνων
    \item 5.9. Μια ιδιότητα του ορθογώνιου τριγώνου.
\end{ylh}
Δίνεται ισοσκελές τρίγωνο $AB\varGamma $ με $AB = A\varGamma $ και $\varDelta  ,  E$ τα μέσα των πλευρών του $AB$ και $ A\varGamma $ αντίστοιχα. Στην προέκταση της $\varDelta  E$ (προς το $E$) θεωρούμε σημείο $\Lambda$ ώστε $ E\Lambda = A E$ και στην προέκταση της $E\varDelta  $ (προς το $\varDelta  $) θεωρούμε σημείο $ K$ τέτοιο ώστε $\varDelta   K = A\varDelta  $.
Να αποδείξετε ότι:
\begin{alist}
    \item $ K\varDelta   = \Lambda E$. \monades{6}
    \item Τα τρίγωνα $A KB$ και $A\Lambda\varGamma $ είναι ορθογώνια. \monades{9}
    \item Τα τρίγωνα $A KB$ και $A\Lambda\varGamma $ είναι ίσα. \monades{10}
\end{alist}

\vspace{4mm}

% --- Άσκηση 1810 (1810.tex) ---
\Askhsh{1810}{4}
\begin{ylh}{1}
    \item 4.2. Τέμνουσα δύο ευθειών - Ευκλείδειο αίτημα
    \item 5.2. Παραλληλόγραμμα.
\end{ylh}
Δίνεται τρίγωνο $AB\varGamma $. Από το μέσο $ M$ του $B\varGamma $ φέρνουμε ευθύγραμμο τμήμα $ M\varDelta  $ ίσο και παράλληλο με το $BA$ και ευθύγραμμο τμήμα $ M E$ ίσο και παράλληλο με το $\varGamma  A$ (τα σημεία $\varDelta  $ και $ E$ είναι στο ημιεπίπεδο που ορίζεται από τη $B\varGamma $ και το σημείο $A$). Να αποδείξετε ότι:
\begin{alist}
    \item Τα σημεία $\varDelta  , A,  E$ είναι συνευθειακά. \monades{10}
    \item Η περίμετρος του τριγώνου $ M\Delta E$ είναι ίση με την περίμετρο του τριγώνου $AB\varGamma $. \monades{9}
    \item Όταν ένας καθηγητής έθεσε στους μαθητές του το ερώτημα αν τα σημεία $\varDelta  , A,  E$ είναι συνευθειακά, ένας από αυτούς έκανε το παρακάτω σχήμα και απάντησε ως εξής:
    
    $\hat{\mathrm{Z}}_1 = \hat{A}_1$ (εντός εναλλάξ των $ $AB$ \parallel  M\Delta$ που τέμνονται από την $ A\mathrm{Z}$)
    $A\hat{\varDelta  } Z = \hat{A}_2$ (εντός εκτός και επί τα αυτά μέρη των $ $AB$ \parallel  M\varDelta  $ που τέμνονται από την $\varDelta   E$)
    Όμως $\hat{ Z}_1 + \hat{A}_3 + A\hat{\varDelta  } Z = 180^\circ$ (άθροισμα γωνιών του τριγώνου $ A\varDelta   Z$). Άρα σύμφωνα με τα προηγούμενα έχουμε: $\hat{ A}_1 + \hat{A}_2 + \hat{A}_3 = 180^\circ$. Οπότε $\varDelta  , A,  E$ συνευθειακά.
    
    Όμως ο καθηγητής είπε ότι υπάρχει λάθος στο συλλογισμό. Μπορείτε να εντοπίσετε το λάθος του μαθητή; \monades{6}
\end{alist}

\vspace{4mm}

% --- Άσκηση 1811 (1811.tex) ---
\Askhsh{1811}{4}
\begin{ylh}{1}
    \item 3.2. 1ο Κριτήριο ισότητας τριγώνων
    \item 4.2. Τέμνουσα δύο ευθειών - Ευκλείδειο αίτημα
    \item 4.6. Άθροισμα γωνιών τριγώνου
    \item 5.9. Μια ιδιότητα του ορθογώνιου τριγώνου.
\end{ylh}
Δίνονται δυο παράλληλες ευθείες $(\varepsilon)$ και $(\eta)$, και μια τρίτη που τις τέμνει στα σημεία $A$ και $B$ αντίστοιχα. Θεωρούμε τις διχοτόμους των εντός και επί τα αυτά μέρη γωνιών που σχηματίζονται, οι οποίες τέμνονται σε σημείο $\varDelta  $. Αν $ M$ είναι το μέσον του $ AB$, να αποδείξετε ότι:
\begin{alist}
    \item Η γωνία $B\hat{\varDelta  }A$ είναι ορθή. \monades{9}
    \item $B\hat{ M}\Delta = 2 \cdot  M\hat{\Delta}A$. \monades{8}
    \item $ M\varDelta   \parallel \varepsilon$. \monades{8}
\end{alist}

\vspace{4mm}

% --- Άσκηση 1812 (1812.tex) ---
\Askhsh{1812}{4}
\begin{ylh}{1}
    \item 5.9. Μια ιδιότητα του ορθογώνιου τριγώνου.
\end{ylh}
% ===================================================================
%  Αρχείο: 1812.tex (Fragment)
%  Αυτόματη μετατροπή μέσω doc2latex.py
% ===================================================================



Δίνεται ορθογώνιο τρίγωνο $AB\varGamma $ με τη γωνία $A$ ορθή και $M$ τυχαίο σημείο της πλευράς $B\varGamma $. Φέρουμε τις διχοτόμους γωνιών $BMA$ και $AM\varGamma $ οι οποίες τέμνουν τις $AB$ και $A\varGamma $ στα σημεία $\varDelta  $ και $E$ αντίστοιχα.

\begin{alist}
\item Να αποδείξετε ότι, η γωνία $\varDelta  \hat{M}E$ είναι ορθή. \monades{12}

\item Αν Κ το μέσο του $\varDelta   E$, να αποδείξετε ότι $MK$ = $KA$ \monades{13}


\begin{center}
\begin{tikzpicture}[scale=1]
\tkzDefPoint(0,4){B}
\tkzDefPoint(0,0){A}
\tkzDefPoint(6,0){C}
\draw[l3] (A)--(B)--(C)--cycle;

% M random on $BC$
\tkzDefPointOnLine[pos=0.7](B,C) \tkzGetPoint{M}

% Bisectors of $BMA$ and $AMC$
\tkzDefLine[bisector](B,M,A) \tkzGetPoint{bis1}
\tkzDefLine[bisector](A,M,C) \tkzGetPoint{bis2}

\tkzInterLL(M,bis1)(A,B) \tkzGetPoint{D}
\tkzInterLL(M,bis2)(A,C) \tkzGetPoint{E}

\tkzDefMidPoint(D,E) \tkzGetPoint{K}

\draw[l3] (A)--(M);
\draw[l3] (M)--(D);
\draw[l3] (M)--(E);
\draw[l3] (D)--(E);
\draw[l3, dashed] (A)--(K);
\draw[l3, dashed] (M)--(K);

\tkzDrawPoints(A,B,C,M,D,E,K)
\tkzLabelPoint[below left](A){$A$}
\tkzLabelPoint[above left](B){$B$}
\tkzLabelPoint[below right](C){$\varGamma $}
\tkzLabelPoint[above right](M){$M$}
\tkzLabelPoint[left](D){$\varDelta  $}
\tkzLabelPoint[below](E){$E$}
\tkzLabelPoint[above](K){$K$}

\tkzMarkRightAngle(C,A,B)
\tkzMarkRightAngle(D,M,E)
\end{tikzpicture}
\end{center}
\end{alist}

\vspace{4mm}

% --- Άσκηση 1814 (1814.tex) ---
\Askhsh{1814}{4}
\begin{ylh}{1}
    \item 3.2. 1ο Κριτήριο ισότητας τριγώνων
    \item 5.5. Τετράγωνο.
\end{ylh}
Δίνεται τετράγωνο $AB\varGamma \varDelta  $ και εντός αυτού ισόπλευρο τρίγωνο $ $MB$\varGamma $. Αν η προέκταση της $ A M$ τέμνει την $B\varDelta  $ στο σημείο $ E$, να αποδείξετε ότι:
\begin{alist}
    \item $\varDelta  \hat{A} E = 15^\circ$. \monades{8}
    \item Τα τρίγωνα $\varDelta   A E$ και $\varDelta  E\varGamma $ είναι ίσα. \monades{8}
    \item Η $\varGamma E$ είναι διχοτόμος της γωνίας $\Delta\hat{\varGamma } M$. \monades{9}
\end{alist}

\vspace{4mm}

% --- Άσκηση 1815 (1815.tex) ---
\Askhsh{1815}{4}
\begin{ylh}{1}
    \item 3.11. Ανισοτικές σχέσεις πλευρών και γωνιών
    \item 4.2. Τέμνουσα δύο ευθειών - Ευκλείδειο αίτημα
    \item 5.10. Τραπέζιο.
\end{ylh}
Δίνεται τραπέζιο $AB\varGamma \varDelta  $ με $AB \parallel \varGamma \varDelta  $ και $AB =  A\varDelta   + B\varGamma $. Αν η διχοτόμος της γωνίας $\hat{\varDelta  }$ τέμνει την $AB$ στο σημείο $ M$, να αποδείξετε ότι:
\begin{alist}
    \item Το τρίγωνο $A\Delta M$ είναι ισοσκελές. \monades{8}
    \item Το τρίγωνο $ $MB$\varGamma $ είναι ισοσκελές. \monades{9}
    \item Η $\varGamma  M$ είναι διχοτόμος της γωνίας $\hat{\varGamma }$ του τραπεζίου. \monades{8}
\end{alist}

\vspace{4mm}

% --- Άσκηση 1816 (1816.tex) ---
\Askhsh{1816}{4}
\begin{ylh}{1}
    \item 3.6. Κριτήρια ισότητας ορθογώνιων τριγώνων
    \item 4.2. Τέμνουσα δύο ευθειών - Ευκλείδειο αίτημα
    \item 5.3. Ορθογώνιο.
\end{ylh}
Δίνεται ισοσκελές τρίγωνο $AB\varGamma $ με $AB = A\varGamma $ και σημείο $\varDelta  $ στην προέκταση της $B\varGamma $. Από το $\varDelta  $ φέρνουμε $\varDelta   K$ κάθετη στην $AB$ και $\varDelta   E$ κάθετη στην προέκταση της $A\varGamma $. Από το σημείο $\varGamma $ φέρνουμε $\varGamma  H$ κάθετη στην $ AB$ και $\varGamma  Z$ κάθετη στην $ K\varDelta  $.
Να αποδείξετε ότι:
\begin{alist}
    \item H γωνία $ Z\hat{\varGamma }\varDelta  $ είναι ίση με τη γωνία $\hat{B}$. \monades{4}
    \item Η $\varGamma \varDelta  $ είναι διχοτόμος της γωνίας $ Z\hat{\varGamma } E$. \monades{4}
    \item Το τρίγωνο $\Delta Z E$ είναι ισοσκελές. \monades{9}
    \item $\varDelta   K - \varDelta   E =  H\varGamma $. \monades{8}
\end{alist}

\vspace{4mm}

% --- Άσκηση 1818 (1818.tex) ---
\Askhsh{1818}{4}
\begin{ylh}{1}
    \item 3.3. 2ο Κριτήριο ισότητας τριγώνων
    \item 3.6. Κριτήρια ισότητας ορθογώνιων τριγώνων
    \item 4.2. Τέμνουσα δύο ευθειών - Ευκλείδειο αίτημα.
\end{ylh}
Δίνεται τρίγωνο $AB\varGamma $ με $AB < A\varGamma $, η διχοτόμος του $A\varDelta  $ και ευθεία $(\varepsilon)$ παράλληλη από το $B$ προς την $ A\varGamma $. Από το μέσο $ M$ της $B\varGamma $ φέρνουμε ευθεία παράλληλη στην $A\varDelta  $ η οποία τέμνει την $ A\varGamma $ στο σημείο $ Z$, την ευθεία $(\varepsilon)$ στο σημείο $\Lambda$ και την προέκταση της $BA$ στο σημείο $ E$.
Να αποδείξετε ότι:
\begin{alist}
    \item Τα τρίγωνα $A E Z$ και $B\Lambda E$ είναι ισοσκελή. \monades{8}
    \item $B\Lambda = \varGamma  Z$. \monades{9}
    \item $A E = A\varGamma  - B\Lambda$. \monades{8}
\end{alist}

\vspace{4mm}

% --- Άσκηση 1819 (1819.tex) ---
\Askhsh{1819}{4}
\begin{ylh}{1}
    \item 3.4. 3ο Κριτήριο ισότητας τριγώνων
    \item 3.6. Κριτήρια ισότητας ορθογώνιων τριγώνων
    \item 4.6. Άθροισμα γωνιών τριγώνου.
\end{ylh}
Δίνεται ισόπλευρο τρίγωνο $AB\varGamma $ και στην προέκταση της $\varGamma  B$ (προς το $B$) θεωρούμε σημείο $\Delta$ τέτοιο ώστε $B\Delta = B\varGamma $, ενώ στην προέκταση της $B\varGamma $ (προς το $\varGamma $) θεωρούμε σημείο $E$ τέτοιο ώστε $\varGamma E = B\varGamma $. Φέρνουμε την κάθετη στην $E\Delta$ στο σημείο $E$, η οποία τέμνει την προέκταση της $\Delta A$ στο $\mathrm{Z}$.
\begin{alist}
    \item Να υπολογίσετε τις γωνίες των τριγώνων $\varGamma  A E$ και $B\varDelta   A$. \monades{8}
    \item Να αποδείξετε ότι η $\varGamma  Z$ είναι μεσοκάθετος του $A E$. \monades{12}
    \item Να αποδείξετε ότι $AB \parallel \varGamma  Z$. \monades{5}
\end{alist}

\vspace{4mm}

% --- Άσκηση 1820 (1820.tex) ---
\Askhsh{1820}{4}
\begin{ylh}{1}
    \item 5.2. Παραλληλόγραμμα
    \item 5.6. Εφαρμογές στα τρίγωνα
    \item 5.7. Βαρύκεντρο τριγώνου.
\end{ylh}
Δίνεται τρίγωνο $AB\varGamma $ και οι διάμεσοί του $ A\varDelta  , B E$ και $\varGamma  Z$. Προεκτείνουμε το τμήμα $ Z E$ (προς το $ E$) κατά τμήμα $ E H =  Z E$.
Να αποδείξετε ότι:
\begin{alist}
    \item Το τετράπλευρο $ E H\varDelta   B$ είναι παραλληλόγραμμο. \monades{8}
    \item Η περίμετρος του τριγώνου $A\Delta H$ είναι ίση με το άθροισμα των διαμέσων του τριγώνου $ $AB$\varGamma $. \monades{9}
    \item Οι ευθείες $B E$ και $\Delta H$ τριχοτομούν το τμήμα $ Z\varGamma $. \monades{8}
\end{alist}

\vspace{4mm}

% --- Άσκηση 1821 (1821.tex) ---
\Askhsh{1821}{4}
\begin{ylh}{1}
    \item 5.2. Παραλληλόγραμμα
    \item 5.4. Ρόμβος
    \item 5.9. Μια ιδιότητα του ορθογώνιου τριγώνου
    \item 5.10. Τραπέζιο.
\end{ylh}
Δίνεται ορθογώνιο τραπέζιο $AB\varGamma \varDelta  $ ($\hat{A} = \hat{\varDelta  } = 90^\circ$) με $B\varGamma  = \varGamma \varDelta   = 2AB$ και $ K, \Lambda$ τα μέσα των $B\varGamma $ και $\varGamma \varDelta  $. Η παράλληλη από το $ K$ προς την $ AB$ τέμνει την $A\Lambda$ στο $ Z$.
Να αποδείξετε ότι:
\begin{alist}
    \item Το $ Z$ είναι μέσο του $A\Lambda$. \monades{6}
    \item $B\varGamma  = 2 \Delta Z$. \monades{6}
    \item Το τετράπλευρο $ Z K\varGamma \Lambda$ είναι ρόμβος. \monades{5}
    \item $A\hat{ K}\Lambda = 90^\circ$. \monades{8}
\end{alist}

\vspace{4mm}

% --- Άσκηση 1823 (1823.tex) ---
\Askhsh{1823}{4}
\begin{ylh}{1}
    \item 3.6. Κριτήρια ισότητας ορθογώνιων τριγώνων
    \item 3.11. Ανισοτικές σχέσεις πλευρών και γωνιών
    \item 3.15. Εφαπτόμενα τμήματα
    \item 5.2. Παραλληλόγραμμα
    \item 5.4. Ρόμβος.
\end{ylh}
Δίνεται κύκλος κέντρου $ O$ και δυο μη αντιδιαμετρικά σημεία του $A$ και $B$. Φέρνουμε τις εφαπτόμενες του κύκλου στα σημεία $ A$ και $B$ οι οποίες τέμνονται στο σημείο $\varGamma $. Φέρνουμε επίσης και τα ύψη $ A\varDelta  $ και $B E$ του τριγώνου $AB\varGamma $ τα οποία τέμνονται στο σημείο $ H$.
Να αποδείξετε ότι:
\begin{alist}
    \item Το τρίγωνο $B HA$ είναι ισοσκελές. \monades{8}
    \item Το τετράπλευρο $ $OB$ HA$ είναι ρόμβος. \monades{9}
    \item Τα σημεία $ O,  H, \varGamma $ είναι συνευθειακά. \monades{8}
\end{alist}

\vspace{4mm}

% --- Άσκηση 1824 (1824.tex) ---
\Askhsh{1824}{4}
\begin{ylh}{1}
    \item 3.3. 2ο Κριτήριο ισότητας τριγώνων
    \item 5.6. Εφαρμογές στα τρίγωνα
    \item 5.9. Μια ιδιότητα του ορθογώνιου τριγώνου.
\end{ylh}
Δίνεται τρίγωνο $AB\varGamma $ και στην προέκταση της $\varGamma  B$ (προς το $B$) θεωρούμε σημείο $\Delta$ τέτοιο ώστε $B\Delta = AB$ ενώ στην προέκταση της $B\varGamma $ (προς το $\varGamma $) θεωρούμε σημείο $E$ τέτοιο ώστε $\varGamma E = \varGamma  A$.
Αν οι εξωτερικοί διχοτόμοι των γωνιών $\hat{B}$ και $\hat{\varGamma }$ τέμνουν τις $A\varDelta  $ και $A E$ στα σημεία $ K$ και $\Lambda$ αντίστοιχα, και η $ K\Lambda$ τέμνει τις $AB$ και $A\varGamma $ στα σημεία $ M$ και $ N$ αντίστοιχα, να αποδείξετε ότι:
\begin{alist}
    \item Τα σημεία $ K$ και $\Lambda$ είναι μέσα των $A\varDelta  $ και $A E$ αντίστοιχα. \monades{8}
    \item Τα τρίγωνα $ K MA$ και $A\mathrm{N}\Lambda$ είναι ισοσκελή. \monades{9}
    \item $ K\Lambda = \frac{AB + A\varGamma  + B\varGamma }{2}$. \monades{8}
\end{alist}

\vspace{4mm}

% --- Άσκηση 1825 (1825.tex) ---
\Askhsh{1825}{4}
\begin{ylh}{1}
    \item 3.6. Κριτήρια ισότητας ορθογώνιων τριγώνων
    \item 3.11. Ανισοτικές σχέσεις πλευρών και γωνιών
    \item 4.6. Άθροισμα γωνιών τριγώνου
    \item 5.5. Τετράγωνο.
\end{ylh}
Δίνεται τετράγωνο $AB\varGamma \varDelta  $ και τυχαίο σημείο $ E$ στην πλευρά $\varDelta  \varGamma $. Φέρνουμε τη διχοτόμο $ A Z$ της γωνίας $E\hat{A}B$ και τη $\varDelta   H$ κάθετη από το $\varDelta  $ προς την $A Z$, η οποία τέμνει την $AE$ στο $ M$ και την $AB$ στο $ N$.
Να αποδείξετε ότι:
\begin{alist}
    \item Τα τρίγωνα $A\varDelta  \mathrm{N}$ και $AB Z$ είναι ίσα. \monades{8}
    \item $A M = A\mathrm{N}$ και $\Delta E = E M$. \monades{10}
    \item $A E = \varDelta  E + B Z$. \monades{7}
\end{alist}

\vspace{4mm}

% --- Άσκηση 1827 (1827.tex) ---
\Askhsh{1827}{4}
\begin{ylh}{1}
    \item 5.2. Παραλληλόγραμμα
    \item 5.7. Βαρύκεντρο τριγώνου.
\end{ylh}
Δίνεται τρίγωνο $AB\varGamma $ και $ E$ το μέσο της διαμέσου $B\varDelta  $. Στην προέκταση της $ AE$ θεωρούμε σημείο $ Z$ τέτοιο ώστε $E Z = A E$ και έστω $\varTheta $ το σημείο τομής της $ A Z$ με την πλευρά $B\varGamma $.
Να αποδείξετε ότι:
\begin{alist}
    \item Το τετράπλευρο $AB Z\varDelta  $ είναι παραλληλόγραμμο. \monades{8}
    \item Το τετράπλευρο $B\varDelta  \varGamma  Z$ είναι παραλληλόγραμμο. \monades{8}
    \item Το σημείο $\varTheta $ είναι βαρύκεντρο του τριγώνου $B\Delta Z$. \monades{9}
\end{alist}

\vspace{4mm}

% --- Άσκηση 1851 (1851.tex) ---
\Askhsh{1851}{4}
\begin{ylh}{1}
    \item 3.11. Ανισοτικές σχέσεις πλευρών και γωνιών
    \item 4.6. Άθροισμα γωνιών τριγώνου.
\end{ylh}
% ===================================================================
%  Αρχείο: 1851.tex (Fragment)
%  Αυτόματη μετατροπή μέσω doc2latex.py
% ===================================================================



Σε τρίγωνο $ABC$, η προέκταση της διχοτόμου της $\hat{\varGamma }$ και της εξωτερικής γωνίας του $\hat{B}$, τέμνονται στο Ε. Δίνεται ότι $A\hat{B}E = 70^{\circ} = 2\varGamma \hat{E}B$.


\begin{alist}
\item Να αποδείξετε ότι το τρίγωνο Γ$BE$ είναι ισοσκελές \monades{12}

\item Να υπολογίσετε τις γωνίες του τριγώνου $ABC$. \monades{13}


\begin{center}
\begin{tikzpicture}[scale=1]
% External bisector of B and internal of C meet at E.
% E is an excenter. E = Ea usually.
% Let's pick some angles. A=40, B=80, C=60.
\tkzDefPoint(0,0){B}
\tkzDefPoint(5,0){C}
\tkzDefTriangle[two angles = 80 and 60](B,C) \tkzGetPoint{A}

% Internal bisector of C
\tkzDefLine[bisector](B,C,A) \tkzGetPoint{bisC}

% External bisector of B.
% Extend $CB$ past B to some point X.
\tkzDefPointBy[homothety=center B ratio -0.5](C) \tkzGetPoint{X}
\tkzDefLine[bisector](A,B,X) \tkzGetPoint{bisBext}

\tkzInterLL(C,bisC)(B,bisBext) \tkzGetPoint{E}

\draw[l3] (A)--(B)--(C)--cycle;
\draw[l3] (C)--(E);
\draw[l3] (B)--(E);
\draw[l3] (B)--(X);

\tkzDrawPoints(A,B,C,E)
\tkzLabelPoint[above](A){$A$}
\tkzLabelPoint[below](B){$B$}
\tkzLabelPoint[below](C){$\varGamma $}
\tkzLabelPoint[right](E){$E$}

\tkzMarkSegments[mark=s|](B,E B,C) % Problem says triangle $GBE$ is isosceles.
% Actually if B is externally bisected and C internally bisected, E is excenter.
% $GB$=$BE$ implies something about the angles.
\end{tikzpicture}
\end{center}
\end{alist}

\vspace{4mm}

% --- Άσκηση 1872 (1872.tex) ---
\Askhsh{1872}{4}
\begin{ylh}{1}
    \item 3.6. Κριτήρια ισότητας ορθογώνιων τριγώνων
    \item 4.6. Άθροισμα γωνιών τριγώνου
    \item 5.4. Ρόμβος
    \item 5.9. Μια ιδιότητα του ορθογώνιου τριγώνου.
\end{ylh}
Έστω ορθογώνιο τρίγωνο $AB\varGamma $ με $\hat{ A} = 90^\circ$ και $\hat{B} = 60^\circ$. Η διχοτόμος της γωνίας $\hat{B}$ τέμνει την $ A\varGamma $ στο $ Z$. Τα σημεία $ M$ και $ K$ είναι τα μέσα των $B Z$ και $B\varGamma $ αντίστοιχα. Αν το τμήμα $\varGamma \Lambda$ είναι κάθετο στη διχοτόμο $B\delta$ να αποδείξετε:
\begin{alist}
    \item Το τρίγωνο $B Z\varGamma $ είναι ισοσκελές. \monades{6}
    \item Το τετράπλευρο $A M K\mathrm{Z}$ είναι ρόμβος. \monades{6}
    \item $\varGamma  Z = 2 ZA$. \monades{7}
    \item $B\Lambda = A\varGamma $. \monades{6}
\end{alist}

\vspace{4mm}

% --- Άσκηση 1874 (1874.tex) ---
\Askhsh{1874}{4}
\begin{ylh}{1}
    \item 3.2. 1ο Κριτήριο ισότητας τριγώνων
    \item 4.6. Άθροισμα γωνιών τριγώνου
    \item 5.9. Μια ιδιότητα του ορθογώνιου τριγώνου.
\end{ylh}
Έστω κύκλος με κέντρο $ O$ και διάμετρο $ K\Lambda$. Έστω $A$ σημείο του κύκλου ώστε η ακτίνα $\mathrm{O}A$ να είναι κάθετη στην $ K\Lambda$. Φέρουμε τις χορδές $ $AB$ = A\varGamma  = \rho$. Έστω $\varDelta  $ και $ E$ τα σημεία τομής των προεκτάσεων των $ AB$ και $A\varGamma $ αντίστοιχα με την ευθεία της διαμέτρου $\mathrm{K}\Lambda$.
Να αποδείξετε ότι:
\begin{alist}
    \item Η γωνία $B\hat{A}\varGamma $ είναι $120^\circ$. \monades{7}
    \item Τα σημεία $B$ και $\varGamma $ είναι μέσα των $A\varDelta  $ και $A E$ αντίστοιχα. \monades{9}
    \item $ K\varGamma  = \Lambda B$. \monades{9}
\end{alist}

\vspace{4mm}

% --- Άσκηση 1875 (1875.tex) ---
\Askhsh{1875}{4}
\begin{ylh}{1}
    \item 3.2. 1ο Κριτήριο ισότητας τριγώνων
    \item 3.6. Κριτήρια ισότητας ορθογώνιων τριγώνων
    \item 3.11. Ανισοτικές σχέσεις πλευρών και γωνιών.
\end{ylh}
Θεωρούμε ισοσκελές τρίγωνο $AB\varGamma $ ($AB=A\varGamma $), και την ευθεία $\varepsilon$ της εξωτερικής διχοτόμου της γωνίας $\hat{ A}$. Η κάθετη στην πλευρά $AB$ στο $B$ τέμνει την $\varepsilon$ στο $ K$ και την ευθεία $A\varGamma $ στο $\mathrm{H}$. Η κάθετη στην πλευρά $A\varGamma $ στο $\varGamma $ τέμνει την $\varepsilon$ στο $\Lambda$ και την ευθεία $ AB$ στο $E$.
\begin{alist}
    \item Να αποδείξετε ότι:
    \begin{rlist}
        \item $A H = A E$. \monades{8}
        \item $A K = A\Lambda$. \monades{9}
    \end{rlist}
    \item Ένας μαθητής κοιτώντας το σχήμα, διατύπωσε την άποψη ότι η $A\varTheta $ είναι διχοτόμος της γωνίας $\hat{ A}$ του τριγώνου $ $AB$\varGamma $, όπου $\varTheta $ το σημείο τομής των $ K H$ και $ E\Lambda$. Συμφωνείτε με την παραπάνω σκέψη του μαθητή ή όχι; Δικαιολογήστε πλήρως την απάντησή σας. \monades{8}
\end{alist}

\vspace{4mm}

% --- Άσκηση 1876 (1876.tex) ---
\Askhsh{1876}{4}
\begin{ylh}{1}
    \item 3.1. Είδη και στοιχεία τριγώνων
    \item 3.11. Ανισοτικές σχέσεις πλευρών και γωνιών
    \item 4.2. Τέμνουσα δύο ευθειών - Ευκλείδειο αίτημα
    \item 4.4. Γωνίες με πλευρές παράλληλες
    \item 5.6. Εφαρμογές στα τρίγωνα
    \item 5.9. Μια ιδιότητα του ορθογώνιου τριγώνου.
\end{ylh}
Δίνονται δυο ίσα ισοσκελή τρίγωνα $AB\varGamma $ ($AB=A\varGamma $) και $AB\Delta$ ($BA=B\Delta$), τέτοια ώστε οι πλευρές τους $A\varGamma $ και $B\Delta$ να τέμνονται κάθετα στο σημείο $E$, όπως φαίνεται στο παρακάτω σχήμα. Τα σημεία $ K$ και $\Lambda$ είναι τα μέσα των τμημάτων $A\Delta$ και $B\varGamma $ αντίστοιχα.
Να αποδείξετε ότι:
\begin{alist}
    \item $ E\varDelta   =  E\varGamma $. \monades{7}
    \item $\varDelta  \varGamma  \parallel AB$. \monades{8}
    \item Το τρίγωνο $ E K\Lambda$ είναι ισοσκελές και $ K\Lambda \parallel AB$. \monades{10}
\end{alist}

\vspace{4mm}

% --- Άσκηση 1877 (1877.tex) ---
\Askhsh{1877}{4}
\begin{ylh}{1}
    \item 3.4. 3ο Κριτήριο ισότητας τριγώνων
    \item 5.2. Παραλληλόγραμμα.
\end{ylh}
Έστω παραλληλόγραμμο $AB\varGamma \varDelta  $ με $ O$ το σημείο τομής των διαγωνίων του και $ K$ το μέσο του $\varGamma \varDelta  $. Προεκτείνουμε το τμήμα $\mathrm{O} K$ κατά τμήμα $ K Z =  K\mathrm{O}$. Η $B Z$ τέμνει τη διαγώνιο $ A\varGamma $ στο $\varTheta $.
Να αποδείξετε ότι:
\begin{alist}
    \item Τα τμήματα $ O\varGamma $ και $B Z$ διχοτομούνται. \monades{8}
    \item $A O = \varDelta   Z$. \monades{9}
    \item Τα τρίγωνα $A OB$ και $\varDelta   Z\varGamma $ είναι ίσα. \monades{8}
\end{alist}

\vspace{4mm}

% --- Άσκηση 1878 (1878.tex) ---
\Askhsh{1878}{4}
\begin{ylh}{1}
    \item 5.2. Παραλληλόγραμμα
    \item 5.6. Εφαρμογές στα τρίγωνα
    \item 5.7. Βαρύκεντρο τριγώνου.
\end{ylh}
Έστω ισοσκελές τρίγωνο $AB\varGamma $ με $AB = A\varGamma $. Προεκτείνουμε το $B\varGamma $ (προς το $\varGamma $) κατά τμήμα $\varGamma \varDelta   = B\varGamma $. Φέρνουμε τις διαμέσους $A E$ και $\varGamma  H$ του τριγώνου $AB\varGamma $ που τέμνονται στο $\varTheta $. Το $B\varTheta $ προεκτεινόμενο, τέμνει το $A\varGamma $ στο $ K$ και το $ A\varDelta  $ στο $\Lambda$.
Να αποδείξετε ότι:
\begin{alist}
    \item Το $ H K\varGamma  E$ είναι παραλληλόγραμμο. \monades{9}
    \item $A H =  K\varGamma $. \monades{9}
    \item $A\varTheta  = 2 H\varTheta $. \monades{7}
\end{alist}

\vspace{4mm}

% --- Άσκηση 1879 (1879.tex) ---
\Askhsh{1879}{4}
\begin{ylh}{1}
    \item 3.4. 3ο Κριτήριο ισότητας τριγώνων
    \item 3.6. Κριτήρια ισότητας ορθογώνιων τριγώνων
    \item 5.2. Παραλληλόγραμμα
    \item 5.3. Ορθογώνιο.
\end{ylh}
Έστω κύκλος με κέντρο $ O$ και διάμετρο $AB$. Φέρνουμε χορδή $\varGamma \varDelta   \parallel AB$ με $ K$ το μέσο της. Από το $\varDelta  $ φέρνουμε το τμήμα $\varDelta   E$ κάθετο στη $\varDelta  \varGamma $.
Να αποδείξετε ότι:
\begin{alist}
    \item Το τετράπλευρο $ K\varGamma  O E$ είναι παραλληλόγραμμο. \monades{8}
    \item $\varDelta  \hat{ E} K = \frac{\varDelta  \hat{ O}\varGamma }{2}$. \monades{12}
    \item $ K E <  KB$. \monades{5}
\end{alist}

\vspace{4mm}

% --- Άσκηση 1893 (1893.tex) ---
\Askhsh{1893}{4}
\begin{ylh}{1}
    \item 3.6. Κριτήρια ισότητας ορθογώνιων τριγώνων
    \item 3.11. Ανισοτικές σχέσεις πλευρών και γωνιών
    \item 5.3. Ορθογώνιο
    \item 5.6. Εφαρμογές στα τρίγωνα
    \item 5.10. Τραπέζιο
    \item 5.11. Ισοσκελές τραπέζιο.
\end{ylh}
Έστω $AB\varGamma \varDelta  $ ορθογώνιο με $ $AB$ > B\varGamma $ τέτοιο ώστε οι διαγώνιοί του να σχηματίζουν γωνία $60^\circ$. Από το $\varDelta  $ φέρνουμε $\varDelta   M$ κάθετη στην $ A\varGamma $.
\begin{alist}
    \item Να αποδείξετε ότι:
    \begin{rlist}
        \item το σημείο $ M$ είναι μέσο του $A O$ όπου $ O$ το κέντρο του ορθογωνίου. \monades{8}
        \item $A M = \frac{1}{4}A\varGamma $. \monades{7}
    \end{rlist}
    \item Αν από το $\varGamma $ φέρνουμε $\varGamma  N$ κάθετη στη $B\varDelta  $, να αποδείξετε ότι το $ M\mathrm{N}\varGamma \varDelta  $ είναι ισοσκελές τραπέζιο. \monades{10}
\end{alist}

\vspace{4mm}

% --- Άσκηση 1894 (1894.tex) ---
\Askhsh{1894}{4}
\begin{ylh}{1}
    \item 3.6. Κριτήρια ισότητας ορθογώνιων τριγώνων
    \item 4.6. Άθροισμα γωνιών τριγώνου.
\end{ylh}
Σε ορθογώνιο τρίγωνο $AB\varGamma $ ($\hat{A} = 90^\circ$) φέρνουμε τη διχοτόμο του $ A\varDelta  $. Έστω $\varDelta   K$ και $\varDelta  \mathrm{P}$ οι προβολές του $\varDelta  $ στις $ AB$ και $A\varGamma $ αντίστοιχα. Η κάθετη της $B\varGamma $ στο σημείο $\varDelta  $ τέμνει την πλευρά $A\varGamma $ στο $ E$ και την προέκταση της πλευράς $AB$ (προς το $B$) στο σημείο $ Z$.
\begin{alist}
    \item Να αποδείξετε ότι:
    \begin{rlist}
        \item $\hat{B} = \varDelta  \hat{ E}\varGamma $. \monades{8}
        \item $\varDelta   E = \varDelta   B$. \monades{8}
    \end{rlist}
    \item Να υπολογίσετε τη γωνία $\varDelta  \hat{\varGamma } Z$. \monades{9}
\end{alist}

\vspace{4mm}

% --- Άσκηση 1895 (1895.tex) ---
\Askhsh{1895}{4}
\begin{ylh}{1}
    \item 3.6. Κριτήρια ισότητας ορθογώνιων τριγώνων
    \item 4.2. Τέμνουσα δύο ευθειών - Ευκλείδειο αίτημα
    \item 5.6. Εφαρμογές στα τρίγωνα
    \item 5.9. Μια ιδιότητα του ορθογώνιου τριγώνου.
\end{ylh}
Δίνεται ισοσκελές τρίγωνο $A\varGamma  B$ ($A\varGamma  = \varGamma  B$). Φέρνουμε τα ύψη του $ A K$ και $\varGamma \Lambda$. Αν $ E$ είναι το μέσο της πλευράς $ A\varGamma $, να αποδείξετε ότι:
\begin{alist}
    \item Tο τρίγωνο $ K E\Lambda$ είναι ισοσκελές. \monades{10}
    \item H $ K\Lambda$ είναι διχοτόμος της γωνίας $B\hat{ K} E$. \monades{15}
\end{alist}

\vspace{4mm}

% --- Άσκηση 1898 (1898.tex) ---
\Askhsh{1898}{4}
\begin{ylh}{1}
    \item 4.2. Τέμνουσα δύο ευθειών - Ευκλείδειο αίτημα
    \item 5.2. Παραλληλόγραμμα
    \item 5.3. Ορθογώνιο
    \item 5.6. Εφαρμογές στα τρίγωνα.
\end{ylh}
Δίνεται τρίγωνο $AB\varGamma $ και η διάμεσός του $ A\varDelta  $. Έστω $ E,  H$ και $\varTheta $ τα μέσα των $B\varDelta  , A\varDelta  $ και $A\varGamma $ αντίστοιχα.
\begin{alist}
    \item Να αποδείξετε ότι το τετράπλευρο $\varDelta   E H\varTheta $ είναι παραλληλόγραμμο. \monades{10}
    \item Να βρείτε τη σχέση των πλευρών $AB$ και $B\varGamma $ του τριγώνου $ $AB$\varGamma $, ώστε το παραλληλόγραμμο $\varDelta   E H\varTheta $ να είναι ρόμβος. \monades{10}
    \item Στην περίπτωση που το τρίγωνο $AB\varGamma $ είναι ορθογώνιο (η γωνία $\hat{B}$ ορθή), να βρείτε το είδος του παραλληλογράμμου $\varDelta   E H\varTheta $. \monades{5}
\end{alist}

\vspace{4mm}

% --- Άσκηση 11882 (11882.tex) ---
\Askhsh{11882}{3}
\begin{ylh}{1}
    \item 3.6. Κριτήρια ισότητας ορθογώνιων τριγώνων
    \item 4.6. Άθροισμα γωνιών τριγώνου.
\end{ylh}
% ===================================================================
%  Αρχείο: 11882.tex (Fragment)
%  Αυτόματη μετατροπή μέσω doc2latex.py
% ===================================================================



Στο παρακάτω σχήμα τα τρίγωνα $AB\varGamma $, $A\varDelta   E$ και $\varDelta   Z$Η είναι ορθογώνια με ορθές γωνίες $A\hat{B}\varGamma $, $A\hat{E}\varDelta  $ και $\varDelta  \hat{Z}H$, αντίστοιχα. Επίσης $A\varGamma $ = $A\varDelta  $ και $B\varGamma $ = $\varDelta   Z$.

Να αποδείξετε ότι:

\begin{alist}
\item Τα ευθύγραμμα τμήματα $B\varGamma $ και $\varDelta   E$ είναι ίσα.

\monades{10}

\item Η $\varDelta   H$ είναι διχοτόμος της γωνίας $E\hat{H}Z$.

\monades{6}

\item Αν, επιπλέον, οι $A\varDelta  $ και $\varDelta   H$ είναι κάθετες, τότε $A\hat{\varDelta  }E$ = $\frac{E\hat{H}Z}{2}$ .

\monades{9}


\begin{center}
\begin{tikzpicture}[scale=1.2]
\tkzDefPoint(1,2){B}
\tkzDefPoint(2,2){A}
\tkzDefPoint(1,0.5){C}
\tkzDefPoint(3,3.5){D}
\tkzDefPoint(3,2){E}
\tkzDefPoint(3.83,4.75){Z}
\tkzDefPoint(7.95,2){H}
\draw[line width=0.3mm] (B)--(C)--(A)--(B);
\draw[line width=0.3mm] (A)--(D)--(E)--(A);
\draw[line width=0.3mm] (D)--(Z)--(H)--(D);
\draw[line width=0.3mm] (B)--(H);
\tkzMarkRightAngle[size=.2](A,B,C)
\tkzMarkRightAngle[size=.2](A,E,D)
\tkzMarkRightAngle[size=.2](D,Z,H)
\tkzMarkSegments[mark=s||,size=3pt](C,A A,D)
\tkzMarkSegments[mark=s|,size=3pt](C,B D,Z)
\tkzDrawPoints(A,B,C,D,E,Z,H)
\tkzLabelPoint[above left](A){$A$}
\tkzLabelPoint[above left](B){$B$}
\tkzLabelPoint[left](C){$\varGamma $}
\tkzLabelPoint[above left](D){$\varDelta  $}
\tkzLabelPoint[below](E){$E$}
\tkzLabelPoint[above](Z){$Z$}
\tkzLabelPoint[right](H){$H$}
\end{tikzpicture}
\end{center}
\end{alist}

\vspace{4mm}

% --- Άσκηση 11896 (11896.tex) ---
\Askhsh{11896}{3}
\begin{ylh}{1}
    \item 5.4. Ρόμβος
    \item 5.6. Εφαρμογές στα τρίγωνα.
\end{ylh}
Στο τετράπλευρο $AB\varGamma \varDelta  $ του σχήματος ισχύει ότι $A\varDelta  \  = B\varGamma $ και τα σημεία $E, Z, H$ και $\varTheta$ είναι τα μέσα των ευθύγραμμων τμημάτων $AB$, $A\varGamma $, $\varGamma \varDelta  $ και $B\varDelta  $ αντίστοιχα.

Να δείξετε ότι:
\begin{alist}
\item  $EZ=H\varTheta $
\monades{15}
\item Το τετράπλευρο $EZH\varTheta $ είναι ρόμβος.\monades{10}
\end{alist}

\begin{center}
\begin{tikzpicture}
\tkzDefPoint(2,4){A}
\tkzDefPoint(7,3.5){B}
\tkzDefPoint(10,2.5){C}
\tkzDefPoint(1,1){D}
\tkzDefMidPoint(A,B)\tkzGetPoint{E}
\tkzDefMidPoint(A,C)\tkzGetPoint{Z}
\tkzDefMidPoint(C,D)\tkzGetPoint{H}
\tkzDefMidPoint(B,D)\tkzGetPoint{Th}
\draw[line width=0.3mm] (A)--(B)--(C)--(D)--cycle;
\draw[line width=0.3mm] (E)--(Z)--(H)--(Th)--cycle;
\draw[dashed, line width=0.3mm] (A)--(C);
\draw[dashed, line width=0.3mm] (B)--(D);
\tkzDrawPoints(A,B,C,D,E,Z,H,Th)
\tkzLabelPoint[above left](A){$A$}
\tkzLabelPoint[above right](B){$B$}
\tkzLabelPoint[right](C){$\varGamma $}
\tkzLabelPoint[left](D){$\varDelta  $}
\tkzLabelPoint[above](E){$E$}
\tkzLabelPoint[below right](Z){$Z$}
\tkzLabelPoint[below](H){$H$}
\tkzLabelPoint[above left](Th){$\varTheta $}
\tkzMarkSegments[mark=s||,size=3pt](A,D B,C)
\end{tikzpicture}
\end{center}

\vspace{4mm}

% --- Άσκηση 11897 (11897.tex) ---
\Askhsh{11897}{3}
\begin{ylh}{1}
    \item 3.3. 2ο Κριτήριο ισότητας τριγώνων
    \item 4.2. Τέμνουσα δύο ευθειών - Ευκλείδειο αίτημα
    \item 4.6. Άθροισμα γωνιών τριγώνου
    \item 5.2. Παραλληλόγραμμα.
\end{ylh}
% ===================================================================
%  Αρχείο: 11897.tex (Fragment)
%  Αυτόματη μετατροπή μέσω doc2latex.py
% ===================================================================



Δίνεται τρίγωνο $AB\varGamma $ και η διάμεσος του $AM$. Στην προέκταση της $A\varGamma $ προς το $\varGamma$ παίρνουμε τμήμα $\varGamma \varDelta  $ = $A\varGamma $. Από το $\varDelta  $ φέρνουμε παράλληλη προς την $AM$ που τέμνει την προέκταση της $B\varGamma $ στο Ε. Να αποδείξετε ότι:

\begin{alist}
\item $M\varGamma $ = $\varGamma  E$. \monades{9}

\item Το τετράπλευρο $AM\varDelta   E$ είναι παραλληλόγραμμο. \monades{7}

\item $\hat{B}+B\hat{A}M=\varGamma \hat{E}\varDelta  $. \monades{9}
\end{alist}

\vspace{4mm}

% --- Άσκηση 12068 (12068.tex) ---
\Askhsh{12068}{3}
\begin{ylh}{1}
    \item 3.12. Τριγωνική ανισότητα
    \item 4.2. Τέμνουσα δύο ευθειών - Ευκλείδειο αίτημα
    \item 5.2. Παραλληλόγραμμα
    \item 5.4. Ρόμβος
    \item 5.5. Τετράγωνο
    \item 5.6. Εφαρμογές στα τρίγωνα.
\end{ylh}
% ===================================================================
%  Αρχείο: 12068.tex (Fragment)
%  Αυτόματη μετατροπή μέσω doc2latex.py
% ===================================================================



Δίνεται ορθογώνιο τρίγωνο $AB\varGamma $ ($\hat{A}$=90°) και $M$ το μέσο της υποτείνουσας του $B\varGamma $. Από το $M$ φέρουμε $M\varDelta  $⊥$AB$ και προεκτείνουμε κατά ίσο τμήμα $\varDelta   Z$.

\begin{alist}
\item Να αποδείξετε ότι:
\end{alist}
\begin{alist}
\item
  Το τρίγωνο $MBZ$ είναι ισοσκελές. \monades{8}
\item
  Το τετράπλευρο $AMBZ$ είναι ρόμβος. \monades{9}

\item Αν το αρχικό τρίγωνο $AB\varGamma $ είναι ορθογώνιο και ισοσκελές, τι είδους τετράπλευρο είναι το $AMBZ$; Δικαιολογήστε την απάντησή σας. \monades{8}
\end{alist}

\vspace{4mm}

% --- Άσκηση 12069 (12069.tex) ---
\Askhsh{12069}{3}
\begin{ylh}{1}
    \item 3.1. Είδη και στοιχεία τριγώνων
    \item 3.2. 1ο Κριτήριο ισότητας τριγώνων
    \item 3.4. 3ο Κριτήριο ισότητας τριγώνων
    \item 3.11. Ανισοτικές σχέσεις πλευρών και γωνιών.
\end{ylh}
Σε ισοσκελές τρίγωνο $AB\varGamma $ ($AB=A\varGamma $) παίρνουμε στην πλευρά $AB$ σημείο $\varDelta  $, ώστε $\varDelta   B$ = 2ΑΔ, και στην πλευρά $A\varGamma $ σημείο $E$, ώστε $E\varGamma $= 2ΑΕ. Το $M$ είναι το μέσο της πλευράς $B\varGamma $ του τριγώνου $AB\varGamma $.

\begin{alist}
\item Να αποδείξετε ότι:
\end{alist}
\begin{alist}
\item
  Τα τμήματα $\varDelta   B$ και $E\varGamma $ είναι ίσα. \monades{6}
\item
  Το τρίγωνο $M\varDelta  E$ είναι ισοσκελές. \monades{6}

\item Αν Ρ το σημείο τομής των τμημάτων $BE$ και $\varGamma \varDelta  $ να δείξετε ότι:
\end{alist}
\begin{alist}
\item
  Οι γωνίες $\varGamma \hat{B}E$ και $B\hat{\varGamma }\varDelta  $ είναι ίσες. \monades{6}
\item
  Το τμήμα $PM$ διχοτομεί τη γωνία $B\hat{Ρ}\varGamma $. \monades{7}
\end{alist}

\vspace{4mm}

% --- Άσκηση 12149 (12149.tex) ---
\Askhsh{12149}{2}
\begin{ylh}{1}
    \item 3.2. 1ο Κριτήριο ισότητας τριγώνων
    \item 3.6. Κριτήρια ισότητας ορθογώνιων τριγώνων.
\end{ylh}

Δίνονται τα αμβλυγώνια τρίγωνα $AB\varGamma $ ($\hat{A} > 90\degree$) και $A'B'\varGamma  '$ ($\hat{A}΄ > 90\degree$) με γ = γ΄ και β = β΄. Αν τα ύψη $BH$ και Β΄Η΄ των τριγώνων $AB\varGamma $ και $A'B'\varGamma  '$ αντίστοιχα είναι ίσα, να αποδείξετε ότι:
\begin{alist}
\item $Β\hat{A}H=B'\hat{A}΄H΄$.
\monades{13}
\item Τα τρίγωνα $AB\varGamma $ και $A'B'\varGamma  '$ είναι ίσα.
\monades{12}
\begin{center}
\begin{tikzpicture}
\tkzDefPoint(3,1){A}
\tkzDefPoint(1,4){B}
\tkzDefPoint(7,1){C}
\tkzDefPoint(1,1){H}
\draw[line width=0.3mm] (A)--(B)--(C)--cycle;
\draw[dashed, line width=0.3mm] (B)--(H)--(A);
\tkzDrawPoints(A,B,C,H)
\tkzLabelPoint[below](A){$A$}
\tkzLabelPoint[above](B){$B$}
\tkzLabelPoint[below](C){$\varGamma $}
\tkzLabelPoint[below](H){$H$}
\tkzMarkRightAngle[size=.2](B,H,A)

\begin{scope}[xshift=8cm]
\tkzDefPoint(3,1){A'}
\tkzDefPoint(1,4){B'}
\tkzDefPoint(7,1){C'}
\tkzDefPoint(1,1){H'}
\draw[line width=0.3mm] (A')--(B')--(C')--cycle;
\draw[dashed, line width=0.3mm] (B')--(H')--(A');
\tkzDrawPoints(A',B',C',H')
\tkzLabelPoint[below](A'){$A'$}
\tkzLabelPoint[above](B'){$B'$}
\tkzLabelPoint[below](C'){$\varGamma '$}
\tkzLabelPoint[below](H'){$H'$}
\tkzMarkRightAngle[size=.2](B',H',A')
\end{scope}
\end{tikzpicture}
\end{center}
\end{alist}

\vspace{4mm}

% --- Άσκηση 12165 (12165.tex) ---
\Askhsh{12165}{3}
\begin{ylh}{1}
    \item 3.11. Ανισοτικές σχέσεις πλευρών και γωνιών
    \item 4.2. Τέμνουσα δύο ευθειών - Ευκλείδειο αίτημα
    \item 5.2. Παραλληλόγραμμα
    \item 5.9. Μια ιδιότητα του ορθογώνιου τριγώνου.
\end{ylh}
% ===================================================================
%  Αρχείο: 12165.tex (Fragment)
%  Αυτόματη μετατροπή μέσω doc2latex.py
% ===================================================================



Δίνεται παραλληλόγραμμο $AB\varGamma \varDelta  $ με $\hat{A} = 120\degree$ και η διχοτόμος της γωνίας $\hat{\varDelta  }$ που τέμνει την $AB$ στο μέσο της Ε.

\begin{alist}
\item Να αποδείξετε ότι $AB$ = 2ΑΔ\textbf{.} \monades{6}

\item Αν το κάθετο ευθύγραμμο τμήμα που φέρνουμε από το σημείο $E$ στην $\varGamma \varDelta  $ την τέμνει στο Η, τότε να αποδείξετε ότι $\frac{\varDelta   E}{HE} =2$. \monades{7}

\item Αν $M$ το μέσο της $\varGamma \varDelta  $, τότε να αποδείξετε ότι το τρίγωνο $MA\varDelta  $ είναι ισόπλευρο.

\monades{6}

\item Να αποδείξετε ότι $\varDelta  \hat{A}\varGamma  = 90\degree$. \monades{6}


\begin{center}
\begin{tikzpicture}[scale=1.2]
\tkzDefPoint(0,0){D}
\tkzDefPoint(4,0){C}
\tkzDefPoint(5, 2.598){B}
\tkzDefPoint(1, 2.598){A}
\tkzDefLine[bisector](C,D,A)\tkzGetPoint{e}
\tkzInterLL(A,B)(D,e)\tkzGetPoint{E}
\draw[l3] (A)--(B)--(C)--(D)--cycle;
\draw[l4,maincolor] (D)--(E);
\tkzMarkAngle[size=0.5,mark=none,fill=maincolor!20](C,D,E)
\tkzLabelAngle[pos=0.7](C,D,E){$1$}
\tkzMarkAngle[size=0.6,mark=none,fill=maincolor!20](E,D,A)
\tkzLabelAngle[pos=0.8](E,D,A){$2$}
\tkzMarkAngle[size=0.5,mark=none,fill=maincolor!20](A,E,D)
\tkzLabelAngle[pos=0.7](A,E,D){$1$}
\tkzDrawPoints(A,B,C,D,E)
\tkzLabelPoint[above left](A){$A$}
\tkzLabelPoint[above right](B){$B$}
\tkzLabelPoint[below right](C){$\varGamma $}
\tkzLabelPoint[below left](D){$\varDelta  $}
\tkzLabelPoint[above](E){$E$}
\end{tikzpicture}
\end{center}
\end{alist}

\vspace{4mm}

% --- Άσκηση 12200 (12200.tex) ---
\Askhsh{12200}{3}
\begin{ylh}{1}
    \item 3.2. 1ο Κριτήριο ισότητας τριγώνων
    \item 4.2. Τέμνουσα δύο ευθειών - Ευκλείδειο αίτημα
    \item 4.6. Άθροισμα γωνιών τριγώνου.
\end{ylh}
\begin{wrapfigure}[5]{r}{4cm}
\begin{tikzpicture}[scale=0.7]
\tkzDefPoint(-2,0){B}
\tkzDefPoint(2,0){C}
\tkzDefTriangle[two angles=72 and 72](B,C)\tkzGetPoint{A}
\tkzDefLine[bisector](C,B,A)\tkzGetPoint{d}
\tkzInterLL(A,C)(B,d)\tkzGetPoint{D}
\tkzInterLC[with nodes](A,B)(A,C,D)\tkzGetFirstPoint{E}
\tkzDefLine[parallel=through B](A,C)\tkzGetPoint{p}
\tkzDrawPoints(A,B,C,D,E)
\tkzDrawPolygon[line width=0.3mm](A,B,C)
\tkzDrawSegments(B,D D,E)
\tkzLabelPoint[above](A){$A$}
\tkzLabelPoint[below left](B){$B$}
\tkzLabelPoint[below right](C){$\varGamma $}
\tkzLabelPoint[right](D){$\varDelta  $}
\tkzLabelPoint[left](E){$E$}
\tkzMarkAngle[size=0.6,mark=none,fill=gray!20](B,A,C)
\tkzLabelAngle[pos=0.8](B,A,C){$36^\circ$}
\tkzMarkSegments[mark=|,size=3pt](A,E C,D)
\end{tikzpicture}
\end{wrapfigure}
Δίνεται ισοσκελές τρίγωνο $AB\varGamma $ με $AB = A\varGamma $ και $\hat{A} = 36\degree$. Έστω $B\varDelta  $ η διχοτόμος της γωνίας $\hat{B}$ και $E$ σημείο της πλευράς $AB$ ώστε $AE = \varGamma \varDelta  $.
\begin{alist}
\item Να αποδείξετε ότι $A\varDelta   = B\varDelta  $. \monades{8}
\item Να αποδείξετε ότι το τρίγωνο $B\varDelta E$ είναι ισοσκελές. \monades{7}
\item Η παράλληλη από το $B$ προς την $A\varGamma $ τέμνει την προέκταση της $\varDelta  E$ (προς το $E$) στο σημείο $Z$. Να αποδείξετε ότι το τρίγωνο $B\varDelta Z$ είναι ισοσκελές. \monades{10}
\end{alist}

\vspace{4mm}

% --- Άσκηση 12417 (12417.tex) ---
\Askhsh{12417}{2}
\begin{ylh}{1}
    \item 3.11. Ανισοτικές σχέσεις πλευρών και γωνιών
    \item 3.16. Σχετικές θέσεις δυο κύκλων.
\end{ylh}
% ===================================================================
%  Αρχείο: 12417.tex (Fragment)
%  Αυτόματη μετατροπή μέσω doc2latex.py
% ===================================================================



Έστω δύο κύκλοι (Κ, R) και (Λ, r), με R=3, r=2 και $K\varLambda $=4. Να αποδείξετε ότι:

\begin{alist}
\item Οι κύκλοι (Κ, R) και (Λ, r) τέμνονται σε δύο σημεία, έστω Α και Β.

\monades{15}

\item K$\hat{A}$Λ$>$Α$\hat{\text{Λ}}$Κ.

\monades{10}
\end{alist}

\vspace{4mm}

% --- Άσκηση 12418 (12418.tex) ---
\Askhsh{12418}{3}
\begin{ylh}{1}
    \item 3.2. 1ο Κριτήριο ισότητας τριγώνων
    \item 5.11. Ισοσκελές τραπέζιο.
\end{ylh}
Δίνεται ισοσκελές τραπέζιο $AB\varGamma \varDelta  $ ($AB\parallel\varGamma \varDelta  $) με $AB>\varGamma\varDelta$. Κατασκευάζουμε εξωτερικά του τραπεζίου $AB\varGamma \varDelta  $ ισοσκελές τρίγωνο $ABE$ με βάση $AB$. Αν $M$ είναι το μέσο της βάσης $\varGamma\varDelta $, να αποδείξετε ότι:

\begin{alist}
\item Τα τρίγωνα $AE\varDelta $ και $BE\varGamma$ είναι ίσα.
\monades{11}
\item Η διάμεσος $EM$ του τριγώνου $E\varDelta  $Γ είναι διχοτόμος της γωνίας $A\hat{E}B$.
\monades{14}
\end{alist}

\vspace{4mm}

% --- Άσκηση 12635 (12635.tex) ---
\Askhsh{12635}{2}
\begin{ylh}{1}
    \item 3.1. Είδη και στοιχεία τριγώνων
    \item 3.2. 1ο Κριτήριο ισότητας τριγώνων.
\end{ylh}
% ===================================================================
%  Αρχείο: 12635.tex (Fragment)
%  Αυτόματη μετατροπή μέσω doc2latex.py
% ===================================================================



Δίνεται ισοσκελές τρίγωνο $AB\varGamma $ με $AB=A\varGamma $ και $M$ είναι το μέσο της βάσης του $B\varGamma $. Στις προεκτάσεις των πλευρών $AB$, $A\varGamma $ προς τα $B,\varGamma$ αντίστοιχα, παίρνουμε τα τμήματα $B\varDelta  $ και $\varGamma E$ ώστε $B\varDelta  $=$\varGamma E$.

\begin{alist}
\item Να αποδείξετε ότι τα τρίγωνα Μ$B\varDelta  $ και $M\varGamma E$ είναι ίσα. \monades{15}

\item Να αποδείξετε ότι η γωνία $M\varDelta  E$ είναι ίση με τη γωνία Μ$E\varDelta  $. \monades{10}
\end{alist}

\vspace{4mm}

% --- Άσκηση 12636 (12636.tex) ---
\Askhsh{12636}{2}
\begin{ylh}{1}
    \item 3.1. Είδη και στοιχεία τριγώνων
    \item 3.2. 1ο Κριτήριο ισότητας τριγώνων
    \item 3.4. 3ο Κριτήριο ισότητας τριγώνων.
\end{ylh}
\begin{wrapfigure}[5]{r}{3cm}
\begin{tikzpicture}[scale=0.7]
\tkzDefPoint(0,5){A}
\tkzDefPoint(-1.5,0){B}
\tkzDefPoint(1.5,0){C}
\tkzDefPoint(0,0){M}
\tkzDefPoint(-2.1,-2){D}
\tkzDefPoint(2.1,-2){E}
\tkzInterLL(A,M)(D,E)\tkzGetPoint{Z}
\draw[line width=0.3mm] (A)--(D)--(E)--(A);
\draw[line width=0.3mm] (B)--(C);
\draw[line width=0.3mm] (M)--(D);
\draw[line width=0.3mm] (M)--(E);
\draw[line width=0.3mm] (A)--(Z);
\tkzDrawPoints(A,B,C,D,E,M,Z)
\tkzLabelPoint[above](A){$A$}
\tkzLabelPoint[above left](B){$B$}
\tkzLabelPoint[above right](C){$\varGamma $}
\tkzLabelPoint[below left](D){$\varDelta  $}
\tkzLabelPoint[below right](E){$E$}
\tkzLabelPoint[above right](M){$M$}
\tkzLabelPoint[below](Z){$Z$}
\tkzMarkSegments[mark=s||,size=3pt](B,D C,E)
\end{tikzpicture}
\end{wrapfigure}
Δίνεται ισοσκελές τρίγωνο $AB\varGamma $ με $AB=A\varGamma $ και $M$ είναι το μέσο της βάσης $B\varGamma $. Στις προεκτάσεις των πλευρών $AB$, $A\varGamma $ παίρνουμε τα τμήματα $B\varDelta  $, $\varGamma E$ αντίστοιχα ώστε $B\varDelta  $=$\varGamma E$.

\begin{alist}
\item Να αποδείξετε ότι τα τρίγωνα Μ$B\varDelta  $ και $M\varGamma E$ είναι ίσα. \monades{12}
\item Να αποδείξετε ότι η γωνία $M\hat{\varDelta}  E$ είναι ίση με τη γωνία $M\hat{E}\varDelta$. \monades{6}
\item Αν η $AM$ τέμνει την $\varDelta  E$ στο σημείο $Z$ να αποδείξετε ότι η $AZ$ είναι κάθετη στην $\varDelta  E$.
\monades{7}
\end{alist}

\vspace{4mm}

% --- Άσκηση 12639 (12639.tex) ---
\Askhsh{12639}{2}
\begin{ylh}{1}
    \item 5.6. Εφαρμογές στα τρίγωνα.
\end{ylh}
% ===================================================================
%  Αρχείο: 12639.tex (Fragment)
%  Αυτόματη μετατροπή μέσω doc2latex.py
% ===================================================================



Από το μέσο $M$ της διαμέσου $A\varDelta  $ τριγώνου $AB\varGamma $, φέρουμε παράλληλη στην $AB$ που τέμνει την $A\varGamma $ στο σημείο $E$. Αν η παράλληλη από το $\varDelta  $ στην $AB$ τέμνει την $A\varGamma $ στο Ζ, να αποδείξετε ότι:

\begin{alist}
\item το $Z$ είναι μέσο της $A\varGamma $. \monades{10}

\item το $AE$ ισούται με το $\frac{1}{4}$ του $A\varGamma $. \monades{15}
\end{alist}

\vspace{4mm}

% --- Άσκηση 12640 (12640.tex) ---
\Askhsh{12640}{2}
\begin{ylh}{1}
    \item 4.6. Άθροισμα γωνιών τριγώνου
    \item 4.8. Άθροισμα γωνιών κυρτού ν-γώνου.
\end{ylh}
% ===================================================================
%  Αρχείο: 12640.tex (Fragment)
%  Αυτόματη μετατροπή μέσω doc2latex.py
% ===================================================================



Στο κυρτό πολύγωνο $AB\varGamma \varDelta  $Ε, οι διχοτόμοι των γωνιών του Α και Β τέμνονται στο Ο. Αν η γωνία του $\varGamma$ ισούται με 100\degree, η γωνία του $\varDelta  $ ισούται με 140\degree και η γωνία του $E$ ισούται με 80\degree τότε, να υπολογίσετε:

\begin{alist}
\item το μέτρο του αθροίσματος $\hat{A}$ + $\hat{B}$. \monades{12}

\item το μέτρο της γωνίας $AOB$. \monades{13}


\begin{center}
\begin{tikzpicture}
    \tkzDefPoint(0,0){A}
    \tkzDefPoint(4,0){B}
    \tkzDefPoint(5.464, 4.022){C}
    \tkzDefPoint(2.000, 6.022){D} 
    \tkzDefPoint(-1.939, 5.327){E} 
    \tkzDefLine[bisector](E,A,B) \tkzGetPoint{dirA}
    \tkzDefLine[bisector](A,B,C) \tkzGetPoint{dirB}
    \tkzInterLL(A,dirA)(B,dirB) \tkzGetPoint{O}
    \tkzDrawPolygon[l3](A,B,C,D,E)
    \tkzDrawSegments[l3, color=maincolor](A,O B,O)
    \tkzMarkAngle[size=0.6](D,C,B)
    \tkzLabelAngle[pos=0.95](D,C,B){\small $100^\circ$}
    \tkzMarkAngle[size=0.6](E,D,C)
    \tkzLabelAngle[pos=0.95](E,D,C){\small $140^\circ$}
    \tkzMarkAngle[size=0.6](A,E,D)
    \tkzLabelAngle[pos=0.95](A,E,D){\small $80^\circ$}
    \tkzMarkAngle[size=0.8, mark=||](B,A,O)
    \tkzMarkAngle[size=0.7, mark=||](O,A,E)
    \tkzMarkAngle[size=0.7, mark=|](C,B,O)
    \tkzMarkAngle[size=0.8, mark=|](O,B,A)
    \tkzDrawPoints(A,B,C,D,E,O)
    \tkzLabelPoint[below left](A){A}
    \tkzLabelPoint[below right](B){B}
    \tkzLabelPoint[right](C){$\Gamma$}
    \tkzLabelPoint[above](D){$\Delta$}
    \tkzLabelPoint[left](E){E}
    \tkzLabelPoint[above](O){O}
\end{tikzpicture}
\end{center}
\end{alist}

\vspace{4mm}

% --- Άσκηση 12644 (12644.tex) ---
\Askhsh{12644}{2}
\begin{ylh}{1}
    \item 4.6. Άθροισμα γωνιών τριγώνου
    \item 4.8. Άθροισμα γωνιών κυρτού ν-γώνου.
\end{ylh}

Στο τετράπλευρο $AB\varGamma \varDelta  $, η εξωτερική γωνία της $\hat{\varGamma}$, ισούται με $130\degree$ και η εξωτερική γωνία της $\varDelta$ ισούται με $110\degree$. Αν οι διχοτόμοι των γωνιών $\hat{A}$ και $\hat{B}$ τέμνονται στο $O$ τότε, να υπολογίσετε:
\begin{alist}
\item τα μέτρα των γωνιών $\varGamma$ και $\varDelta$ του τετραπλεύρου. \monades{9}
\item το μέτρο του αθροίσματος $\hat{A}$ + $\hat{B}$. \monades{9}
\item το μέτρο της γωνίας $AOB$. \monades{7}
\begin{center}
\begin{tikzpicture}[scale=0.9]
\tkzDefPoint(0,0){Dext_pt}
\tkzDefPoint(10,0){Gext_pt}
\tkzDefPoint(2,0){D}
\tkzDefPoint(8,0){C}
\tkzDefShiftPoint[D](70:3.5){A}
\tkzDefShiftPoint[C](130:3.5){B}
\tkzDefLine[bisector](D,A,B)\tkzGetPoint{o1}
\tkzDefLine[bisector](A,B,C)\tkzGetPoint{o2}
\tkzInterLL(A,o1)(B,o2)\tkzGetPoint{O}
\draw[line width=0.3mm] (Dext_pt)--(Gext_pt);
\draw[line width=0.3mm] (D)--(A)--(B)--(C);
\draw[line width=0.3mm] (A)--(O)--(B);
\tkzDrawPoints(A,B,C,D,O)
\tkzLabelPoint[above left](A){$A$}
\tkzLabelPoint[above right](B){$B$}
\tkzLabelPoint[below](C){$\varGamma $}
\tkzLabelPoint[below](D){$\varDelta  $}
\tkzLabelPoint[below right](O){$O$}
\tkzMarkAngle[size=0.6,mark=none,fill=gray!20](A,D,Dext_pt)
\tkzLabelAngle[pos=0.9](A,D,Dext_pt){$110^\circ$}
\tkzMarkAngle[size=0.6,mark=none,fill=gray!20](Gext_pt,C,B)
\tkzLabelAngle[pos=0.9](Gext_pt,C,B){$130^\circ$}
\end{tikzpicture}
\end{center}
\end{alist}

\vspace{4mm}

% --- Άσκηση 12704 (12704.tex) ---
\Askhsh{12704}{2}
\begin{ylh}{1}
    \item 3.1. Είδη και στοιχεία τριγώνων
    \item 3.2. 1ο Κριτήριο ισότητας τριγώνων
    \item 4.6. Άθροισμα γωνιών τριγώνου.
\end{ylh}

Δίνεται τρίγωνο $AB\varGamma $ με $\hat{B}=70\degree$ και $\hat{\varGamma}_{\text{εξ}}=140\degree$. Στην πλευρά $B\varGamma $ θεωρούμε εσωτερικό σημείο $\varDelta  $, ώστε $A\varDelta=AB$. Να αποδείξετε ότι:

\begin{alist}
\item $B\hat{A}\varDelta  =40\degree$ . \monades{9}
\item $Α\hat{\varDelta  }\varGamma  = 110^{0}$. \monades{7}
\item Το τρίγωνο $AB\varGamma $ είναι ισοσκελές. \monades{9}

\begin{center}
\begin{tikzpicture}
\tkzDefPoint(0,0){B}
\tkzDefPoint(6,0){C}
\tkzDefTriangle[two angles=70 and 40](B,C)\tkzGetPoint{A}
\tkzInterLC(B,C)(A,B)\tkzGetSecondPoint{D}
\tkzDefPoint(8, 0){Gext_pt}
\draw[line width=0.3mm] (A)--(B)--(Gext_pt);
\draw[line width=0.3mm] (A)--(C);
\draw[line width=0.3mm] (A)--(D);
\tkzDrawPoints(A,B,C,D)
\tkzLabelPoint[above](A){$A$}
\tkzLabelPoint[below](B){$B$}
\tkzLabelPoint[below](C){$\varGamma $}
\tkzLabelPoint[below](D){$\varDelta  $}
\tkzMarkAngle[size=0.6,mark=none,fill=gray!20](Gext_pt,C,A)
\tkzLabelAngle[pos=0.9](Gext_pt,C,A){$140^\circ$}
\tkzMarkAngle[size=0.6,mark=none,fill=gray!20](C,B,A)
\tkzLabelAngle[pos=0.9](C,B,A){$70^\circ$}
\end{tikzpicture}
\end{center}
\end{alist}

\vspace{4mm}

% --- Άσκηση 12705 (12705.tex) ---
\Askhsh{12705}{2}
\begin{ylh}{1}
    \item 3.1. Είδη και στοιχεία τριγώνων
    \item 3.2. 1ο Κριτήριο ισότητας τριγώνων.
\end{ylh}
% ===================================================================
%  Αρχείο: 12705.tex (Fragment)
%  Αυτόματη μετατροπή μέσω doc2latex.py
% ===================================================================

Δίνεται τρίγωνο $AB\varGamma $ τέτοιο, ώστε $A\varGamma =2AB$. Η διχοτόμος του $A\varDelta  $ τέμνει την διάμεσο $BE$ στο σημείο $Z$. Να αποδείξετε ότι:
\begin{alist}
\item $AB = AE = \frac{A\varGamma }{2}$. \monades{7}
\item $\varDelta   B$ = $\varDelta  E$. \monades{8}
\item $A Z \perp B E$ \monades{10}

\begin{center}
\begin{tikzpicture}
\tkzDefPoint(0,5){A}
\tkzDefPoint(-2,0){B}
\tkzDefPoint(4,0){C}
\tkzDefMidPoint(A,C)\tkzGetPoint{E}
\tkzDefLine[bisector](B,A,C)\tkzGetPoint{a_dir}
\tkzInterLL(A,a_dir)(B,C)\tkzGetPoint{D}
\tkzInterLL(A,D)(B,E)\tkzGetPoint{Z}
\draw[line width=0.3mm] (A)--(B)--(C)--cycle;
\draw[line width=0.3mm] (A)--(D);
\draw[line width=0.3mm] (B)--(E);
\tkzDrawPoints(A,B,C,D,E,Z)
\tkzLabelPoint[above](A){$A$}
\tkzLabelPoint[left](B){$B$}
\tkzLabelPoint[right](C){$\varGamma $}
\tkzLabelPoint[below](D){$\varDelta  $}
\tkzLabelPoint[right](E){$E$}
\tkzLabelPoint[below right](Z){$Z$}
\tkzMarkAngle[size=0.6,mark=none](B,A,D)
\tkzLabelAngle[pos=0.8](B,A,D){$1$}
\tkzMarkAngle[size=0.6,mark=none](D,A,E)
\tkzLabelAngle[pos=0.8](D,A,E){$2$}
\end{tikzpicture}
\end{center}
\end{alist}

\vspace{4mm}

% --- Άσκηση 12707 (12707.tex) ---
\Askhsh{12707}{2}
\begin{ylh}{1}
    \item 3.2. 1ο Κριτήριο ισότητας τριγώνων
    \item 4.6. Άθροισμα γωνιών τριγώνου.
\end{ylh}
Δίνεται το τρίγωνο $AB\varGamma $ με $\hat{A}=70\degree$ και $\ \hat{\varGamma }=55\degree$. Προεκτείνουμε την πλευρά $BA$ προς το σημείο $A$ και παίρνουμε στην προέκταση σημείο $Z$ ώστε $\ B\hat{Z}\varDelta=35\degree$, όπου $\varDelta  $ εσωτερικό σημείο της $B\varGamma $. Η $Z\varDelta  $ τέμνει την $A\varGamma $ στο σημείο $E$. Να αποδείξετε ότι:

\begin{alist}
\item το τρίγωνο $AB\varGamma $ είναι ισοσκελές. \monades{7}

\item $Z\hat{\varDelta  }B\ $= 90$\degree$. \monades{8}

\item το τρίγωνο $AZE$ είναι ισοσκελές. \monades{10}


\begin{center}
\begin{tikzpicture}[scale=1.2]
\tkzDefPoint(0,0){B}
\tkzDefPoint(4,0){C}
\tkzDefTriangle[two angles = 55 and 55](B,C) \tkzGetPoint{A}
\tkzDefPointOnLine[pos=1.5](B,A)\tkzGetPoint{Z} 
\tkzDefPointBy[projection=onto B--C](Z) \tkzGetPoint{D}
\tkzInterLL(Z,D)(A,C) \tkzGetPoint{E}

\draw[l3] (A)--(B)--(C)--cycle;
\draw[l3] (A)--(Z);
\draw[l3] (Z)--(D);
\tkzDrawPoints(A,B,C,D,E,Z)
\tkzLabelPoint[above left](A){$A$}
\tkzLabelPoint[left](B){$B$}
\tkzLabelPoint[right](C){$\varGamma $}
\tkzLabelPoint[below](D){$\varDelta  $}
\tkzLabelPoint[right](E){$E$}
\tkzLabelPoint[above](Z){$Z$}
\tkzMarkAngle[size=0.6,mark=none](B,A,C) \tkzLabelAngle[pos=0.8](B,A,C){$70^\circ$}
\tkzMarkAngle[size=0.6,mark=none](A,C,B) \tkzLabelAngle[pos=0.8](A,C,B){$55^\circ$}
\tkzMarkAngle[size=0.6,mark=none](B,Z,D) \tkzLabelAngle[pos=0.8](B,Z,D){$35^\circ$}
\tkzMarkRightAngle(B,D,Z)
\end{tikzpicture}
\end{center}
\end{alist}

\vspace{4mm}

% --- Άσκηση 12708 (12708.tex) ---
\Askhsh{12708}{2}
\begin{ylh}{1}
    \item 3.2. 1ο Κριτήριο ισότητας τριγώνων
    \item 4.6. Άθροισμα γωνιών τριγώνου.
\end{ylh}

Δίνεται το ισόπλευρο τρίγωνο $AB\varGamma $. Στις πλευρές $B\varGamma $ και $\varGamma A$ θεωρούμε σημεία $E$ και $Z$ αντίστοιχα ώστε $BE$ = $\varGamma Z$. Στις πλευρές $AB$ και $\varGamma  B$ θεωρούμε σημεία $H$ και $\varDelta $ αντίστοιχα ώστε $BH$ = $\varGamma \varDelta  $. Τα ευθύγραμμα τμήματα $\varDelta   Z$ και $EH$ τέμνονται στο σημείο $K$ το οποίο είναι εσωτερικό σημείο του τριγώνου $AB\varGamma $. Να αποδείξετε ότι:

\begin{alist}
\item $EH$ = $\varDelta Z$ και $B\hat{H}E$ = $\varGamma \hat{\varDelta }Z$. \monades{12}
\item τα τρίγωνα $BEH$ και $KE\varDelta  $ έχουν ίσες γωνίες μία προς μία. \monades{13}

\begin{center}
\begin{tikzpicture}[scale=1.2]
\tkzDefPoint(0,0){B}
\tkzDefPoint(5,0){C}
\tkzDefTriangle[equilateral](B,C) \tkzGetPoint{A}
\tkzDefPointOnLine[pos=0.7](B,C) \tkzGetPoint{E}
\tkzDefPointOnLine[pos=0.7](C,A) \tkzGetPoint{Z}
\tkzDefPointOnLine[pos=0.6](B,A) \tkzGetPoint{H}
\tkzDefPointOnLine[pos=0.6](C,B) \tkzGetPoint{D}
\tkzInterLL(D,Z)(E,H) \tkzGetPoint{K}

\draw[l3] (A)--(B)--(C)--cycle;
\draw[l3] (E)--(H);
\draw[l3] (D)--(Z);
\tkzDrawPoints(A,B,C,D,E,Z,H,K)
\tkzLabelPoint[above](A){$A$}
\tkzLabelPoint[left](B){$B$}
\tkzLabelPoint[right](C){$\varGamma $}
\tkzLabelPoint[below](D){$\varDelta  $}
\tkzLabelPoint[below](E){$E$}
\tkzLabelPoint[right](Z){$Z$}
\tkzLabelPoint[left](H){$H$}
\tkzLabelPoint[right](K){$K$}
\end{tikzpicture}
\end{center}
\end{alist}

\vspace{4mm}

% --- Άσκηση 12709 (12709.tex) ---
\Askhsh{12709}{2}
\begin{ylh}{1}
    \item 3.2. 1ο Κριτήριο ισότητας τριγώνων
    \item 4.6. Άθροισμα γωνιών τριγώνου.
\end{ylh}
Δίνεται το ορθογώνιο και ισοσκελές τρίγωνο $AB\varGamma $ με $AB=A\varGamma $ και $\hat{A}=90\degree$. Εξωτερικά του τριγώνου $AB\varGamma $ κατασκευάζουμε το ισόπλευρο τρίγωνο $A\varGamma \varDelta  $.
\begin{alist}
\item Nα υπολογίσετε το μέτρο των γωνιών $\hat{B},\hat{\varGamma}$ του τριγώνου $AB\varGamma $. \monades{5}
\item Να αποδείξετε ότι το τρίγωνο $AB\varDelta$ είναι ισοσκελές. \monades{8}
\item Να υπολογίσετε το μέτρο της γωνίας $A\hat{B}\varDelta  $. \monades{12}
\begin{center}
\begin{tikzpicture}[scale=0.7]
\tkzDefPoint(0,0){A}
\tkzDefPoint(0,3.5){B}
\tkzDefPoint(3.5,0){C}
\tkzDefTriangle[equilateral](C,A) \tkzGetPoint{D}
\draw[l3] (A)--(B)--(C)--cycle;
\draw[l3] (A)--(D)--(C);
\draw[l3] (B)--(D);
\tkzDrawPoints(A,B,C,D)
\tkzLabelPoint[below left](A){$A$}
\tkzLabelPoint[left](B){$B$}
\tkzLabelPoint[right](C){$\varGamma $}
\tkzLabelPoint[right](D){$\varDelta  $}
\tkzMarkRightAngle[size=0.3](B,A,C)
\end{tikzpicture}
\end{center}
\end{alist}

\vspace{4mm}

% --- Άσκηση 12710 (12710.tex) ---
\Askhsh{12710}{2}
\begin{ylh}{1}
    \item 3.2. 1ο Κριτήριο ισότητας τριγώνων
    \item 4.2. Τέμνουσα δύο ευθειών - Ευκλείδειο αίτημα.
\end{ylh}
Δίνεται το ισόπλευρο τρίγωνο $AB\varGamma $ και η διχοτόμος του $BE$. Εξωτερικά του τριγώνου $AB\varGamma $ κατασκευάζουμε το ορθογώνιο και ισοσκελές τρίγωνο $A\varGamma \varDelta  $ με υποτείνουσα τη $\varGamma \varDelta  $ έτσι, ώστε τα σημεία $B$ και $\varDelta  $ να βρίσκονται εκατέρωθεν της ευθείας $A\varGamma $. Να αποδείξετε ότι:
\begin{alist}
\item $BE\parallel A\varDelta  $. \monades{10}
\item οι γωνίες $E\hat{B}\varDelta $ και $A\hat{\varDelta} B$ είναι ίσες. \monades{7}
\item το τρίγωνο $BA\varDelta $ είναι ισοσκελές. \monades{8}

\begin{center}
\begin{tikzpicture}[scale=0.7]
    \tkzDefPoint(0,0){B}
    \tkzDefPoint(4,0){C}
    \tkzDefTriangle[equilateral](B,C) \tkzGetPoint{A}
    \tkzDefMidPoint(A,C) \tkzGetPoint{E}
    \tkzDefPointBy[rotation=center A angle 90](C) \tkzGetPoint{D}
    \tkzDrawPolygon[thick](A,B,C)
    \tkzDrawSegment[thick](B,E)
    \tkzDrawPolygon[thick](A,C,D)
    \tkzMarkRightAngle[size=0.3](C,A,D)
    \tkzMarkRightAngle[size=0.25](A,E,B)
    \tkzMarkSegments[mark=|, size=3pt](A,B B,C C,A A,D)
    \tkzDrawPoints(A,B,C,D,E)
    \tkzLabelPoint[above=2pt](A){$A$}
    \tkzLabelPoint[below left](B){$B$}
    \tkzLabelPoint[below right](C){$\varGamma$}
    \tkzLabelPoint[right](D){$\varDelta$}
    \tkzLabelPoint[above right](E){$E$}
\end{tikzpicture}
\end{center}
\end{alist}

\vspace{4mm}

% --- Άσκηση 13442 (13442.tex) ---
\Askhsh{13442}{2}
\begin{ylh}{1}
    \item 3.1. Είδη και στοιχεία τριγώνων
    \item 3.2. 1ο Κριτήριο ισότητας τριγώνων
    \item 4.6. Άθροισμα γωνιών τριγώνου.
\end{ylh}
% ===================================================================
%  Αρχείο: 13442.tex (Fragment)
%  Αυτόματη μετατροπή μέσω doc2latex.py
% ===================================================================



Δίνεται ορθογώνιο τρίγωνο $AB\varGamma $ με $\hat{A} = 90\degree$. Στην πλευρά του $AB$ θεωρούμε σημείο $\varDelta  $ ώστε $B\varDelta   = \varDelta  \varGamma $ και $A\varDelta   = A\varGamma $.

\begin{alist}
\item Να αποδείξετε ότι $A\hat{\varDelta  }\varGamma  = 45\degree$ . \monades{12}
\item Να υπολογίσετε τη γωνία $\hat{B}$. \monades{13}
\begin{center}
\begin{tikzpicture}[scale=1.2,rotate=90]
\tkzDefPoint(0,0){A}
\tkzDefPoint(2.5,0){C}
\tkzDefPoint(0,2.5){D}
\tkzDefPoint(0, {2.5 + 2.5*sqrt(2)}){B}
\draw[l3] (A)--(B)--(C)--cycle;
\draw[l3] (D)--(C);
\tkzDrawPoints(A,B,C,D)
\tkzLabelPoint[right](A){$A$}
\tkzLabelPoint[left](B){$B$}
\tkzLabelPoint[right](C){$\varGamma $}
\tkzLabelPoint[below](D){$\varDelta  $}
\tkzMarkRightAngle(B,A,C)
\tkzMarkSegments[mark=s|](A,D A,C)
\tkzMarkSegments[mark=s||](B,D D,C)
\end{tikzpicture}
\end{center}
\end{alist}

\vspace{4mm}

% --- Άσκηση 13443 (13443.tex) ---
\Askhsh{13443}{2}
\begin{ylh}{1}
    \item 3.1. Είδη και στοιχεία τριγώνων
    \item 3.11. Ανισοτικές σχέσεις πλευρών και γωνιών
    \item 4.2. Τέμνουσα δύο ευθειών - Ευκλείδειο αίτημα
    \item 4.6. Άθροισμα γωνιών τριγώνου.
\end{ylh}
% ===================================================================
%  Αρχείο: 13443.tex (Fragment)
%  Αυτόματη μετατροπή μέσω doc2latex.py
% ===================================================================



Δίνεται τρίγωνο $AB\varGamma $ με $\hat{A} = 60\degree$ και $\hat{\varGamma } = 40\degree$. Στην πλευρά $A\varGamma $ θεωρούμε σημείο $\varDelta  $, ώστε $\varGamma \hat{B}\varDelta   = 20^{ο}$.

\begin{alist}
\item Να αποδείξετε ότι το τρίγωνο Α$B\varDelta  $ είναι ισόπλευρο. \monades{10}

\item Η παράλληλη από το $\varDelta  $ προς την $AB$ τέμνει την πλευρά $B\varGamma $ στο σημείο $E$. Να αποδείξετε ότι:

\begin{rlist}
\item
  $B\hat{\varDelta  }E = 60\degree$. \monades{8}
\item
  Η $\varDelta  E$ είναι διχοτόμος της γωνίας $B\hat{\varDelta  }\varGamma $. \monades{7}
\begin{center}
\begin{tikzpicture}[scale=1.5]
\tkzDefPoint(0,0){B}
\tkzDefPoint(4,0){C}
\tkzDefTriangle[two angles = 80 and 40](B,C) \tkzGetPoint{A}
\tkzDefPointBy[rotation=center B angle 20](C) \tkzGetPoint{d_dir}
\tkzInterLL(B,d_dir)(A,C) \tkzGetPoint{D}
\tkzDefLine[parallel=through D](A,B) \tkzGetPoint{e_dir}
\tkzInterLL(D,e_dir)(B,C) \tkzGetPoint{E}

\draw[l3] (A)--(B)--(C)--cycle;
\draw[l3] (B)--(D);
\draw[l3] (D)--(E);
\tkzDrawPoints(A,B,C,D,E)
\tkzLabelPoint[above](A){$A$}
\tkzLabelPoint[left](B){$B$}
\tkzLabelPoint[right](C){$\varGamma $}
\tkzLabelPoint[above right](D){$\varDelta  $}
\tkzLabelPoint[below](E){$E$}
\tkzMarkAngle[size=0.7,mark=none](C,B,D) \tkzLabelAngle[pos=0.9](C,B,D){$20^\circ$}
\end{tikzpicture}
\end{center}
\end{rlist}
\end{alist}
\vspace{4mm}

% --- Άσκηση 13497 (13497.tex) ---
\Askhsh{13497}{2}
\begin{ylh}{1}
    \item 3.2. 1ο Κριτήριο ισότητας τριγώνων
    \item 4.2. Τέμνουσα δύο ευθειών - Ευκλείδειο αίτημα
    \item 5.10. Τραπέζιο
    \item 5.11. Ισοσκελές τραπέζιο.
\end{ylh}
% ===================================================================
%  Αρχείο: 13497.tex (Fragment)
%  Αυτόματη μετατροπή μέσω doc2latex.py
% ===================================================================



Δίνεται ισοσκελές τραπέζιο $AB\varGamma \varDelta  $ ($A\varDelta  $//$B\varGamma $) με $B\varGamma $ > $\varDelta  \varGamma $.

Στην πλευρά $B\varGamma $ θεωρούμε σημείο $E$, τέτοιο ώστε $\varGamma E$ = $\varGamma \varDelta  $.

\begin{alist}
\item Να αποδείξετε ότι η $\varDelta  E$ είναι διχοτόμος της $A\hat{\varDelta  }\varGamma $. \monades{12}

\item Αν $\hat{A}=\text{120}\degree$, να αποδείξετε ότι το τρίγωνο $\varDelta  E\varGamma $ είναι ισόπλευρο. \monades{13}


\begin{center}
\begin{tikzpicture}[scale=1.2]
\tkzDefPoint(0,0){B}
\tkzDefPoint(6,0){C}
\tkzDefPoint(1.5,2.598){A}
\tkzDefPoint(4.5,2.598){D}
\tkzInterLC(B,C)(C,D) \tkzGetPoints{E}{e}
\draw[l3] (A)--(B)--(C)--(D)--cycle;
\draw[l3, maincolor] (D)--(E);
\tkzDrawPoints(A,B,C,D,E)
\tkzLabelPoint[above](A){$A$}
\tkzLabelPoint[below](B){$B$}
\tkzLabelPoint[below](C){$\varGamma $}
\tkzLabelPoint[above](D){$\varDelta  $}
\tkzLabelPoint[below](E){$E$}
\tkzMarkSegments[mark=s||](A,B D,C C,E)
\end{tikzpicture}
\end{center}
\end{alist}

\vspace{4mm}

% --- Άσκηση 13499 (13499.tex) ---
\Askhsh{13499}{4}
\begin{ylh}{1}
    \item 3.1. Είδη και στοιχεία τριγώνων
    \item 3.2. 1ο Κριτήριο ισότητας τριγώνων
    \item 3.5. Ύπαρξη και μοναδικότητα καθέτου
    \item 3.6. Κριτήρια ισότητας ορθογώνιων τριγώνων
    \item 3.7. Κύκλος - Μεσοκάθετος - Διχοτόμος
    \item 4.6. Άθροισμα γωνιών τριγώνου.
\end{ylh}

Δίνεται ορθογώνιο τρίγωνο $AB\varGamma $ ($\hat{A} = \ $90\textsuperscript{ο}) με $AB < A\varGamma $ και $AH$ το ύψος προς την υποτείνουσα. Στην πλευρά $B\varGamma $ θεωρούμε τα σημεία $\varDelta  $ και $E$ τέτοια ώστε $\varDelta   B=AB$ και $\varGamma  E = \varGamma  A$. Αν $\varDelta   Z$ και $E\varTheta $ είναι οι αποστάσεις των $\varDelta  $ και $E$ από τις πλευρές $A\varGamma $ και $AB$ αντίστοιχα, να αποδείξετε ότι:

\begin{alist}
\item $\varGamma\hat{H}\varDelta=\varDelta\hat{A}H$ και $E\hat{A}B= H\hat{A}E$.\monades{14}
\item $\varDelta E=\varDelta Z+E\varTheta$.
\monades{11}
\end{alist}
\vspace{4mm}

% --- Άσκηση 13517 (13517.tex) ---
\Askhsh{13517}{2}
\begin{ylh}{1}
    \item 3.1. Είδη και στοιχεία τριγώνων
    \item 3.6. Κριτήρια ισότητας ορθογώνιων τριγώνων.
\end{ylh}
% ===================================================================
%  Αρχείο: 13517.tex (Fragment)
%  Αυτόματη μετατροπή μέσω doc2latex.py
% ===================================================================



Δίνονται δύο οξυγώνια τρίγωνα $AB\varGamma $ και $\varDelta  EZ$ με $\hat{A}=\hat{\text{Δ}},\ \ A\hat{B}\varGamma  = \varDelta  \hat{Ε}Ζ$. Αν τα ύψη τους $BH$ και $E\varTheta $ είναι ίσα τότε να αποδείξετε ότι:

\begin{alist}
\item $AB=\varDelta   E$. \monades{13}

\item Τα τρίγωνα $AB\varGamma $ και $\varDelta  EZ$ είναι ίσα . \monades{12}


\begin{center}
\begin{tikzpicture}[scale=1]
\tkzDefPoint(0,0){B}
\tkzDefPoint(4,0){C}
\tkzDefTriangle[two angles = 70 and 50](B,C) \tkzGetPoint{A}
\tkzDefPointBy[projection=onto A--C](B) \tkzGetPoint{H}
\draw[l3] (A)--(B)--(C)--cycle;
\draw[l3, maincolor] (B)--(H);
\tkzDrawPoints(A,B,C,H)
\tkzLabelPoint[above](A){$A$}
\tkzLabelPoint[left](B){$B$}
\tkzLabelPoint[right](C){$\varGamma $}
\tkzLabelPoint[right](H){$H$}
\tkzMarkRightAngle(B,H,A)

\begin{scope}[xshift=6cm]
\tkzDefPoint(0,0){E}
\tkzDefPoint(4,0){Z}
\tkzDefTriangle[two angles = 70 and 50](E,Z) \tkzGetPoint{D}
\tkzDefPointBy[projection=onto D--Z](E) \tkzGetPoint{T}
\draw[l3] (D)--(E)--(Z)--cycle;
\draw[l3, maincolor] (E)--(T);
\tkzDrawPoints(D,E,Z,T)
\tkzLabelPoint[above](D){$\varDelta  $}
\tkzLabelPoint[left](E){$E$}
\tkzLabelPoint[right](Z){$Z$}
\tkzLabelPoint[right](T){$\varTheta $}
\tkzMarkRightAngle(E,T,D)
\end{scope}
\end{tikzpicture}
\end{center}
\end{alist}

\vspace{4mm}

% --- Άσκηση 13518 (13518.tex) ---
\Askhsh{13518}{2}
\begin{ylh}{1}
    \item 3.2. 1ο Κριτήριο ισότητας τριγώνων
    \item 3.4. 3ο Κριτήριο ισότητας τριγώνων.
\end{ylh}
% ===================================================================
%  Αρχείο: 13518.tex (Fragment)
%  Αυτόματη μετατροπή μέσω doc2latex.py
% ===================================================================



Δίνονται τα τρίγωνα $AB\varGamma $ και $A'B'\varGamma  '$ του σχήματος με $A\varGamma =A'\varGamma '$ και $AB=A'B'$. Αν οι διάμεσοι $B\varDelta  $ και Β΄Δ΄ είναι ίσες, να αποδείξετε ότι:

\begin{alist}
\item $\hat{A}=\hat{A}\text{΄}$ \monades{15}

\item Τα τρίγωνα $AB\varGamma $ και $A'B'\varGamma  '$ είναι ίσα. \monades{10}


\begin{center}
\begin{tikzpicture}[scale=1]
\tkzDefPoint(0,0){B}
\tkzDefPoint(4,0){C}
\tkzDefTriangle[two angles = 60 and 40](B,C) \tkzGetPoint{A}
\tkzDefMidPoint(A,C)\tkzGetPoint{D}
\draw[l3] (A)--(B)--(C)--cycle;
\draw[l3] (B)--(D);
\tkzDrawPoints(A,B,C,D)
\tkzLabelPoint[above](A){$A$}
\tkzLabelPoint[left](B){$B$}
\tkzLabelPoint[right](C){$\varGamma $}
\tkzLabelPoint[right](D){$\varDelta  $}

\begin{scope}[xshift=5.5cm]
\tkzDefPoint(0,0){B'}
\tkzDefPoint(4,0){C'}
\tkzDefTriangle[two angles = 60 and 40](B',C') \tkzGetPoint{A'}
\tkzDefMidPoint(A',C')\tkzGetPoint{D'}
\draw[l3] (A')--(B')--(C')--cycle;
\draw[l3] (B')--(D');
\tkzDrawPoints(A',B',C',D')
\tkzLabelPoint[above](A'){$A'$}
\tkzLabelPoint[left](B'){$B'$}
\tkzLabelPoint[right](C'){$\varGamma'$}
\tkzLabelPoint[right](D'){$\varDelta'$}
\end{scope}
\end{tikzpicture}
\end{center}
\end{alist}

\vspace{4mm}

% --- Άσκηση 13519 (13519.tex) ---
\Askhsh{13519}{4}
\begin{ylh}{1}
    \item 3.2. 1ο Κριτήριο ισότητας τριγώνων
    \item 4.2. Τέμνουσα δύο ευθειών - Ευκλείδειο αίτημα
    \item 5.6. Εφαρμογές στα τρίγωνα
    \item 5.10. Τραπέζιο.
\end{ylh}
% ===================================================================
%  Αρχείο: 13519.tex (Fragment)
%  Αυτόματη μετατροπή μέσω doc2latex.py
% ===================================================================



Δίνεται ορθογώνιο παραλληλόγραμμο $AB\varGamma \varDelta  $ με $AB>A\varDelta  $. Στην $AB$ θεωρούμε σημείο $E$ τέτοιο, ώστε $AE=A\varDelta  $. Από το μέσο $M$ της $\varDelta  E$ φέρουμε παράλληλη προς την $\varDelta  \varGamma $ που τέμνει την $B\varGamma $ στο Κ.

\begin{alist}
\item Να αποδείξετε $AM\perp\varDelta   E$. \monades{7}

\item Να αποδείξετε ότι $2MK=2AB-A\varDelta  $. \monades{9}

\item Φέρνουμε την $EK$ που τέμνει την προέκταση της $\varDelta  \varGamma $ στο $Z$.

Να αποδείξετε ότι $\varGamma  Z=$AB$-A\varDelta  $. \monades{9}
\end{alist}

\vspace{4mm}

% --- Άσκηση 13520 (13520.tex) ---
\Askhsh{13520}{4}
\begin{ylh}{1}
    \item 3.7. Κύκλος - Μεσοκάθετος - Διχοτόμος
    \item 3.15. Εφαπτόμενα τμήματα
    \item 4.6. Άθροισμα γωνιών τριγώνου
    \item 5.4. Ρόμβος
    \item 5.9. Μια ιδιότητα του ορθογώνιου τριγώνου.
\end{ylh}
% ===================================================================
%  Αρχείο: 13520.tex (Fragment)
%  Αυτόματη μετατροπή μέσω doc2latex.py
% ===================================================================



Δίνεται κύκλος (Ο,ρ) και σημείο $P$ εκτός του κύκλου. Από το Ρ φέρνουμε τα εφαπτόμενα τμήματα $PA$ και $PB$. Η $PO$ τέμνει το μικρότερο του ημικυκλίου τόξο $AB$ στο $\varGamma$ και $A\hat{P}B=60\degree$. Να αποδείξετε ότι:
\begin{alist}
\item $OP=2\rho$. \monades{10}
\item $A\hat{\varGamma }B=120\degree$. \monades{10}
\item Ένας μαθητής ισχυρίζεται ότι το τετράπλευρο $OA\varGamma  B$ είναι ρόμβος. Συμφωνείτε μαζί του; Να δικαιολογήσετε την απάντηση σας. \monades{05}


\begin{center}
\begin{tikzpicture}[scale=1]
\tkzDefPoint(0,0){O}
\tkzDefPoint(3,0){P}
\tkzDefPoint(1.5,0){temp} 
\tkzDrawCircle[l3](O,temp)
\tkzDefTangent[from=P](O,temp) \tkzGetPoints{A}{B}
\tkzInterLC(O,P)(O,temp) \tkzGetPoints{G_fake}{C}
\draw[l4,maincolor] (P)--(A);
\draw[l4,maincolor] (P)--(B);
\draw[l3] (P)--(O);
\draw[l3] (O)--(A);
\draw[l3] (O)--(B);
\tkzDrawPoints(O,P,A,B,C)
\tkzLabelPoint[left](O){$O$}
\tkzLabelPoint[right](P){$P$}
\tkzLabelPoint[below](A){$A$}
\tkzLabelPoint[above](B){$B$}
\tkzLabelPoint[above left](C){$\varGamma $}
\tkzMarkRightAngle(O,A,P)
\tkzMarkRightAngle(O,B,P)
\end{tikzpicture}
\end{center}
\end{alist}

\vspace{4mm}

% --- Άσκηση 13522 (13522.tex) ---
\Askhsh{13522}{4}
\begin{ylh}{1}
    \item 3.6. Κριτήρια ισότητας ορθογώνιων τριγώνων
    \item 3.7. Κύκλος - Μεσοκάθετος - Διχοτόμος
    \item 5.9. Μια ιδιότητα του ορθογώνιου τριγώνου.
\end{ylh}
% ===================================================================
%  Αρχείο: 13522.tex (Fragment)
%  Αυτόματη μετατροπή μέσω doc2latex.py
% ===================================================================



Δίνεται τρίγωνο $AB\varGamma $ με $AB$ < $A\varGamma $. Η διχοτόμος της γωνίας Α τέμνει την μεσοκάθετο (ε) της $B\varGamma $ στο $\varDelta  $. Από το $\varDelta  $ φέρνουμε τα κάθετα τμήματα $\varDelta   Z$ και $\varDelta  H$ προς τις $AB$ και $A\varGamma $ αντίστοιχα.

\begin{alist}
\item Να συγκρίνετε τα τρίγωνα $AZ\varDelta  $ και $AH\varDelta  $. \monades{08}

\item Να αποδείξετε ότι $BZ$ = $H\varGamma $. \monades{09}

\item Αν η γωνία $A$ = 60\textsuperscript{ο} και $M$ το μέσο ης $A\varDelta  $, να αποδείξετε ότι $HM$ = $Z\varDelta  $.

\monades{08}


\begin{center}
\begin{tikzpicture}[scale=1]
\tkzDefPoint(0,0){B}
\tkzDefPoint(6,0){C}
\tkzDefTriangle[two angles = 70 and 50](B,C) \tkzGetPoint{A}
\tkzDefLine[bisector](B,A,C) \tkzGetPoint{a_dir}
\tkzDefLine[mediator](B,C) \tkzGetPoints{m1}{m2}
\tkzInterLL(A,a_dir)(m1,m2) \tkzGetPoint{D}
\tkzDefPointBy[projection=onto A--B](D)\tkzGetPoint{Z}
\tkzDefPointBy[projection=onto A--C](D)\tkzGetPoint{H}

\draw[l3] (A)--(B)--(C)--cycle;
\draw[l3] (A)--(D);
\draw[l3, dashed, maincolor] (m1)--(m2);
\draw[l3] (D)--(Z);
\draw[l3] (D)--(H);
\tkzDrawPoints(A,B,C,D,Z,H)
\tkzLabelPoint[above](A){$A$}
\tkzLabelPoint[left](B){$B$}
\tkzLabelPoint[right](C){$\varGamma $}
\tkzLabelPoint[below](D){$\varDelta  $}
\tkzLabelPoint[above left](Z){$Z$}
\tkzLabelPoint[above right](H){$H$}
\tkzMarkRightAngle(A,Z,D)
\tkzMarkRightAngle(A,H,D)
\tkzMarkSegments[mark=s||](B,D C,D) % implicit from mediator
\end{tikzpicture}
\end{center}
\end{alist}

\vspace{4mm}

% --- Άσκηση 13532 (13532.tex) ---
\Askhsh{13532}{2}
\begin{ylh}{1}
    \item 5.2. Παραλληλόγραμμα
    \item 5.6. Εφαρμογές στα τρίγωνα.
\end{ylh}
% ===================================================================
%  Αρχείο: 13532.tex (Fragment)
%  Αυτόματη μετατροπή μέσω doc2latex.py
% ===================================================================



Δίνεται παραλληλόγραμμο $AB\varGamma \varDelta  $ και τα μέσα Ε, Ζ και Η των $AB$, $A\varGamma $ και $A\varDelta  $ αντίστοιχα. Να αποδείξετε ότι:

\begin{alist}
\item $ΖH = \frac{AB}{2}$. \monades{15}

\item Το τετράπλευρο $AEZH$ είναι παραλληλόγραμμο. \monades{10}


\begin{center}
\begin{tikzpicture}[scale=1.2]
\tkzDefPoint(0,0){B}
\tkzDefPoint(5,0){C}
\tkzDefPoint(1.5,3){A}
\tkzDefPointBy[translation=from B to C](A) \tkzGetPoint{D}
\tkzDefMidPoint(A,B) \tkzGetPoint{E}
\tkzDefMidPoint(A,C) \tkzGetPoint{Z}
\tkzDefMidPoint(A,D) \tkzGetPoint{H}

\draw[l3] (A)--(B)--(C)--(D)--cycle;
\draw[l3] (A)--(C);
\draw[l3, secondarycolor] (H)--(Z);
\draw[l3, secondarycolor] (A)--(E)--(Z)--(H)--cycle;
\tkzDrawPoints(A,B,C,D,E,Z,H)
\tkzLabelPoint[above left](A){$A$}
\tkzLabelPoint[below left](B){$B$}
\tkzLabelPoint[below right](C){$\varGamma $}
\tkzLabelPoint[above right](D){$\varDelta  $}
\tkzLabelPoint[left](E){$E$}
\tkzLabelPoint[above right](Z){$Z$}
\tkzLabelPoint[above](H){$H$}
\end{tikzpicture}
\end{center}
\end{alist}

\vspace{4mm}

% --- Άσκηση 13533 (13533.tex) ---
\Askhsh{13533}{2}
\begin{ylh}{1}
    \item 3.1. Είδη και στοιχεία τριγώνων
    \item 3.2. 1ο Κριτήριο ισότητας τριγώνων
    \item 3.6. Κριτήρια ισότητας ορθογώνιων τριγώνων.
\end{ylh}

Δίνεται ισοσκελές και αμβλυγώνιο τρίγωνο $AB\varGamma $ με $AB = A\varGamma $. Η κάθετη στην $AB$ στο σημείο $A$ τέμνει την πλευρά $B\varGamma $ στο σημείο $\varDelta  $ και η κάθετη στην $A\varGamma $ στο σημείο $A$ τέμνει την πλευρά $B\varGamma $ στο σημείο $E$. Να αποδείξετε ότι:

\begin{alist}
\item τα τρίγωνα Α$B\varDelta  $ και $A\varGamma E$ είναι ίσα. \monades{10}

\item το τρίγωνο $A\varDelta  E$ είναι ισοσκελές. \monades{7}

\item $ΒE = \varGamma \varDelta  $. \monades{8}


\begin{center}
\begin{tikzpicture}[scale=1]
\tkzDefPoint(0,0){B}
\tkzDefPoint(6,0){C}
\tkzDefTriangle[two angles = 30 and 30](B,C) \tkzGetPoint{A}
\tkzDefPointBy[rotation=center A angle 90](B) \tkzGetPoint{d_temp}
\tkzInterLL(A,d_temp)(B,C) \tkzGetPoint{D}
\tkzDefPointBy[rotation=center A angle -90](C) \tkzGetPoint{e_temp}
\tkzInterLL(A,e_temp)(B,C) \tkzGetPoint{E}

\draw[l3] (B)--(A)--(C)--cycle;
\draw[l3] (B)--(C);
\draw[l3, maincolor] (A)--(D);
\draw[l3, maincolor] (A)--(E);
\tkzDrawPoints(A,B,C,D,E)
\tkzLabelPoint[above](A){$A$}
\tkzLabelPoint[left](B){$B$}
\tkzLabelPoint[right](C){$\varGamma $}
\tkzLabelPoint[below](D){$\varDelta  $}
\tkzLabelPoint[below](E){$E$}
\tkzMarkRightAngle(B,A,D)
\tkzMarkRightAngle(C,A,E)
\end{tikzpicture}
\end{center}
\end{alist}

\vspace{4mm}

% --- Άσκηση 13534 (13534.tex) ---
\Askhsh{13534}{2}
\begin{ylh}{1}
    \item 3.1. Είδη και στοιχεία τριγώνων
    \item 3.2. 1ο Κριτήριο ισότητας τριγώνων
    \item 4.2. Τέμνουσα δύο ευθειών - Ευκλείδειο αίτημα.
\end{ylh}
% ===================================================================
%  Αρχείο: 13534.tex (Fragment)
%  Αυτόματη μετατροπή μέσω doc2latex.py
% ===================================================================



Δίνεται τρίγωνο $AB\varGamma $ με $AB < A\varGamma $. Η μεσοκάθετος της πλευράς $B\varGamma $ τέμνει την πλευρά $A\varGamma $ στο σημείο $\varDelta  $ και η παράλληλη από το $\varDelta  $ προς τη $B\varGamma $ τέμνει την πλευρά $AB$ στο σημείο $E$. Να αποδείξετε ότι:

\begin{alist}
\item το τρίγωνο $B\varGamma $Δ είναι ισοσκελές. \monades{12}

\item η $\varDelta  E$ είναι διχοτόμος της γωνίας $Α\hat{\varDelta  }Β$. \monades{13}


\begin{center}
\begin{tikzpicture}[scale=1]
\tkzDefPoint(0,0){B}
\tkzDefPoint(6,0){C}
\tkzDefTriangle[two angles = 70 and 50](B,C) \tkzGetPoint{A}
\tkzDefLine[mediator](B,C) \tkzGetPoints{m1}{m2}
\tkzInterLL(A,C)(m1,m2) \tkzGetPoint{D}
\tkzDefLine[parallel=through D](B,C) \tkzGetPoint{e_dir}
\tkzInterLL(D,e_dir)(A,B) \tkzGetPoint{E}

\draw[l3] (A)--(B)--(C)--cycle;
\draw[l3] (D)--(B);
\draw[l3] (D)--(E);
\draw[l3, dashed, maincolor] (m1)--(m2);
\tkzDrawPoints(A,B,C,D,E)
\tkzLabelPoint[above](A){$A$}
\tkzLabelPoint[below left](B){$B$}
\tkzLabelPoint[below right](C){$\varGamma $}
\tkzLabelPoint[above right](D){$\varDelta  $}
\tkzLabelPoint[above left](E){$E$}
\end{tikzpicture}
\end{center}
\end{alist}

\vspace{4mm}

% --- Άσκηση 13535 (13535.tex) ---
\Askhsh{13535}{2}
\begin{ylh}{1}
    \item 3.11. Ανισοτικές σχέσεις πλευρών και γωνιών
    \item 4.2. Τέμνουσα δύο ευθειών - Ευκλείδειο αίτημα
    \item 4.6. Άθροισμα γωνιών τριγώνου.
\end{ylh}
% ===================================================================
%  Αρχείο: 13535.tex (Fragment)
%  Αυτόματη μετατροπή μέσω doc2latex.py
% ===================================================================



Στο παρακάτω σχήμα η ευθεία $\varepsilon_{1}$ διέρχεται από την κορυφή $\varGamma $ του τριγώνου $AB\varGamma $ και είναι παράλληλη στην ευθεία $\varepsilon_{2}$ που ορίζεται από τις κορυφές του Α και Β. Αξιοποιώντας τα δεδομένα του σχήματος:

\begin{alist}
\item να υπολογίσετε τις γωνίες του τριγώνου $AB\varGamma $. \monades{15}

\item να δικαιολογήσετε γιατί το τρίγωνο $AB\varGamma $ είναι ισοσκελές και να εξηγήσετε ποιες είναι οι ίσες πλευρές του. \monades{10}


\begin{center}
\begin{tikzpicture}[scale=1.2]
\tkzDefPoint(0,0){B}
\tkzDefPoint(5,0){A}
\tkzDefTriangle[two angles = 50 and 80](B,A)\tkzGetPoint{C}
\tkzDefLine[parallel=through C](A,B) \tkzGetPoint{e1_dir}
\tkzDefLine[parallel=through A](e1_dir,C) \tkzGetPoint{e2_dir}

\draw[l3] (C)--($(C)-(3,0)$) node[left] {$\varepsilon_1$};
\draw[l3] (C)--($(C)+(3,0)$);
\draw[l3] ($(B)-(1,0)$)--($(A)+(1,0)$) node[right] {$\varepsilon_2$};

\draw[l3] (A)--(B)--(C)--cycle;
\tkzDrawPoints(A,B,C)
\tkzLabelPoint[below](A){$A$}
\tkzLabelPoint[below](B){$B$}
\tkzLabelPoint[above](C){$\varGamma $}

\tkzDefPointBy[translation=from B to A](C) \tkzGetPoint{C_right}
\tkzDefPointBy[translation=from A to B](C) \tkzGetPoint{C_left}
\tkzMarkAngle[size=0.6,mark=none](A,C,C_right) \tkzLabelAngle[pos=0.8](A,C,C_right){$50^\circ$}
\tkzMarkAngle[size=0.6,mark=none](C_left,C,B) \tkzLabelAngle[pos=0.8](C_left,C,B){$50^\circ$}
\tkzMarkAngle[size=0.6,mark=none](B,C,A) \tkzLabelAngle[pos=0.8](B,C,A){$80^\circ$}
\end{tikzpicture}
\end{center}
\end{alist}

\vspace{4mm}

% --- Άσκηση 13536 (13536.tex) ---
\Askhsh{13536}{2}
\begin{ylh}{1}
    \item 3.1. Είδη και στοιχεία τριγώνων
    \item 3.6. Κριτήρια ισότητας ορθογώνιων τριγώνων
    \item 4.6. Άθροισμα γωνιών τριγώνου
    \item 5.5. Τετράγωνο.
\end{ylh}
% ===================================================================
%  Αρχείο: 13536.tex (Fragment)
%  Αυτόματη μετατροπή μέσω doc2latex.py
% ===================================================================



Δίνεται τετράγωνο $AB\varGamma \varDelta  $. Στις προεκτάσεις των πλευρών $AB$ προς το Β, $B\varGamma $ προς το $\varGamma$ και $\varGamma \varDelta  $ προς το $\varDelta  $ θεωρούμε σημεία $E$, Ζ και Η αντίστοιχα, ώστε $ΒE = \varGamma  Z = \varDelta   H$.

\begin{alist}
\item Να αποδείξετε ότι $ΖE = ΖH$. \monades{15}

\item Να αποδείξετε ότι $E\hat{Z}Η = 90\degree$. \monades{10}


\begin{center}
\begin{tikzpicture}[scale=1]
\tkzDefPoint(0,0){B}
\tkzDefPoint(4,0){C}
\tkzDefSquare(B,C) \tkzGetPoints{D}{A}
\tkzDefPointOnLine[pos=1.4](A,B)\tkzGetPoint{E}
\tkzDefPointOnLine[pos=1.4](B,C)\tkzGetPoint{Z}
\tkzDefPointOnLine[pos=1.4](C,D)\tkzGetPoint{H}

\draw[l3] (A)--(B)--(C)--(D)--cycle;
\draw[l3] (B)--(E);
\draw[l3] (C)--(Z);
\draw[l3] (D)--(H);
\draw[l3, maincolor] (E)--(Z)--(H);
\tkzDrawPoints(A,B,C,D,E,Z,H)
\tkzLabelPoint[above left](A){$A$}
\tkzLabelPoint[below left](B){$B$}
\tkzLabelPoint[below right](C){$\varGamma $}
\tkzLabelPoint[above right](D){$\varDelta  $}
\tkzLabelPoint[below](E){$E$}
\tkzLabelPoint[right](Z){$Z$}
\tkzLabelPoint[above](H){$H$}
\tkzMarkSegments[mark=s|](B,E C,Z D,H)
\end{tikzpicture}
\end{center}
\end{alist}

\vspace{4mm}

% --- Άσκηση 13537 (13537.tex) ---
\Askhsh{13537}{4}
\begin{ylh}{1}
    \item 3.1. Είδη και στοιχεία τριγώνων
    \item 3.2. 1ο Κριτήριο ισότητας τριγώνων
    \item 4.6. Άθροισμα γωνιών τριγώνου.
\end{ylh}
% ===================================================================
%  Αρχείο: 13537.tex (Fragment)
%  Αυτόματη μετατροπή μέσω doc2latex.py
% ===================================================================



Στο παρακάτω σχήμα δίνεται ισοσκελές τρίγωνο $AB\varGamma $ με $AB = A\varGamma \ $, σημείο $\varDelta  $ της πλευράς $A\varGamma $, ώστε $Α\varDelta   = Β\varDelta   = Β\varGamma $ και σημείο $E$ της πλευράς $AB$, ώστε $ΑE = \varGamma \varDelta  $.

\begin{alist}
\item Να αποδείξετε ότι:

\begin{alist}
\item
  $\hat{\varGamma } = 2\hat{A}$ . \monades{6}
\item
  $\hat{A} = 36\degree$. \monades{6}
\item
  Το τρίγωνο $A\varDelta  E$ είναι ισοσκελές. \monades{6}

\item Στην προέκταση της $\varDelta  E$ προς το $E$ θεωρούμε σημείο $Z$, ώστε $\varDelta   Z = Α\varGamma $. Να αποδείξετε ότι το τρίγωνο $B\varDelta   Z$ είναι ισοσκελές. \monades{7}


\begin{center}
\begin{tikzpicture}[scale=1.5]
\tkzDefPoint(0,0){B}
\tkzDefPoint(4,0){C}
\tkzDefTriangle[two angles = 72 and 72](B,C) \tkzGetPoint{A}
\tkzDefPointOnLine[pos=0.618](A,C)\tkzGetPoint{D} % Golden ratio approximation since $BC$=$BD$=$AD$
\tkzDefPointOnLine[pos=0.382](A,B)\tkzGetPoint{E} % $AE$=$DC$
\tkzDefPointOnLine[pos=2.618](D,E)\tkzGetPoint{Z} % $DZ$ = $AC$

\draw[l3] (A)--(B)--(C)--cycle;
\draw[l3] (B)--(D);
\draw[l3] (D)--(Z);
\draw[l3] (B)--(Z);
\tkzDrawPoints(A,B,C,D,E,Z)
\tkzLabelPoint[above](A){$A$}
\tkzLabelPoint[below left](B){$B$}
\tkzLabelPoint[below right](C){$\varGamma $}
\tkzLabelPoint[right](D){$\varDelta  $}
\tkzLabelPoint[left](E){$E$}
\tkzLabelPoint[below](Z){$Z$}
\tkzMarkSegments[mark=s||](A,E D,C)
\tkzMarkSegments[mark=s|](B,C B,D A,D)
\end{tikzpicture}
\end{center}
\end{alist}
\end{alist}

\vspace{4mm}

% --- Άσκηση 13539 (13539.tex) ---
\Askhsh{13539}{4}
\begin{ylh}{1}
    \item 3.11. Ανισοτικές σχέσεις πλευρών και γωνιών
    \item 4.2. Τέμνουσα δύο ευθειών - Ευκλείδειο αίτημα
    \item 5.2. Παραλληλόγραμμα
    \item 5.4. Ρόμβος
    \item 5.10. Τραπέζιο
    \item 5.11. Ισοσκελές τραπέζιο.
\end{ylh}
% ===================================================================
%  Αρχείο: 13539.tex (Fragment)
%  Αυτόματη μετατροπή μέσω doc2latex.py
% ===================================================================



Στο παρακάτω σχήμα δίνεται ισοσκελές τραπέζιο $AB\varGamma \varDelta  $ με $AB$//$\varGamma \varDelta  $ και $\hat{A} = 108\degree$. Στη βάση $\varGamma \varDelta  $ θεωρούμε σημείο $E$, ώστε οι $A\varGamma $, $AE$ να τριχοτομούν τη γωνία $\hat{A}$.

\begin{alist}
\item Να υπολογίσετε τις γωνίες του τριγώνου $A\varDelta  E$. \monades{10}

\item Να αποδείξετε ότι:

\begin{alist}
\item
  Το τρίγωνο $A\varDelta  E$ είναι ισοσκελές. \monades{5}
\item
  Το τετράπλευρο $AB\varGamma $Ε είναι ρόμβος. \monades{10}

\begin{center}
\begin{tikzpicture}[scale=1.2]
\tkzDefPoint(0,0){B}
\tkzDefPoint(5,0){C}
\tkzDefTriangle[two angles = 72 and 72](B,C) \tkzGetPoint{A_temp}
\tkzDefPointBy[rotation=center B angle 72](C) \tkzGetPoint{A_dir}
\tkzDefPointOnLine[pos=0.6](B,A_dir)\tkzGetPoint{A}
\tkzDefLine[parallel=through A](B,C)\tkzGetPoint{d_dir}
\tkzDefPointBy[rotation=center C angle -72](B) \tkzGetPoint{C_dir}
\tkzInterLL(A,d_dir)(C,C_dir)\tkzGetPoint{D}

\tkzDefPointBy[rotation=center A angle -36](D)\tkzGetPoint{E_dir}
\tkzInterLL(A,E_dir)(C,D)\tkzGetPoint{E}

\draw[l3] (A)--(B)--(C)--(D)--cycle;
\draw[l3] (A)--(E);
\draw[l3] (A)--(C);
\tkzDrawPoints(A,B,C,D,E)
\tkzLabelPoint[above left](A){$A$}
\tkzLabelPoint[below left](B){$B$}
\tkzLabelPoint[below right](C){$\varGamma $}
\tkzLabelPoint[above right](D){$\varDelta  $}
\tkzLabelPoint[above](E){$E$}
\tkzMarkAngle[size=0.6,mark=none](E,A,D) \tkzLabelAngle[pos=0.8](E,A,D){$36^\circ$}
\tkzMarkAngle[size=0.6,mark=none](C,A,E) \tkzLabelAngle[pos=0.8](C,A,E){$36^\circ$}
\tkzMarkAngle[size=0.6,mark=none](B,A,C) \tkzLabelAngle[pos=0.8](B,A,C){$36^\circ$}
\end{tikzpicture}
\end{center}
\end{alist}
\end{alist}

\vspace{4mm}

% --- Άσκηση 13540 (13540.tex) ---
\Askhsh{13540}{4}
\begin{ylh}{1}
    \item 5.2. Παραλληλόγραμμα
    \item 5.6. Εφαρμογές στα τρίγωνα
    \item 5.9. Μια ιδιότητα του ορθογώνιου τριγώνου.
\end{ylh}
% ===================================================================
%  Αρχείο: 13540.tex (Fragment)
%  Αυτόματη μετατροπή μέσω doc2latex.py
% ===================================================================



Δίνεται παραλληλόγραμμο $AB\varGamma \varDelta  $ τέτοιο, ώστε η διαγώνιός του $A\varGamma $ να είναι κάθετη στη $B\varGamma $. Θεωρούμε τα μέσα Ε, Ζ και Η των $AB$, $A\varGamma $ και $A\varDelta  $ αντίστοιχα.

\begin{alist}
\item Να αποδείξετε ότι:

\begin{alist}
\item
  $\varGamma  E = ΖH$. \monades{9}
\item
  Η $\varGamma  A$ είναι διχοτόμος της γωνίας $\varDelta  \hat{\varGamma }E$. \monades{9}

\item Αν $\varDelta   H = \frac{AB}{4}$, να αποδείξετε ότι το τρίγωνο $B\varGamma $Ε είναι ισόπλευρο. \monades{7}


\begin{center}
\begin{tikzpicture}[scale=1]
\tkzDefPoint(0,0){B}
\tkzDefPoint(0,4){C}
\tkzDefPoint(2.5,0){A}
\tkzDefPointBy[translation=from B to C](A) \tkzGetPoint{D}
\tkzDefMidPoint(A,B) \tkzGetPoint{E}
\tkzDefMidPoint(A,C) \tkzGetPoint{Z}
\tkzDefMidPoint(A,D) \tkzGetPoint{H}

\draw[l3] (A)--(B)--(C)--(D)--cycle;
\draw[l3] (A)--(C);
\draw[l3] (C)--(E);
\draw[l3] (Z)--(H);
\tkzDrawPoints(A,B,C,D,E,Z,H)
\tkzLabelPoint[right](A){$A$}
\tkzLabelPoint[left](B){$B$}
\tkzLabelPoint[left](C){$\varGamma $}
\tkzLabelPoint[right](D){$\varDelta  $}
\tkzLabelPoint[below](E){$E$}
\tkzLabelPoint[below](Z){$Z$}
\tkzLabelPoint[right](H){$H$}
\tkzMarkRightAngle(C,B,A) % It's a rectangle actually, $AC$ perp $BC$ means $ABC$ is right triangle that translates
\tkzMarkRightAngle(B,C,A) % Oh wait $AC$ perp $BC$ means Angle C is 90 in $ABC$
\end{tikzpicture}
\end{center}
\end{alist}
\end{alist}

\vspace{4mm}

% --- Άσκηση 13619 (13619.tex) ---
\Askhsh{13619}{2}
\begin{ylh}{1}
    \item 4.6. Άθροισμα γωνιών τριγώνου
    \item 4.8. Άθροισμα γωνιών κυρτού ν-γώνου.
\end{ylh}
% ===================================================================
%  Αρχείο: 13619.tex (Fragment)
%  Αυτόματη μετατροπή μέσω doc2latex.py
% ===================================================================



Θεωρούμε το τετράπλευρο $AB\varGamma \varDelta  $ του σχήματος με ${\hat{A}}_{\varepsilon\xi} = 100\degree$ και $\hat{\text{Β}} + \hat{\varGamma} = 220\degree$.

Αν οι διχοτόμοι των γωνιών $\hat{\text{Β}}$ και $\hat{\varGamma}$ τέμνονται στο Ο, τότε:

\begin{alist}
\item Να υπολογίσετε τις γωνίες $\hat{A}$ και $\hat{\text{Δ}}$ του τετραπλεύρου $AB\varGamma \varDelta  $.

\monades{10}

\item Να αποδείξετε ότι $B\hat{\text{Ο}}$Γ$\  = 70\degree$.

\monades{15}


\begin{center}
\begin{tikzpicture}[scale=0.8]
% Quadrilateral $AB\varGamma \varDelta  $: Â_ext = 100° → Â = 80°, B̂+Γ̂ = 220°
\tkzDefPoint(0,3){A}
\tkzDefPoint(-1.5,0){B}
\tkzDefPoint(4,0){Gamma}
\tkzDefPoint(3.5,3.5){Delta}

% Draw quadrilateral
\draw[l3] (A)--(B)--(Gamma)--(Delta)--cycle;

% Bisectors of B and Γ
\tkzDefLine[bisector](A,B,Gamma) \tkzGetPoint{bB}
\tkzDefLine[bisector](B,Gamma,Delta) \tkzGetPoint{bG}

% O is intersection of bisectors
\tkzInterLL(B,bB)(Gamma,bG) \tkzGetPoint{O}

% Draw bisectors
\draw[l3, dashed] (B)--(O);
\draw[l3, dashed] (Gamma)--(O);

% Points and labels
\tkzDrawPoints(A,B,Gamma,Delta,O)
\tkzLabelPoint[above left](A){$A$}
\tkzLabelPoint[below left](B){$B$}
\tkzLabelPoint[below right](Gamma){$\varGamma $}
\tkzLabelPoint[above right](Delta){$\varDelta  $}
\tkzLabelPoint[above](O){$ O$}
\end{tikzpicture}
\end{center}
\end{alist}

\vspace{4mm}

% --- Άσκηση 13653 (13653.tex) ---
\Askhsh{13653}{2}
\begin{ylh}{1}
    \item 3.2. 1ο Κριτήριο ισότητας τριγώνων
    \item 3.7. Κύκλος - Μεσοκάθετος - Διχοτόμος
    \item 4.6. Άθροισμα γωνιών τριγώνου
    \item 5.9. Μια ιδιότητα του ορθογώνιου τριγώνου.
\end{ylh}
% ===================================================================
%  Αρχείο: 13653.tex (Fragment)
%  Αυτόματη μετατροπή μέσω doc2latex.py
% ===================================================================



Σχεδιάζουμε γωνία $x\hat{\text{O}}$y$\  = \ $60\textsuperscript{ο} και παίρνουμε σημείο $A$ επί της πλευράς Οx, τέτοιο ώστε AO$\  =$ 2. Φέρουμε τη διχοτόμο Οδ της γωνίας $x\hat{\text{O}}$y και θεωρούμε σημείο $M$ στην Οδ, τέτοιο ώστε AΜ$\  = \ $ΑΟ. Να υπολογίσετε:

\begin{alist}
\item Τη γωνία $\text{δ}\hat{\text{O}}$y.

\monades{6}

\item Τις γωνίες του τριγώνου $AOM$.

\monades{9}

\item Το μήκος του ύψους $AB$ που αντιστοιχεί στη βάση $OM$ του ισοσκελούς τριγώνου $AOM$.

\monades{10}


\begin{center}
\begin{tikzpicture}[scale=1.1]
% Angle xÔy = 60°, A on Ox, $AO$ = 2
\tkzDefPoint(0,0){O}
\tkzDefPoint(3,0){x_dir}
\tkzDefPointBy[rotation=center O angle 60](x_dir) \tkzGetPoint{y_dir}

% Bisector Oδ at 30°
\tkzDefLine[bisector](x_dir,O,y_dir) \tkzGetPoint{d_dir}

% A on Ox with $AO$ = 2
\tkzDefPointOnLine[pos=0.65](O,x_dir) \tkzGetPoint{A}

% M on bisector Oδ with $AM$ = $AO$
\tkzInterLC(O,d_dir)(A,O) \tkzGetPoints{M_temp}{M}

% B is foot of altitude from A to $OM$
\tkzDefPointBy[projection=onto O--M](A) \tkzGetPoint{B_alt}

% Draw rays
\draw[l3, ->] (O)--(x_dir) node[right]{$x$};
\draw[l3, ->] (O)--(y_dir) node[above]{$y$};
\draw[l3, ->] (O)--($(O)!1.2!(d_dir)$) node[above right]{$\delta$};

% Draw triangle $AOM$
\draw[l3] (A)--(O)--(M)--cycle;

% Draw altitude $AB$
\draw[l3, dashed] (A)--(B_alt);

% Angle and right angle marks
\tkzMarkAngle[size=0.6](x_dir,O,y_dir)
\tkzLabelAngle[pos=0.9](x_dir,O,y_dir){$60\degree$}
\tkzMarkRightAngle[size=0.2](A,B_alt,O)

% Equal side marks
\tkzMarkSegment[mark=||](A,O)
\tkzMarkSegment[mark=||](A,M)

% Points and labels
\tkzDrawPoints(O,A,M,B_alt)
\tkzLabelPoint[below left](O){$ O$}
\tkzLabelPoint[below](A){$A$}
\tkzLabelPoint[above](M){$\mathrm{M}$}
\tkzLabelPoint[left](B_alt){$B$}
\end{tikzpicture}
\end{center}
\end{alist}

\vspace{4mm}

% --- Άσκηση 13654 (13654.tex) ---
\Askhsh{13654}{2}
\begin{ylh}{1}
    \item 3.1. Είδη και στοιχεία τριγώνων
    \item 4.6. Άθροισμα γωνιών τριγώνου.
\end{ylh}
% ===================================================================
%  Αρχείο: 13654.tex (Fragment)
%  Αυτόματη μετατροπή μέσω doc2latex.py
% ===================================================================



Στο ακόλουθο σχήμα είναι $\hat{A} = \text{9}$0\textsuperscript{ο}, $\text{A}\hat{\text{Β}}\varGamma - A\hat{\varGamma}\text{Β} = 5$0\textsuperscript{ο} και $\text{A}\hat{\text{Δ}}$Γ$\  = 5$0\textsuperscript{ο}.

\begin{alist}
\item Να υπολογίσετε τις οξείες γωνίες $A\hat{\text{Β}}\varGamma$ και $A\hat{\varGamma}\text{Β}$ του ορθογωνίου τριγώνου $AB\varGamma $.

\monades{10}

\item Nα αποδείξετε ότι η $\varGamma B$ είναι διχοτόμος της γωνίας Α$\hat{\varGamma}\text{Δ}$.

\monades{15}


\begin{center}
\begin{tikzpicture}[scale=0.9]
% Right triangle $AB\varGamma $, Â=90°, ÂBΓ−ÂΓB=50°
% → ÂBΓ=70°, ÂΓB=20°
\tkzDefPoint(0,0){A}
\tkzDefPoint(0,3){B}       % $AB$ vertical
\tkzDefPoint(1.1,0){Gamma} % $A\varGamma $ horizontal, small (angle 20° at Γ)

% Δ on extension of $\varGamma A$ past A such that ÂΔΓ=50°
\tkzDefPointOnLine[pos=-1.5](Gamma,A) \tkzGetPoint{Delta}

% Draw right triangle
\draw[l3] (A)--(B)--(Gamma)--cycle;

% Draw extension to Δ
\draw[l3] (Gamma)--(Delta);
\draw[l3] (B)--(Delta);

% Right angle at A
\tkzMarkRightAngle[size=0.2](B,A,Gamma)

% Points and labels
\tkzDrawPoints(A,B,Gamma,Delta)
\tkzLabelPoint[below right](A){$A$}
\tkzLabelPoint[above](B){$B$}
\tkzLabelPoint[below right](Gamma){$\varGamma $}
\tkzLabelPoint[below left](Delta){$\varDelta  $}
\end{tikzpicture}
\end{center}
\end{alist}

\vspace{4mm}

% --- Άσκηση 13672 (13672.tex) ---
\Askhsh{13672}{4}
\begin{ylh}{1}
    \item 3.1. Είδη και στοιχεία τριγώνων
    \item 3.2. 1ο Κριτήριο ισότητας τριγώνων
    \item 4.2. Τέμνουσα δύο ευθειών - Ευκλείδειο αίτημα
    \item 5.9. Μια ιδιότητα του ορθογώνιου τριγώνου.
\end{ylh}
% ===================================================================
%  Αρχείο: 13672.tex (Fragment)
%  Αυτόματη μετατροπή μέσω doc2latex.py
% ===================================================================



Δίνεται ορθογώνιο τρίγωνο $AB\varGamma $ με $\hat{A} = 90\degree$ και $AB>A\varGamma $. Από το μέσο $\varDelta  $ της πλευράς $B\varGamma $ φέρουμε κάθετη στη $B\varGamma $ όπως φαίνεται στο παρακάτω σχήμα, η οποία τέμνει τη διχοτόμο $AH$ της γωνίας $\hat{A}$ στο σημείο $E$. Έστω $AZ$ το ύψος στην υποτείνουσα. Να αποδείξετε ότι:

\begin{alist}
\item $\varGamma\hat{A}\text{Ζ} =$ $\text{Δ}\hat{A}\text{Β}$.

\monades{8}

\item $\text{$A\varDelta  $} = \text{$\varDelta   E$}$.

\monades{9}

\item $\text{Ζ}\hat{A}\text{Δ} = \hat{\varGamma} - \hat{\text{Β}}$.

\monades{8}


\begin{center}
\begin{tikzpicture}[scale=0.9]
% Right triangle $AB\varGamma $: Â=90°, $AB$ > $A\varGamma $
\tkzDefPoint(0,0){A}
\tkzDefPoint(0,4){B}    % $AB$ = 4 (με$\varGamma  A$λύτερη κάθετη)
\tkzDefPoint(2.5,0){Gamma} % $A\varGamma $ = 2.5

% Midpoint $\varDelta$ of $B\varGamma $
\tkzDefMidPoint(B,Gamma) \tkzGetPoint{Delta}

% Bisector of angle A: $AH$ meets $B\varGamma $ at H
\tkzDefLine[bisector](B,A,Gamma) \tkzGetPoint{bisect_dir}
\tkzInterLL(A,bisect_dir)(B,Gamma) \tkzGetPoint{H}

% E: perpendicular from $\varDelta$ to $B\varGamma $ meets bisector $AH$
\tkzDefLine[orthogonal=through Delta](B,Gamma) \tkzGetPoint{perp_dir}
\tkzInterLL(Delta,perp_dir)(A,H) \tkzGetPoint{E}

% Z: foot of altitude from A to $B\varGamma $
\tkzDefPointBy[projection=onto B--Gamma](A) \tkzGetPoint{Z}

% Draw triangle
\draw[l3] (A)--(B)--(Gamma)--cycle;

% Draw bisector $AH$
\draw[l3] (A)--(H);

% Draw perpendicular from Δ
\draw[l3, dashed] (Delta)--(E);

% Draw altitude $AZ$
\draw[l3, dashed] (A)--(Z);

% Right angle marks
\tkzMarkRightAngle[size=0.25](B,A,Gamma)
\tkzMarkRightAngle[size=0.25](A,Z,B)
\tkzMarkRightAngle[size=0.2](A,Delta,E)

% Points and labels
\tkzDrawPoints(A,B,Gamma,Delta,H,E,Z)
\tkzLabelPoint[below left](A){$A$}
\tkzLabelPoint[above left](B){$B$}
\tkzLabelPoint[below right](Gamma){$\varGamma $}
\tkzLabelPoint[right](Delta){$\varDelta  $}
\tkzLabelPoint[below](H){$ H$}
\tkzLabelPoint[above right](E){$ E$}
\tkzLabelPoint[below](Z){$ Z$}
\end{tikzpicture}
\end{center}
\end{alist}

\vspace{4mm}

% --- Άσκηση 13687 (13687.tex) ---
\Askhsh{13687}{2}
\begin{ylh}{1}
    \item 3.2. 1ο Κριτήριο ισότητας τριγώνων
    \item 4.2. Τέμνουσα δύο ευθειών - Ευκλείδειο αίτημα
    \item 4.6. Άθροισμα γωνιών τριγώνου.
\end{ylh}
% ===================================================================
%  Αρχείο: 13687.tex (Fragment)
%  Αυτόματη μετατροπή μέσω doc2latex.py
% ===================================================================



Σε κύκλο με κέντρο Ο και ακτίνα R θεωρούμε επίκεντρη γωνία Α$\hat{\text{Ο}}$Β$\  = \ $40\textsuperscript{ο}. Προεκτείνουμε τις ακτίνες $OA$ και $OB$ κατά τμήματα $A\varGamma $ και $B\varDelta  $ αντίστοιχα, έτσι ώστε ΑΓ$\  = \ $ΟΑ και $B\varDelta  $$\  = \ $ΟΒ, όπως φαίνεται στο παρακάτω σχήμα.

\begin{alist}
\item Να αποδείξετε ότι Ο$\hat{A}$Β$\  = \ $Ο$\hat{\text{Β}}$Α$\  = 7\text{0}\text{ο}$.

\monades{10}

\item Να υπολογίσετε τις γωνίες $\text{O}\hat{\varGamma}$Δ και $\text{O}\hat{\text{Δ}}$Γ.

\monades{10}

\item Να αποδείξετε ότι $AB$//$\varGamma \varDelta  $.

\monades{5}


\begin{center}
\begin{tikzpicture}[scale=0.9]
% Circle center O, radius R
\tkzDefPoint(0,0){O}
\def\R{2}

% Points A and B on circle, central angle 40°
\tkzDefPoint(20:\R){B}
\tkzDefPoint(60:\R){A}

% Extend $OA$ to $\varGamma$ so $A\varGamma $ = $OA$ (i.e. $O\varGamma $ = 2·$OA$)
\tkzDefPointOnLine[pos=2](O,A) \tkzGetPoint{Gamma}

% Extend $OB$ to $\varDelta$ so $B\varDelta  $ = $OB$ (i.e. $O\varDelta  $ = 2·$OB$)
\tkzDefPointOnLine[pos=2](O,B) \tkzGetPoint{Delta}

% Draw circle
\tkzDrawCircle(O,A)

% Draw radii and extensions
\draw[l3] (O)--(Gamma);
\draw[l3] (O)--(Delta);

% Draw $AB$ and $\varGamma \varDelta  $
\draw[l3] (A)--(B);
\draw[l3] (Gamma)--(Delta);

% Draw $\varGamma \varDelta  $ connection to show parallel
\draw[l3, dashed] (A)--(Delta);
\draw[l3, dashed] (B)--(Gamma);

% Central angle mark
\tkzMarkAngle[size=0.5](B,O,A)
\tkzLabelAngle[pos=0.8](B,O,A){$40^{\circ}$}

% Points and labels
\tkzDrawPoints(O,A,B,Gamma,Delta)
\tkzLabelPoint[above right](A){$A$}
\tkzLabelPoint[right](B){$B$}
\tkzLabelPoint[above left](Gamma){$\varGamma $}
\tkzLabelPoint[below right](Delta){$\varDelta  $}
\tkzLabelPoint[below left](O){$ O$}
\end{tikzpicture}
\end{center}
\end{alist}

\vspace{4mm}

% --- Άσκηση 13697 (13697.tex) ---
\Askhsh{13697}{4}
\begin{ylh}{1}
    \item 3.6. Κριτήρια ισότητας ορθογώνιων τριγώνων
    \item 4.6. Άθροισμα γωνιών τριγώνου.
\end{ylh}
% ===================================================================
%  Αρχείο: 13697.tex (Fragment)
%  Αυτόματη μετατροπή μέσω doc2latex.py
% ===================================================================



Στο παρακάτω σχήμα, τα τμήματα $AE$, $BZ$, $B\varDelta  $ και $\varGamma  E$ αναπαριστάνουν τέσσερεις ίσους ράβδους μήκους 40 cm οι οποίες αποτελούν μέρη μιας κρεμάστρας τοίχου.

Οι ράβδοι συνδέονται με τέτοιο τρόπο ώστε ανά δύο απέναντι να είναι παράλληλες, δηλαδή $AE$ // $B\varDelta  $ και $BZ$ // $\varGamma  E$, και ανά δύο να έχουν κοινό μέσο, δηλαδή Κ κοινό μέσο των $AE$, $BZ$ και Λ κοινό μέσο των $B\varDelta  $, $\varGamma  E$. Έστω ότι η μία από τις γωνίες που σχηματίζουν οι τεμνόμενες ράβδοι $AE$ και $BZ$ με κορυφή το κοινό τους μέσο $K$, η γωνία $B\hat{Κ}$Ε, είναι ίση με 120\degree.

\begin{alist}
\item Να αποδείξετε ότι Α$\hat{Κ}$Β = Κ$\hat{B}$Λ = $B\hat{\Lambda}$Γ = 60\degree. \monades{9}

\item Ένας μαθητής ισχυρίζεται ότι τα τρίγωνα $AKB$ και $B\varLambda \varGamma $ είναι ίσα και ισόπλευρα. Να εξετάσετε αν ο ισχυρισμός του είναι αληθής. Να αιτιολογήσετε την απάντησή σας.

\monades{10}

\item Να αποδείξετε ότι τα σημεία $A$, Β και $\varGamma$ ανήκουν στην ίδια ευθεία. \monades{6}


\begin{center}
\begin{tikzpicture}[scale=1.5]
\tkzDefPoint(0,0){A}
\tkzDefPoint(2,0){B}
\tkzDefPoint(4,0){C}
\tkzDefPoint(0,3.464){Z}
\tkzDefPoint(2,3.464){E}
\tkzDefPoint(4,3.464){D}

\tkzDefMidPoint(A,E) \tkzGetPoint{K}
\tkzDefMidPoint(B,D) \tkzGetPoint{L}

\draw[l3] (A)--(E);
\draw[l3] (B)--(Z);
\draw[l3] (B)--(D);
\draw[l3] (C)--(E);

\draw[l3, dashed] (A)--(B)--(C);

\tkzDrawPoints(A,B,C,Z,E,D,K,L)
\tkzLabelPoint[below](A){$A$}
\tkzLabelPoint[below](B){$B$}
\tkzLabelPoint[below](C){$\varGamma $}
\tkzLabelPoint[above](Z){$Z$}
\tkzLabelPoint[above](E){$E$}
\tkzLabelPoint[above](D){$\varDelta  $}
\tkzLabelPoint[left](K){$K$}
\tkzLabelPoint[right](L){$\Lambda$}

\tkzMarkAngle[size=0.4,mark=none](E,K,B) \tkzLabelAngle[pos=0.6](E,K,B){$120^\circ$}
\end{tikzpicture}
\end{center}
\end{alist}

\vspace{4mm}

% --- Άσκηση 13699 (13699.tex) ---
\Askhsh{13699}{4}
\begin{ylh}{1}
    \item 3.11. Ανισοτικές σχέσεις πλευρών και γωνιών
    \item 3.14. Σχετικές θέσεις ευθείας και κύκλου
    \item 4.2. Τέμνουσα δύο ευθειών - Ευκλείδειο αίτημα
    \item 5.3. Ορθογώνιο.
\end{ylh}
% ===================================================================
%  Αρχείο: 13699.tex (Fragment)
%  Αυτόματη μετατροπή μέσω doc2latex.py
% ===================================================================



Δίνονται δυο κύκλοι (Κ, ρ\textsubscript{1}) και (Λ, ρ\textsubscript{2}) που εφάπτονται εξωτερικά σε σημείο $A$. Έστω ότι μια ευθεία (ε) εφάπτεται εξωτερικά στους δυο κύκλους σε σημεία τους Β και $\varGamma$ αντίστοιχα και ότι η εσωτερική εφαπτομένη (ζ) των κύκλων στο σημείο επαφής τους Α τέμνει την ευθεία (ε) σε σημείο $M$.

\begin{alist}
\item Να αποδείξετε ότι:


\end{alist}
\begin{alist}
\item
  οι ευθείες $KB$ και $\varLambda  M$ τέμνονται σε σημείο, έστω $\varDelta  $. \monades{10}
\item
  το τρίγωνο $\varDelta   K\varLambda $ είναι ισοσκελές. \monades{10}

\item Με ποια σχέση πρέπει να συνδέονται οι ακτίνες ρ\textsubscript{1} και ρ\textsubscript{2} των δύο κύκλων ώστε το ισοσκελές τρίγωνο $\varDelta  K\varLambda $ να είναι ορθογώνιο; Να αιτιολογήσετε την απάντησή σας.

\monades{5}
\end{alist}

\vspace{4mm}

% --- Άσκηση 13702 (13702.tex) ---
\Askhsh{13702}{3}
\begin{ylh}{1}
    \item 3.7. Κύκλος - Μεσοκάθετος - Διχοτόμος
    \item 3.14. Σχετικές θέσεις ευθείας και κύκλου
    \item 3.15. Εφαπτόμενα τμήματα
    \item 3.16. Σχετικές θέσεις δυο κύκλων.
\end{ylh}
% ===================================================================
%  Αρχείο: 13702.tex (Fragment)
%  Αυτόματη μετατροπή μέσω doc2latex.py
% ===================================================================



Δίνονται δυο κύκλοι (Κ, ρ\textsubscript{1}) και (Λ, ρ\textsubscript{2}) που εφάπτονται εξωτερικά σε σημείο $A$. Μια ευθεία (ε) εφάπτεται εξωτερικά στους δυο κύκλους σε σημεία $B$ και $\varGamma$ αντίστοιχα. Αν η εσωτερική εφαπτομένη των κύκλων στο σημείο επαφής τους Α τέμνει την ευθεία (ε) σε σημείο $M$, να αποδείξετε ότι:

\begin{alist}
\item τα σημεία $A$, Β και $\varGamma$ ανήκουν σε κύκλο του οποίου να προσδιορίσετε το κέντρο και την ακτίνα. \monades{12}

\item ο κύκλος που διέρχεται από τα σημεία $A$, Β και $\varGamma$ εφάπτεται στη διάκεντρο $K\varLambda $ των κύκλων (Κ, ρ\textsubscript{1}) και (Λ, ρ\textsubscript{2}). \monades{13}
\end{alist}

\vspace{4mm}

% --- Άσκηση 13741 (13741.tex) ---
\Askhsh{13741}{2}
\begin{ylh}{1}
    \item 3.2. 1ο Κριτήριο ισότητας τριγώνων.
\end{ylh}
% ===================================================================
%  Αρχείο: 13741.tex (Fragment)
%  Αυτόματη μετατροπή μέσω doc2latex.py
% ===================================================================



Στο σχήμα που ακολουθεί οι ευθείες ε και ζ είναι παράλληλες. Αν είναι $\hat{\alpha}$ = 76\textsuperscript{ο} και

$\hat{\varGamma }$ = 120\textsuperscript{ο} , να υπολογίσετε :

\begin{alist}
\item Τη γωνία $\hat{\beta}$ . \monades{5}

\item Τις γωνίες του τετράπλευρου $AB\varGamma \varDelta  $ . \monades{12}

\item Τη γωνία $\hat{E\ }$ του τριγώνου $EA\varDelta  $. \monades{8}


\begin{center}
\begin{tikzpicture}[scale=1]
\tkzDefPoint(0,3){A_line}
\tkzDefPoint(6,3){B_line}
\tkzDefPoint(0,0){C_line}
\tkzDefPoint(6,0){D_line}

\tkzDefPoint(1,3){A}
\tkzDefPoint(4,3){B}
\tkzDefPoint(1.5,0){D}
\tkzDefPoint(5,0){C}

\tkzInterLL(A,D)(B,C) \tkzGetPoint{E}

\draw[l3] (A_line)--(B_line) node[right] {$\varepsilon$};
\draw[l3] (C_line)--(D_line) node[right] {$\zeta$};

\draw[l3] (E)--(D);
\draw[l3] (E)--(C);
\draw[l3] (A)--(B);
\draw[l3] (C)--(D);

\tkzDrawPoints(A,B,C,D,E)
\tkzLabelPoint[above left](A){$A$}
\tkzLabelPoint[above right](B){$B$}
\tkzLabelPoint[below right](C){$\varGamma $}
\tkzLabelPoint[below left](D){$\varDelta  $}
\tkzLabelPoint[above](E){$E$}

\tkzMarkAngle[size=0.5,mark=none](D,A,B) \tkzLabelAngle[pos=0.7](D,A,B){$\alpha$}
\tkzMarkAngle[size=0.5,mark=none](A,B,C) \tkzLabelAngle[pos=0.7](A,B,C){$\beta$}
\tkzMarkAngle[size=0.5,mark=none](B,C,D) \tkzLabelAngle[pos=0.8](B,C,D){$\varGamma $}
\end{tikzpicture}
\end{center}
\end{alist}

\vspace{4mm}

% --- Άσκηση 13742 (13742.tex) ---
\Askhsh{13742}{4}
\begin{ylh}{1}
    \item 3.4. 3ο Κριτήριο ισότητας τριγώνων
    \item 4.2. Τέμνουσα δύο ευθειών - Ευκλείδειο αίτημα
    \item 4.6. Άθροισμα γωνιών τριγώνου
    \item 5.2. Παραλληλόγραμμα
    \item 5.10. Τραπέζιο.
\end{ylh}
% ===================================================================
%  Αρχείο: 13742.tex (Fragment)
%  Αυτόματη μετατροπή μέσω doc2latex.py
% ===================================================================



Δίνεται ισοσκελές τρίγωνο $AB\varGamma $ με $AB=Α\varGamma $ και $M$ το μέσο της βάσης του $B\varGamma $. Φέρουμε ΒΚ$\ \bot$ $B\varGamma $ έτσι ώστε $BK$ = $A\varGamma $ (το σημείο $K$ είναι στο ημιεπίπεδο που δεν ανήκει το Α).

\begin{alist}
\item Να αποδείξετε ότι $AM$//$BK$ και $AB$ = $BK$. \monades{8}

\item Να δείξετε ότι η $AK$ είναι διχοτόμος της γωνίας $BAM$. \monades{5}

\item Να αποδείξετε ότι B$\hat{Κ}$A= 45\textsuperscript{ο \_} $\frac{\hat{\varGamma }}{2}$ \monades{6}

\item Μπορεί το τετράπλευρο $ABKM$ να είναι παραλληλόγραμμο; Να αιτιολογήσετε την απάντησή σας. \monades{6}
\end{alist}

\vspace{4mm}

% --- Άσκηση 13743 (13743.tex) ---
\Askhsh{13743}{4}
\begin{ylh}{1}
    \item 3.1. Είδη και στοιχεία τριγώνων
    \item 3.2. 1ο Κριτήριο ισότητας τριγώνων
    \item 4.2. Τέμνουσα δύο ευθειών - Ευκλείδειο αίτημα
    \item 5.2. Παραλληλόγραμμα
    \item 5.6. Εφαρμογές στα τρίγωνα.
\end{ylh}
% ===================================================================
%  Αρχείο: 13743.tex (Fragment)
%  Αυτόματη μετατροπή μέσω doc2latex.py
% ===================================================================



Δίνεται τρίγωνο $AB\varGamma $ και σημείο $M$ στην πλευρά $AB$. Από το $M$ φέρουμε παράλληλη στη $B\varGamma $ που τέμνει την $A\varGamma $ στο σημείο $\varDelta  $.

\begin{alist}
\item Να αποδείξετε ότι Δ$\hat{Μ}$Γ = $B\hat{\varGamma }$Μ. \monades{05}

\item Αν το τρίγωνο $\varGamma  A$Β είναι ισοσκελές με βάση $AB$, να προσδιορίσετε τη θέση του σημείου $M$ στην $AB$ ώστε το τρίγωνο $\varDelta  M\varGamma $ να είναι ισοσκελές με $\varDelta  M$ = $\varDelta  \varGamma $ και να δικαιολογήσετε τους ισχυρισμούς σας. \monades{10}

\item Αν $M$ είναι το μέσο του τμήματος $AB$ και $E$ το μέσο του τμήματος $B\varGamma $ να δικαιολογήσετε γιατί το τετράπλευρο $M\varDelta  EB$ είναι παραλληλόγραμμο. \monades{10}
\end{alist}

\vspace{4mm}

% --- Άσκηση 13744 (13744.tex) ---
\Askhsh{13744}{4}
\begin{ylh}{1}
    \item 3.1. Είδη και στοιχεία τριγώνων
    \item 3.2. 1ο Κριτήριο ισότητας τριγώνων
    \item 3.6. Κριτήρια ισότητας ορθογώνιων τριγώνων
    \item 3.11. Ανισοτικές σχέσεις πλευρών και γωνιών
    \item 4.6. Άθροισμα γωνιών τριγώνου
    \item 5.5. Τετράγωνο.
\end{ylh}
% ===================================================================
%  Αρχείο: 13744.tex (Fragment)
%  Αυτόματη μετατροπή μέσω doc2latex.py
% ===================================================================



Δίνεται τετράγωνο $AB\varGamma \varDelta  $. Στις προεκτάσεις των πλευρών του $AB$ και $B\varGamma $ προς το Β και προς το $\varGamma$ αντίστοιχα, παίρνουμε τα σημεία $E$ και $Z$ τέτοια ώστε $BE$ = $\varGamma Z$. Αν Ρ είναι το σημείο τομής των $AZ$ και $\varDelta  E$, τότε:

\begin{alist}
\item Να αποδείξετε ότι:


\end{alist}
\begin{alist}
\item
  Οι γωνίες Α$\hat{E}$Δ και $B\hat{Ζ}$Α είναι ίσες.
\item
  Τα τμήματα $AZ$ και $\varDelta  E$ είναι κάθετα.

\monades{18}

\item Αν γνωρίζετε ότι το σημείο τομής Ρ των $AZ$ και $\varDelta  E$ είναι τέτοιο ώστε $PB$ = $AB$, να προσδιορίσετε τη θέση του σημείου $E$ στην προέκταση του τμήματος $AB$. \monades{07}
\end{alist}

\vspace{4mm}

% --- Άσκηση 13745 (13745.tex) ---
\Askhsh{13745}{4}
\begin{ylh}{1}
    \item 3.1. Είδη και στοιχεία τριγώνων
    \item 3.2. 1ο Κριτήριο ισότητας τριγώνων
    \item 3.11. Ανισοτικές σχέσεις πλευρών και γωνιών
    \item 4.2. Τέμνουσα δύο ευθειών - Ευκλείδειο αίτημα
    \item 5.4. Ρόμβος
    \item 5.6. Εφαρμογές στα τρίγωνα.
\end{ylh}
% ===================================================================
%  Αρχείο: 13745.tex (Fragment)
%  Αυτόματη μετατροπή μέσω doc2latex.py
% ===================================================================



Δίνεται ισοσκελές τρίγωνο $AB\varGamma $, το μέσο $M$ της βάσης $B\varGamma $ και τυχαίο εσωτερικό σημείο $\varDelta  $ στη βάση του.

\begin{alist}
\item Αν από το μέσο $M$ φέρουμε παράλληλες προς τις πλευρές $AB$ και $A\varGamma $ του τριγώνου, που τις τέμνουν στα σημεία $E$ και $Z$ αντίστοιχα να αποδείξετε ότι:
\end{alist}
\begin{alist}
\item
  $ME$ = $MZ$. \monades{6}
\item
  Το $AEMZ$ είναι ρόμβος με περίμετρο ίση με 2ΑΒ. \monades{7}

\item Αν πάρουμε τυχαίο εσωτερικό σημείο $\varDelta  $ στο ευθύγραμμο τμήμα $B\varGamma $, διαφορετικό από το μέσο $M$, και φέρουμε τις παράλληλες προς τις πλευρές $AB$ και $A\varGamma $ του τριγώνου, που τις τέμνουν στα σημεία $K$ και Λ αντίστοιχα, τότε:
\end{alist}
\begin{alist}
\item
  Ποιο είναι το είδος του τετράπλευρου $AK\varDelta  \varLambda $;
\item
  Να συγκρίνετε την περίμετρο του τετράπλευρου $AK\varDelta  \varLambda $ με την περίμετρο του ρόμβου $AEMZ$ του ερωτήματος α ii) και να διατυπώστε λεκτικά το συμπέρασμα που προκύπτει.

\monades{12}
\end{alist}

\vspace{4mm}

% --- Άσκηση 13746 (13746.tex) ---
\Askhsh{13746}{4}
\begin{ylh}{1}
    \item 3.1. Είδη και στοιχεία τριγώνων
    \item 3.2. 1ο Κριτήριο ισότητας τριγώνων
    \item 3.12. Τριγωνική ανισότητα
    \item 5.2. Παραλληλόγραμμα
    \item 5.3. Ορθογώνιο.
\end{ylh}
% ===================================================================
%  Αρχείο: 13746.tex (Fragment)
%  Αυτόματη μετατροπή μέσω doc2latex.py
% ===================================================================



Δίνεται τρίγωνο $AB\varGamma $ και η διάμεσός του $A\varDelta  $. Στην προέκταση της διαμέσου $A\varDelta  $ προς το $\varDelta$ παίρνουμε σημείο $E$, έτσι ώστε $A\varDelta  $ = $\varDelta  E$.

\begin{alist}
\item Να αποδείξετε ότι :
\end{alist}
\begin{alist}
\item
  Τα τρίγωνα Α$B\varDelta  $ και $E\varGamma \varDelta  $ είναι ίσα. \monades{07}
\item
  Η διάμεσος $A\varDelta  $ είναι μικρότερη από το ημιάθροισμα των πλευρών $AB$ και $A\varGamma $ που την περιέχουν. \monades{08}

\item Αν στο τρίγωνο $AB\varGamma $ το διπλάσιο της διαμέσου $A\varDelta  $ ισούται με την πλευρά $B\varGamma $, να χαρακτηρίσετε το είδος του τετράπλευρου Α$BE$Γ και το είδος του τριγώνου $AB\varGamma $ και να αιτιολογήσετε τις απαντήσεις σας.

\monades{10}
\end{alist}

\vspace{4mm}

% --- Άσκηση 13748 (13748.tex) ---
\Askhsh{13748}{2}
\begin{ylh}{1}
    \item 3.2. 1ο Κριτήριο ισότητας τριγώνων
    \item 4.2. Τέμνουσα δύο ευθειών - Ευκλείδειο αίτημα.
\end{ylh}
% ===================================================================
%  Αρχείο: 13748.tex (Fragment)
%  Αυτόματη μετατροπή μέσω doc2latex.py
% ===================================================================



Σε τρίγωνο $AB\varGamma $ θεωρούμε το μέσο $\varDelta  $ της πλευράς $A\varGamma $. Φέρουμε τμήμα $\varDelta  E$ ίσο και παράλληλο με την πλευρά $B\varGamma $ όπως φαίνεται στο σχήμα. Προεκτείνουμε την $A\varGamma $ προς το μέρος του $\varGamma$ και παίρνουμε σημείο $Z$ τέτοιο ώστε $\varGamma Z$ = $\varDelta  \varGamma $. Να αποδείξετε ότι:

\begin{alist}
\item Τα τρίγωνα $AB\varGamma $ και Ζ$E\varDelta  $ είναι ίσα. \monades{10}

\item $AB$ // $EZ$. \monades{15}


\begin{center}
\begin{tikzpicture}[scale=1]
\tkzDefPoint(0,4){A}
\tkzDefPoint(-2,0){B}
\tkzDefPoint(3,0){C}
\tkzDefMidPoint(A,C) \tkzGetPoint{D}
\tkzDefPointBy[translation=from B to C](D) \tkzGetPoint{E}
\tkzDefPointBy[translation=from A to D](C) \tkzGetPoint{Z}

\draw[l3] (A)--(B)--(C)--cycle;
\draw[l3] (D)--(E);
\draw[l3] (Z)--(E);
\draw[l3] (D)--(Z);
\draw[l3, dashed] (Z)--(C);

\tkzDrawPoints(A,B,C,D,E,Z)
\tkzLabelPoint[above](A){$A$}
\tkzLabelPoint[left](B){$B$}
\tkzLabelPoint[below](C){$\varGamma $}
\tkzLabelPoint[above right](D){$\varDelta  $}
\tkzLabelPoint[right](E){$E$}
\tkzLabelPoint[below right](Z){$Z$}
\end{tikzpicture}
\end{center}
\end{alist}

\vspace{4mm}

% --- Άσκηση 13749 (13749.tex) ---
\Askhsh{13749}{2}
\begin{ylh}{1}
    \item 3.2. 1ο Κριτήριο ισότητας τριγώνων
    \item 4.6. Άθροισμα γωνιών τριγώνου
    \item 4.8. Άθροισμα γωνιών κυρτού ν-γώνου.
\end{ylh}
% ===================================================================
%  Αρχείο: 13749.tex (Fragment)
%  Αυτόματη μετατροπή μέσω doc2latex.py
% ===================================================================



Στο παρακάτω σχήμα το $AB\varGamma \varDelta  $Ε είναι ένα πεντάγωνο στο οποίο η διαγώνιος $A\varDelta  $ είναι ίση με την πλευρά $AE$ και η ημιευθεία Αx είναι προέκταση της $BA$ προς το Α. Να υπολογίσετε δικαιολογώντας τις απαντήσεις σας:

\begin{alist}
\item Τη γωνία $\hat{\alpha}$. \monades{08}

\item Τη γωνία $\hat{\omega}$. \monades{08}

\item Τη γωνία $\hat{\varphi}$. \monades{09}


\begin{center}
\begin{tikzpicture}[scale=1.5]
\tkzDefPoint(0,3){A}
\tkzDefPoint(-1.5,1.5){B}
\tkzDefPoint(1.5,1.5){E}
\tkzDefPoint(-1,0){C}
\tkzDefPoint(1,0){D}
\tkzDefPointOnLine[pos=-0.5](B,A) \tkzGetPoint{x}

\draw[l3] (A)--(B)--(C)--(D)--(E)--cycle;
\draw[l3] (A)--(D);
\draw[l3] (A)--(E); % They are adjacent actually, but A,D and A,E are mentioned. "η διαγώνιος $A\varDelta  $ είναι ίση με την πλευρά $AE$" - oh $AE$ is a side! $AD$ is diagonal.
\draw[l3, ->] (B)--(x) node[above] {$x$};

\tkzDrawPoints(A,B,C,D,E)
\tkzLabelPoint[above right](A){$A$}
\tkzLabelPoint[left](B){$B$}
\tkzLabelPoint[below left](C){$\varGamma $}
\tkzLabelPoint[below right](D){$\varDelta  $}
\tkzLabelPoint[right](E){$E$}

\tkzMarkSegments[mark=s|](A,D A,E)
\end{tikzpicture}
\end{center}
\end{alist}

\vspace{4mm}

% --- Άσκηση 13750 (13750.tex) ---
\Askhsh{13750}{4}
\begin{ylh}{1}
    \item 3.6. Κριτήρια ισότητας ορθογώνιων τριγώνων
    \item 3.15. Εφαπτόμενα τμήματα
    \item 4.2. Τέμνουσα δύο ευθειών - Ευκλείδειο αίτημα
    \item 4.6. Άθροισμα γωνιών τριγώνου.
\end{ylh}
% ===================================================================
%  Αρχείο: 13750.tex (Fragment)
%  Αυτόματη μετατροπή μέσω doc2latex.py
% ===================================================================



Από σημείο $B$ εξωτερικό ενός κύκλου (Ο,R) φέρουμε το εφαπτόμενο τμήμα $BA$. Ενώνουμε το σημείο $B$ με το κέντρο Ο του κύκλου και προεκτείνουμε κατά ίσο τμήμα $O\varGamma $ = $BO$. Από το σημείο $\varGamma $ φέρουμε το εφαπτόμενο τμήμα $\varGamma \varDelta  $, όπως στο σχήμα.

\begin{alist}
\item Να αποδείξετε ότι:

\begin{alist}
\item
  $AB$ = $\varDelta  \varGamma $ \monades{08}
\item
  $A\varDelta  $ // $B\varGamma $ \monades{10}

\item Αν το μήκος του εφαπτόμενου τμήματος $BA$ είναι ίσο με την ακτίνα R, τι είδους τρίγωνο είναι το τρίγωνο $AO\varDelta  $; Να δικαιολογήσετε την απάντησή σας. \monades{07}


\begin{center}
\begin{tikzpicture}[scale=1]
\tkzDefPoint(0,0){O}
\tkzDefPoint(4,0){B}
\tkzDefPoint(-4,0){C}
\tkzDefPoint(2,0){temp}

\tkzDrawCircle(O,temp)
\tkzDefTangent[from=B](O,temp) \tkzGetPoints{A_temp}{A}
\tkzDefTangent[from=C](O,temp) \tkzGetPoints{D}{D_temp}

\draw[l3] (B)--(C);
\draw[l3] (B)--(A);
\draw[l3] (C)--(D);
\draw[l3] (A)--(D);
\draw[l3] (O)--(A);
\draw[l3] (O)--(D);

\tkzDrawPoints(O,A,B,C,D)
\tkzLabelPoint[below](O){$O$}
\tkzLabelPoint[right](B){$B$}
\tkzLabelPoint[left](C){$\varGamma $}
\tkzLabelPoint[above right](A){$A$}
\tkzLabelPoint[above left](D){$\varDelta  $}

\tkzMarkRightAngle(O,A,B)
\tkzMarkRightAngle(O,D,C)
\tkzMarkSegments[mark=s||](O,B O,C)
\end{tikzpicture}
\end{center}
\end{alist}
\end{alist}

\vspace{4mm}

% --- Άσκηση 13751 (13751.tex) ---
\Askhsh{13751}{4}
\begin{ylh}{1}
    \item 3.6. Κριτήρια ισότητας ορθογώνιων τριγώνων
    \item 5.6. Εφαρμογές στα τρίγωνα.
\end{ylh}
% ===================================================================
%  Αρχείο: 13751.tex (Fragment)
%  Αυτόματη μετατροπή μέσω doc2latex.py
% ===================================================================



Στο παρακάτω σχήμα οι ευθείες ε\textsubscript{1} και ε\textsubscript{2} είναι παράλληλες. Το τρίγωνο $AB\varGamma $ είναι οξυγώνιο, ενώ το $\varDelta  EZ$ είναι αμβλυγώνιο με $\hat{E}$ > 90\textsuperscript{ο} . Ισχύει επίσης ότι $A\varGamma $ = $\varDelta   Z$.

\begin{alist}
\item

\begin{alist}
\item
  Να σχεδιάσετε τα ύψη των τριγώνων από τις κορυφές Α και $\varDelta$ ονομάζοντάς τα $AH$ και $\varDelta  \varTheta $ αντίστοιχα. \monades{05}
\item
  Να αποδείξετε ότι $H\varGamma $ = $\varTheta Z$. \monades{12}

\item Να δικαιολογήσετε γιατί $EZ$ < $B\varGamma $. \monades{08}


\begin{center}
\begin{tikzpicture}[scale=1]
\tkzDefPoint(-1,3){L1_left}
\tkzDefPoint(8,3){L1_right}
\tkzDefPoint(-1,0){L2_left}
\tkzDefPoint(8,0){L2_right}

\tkzDefPoint(1,3){A}
\tkzDefPoint(0,0){B}
\tkzDefPoint(2.5,0){C}
\tkzDefPoint(5,3){D}
\tkzDefPoint(4.5,0){E}
\tkzDefPoint(6.5,0){Z}

\tkzDefPointBy[projection=onto L2_left--L2_right](A) \tkzGetPoint{H}
\tkzDefPointBy[projection=onto L2_left--L2_right](D) \tkzGetPoint{Theta}

\draw[l3] (A)--(B)--(C)--cycle;
\draw[l3] (D)--(E)--(Z)--cycle;
\draw[l3, dashed] (A)--(H);
\draw[l3, dashed] (D)--(Theta);

\draw[l3] (L1_left)--(L1_right) node[right] {$\varepsilon_1$};
\draw[l3] (L2_left)--(L2_right) node[right] {$\varepsilon_2$};

\tkzDrawPoints(A,B,C,D,E,Z,H,Theta)
\tkzLabelPoint[above](A){$A$}
\tkzLabelPoint[below](B){$B$}
\tkzLabelPoint[below](C){$\varGamma $}
\tkzLabelPoint[above](D){$\varDelta  $}
\tkzLabelPoint[below](E){$E$}
\tkzLabelPoint[below](Z){$Z$}
\tkzLabelPoint[below](H){$H$}
\tkzLabelPoint[below](Theta){$\varTheta $}

\tkzMarkRightAngle(A,H,C)
\tkzMarkRightAngle(D,Theta,E)
\end{tikzpicture}
\end{center}
\end{alist}
\end{alist}

\vspace{4mm}

% --- Άσκηση 13752 (13752.tex) ---
\Askhsh{13752}{4}
\begin{ylh}{1}
    \item 3.2. 1ο Κριτήριο ισότητας τριγώνων
    \item 3.12. Τριγωνική ανισότητα
    \item 4.4. Γωνίες με πλευρές παράλληλες.
\end{ylh}
% ===================================================================
%  Αρχείο: 13752.tex (Fragment)
%  Αυτόματη μετατροπή μέσω doc2latex.py
% ===================================================================



Σε τρίγωνο $AB\varGamma $ με $\hat{B}$ < 90\textsuperscript{ο} θεωρούμε τυχαίο σημείο $\varDelta  $ της πλευράς $A\varGamma $. Φέρουμε τμήμα $\varDelta  E$ ίσο και παράλληλο με την πλευρά $B\varGamma $ και από το σημείο $E$ φέρουμε τμήμα $EZ$ ίσο και παράλληλο με την πλευρά $AB$, όπως φαίνεται στο σχήμα.

\begin{alist}
\item Ένας μαθητής κάνει τους παρακάτω διαδοχικούς συλλογισμούς. Να χαρακτηρίσετε Σ (Σωστό) ή Λ (Λάθος) κάθε έναν από αυτούς.

\begin{alist}
\item
  Οι γωνίες Δ$\hat{E}$Ζ και Α$\hat{B}$Γ είναι γωνίες με πλευρές παράλληλες.
\item
  Οπότε Δ$\hat{E}$Ζ = Α$\hat{B}$Γ.
\item
  Τα τρίγωνα $\varDelta  EZ$ και $AB\varGamma $ είναι ίσα.
\item
  Το τμήμα $\varDelta   Z$ είναι ίσο με το τμήμα $A\varGamma $.

\monades{08}

\item Να αιτιολογήσετε τους χαρακτηρισμούς σας (Σ ή Λ) που αφορούν τους ισχυρισμούς 2. και 3. \monades{10}

\item Αν στα δεδομένα παραλείψουμε τη συνθήκη $\hat{B}$ < 90\textsuperscript{ο}, να συγκρίνετε τα τμήματα $A\varGamma $ και $\varDelta   Z$ για τα διάφορα είδη της γωνίας $\hat{B}$ και να δικαιολογήσετε τις απαντήσεις σας. \monades{07}


\begin{center}
\begin{tikzpicture}[scale=1]
\tkzDefPoint(0,4){A}
\tkzDefPoint(-2,0){B}
\tkzDefPoint(2,0){C}
\tkzDefPointOnLine[pos=0.4](A,C) \tkzGetPoint{D}
\tkzDefPointBy[translation=from B to C](D) \tkzGetPoint{E}
\tkzDefPointBy[translation=from B to A](E) \tkzGetPoint{Z}

\draw[l3] (A)--(B)--(C)--cycle;
\draw[l3] (D)--(E);
\draw[l3] (Z)--(E);
\draw[l3] (D)--(Z);

\tkzDrawPoints(A,B,C,D,E,Z)
\tkzLabelPoint[above](A){$A$}
\tkzLabelPoint[left](B){$B$}
\tkzLabelPoint[below](C){$\varGamma $}
\tkzLabelPoint[above right](D){$\varDelta  $}
\tkzLabelPoint[right](E){$E$}
\tkzLabelPoint[above](Z){$Z$}

\tkzMarkSegments[mark=s|](B,C D,E)
\tkzMarkSegments[mark=s||](A,B Z,E)
\end{tikzpicture}
\end{center}
\end{alist}
\end{alist}

\vspace{4mm}

% --- Άσκηση 13755 (13755.tex) ---
\Askhsh{13755}{2}
\begin{ylh}{1}
    \item 3.4. 3ο Κριτήριο ισότητας τριγώνων
    \item 5.2. Παραλληλόγραμμα.
\end{ylh}
% ===================================================================
%  Αρχείο: 13755.tex (Fragment)
%  Αυτόματη μετατροπή μέσω doc2latex.py
% ===================================================================



Δίνεται τρίγωνο $AB\varGamma $. Από την κορυφή $A$ φέρουμε ευθεία (ε) παράλληλη προς τη $B\varGamma $. Από το τυχαίο σημείο $\varDelta  $ της πλευράς $B\varGamma $ φέρουμε τις παράλληλες προς την $AB$ και

$A\varGamma $, οι οποίες τέμνουν την ευθεία (ε) στα σημεία $E$ και $Z$ αντίστοιχα.

Να αποδείξετε ότι :

\begin{alist}
\item τα τετράπλευρα $ZA\varGamma \varDelta  $ και Α$B\varDelta  $Ε είναι παραλληλόγραμμα. \monades{10}

\item τα τρίγωνα $AB\varGamma $ και $\varDelta  EZ$ είναι ίσα. \monades{15}


\begin{center}
\begin{tikzpicture}[scale=1]
\tkzDefPoint(0,3){A}
\tkzDefPoint(-2,0){B}
\tkzDefPoint(2,0){C}
\tkzDefPointOnLine[pos=0.3](B,C) \tkzGetPoint{D}

\tkzDefLine[parallel=through A](B,C) \tkzGetPoint{e_dir}
\tkzDefLine[parallel=through D](A,B) \tkzGetPoint{ab_dir}
\tkzDefLine[parallel=through D](A,C) \tkzGetPoint{ac_dir}

\tkzInterLL(A,e_dir)(D,ab_dir) \tkzGetPoint{E}
\tkzInterLL(A,e_dir)(D,ac_dir) \tkzGetPoint{Z}

\draw[l3] ($(Z)-(1,0)$)--($(E)+(1,0)$) node[right] {$\varepsilon$};
\draw[l3] (A)--(B)--(C)--cycle;
\draw[l3, dashed] (D)--(E);
\draw[l3, dashed] (D)--(Z);

\tkzDrawPoints(A,B,C,D,E,Z)
\tkzLabelPoint[above](A){$A$}
\tkzLabelPoint[below](B){$B$}
\tkzLabelPoint[below](C){$\varGamma $}
\tkzLabelPoint[below](D){$\varDelta  $}
\tkzLabelPoint[above](E){$E$}
\tkzLabelPoint[above](Z){$Z$}
\end{tikzpicture}
\end{center}
\end{alist}

\vspace{4mm}

% --- Άσκηση 13757 (13757.tex) ---
\Askhsh{13757}{2}
\begin{ylh}{1}
    \item 3.16. Σχετικές θέσεις δυο κύκλων.
\end{ylh}
% ===================================================================
%  Αρχείο: 13757.tex (Fragment)
%  Αυτόματη μετατροπή μέσω doc2latex.py
% ===================================================================



Δίνονται δύο κύκλοι (Κ,2) και (Λ,5).

\begin{alist}
\item Να υπολογίσετε το μήκος της διακέντρου $K\varLambda $, αν οι κύκλοι εφάπτονται εξωτερικά.

\monades{6}

\item Να υπολογίσετε το μήκος της διακέντρου $K\varLambda $, αν οι κύκλοι εφάπτονται

εσωτερικά. \monades{6}

\item Μεταξύ ποιών τιμών βρίσκεται το μήκος της διακέντρου $K\varLambda $, αν ο κύκλος (Κ,2) βρίσκεται στο εσωτερικό του κύκλου (Λ,5); Να αιτιολογήσετε την απάντησή σας.

\monades{6}

\item Μεταξύ ποιών τιμών βρίσκεται το μήκος της διακέντρου $K\varLambda $, αν οι κύκλοι

τέμνονται; Να αιτιολογήσετε την απάντησή σας. \monades{7}
\end{alist}

\vspace{4mm}

% --- Άσκηση 13758 (13758.tex) ---
\Askhsh{13758}{2}
\begin{ylh}{1}
    \item 3.16. Σχετικές θέσεις δυο κύκλων.
\end{ylh}
% ===================================================================
%  Αρχείο: 13758.tex (Fragment)
%  Αυτόματη μετατροπή μέσω doc2latex.py
% ===================================================================



Δίνονται δύο κύκλοι (Κ,3) και (Λ,8). Να βρείτε τη σχετική θέση των δύο κύκλων,

αιτιολογώντας την απάντησή σας, όταν:

\begin{alist}
\item $K\varLambda $ = 13. \monades{5}
\item $K\varLambda $ = 2. \monades{5}
\item $K\varLambda $ = 5. \monades{5}
\item $K\varLambda $ = 11. \monades{5}
\item $K\varLambda $ = 9. \monades{5}
\end{alist}

\vspace{4mm}

% --- Άσκηση 13759 (13759.tex) ---
\Askhsh{13759}{2}
\begin{ylh}{1}
    \item 3.14. Σχετικές θέσεις ευθείας και κύκλου.
\end{ylh}
% ===================================================================
%  Αρχείο: 13759.tex (Fragment)
%  Αυτόματη μετατροπή μέσω doc2latex.py
% ===================================================================



Δίνεται κύκλος με κέντρο Ο και ακτίνα ρ = 6. Έστω d η απόσταση του κέντρου Ο του κύκλου από μια ευθεία (ε). Να βρείτε τη σχετική θέση του κύκλου και της ευθείας (ε) στις εξής περιπτώσεις:

\begin{alist}
\item d = 3. \monades{9}

\item d = 6. \monades{8}

\item d = 9. \monades{8}
\end{alist}

\vspace{4mm}

% --- Άσκηση 13767 (13767.tex) ---
\Askhsh{13767}{2}
\begin{ylh}{1}
    \item 3.1. Είδη και στοιχεία τριγώνων
    \item 3.2. 1ο Κριτήριο ισότητας τριγώνων
    \item 4.6. Άθροισμα γωνιών τριγώνου
    \item 5.4. Ρόμβος.
\end{ylh}
% ===================================================================
%  Αρχείο: 13767.tex (Fragment)
%  Αυτόματη μετατροπή μέσω doc2latex.py
% ===================================================================



Σε ευθεία ε θεωρούμε τα διαδοχικά σημεία $A$, Β και $\varGamma$ έτσι ώστε ΑΒ$\  = \ $$B\varGamma $. Στη συνέχεια, σχεδιάζουμε τα ισόπλευρα τρίγωνα Α$B\varDelta  $ και $B\varGamma $Ε προς το ίδιο ημιεπίπεδο ως προς την ευθεία ε όπως φαίνεται στο παρακάτω σχήμα.

\begin{alist}
\item Να υπολογίσετε τη γωνία $\text{Δ}\hat{B}\text{E}$.

\monades{7}

\item Να αποδείξετε ότι το τρίγωνο $B\varDelta  $Ε είναι ισόπλευρο.

\monades{10}

\item Να αποδείξετε ότι το τετράπλευρο $A\varDelta  EB$ είναι ρόμβος.

\monades{8}


\begin{center}
\begin{tikzpicture}[scale=0.9]
% Line ε with A, B, Γ consecutive, $AB$ = $B\varGamma $
\tkzDefPoint(0,0){A}
\tkzDefPoint(3,0){B}
\tkzDefPoint(6,0){Gamma}

% Equilateral triangle $AB\varDelta  $ (above the line)
\tkzDefTriangle[equilateral](A,B) \tkzGetPoint{Delta}

% Equilateral triangle $B\varGamma E$ (above the line)
\tkzDefTriangle[equilateral](B,Gamma) \tkzGetPoint{E}

% Draw line ε
\draw[l3] (-0.5,0)--(6.5,0);

% Draw equilateral triangles
\draw[l3] (A)--(Delta)--(B);
\draw[l3] (B)--(E)--(Gamma);

% Draw segment $\varDelta  E$
\draw[l3] (Delta)--(E);

% Points and labels
\tkzDrawPoints(A,B,Gamma,Delta,E)
\tkzLabelPoint[below](A){$A$}
\tkzLabelPoint[below](B){$B$}
\tkzLabelPoint[below](Gamma){$\varGamma $}
\tkzLabelPoint[above left](Delta){$\varDelta  $}
\tkzLabelPoint[above right](E){$ E$}

% Label ε
\node[below] at (6.5,0) {$\varepsilon$};
\end{tikzpicture}
\end{center}
\end{alist}

\vspace{4mm}

% --- Άσκηση 13816 (13816.tex) ---
\Askhsh{13816}{2}
\begin{ylh}{1}
    \item 5.2. Παραλληλόγραμμα.
\end{ylh}
% ===================================================================
%  Αρχείο: 13816.tex (Fragment)
%  Αυτόματη μετατροπή μέσω doc2latex.py
% ===================================================================



Δίνεται παραλληλόγραμμο $AB\varGamma \varDelta  $ τέτοιο, ώστε $A\varDelta  $ < $AB$ και $M$ το μέσο της $B\varGamma $.

Προεκτείνουμε την $AM$ προς το $M$ κατά τμήμα $ME$ = $AM$.

Να αποδείξετε ότι :

\begin{alist}
\item το τετράπλευρο Α$BE$Γ είναι παραλληλόγραμμο. \monades{10}

\item τα σημεία $\varDelta  $, Γ και $E$ είναι συνευθειακά. \monades{15}


\begin{center}
\begin{tikzpicture}[scale=1]
\tkzDefPoint(0,0){B}
\tkzDefPoint(3,0){C}
\tkzDefPoint(1,3){A}
\tkzDefPointBy[translation=from B to C](A) \tkzGetPoint{D}

\tkzDefMidPoint(B,C) \tkzGetPoint{M}
\tkzDefPointBy[symmetry=center M](A) \tkzGetPoint{E}

\draw[l3] (A)--(B)--(C)--(D)--cycle;
\draw[l3] (A)--(E);
\draw[l3] (B)--(E);
\draw[l3] (C)--(E);
\draw[l3, dashed] (D)--(E);

\tkzDrawPoints(A,B,C,D,M,E)
\tkzLabelPoint[above left](A){$A$}
\tkzLabelPoint[below left](B){$B$}
\tkzLabelPoint[left](M){$M$}
\tkzLabelPoint[above right](C){$\varGamma $}
\tkzLabelPoint[above right](D){$\varDelta  $}
\tkzLabelPoint[below](E){$E$}

\tkzMarkSegments[mark=s|](B,M M,C)
\tkzMarkSegments[mark=s||](A,M M,E)
\end{tikzpicture}
\end{center}
\end{alist}

\vspace{4mm}

% --- Άσκηση 13817 (13817.tex) ---
\Askhsh{13817}{2}
\begin{ylh}{1}
    \item 3.15. Εφαπτόμενα τμήματα.
\end{ylh}
% ===================================================================
%  Αρχείο: 13817.tex (Fragment)
%  Αυτόματη μετατροπή μέσω doc2latex.py
% ===================================================================



Δίνεται κύκλος με κέντρο Ο και ακτίνα ρ. Σε σημείο $B$ του κύκλου φέρουμε

εφαπτόμενη ευθεία (ε). Θεωρούμε στην ευθεία (ε) δύο σημεία $A$ και $\varGamma$ εκατέρωθεν του Β έτσι ώστε $BA$ < $B\varGamma $ και από τα σημεία αυτά, φέρουμε τα εφαπτόμενα τμήματα $AZ$ και $\varGamma M$ στον κύκλο.

\begin{alist}
\item Να γράψετε τα ευθύγραμμα τμήματα τα οποία είναι ίσα, αιτιολογώντας την

απάντησή σας. \monades{15}

\item Να αποδείξετε ότι $A\varGamma $ = $AZ$ + $M\varGamma $. \monades{10}


\begin{center}
\begin{tikzpicture}[scale=1]
\tkzDefPoint(0,3){O}
\tkzDefPoint(0,0){B}
\tkzDefPoint(-2,0){A}
\tkzDefPoint(4,0){C}

\tkzDrawCircle(O,B)
\tkzDefTangent[from=A](O,B) \tkzGetPoints{Z}{Z_temp} % B is one tangent point, Z is the other
\tkzDefTangent[from=C](O,B) \tkzGetPoints{M_temp}{M} % B is one tangent point, M is the other

\draw[l3] ($(A)-(1,0)$)--($(C)+(1,0)$) node[right] {$\varepsilon$};
\draw[l3] (A)--(Z);
\draw[l3] (C)--(M);
\draw[l3, dashed] (O)--(B);
\draw[l3, dashed] (O)--(Z);
\draw[l3, dashed] (O)--(M);

\tkzDrawPoints(O,B,A,C,Z,M)
\tkzLabelPoint[above](O){$O$}
\tkzLabelPoint[below](B){$B$}
\tkzLabelPoint[below](A){$A$}
\tkzLabelPoint[below](C){$\varGamma $}
\tkzLabelPoint[left](Z){$Z$}
\tkzLabelPoint[right](M){$M$}

\tkzMarkRightAngle(O,B,A)
\tkzMarkRightAngle(O,Z,A)
\tkzMarkRightAngle(O,M,C)
\end{tikzpicture}
\end{center}
\end{alist}

\vspace{4mm}

% --- Άσκηση 13822 (13822.tex) ---
\Askhsh{13822}{4}
\begin{ylh}{1}
    \item 4.2. Τέμνουσα δύο ευθειών - Ευκλείδειο αίτημα.
\end{ylh}
% ===================================================================
%  Αρχείο: 13822.tex (Fragment)
%  Αυτόματη μετατροπή μέσω doc2latex.py
% ===================================================================



Δίνονται οι ευθείες (ε) και (ψ).

\begin{alist}
\item Αν η γωνία $B\hat{A}\varGamma $ είναι με$\varGamma  A$λύτερη από την $A\hat{B}\psi$:

\begin{rlist}
\item Να αποδείξετε ότι $B\hat{A}\varepsilon + A\hat{B}\psi < \ 180\degree$. \monades{6}

\item Να αποδείξετε ότι οι ευθείες ε και ψ τέμνονται. Σε ποιο από τα ημιεπίπεδα που χωρίζει το επίπεδο η $AB$ βρίσκεται το σημείο τομής των ε και ψ και γιατί;

\monades{6}

\end{rlist}
\item Να διατυπώσετε την πρόταση που αποδείχθηκε στο α) για τις εντός και εναλλάξ γωνίες δύο ευθειών που τέμνονται από τρίτη και το σημείο τομής των ευθειών αυτών.

\monades{7}

\item Αν ισχύει $B\hat{A}\varGamma  < Α\hat{B}\psi$, τότε σε ποιο από τα ημιεπίπεδα που χωρίζει το επίπεδο η $AB$ βρίσκεται το σημείο τομής των ε και ψ και γιατί; \monades{6}


\begin{center}
\begin{tikzpicture}[scale=0.8]
% Two lines ε and ψ cut by transversal $AB$
\tkzDefPoint(-3,0){e1}
\tkzDefPoint(3,0){e2}
\tkzDefPoint(-3.5,2.5){p1}
\tkzDefPoint(3.5,2.5){p2}

% Transversal
\tkzDefPoint(-1,0){A}
\tkzDefPoint(1,2.5){B}

% Point $\varGamma$ on one side
\tkzDefPoint(3,0){Gamma}

% Draw lines
\draw[l3] (e1)--(e2) node[right]{$\varepsilon$};
\draw[l3] (p1)--(p2) node[right]{$\psi$};
\draw[l3] ($(A)!-0.2!(B)$)--($(B)!-0.2!(A)$);

% Points and labels
\tkzDrawPoints(A,B)
\tkzLabelPoint[below left](A){$A$}
\tkzLabelPoint[above left](B){$B$}

% Angle marks
\tkzMarkAngle[size=0.6](Gamma,A,B)
\end{tikzpicture}
\end{center}
\end{alist}

\vspace{4mm}

% --- Άσκηση 13823 (13823.tex) ---
\Askhsh{13823}{4}
\begin{ylh}{1}
    \item 3.16. Σχετικές θέσεις δυο κύκλων.
\end{ylh}
% ===================================================================
%  Αρχείο: 13823.tex (Fragment)
%  Αυτόματη μετατροπή μέσω doc2latex.py
% ===================================================================



\begin{alist}
\item Στο παρακάτω σχήμα για τους κύκλους (Α, ρ) και (Β, R) ισχύει $\rho < R$ και $AB$ = 6.


\begin{center}
\begin{tikzpicture}[scale=0.65]
% Two circles (A,ρ) and (B,R) with ρ<R, $AB$=6
\tkzDefPoint(0,0){A}
\tkzDefPoint(6,0){B}

% Circle (A,ρ) with ρ=1.5
\tkzDefPoint(1.5,0){rA}
\tkzDrawCircle(A,rA)

% Circle (B,R) with R=2.5
\tkzDefPoint(8.5,0){rB}
\tkzDrawCircle(B,rB)

% Γ on circle (A,ρ), K on circle (B,R) — on line $AB$
\tkzDefPoint(-1.5,0){Gamma}
\tkzDefPoint(8.5,0){K}

% Draw $AB$ line
\draw[l3] (Gamma)--(K);

% Points and labels
\tkzDrawPoints(A,B,Gamma,K)
\tkzLabelPoint[below](A){$A$}
\tkzLabelPoint[below](B){$B$}
\tkzLabelPoint[below left](Gamma){$\varGamma $}
\tkzLabelPoint[below right](K){$ K$}
\tkzLabelSegment[above](A,rA){$\rho$}
\tkzLabelSegment[above](B,rB){$R$}
\end{tikzpicture}
\end{center}

\begin{rlist}
\item Να αποδείξετε ότι $ΒK - Α\varGamma  < $AB$ < $BK$ + Α\varGamma $.

\item Παρακάτω γράφονται οι ιδιότητες 1 και 2. Ποιο σημείο από τα Κ και $\varGamma$ έχει την ιδιότητα 1, ποιο την ιδιότητα 2 και ποιο έχει και τις δύο;

Ιδιότητα 1: «Το σημείο απέχει R από το Β.»

Ιδιότητα 2: «Το σημείο απέχει ρ από το Α.»

Να αιτιολογήσετε τις απαντήσεις σας.

\monades{16}

\end{rlist}
\item Ο χάρτης ενός κρυμμένου θησαυρού έχει δύο σταθερά σημεία $A$ και Β, τα οποία απέχουν μεταξύ τους 6. Επίσης γράφει ότι ο θησαυρός είναι κρυμμένος σε ένα σημείο το οποίο απέχει 3 από το Α του χάρτη και 2 από το Β του χάρτη. Μπορεί να είναι σωστή η πληροφορία που δίνει ο χάρτης για να βρει κανείς το θησαυρό;

\monades{9}
\end{alist}

\vspace{4mm}

% --- Άσκηση 13824 (13824.tex) ---
\Askhsh{13824}{2}
\begin{ylh}{1}
    \item 3.3. 2ο Κριτήριο ισότητας τριγώνων
    \item 4.2. Τέμνουσα δύο ευθειών - Ευκλείδειο αίτημα
    \item 5.2. Παραλληλόγραμμα.
\end{ylh}
% ===================================================================
%  Αρχείο: 13824.tex (Fragment)
%  Αυτόματη μετατροπή μέσω doc2latex.py
% ===================================================================



Δίνεται τραπέζιο $AB\varGamma \varDelta  $ με βάσεις $AB$ και $\varGamma \varDelta  $. Αν $E$ και $Z$ τα μέσα των $\varGamma \varDelta  $ και $BE$ αντίστοιχα και Θ το σημείο τομής της $AB$ και της προέκτασης της $\varGamma Z$, να αποδείξετε ότι:

\begin{alist}
\item Τα τρίγωνα $\varGamma EZ$, $\varTheta BZ$ είναι ίσα. \monades{13}

\item $E\varGamma $=$\varTheta B$. \monades{5}

\item Το τετράπλευρο $EB\varTheta\varDelta  $ είναι παραλληλόγραμμο. \monades{7}


\begin{center}
\begin{tikzpicture}[scale=1]
\tkzDefPoint(0,0){A}
\tkzDefPoint(6,0){B}
\tkzDefPoint(4,3){C}
\tkzDefPoint(1,3){D}
\tkzDefMidPoint(D,C) \tkzGetPoint{E}

\tkzDefMidPoint(B,E) \tkzGetPoint{Z}
\tkzInterLL(C,Z)(A,B) \tkzGetPoint{Theta}

\draw[l3] (A)--(B)--(C)--(D)--cycle;
\draw[l3] (C)--(Theta);
\draw[l3] (B)--(E);
\draw[l3, dashed] (B)--(Theta); % $AB$ extension

\tkzDrawPoints(A,B,C,D,E,Z,Theta)
\tkzLabelPoint[below left](A){$A$}
\tkzLabelPoint[below](B){$B$}
\tkzLabelPoint[above right](C){$\varGamma $}
\tkzLabelPoint[above left](D){$\varDelta  $}
\tkzLabelPoint[above](E){$E$}
\tkzLabelPoint[above right](Z){$Z$}
\tkzLabelPoint[below right](Theta){$\varTheta $}

\tkzMarkSegments[mark=s|](C,E E,D)
\tkzMarkSegments[mark=s||](B,Z Z,E)
\end{tikzpicture}
\end{center}
\end{alist}

\vspace{4mm}

% --- Άσκηση 13825 (13825.tex) ---
\Askhsh{13825}{2}
\begin{ylh}{1}
    \item 5.2. Παραλληλόγραμμα.
\end{ylh}
% ===================================================================
%  Αρχείο: 13825.tex (Fragment)
%  Αυτόματη μετατροπή μέσω doc2latex.py
% ===================================================================



Δίνεται τρίγωνο $AB\varGamma $. Από το μέσο $M$ της $B\varGamma $ γράφουμε ευθύγραμμο τμήμα $M\varDelta  $ ίσο και παράλληλο προς την $BA$ και ένα άλλο ευθύγραμμο τμήμα $ME$ ίσο και παράλληλο προς την $\varGamma  A$ (τα σημεία $\varDelta  $ και $E$ βρίσκονται στο ημιεπίπεδο που ορίζεται από τη $B\varGamma $ και το σημείο $A$). Να αποδείξετε ότι:

\begin{alist}
\item Τα τετράπλευρα $A\varDelta  MB$ και $A\varGamma ME$ είναι παραλληλόγραμμα. \monades{12}

\item $\varDelta  A$=$AE$. \monades{13}
\end{alist}

\vspace{4mm}

% --- Άσκηση 13826 (13826.tex) ---
\Askhsh{13826}{2}
\begin{ylh}{1}
    \item 3.1. Είδη και στοιχεία τριγώνων
    \item 3.2. 1ο Κριτήριο ισότητας τριγώνων
    \item 3.4. 3ο Κριτήριο ισότητας τριγώνων.
\end{ylh}
% ===================================================================
%  Αρχείο: 13826.tex (Fragment)
%  Αυτόματη μετατροπή μέσω doc2latex.py
% ===================================================================



Τα τρίγωνα $ABK$ και $\varGamma \varDelta  \varLambda $ του σχήματος έχουν $AB$ = $\varGamma \varDelta  $ = $AK$= $\varGamma \varLambda $ και $\hat{A}$ = $\hat{\varGamma }$.

\begin{alist}
\item Να αποδείξετε ότι τα τρίγωνα $ABK$ και $\varGamma \varDelta  \varLambda $ είναι ίσα και ότι έχουν $BK$ = $\varDelta  \varLambda $.

\monades{12}

\item Έστω ότι Λ και Κ είναι τα μέσα των $BK$ και $\varDelta  \varLambda $ αντίστοιχα:

\begin{alist}
\item
  Να εξετάσετε αν τα τμήματα $B\varLambda $, $\varLambda K$ και $K\varDelta  $ είναι ίσα. Να αιτιολογήσετε την απάντησή σας. \monades{5}
\item
  Να αποδείξετε ότι οι $A\varLambda $ και $\varGamma K$ είναι κάθετες στην ευθεία $K\varLambda $.

\monades{8}


\begin{center}
\begin{tikzpicture}[scale=1.5]
\tkzDefPoint(0,0){L_line_start}
\tkzDefPoint(6,0){L_line_end}
\tkzDefPoint(2,0){K}
\tkzDefPoint(4,0){L}

% Define A, B such that $ABK$ is a triangle. B is on line, A is above.
\tkzDefPoint(0.5,0){B}
\tkzDefPoint(1.5,2){A}

% Define \varGamma , \Delta. Since $AB$=$CD$, $AK$=$CL$, \hat{A}=\hat{\varGamma }.
% Congruent triangles $ABK$ and \varGamma  \Delta L
% \varGamma  corresponds to A, \Delta corresponds to B, L corresponds to K.
% \Delta is on the line, \varGamma  is above.
\tkzDefPoint(5.5,0){D} % \Delta
\tkzDefPoint(4.5,2){C} % \varGamma 

\draw[l3] (B)--(L_line_end);
\draw[l3] (A)--(B);
\draw[l3] (A)--(K);
\draw[l3] (C)--(D);
\draw[l3] (C)--(L);

\draw[l3, dashed] (A)--(L);
\draw[l3, dashed] (C)--(K);

\tkzDrawPoints(A,B,C,D,K,L)
\tkzLabelPoint[above](A){$A$}
\tkzLabelPoint[below](B){$B$}
\tkzLabelPoint[above](C){$\varGamma $}
\tkzLabelPoint[below](D){$\varDelta  $}
\tkzLabelPoint[below](K){$K$}
\tkzLabelPoint[below](L){$\Lambda$}

\tkzMarkSegments[mark=s|](A,B C,D)
\tkzMarkSegments[mark=s||](A,K C,L)
\tkzMarkAngle[size=0.4,mark=none](B,A,K)
\tkzMarkAngle[size=0.4,mark=none](D,C,L)
\end{tikzpicture}
\end{center}
\end{alist}
\end{alist}

\vspace{4mm}

% --- Άσκηση 13828 (13828.tex) ---
\Askhsh{13828}{2}
\begin{ylh}{1}
    \item 3.10. Σχέση εξωτερικής και απέναντι γωνίας
    \item 4.2. Τέμνουσα δύο ευθειών - Ευκλείδειο αίτημα
    \item 5.9. Μια ιδιότητα του ορθογώνιου τριγώνου.
\end{ylh}
% ===================================================================
%  Αρχείο: 13828.tex (Fragment)
%  Αυτόματη μετατροπή μέσω doc2latex.py
% ===================================================================



Σε τραπέζιο $AB\varGamma \varDelta  $ η διαγώνιος $B\varDelta  $ είναι κάθετη στην πλευρά $A\varDelta  $ και η πλευρά $\varGamma B$ κάθετη στη βάση $AB$. Αν $A\varDelta  $=4 και $AB$=8 τότε:

\begin{alist}
\item Nα υπολογιστεί η γωνία Δ$\hat{A}$Β. \monades{12}

\item Να αποδείξετε ότι η διαγώνιος $B\varDelta  $ του τραπεζίου $AB\varGamma \varDelta  $ είναι διπλάσια της πλευράς του $B\varGamma $. \monades{13}


\begin{center}
\begin{tikzpicture}[scale=0.8]
\tkzDefPoint(0,0){A}
\tkzDefPoint(8,0){B} % $AB$ = 8
\tkzDefPoint(2,3.464){D} % $AD$ = 4. Triangle $ABD$ has angle D=90. $AB$=8, $AD$=4 implies angle A=60. x=4*cos(60)=2, y=4*sin(60)=3.464

\tkzDefLine[parallel=through D](A,B) \tkzGetPoint{cd_dir}
\tkzDefLine[orthogonal=through B](A,B) \tkzGetPoint{bc_dir}
\tkzInterLL(D,cd_dir)(B,bc_dir) \tkzGetPoint{C}

\draw[l3] (A)--(B)--(C)--(D)--cycle;
\draw[l3] (B)--(D);

\tkzDrawPoints(A,B,C,D)
\tkzLabelPoint[left](A){$A$}
\tkzLabelPoint[right](B){$B$}
\tkzLabelPoint[above right](C){$\varGamma $}
\tkzLabelPoint[above left](D){$\varDelta  $}

\tkzMarkRightAngle(A,D,B)
\tkzMarkRightAngle(A,B,C)
\end{tikzpicture}
\end{center}
\end{alist}

\vspace{4mm}

% --- Άσκηση 13829 (13829.tex) ---
\Askhsh{13829}{2}
\begin{ylh}{1}
    \item 3.3. 2ο Κριτήριο ισότητας τριγώνων
    \item 4.2. Τέμνουσα δύο ευθειών - Ευκλείδειο αίτημα
    \item 5.2. Παραλληλόγραμμα.
\end{ylh}
% ===================================================================
%  Αρχείο: 13829.tex (Fragment)
%  Αυτόματη μετατροπή μέσω doc2latex.py
% ===================================================================



Δίνεται παραλληλόγραμμο $AB\varGamma \varDelta  $ και Ο το σημείο τομής των διαγωνίων του. Θεωρούμε τα σημεία $E$ και $Z$ των τμημάτων $AO$ και $\varGamma O$ αντίστοιχα, τέτοια ώστε $AE$=$\varGamma Z$.

\begin{alist}
\item Να αποδείξετε ότι τα τρίγωνα Α$E\varDelta  $ και $\varGamma ZB$ είναι ίσα. \monades{12}

\item Να αποδείξετε ότι το τετράπλευρο $\varDelta   EBZ$ είναι παραλληλόγραμμο. \monades{13}
\end{alist}

\vspace{4mm}

% --- Άσκηση 13831 (13831.tex) ---
\Askhsh{13831}{2}
\begin{ylh}{1}
    \item 4.6. Άθροισμα γωνιών τριγώνου
    \item 5.9. Μια ιδιότητα του ορθογώνιου τριγώνου.
\end{ylh}
% ===================================================================
%  Αρχείο: 13831.tex (Fragment)
%  Αυτόματη μετατροπή μέσω doc2latex.py
% ===================================================================



Ένα τρίγωνο $AB\varGamma $ έχει $\hat{A} = 90^{Ο}$.

\begin{alist}
\item Να σχεδιάσετε ένα τέτοιο τρίγωνο αν επιπλέον γνωρίζετε ότι $AB$ > $A\varGamma $. Ποια είναι η μικρότερη γωνία του τριγώνου και γιατί; \monades{10}

\item Αν για το τρίγωνο που σας ζητήθηκε να σχεδιάσετε στο α) ερώτημα γνωρίζετε επιπλέον ότι η μια από τις οξείες γωνίες του είναι ίση με $30\degree$, τότε να απαντήσετε στα παρακάτω:

\begin{rlist}
\item Πόσες μοίρες θα είναι η γωνία $\hat{B}$ και πόσες η γωνία $\hat{\varGamma }$; \monades{8}

\item Ποια πλευρά του τριγώνου είναι ίση με το μισό της υποτείνουσας; \monades{7}
\end{rlist}
\end{alist}

\vspace{4mm}

% --- Άσκηση 13832 (13832.tex) ---
\Askhsh{13832}{2}
\begin{ylh}{1}
    \item 5.4. Ρόμβος.
\end{ylh}
% ===================================================================
%  Αρχείο: 13832.tex (Fragment)
%  Αυτόματη μετατροπή μέσω doc2latex.py
% ===================================================================



Στο σχήμα το $M$ είναι μέσο των τμημάτων $A\varGamma $ και $B\varDelta  $. Επίσης $A\hat{Μ}B = \varGamma \hat{Μ}B$.

\begin{alist}
\item Να αποδείξετε ότι:

\begin{rlist}
\item Οι $A\varGamma $ και $B\varDelta  $ είναι κάθετες. \monades{10}

\item Το $AB\varGamma \varDelta  $ είναι ρόμβος. \monades{8}

\end{rlist}
\item Το $AB\varGamma \varDelta  $ είναι η κάτοψη ενός κήπου. Για να περιφράξουμε τον κήπο χρειαζόμαστε 30 μέτρα φράχτη. Αν αφήσουμε την πλευρά $AB$ του κήπου χωρίς φράχτη πόσα μέτρα φράχτη θα χρειαστούμε για τις υπόλοιπες πλευρές;

\monades{7}


\begin{center}
\begin{tikzpicture}[scale=0.8]
% Quadrilateral $AB\varGamma \varDelta  $ with M midpoint of $A\varGamma $ and $B\varDelta  $
% For the figure, place M at origin, $A\varGamma $ horizontal, $B\varDelta  $ at an angle
\tkzDefPoint(0,0){M}
\tkzDefPoint(-3,0){A}
\tkzDefPoint(3,0){Gamma}
\tkzDefPoint(0,2.3){B}
\tkzDefPoint(0,-2.3){Delta}

% Draw quadrilateral
\draw[l3] (A)--(B)--(Gamma)--(Delta)--cycle;

% Draw diagonals
\draw[l3] (A)--(Gamma);
\draw[l3] (B)--(Delta);

% Angle marks at M
\tkzMarkAngle[size=0.5](B,M,A)
\tkzMarkAngle[size=0.5](Gamma,M,B)

% Points and labels
\tkzDrawPoints(A,B,Gamma,Delta,M)
\tkzLabelPoint[left](A){$A$}
\tkzLabelPoint[above](B){$B$}
\tkzLabelPoint[right](Gamma){$\varGamma $}
\tkzLabelPoint[below](Delta){$\varDelta  $}
\tkzLabelPoint[below right](M){$\mathrm{M}$}
\end{tikzpicture}
\end{center}
\end{alist}

\vspace{4mm}

% --- Άσκηση 13833 (13833.tex) ---
\Askhsh{13833}{2}
\begin{ylh}{1}
    \item 3.6. Κριτήρια ισότητας ορθογώνιων τριγώνων
    \item 5.2. Παραλληλόγραμμα.
\end{ylh}
% ===================================================================
%  Αρχείο: 13833.tex (Fragment)
%  Αυτόματη μετατροπή μέσω doc2latex.py
% ===================================================================



Στο παρακάτω σχήμα το $\varGamma \varDelta  $ είναι ύψος του τριγώνου $AB\varGamma $, το $EH$ είναι ύψος του τριγώνου Α$BE$ και η $BA$ είναι διάμεσος του τριγώνου $BE$Γ.

\begin{alist}
\item Να αποδείξετε ότι τα τρίγωνα $A\varGamma \varDelta  $ και $AEH$ είναι ίσα. \monades{10}

\item Να αποδείξετε ότι $AH$=$A\varDelta  $. \monades{5}

\item Να αποδείξετε ότι το τετράπλευρο $\varGamma \varDelta  EH$ είναι παραλληλόγραμμο. \monades{10}


\begin{center}
\begin{tikzpicture}[scale=0.8]
\tkzDefPoint(0,0){B}
\tkzDefPoint(5,0){C} % \varGamma 
\tkzDefPoint(2,4){A}

% "$BA$ είναι διάμεσος του τριγώνου $BE\varGamma $", so A is the midpoint of $E\varGamma $
\tkzDefPointBy[symmetry=center A](C) \tkzGetPoint{E}

\tkzDefPointBy[projection=onto A--B](C) \tkzGetPoint{D}
\tkzDefPointBy[projection=onto A--B](E) \tkzGetPoint{H}

\draw[l3] (A)--(B)--(C)--cycle;
\draw[l3] (A)--(E);
\draw[l3] (B)--(E);
\draw[l3] (C)--(D);
\draw[l3] (E)--(H);
\draw[l3, dashed] (B)--(H); % Extending $AB$ to H if needed
\draw[l3, dashed] (B)--(D); % Extending $AB$ to D if needed
\draw[l3, dashed] (D)--(H); % For parallelogram $CDHE$

% If H is on the extension of $AB$ towards A, and D is on $AB$
% Let's connect C to H and D to E to form the parallelogram.
\draw[l3, dashed] (C)--(H);
\draw[l3, dashed] (E)--(D);

\tkzDrawPoints(A,B,C,E,D,H)
\tkzLabelPoint[left](A){$A$}
\tkzLabelPoint[below](B){$B$}
\tkzLabelPoint[below](C){$\varGamma $}
\tkzLabelPoint[right](E){$E$}
\tkzLabelPoint[below right](D){$\varDelta  $}
\tkzLabelPoint[above left](H){$H$}

\tkzMarkRightAngle(C,D,B)
\tkzMarkRightAngle(E,H,A)
\tkzMarkSegments[mark=s|](E,A A,C)
\end{tikzpicture}
\end{center}
\end{alist}

\vspace{4mm}

% --- Άσκηση 13834 (13834.tex) ---
\Askhsh{13834}{2}
\begin{ylh}{1}
    \item 3.2. 1ο Κριτήριο ισότητας τριγώνων
    \item 5.2. Παραλληλόγραμμα.
\end{ylh}
% ===================================================================
%  Αρχείο: 13834.tex (Fragment)
%  Αυτόματη μετατροπή μέσω doc2latex.py
% ===================================================================



Σε τυχαίο τρίγωνο $AB\varGamma $ φέρουμε τη διάμεσό του $AM$. Προεκτείνουμε την πλευρά $B\varGamma $ προς το μέρος του Β κατά τμήμα $BZ$=$B\varGamma $ και προς το μέρος του $\varGamma$ κατά τμήμα $\varGamma H$=$B\varGamma $, επίσης προεκτείνουμε τη διάμεσο $AM$ κατά τμήμα $ME$=$AM$.

\begin{alist}
\item Να αποδείξετε ότι τα τρίγωνα $AMZ$ και $EMH$ είναι ίσα. \monades{12}

\item Να αποδείξετε ότι το τετράπλευρο $AHEZ$ είναι παραλληλόγραμμο. \monades{13}
\end{alist}

\vspace{4mm}

% --- Άσκηση 13835 (13835.tex) ---
\Askhsh{13835}{2}
\begin{ylh}{1}
    \item 3.12. Τριγωνική ανισότητα
    \item 3.16. Σχετικές θέσεις δυο κύκλων.
\end{ylh}
% ===================================================================
%  Αρχείο: 13835.tex (Fragment)
%  Αυτόματη μετατροπή μέσω doc2latex.py
% ===================================================================



Τα σημεία $A$, Κ και Λ δε βρίσκονται στην ίδια ευθεία. Το σημείο $A$ απέχει 4 από το Κ και 5 από το Λ.

\begin{alist}
\item Να αποδείξετε ότι 1< $K\varLambda $ < 9. \monades{12}

\item Να βρείτε ένα σημείο $B$ του επιπέδου διαφορετικό από το Α, που να απέχει 4 από το Κ και 5 από το Λ. \monades{13}


\begin{center}
\begin{tikzpicture}[scale=0.7]
% Points A, K, Λ not collinear
\tkzDefPoint(0,0){K}
\tkzDefPoint(6,0){L}
\tkzDefPoint(2,3.5){A}

% Draw circles (K,4) and (Λ,5)
\tkzDrawCircle[dashed, l3](K,A)  % radius $AK$=4
\tkzDrawCircle[dashed, l3](L,A)  % radius $A\varLambda $=5

% Draw triangle $AK\varLambda $
\draw[l3] (A)--(K)--(L)--(A);

% Labels for segments
\tkzLabelSegment[above left](A,K){$4$}
\tkzLabelSegment[above right](A,L){$5$}

% Points and labels
\tkzDrawPoints(A,K,L)
\tkzLabelPoint[above](A){$A$}
\tkzLabelPoint[below left](K){$ K$}
\tkzLabelPoint[below right](L){$\Lambda$}
\end{tikzpicture}
\end{center}
\end{alist}

\vspace{4mm}

% --- Άσκηση 13836 (13836.tex) ---
\Askhsh{13836}{2}
\begin{ylh}{1}
    \item 3.16. Σχετικές θέσεις δυο κύκλων.
\end{ylh}
% ===================================================================
%  Αρχείο: 13836.tex (Fragment)
%  Αυτόματη μετατροπή μέσω doc2latex.py
% ===================================================================



\begin{alist}
\item Στο παρακάτω σχήμα για τους κύκλους (Α, ρ) και (Β, R) ισχύει $\rho < R$.


\begin{center}
\begin{tikzpicture}[scale=0.7]
% Two circles (A,ρ) and (B,R) with ρ < R
\tkzDefPoint(0,0){A}
\tkzDefPoint(6,0){B}

% Circle (A,ρ) with ρ=1.5
\tkzDefPoint(1.5,0){Gamma_r}
\tkzDrawCircle(A,Gamma_r)

% Circle (B,R) with R=2.5
\tkzDefPoint(8.5,0){Delta_r}
\tkzDrawCircle(B,Delta_r)

% $A\varGamma $ = radius segment of (A,ρ) extended to left
\tkzDefPoint(-1.5,0){Gamma}

% $B\varDelta  $ = radius segment of (B,R) extended to right
\tkzDefPoint(8.5,0){Delta}

% Draw segment $AB$
\draw[l3] (Gamma)--(Delta);

% Draw labels ρ and R
\tkzLabelSegment[above](A,Gamma_r){$\rho$}
\tkzLabelSegment[above](B,Delta_r){$R$}

% Points and labels
\tkzDrawPoints(A,B,Gamma,Delta)
\tkzLabelPoint[below](A){$A$}
\tkzLabelPoint[below](B){$B$}
\tkzLabelPoint[below left](Gamma){$\varGamma $}
\tkzLabelPoint[below right](Delta){$\varDelta  $}
\end{tikzpicture}
\end{center}

Να αποδείξετε ότι $Β\varDelta   - Α\varGamma  < $AB$ < Α\varGamma  + Β\varDelta  $. \monades{10} β) Ο χάρτης ενός κρυμμένου θησαυρού έχει δύο σταθερά σημεία $A$ και Β, τα οποία απέχουν μεταξύ τους 6. Επίσης γράφει ότι ο θησαυρός είναι κρυμμένος σε ένα σημείο το οποίο απέχει 3 από το Α του χάρτη και 5 από το Β του χάρτη. Ποια είναι τα σημεία του χάρτη στα οποία μπορεί να είναι κρυμμένος ο θησαυρός;

\monades{15}
\end{alist}

\vspace{4mm}

% --- Άσκηση 13837 (13837.tex) ---
\Askhsh{13837}{2}
\begin{ylh}{1}
    \item 3.2. 1ο Κριτήριο ισότητας τριγώνων
    \item 4.6. Άθροισμα γωνιών τριγώνου
    \item 5.2. Παραλληλόγραμμα
    \item 5.9. Μια ιδιότητα του ορθογώνιου τριγώνου.
\end{ylh}
% ===================================================================
%  Αρχείο: 13837.tex (Fragment)
%  Αυτόματη μετατροπή μέσω doc2latex.py
% ===================================================================



Δίνεται ορθογώνιο τρίγωνο $AB\varGamma $ με $\hat{A}$=90\textsuperscript{ο} και $\hat{B}$=60\textsuperscript{ο} . Θεωρούμε τα σημεία $\varDelta  $ και $E$ που είναι τα μέσα των πλευρών $B\varGamma $ και $A\varGamma $ αντίστοιχα. Προεκτείνουμε την $\varDelta  E$ κατά τμήμα $EZ$=$\varDelta  E$.

\begin{alist}
\item Να αποδείξετε ότι $\varGamma \varDelta  $=$AZ$. \monades{12}

\item Να αποδείξετε ότι το τρίγωνο $ZAB$ είναι ισοσκελές. \monades{13}


\begin{center}
\begin{tikzpicture}[scale=1]
\tkzDefPoint(0,0){A}
\tkzDefPoint(2,0){B}
\tkzDefPoint(0,3.464){C} % Angle B is 60, A is 90

\tkzDefMidPoint(B,C) \tkzGetPoint{D}
\tkzDefMidPoint(A,C) \tkzGetPoint{E}

% Extend \Delta E to Z such that $EZ$ = \Delta E.
% So Z is symmetric of \Delta w.r.t E
\tkzDefPointBy[symmetry=center E](D) \tkzGetPoint{Z}

\draw[l3] (A)--(B)--(C)--cycle;
\draw[l3] (D)--(Z);
\draw[l3] (A)--(Z);
\draw[l3] (Z)--(B);

\tkzDrawPoints(A,B,C,D,E,Z)
\tkzLabelPoint[below](A){$A$}
\tkzLabelPoint[below](B){$B$}
\tkzLabelPoint[above](C){$\varGamma $}
\tkzLabelPoint[above right](D){$\varDelta  $}
\tkzLabelPoint[left](E){$E$}
\tkzLabelPoint[left](Z){$Z$}

\tkzMarkRightAngle(C,A,B)
\tkzMarkSegments[mark=s|](C,D D,B)
\tkzMarkSegments[mark=s||](C,E E,A)
\tkzMarkSegments[mark=s|||](D,E E,Z)
\end{tikzpicture}
\end{center}
\end{alist}

\vspace{4mm}

% --- Άσκηση 13838 (13838.tex) ---
\Askhsh{13838}{4}
\begin{ylh}{1}
    \item 4.2. Τέμνουσα δύο ευθειών - Ευκλείδειο αίτημα
    \item 5.6. Εφαρμογές στα τρίγωνα
    \item 5.10. Τραπέζιο
    \item 5.11. Ισοσκελές τραπέζιο.
\end{ylh}
% ===================================================================
%  Αρχείο: 13838.tex (Fragment)
%  Αυτόματη μετατροπή μέσω doc2latex.py
% ===================================================================



Δίνεται ισοσκελές τρίγωνο $AB\varGamma $ ($AB=A\varGamma $), με Κ, Μ τα μέσα των πλευρών $AB$ και $A\varGamma $ αντίστοιχα. Η ευθεία που διέρχεται από τα σημεία $K$ και $M$ τέμνει τις εξωτερικές διχοτόμους των γωνιών Β και $\varGamma$ στα σημεία $H$ και $Z$ αντίστοιχα, όπως φαίνεται στο παρακάτω σχήμα.

\begin{alist}
\item Να αποδείξετε ότι το τετράπλευρο $KM\varGamma B$ είναι ισοσκελές τραπέζιο. \monades{11}

\item Να αποδείξετε ότι το τετράπλευρο $B\varGamma $ΖΗ είναι ισοσκελές τραπέζιο. \monades{14}


\begin{center}
\begin{tikzpicture}[scale=1]
\tkzDefPoint(0,0){B}
\tkzDefPoint(4,0){C}
\tkzDefPoint(2,4){A}
\tkzDefMidPoint(A,B) \tkzGetPoint{K}
\tkzDefMidPoint(A,C) \tkzGetPoint{M}

\tkzDefLine[bisector out](A,B,C) \tkzGetPoint{extB}
\tkzDefLine[bisector out](A,C,B) \tkzGetPoint{extC}

% $KM$ intersects external bisectors at H, Z
\tkzInterLL(K,M)(B,extB) \tkzGetPoint{H_temp}
\tkzInterLL(K,M)(C,extC) \tkzGetPoint{Z_temp}

% Wait! "εξωτερικές διχοτόμους των γωνιών Β και $\varGamma$ στα σημεία $H$ και Ζ", let's make sure H is near B and Z is near C.
\tkzDefPointBy[projection=onto B--extB](K) \tkzGetPoint{H_verify} % Should just be intersecting lines.
% Let's manually draw it to avoid messy intersection of parallel or diverging lines if defined wrong
\tkzInterLL(K,M)(B,extB) \tkzGetPoint{H}
\tkzInterLL(K,M)(C,extC) \tkzGetPoint{Z}

\draw[l3] (A)--(B)--(C)--cycle;
\draw[l3] (H)--(Z);
\draw[l3] (B)--(H);
\draw[l3] (C)--(Z);

\tkzDefPointOnLine[pos=-0.5](B,A) \tkzGetPoint{Bx}
\tkzDefPointOnLine[pos=-0.5](C,A) \tkzGetPoint{Cx}
\draw[l3, ->] (B)--(Bx) node[left] {$x$};
\draw[l3, ->] (C)--(Cx) node[right] {$y$};

\tkzDrawPoints(A,B,C,K,M,H,Z)
\tkzLabelPoint[above](A){$A$}
\tkzLabelPoint[below](B){$B$}
\tkzLabelPoint[below](C){$\varGamma $}
\tkzLabelPoint[above left](K){$K$}
\tkzLabelPoint[above right](M){$M$}
\tkzLabelPoint[left](H){$H$}
\tkzLabelPoint[right](Z){$Z$}
\end{tikzpicture}
\end{center}
\end{alist}

\vspace{4mm}

% --- Άσκηση 13839 (13839.tex) ---
\Askhsh{13839}{4}
\begin{ylh}{1}
    \item 3.2. 1ο Κριτήριο ισότητας τριγώνων
    \item 3.6. Κριτήρια ισότητας ορθογώνιων τριγώνων.
\end{ylh}
% ===================================================================
%  Αρχείο: 13839.tex (Fragment)
%  Αυτόματη μετατροπή μέσω doc2latex.py
% ===================================================================



Τα ευθύγραμμα τμήματα $A\varDelta  $ και $B\varGamma $ τέμνονται στο σημείο $E$ έτσι ώστε $AE$=$\varGamma E$ και $BE$=$E\varDelta  $.

\begin{alist}
\item Να αποδείξετε ότι τα τρίγωνα Α$BE$ και $\varGamma \varDelta  E$ είναι ίσα. \monades{8}

\item Να αποδείξετε ότι οι αποστάσεις $EH$ και $E\varTheta $ του σημείου $E$ από τις πλευρές $AB$ και $\varGamma \varDelta  $, αντίστοιχα, είναι ίσες. \monades{5}

\item Αν οι προεκτάσεις των $AB$ και $\varGamma \varDelta  $ προς τα Α και $\varGamma$ αντίστοιχα τέμνονται στο Ζ, να αποδείξετε ότι το τρίγωνο Β$\varDelta   Z$ είναι ισοσκελές. \monades{12}


\begin{center}
\begin{tikzpicture}[scale=1.5]
\tkzDefPoint(0,0){Z}
\tkzDefPoint(1,2){E}
\tkzDefPoint(2,4){D} % \Delta is (2,4)
\tkzDefPoint(2,0){B} % B is (2,0)

% E is midpoint of \Delta B? $BE$ = \Delta E. Let's make B=(-1,0), \Delta=(1,0), E=(0,0).
% Or Z is intersection. "Αν οι προεκτάσεις των $AB$, Γ\Delta προς τα Α, Γ τέμνονται στο Ζ"
% Let's build it backward from Z.
\end{tikzpicture}
\end{center}
\end{alist}

\vspace{4mm}

% --- Άσκηση 13841 (13841.tex) ---
\Askhsh{13841}{4}
\begin{ylh}{1}
    \item 3.3. 2ο Κριτήριο ισότητας τριγώνων
    \item 4.2. Τέμνουσα δύο ευθειών - Ευκλείδειο αίτημα
    \item 5.2. Παραλληλόγραμμα
    \item 5.4. Ρόμβος
    \item 5.5. Τετράγωνο.
\end{ylh}
% ===================================================================
%  Αρχείο: 13841.tex (Fragment)
%  Αυτόματη μετατροπή μέσω doc2latex.py
% ===================================================================



Σε τρίγωνο $AB\varGamma $, $B\varDelta  $ η διχοτόμος της γωνίας Β και $M$ το μέσο της. Από το σημείο $\varDelta  $ φέρουμε παράλληλη προς τη $B\varGamma $, η οποία τέμνει την πλευρά $AB$ στο σημείο $E$. Αν η $EM$ τέμνει τη $B\varGamma $ στο σημείο $Z$ τότε:

\begin{alist}
\item Να αποδείξετε ότι $BE=E\varDelta  $. \monades{7}

\item Να αποδείξετε ότι $BE\parallel Z\varDelta  $. \monades{8}

\item Να αποδείξετε ότι το τετράπλευρο $\varDelta   EBZ$ είναι ρόμβος. \monades{5}

\item Ποιο θα έπρεπε να είναι το είδος του τριγώνου $AB\varGamma $ ώστε το τετράπλευρο $\varDelta   EBZ$ να είναι τετράγωνο; Δικαιολογήστε πλήρως την απάντησή σας. \monades{5}
\end{alist}

\vspace{4mm}

% --- Άσκηση 13842 (13842.tex) ---
\Askhsh{13842}{2}
\begin{ylh}{1}
    \item 4.2. Τέμνουσα δύο ευθειών - Ευκλείδειο αίτημα
    \item 4.6. Άθροισμα γωνιών τριγώνου
    \item 5.4. Ρόμβος.
\end{ylh}
% ===================================================================
%  Αρχείο: 13842.tex (Fragment)
%  Αυτόματη μετατροπή μέσω doc2latex.py
% ===================================================================




\begin{center}
\begin{tikzpicture}[scale=0.9]
% Rhombus $ABZH$: $ZB$̂H = 51° means angle at B in triangle $ZBH$
% Place B at origin, rhombus with acute angle 51° at B (between Z and H)
\tkzDefPoint(0,0){B}
\tkzDefPoint(3,0){Z}  % $BZ$ along x-axis

% H is at angle 51° from $BZ$ direction
\tkzDefPointBy[rotation=center B angle 51](Z) \tkzGetPoint{H_dir}
\tkzDefPointOnLine[pos=1](B,H_dir) \tkzGetPoint{H}

% A = Z + $BH$ - B (parallelogram rule)
\pgfmathsetmacro{\Ax}{3+3*cos(51)-0}
\pgfmathsetmacro{\Ay}{0+3*sin(51)-0}
\tkzDefPoint(\Ax,\Ay){A}

% Draw rhombus
\draw[l3] (A)--(B)--(Z)--(H)--cycle;
\draw[l3] (A)--(B);

% Diagonal
\draw[l3, dashed] (A)--(Z);
\draw[l3, dashed] (B)--(H);

% Triangle $\varGamma \varDelta  E$ connected via line from A through Γ
% Γ is on extension of a side, Δ below, $A\varGamma \varDelta  $ = 90°
\tkzDefPointOnLine[pos=1.6](A,Z) \tkzGetPoint{Gamma}
\tkzDefLine[orthogonal=through Gamma](A,Gamma) \tkzGetPoint{perp_d}
\tkzDefPointOnLine[pos=1.5](Gamma,perp_d) \tkzGetPoint{Delta}

% E: in triangle $\varGamma \varDelta  E$, Ê = 56°
\tkzDefPointBy[rotation=center Delta angle 56](Gamma) \tkzGetPoint{E_dir}
\tkzInterLL(Gamma,E_dir)(Delta,E_dir) \tkzGetPoint{E_temp}
\tkzDefPointOnLine[pos=1.8](Delta,E_dir) \tkzGetPoint{E}

% Draw triangle $\varGamma \varDelta  E$
\draw[l3] (Gamma)--(Delta)--(E)--cycle;

% Draw line A-Γ
\draw[l3] (A)--(Gamma);

% Right angle mark at Γ
\tkzMarkRightAngle[size=0.25](A,Gamma,Delta)

% Angle mark $ZB$̂H = 51°
\tkzMarkAngle[size=0.5](Z,B,H)
\tkzLabelAngle[pos=0.8](Z,B,H){\small $51\degree$}

% Points and labels
\tkzDrawPoints(A,B,Z,H,Gamma,Delta,E)
\tkzLabelPoint[above](A){$A$}
\tkzLabelPoint[below left](B){$B$}
\tkzLabelPoint[below right](Z){$ Z$}
\tkzLabelPoint[above](H){$ H$}
\tkzLabelPoint[right](Gamma){$\varGamma $}
\tkzLabelPoint[below right](Delta){$\varDelta  $}
\tkzLabelPoint[right](E){$ E$}
\end{tikzpicture}
\end{center}

Στο παραπάνω σχήμα το τετράπλευρο $ABZH$ είναι ρόμβος.

Επίσης δίνεται ότι $Ζ\hat{B}Η = 51\degree$ και ότι η $A\hat{\varGamma }\varDelta  $ είναι ορθή.

\begin{alist}
\item Να υπολογίσετε τη γωνία $A\hat{B}H$. \monades{9}

\item Να υπολογίσετε τη γωνία $A\hat{H}B$. \monades{6}

\item Αν η γωνία $\hat{E}$ του τριγώνου $\varGamma \varDelta  E$ είναι ίση με 56\textsuperscript{ο}, να υπολογίσετε τη γωνία $\hat{\varGamma }$ του τριγώνου $\varGamma \varDelta  E$. \monades{10}
\end{alist}

\vspace{4mm}

% --- Άσκηση 13843 (13843.tex) ---
\Askhsh{13843}{4}
\begin{ylh}{1}
    \item 3.2. 1ο Κριτήριο ισότητας τριγώνων
    \item 3.4. 3ο Κριτήριο ισότητας τριγώνων
    \item 3.15. Εφαπτόμενα τμήματα
    \item 4.2. Τέμνουσα δύο ευθειών - Ευκλείδειο αίτημα.
\end{ylh}
% ===================================================================
%  Αρχείο: 13843.tex (Fragment)
%  Αυτόματη μετατροπή μέσω doc2latex.py
% ===================================================================



Έστω ότι οι ευθείες x΄x και y΄y εφάπτονται στον κύκλο (Ο,R) στα άκρα μιας

διαμέτρου του $AB$. Να αποδείξετε ότι:

\begin{alist}
\item οι ευθείες x΄x και y΄y είναι παράλληλες. \monades{4}

\item οι διχοτόμοι των γωνιών ΒΑx και ΑΒy τέμνονται σε σημείο $M$. \monades{6}

\item το σημείο $M$ είναι το μέσο του ημικυκλίου $AB$. \monades{10}

\item αν η διχοτόμος της γωνίας ΒΑx τέμνει την y΄y στο σημείο $\varGamma $ και η διχοτόμος της γωνίας ΑΒy τέμνει την x`x στο σημείο $\varDelta  $, τότε $M\varGamma $ = $M\varDelta  $. \monades{5}
\end{alist}

\vspace{4mm}

% --- Άσκηση 13845 (13845.tex) ---
\Askhsh{13845}{4}
\begin{ylh}{1}
    \item 3.2. 1ο Κριτήριο ισότητας τριγώνων
    \item 3.16. Σχετικές θέσεις δυο κύκλων
    \item 4.2. Τέμνουσα δύο ευθειών - Ευκλείδειο αίτημα
    \item 5.2. Παραλληλόγραμμα.
\end{ylh}
% ===================================================================
%  Αρχείο: 13845.tex (Fragment)
%  Αυτόματη μετατροπή μέσω doc2latex.py
% ===================================================================



Οι κύκλοι (Κ,R), (Λ,ρ) εφάπτονται εξωτερικά στο σημείο $A$. Φέρουμε τυχαία ευθεία η οποία διέρχεται από το Α και δεν περνάει από τα κέντρα των κύκλων, τέμνει τους κύκλους αντίστοιχα στα σημεία $B$ και Γ. Φέρουμε τις εφαπτόμενες (ε) και (δ) στα σημεία $B$ και Γ.

Να αποδείξετε ότι:

\begin{alist}
\item $Κ\hat{B}A$ = $\Lambda\hat{\varGamma }A$. \monades{8}

\item (ε) // (δ). \monades{10}

\item Να εξετάσετε σε ποια περίπτωση το τετράπλευρο $K\varGamma \varLambda B$ θα είναι

παραλληλόγραμμο; Να αιτιολογήσετε την απάντηση σας. \monades{7}


\begin{center}
\begin{tikzpicture}[scale=1]
\tkzDefPoint(0,0){K}
\tkzDefPoint(5,0){L}
\tkzDefPoint(3,0){A}

\tkzDrawCircle(K,A)
\tkzDrawCircle(L,A)

\tkzDefPoint(1.5, 2.598){B} % B is on circle K (R=3: dx=1.5, dy=2.598 => 1.5^2+2.598^2=9)
% $BA$ determines \varGamma 
\tkzInterLC(B,A)(L,A) \tkzGetPoints{Gamma_temp}{Gamma} % Intersection with circle (L, \rho=2)

\tkzDefTangent[at=B](K) \tkzGetPoint{e_dir}
\tkzDefTangent[at=Gamma](L) \tkzGetPoint{d_dir}

\draw[l3] ($(B)-0.5*(e_dir)$)--($(B)+0.5*(e_dir)$) node[right] {$\varepsilon$};
\draw[l3] ($(Gamma)-0.5*(d_dir)$)--($(Gamma)+0.5*(d_dir)$) node[left] {$\delta$};

\draw[l3] (B)--(Gamma);
\draw[l3, dashed] (K)--(B);
\draw[l3, dashed] (L)--(Gamma);
\draw[l3, dashed] (K)--(L);

\tkzDrawPoints(K,L,A,B,Gamma)
\tkzLabelPoint[below](K){$K$}
\tkzLabelPoint[below](L){$\Lambda$}
\tkzLabelPoint[below left](A){$A$}
\tkzLabelPoint[above](B){$B$}
\tkzLabelPoint[below right](Gamma){$\varGamma $}

\tkzMarkRightAngle(K,B,e_dir)
\tkzMarkRightAngle(L,Gamma,d_dir)
\end{tikzpicture}
\end{center}
\end{alist}

\vspace{4mm}

% --- Άσκηση 13846 (13846.tex) ---
\Askhsh{13846}{4}
\begin{ylh}{1}
    \item 3.12. Τριγωνική ανισότητα
    \item 3.16. Σχετικές θέσεις δυο κύκλων.
\end{ylh}
% ===================================================================
%  Αρχείο: 13846.tex (Fragment)
%  Αυτόματη μετατροπή μέσω doc2latex.py
% ===================================================================



Δίνεται το παρακάτω σχήμα με τους κύκλους (Α, ρ) και (Β, R) με R > ρ. Επίσης $AB$ = 9.

\begin{alist}
\item Να αποδείξετε ότι $R + \rho < 9$. \monades{7}

\item Να σχεδιάσετε ένα τρίγωνο $K\varLambda M$ με $K\varLambda $ να είναι ίση με ρ και η πλευρά $\varLambda M$ να είναι ίση με R. Να περιγράψετε τον τρόπο που το σχεδιάσατε και να αποδείξετε ότι η τρίτη πλευρά του είναι μικρότερη από 9.

\monades{10}

\item Έστω το τρίγωνο $K\varLambda M$ που σχεδιάσατε στο β) ερώτημα. Πόσα σημεία του επιπέδου έχουν και τις δύο ιδιότητες Ι1 και Ι2 που περιγράφονται παρακάτω;

Ι1: «Η απόσταση των σημείων από το Κ είναι ίση με ρ».

Ι2: «Η απόσταση των σημείων από το $M$ είναι ίση με R».

Να αιτιολογήσετε την απάντησή σας.

\monades{8}


\begin{center}
\begin{tikzpicture}[scale=0.6]
% Two circles (A,ρ) and (B,R) with R > ρ, $AB$ = 9
\tkzDefPoint(0,0){A}
\tkzDefPoint(6.3,0){B}

% Circle (A,ρ) with ρ=1.2
\tkzDefPoint(1.2,0){rA}
\tkzDrawCircle(A,rA)

% Circle (B,R) with R=2.2
\tkzDefPoint(8.5,0){rB}
\tkzDrawCircle(B,rB)

% Draw segment $AB$
\draw[l3] (A)--(B);

% Labels
\tkzDrawPoints(A,B)
\tkzLabelPoint[below](A){$A$}
\tkzLabelPoint[below](B){$B$}
\tkzLabelSegment[above](A,rA){$\rho$}
\tkzLabelSegment[above](B,rB){$R$}
\end{tikzpicture}
\end{center}
\end{alist}

\vspace{4mm}

% --- Άσκηση 13848 (13848.tex) ---
\Askhsh{13848}{4}
\begin{ylh}{1}
    \item 5.2. Παραλληλόγραμμα
    \item 5.3. Ορθογώνιο
    \item 5.4. Ρόμβος
    \item 5.5. Τετράγωνο.
\end{ylh}
% ===================================================================
%  Αρχείο: 13848.tex (Fragment)
%  Αυτόματη μετατροπή μέσω doc2latex.py
% ===================================================================



Στο παρακάτω σχήμα οι κύκλοι έχουν κέντρο Κ και οι $A\varGamma $ και $B\varDelta  $ είναι διάμετροί τους.


\begin{center}
\begin{tikzpicture}[scale=0.8]
% Two concentric circles at K with diameters $A\varGamma $ and $B\varDelta  $
\tkzDefPoint(0,0){K}

% Outer circle: diameter $A\varGamma $
\tkzDefPoint(-3,0){A}
\tkzDefPoint(3,0){Gamma}
\tkzDrawCircle(K,Gamma)

% Inner circle: diameter $B\varDelta  $
\tkzDefPoint(0,2){B}
\tkzDefPoint(0,-2){Delta}
\tkzDrawCircle(K,B)

% Draw diameters
\draw[l3] (A)--(Gamma);
\draw[l3] (B)--(Delta);

% Draw quadrilateral $AB\varGamma \varDelta  $
\draw[l3] (A)--(B)--(Gamma)--(Delta)--cycle;

% Points and labels
\tkzDrawPoints(K,A,B,Gamma,Delta)
\tkzLabelPoint[below right](K){$ K$}
\tkzLabelPoint[left](A){$A$}
\tkzLabelPoint[above](B){$B$}
\tkzLabelPoint[right](Gamma){$\varGamma $}
\tkzLabelPoint[below](Delta){$\varDelta  $}
\end{tikzpicture}
\end{center}

\begin{alist}
\item Αν ισχύει $A\varGamma $\textgreater $B\varDelta  $:

\begin{rlist}
\item να σχεδιάσετε το τετράπλευρο $AB\varGamma \varDelta  $ και να αποδείξετε ότι είναι παραλληλόγραμμο.

\monades{8}

\item να διατυπώσετε μια επιπλέον υπόθεση για τις $A\varGamma $ και $B\varDelta  $, ώστε το $AB\varGamma \varDelta  $ να είναι ρόμβος. Να αιτιολογήσετε την απάντησή σας.

\monades{9}

\end{rlist}
\item Αν οι δύο κύκλοι ταυτίζονται, τότε να εξετάσετε αν ο ακόλουθος ισχυρισμός είναι αληθής:

«To $AB\varGamma \varDelta  $ είναι τετράγωνο».

Να αιτιολογήσετε την απάντησή σας.

\monades{8}
\end{alist}

\vspace{4mm}

% --- Άσκηση 13850 (13850.tex) ---
\Askhsh{13850}{4}
\begin{ylh}{1}
    \item 3.3. 2ο Κριτήριο ισότητας τριγώνων
    \item 3.4. 3ο Κριτήριο ισότητας τριγώνων
    \item 4.2. Τέμνουσα δύο ευθειών - Ευκλείδειο αίτημα
    \item 5.2. Παραλληλόγραμμα
    \item 5.3. Ορθογώνιο
    \item 5.4. Ρόμβος
    \item 5.5. Τετράγωνο.
\end{ylh}
% ===================================================================
%  Αρχείο: 13850.tex (Fragment)
%  Αυτόματη μετατροπή μέσω doc2latex.py
% ===================================================================



Δίνεται το παραλληλόγραμμο $AB\varGamma \varDelta  $ του σχήματος και η $BZ$ διχοτόμος της γωνίας $\hat{B}$. Φέρουμε $\varGamma O$ κάθετη στη $BZ$ και την προεκτείνουμε έτσι ώστε να τέμνει την $AB$ στο σημείο $E$.

\begin{alist}
\item Να αποδείξετε ότι το τρίγωνο Ε$B\varGamma $ είναι ισοσκελές. \monades{8}

\item Να αποδείξετε ότι τα τρίγωνα $OZ\varGamma $ και Ο$BE$ είναι ίσα. \monades{7}

\item Να αποδείξετε ότι το τετράπλευρο Ε$B\varGamma $Ζ είναι ρόμβος \monades{6}

\item Πόσο πρέπει να είναι το μέτρο της γωνίας $\hat{B}$ ώστε το τετράπλευρο Ε$B\varGamma $Ζ να είναι τετράγωνο; \monades{4}


\begin{center}
\begin{tikzpicture}[scale=1]
\tkzDefPoint(0,0){B}
\tkzDefPoint(4,0){C}
\tkzDefPoint(2,3){A}
\tkzDefPointBy[translation=from B to C](A) \tkzGetPoint{D}

\tkzDefLine[bisector](C,B,A) \tkzGetPoint{Z_dir}
\tkzInterLL(B,Z_dir)(A,D) \tkzGetPoint{Z}

\tkzDefPointBy[projection=onto B--Z](C) \tkzGetPoint{O}
\tkzInterLL(C,O)(A,B) \tkzGetPoint{E}

\draw[l3] (A)--(B)--(C)--(D)--cycle;
\draw[l3] (B)--(Z);
\draw[l3] (C)--(E);
\draw[l3, dashed] (E)--(Z);

\tkzDrawPoints(A,B,C,D,Z,E,O)
\tkzLabelPoint[above left](A){$A$}
\tkzLabelPoint[below left](B){$B$}
\tkzLabelPoint[below right](C){$\varGamma $}
\tkzLabelPoint[above right](D){$\varDelta  $}
\tkzLabelPoint[above](Z){$Z$}
\tkzLabelPoint[left](E){$E$}
\tkzLabelPoint[above left](O){$O$}

\tkzMarkRightAngle(C,O,Z)
\end{tikzpicture}
\end{center}
\end{alist}

\vspace{4mm}

% --- Άσκηση 13851 (13851.tex) ---
\Askhsh{13851}{4}
\begin{ylh}{1}
    \item 4.2. Τέμνουσα δύο ευθειών - Ευκλείδειο αίτημα
    \item 5.2. Παραλληλόγραμμα
    \item 5.3. Ορθογώνιο
    \item 5.9. Μια ιδιότητα του ορθογώνιου τριγώνου.
\end{ylh}
% ===================================================================
%  Αρχείο: 13851.tex (Fragment)
%  Αυτόματη μετατροπή μέσω doc2latex.py
% ===================================================================



Δίνεται ορθογώνιο $AB\varGamma \varDelta  $. Προεκτείνουμε την πλευρά $\varDelta  \varGamma $ προς το μέρος του $\varGamma$ κατά τμήμα $\varGamma E$=$\varDelta  \varGamma $.

\begin{alist}
\item Να αποδείξετε ότι το τετράπλευρο $A\varGamma EB$ είναι παραλληλόγραμμο. \monades{7}

\item Να αποδείξετε ότι το τρίγωνο $B\varDelta  $Ε είναι ισοσκελές. \monades{9}

\item Αν Δ$\hat{B}$Ε=120\textsuperscript{ο} να αποδείξετε ότι $B\varDelta  $=2ΑΔ. \monades{9}


\begin{center}
\begin{tikzpicture}[scale=1]
\tkzDefPoint(0,0){A}
\tkzDefPoint(4,0){B}
\tkzDefPoint(4,2){C}
\tkzDefPoint(0,2){D}

\tkzDefPointBy[symmetry=center C](D) \tkzGetPoint{E}

\draw[l3] (A)--(B)--(C)--(D)--cycle;
\draw[l3] (C)--(E);
\draw[l3] (A)--(C);
\draw[l3] (B)--(E);
\draw[l3] (B)--(D);

\draw[l3, dashed] (D)--(E);

\tkzDrawPoints(A,B,C,D,E)
\tkzLabelPoint[below left](A){$A$}
\tkzLabelPoint[below right](B){$B$}
\tkzLabelPoint[above right](C){$\varGamma $}
\tkzLabelPoint[above left](D){$\varDelta  $}
\tkzLabelPoint[above right](E){$E$}

\tkzMarkSegments[mark=s|](D,C C,E)
\end{tikzpicture}
\end{center}
\end{alist}

\vspace{4mm}

% --- Άσκηση 13852 (13852.tex) ---
\Askhsh{13852}{4}
\begin{ylh}{1}
    \item 3.6. Κριτήρια ισότητας ορθογώνιων τριγώνων
    \item 5.2. Παραλληλόγραμμα
    \item 5.9. Μια ιδιότητα του ορθογώνιου τριγώνου.
\end{ylh}
% ===================================================================
%  Αρχείο: 13852.tex (Fragment)
%  Αυτόματη μετατροπή μέσω doc2latex.py
% ===================================================================



Δίνεται ορθογώνιο $AB\varGamma \varDelta  $ με $AB$\textgreater $A\varDelta  $ και με κέντρο Ο. Αν $BZ$ και $\varDelta  E$ είναι οι αποστάσεις των κορυφών Β και $\varDelta$ από τη διαγώνιο $A\varGamma $, τότε:

\begin{alist}
\item Να αποδείξετε ότι τα τρίγωνα $\varDelta  EO$ και $BZO$ είναι ίσα. \monades{5}

\item Να αποδείξετε ότι το τετράπλευρο $EBZ\varDelta  $ είναι παραλληλόγραμμο. \monades{8}

\item Αν Δ$\hat{A}$Ε=60\textsuperscript{ο} και $OE$=5, να βρείτε το μήκος της πλευράς $A\varDelta  $. \monades{12}


\begin{center}
\begin{tikzpicture}[scale=1]
\tkzDefPoint(0,0){D}
\tkzDefPoint(6,0){C}
\tkzDefPoint(6,3){B}
\tkzDefPoint(0,3){A}

\tkzInterLL(A,C)(B,D) \tkzGetPoint{O}

\tkzDefPointBy[projection=onto A--C](B) \tkzGetPoint{Z}
\tkzDefPointBy[projection=onto A--C](D) \tkzGetPoint{E}

\draw[l3] (A)--(B)--(C)--(D)--cycle;
\draw[l3] (A)--(C);
\draw[l3] (B)--(Z);
\draw[l3] (D)--(E);
\draw[l3, dashed] (B)--(D);

\tkzDrawPoints(A,B,C,D,O,Z,E)
\tkzLabelPoint[above left](A){$A$}
\tkzLabelPoint[above right](B){$B$}
\tkzLabelPoint[below right](C){$\varGamma $}
\tkzLabelPoint[below left](D){$\varDelta  $}
\tkzLabelPoint[above](O){$O$}
\tkzLabelPoint[below right](Z){$Z$}
\tkzLabelPoint[above left](E){$E$}

\tkzMarkRightAngle(B,Z,C)
\tkzMarkRightAngle(D,E,A)
\end{tikzpicture}
\end{center}
\end{alist}

\vspace{4mm}

% --- Άσκηση 13853 (13853.tex) ---
\Askhsh{13853}{4}
\begin{ylh}{1}
    \item 4.2. Τέμνουσα δύο ευθειών - Ευκλείδειο αίτημα
    \item 5.2. Παραλληλόγραμμα
    \item 5.6. Εφαρμογές στα τρίγωνα
    \item 5.9. Μια ιδιότητα του ορθογώνιου τριγώνου.
\end{ylh}
% ===================================================================
%  Αρχείο: 13853.tex (Fragment)
%  Αυτόματη μετατροπή μέσω doc2latex.py
% ===================================================================




\begin{center}
\begin{tikzpicture}[scale=0.9]
% Right triangle $AB\varGamma $: Â=90°, with equilateral $AB\varDelta  $, $A\varDelta  $//$B\varGamma $
% Since $A\varDelta  $//$B\varGamma $ and $AB\varDelta  $ equilateral → angle $AB\varGamma $=60°, angle $A\varGamma B$=30°
\tkzDefPoint(0,0){A}
\tkzDefPoint(4,0){B}

% Γ: right angle at A, angle at B = 60° → Γ above on perpendicular from A
\tkzDefPoint(0,{4*tan(30)}){Gamma}

% Equilateral triangle $AB\varDelta  $ below the line $AB$
\tkzDefTriangle[equilateral](B,A) \tkzGetPoint{Delta}

% K: midpoint of $B\varGamma $ (for parallelogram $A\varDelta  BK$)
\tkzDefMidPoint(B,Gamma) \tkzGetPoint{K}

% Draw right triangle
\draw[l3] (A)--(B)--(Gamma)--cycle;

% Draw equilateral triangle
\draw[l3] (A)--(Delta)--(B);

% Draw $A\varDelta  $ and $B\varGamma $ parallel indication
\draw[l3, dashed] (Delta)--(K);

% Right angle at A
\tkzMarkRightAngle[size=0.25](B,A,Gamma)

% Equal side marks for equilateral
\tkzMarkSegment[mark=||](A,B)
\tkzMarkSegment[mark=||](A,Delta)
\tkzMarkSegment[mark=||](B,Delta)

% Points and labels
\tkzDrawPoints(A,B,Gamma,Delta,K)
\tkzLabelPoint[left](A){$A$}
\tkzLabelPoint[right](B){$B$}
\tkzLabelPoint[above left](Gamma){$\varGamma $}
\tkzLabelPoint[below](Delta){$\varDelta  $}
\tkzLabelPoint[right](K){$ K$}
\end{tikzpicture}
\end{center}

Στο παραπάνω σχήμα το τρίγωνο $AB\varGamma $ είναι ορθογώνιο με $\hat{A} = 90\degree$. Επίσης οι $A\varDelta  $ και $B\varGamma $ είναι παράλληλες και το τρίγωνο Α$B\varDelta  $ είναι ισόπλευρο.

\begin{alist}
\item Να υπολογίσετε τις γωνίες $\hat{B}$ και $\hat{\varGamma }$ του τριγώνου $AB\varGamma $.

\monades{8}

\item Αν η περίμετρος του Α$B\varDelta  $ είναι 12 να βρείτε το μήκος της υποτείνουσας του $AB\varGamma $.

\monades{7}

\item Αν το σημείο $K$ είναι σημείο της υποτείνουσας τέτοιο ώστε το $A\varDelta  BK$ να είναι παραλληλόγραμμο, τότε να βρείτε τη θέση του σημείου Κ. Τι είδους παραλληλόγραμμο είναι το $A\varDelta  BK$; Να αιτιολογήσετε τις απαντήσεις σας.

\monades{10}
\end{alist}

\vspace{4mm}

% --- Άσκηση 13854 (13854.tex) ---
\Askhsh{13854}{4}
\begin{ylh}{1}
    \item 3.3. 2ο Κριτήριο ισότητας τριγώνων
    \item 3.6. Κριτήρια ισότητας ορθογώνιων τριγώνων.
\end{ylh}
% ===================================================================
%  Αρχείο: 13854.tex (Fragment)
%  Αυτόματη μετατροπή μέσω doc2latex.py
% ===================================================================



Δίνεται ισοσκελές τρίγωνο $AB\varGamma $ ($AB=A\varGamma $). Θεωρούμε τις διχοτόμους $B\varDelta  $ και $\varGamma E$ των γωνιών Β και Γ, αντίστοιχα.

\begin{alist}
\item Να αποδείξετε ότι $B\varDelta  $=$\varGamma E$. \monades{9}

\item Από τα σημεία $E$ και $\varDelta$ φέρνουμε κάθετες $E\varLambda $ και $\varDelta  K$ στις πλευρές $A\varGamma $ και $B\varGamma $ αντίστοιχα. Να αποδείξετε ότι: $\varDelta  K$=$E\varLambda $. \monades{9}

\item Να εντοπίσετε και να σχεδιάσετε σημείο $Z$ της πλευράς $B\varGamma $ που η απόστασή του από το σημείο $E$ να ισούται με την απόσταση των σημείων $\varDelta$ και Κ αιτιολογώντας πλήρως την απάντησή σας. \monades{7}


\begin{center}
\begin{tikzpicture}[scale=0.9]
% Isosceles triangle $AB\varGamma $ with $AB$ = $A\varGamma $
\tkzDefPoint(3,5){A}
\tkzDefPoint(0,0){B}
\tkzDefPoint(6,0){Gamma}

% Bisector $B\varDelta  $ of angle B: Δ on $A\varGamma $
\tkzDefLine[bisector](A,B,Gamma) \tkzGetPoint{BD_dir}
\tkzInterLL(B,BD_dir)(A,Gamma) \tkzGetPoint{Delta}

% Bisector $\varGamma E$ of angle Γ: Ε on $AB$
\tkzDefLine[bisector](B,Gamma,A) \tkzGetPoint{GE_dir}
\tkzInterLL(Gamma,GE_dir)(A,B) \tkzGetPoint{E}

% $\varDelta  K$ perpendicular from $\varDelta$ to $B\varGamma $
\tkzDefPointBy[projection=onto B--Gamma](Delta) \tkzGetPoint{K}

% $E\varLambda $ perpendicular from $E$ to $A\varGamma $
\tkzDefPointBy[projection=onto A--Gamma](E) \tkzGetPoint{L}

% Draw triangle
\draw[l3] (A)--(B)--(Gamma)--cycle;

% Draw bisectors
\draw[l3] (B)--(Delta);
\draw[l3] (Gamma)--(E);

% Draw perpendiculars
\draw[l3, dashed] (Delta)--(K);
\draw[l3, dashed] (E)--(L);

% Right angle marks
\tkzMarkRightAngle[size=0.2](Delta,K,B)
\tkzMarkRightAngle[size=0.2](E,L,A)

% Equal side marks
\tkzMarkSegment[mark=||](A,B)
\tkzMarkSegment[mark=||](A,Gamma)

% Points and labels
\tkzDrawPoints(A,B,Gamma,Delta,E,K,L)
\tkzLabelPoint[above](A){$A$}
\tkzLabelPoint[below left](B){$B$}
\tkzLabelPoint[below right](Gamma){$\varGamma $}
\tkzLabelPoint[right](Delta){$\varDelta  $}
\tkzLabelPoint[left](E){$ E$}
\tkzLabelPoint[below](K){$ K$}
\tkzLabelPoint[right](L){$\Lambda$}
\end{tikzpicture}
\end{center}
\end{alist}

\vspace{4mm}

% --- Άσκηση 13855 (13855.tex) ---
\Askhsh{13855}{4}
\begin{ylh}{1}
    \item 4.2. Τέμνουσα δύο ευθειών - Ευκλείδειο αίτημα
    \item 5.2. Παραλληλόγραμμα
    \item 5.3. Ορθογώνιο
    \item 5.6. Εφαρμογές στα τρίγωνα
    \item 5.9. Μια ιδιότητα του ορθογώνιου τριγώνου.
\end{ylh}
% ===================================================================
%  Αρχείο: 13855.tex (Fragment)
%  Αυτόματη μετατροπή μέσω doc2latex.py
% ===================================================================



Δίνεται ορθογώνιο τρίγωνο $AB\varGamma $ με $\hat{A}$=90\textsuperscript{ο}. Από το μέσο $\varDelta  $ της πλευράς $B\varGamma $ φέρουμε παράλληλη προς την πλευρά $A\varGamma $ που τέμνει την πλευρά $AB$ στο σημείο $E$. Αν επιπλέον γνωρίζουμε ότι Ε$\hat{\varDelta  }$Γ=120\textsuperscript{ο}, τότε:

\begin{alist}
\item Να υπολογίσετε το μέτρο της γωνία Δ$\hat{\varGamma }$Α. \monades{5}

\item Να αποδείξετε ότι το τρίγωνο $A\varDelta  \varGamma $ είναι ισόπλευρο. \monades{8}

\item Προεκτείνουμε την πλευρά $A\varGamma $ προς το $\varGamma$ κατά τμήμα $\varGamma Z$=$A\varGamma $ και την πλευρά $B\varGamma $ προς το $\varGamma$ κατά τμήμα $\varGamma H$=$\frac{Β\varGamma }{2}$. Να αποδείξετε ότι Α$\hat{H}$Ζ=90\textsuperscript{ο}. \monades{12}


\begin{center}
\begin{tikzpicture}[scale=0.8]
\tkzDefPoint(0,0){A}
\tkzDefPoint(0,-4.5){B}
% Angle $EDC$ = 120 means Angle C = 60, since $DE$ || $AC$.
% Wait, $DE$ || $AC$. If angle C is 60, then angle B is 30.
\tkzDefPoint(2.598,0){C} % A=90, B=30, C=60 => $AB$=4.5, $AC$=2.598

\tkzDefMidPoint(B,C) \tkzGetPoint{D}
\tkzDefLine[parallel=through D](A,C) \tkzGetPoint{E_dir}
\tkzInterLL(D,E_dir)(A,B) \tkzGetPoint{E}

\tkzDefPointBy[symmetry=center C](A) \tkzGetPoint{Z}
\tkzDefPointOnLine[pos=-0.5](C,B) \tkzGetPoint{H}

\draw[l3] (A)--(B)--(C)--cycle;
\draw[l3] (D)--(E);
\draw[l3] (A)--(D);
\draw[l3, dashed] (C)--(Z);
\draw[l3, dashed] (C)--(H);
\draw[l3, dashed] (A)--(H);
\draw[l3, dashed] (H)--(Z);

\tkzDrawPoints(A,B,C,D,E,Z,H)
\tkzLabelPoint[left](A){$A$}
\tkzLabelPoint[below left](B){$B$}
\tkzLabelPoint[above left](C){$\varGamma $}
\tkzLabelPoint[below right](D){$\varDelta  $}
\tkzLabelPoint[left](E){$E$}
\tkzLabelPoint[right](Z){$Z$}
\tkzLabelPoint[above right](H){$H$}

\tkzMarkSegments[mark=s|](C,D D,B H,C)
\tkzMarkSegments[mark=s||](A,C C,Z)
\tkzMarkAngle[size=0.4,mark=none](C,D,E)
\tkzLabelAngle[pos=0.6](C,D,E){$120^\circ$}
\tkzMarkRightAngle(C,A,B)
\end{tikzpicture}
\end{center}
\end{alist}

\vspace{4mm}

% --- Άσκηση 13856 (13856.tex) ---
\Askhsh{13856}{4}
\begin{ylh}{1}
    \item 3.2. 1ο Κριτήριο ισότητας τριγώνων
    \item 5.2. Παραλληλόγραμμα
    \item 5.6. Εφαρμογές στα τρίγωνα.
\end{ylh}
% ===================================================================
%  Αρχείο: 13856.tex (Fragment)
%  Αυτόματη μετατροπή μέσω doc2latex.py
% ===================================================================



Σε τρίγωνο $\varDelta   EZ$, φέρουμε τη διάμεσο $\varDelta   M$ και στην προέκτασή της προς το μέρος του $M$ παίρνουμε σημείο $\varTheta $ έτσι ώστε $\varDelta   M$=$M\varTheta $. Προεκτείνουμε την πλευρά $EZ$ προς το $E$ κατά τμήμα $EA$=$EZ$ και προς το $Z$ κατά τμήμα $Z\varGamma $=$EZ$.

\begin{alist}
\item Να αποδείξετε ότι τα τρίγωνα $\varDelta   AM$ και $\varTheta \varGamma  M$ είναι ίσα. \monades{8}

\item Να αποδείξετε ότι το τετράπλευρο $\varTheta  A\varDelta  \varGamma $ είναι παραλληλόγραμμο. \monades{8}

\item Στο σχήμα της άσκησης που κατασκεύασε στο τετράδιό του ο Γιάννης είναι $A\varDelta  $=12. Πόσο θα είναι το μήκος της διαμέσου $EH$ του τριγώνου $\varDelta   EZ$ στο σχήμα του Γιάννη; \monades{9}
\end{alist}

\vspace{4mm}

% --- Άσκηση 13857 (13857.tex) ---
\Askhsh{13857}{4}
\begin{ylh}{1}
    \item 5.2. Παραλληλόγραμμα
    \item 5.4. Ρόμβος.
\end{ylh}
% ===================================================================
%  Αρχείο: 13857.tex (Fragment)
%  Αυτόματη μετατροπή μέσω doc2latex.py
% ===================================================================



\begin{alist}
\item Στο σχήμα η $B\varDelta  $ είναι μεσοκάθετος του τμήματος $A\varGamma $ και διάμετρος του κύκλου με κέντρο Μ. Να αποδείξετε ότι το $AB\varGamma \varDelta  $ είναι ρόμβος. \monades{8}


\begin{center}
\begin{tikzpicture}[scale=0.8]
% $B\varDelta  $ perpendicular bisector of $A\varGamma $, $B\varDelta  $ diameter of circle(M)
\tkzDefPoint(0,0){M}
\tkzDefPoint(0,2.5){B}
\tkzDefPoint(0,-2.5){Delta}
\tkzDefPoint(-1.8,0){A}
\tkzDefPoint(1.8,0){Gamma}

% Circle with center M passing through B and Δ
\tkzDrawCircle(M,B)

% Draw quadrilateral $AB\varGamma \varDelta  $
\draw[l3] (A)--(B)--(Gamma)--(Delta)--cycle;

% Draw diagonals
\draw[l3] (A)--(Gamma);
\draw[l3] (B)--(Delta);

% Right angle mark
\tkzMarkRightAngle[size=0.25](A,M,B)

% Points and labels
\tkzDrawPoints(A,B,Gamma,Delta,M)
\tkzLabelPoint[left](A){$A$}
\tkzLabelPoint[above](B){$B$}
\tkzLabelPoint[right](Gamma){$\varGamma $}
\tkzLabelPoint[below](Delta){$\varDelta  $}
\tkzLabelPoint[below right](M){$\mathrm{M}$}
\end{tikzpicture}
\end{center}

\item Χαρακτηρίστε κάθε μια από τις παρακάτω προτάσεις ως αληθή ή ψευδή.

Πρόταση 1: «Αν η διαγώνιος ενός τυχαίου τετραπλεύρου είναι μεσοκάθετος της άλλης διαγωνίου και διάμετρος κύκλου με κέντρο το σημείο τομής των διαγωνίων, τότε το τετράπλευρο είναι ρόμβος».

Πρόταση 2: «Αν η διαγώνιος ενός τυχαίου τετραπλεύρου είναι κάθετη στην άλλη διαγώνιο και διάμετρος κύκλου με κέντρο το σημείο τομής των διαγωνίων, τότε το τετράπλευρο είναι ρόμβος».

Να αιτιολογήσετε την απάντησή σας σε κάθε περίπτωση.

\monades{10}

\item Στο παρακάτω σχήμα τα ευθύγραμμα τμήματα $PT$ και $\varPi\varSigma$ τέμνονται κάθετα στο $N$ και $\varPi N$ = $N\varSigma$. Επίσης η $PT$ είναι διάμετρος του κύκλου με κέντρο το Ν.

Να αποδείξετε ότι $\varPi P$ = $P\varSigma$ = $\varSigma T$ = $T\varPi$. \monades{7}


\begin{center}
\begin{tikzpicture}[scale=0.8]
% $PT$ and $\varPi\varSigma$ perpendicular at N, $\varPiN$=$N\varSigma$, $PT$ diameter of circle(N)
\tkzDefPoint(0,0){N}
\tkzDefPoint(0,2){R}
\tkzDefPoint(0,-2){T}
\tkzDefPoint(-1.5,0){Pi}
\tkzDefPoint(1.5,0){S}

% Circle with center N through R
\tkzDrawCircle(N,R)

% Draw $\varPiP\varSigmaT$ quadrilateral
\draw[l3] (Pi)--(R)--(S)--(T)--cycle;

% Draw diagonals
\draw[l3] (R)--(T);
\draw[l3] (Pi)--(S);

% Right angle mark
\tkzMarkRightAngle[size=0.2](Pi,N,R)

% Equal marks $\varPiN$ = $N\varSigma$
\tkzMarkSegment[mark=||](Pi,N)
\tkzMarkSegment[mark=||](N,S)

% Points and labels
\tkzDrawPoints(N,R,T,Pi,S)
\tkzLabelPoint[above](R){$\mathrm{P}$}
\tkzLabelPoint[below](T){$ T$}
\tkzLabelPoint[left](Pi){$\Pi$}
\tkzLabelPoint[right](S){$\Sigma$}
\tkzLabelPoint[below right](N){$\mathrm{N}$}
\end{tikzpicture}
\end{center}
\end{alist}

\vspace{4mm}

% --- Άσκηση 14566 (14566.tex) ---
\Askhsh{14566}{4}
\begin{ylh}{1}
    \item 3.2. 1ο Κριτήριο ισότητας τριγώνων
    \item 5.2. Παραλληλόγραμμα
    \item 5.6. Εφαρμογές στα τρίγωνα.
\end{ylh}
% ===================================================================
%  Αρχείο: 14566.tex (Fragment)
%  Αυτόματη μετατροπή μέσω doc2latex.py
% ===================================================================



Δίνεται τετράπλευρο $AB\varGamma \varDelta  $ με $\hat{\varDelta  } = \hat{Β\ } = 90^{0}.\ \ $ Αν τα σημεία $E$, Ζ, Μ είναι τα μέσα των $AB$, $\varGamma \varDelta  $, και $A\varGamma $ αντιστοίχως και το $MK$ είναι κάθετο στην $B\varDelta  $, να αποδείξετε ότι :

\begin{alist}
\item Το τρίγωνο $BM\varDelta  $ είναι ισοσκελές (μον. 5) και το Κ είναι το μέσο του $B\varDelta  $ (μον. 2)

\monades{7}

\begin{alist}
\item
  $ΕΚ = \frac{Α\varDelta  }{2}$. \monades{6}
\item
  $MZ$=$EK$. \monades{6}

\item Το τετράπλευρο $KEMZ$ είναι παραλληλόγραμμο. \monades{6}


\begin{center}
\begin{tikzpicture}[scale=1]
\tkzDefPoint(0,0){B}
\tkzDefPoint(0,4){A}
\tkzDefPoint(6,0){C}
\tkzDefPoint(2,6){D} % make angle D = 90

% We need angle D = 90. Let D = (x,y). Vector $AD$=(x, y-4), $CD$=(x-6, y).
% $AD$ dot $CD$ = x(x-6) + y(y-4) = 0 => x^2 - 6x + y^2 - 4y = 0.
% Let D=(2,4). x^2-6x=-8. y^2-4y=0. Sum is -8. No.
% Let D=(3, 5). x^2-6x=-9, y^2-4y=5. Sum = -4.
% Let D=(2, 4.8). Too complex. Let's just define orthogonal from A to $CD$?
% No, angle D = 90 and angle B = 90. This means $ABCD$ is inscribed in a circle with diameter $AC$.
% Center of circle is M, midpoint of $AC$. So M is (3,2).
% Circle radius is sqrt(3^2+2^2) = sqrt(13) ~ 3.6
% A = (0,4), C = (6,0). B = (0,0) is on circle! Since angle B=90.
% We can just pick D on the circle. Let M be (3,2). Radius is sqrt(13).
\tkzDefMidPoint(A,C) \tkzGetPoint{M}
\tkzDefPoint(4.66, 5.58){D} % Approximate point on circle. Or just:
\tkzDefPointBy[rotation=center M angle -70](A) \tkzGetPoint{D} % rotate A around M to get D on circle!

\tkzDefMidPoint(A,B) \tkzGetPoint{E}
\tkzDefMidPoint(C,D) \tkzGetPoint{Z}

\tkzDefPointBy[projection=onto B--D](M) \tkzGetPoint{K}

\draw[l3] (A)--(B)--(C)--(D)--cycle;
\draw[l3] (A)--(C);
\draw[l3] (B)--(D);
\draw[l3] (M)--(K);
\draw[l3, dashed] (K)--(E)--(M)--(Z)--cycle;

\tkzDrawPoints(A,B,C,D,M,E,Z,K)
\tkzLabelPoint[above left](A){$A$}
\tkzLabelPoint[below left](B){$B$}
\tkzLabelPoint[below right](C){$\varGamma $}
\tkzLabelPoint[above right](D){$\varDelta  $}
\tkzLabelPoint[above right](M){$M$}
\tkzLabelPoint[left](E){$E$}
\tkzLabelPoint[right](Z){$Z$}
\tkzLabelPoint[below](K){$K$}

\tkzMarkRightAngle(C,B,A)
\tkzMarkRightAngle(A,D,C)
\tkzMarkRightAngle(M,K,B)
\end{tikzpicture}
\end{center}
\end{alist}
\end{alist}

\vspace{4mm}

% --- Άσκηση 14876 (14876.tex) ---
\Askhsh{14876}{2}
\begin{ylh}{1}
    \item 4.2. Τέμνουσα δύο ευθειών - Ευκλείδειο αίτημα
    \item 5.2. Παραλληλόγραμμα
    \item 5.9. Μια ιδιότητα του ορθογώνιου τριγώνου.
\end{ylh}
% ===================================================================
%  Αρχείο: 14876.tex (Fragment)
%  Αυτόματη μετατροπή μέσω doc2latex.py
% ===================================================================



Δίνεται παραλληλόγραμμο $AB\varGamma \varDelta  $ με γωνία $\hat{A}$ = 120\textsuperscript{ο} και $AB$=2ΑΔ. Φέρουμε τη διχοτόμο της γωνίας $\varDelta$ του παραλληλογράμμου, η οποία τέμνει την $AB$ στο Ε, και στη συνέχεια το κάθετο τμήμα $AZ$ στη $\varDelta  E$. Να αποδείξετε ότι:

\begin{alist}
\item γωνία Α$\hat{\varDelta  }$Ε = 30\textsuperscript{ο} \monades{10}

\item $AZ$ = $\frac{AB}{4}$ \monades{15}


\begin{center}
\begin{tikzpicture}[scale=0.8]
\tkzDefPoint(0,0){D}
\tkzDefPoint(8,0){C}
\tkzDefPoint(2, 3.464){A} % A=120 means $ADC$=60, $AD$=4 => dx=4cos60=2, dy=4sin60=3.464
\tkzDefPointBy[translation=from D to C](A) \tkzGetPoint{B}
% Wait, A=120 means $DAB$ = 120. If $AB$ is horizontal, then $AD$ is at 120 degrees.
% Let's keep $AB$=8, $AD$=4.
\tkzDefPoint(0,0){A}
\tkzDefPoint(8,0){B}
\tkzDefPointBy[rotation=center A angle 120](B) \tkzGetPoint{D_temp}
\tkzDefPointOnLine[pos=0.5](A,D_temp) \tkzGetPoint{D}
\tkzDefPointBy[translation=from A to D](B) \tkzGetPoint{C}

\tkzDefLine[bisector](A,D,C) \tkzGetPoint{bis_dir}
\tkzInterLL(D,bis_dir)(A,B) \tkzGetPoint{E}

\tkzDefPointBy[projection=onto D--E](A) \tkzGetPoint{Z}

\draw[l3] (A)--(B)--(C)--(D)--cycle;
\draw[l3] (D)--(E);
\draw[l3] (A)--(Z);

\tkzDrawPoints(A,B,C,D,E,Z)
\tkzLabelPoint[below left](A){$A$}
\tkzLabelPoint[below right](B){$B$}
\tkzLabelPoint[above right](C){$\varGamma $}
\tkzLabelPoint[above left](D){$\varDelta  $}
\tkzLabelPoint[below](E){$E$}
\tkzLabelPoint[right](Z){$Z$}

\tkzMarkRightAngle(A,Z,D)
\end{tikzpicture}
\end{center}
\end{alist}

\vspace{4mm}

% --- Άσκηση 14877 (14877.tex) ---
\Askhsh{14877}{2}
\begin{ylh}{1}
    \item 3.12. Τριγωνική ανισότητα
    \item 5.6. Εφαρμογές στα τρίγωνα
    \item 5.10. Τραπέζιο.
\end{ylh}
% ===================================================================
%  Αρχείο: 14877.tex (Fragment)
%  Αυτόματη μετατροπή μέσω doc2latex.py
% ===================================================================



Στο τρίγωνο $AB\varGamma $ του παρακάτω σχήματος τα σημεία $\varDelta  $ και $E$ είναι τα μέσα των πλευρών $AB$ και $A\varGamma $ αντίστοιχα, $AE$=8 , $E\varDelta  $=9 και $\varDelta   B$=10 .

\begin{alist}
\item Να αποδείξετε ότι το τετράπλευρο $\varDelta  E\varGamma B$ είναι τραπέζιο. \monades{8}

\item Να υπολογίσετε το μήκος της πλευράς $B\varGamma $. \monades{8}

\item Να συγκρίνετε τις περιμέτρους του τριγώνου $AB\varGamma $ και του τετραπλεύρου $\varDelta  E\varGamma B$. \monades{9}


\begin{center}
\begin{tikzpicture}[scale=0.3]
% Triangle with sides 20, 18, 16.
% Let A be at origin. $AB$ = 20, $BC$ = 18, $AC$ = 16.
% But conventionally $BC$ is base. Let C=(0,0), B=(18,0).
% Or just use standard construction:
\tkzDefPoint(0,0){B}
\tkzDefPoint(18,0){C} % B\varGamma  = 18
% $AB$ = 20, $AC$ = 16
\tkzInterCC[R](B,20)(C,16) \tkzGetPoints{A_temp}{A}

\tkzDefMidPoint(A,B) \tkzGetPoint{D}
\tkzDefMidPoint(A,C) \tkzGetPoint{E}

\draw[l3] (A)--(B)--(C)--cycle;
\draw[l3] (D)--(E);

\tkzDrawPoints(A,B,C,D,E)
\tkzLabelPoint[above](A){$A$}
\tkzLabelPoint[below left](B){$B$}
\tkzLabelPoint[below right](C){$\varGamma $}
\tkzLabelPoint[above left](D){$\varDelta  $}
\tkzLabelPoint[above right](E){$E$}

\tkzLabelSegment[left](A,D){$10$}
\tkzLabelSegment[left](D,B){$10$}
\tkzLabelSegment[right](A,E){$8$}
\tkzLabelSegment[right](E,C){$8$}
\tkzLabelSegment[above](D,E){$9$}
\end{tikzpicture}
\end{center}
\end{alist}

\vspace{4mm}

% --- Άσκηση 14879 (14879.tex) ---
\Askhsh{14879}{4}
\begin{ylh}{1}
    \item 5.3. Ορθογώνιο
    \item 5.6. Εφαρμογές στα τρίγωνα
    \item 5.9. Μια ιδιότητα του ορθογώνιου τριγώνου.
\end{ylh}
% ===================================================================
%  Αρχείο: 14879.tex (Fragment)
%  Αυτόματη μετατροπή μέσω doc2latex.py
% ===================================================================



Δίνεται ορθογώνιο παραλληλόγραμμο $AB\varGamma \varDelta  $. Από την κορυφή $A$ φέρουμε $AE$ κάθετη στη $B\varDelta  $. Έστω Κ, Λ τα μέσα των $AB$ και $A\varDelta  $ αντίστοιχα.

\begin{alist}
\item Να αποδείξετε ότι :

\begin{alist}
\item
  Κ$\hat{Ε}$Λ = 90\textsuperscript{ο} .
\item
  $K\varLambda $ = $\frac{A\varGamma }{2}$ .

\monades{16}

\item Αν Β$\hat{A}$Γ = 30\textsuperscript{ο} , να αποδείξετε ότι $K\varLambda $ = $B\varGamma $. \monades{9}


\begin{center}
\begin{tikzpicture}[scale=1]
\tkzDefPoint(0,0){D}
\tkzDefPoint(6,0){C}
\tkzDefPoint(6,4){B}
\tkzDefPoint(0,4){A}

\tkzDefPointBy[projection=onto B--D](A) \tkzGetPoint{E}
\tkzDefMidPoint(A,B) \tkzGetPoint{K}
\tkzDefMidPoint(A,D) \tkzGetPoint{L}

\draw[l3] (A)--(B)--(C)--(D)--cycle;
\draw[l3] (B)--(D);
\draw[l3] (A)--(E);
\draw[l3] (K)--(L);
\draw[l3] (E)--(K);
\draw[l3] (E)--(L);

\tkzDrawPoints(A,B,C,D,E,K,L)
\tkzLabelPoint[above left](A){$A$}
\tkzLabelPoint[above right](B){$B$}
\tkzLabelPoint[below right](C){$\varGamma $}
\tkzLabelPoint[below left](D){$\varDelta  $}
\tkzLabelPoint[below](E){$E$}
\tkzLabelPoint[above](K){$K$}
\tkzLabelPoint[left](L){$\Lambda$}

\tkzMarkRightAngle(A,E,B)
\tkzMarkSegments[mark=s|](A,K K,B)
\tkzMarkSegments[mark=s||](A,L L,D)
\end{tikzpicture}
\end{center}
\end{alist}
\end{alist}

\vspace{4mm}

% --- Άσκηση 14880 (14880.tex) ---
\Askhsh{14880}{4}
\begin{ylh}{1}
    \item 3.1. Είδη και στοιχεία τριγώνων
    \item 3.2. 1ο Κριτήριο ισότητας τριγώνων
    \item 3.3. 2ο Κριτήριο ισότητας τριγώνων
    \item 3.4. 3ο Κριτήριο ισότητας τριγώνων.
\end{ylh}
% ===================================================================
%  Αρχείο: 14880.tex (Fragment)
%  Αυτόματη μετατροπή μέσω doc2latex.py
% ===================================================================



Δίνεται τετράπλευρο $AB\varGamma \varDelta  $ με $AB$ = $A\varDelta  $ και $\varGamma  B$ = $\varGamma \varDelta  $. Αν $E$ είναι το σημείο τομής των προεκτάσεων των $BA$ και $\varGamma \varDelta  $ και $Z$ το σημείο τομής των προεκτάσεων των $\varDelta   A$ και $\varGamma  B$ να αποδείξετε ότι:

\begin{alist}
\item Η $\varGamma  A$ είναι διχοτόμος της γωνίας $B\varGamma $Δ. \monades{7}

\item $\varGamma  Z$ = $\varGamma  E$ \monades{9}

\item $EZ$ // $B\varDelta  $ \monades{9}


\begin{center}
\begin{tikzpicture}[scale=0.8]
\tkzDefPoint(0,0){C}
\tkzDefPoint(2,3){D}
\tkzDefPointBy[symmetry=center C](D) \tkzGetPoint{B_sym}
% Actually $CB$ = $CD$, $AB$ = $AD$. Let C be at origin. D at (r, \theta), B at (r, -\theta).
\tkzDefPoint(5,0){A} % A is on x axis
\tkzDefPoint(2,2){D}
\tkzDefPoint(2,-2){B}
% Now $AB$ = $AD$ and $CB$ = $CD$.

% E is intersection of $BA$ and $CD$.
\tkzInterLL(B,A)(C,D) \tkzGetPoint{E}
% Z is intersection of $DA$ and $CB$.
\tkzInterLL(D,A)(C,B) \tkzGetPoint{Z}

\draw[l3] (A)--(B)--(C)--(D)--cycle;
\draw[l3] (C)--(A);
\draw[l3] (B)--(E);
\draw[l3] (D)--(E);
\draw[l3] (D)--(Z);
\draw[l3] (B)--(Z);
\draw[l3, dashed] (E)--(Z);

\tkzDrawPoints(A,B,C,D,E,Z)
\tkzLabelPoint[right](A){$A$}
\tkzLabelPoint[below](B){$B$}
\tkzLabelPoint[left](C){$\varGamma $}
\tkzLabelPoint[above](D){$\varDelta  $}
\tkzLabelPoint[above right](E){$E$}
\tkzLabelPoint[below right](Z){$Z$}

\tkzMarkSegments[mark=s|](C,B C,D)
\tkzMarkSegments[mark=s||](A,B A,D)
\end{tikzpicture}
\end{center}
\end{alist}

\vspace{4mm}

% --- Άσκηση 14881 (14881.tex) ---
\Askhsh{14881}{4}
\begin{ylh}{1}
    \item 3.2. 1ο Κριτήριο ισότητας τριγώνων
    \item 5.6. Εφαρμογές στα τρίγωνα
    \item 5.9. Μια ιδιότητα του ορθογώνιου τριγώνου.
\end{ylh}
% ===================================================================
%  Αρχείο: 14881.tex (Fragment)
%  Αυτόματη μετατροπή μέσω doc2latex.py
% ===================================================================



Δίνεται ορθογώνιο τρίγωνο $AB\varGamma $ με τη γωνία $A$ ορθή και $AM$ η διάμεσός του. Από το $M$ φέρουμε $MK$ κάθετη στην $AB$ και $M\varLambda $ κάθετη στην $A\varGamma $. Αν Ν, Ρ είναι τα μέσα των $BM$ και $\varGamma  M$ αντίστοιχα, να αποδείξετε ότι:

\begin{alist}
\item N$\hat{K}$M = N$\hat{M}$K \monades{7}

\item Η $MK$ είναι διχοτόμος της γωνίας $NMA$. \monades{9}

\item $AM$ = $KN$ + $\varLambda  P$. \monades{9}
\end{alist}

\vspace{4mm}

% --- Άσκηση 14882 (14882.tex) ---
\Askhsh{14882}{4}
\begin{ylh}{1}
    \item 3.2. 1ο Κριτήριο ισότητας τριγώνων
    \item 4.2. Τέμνουσα δύο ευθειών - Ευκλείδειο αίτημα
    \item 4.6. Άθροισμα γωνιών τριγώνου
    \item 5.6. Εφαρμογές στα τρίγωνα
    \item 5.10. Τραπέζιο
    \item 5.11. Ισοσκελές τραπέζιο.
\end{ylh}
% ===================================================================
%  Αρχείο: 14882.tex (Fragment)
%  Αυτόματη μετατροπή μέσω doc2latex.py
% ===================================================================



Δίνεται ισόπλευρο τρίγωνο $AB\varGamma $. Στην προέκταση της $B\varGamma $ (προς το Γ) θεωρούμε τμήμα $\varGamma \varDelta  $ = $B\varGamma $. Αν Μ, Κ και Λ είναι τα μέσα των πλευρών $B\varGamma $, $AB$ και $A\varDelta  $ αντίστοιχα τότε:

\begin{alist}
\item Να υπολογίσετε τις γωνίες του τριγώνου $BA\varDelta  $. \monades{7}

\item Να αποδείξετε ότι:

\begin{alist}
\item
  Το τετράπλευρο $K\varLambda \varGamma  M$ είναι ισοσκελές τραπέζιο με τη μεγάλη βάση διπλάσια από τη μικρή. \monades{8}
\item
  Το τρίγωνο $KM\varLambda $ είναι ορθογώνιο. \monades{10}


\begin{center}
\begin{tikzpicture}[scale=0.8]
\tkzDefPoint(0,0){B}
\tkzDefPoint(4,0){C}
\tkzDefTriangle[equilateral](B,C) \tkzGetPoint{A}
\tkzDefPointBy[symmetry=center C](B) \tkzGetPoint{D} % C is midpoint of $BD$, so $CD$=$BC$ and B-C-D collinear! Wait. "Στην προέκταση της $B\varGamma $ (προς το Γ) θεωρούμε τμήμα $\varGamma \varDelta  $ = $B\varGamma $". Yes.
\tkzDefPointOnLine[pos=2](B,C) \tkzGetPoint{D_alt} % Should be same as symmetry

\tkzDefMidPoint(B,C) \tkzGetPoint{M}
\tkzDefMidPoint(A,B) \tkzGetPoint{K}
\tkzDefMidPoint(A,D_alt) \tkzGetPoint{L}

\draw[l3] (A)--(B)--(C)--cycle;
\draw[l3] (A)--(D_alt);
\draw[l3] (C)--(D_alt);
\draw[l3] (K)--(L)--(C)--(M)--cycle;
\draw[l3, dashed] (K)--(M);
\draw[l3, dashed] (M)--(L);

\tkzDrawPoints(A,B,C,D_alt,M,K,L)
\tkzLabelPoint[above](A){$A$}
\tkzLabelPoint[below left](B){$B$}
\tkzLabelPoint[below](C){$\varGamma $}
\tkzLabelPoint[below right](D_alt){$\varDelta  $}
\tkzLabelPoint[below](M){$M$}
\tkzLabelPoint[above left](K){$K$}
\tkzLabelPoint[above right](L){$\Lambda$}

\tkzMarkSegments[mark=s|](A,B B,C C,A C,D_alt)
\end{tikzpicture}
\end{center}
\end{alist}
\end{alist}

\vspace{4mm}

% --- Άσκηση 14883 (14883.tex) ---
\Askhsh{14883}{2}
\begin{ylh}{1}
    \item 5.3. Ορθογώνιο
    \item 5.5. Τετράγωνο.
\end{ylh}
% ===================================================================
%  Αρχείο: 14883.tex (Fragment)
%  Αυτόματη μετατροπή μέσω doc2latex.py
% ===================================================================



Σε κύκλο κέντρου Ο φέρουμε τις διαμέτρους του $A\varGamma $ και $B\varDelta  $.

\begin{alist}
\item Να αποδείξετε ότι το τετράπλευρο $AB\varGamma \varDelta  $ είναι ορθογώνιο. \monades{13}

\item Τί είδους γωνία σχηματίζουν οι διάμετροι $A\varGamma $ και $B\varDelta  $ αν το $AB\varGamma \varDelta  $ είναι τετράγωνο; Να δικαιολογήσετε την απάντησή σας. \monades{12}
\end{alist}

\vspace{4mm}

% --- Άσκηση 14884 (14884.tex) ---
\Askhsh{14884}{2}
\begin{ylh}{1}
    \item 3.2. 1ο Κριτήριο ισότητας τριγώνων
    \item 4.2. Τέμνουσα δύο ευθειών - Ευκλείδειο αίτημα
    \item 4.6. Άθροισμα γωνιών τριγώνου.
\end{ylh}
% ===================================================================
%  Αρχείο: 14884.tex (Fragment)
%  Αυτόματη μετατροπή μέσω doc2latex.py
% ===================================================================



Α$BE$ και Α$\varDelta   Z$.

Να αποδείξετε ότι:

\begin{alist}
\item Τα τρίγωνα $AEZ$ και Α$B\varDelta  $ είναι ίσα. \monades{13}

\item Το τμήμα $\varDelta   Z$ είναι παράλληλο στο $BE$. \monades{12}


\begin{center}
\begin{tikzpicture}[scale=0.8]
\tkzDefPoint(0,0){A}
\tkzDefPoint(4,0){B}
\tkzDefPointBy[rotation=center A angle 120](B) \tkzGetPoint{D_dir}
\tkzDefPointOnLine[pos=0.7](A,D_dir) \tkzGetPoint{D} % $AD$ < $AB$ for visual clarity

% Equilateral $ABE$ exterior to $ABD$. A=0, B=4. E should be below
\tkzDefPointBy[rotation=center B angle -60](A) \tkzGetPoint{E}

% Equilateral $ADZ$ exterior to $ABD$.
\tkzDefPointBy[rotation=center A angle -60](D) \tkzGetPoint{Z}
% Wait, "exterior". $ABE$ is on side $AB$. Order is A-B-E. For $ABE$ to be equilateral and exterior:
\tkzDefTriangle[equilateral](A,B) \tkzGetPoint{E_1}
\tkzDefTriangle[equilateral](B,A) \tkzGetPoint{E_2}
% D is above $AB$. So exterior $ABE$ is E_2 (below $AB$).
\tkzDefPointBy[rotation=center A angle -60](B) \tkzGetPoint{E}

% For $ADZ$, D is at 120 degrees from B (y>0). Triangle $ADZ$ should be exterior, so Z must be further left/down.
\tkzDefPointBy[rotation=center D angle -60](A) \tkzGetPoint{Z}

\draw[l3] (A)--(B)--(D)--cycle;
\draw[l3] (A)--(B)--(E)--cycle;
\draw[l3] (A)--(D)--(Z)--cycle;
\draw[l3, dashed] (A)--(E);
\draw[l3, dashed] (A)--(Z);
\draw[l3, dashed] (D)--(Z);
\draw[l3, dashed] (B)--(E);
% Wait, triangle matches the definition of $ADZ$ equilateral, $ABE$ equilateral.
\tkzDrawPoints(A,B,D,E,Z)
\tkzLabelPoint[below](A){$A$}
\tkzLabelPoint[right](B){$B$}
\tkzLabelPoint[above](D){$\varDelta  $}
\tkzLabelPoint[below right](E){$E$}
\tkzLabelPoint[left](Z){$Z$}
\end{tikzpicture}
\end{center}
\end{alist}

\vspace{4mm}

% --- Άσκηση 14885 (14885.tex) ---
\Askhsh{14885}{4}
\begin{ylh}{1}
    \item 3.2. 1ο Κριτήριο ισότητας τριγώνων
    \item 3.6. Κριτήρια ισότητας ορθογώνιων τριγώνων
    \item 5.6. Εφαρμογές στα τρίγωνα
    \item 5.11. Ισοσκελές τραπέζιο.
\end{ylh}
% ===================================================================
%  Αρχείο: 14885.tex (Fragment)
%  Αυτόματη μετατροπή μέσω doc2latex.py
% ===================================================================



Δίνεται το οξυγώνιο και σκαληνό τρίγωνο $AB\varGamma $. Προεκτείνουμε το ύψος του $AH$ κατά τμήμα $H\varDelta  $=$AH$ και τη διάμεσό του $AM$ κατά τμήμα $ME$=$AM$. Να αποδείξετε ότι:\\
\begin{alist}
\item $AB$ = $\varGamma  E$
\item $AB$ = $B\varDelta  $ \monades{8}

\item Γ$\hat{B}$Δ = Β$\hat{\varGamma }$Ε \monades{8}

\item Εξετάστε αν το τμήμα $B\varDelta  $ μπορεί να είναι παράλληλο στο τμήμα $\varGamma  E$. \monades{5}

\item Ποιο είναι το είδος του τετραπλεύρου $B\varGamma $$E\varDelta  $; Δικαιολογήστε την απάντησή σας. \monades{4}


\begin{center}
\begin{tikzpicture}[scale=0.8]
\tkzDefPoint(0,0){B}
\tkzDefPoint(6,0){C}
\tkzDefPoint(2,4){A} % Scalene triangle A=(2,4)

\tkzDefPointBy[projection=onto B--C](A) \tkzGetPoint{H_temp}
% Wait, "προεκτείνουμε το ύψος $AH$ κατά τμήμα $H\varDelta  $=$AH$". Let H be intersection with $BC$.
\tkzInterLL(A,H_temp)(B,C) \tkzGetPoint{H}
\tkzDefPointBy[symmetry=center H](A) \tkzGetPoint{D}

\tkzDefMidPoint(B,C) \tkzGetPoint{M}
\tkzDefPointBy[symmetry=center M](A) \tkzGetPoint{E}

\draw[l3] (A)--(B)--(C)--cycle;
\draw[l3] (A)--(D);
\draw[l3] (A)--(E);
\draw[l3] (B)--(D);
\draw[l3] (C)--(E);
\draw[l3, dashed] (D)--(C); % Quad $BCED$
\draw[l3, dashed] (D)--(E);

\tkzDrawPoints(A,B,C,H,D,M,E)
\tkzLabelPoint[above](A){$A$}
\tkzLabelPoint[below left](B){$B$}
\tkzLabelPoint[below right](C){$\varGamma $}
\tkzLabelPoint[above right](H){$H$}
\tkzLabelPoint[below](D){$\varDelta  $}
\tkzLabelPoint[below](M){$M$}
\tkzLabelPoint[below](E){$E$}

\tkzMarkRightAngle(A,H,C)
\tkzMarkSegments[mark=s|](A,H H,D)
\tkzMarkSegments[mark=s||](A,M M,E)
\end{tikzpicture}
\end{center}
\end{alist}

\vspace{4mm}

% --- Άσκηση 14886 (14886.tex) ---
\Askhsh{14886}{4}
\begin{ylh}{1}
    \item 4.6. Άθροισμα γωνιών τριγώνου
    \item 5.3. Ορθογώνιο
    \item 5.6. Εφαρμογές στα τρίγωνα
    \item 5.9. Μια ιδιότητα του ορθογώνιου τριγώνου.
\end{ylh}
% ===================================================================
%  Αρχείο: 14886.tex (Fragment)
%  Αυτόματη μετατροπή μέσω doc2latex.py
% ===================================================================



Θεωρούμε ένα ορθογώνιο τρίγωνο $AB\varGamma $ ($\hat{A} = 90\degree$), τα μέσα Δ, Ε, Ζ των πλευρών του και το ύψος του $AK$. Αν Θ είναι το σημείο τομής των $AZ$ και $\varDelta   E$, τότε:

\begin{alist}
\item Να αποδείξετε ότι:

\end{alist}
\begin{alist}
\item
  Το τετράπλευρο Α$\varDelta   Z$Ε είναι ορθογώνιο. \monades{8}
\item
  $A\varTheta $=$\varTheta  E$= $\frac{Β\varGamma }{4}$ \monades{7}

\item Αν επιπλέον είναι $\hat{\varGamma }$ = 30\textsuperscript{ο} , τότε:
\end{alist}
\begin{alist}
\item
  να βρείτε τη γωνία Α$\hat{Ζ}$Β. \monades{5}
\item
  να αποδείξετε ότι $BK$ = $\frac{Β\varGamma }{4}$ . \monades{5}
\end{alist}

\vspace{4mm}

% --- Άσκηση 14887 (14887.tex) ---
\Askhsh{14887}{4}
\begin{ylh}{1}
    \item 3.2. 1ο Κριτήριο ισότητας τριγώνων
    \item 3.6. Κριτήρια ισότητας ορθογώνιων τριγώνων
    \item 4.2. Τέμνουσα δύο ευθειών - Ευκλείδειο αίτημα
    \item 5.3. Ορθογώνιο.
\end{ylh}
% ===================================================================
%  Αρχείο: 14887.tex (Fragment)
%  Αυτόματη μετατροπή μέσω doc2latex.py
% ===================================================================



Δίνεται ισοσκελές τρίγωνο $AB\varGamma $ ($AB=A\varGamma $), και τυχαίο σημείο $M$ της πλευράς $B\varGamma $. Από το σημείο $M$ φέρουμε ευθεία κάθετη στην πλευρά $B\varGamma $ που τέμνει τις ευθείες $AB$ και $A\varGamma $ στα σημεία $E$ και Θ αντίστοιχα. Αν $A\varDelta  $ και $AH$ τα ύψη των τριγώνων $AB\varGamma $ και $A\varTheta E$ αντίστοιχα, να αποδείξετε ότι:

\begin{alist}
\item Δ$\hat{A}$Η = 90\textsuperscript{ο} . \monades{8}

\item Το τρίγωνο $A\varTheta E$ είναι ισοσκελές. \monades{8}

\item $M\varTheta $+$ME$ = 2ΑΔ. \monades{9}


\begin{center}
\begin{tikzpicture}[scale=0.8]
\tkzDefPoint(0,0){B}
\tkzDefPoint(6,0){C}
\tkzDefPoint(3,4){A}
\tkzDefPoint(3,0){D}

\tkzDefPoint(1.5,0){M}
\tkzDefLine[orthogonal=through M](B,C) \tkzGetPoint{M_perp}
\tkzInterLL(M,M_perp)(A,B) \tkzGetPoint{E}
\tkzInterLL(M,M_perp)(A,C) \tkzGetPoint{Theta}

\tkzDefPointBy[projection=onto E--Theta](A) \tkzGetPoint{H}

\draw[l3] (A)--(B)--(C)--cycle;
\draw[l3] (A)--(D);
\draw[l3] (M)--(Theta);
\draw[l3] (A)--(E);
\draw[l3] (A)--(Theta);
\draw[l3] (A)--(H);

\tkzDrawPoints(A,B,C,D,M,E,Theta,H)
\tkzLabelPoint[above](A){$A$}
\tkzLabelPoint[below left](B){$B$}
\tkzLabelPoint[below right](C){$\varGamma $}
\tkzLabelPoint[below](D){$\varDelta  $}
\tkzLabelPoint[below](M){$M$}
\tkzLabelPoint[left](E){$E$}
\tkzLabelPoint[right](Theta){$\varTheta $}
\tkzLabelPoint[left](H){$H$}

\tkzMarkRightAngle(A,D,C)
\tkzMarkRightAngle(Theta,M,C)
\tkzMarkRightAngle(A,H,Theta)
\tkzMarkSegments[mark=s|](A,B A,C)
\end{tikzpicture}
\end{center}
\end{alist}

\vspace{4mm}

% --- Άσκηση 14888 (14888.tex) ---
\Askhsh{14888}{4}
\begin{ylh}{1}
    \item 5.2. Παραλληλόγραμμα
    \item 5.6. Εφαρμογές στα τρίγωνα
    \item 5.9. Μια ιδιότητα του ορθογώνιου τριγώνου
    \item 5.10. Τραπέζιο.
\end{ylh}
% ===================================================================
%  Αρχείο: 14888.tex (Fragment)
%  Αυτόματη μετατροπή μέσω doc2latex.py
% ===================================================================



Έστω ορθογώνιο τρίγωνο $AB\varGamma $ με $\hat{A}$ = 90\textsuperscript{ο}. Στην πλευρά $B\varGamma $ θεωρούμε τα σημεία $K$, Μ, Λ ώστε $BK$=$KM$=$M\varLambda $=$\varLambda \varGamma $. Αν τα σημεία $\varDelta  $ και $E$ είναι τα μέσα των πλευρών $AB$ και $A\varGamma $ αντίστοιχα, να αποδείξετε ότι:

\begin{alist}
\item Το τετράπλευρο $\varDelta  E\varLambda K$ είναι παραλληλόγραμμο. \monades{13}

\item Το τετράπλευρο $K\varDelta  AM$ είναι τραπέζιο και η διάμεσός του ισούται με $\frac{3}{8}$ $B\varGamma $ .

\monades{12}


\begin{center}
\begin{tikzpicture}[scale=1]
\tkzDefPoint(0,4){B}
\tkzDefPoint(0,0){A}
\tkzDefPoint(6,0){C}

\tkzDefPointOnLine[pos=0.25](B,C) \tkzGetPoint{K}
\tkzDefPointOnLine[pos=0.5](B,C) \tkzGetPoint{M}
\tkzDefPointOnLine[pos=0.75](B,C) \tkzGetPoint{L}

\tkzDefMidPoint(A,B) \tkzGetPoint{D}
\tkzDefMidPoint(A,C) \tkzGetPoint{E}

\draw[l3] (A)--(B)--(C)--cycle;
\draw[l3] (D)--(E);
\draw[l3] (D)--(K);
\draw[l3] (E)--(L);
\draw[l3, dashed] (A)--(M);
\draw[l3, dashed] (D)--(M);

\tkzDrawPoints(A,B,C,K,M,L,D,E)
\tkzLabelPoint[below left](A){$A$}
\tkzLabelPoint[above left](B){$B$}
\tkzLabelPoint[below right](C){$\varGamma $}
\tkzLabelPoint[above right](K){$K$}
\tkzLabelPoint[above right](M){$M$}
\tkzLabelPoint[above right](L){$\Lambda$}
\tkzLabelPoint[left](D){$\varDelta  $}
\tkzLabelPoint[below](E){$E$}

\tkzMarkRightAngle(C,A,B)
\tkzMarkSegments[mark=s|](B,K K,M M,L L,C)
\end{tikzpicture}
\end{center}
\end{alist}

\vspace{4mm}

% --- Άσκηση 34313 (34313.tex) ---
\Askhsh{34313}{2}
\begin{ylh}{1}
    \item 3.1. Είδη και στοιχεία τριγώνων
    \item 3.2. 1ο Κριτήριο ισότητας τριγώνων
    \item 3.14. Σχετικές θέσεις ευθείας και κύκλου
    \item 4.6. Άθροισμα γωνιών τριγώνου
    \item 5.9. Μια ιδιότητα του ορθογώνιου τριγώνου.
\end{ylh}
% ===================================================================
%  Αρχείο: 34313.tex (Fragment)
%  Αυτόματη μετατροπή μέσω doc2latex.py
% ===================================================================



Δίνεται κύκλος (Ο,R) διαμέτρου $AB$, και χορδή του $A\varGamma $ τέτοια ώστε Β$\hat{A}$Γ = 30\degree. Στο σημείο $\varGamma $ του κύκλου φέρουμε εφαπτομένη, η οποία τέμνει την προέκταση της διαμέτρου $AB$ (προς το Β) σε σημείο $\varDelta  $.

\begin{alist}
\item Να αποδείξετε ότι το τρίγωνο $\varGamma O\varDelta  $ είναι ορθογώνιο και να υπολογίσετε το μέτρο της γωνίας του Γ$\hat{O}$Δ. \monades{9}

\item Να υπολογίσετε το μέτρο της γωνίας $\hat{\omega}$. \monades{8}

\item Να αποδείξετε ότι $O\varDelta  $ = $2R$. \monades{8}


\begin{center}
\begin{tikzpicture}[scale=1]
\tkzDefPoint(0,0){O}
\tkzDefPoint(-2,0){A}
\tkzDefPoint(2,0){B}
\tkzDrawCircle(O,A)

% $BAC$ = 30. A is (-2,0), B is (2,0). $AC$ makes 30 deg with $AB$.
% C is on circle. Angle $BOC$ = 2*$BAC$ = 60 degrees.
% So C = (2*cos(120), 2*sin(120)) = (-1, 1.732)
\tkzDefPointBy[rotation=center O angle -120](A) \tkzGetPoint{C}

% Tangent at C
\tkzDefLine[orthogonal=through C](O,C) \tkzGetPoint{tangent_dir}
\tkzInterLL(C,tangent_dir)(A,B) \tkzGetPoint{D}

\draw[l3] (A)--(D);
\draw[l3] (A)--(C);
\draw[l3] (O)--(C);
\draw[l3] (C)--(D);

\tkzDrawPoints(O,A,B,C,D)
\tkzLabelPoint[below](O){$O$}
\tkzLabelPoint[left](A){$A$}
\tkzLabelPoint[below right](B){$B$}
\tkzLabelPoint[above](C){$\varGamma $}
\tkzLabelPoint[right](D){$\varDelta  $}

\tkzMarkRightAngle(O,C,D)
\tkzMarkAngle[size=0.5,mark=none](B,A,C)
\tkzLabelAngle[pos=0.8](B,A,C){$30^\circ$}
\tkzMarkAngle[size=0.5,mark=none](C,D,O)
\tkzLabelAngle[pos=0.8](C,D,O){$\omega$}
\end{tikzpicture}
\end{center}
\end{alist}

\vspace{4mm}

% --- Άσκηση 34314 (34314.tex) ---
\Askhsh{34314}{2}
\begin{ylh}{1}
    \item 3.1. Είδη και στοιχεία τριγώνων
    \item 4.6. Άθροισμα γωνιών τριγώνου
    \item 5.9. Μια ιδιότητα του ορθογώνιου τριγώνου.
\end{ylh}
% ===================================================================
%  Αρχείο: 34314.tex (Fragment)
%  Αυτόματη μετατροπή μέσω doc2latex.py
% ===================================================================



Θεωρούμε δυο ίσους κύκλους (Ο,ρ) και (Κ,ρ) τεμνόμενους στα σημεία $A$ και Δ, με τα κέντρα τους Ο και Κ να βρίσκονται σε ευθεία (ε) που τέμνει τον κύκλο (Ο,ρ) σε σημείο $B$, και τη διάκεντρό τους $OK$ να είναι ίση με ρ.

\begin{alist}
\item Να αποδείξετε ότι:

\begin{alist}
\item
  το τρίγωνο $OAK$ είναι ισόπλευρο, \monades{7}
\item
  το τρίγωνο $ABK$ είναι ορθογώνιο. \monades{9}

\item Να υπολογίσετε τη γωνία $ABK$. \monades{9}


\begin{center}
\begin{tikzpicture}[scale=0.8]
\tkzDefPoint(0,0){O}
\tkzDefPoint(3,0){K} % radius is 3
\tkzDrawCircle(O,K)
\tkzDrawCircle(K,O)

\tkzInterCC(O,K)(K,O) \tkzGetPoints{A}{D}

% B is the intersection of $OK$ with circle O, other than K.
\tkzDefPointBy[symmetry=center O](K) \tkzGetPoint{B}
% Draw line epsilon
\tkzDrawLine[add=0.2 and 0.5](B,K) % extends past B and K

\draw[l3] (O)--(A)--(K)--cycle; % Triangle $OAK$
\draw[l3] (A)--(B);
\draw[l3] (A)--(K);
\draw[l3] (B)--(K);

\tkzDrawPoints(O,K,A,D,B)
\tkzLabelPoint[below](O){$O$}
\tkzLabelPoint[below right](K){$K$}
\tkzLabelPoint[above](A){$A$}
\tkzLabelPoint[below](D){$\varDelta  $}
\tkzLabelPoint[below left](B){$B$}
\tkzLabelPoint[right](K){$\epsilon$} % approximate

\tkzMarkRightAngle(B,A,K)
\tkzMarkSegments[mark=s|](O,A A,K K,O)
\end{tikzpicture}
\end{center}
\end{alist}
\end{alist}

\vspace{4mm}

% --- Άσκηση 34315 (34315.tex) ---
\Askhsh{34315}{2}
\begin{ylh}{1}
    \item 3.1. Είδη και στοιχεία τριγώνων
    \item 3.15. Εφαπτόμενα τμήματα
    \item 4.8. Άθροισμα γωνιών κυρτού ν-γώνου.
\end{ylh}
% ===================================================================
%  Αρχείο: 34315.tex (Fragment)
%  Αυτόματη μετατροπή μέσω doc2latex.py
% ===================================================================



Δίνεται κύκλος κέντρου Ο και ακτίνας 4 cm και εξωτερικό του σημείο $P$. Έστω $PA$, $PB$ τα εφαπτόμενα τμήματα που φέρονται από το Ρ, τα σημεία επαφής τους A, B με τον κύκλο αντίστοιχα και τέτοια ώστε η γωνία Α$\hat{Ρ}$Β να ισούται με 60\degree.

\begin{alist}
\item Να αποδείξετε ότι το μέτρο της γωνίας Α$\hat{O}$Β είναι ίσο με 120\degree. \monades{7}

\item Αν $PO$ η διακεντρική ευθεία του σημείου Ρ, τότε να υπολογίσετε:

\begin{alist}
\item
  το μέτρο της γωνίας Α$\hat{Ρ}$Ο, \monades{9}
\item
  το μήκος του τμήματος $OP$. \monades{9}


\begin{center}
\begin{tikzpicture}[scale=0.8]
\tkzDefPoint(0,0){O}
% Radius is 4, R is at distance 8.
\tkzDefPoint(6,0){P} % Let's use P = (-6,0) for left side placement? No, P=Ρ
% $PA$ and $PB$ are tangents.
\tkzDefPointBy[rotation=center O angle 60](P) \tkzGetPoint{A_dir} % Wait. Angle $AOP$ = 60.
% Distance $OP$ = 8. Angle $APO$ = 30. Angle $AOP$ = 60.
\tkzDefPoint(3, 5.196){A_temp} % 6 * cos60 = 3, 6 * sin60 = 5.196 (if $OP$=6, radius=3)
% Let's just use tkz-euclide tangent macro
\tkzDefPoint(2,0){O_radius}
\tkzDrawCircle(O,O_radius)

% P is outside. Let's place it on the left at distance 4 from O (radius is 2 in scale)
\tkzDefPoint(-4,0){P}
%\tkzExtCircularAlley[R](O,2)(P)
\tkzDefTangent[from=P](O,O_radius) \tkzGetPoints{B}{A} % B is top, A is bottom usually, let's see. A and B are the points.

\draw[l3] (P)--(A);
\draw[l3] (P)--(B);
\draw[l3] (O)--(A);
\draw[l3] (O)--(B);
\draw[l3, dashed] (O)--(P);

\tkzDrawPoints(O,P,A,B)
\tkzLabelPoint[right](O){$O$}
\tkzLabelPoint[left](P){$P$}
\tkzLabelPoint[above](A){$A$}
\tkzLabelPoint[below](B){$B$}

\tkzMarkRightAngle(O,A,P)
\tkzMarkRightAngle(O,B,P)
\tkzMarkAngle[size=0.6,mark=none](B,P,A)
\tkzLabelAngle[pos=1](B,P,A){$60^\circ$}
\end{tikzpicture}
\end{center}
\end{alist}
\end{alist}

\vspace{4mm}

% --- Άσκηση 34316 (34316.tex) ---
\Askhsh{34316}{4}
\begin{ylh}{1}
    \item 3.2. 1ο Κριτήριο ισότητας τριγώνων
    \item 3.4. 3ο Κριτήριο ισότητας τριγώνων
    \item 3.14. Σχετικές θέσεις ευθείας και κύκλου
    \item 4.8. Άθροισμα γωνιών κυρτού ν-γώνου.
\end{ylh}
% ===================================================================
%  Αρχείο: 34316.tex (Fragment)
%  Αυτόματη μετατροπή μέσω doc2latex.py
% ===================================================================



Δίνεται κύκλος κέντρου Ο και διάμετρός του $AB$. Έστω $\varGamma$ το μέσο ενός ημικυκλίου του, Δ τυχαίο σημείο του άλλου ημικυκλίου του και Α$\hat{\varGamma }$Β = Α$\hat{\varDelta  }$Β = 90\degree. Στην προέκταση της $\varDelta   B$ προς το μέρος του Β θεωρούμε σημείο $E$ ώστε $BE$ = $A\varDelta  $.

\begin{alist}
\item Να αποδείξετε ότι:

\begin{alist}
\item
  οι γωνίες Γ$\hat{A}$Δ και Γ$\hat{B}$Ε είναι ίσες, \monades{6}
\item
  τα τρίγωνα $A\varDelta  \varGamma $ και $BE$Γ είναι ίσα, \monades{7}
\item
  η $\varGamma \varDelta  $ είναι κάθετη στην $\varGamma E$. \monades{6}

\item Να αιτιολογήσετε γιατί, στην περίπτωση που το σημείο $\varDelta  $ είναι αντιδιαμετρικό του Γ, η $\varGamma E$ είναι εφαπτόμενη του κύκλου. \monades{6}


\begin{center}
\begin{tikzpicture}[scale=1]
\tkzDefPoint(0,0){O}
\tkzDefPoint(-2,0){A}
\tkzDefPoint(2,0){B}
\tkzDrawCircle(O,A)
\tkzDefPoint(0,2){C} % Γ is midpoint of top semicircle

% Δ is random point on bottom semicircle
\tkzDefPointBy[rotation=center O angle -50](B) \tkzGetPoint{D}

% E is on extension of $DB$ past B, with $BE$ = $AD$
% E = B + (B-D)/|B-D| * |$AD$|
% Use tkzDefPointWith to define E along $BD$ direction from B
\tkzDefPointWith[colinear= at B](D,B) \tkzGetPoint{E_dir}
\tkzCalcLength(A,D) \tkzGetLength{lenAD}
\tkzCalcLength(D,B) \tkzGetLength{lenDB}
\pgfmathsetmacro{\scaleRatio}{\lenAD/\lenDB}
\tkzDefPointOnLine[pos=\scaleRatio](B,E_dir) \tkzGetPoint{E}

\draw[l3] (A)--(B);
\draw[l3] (D)--(E);
\draw[l3] (A)--(D);
\draw[l3] (A)--(C);
\draw[l3] (B)--(C);
\draw[l3] (D)--(C);
\draw[l3] (C)--(E);

\tkzDrawPoints(O,A,B,C,D,E)
\tkzLabelPoint[below](O){$ O$}
\tkzLabelPoint[left](A){$A$}
\tkzLabelPoint[below right](B){$B$}
\tkzLabelPoint[above](C){$\varGamma $}
\tkzLabelPoint[below left](D){$\varDelta  $}
\tkzLabelPoint[above right](E){$ E$}

\tkzMarkSegments[mark=s|](A,D B,E)
\end{tikzpicture}
\end{center}
\end{alist}
\end{alist}

\vspace{4mm}

% --- Άσκηση 34318 (34318.tex) ---
\Askhsh{34318}{4}
\begin{ylh}{1}
    \item 3.1. Είδη και στοιχεία τριγώνων
    \item 3.4. 3ο Κριτήριο ισότητας τριγώνων
    \item 3.11. Ανισοτικές σχέσεις πλευρών και γωνιών
    \item 3.14. Σχετικές θέσεις ευθείας και κύκλου
    \item 4.2. Τέμνουσα δύο ευθειών - Ευκλείδειο αίτημα
    \item 4.6. Άθροισμα γωνιών τριγώνου
    \item 5.9. Μια ιδιότητα του ορθογώνιου τριγώνου.
\end{ylh}
% ===================================================================
%  Αρχείο: 34318.tex (Fragment)
%  Αυτόματη μετατροπή μέσω doc2latex.py
% ===================================================================



Θεωρούμε κύκλο (Ο, ρ) και $E$ το μέσο του τόξου του $B\varGamma $. Μια ευθεία (ε) εφάπτεται στον κύκλο στο Ε. Οι προεκτάσεις των $OB$, $O\varGamma $ (προς το Β και το $\varGamma$ αντίστοιχα) τέμνουν την ευθεία (ε) στα σημεία $Z$ και Η αντίστοιχα.

\begin{alist}
\item Να αποδείξετε ότι:

\begin{alist}
\item
  $B\varGamma $ // (ε) \monades{9}
\item
  $OZ$ = $OH$ \monades{9}

\item Αν το σημείο $B$ είναι το μέσο του $OZ$ , να υπολογίσετε τις γωνίες του τριγώνου $ZOH$.

\monades{7}


\begin{center}
\begin{tikzpicture}[scale=1]
\tkzDefPoint(0,0){O}
\tkzDefPoint(0,-2){E} % E is at the bottom
\tkzDrawCircle(O,E)

% E is midpoint of arc $BC$. Angles $BOE$ and $EOC$ must be equal.
\tkzDefPointBy[rotation=center O angle -40](E) \tkzGetPoint{B}
\tkzDefPointBy[rotation=center O angle 40](E) \tkzGetPoint{C}

% Tangent at E
\tkzDefLine[orthogonal=through E](O,E) \tkzGetPoint{tangent_dir}

% Extents of $OB$ and $OC$
\tkzInterLL(E,tangent_dir)(O,B) \tkzGetPoint{Z}
\tkzInterLL(E,tangent_dir)(O,C) \tkzGetPoint{H}

\draw[l3] (B)--(C);
\draw[l3] (O)--(Z);
\draw[l3] (O)--(H);
\draw[l3] (O)--(E);
\draw[l3] (Z)--(H); % Tangent line segment

\tkzDrawPoints(O,A,B,C,E,Z,H) % A doesn't exist but ok
\tkzDrawPoints(O,B,C,E,Z,H)
\tkzLabelPoint[above](O){$O$}
\tkzLabelPoint[left](B){$B$}
\tkzLabelPoint[right](C){$\varGamma $}
\tkzLabelPoint[below](E){$E$}
\tkzLabelPoint[below left](Z){$Z$}
\tkzLabelPoint[below right](H){$H$}

\tkzMarkRightAngle(O,E,H)
\tkzMarkSegments[mark=s|](O,B O,C)
\end{tikzpicture}
\end{center}
\end{alist}
\end{alist}

\vspace{4mm}

% --- Άσκηση 34320 (34320.tex) ---
\Askhsh{34320}{4}
\begin{ylh}{1}
    \item 3.1. Είδη και στοιχεία τριγώνων
    \item 4.2. Τέμνουσα δύο ευθειών - Ευκλείδειο αίτημα
    \item 4.6. Άθροισμα γωνιών τριγώνου
    \item 5.10. Τραπέζιο.
\end{ylh}
% ===================================================================
%  Αρχείο: 34320.tex (Fragment)
%  Αυτόματη μετατροπή μέσω doc2latex.py
% ===================================================================



Έστω Α, Β και $\varGamma$ συνευθειακά σημεία με $AB$ = 2$B\varGamma $. Θεωρούμε το μέσο $M$ της $AB$. Προς το ίδιο ημιεπίεδο κατασκευάζουμε τα ισόπλευρα τρίγωνα $A\varDelta  B$ και $BE$Γ.

Να αποδείξετε ότι:

\begin{alist}
\item το τετράπλευρο $A\varDelta  EB$ είναι τραπέζιο με βάσεις τα τμήματα $A\varDelta  $ και $BE$, \monades{9}

\item τα τρίγωνα $\varDelta  MB$ και $\varDelta  EB$ είναι ίσα, \monades{8}

\item Δ$\hat{Μ}$Β + Δ$\hat{Ε}$Β = 180\degree. \monades{8}


\begin{center}
\begin{tikzpicture}[scale=0.8]
\tkzDefPoint(0,0){A}
\tkzDefPoint(4,0){B}
\tkzDefPoint(6,0){C}
\tkzDefMidPoint(A,B) \tkzGetPoint{M}

\tkzDefTriangle[equilateral](A,B) \tkzGetPoint{D}
\tkzDefTriangle[equilateral](B,C) \tkzGetPoint{E}

\draw[l3] (A)--(C);
\draw[l3] (A)--(D)--(B)--cycle;
\draw[l3] (B)--(E)--(C)--cycle;
\draw[l3, dashed] (D)--(E);
\draw[l3, dashed] (D)--(M);

\tkzDrawPoints(A,B,C,M,D,E)
\tkzLabelPoint[below](A){$A$}
\tkzLabelPoint[below](B){$B$}
\tkzLabelPoint[below](C){$\varGamma $}
\tkzLabelPoint[below](M){$M$}
\tkzLabelPoint[above](D){$\varDelta  $}
\tkzLabelPoint[above right](E){$E$}

\tkzMarkSegments[mark=s|](A,M M,B B,C B,E E,C)
\tkzMarkSegments[mark=s||](A,D D,B)
\end{tikzpicture}
\end{center}
\end{alist}

\vspace{4mm}

% --- Άσκηση 34321 (34321.tex) ---
\Askhsh{34321}{4}
\begin{ylh}{1}
    \item 3.1. Είδη και στοιχεία τριγώνων
    \item 3.2. 1ο Κριτήριο ισότητας τριγώνων
    \item 4.6. Άθροισμα γωνιών τριγώνου.
\end{ylh}
% ===================================================================
%  Αρχείο: 34321.tex (Fragment)
%  Αυτόματη μετατροπή μέσω doc2latex.py
% ===================================================================



Δίνεται ισοσκελές τρίγωνο $AB\varGamma $ με $AB=A\varGamma $ και $A\varDelta  $, $BE$ τα ύψη του.

Να αποδείξετε ότι:

\begin{alist}
\item $B\varGamma $ = 2$E\varDelta  $, \monades{7}

\item $\hat{E}$\textsubscript{1} = $\frac{\hat{A}}{2}$ , \monades{9}

\item $\hat{B}$\textsubscript{2} = $\hat{\varDelta  }$\textsubscript{2} \monades{9}


\begin{center}
\begin{tikzpicture}[scale=1]
\tkzDefPoint(0,0){B}
\tkzDefPoint(6,0){C}
\tkzDefPoint(3,5){A}

\tkzDefPointBy[projection=onto B--C](A) \tkzGetPoint{D}
\tkzDefPointBy[projection=onto A--C](B) \tkzGetPoint{E}

\draw[l3] (A)--(B)--(C)--cycle;
\draw[l3] (A)--(D);
\draw[l3] (B)--(E);
\draw[l3, dashed] (D)--(E);

\tkzDrawPoints(A,B,C,D,E)
\tkzLabelPoint[above](A){$A$}
\tkzLabelPoint[below left](B){$B$}
\tkzLabelPoint[below right](C){$\varGamma $}
\tkzLabelPoint[below](D){$\varDelta  $}
\tkzLabelPoint[above right](E){$E$}

\tkzMarkRightAngle(A,D,C)
\tkzMarkRightAngle(B,E,A)
\tkzMarkSegments[mark=s|](A,B A,C)
\tkzMarkSegments[mark=s||](E,D B,D D,C) % $ED$=$BD$=$DC$=$BC$/2 Since $BED$ is right triangle median
\end{tikzpicture}
\end{center}
\end{alist}

\vspace{4mm}

% --- Άσκηση 34324 (34324.tex) ---
\Askhsh{34324}{4}
\begin{ylh}{1}
    \item 3.11. Ανισοτικές σχέσεις πλευρών και γωνιών
    \item 3.14. Σχετικές θέσεις ευθείας και κύκλου
    \item 3.15. Εφαπτόμενα τμήματα
    \item 4.2. Τέμνουσα δύο ευθειών - Ευκλείδειο αίτημα
    \item 5.9. Μια ιδιότητα του ορθογώνιου τριγώνου
    \item 5.10. Τραπέζιο
    \item 5.11. Ισοσκελές τραπέζιο.
\end{ylh}
% ===================================================================
%  Αρχείο: 34324.tex (Fragment)
%  Αυτόματη μετατροπή μέσω doc2latex.py
% ===================================================================



Δίνεται κύκλος κέντρου Ο και ακτίνας ρ. Από σημείο $A$ εξωτερικό του κύκλου θεωρούμε τις εφαπτόμενες του κύκλου που εφάπτονται σε αυτόν στα σημεία $B$, Γ και τέτοιες ώστε, η γωνία $B\hat{A}\varGamma $ που σχηματίζουν τα εφαπτόμενα τμήματα $AB$ και $A\varGamma $ να είναι $60\degree$. Έστω ότι η ευθεία $AO$ τέμνει τον κύκλο στο σημείο $\varDelta  $ και η εφαπτόμενη του κύκλου στο $\varDelta$ τέμνει τα τμήματα $AB$ και $A\varGamma $ στα σημεία $Z$ και $E$ αντίστοιχα.

Να αποδείξετε ότι:

\begin{alist}
\item $OA$=2ρ, \monades{6}

\item το τρίγωνο $AZE$ είναι ισόπλευρο, \monades{5}

\item $AZ$ = 2ΖΒ, \monades{7}

\item το τετράπλευρο ΕΖ$B\varGamma $ είναι ισοσκελές τραπέζιο. \monades{7}


\begin{center}
\begin{tikzpicture}[scale=1]
\tkzDefPoint(0,0){O}
\tkzDefPoint(4,0){A}
\tkzDefPoint(2,0){D} % Point D is on $OA$, radius is 2 since $OA$=2R and angle is 30.
\tkzDrawCircle(O,D)

\tkzDefTangent[from=A](O,D) \tkzGetPoints{B}{C}

% Tangent at D
\tkzDefLine[orthogonal=through D](O,A) \tkzGetPoint{tangent_D_dir}
\tkzInterLL(D,tangent_D_dir)(A,B) \tkzGetPoint{Z}
\tkzInterLL(D,tangent_D_dir)(A,C) \tkzGetPoint{E}

\draw[l3] (A)--(O);
\draw[l3] (A)--(B);
\draw[l3] (A)--(C);
\draw[l3] (Z)--(E);
\draw[l3] (O)--(B);
\draw[l3] (O)--(C);
\draw[l3, dashed] (B)--(C);

\tkzDrawPoints(O,A,B,C,D,Z,E)
\tkzLabelPoint[left](O){$O$}
\tkzLabelPoint[right](A){$A$}
\tkzLabelPoint[above](B){$B$}
\tkzLabelPoint[below](C){$\varGamma $}
\tkzLabelPoint[above right](D){$\varDelta  $}
\tkzLabelPoint[above right](Z){$Z$}
\tkzLabelPoint[below right](E){$E$}

\tkzMarkRightAngle(O,B,A)
\tkzMarkRightAngle(O,C,A)
\tkzMarkRightAngle(A,D,Z)
\tkzMarkAngle[size=0.6,mark=none](C,A,B)
\tkzLabelAngle[pos=1](C,A,B){$60^\circ$}
\end{tikzpicture}
\end{center}
\end{alist}

\vspace{4mm}

% --- Άσκηση 34326 (34326.tex) ---
\Askhsh{34326}{4}
\begin{ylh}{1}
    \item 3.1. Είδη και στοιχεία τριγώνων
    \item 3.4. 3ο Κριτήριο ισότητας τριγώνων
    \item 3.6. Κριτήρια ισότητας ορθογώνιων τριγώνων
    \item 4.2. Τέμνουσα δύο ευθειών - Ευκλείδειο αίτημα
    \item 5.2. Παραλληλόγραμμα
    \item 5.3. Ορθογώνιο.
\end{ylh}
% ===================================================================
%  Αρχείο: 34326.tex (Fragment)
%  Αυτόματη μετατροπή μέσω doc2latex.py
% ===================================================================



Έστω κύκλος κέντρου Ο και $A\varGamma $ μια διάμετρός του. Θεωρούμε δυο ίσες χορδές $A\varDelta  $, $B\varGamma $ και χορδές $\varDelta  \varGamma $, $AB$ τέτοιες ώστε να είναι κάθετες στις $A\varDelta  $, $B\varGamma $ αντίστοιχα. Έστω Κ και Λ τα μέσα των χορδών $\varDelta  \varGamma $ και $B\varGamma $ αντίστοιχα.

Να αποδείξετε ότι:

\begin{alist}
\item οι χορδές $AB$ και $\varDelta  \varGamma $ είναι παράλληλες, \monades{7}

\item το τετράπλευρο $AB\varGamma \varDelta  $ είναι ορθογώνιο παραλληλόγραμμο, \monades{7}

\item η $B\varDelta  $ είναι διάμετρος του κύκλου, \monades{6}

\item το τετράπλευρο $OK\varGamma \varLambda $ είναι ορθογώνιο. \monades{5}


\begin{center}
\begin{tikzpicture}[scale=1]
\tkzDefPoint(0,0){O}
\tkzDefPoint(-3,0){A}
\tkzDefPoint(3,0){C}
\tkzDrawCircle(O,A)

% $AD$ = $BC$. $ABCD$ is a rectangle inscribed in circle.
% Let's put B at (-1.5, -2.598) and D at (1.5, 2.598).
% Wait, A is (-3,0), C is (3,0). $AC$ is diameter.
% If D is on the top semicircle, angle $ADC$ = 90.
\tkzDefPointBy[rotation=center O angle 120](C) \tkzGetPoint{D}
% B is diametrically opposite to D? No, $AC$ is diameter.
% For $AD$=$BC$, $ABCD$ is an isosceles trapezoid or rectangle. Since angle is 90, it's a rectangle.
\tkzDefPointBy[symmetry=center O](D) \tkzGetPoint{B}

\tkzDefMidPoint(D,C) \tkzGetPoint{K}
\tkzDefMidPoint(B,C) \tkzGetPoint{L}

\draw[l3] (A)--(B)--(C)--(D)--cycle;
\draw[l3] (A)--(C);
\draw[l3] (B)--(D);
\draw[l3] (O)--(K);
\draw[l3] (O)--(L);
\draw[l3] (K)--(C)--(L);

\tkzDrawPoints(O,A,B,C,D,K,L)
\tkzLabelPoint[left](A){$A$}
\tkzLabelPoint[below](B){$B$}
\tkzLabelPoint[right](C){$\varGamma $}
\tkzLabelPoint[above](D){$\varDelta  $}
\tkzLabelPoint[below](O){$O$}
\tkzLabelPoint[above right](K){$K$}
\tkzLabelPoint[below right](L){$\Lambda$}

\tkzMarkRightAngle(A,D,C)
\tkzMarkRightAngle(A,B,C)
\end{tikzpicture}
\end{center}
\end{alist}

\vspace{4mm}

% --- Άσκηση 34328 (34328.tex) ---
\Askhsh{34328}{4}
\begin{ylh}{1}
    \item 3.1. Είδη και στοιχεία τριγώνων
    \item 3.6. Κριτήρια ισότητας ορθογώνιων τριγώνων
    \item 3.11. Ανισοτικές σχέσεις πλευρών και γωνιών
    \item 4.2. Τέμνουσα δύο ευθειών - Ευκλείδειο αίτημα
    \item 4.8. Άθροισμα γωνιών κυρτού ν-γώνου
    \item 5.4. Ρόμβος
    \item 5.9. Μια ιδιότητα του ορθογώνιου τριγώνου.
\end{ylh}
% ===================================================================
%  Αρχείο: 34328.tex (Fragment)
%  Αυτόματη μετατροπή μέσω doc2latex.py
% ===================================================================


Δίνεται τετράπλευρο με κορυφές σημεία $A$, Β, Γ και $\varDelta$ κύκλου (Ο, ρ), διαγώνιες $A\varGamma $ και $B\varDelta  $, με τη $B\varDelta  $ να διέρχεται από το κέντρο Ο του κύκλου, και με ίσες πλευρές τις $AB$ και $B\varGamma $. Έστω ότι η κάθετη στη $B\varDelta  $ στο σημείο $O$ τέμνει τις πλευρές $A\varDelta  $ και $\varGamma \varDelta  $ του τετράπλευρου $AB\varGamma \varDelta  $ στα σημεία $E$ και $Z$ αντίστοιχα, οι γωνίες του $\hat{A}$ και $\hat{\varGamma }$ είναι ορθές και η γωνία του $\hat{B}$ είναι διπλάσια της γωνίας του $\hat{\varDelta  }$.

\begin{alist}
\item Να υπολογίσετε τα μέτρα των γωνιών $\hat{B}$ και $\hat{\varDelta  }$ του $AB\varGamma \varDelta  $. \monades{6}

\item Να αποδείξετε ότι:

\begin{alist}
\item
  η διαγώνιος $B\varDelta  $ διχοτομεί τη γωνία $\hat{\varDelta  }$ του $AB\varGamma \varDelta  $, \monades{7}
\item
  το τετράπλευρο $AB\varGamma $Ο είναι ρόμβος, \monades{6}
\item
  το τετράπλευρο $A\varGamma  ZE$ είναι τραπέζιο. \monades{6}


\begin{center}
\begin{tikzpicture}[scale=1]
\tkzDefPoint(0,0){O}
\tkzDefPoint(2,0){D}
\tkzDefPoint(-2,0){B}
\tkzDrawCircle(O,D)

% Angle B is 120, D is 60. $AC$ is orthogonal to $BD$ and bisects it? No, $AB$=$BC$.
% Since $ABCD$ is cyclic and $AB$=$BC$, B is on the perpendicular bisector of $AC$.
% Wait, $BD$ is diameter. So triangle $DAB$ has A=90, triangle $DCB$ has C=90.
% Also $AB$=$BC$, so right triangles $DAB$ and $DCB$ are congruent.
% So $AD$=$CD$.
% Since Angle B = 2*Angle D. B+D = 180 => B=120, D=60.
% In right triangle $DAB$, angle $ABD$ = 60, angle $ADB$ = 30.
\tkzDefPointBy[rotation=center O angle 120](D) \tkzGetPoint{A}
\tkzDefPointBy[rotation=center O angle -120](D) \tkzGetPoint{C}

% The perpendicular to $BD$ at O intersects $AD$ at E and $CD$ at Z.
\tkzDefLine[orthogonal=through O](B,D) \tkzGetPoint{perp_dir}
\tkzInterLL(O,perp_dir)(A,D) \tkzGetPoint{E}
\tkzInterLL(O,perp_dir)(C,D) \tkzGetPoint{Z}

\draw[l3] (A)--(B)--(C)--(D)--cycle;
\draw[l3] (B)--(D);
\draw[l3] (A)--(C);
\draw[l3] (E)--(Z);
\draw[l3] (A)--(O);
\draw[l3] (C)--(O);

\tkzDrawPoints(O,A,B,C,D,E,Z)
\tkzLabelPoint[left](B){$B$}
\tkzLabelPoint[right](D){$\varDelta  $}
\tkzLabelPoint[above](A){$A$}
\tkzLabelPoint[below](C){$\varGamma $}
\tkzLabelPoint[above](E){$E$}
\tkzLabelPoint[below](Z){$Z$}
\tkzLabelPoint[below left](O){$O$}

\tkzMarkRightAngle(B,A,D)
\tkzMarkRightAngle(B,C,D)
\tkzMarkRightAngle(Z,O,D)
\end{tikzpicture}
\end{center}
\end{alist}
\end{alist}

\vspace{4mm}

% --- Άσκηση 34329 (34329.tex) ---
\Askhsh{34329}{4}
\begin{ylh}{1}
    \item 3.1. Είδη και στοιχεία τριγώνων
    \item 3.2. 1ο Κριτήριο ισότητας τριγώνων
    \item 3.11. Ανισοτικές σχέσεις πλευρών και γωνιών
    \item 3.14. Σχετικές θέσεις ευθείας και κύκλου
    \item 3.15. Εφαπτόμενα τμήματα
    \item 4.6. Άθροισμα γωνιών τριγώνου
    \item 5.9. Μια ιδιότητα του ορθογώνιου τριγώνου.
\end{ylh}
% ===================================================================
%  Αρχείο: 34329.tex (Fragment)
%  Αυτόματη μετατροπή μέσω doc2latex.py
% ===================================================================



$\hat{A}=90^{\circ}$
). Με διάμετρο την κάθετη πλευρά του $A\varGamma $ φέρουμε κύκλο κέντρου Ο, ο οποίος τέμνει την πλευρά $B\varGamma $ του τριγώνου σε σημείο $\varDelta  $. Έστω ότι η εφαπτόμενη του κύκλου στο σημείο $\varDelta  $ τέμνει την πλευρά $AB$ σε σημείο $M$.

Να αποδείξετε ότι:

\begin{alist}
\item $\varGamma \hat{A}\varDelta   = \hat{B}$, \monades{9}
, \monades{9}

\item $Μ\hat{\varDelta  }Β = 90\degree - \hat{\varGamma }$
 και το τρίγωνο $\varDelta   MB$ είναι ισοσκελές, \monades{9}

\item το $M$ είναι το μέσο του $AB$. \monades{7}


\begin{center}
\begin{tikzpicture}[scale=1]
\tkzDefPoint(0,4){A}
\tkzDefPoint(0,0){B}
\tkzDefPoint(6,4){C}
\tkzDefMidPoint(A,C) \tkzGetPoint{O}

% Circle with diameter $AC$. Radius is 3. Center O is (3,4).
\tkzDrawCircle(O,A)

% D is intersection of circle with $BC$.
% Angle $ADC$ = 90 since $AC$ is diameter. So $AD$ is altitude to $BC$.
\tkzDefPointBy[projection=onto B--C](A) \tkzGetPoint{D}

% Tangent at D. It's perpendicular to $OD$.
\tkzDefLine[orthogonal=through D](O,D) \tkzGetPoint{tangent_D_dir}
\tkzInterLL(D,tangent_D_dir)(A,B) \tkzGetPoint{M}

\draw[l3] (A)--(B)--(C)--cycle;
\draw[l3] (A)--(D);
\draw[l3] (O)--(D);
\draw[l3] (M)--(D);

\tkzDrawPoints(A,B,C,O,D,M)
\tkzLabelPoint[above left](A){$A$}
\tkzLabelPoint[below left](B){$B$}
\tkzLabelPoint[above right](C){$\varGamma $}
\tkzLabelPoint[above](O){$O$}
\tkzLabelPoint[below right](D){$\varDelta  $}
\tkzLabelPoint[left](M){$M$}

\tkzMarkRightAngle(C,A,B)
\tkzMarkRightAngle(A,D,C)
\tkzMarkRightAngle(O,D,M)
\end{tikzpicture}
\end{center}
\end{alist}

\vspace{4mm}

% --- Άσκηση 34330 (34330.tex) ---
\Askhsh{34330}{4}
\begin{ylh}{1}
    \item 3.1. Είδη και στοιχεία τριγώνων
    \item 3.2. 1ο Κριτήριο ισότητας τριγώνων
    \item 3.7. Κύκλος - Μεσοκάθετος - Διχοτόμος
    \item 3.10. Σχέση εξωτερικής και απέναντι γωνίας
    \item 5.9. Μια ιδιότητα του ορθογώνιου τριγώνου.
\end{ylh}
% ===================================================================
%  Αρχείο: 34330.tex (Fragment)
%  Αυτόματη μετατροπή μέσω doc2latex.py
% ===================================================================



Δίνονται τα ορθογώνια τρίγωνα $AB\varGamma $ ($\hat{A}$ = 90\degree ) και Δ$B\varGamma $ ($\hat{\varDelta  }$ = 90\degree ) με τις κορυφές τους Α και $\varDelta$ εκατέρωθεν της $B\varGamma $ και το μέσο $M$ της $B\varGamma $.

Να αποδείξετε ότι:

\begin{alist}
\item το τρίγωνο $AM\varDelta  $ είναι ισοσκελές, \monades{9}

\item Α$\hat{Μ}$Δ = 2Α$\hat{\varGamma }$Δ, \monades{9}

\item τα Α, Β, Δ και $\varGamma$ είναι σημεία ενός κύκλου, τον οποίο και να κατασκευάσετε.

\monades{7}


\begin{center}
\begin{tikzpicture}[scale=1]
\tkzDefPoint(0,0){B}
\tkzDefPoint(4,0){C}
\tkzDefMidPoint(B,C) \tkzGetPoint{M}

% The circle has diameter $BC$.
\tkzDrawCircle(M,B)
% A and D are on the circle, opposite sides of $BC$.
\tkzDefPointBy[rotation=center M angle 110](C) \tkzGetPoint{A}
\tkzDefPointBy[rotation=center M angle -50](C) \tkzGetPoint{D}

\draw[l3] (A)--(B)--(C)--cycle;
\draw[l3] (D)--(B)--(C)--cycle;
\draw[l3] (A)--(M);
\draw[l3] (D)--(M);
\draw[l3] (A)--(D);
\draw[l3] (A)--(C);
\draw[l3] (D)--(C);

\tkzDrawPoints(A,B,C,D,M)
\tkzLabelPoint[above left](A){$A$}
\tkzLabelPoint[left](B){$B$}
\tkzLabelPoint[right](C){$\varGamma $}
\tkzLabelPoint[below right](D){$\varDelta  $}
\tkzLabelPoint[below](M){$M$}

\tkzMarkRightAngle(B,A,C)
\tkzMarkRightAngle(B,D,C)
\end{tikzpicture}
\end{center}
\end{alist}

\vspace{4mm}

% --- Άσκηση 34331 (34331.tex) ---
\Askhsh{34331}{4}
\begin{ylh}{1}
    \item 3.1. Είδη και στοιχεία τριγώνων
    \item 3.4. 3ο Κριτήριο ισότητας τριγώνων
    \item 5.6. Εφαρμογές στα τρίγωνα
    \item 5.9. Μια ιδιότητα του ορθογώνιου τριγώνου
    \item 5.10. Τραπέζιο
    \item 5.11. Ισοσκελές τραπέζιο.
\end{ylh}
% ===================================================================
%  Αρχείο: 34331.tex (Fragment)
%  Αυτόματη μετατροπή μέσω doc2latex.py
% ===================================================================



Σε οξυγώνιο τρίγωνο $AB\varGamma $ ($AB$ < $A\varGamma $) φέρουμε το ύψος $A\varDelta  $. Έστω Κ,Λ,Μ τα μέσα των $AB$, $A\varGamma $, $B\varGamma $ αντίστοιχα.

Να αποδείξετε ότι:

\begin{alist}
\item $K\varLambda $ // $B\varGamma $ \monades{5}

\item

\end{alist}
\begin{alist}
\item
  $M\varLambda $ = $K\varDelta  $ \monades{6}
\item
  Το τετράπλευρο $K\varLambda  M\varDelta  $ είναι ισοσκελές τραπέζιο. \monades{8}

\item Οι γωνίες Κ$\hat{\varDelta  }\text{Λ}$ και $\text{Κ}\hat{Μ}\Lambda$ είναι ίσες. \monades{6}
\end{alist}

\vspace{4mm}

% --- Άσκηση 34332 (34332.tex) ---
\Askhsh{34332}{2}
\begin{ylh}{1}
    \item 3.4. 3ο Κριτήριο ισότητας τριγώνων
    \item 4.2. Τέμνουσα δύο ευθειών - Ευκλείδειο αίτημα
    \item 5.6. Εφαρμογές στα τρίγωνα.
\end{ylh}
% ===================================================================
%  Αρχείο: 34332.tex (Fragment)
%  Αυτόματη μετατροπή μέσω doc2latex.py
% ===================================================================



Στο παρακάτω σχήμα δίνονται τα ισοσκελή τρίγωνα $AB\varGamma $ και $E\varGamma \varDelta  $ με $ΑΒ = Α\varGamma  = Ε\varGamma  = Ε\varDelta  $, όπου $\varDelta$ είναι το μέσο της $A\varGamma $ και $B\varGamma  = \frac{AB}{2}$. Έστω $Z$ το σημείο στο οποίο η προέκταση της $E\varDelta  $ προς το $\varDelta$ τέμνει την $AB$.

Να αποδείξετε ότι:

\begin{alist}
\item τα τρίγωνα $AB\varGamma $ και $\varGamma \varDelta  E$ είναι ίσα, \monades{12}

\item το σημείο $Z$ είναι το μέσο της $AB$. \monades{13}


\begin{center}
\begin{tikzpicture}[scale=1]
% $BC$ = $AB$/2. Let's make $BC$=2, $AB$=$AC$=4.
\tkzDefPoint(0,0){B}
\tkzDefPoint(2,0){C}
\tkzDefTriangle[two angles = 75 and 75](B,C) \tkzGetPoint{A} % wait, if $AB$=4, $BC$=2. Base is 2, legs 4.
% actually we can construct it exactly
\tkzInterCC[R](B,4)(C,4) \tkzGetPoints{A}{A2}

\tkzDefMidPoint(A,C) \tkzGetPoint{D}

% Triangle $CDE$ is isosceles with $EC$=$ED$=$AC$=4.
% E is some point such that $EC$=$ED$=4.
% Wait: "$E\varGamma \varDelta  $ με $AB$=$AC$=$E\varGamma $=$E\varDelta  $". So $EC$=$ED$=4. D is midpoint, so $CD$=2.
% We have triangle $CDE$ with base $CD$=2, legs $EC$=$ED$=4.
% Since $AC$=4 and $CD$=2, this is similar to $ABC$!
% Actually, it is "τα ισοσκελή τρίγωνα $AB\varGamma $ και $E\varGamma \varDelta  $".
% Where does E lie? Usually on the opposite side of $AC$ or something. Let's place E.
\tkzInterCC[R](C,4)(D,4) \tkzGetPoints{E_alt}{E} % E to the right

% Extension of $ED$ towards D intersects $AB$ at Z.
\tkzInterLL(E,D)(A,B) \tkzGetPoint{Z}

\draw[l3] (A)--(B)--(C)--cycle;
\draw[l3] (E)--(C)--(D)--cycle;
\draw[l3] (E)--(Z);
\draw[l3] (A)--(B); % $AZB$ line

\tkzDrawPoints(A,B,C,D,E,Z)
\tkzLabelPoint[above](A){$A$}
\tkzLabelPoint[left](B){$B$}
\tkzLabelPoint[below](C){$\varGamma $}
\tkzLabelPoint[right](D){$\varDelta  $}
\tkzLabelPoint[right](E){$E$}
\tkzLabelPoint[left](Z){$Z$}

\tkzMarkSegments[mark=s|](A,B A,C E,C E,D)
\tkzMarkSegments[mark=s||](A,D D,C B,C) % since $AD$=$DC$=2, $BC$=2
\end{tikzpicture}
\end{center}
\end{alist}

\vspace{4mm}

% --- Άσκηση 34333 (34333.tex) ---
\Askhsh{34333}{4}
\begin{ylh}{1}
    \item 3.1. Είδη και στοιχεία τριγώνων
    \item 3.2. 1ο Κριτήριο ισότητας τριγώνων
    \item 4.2. Τέμνουσα δύο ευθειών - Ευκλείδειο αίτημα
    \item 5.6. Εφαρμογές στα τρίγωνα
    \item 5.10. Τραπέζιο
    \item 5.11. Ισοσκελές τραπέζιο.
\end{ylh}
% ===================================================================
%  Αρχείο: 34333.tex (Fragment)
%  Αυτόματη μετατροπή μέσω doc2latex.py
% ===================================================================



Θεωρούμε τρίγωνο $AB\varGamma $ με $B\varGamma $$\  = \ $2ΑΒ. Έστω $\varDelta$ το μέσο της πλευράς $B\varGamma $ και $E$ το μέσο του τμήματος $B\varDelta  $. Από το σημείο $\varDelta  $ φέρουμε ευθεία παράλληλη προς την $A\varGamma $, η οποία τέμνει την πλευρά $AB$ στο σημείο $Z$.

Να αποδείξετε ότι:

\begin{alist}
\item τα τρίγωνα Α$BE$ και $BZ\varDelta  $ είναι ίσα, \monades{7}

\item η $A\varDelta  $ είναι διχοτόμος της γωνίας Ε$\hat{A}$Γ, \monades{9}

\item το τετράπλευρο $A\varDelta  EZ$ είναι ισοσκελές τραπέζιο. \monades{9}
\end{alist}

\vspace{4mm}

% --- Άσκηση 34334 (34334.tex) ---
\Askhsh{34334}{4}
\begin{ylh}{1}
    \item 3.1. Είδη και στοιχεία τριγώνων
    \item 3.2. 1ο Κριτήριο ισότητας τριγώνων
    \item 3.11. Ανισοτικές σχέσεις πλευρών και γωνιών
    \item 5.2. Παραλληλόγραμμα
    \item 5.6. Εφαρμογές στα τρίγωνα
    \item 5.9. Μια ιδιότητα του ορθογώνιου τριγώνου.
\end{ylh}
% ===================================================================
%  Αρχείο: 34334.tex (Fragment)
%  Αυτόματη μετατροπή μέσω doc2latex.py
% ===================================================================



Σε οξυγώνιο τρίγωνο $AB\varGamma $ ($\hat{\text{Β}} < \hat{\varGamma}$) θεωρούμε τα μέσα Δ, Ε, Ζ των πλευρών $AB$, $A\varGamma $, $B\varGamma $ αντίστοιχα. Έστω Η η προβολή της κορυφής $\varGamma$ πάνω στην πλευρά $AB$.

Να αποδείξετε ότι:

\begin{alist}
\item ΗΕ$\  = \ $$E\varGamma $ και ΗΖ$\  = \ $ΖΓ, \monades{8}

\item το τετράπλευρο $\varDelta  E\varGamma Z$ είναι παραλληλόγραμμο, \monades{8}

\item Ζ$\hat{\text{Δ}}$Ε$\  = \ $Ζ$\hat{\text{H}}$Ε. \monades{9}
\end{alist}

\vspace{4mm}

% --- Άσκηση 34335 (34335.tex) ---
\Askhsh{34335}{4}
\begin{ylh}{1}
    \item 3.1. Είδη και στοιχεία τριγώνων
    \item 3.4. 3ο Κριτήριο ισότητας τριγώνων
    \item 3.10. Σχέση εξωτερικής και απέναντι γωνίας
    \item 3.14. Σχετικές θέσεις ευθείας και κύκλου
    \item 3.15. Εφαπτόμενα τμήματα
    \item 4.2. Τέμνουσα δύο ευθειών - Ευκλείδειο αίτημα.
\end{ylh}
% ===================================================================
%  Αρχείο: 34335.tex (Fragment)
%  Αυτόματη μετατροπή μέσω doc2latex.py
% ===================================================================



Δίνεται κύκλος (Ο,R) και μία ευθεία x΄x η οποία έχει μοναδικό κοινό σημείο με τον κύκλο το σημείο $A$. Θεωρούμε τυχαίο σημείο $M$ της ημιευθείας Αx. Aν για κάποιο σημείο $B$ του κύκλου ισχύει η σχέση $MA$ = $MB$, να αποδείξετε ότι:

\begin{alist}
\item το $MB$ είναι εφαπτόμενο τμήμα του κύκλου (Ο,R) \monades{9}

\item η διχοτόμος της γωνίας ΒΜx είναι κάθετη στη $MO$, \monades{9}

\item το ευθύγραμμο τμήμα $OB$ τέμνει τη διχοτόμο της γωνίας ΒΜx. \monades{7}
\end{alist}

\vspace{4mm}

% --- Άσκηση 34336 (34336.tex) ---
\Askhsh{34336}{4}
\begin{ylh}{1}
    \item 3.1. Είδη και στοιχεία τριγώνων
    \item 3.6. Κριτήρια ισότητας ορθογώνιων τριγώνων
    \item 4.6. Άθροισμα γωνιών τριγώνου
    \item 5.3. Ορθογώνιο
    \item 5.4. Ρόμβος
    \item 5.5. Τετράγωνο.
\end{ylh}
% ===================================================================
%  Αρχείο: 34336.tex (Fragment)
%  Αυτόματη μετατροπή μέσω doc2latex.py
% ===================================================================



Δίνεται το τετράγωνο $AB\varGamma \varDelta  $. Προεκτείνουμε την πλευρά $AB$ προς το Β κατά τμήμα $BZ$ και την πλευρά $B\varGamma $ προς το $\varGamma$ κατά τμήμα $\varGamma M$ = $AZ$. Θεωρούμε σημείο $E$ τέτοιο, ώστε το τετράπλευρο $\varDelta  MEZ$ να είναι παραλληλόγραμμο.

Να αποδείξετε ότι:

\begin{alist}
\item τα τρίγωνα Α$\varDelta   Z$ και $\varGamma \varDelta  M$ είναι ίσα και οι γωνίες Α$\hat{\varDelta  }$Ζ και Γ$\hat{\varDelta  }$Μ είναι ίσες, \monades{9}

\item το τετράπλευρο $\varDelta  MEZ$ είναι τετράγωνο, \monades{9}

\item οι γωνίες Β$\hat{Ζ}$Ε και Ε$\hat{Μ}$Β είναι παραπληρωματικές. \monades{7}


\begin{center}
\begin{tikzpicture}[scale=0.8]
\tkzDefPoint(0,4){A}
\tkzDefPoint(0,0){B}
\tkzDefPoint(4,0){C}
\tkzDefPoint(4,4){D}

% Extend $AB$ towards B by $BZ$
\tkzDefPointOnLine[pos=1.4](A,B) \tkzGetPoint{Z}
\tkzCalcLength(B,Z) \tkzGetLength{rBZ}
\tkzCalcLength(A,Z) \tkzGetLength{rAZ}

% Extend $BC$ towards C by $CM$ = $AZ$
% $AZ$ = 5.6 (A to B = 4, $BZ$ = 1.6), $BC$ = 4, so M is at pos = (4+5.6)/4 = 2.4 from B
\tkzDefPointOnLine[pos=2.4](B,C) \tkzGetPoint{M}
% E is such that $DMEZ$ is a parallelogram. 
% So Vector($ZE$) = Vector($MD$). E = Z + D - M
\tkzDefPointBy[translation=from M to D](Z) \tkzGetPoint{E}

\draw[l3] (A)--(B)--(C)--(D)--cycle;
\draw[l3] (B)--(Z);
\draw[l3] (C)--(M);
\draw[l3] (D)--(M)--(E)--(Z)--cycle; % $DMEZ$

\tkzDrawPoints(A,B,C,D,Z,M,E)
\tkzLabelPoint[above left](A){$A$}
\tkzLabelPoint[below left](B){$B$}
\tkzLabelPoint[below](C){$\varGamma $}
\tkzLabelPoint[above](D){$\varDelta  $}
\tkzLabelPoint[below](Z){$Z$}
\tkzLabelPoint[right](M){$M$}
\tkzLabelPoint[below right](E){$E$}

\tkzMarkRightAngle(A,B,C)
\tkzMarkSegments[mark=s||](A,Z C,M)
\end{tikzpicture}
\end{center}
\end{alist}

\vspace{4mm}

% --- Άσκηση 34385 (34385.tex) ---
\Askhsh{34385}{2}
\begin{ylh}{1}
    \item 3.1. Είδη και στοιχεία τριγώνων
    \item 3.2. 1ο Κριτήριο ισότητας τριγώνων
    \item 3.11. Ανισοτικές σχέσεις πλευρών και γωνιών
    \item 5.10. Τραπέζιο
    \item 5.11. Ισοσκελές τραπέζιο.
\end{ylh}
% ===================================================================
%  Αρχείο: 34385.tex (Fragment)
%  Αυτόματη μετατροπή μέσω doc2latex.py
% ===================================================================



Θεωρούμε ισοσκελές τρίγωνο $AB\varGamma $ ($AB=A\varGamma $). Στο μέσο $\varDelta  $ της πλευράς $AB$ φέρουμε κάθετη ευθεία που τέμνει την $A\varGamma $ στο Ε. Από το $E$ φέρουμε ευθεία παράλληλη στη βάση $B\varGamma $ που τέμνει την $AB$ στο Ζ.

Να αποδείξετε ότι:

\begin{alist}
\item $AE$= $BE$, \monades{15}

\item το τετράπλευρο $B\varGamma $ΕΖ είναι ισοσκελές τραπέζιο. \monades{10}


\begin{center}
\begin{tikzpicture}[scale=1]
\tkzDefPoint(0,0){B}
\tkzDefPoint(4,0){C}
\tkzDefTriangle[isosceles right](B,C) \tkzGetPoint{A} % Wait, just any isosceles.
% Actually let's just make it tall.
\tkzDefPoint(2,4){A}

\tkzDefMidPoint(A,B) \tkzGetPoint{D}

\tkzDefLine[orthogonal=through D](A,B) \tkzGetPoint{perp_dir}
\tkzInterLL(D,perp_dir)(A,C) \tkzGetPoint{E}

\tkzDefLine[parallel=through E](B,C) \tkzGetPoint{par_dir}
\tkzInterLL(E,par_dir)(A,B) \tkzGetPoint{Z}

\draw[l3] (A)--(B)--(C)--cycle;
\draw[l3] (D)--(E);
\draw[l3] (E)--(Z);
\draw[l3, dashed] (B)--(E);

\tkzDrawPoints(A,B,C,D,E,Z)
\tkzLabelPoint[above](A){$A$}
\tkzLabelPoint[below left](B){$B$}
\tkzLabelPoint[below right](C){$\varGamma $}
\tkzLabelPoint[above left](D){$\varDelta  $}
\tkzLabelPoint[above right](E){$E$}
\tkzLabelPoint[left](Z){$Z$}

\tkzMarkRightAngle(A,D,E)
\tkzMarkSegments[mark=s|](A,D D,B A,E E,B)
\end{tikzpicture}
\end{center}
\end{alist}

\vspace{4mm}

% --- Άσκηση 34386 (34386.tex) ---
\Askhsh{34386}{2}
\begin{ylh}{1}
    \item 3.1. Είδη και στοιχεία τριγώνων
    \item 3.11. Ανισοτικές σχέσεις πλευρών και γωνιών
    \item 5.2. Παραλληλόγραμμα.
\end{ylh}
% ===================================================================
%  Αρχείο: 34386.tex (Fragment)
%  Αυτόματη μετατροπή μέσω doc2latex.py
% ===================================================================



Δίνεται παραλληλόγραμμο $AB\varGamma \varDelta  $ με $AB$=2$B\varGamma $. Προεκτείνουμε την πλευρά $A\varDelta  $ (προς το μέρος του Δ) κατά τμήμα $\varDelta  E$=$A\varDelta  $ και φέρουμε την $BE$ που τέμνει τη $\varDelta  \varGamma $ στο σημείο $H$.

Να αποδείξετε ότι:

\begin{alist}
\item το τρίγωνο $BAE$ είναι ισοσκελές, \monades{7}

\item το $\varDelta  E\varGamma B$ είναι παραλληλόγραμμο, \monades{9}

\item η $AH$ είναι διάμεσος του $BAE$ τριγώνου. \monades{9}
\end{alist}

\vspace{4mm}

% --- Άσκηση 34387 (34387.tex) ---
\Askhsh{34387}{2}
\begin{ylh}{1}
    \item 3.2. 1ο Κριτήριο ισότητας τριγώνων
    \item 3.3. 2ο Κριτήριο ισότητας τριγώνων
    \item 3.6. Κριτήρια ισότητας ορθογώνιων τριγώνων.
\end{ylh}
% ===================================================================
%  Αρχείο: 34387.tex (Fragment)
%  Αυτόματη μετατροπή μέσω doc2latex.py
% ===================================================================



Δίνεται ισοσκελές τρίγωνο $AB\varGamma $ ($AB=A\varGamma $) και οι διχοτόμοι του $B\varDelta  $ και $\varGamma  E$ των γωνιών Β και $\varGamma$ αντίστοιχα.

Να αποδείξετε ότι:

\begin{alist}
\item Να αποδείξετε ότι τα τρίγωνα $B\varGamma $Δ και Γ$BE$ είναι ίσα. \monades{13}

\item Έστω $EH$ και $\varDelta   Z$ οι κάθετες από τα σημεία $E$ και $\varDelta$ αντίστοιχα στη $B\varGamma $. Να αποδείξετε ότι

$EH$ = $\varDelta   Z$. \monades{12}
\end{alist}

\vspace{4mm}

% --- Άσκηση 34388 (34388.tex) ---
\Askhsh{34388}{2}
\begin{ylh}{1}
    \item 3.1. Είδη και στοιχεία τριγώνων
    \item 4.2. Τέμνουσα δύο ευθειών - Ευκλείδειο αίτημα
    \item 5.2. Παραλληλόγραμμα.
\end{ylh}
% ===================================================================
%  Αρχείο: 34388.tex (Fragment)
%  Αυτόματη μετατροπή μέσω doc2latex.py
% ===================================================================



Δίνεται τρίγωνο $AB\varGamma $, στο οποίο φέρουμε τις διαμέσους του $BM$ και $\varGamma N$. Προεκτείνουμε την $BM$ (προς το Μ) κατά τμήμα $M\varDelta  $=$BM$ και την $\varGamma N$ (προς το Ν) κατά τμήμα $NE$=$\varGamma N$.

\begin{alist}
\item Να αποδείξετε ότι $A\varDelta  $//$B\varGamma $ και $AE$//$B\varGamma $. \monades{13}

\item Είναι τα σημεία $E$, Α και $\varDelta$ συνευθειακά; Να αιτιολογήσετε την απάντησή σας.

\monades{12}
\end{alist}

\vspace{4mm}

% --- Άσκηση 34389 (34389.tex) ---
\Askhsh{34389}{2}
\begin{ylh}{1}
    \item 3.6. Κριτήρια ισότητας ορθογώνιων τριγώνων
    \item 4.2. Τέμνουσα δύο ευθειών - Ευκλείδειο αίτημα
    \item 5.2. Παραλληλόγραμμα.
\end{ylh}
% ===================================================================
%  Αρχείο: 34389.tex (Fragment)
%  Αυτόματη μετατροπή μέσω doc2latex.py
% ===================================================================



Δίνεται παραλληλόγραμμο $AB\varGamma \varDelta  $ και η διαγώνιός του $B\varDelta  $. Από τις κορυφές Α και $\varGamma$ φέρουμε τις κάθετες $AE$ και $\varGamma Z$ στη $B\varDelta  $, που την τέμνουν στα σημεία $E$ και $Z$ αντίστοιχα.

Να αποδείξετε ότι:

\begin{alist}
\item τα τρίγωνα $A\varDelta  E$ και $\varGamma BZ$ είναι ίσα, \monades{10}

\item το τετράπλευρο $AE\varGamma Z$ είναι παραλληλόγραμμο. \monades{15}
\end{alist}

\vspace{4mm}

% --- Άσκηση 34391 (34391.tex) ---
\Askhsh{34391}{2}
\begin{ylh}{1}
    \item 4.2. Τέμνουσα δύο ευθειών - Ευκλείδειο αίτημα
    \item 5.2. Παραλληλόγραμμα.
\end{ylh}
% ===================================================================
%  Αρχείο: 34391.tex (Fragment)
%  Αυτόματη μετατροπή μέσω doc2latex.py
% ===================================================================



Δίνεται τρίγωνο $AB\varGamma $. Από το μέσο $M$ της πλευράς $B\varGamma $ φέρουμε ευθύγραμμο τμήμα $M\varDelta  $ ίσο και παράλληλο προς την πλευρά $BA$ και ευθύγραμμο τμήμα $ME$ ίσο και παράλληλο προς την πλευρά $\varGamma  A$.

Να αποδείξετε ότι:

\begin{alist}
\item $\varDelta  A$=$AE$, \monades{8}

\item τα σημεία $\varDelta  $, Α και $E$ βρίσκονται στην ίδια ευθεία, \monades{9}

\item $\varDelta  E$=$B\varGamma $. \monades{8}


\begin{center}
\begin{tikzpicture}[scale=1]
\tkzDefPoint(0,0){B}
\tkzDefPoint(4,0){C}
\tkzDefPoint(1,3){A}
\tkzDefMidPoint(B,C) \tkzGetPoint{M}

% $MD$ equal and parallel to $BA$. So Vector($MD$) = Vector($BA$). D = M + A - B
\tkzDefPointBy[translation=from B to A](M) \tkzGetPoint{D}

% $ME$ equal and parallel to $CA$. So Vector($ME$) = Vector($CA$). E = M + A - C
\tkzDefPointBy[translation=from C to A](M) \tkzGetPoint{E}

\draw[l3] (A)--(B)--(C)--cycle;
\draw[l3] (M)--(D);
\draw[l3] (M)--(E);
\draw[l3, dashed] (D)--(A)--(E);
\draw[l3, dashed] (D)--(E);
\draw[l3, dashed] (M)--(A);

\tkzDrawPoints(A,B,C,M,D,E)
\tkzLabelPoint[above](A){$A$}
\tkzLabelPoint[below left](B){$B$}
\tkzLabelPoint[below right](C){$\varGamma $}
\tkzLabelPoint[below](M){$M$}
\tkzLabelPoint[above left](D){$\varDelta  $}
\tkzLabelPoint[above right](E){$E$}

\tkzMarkSegments[mark=s|](B,M M,C)
\tkzMarkSegments[mark=s||](A,B M,D)
\tkzMarkSegments[mark=s|||](A,C M,E)
\end{tikzpicture}
\end{center}
\end{alist}

\vspace{4mm}

% --- Άσκηση 34392 (34392.tex) ---
\Askhsh{34392}{2}
\begin{ylh}{1}
    \item 3.11. Ανισοτικές σχέσεις πλευρών και γωνιών
    \item 5.6. Εφαρμογές στα τρίγωνα
    \item 5.10. Τραπέζιο
    \item 5.11. Ισοσκελές τραπέζιο.
\end{ylh}
% ===================================================================
%  Αρχείο: 34392.tex (Fragment)
%  Αυτόματη μετατροπή μέσω doc2latex.py
% ===================================================================



Δίνεται τρίγωνο $AB\varGamma $ και $\varDelta$ το μέσο της πλευράς $AB$. Από το $\varDelta$ διέρχεται μια τυχαία ευθεία (ε) που τέμνει την πλευρά $A\varGamma $ σε εσωτερικό της σημείο $E$. Η ευθεία (ε) χωρίζει το τρίγωνο $AB\varGamma $ σε ένα τρίγωνο $A\varDelta  E$ και σε ένα τετράπλευρο $B\varDelta  $ΕΓ .

\begin{alist}
\item Ποια πρέπει να είναι η θέση του σημείου Ε, ώστε το τετράπλευρο $B\varDelta  $ΕΓ να είναι τραπέζιο; Να αιτιολογήσετε την απάντησή σας. \monades{12}

\item Ποιο πρέπει να είναι το είδος του $AB\varGamma $ τριγώνου, ώστε το τραπέζιο του ερωτήματος (α) να είναι ισοσκελές τραπέζιο; Να αιτιολογήσετε την απάντησή σας. \monades{13}


\begin{center}
\begin{tikzpicture}[scale=1]
\tkzDefPoint(0,0){B}
\tkzDefPoint(5,0){C}
\tkzDefPoint(2,4){A}

\tkzDefMidPoint(A,B) \tkzGetPoint{D}

% E is a random point on $AC$. Let's make it such that $DE$ is parallel to $BC$ for the trapezoid case (as it's the answer probably)
% Or just a random point to show the generic case.
\tkzDefPointOnLine[pos=0.6](A,C) \tkzGetPoint{E}

\draw[l3] (A)--(B)--(C)--cycle;
\draw[l3] (D)--(E);

\tkzDrawPoints(A,B,C,D,E)
\tkzLabelPoint[above](A){$A$}
\tkzLabelPoint[below left](B){$B$}
\tkzLabelPoint[below right](C){$\varGamma $}
\tkzLabelPoint[above left](D){$\varDelta  $}
\tkzLabelPoint[above right](E){$E$}

\tkzMarkSegments[mark=s|](A,D D,B)
\end{tikzpicture}
\end{center}
\end{alist}

\vspace{4mm}

% --- Άσκηση 34393 (34393.tex) ---
\Askhsh{34393}{2}
\begin{ylh}{1}
    \item 3.11. Ανισοτικές σχέσεις πλευρών και γωνιών
    \item 4.2. Τέμνουσα δύο ευθειών - Ευκλείδειο αίτημα
    \item 5.1 Εισαγωγή
    \item 5.2. Παραλληλόγραμμα
    \item 5.6. Εφαρμογές στα τρίγωνα.
\end{ylh}
% ===================================================================
%  Αρχείο: 34393.tex (Fragment)
%  Αυτόματη μετατροπή μέσω doc2latex.py
% ===================================================================



Σε παραλληλόγραμμο $AB\varGamma \varDelta  $, προεκτείνουμε την πλευρά $\varDelta  A$ (προς το Α) κατά τμήμα $AH$=$\varDelta  A$. Φέρουμε τη διχοτόμο της γωνίας $\hat{\varDelta  }$, η οποία τέμνει την $AB$ στο σημείο $Z$.

Να αποδείξετε ότι:

\begin{alist}
\item το τρίγωνο Α$\varDelta   Z$ είναι ισοσκελές, \monades{12}

\item το τρίγωνο $\varDelta   Z$Η είναι ορθογώνιο με ορθή τη γωνία $\hat{Ζ}$. \monades{13}

\begin{center}
\begin{tikzpicture}[scale=0.8]
\tkzDefPoint(0,0){B}
\tkzDefPoint(5,0){C}
\tkzDefPoint(1.5,3){A}

% D is C + A - B
\tkzDefPointBy[translation=from B to A](C) \tkzGetPoint{D}

% Extend $DA$ towards A by $AH$=$DA$. So H = A + A - D
\tkzDefPointBy[translation=from D to A](A) \tkzGetPoint{H}

% Bisector of angle D (angle $ADC$)
\tkzDefLine[bisector](A,D,C) \tkzGetPoint{bis_dir}
\tkzInterLL(D,bis_dir)(A,B) \tkzGetPoint{Z}

\draw[l3] (A)--(B)--(C)--(D)--cycle;
\draw[l3] (A)--(H);
\draw[l3] (D)--(Z);
\draw[l3] (H)--(Z);

\tkzDrawPoints(A,B,C,D,H,Z)
\tkzLabelPoint[above left](A){$A$}
\tkzLabelPoint[below left](B){$B$}
\tkzLabelPoint[below right](C){$\varGamma $}
\tkzLabelPoint[above right](D){$\varDelta  $}
\tkzLabelPoint[left](H){$H$}
\tkzLabelPoint[left](Z){$Z$}

\tkzMarkSegments[mark=s|](A,D A,H)
\tkzMarkAngles[size=0.6,mark=none](A,D,Z Z,D,C)
\tkzMarkRightAngle(H,Z,D)
\end{tikzpicture}
\end{center}
\end{alist}

\vspace{4mm}

% --- Άσκηση 34394 (34394.tex) ---
\Askhsh{34394}{2}
\begin{ylh}{1}
    \item 3.11. Ανισοτικές σχέσεις πλευρών και γωνιών
    \item 4.2. Τέμνουσα δύο ευθειών - Ευκλείδειο αίτημα
    \item 5.1 Εισαγωγή
    \item 5.2. Παραλληλόγραμμα.
\end{ylh}
% ===================================================================
%  Αρχείο: 34394.tex (Fragment)
%  Αυτόματη μετατροπή μέσω doc2latex.py
% ===================================================================



Δίνεται $AB\varGamma \varDelta  $ παραλληλόγραμμο με $AB$=2ΑΔ. Φέρουμε τη διχοτόμο της γωνίας $\hat{\varDelta  }$ του παραλληλογράμμου, η οποία τέμνει την $AB$ στο Ε.

\begin{center}
\begin{tikzpicture}[scale=1]
\tkzDefPoint(0,0){A}
\tkzDefPoint(6,0){B}
\tkzDefPoint(1.5,2.5){D}
\tkzDefPointBy[translation=from A to B](D) \tkzGetPoint{C}

% Bisector of angle D
\tkzDefLine[bisector](A,D,C) \tkzGetPoint{bisD}
\tkzInterLL(D,bisD)(A,B) \tkzGetPoint{E}

\draw[l3] (A)--(B)--(C)--(D)--cycle;
\draw[l3] (D)--(E);

\tkzDrawPoints(A,B,C,D,E)
\tkzLabelPoint[below left](A){$A$}
\tkzLabelPoint[below right](B){$B$}
\tkzLabelPoint[above right](C){$\varGamma $}
\tkzLabelPoint[above left](D){$\varDelta  $}
\tkzLabelPoint[below](E){$E$}

\tkzMarkSegments[mark=s|](A,D D,E A,E) % $ADE$ is isosceles
\tkzMarkSegments[mark=s|](E,B) % E is midpoint since $AB$=2AD
\tkzMarkAngles[size=0.6,mark=none](A,D,E E,D,C)
\end{tikzpicture}
\end{center}

\begin{alist}
\item Να αποδείξετε ότι το τρίγωνο $A\varDelta  E$ είναι ισοσκελές. \monades{12}

\item Είναι το σημείο $E$ μέσο της πλευράς $AB$; Να αιτιολογήσετε την απάντησή σας.

\monades{13}
\end{alist}

\vspace{4mm}

% --- Άσκηση 34395 (34395.tex) ---
\Askhsh{34395}{2}
\begin{ylh}{1}
    \item 3.2. 1ο Κριτήριο ισότητας τριγώνων
    \item 5.1 Εισαγωγή
    \item 5.2. Παραλληλόγραμμα.
\end{ylh}
% ===================================================================
%  Αρχείο: 34395.tex (Fragment)
%  Αυτόματη μετατροπή μέσω doc2latex.py
% ===================================================================



Δίνεται παραλληλόγραμμο $AB\varGamma \varDelta  $ και Ο το σημείο τομής των διαγωνίων του. Θεωρούμε σημείο $E$ του τμήματος $AO$ και σημείο $Z$ του τμήματος $O\varGamma $, ώστε $OE$=$OZ$.

Να αποδείξετε ότι:

\begin{alist}
\item $\varDelta  E$=$BZ$, \monades{12}

\item το $\varDelta   EBZ$ είναι παραλληλόγραμμο. \monades{13}
\end{alist}

\vspace{4mm}

% --- Άσκηση 34396 (34396.tex) ---
\Askhsh{34396}{2}
\begin{ylh}{1}
    \item 3.6. Κριτήρια ισότητας ορθογώνιων τριγώνων
    \item 3.11. Ανισοτικές σχέσεις πλευρών και γωνιών.
\end{ylh}
% ===================================================================
%  Αρχείο: 34396.tex (Fragment)
%  Αυτόματη μετατροπή μέσω doc2latex.py
% ===================================================================



Σε ορθογώνιο τρίγωνο $AB\varGamma $ (Α=90˚), η διχοτόμος τη γωνίας $\hat{\varGamma }$ τέμνει την πλευρά $AB$ σε σημείο $\varDelta  $. Από το $\varDelta$ φέρουμε προς την πλευρά $B\varGamma $ μια κάθετη ευθεία, η οποία τέμνει την πλευρά $B\varGamma $ σε σημείο $E$.

\begin{center}
\begin{tikzpicture}[scale=1]
\tkzDefPoint(0,0){A}
\tkzDefPoint(0,4){B}
\tkzDefPoint(6,0){C}
\draw[l3] (A)--(B)--(C)--cycle;

\tkzDefLine[bisector](B,C,A) \tkzGetPoint{bisC}
\tkzInterLL(C,bisC)(A,B) \tkzGetPoint{D}

\tkzDefPointBy[projection=onto B--C](D) \tkzGetPoint{E}

\draw[l3] (C)--(D);
\draw[l3] (D)--(E);

\tkzDrawPoints(A,B,C,D,E)
\tkzLabelPoint[below left](A){$A$}
\tkzLabelPoint[above left](B){$B$}
\tkzLabelPoint[below right](C){$\varGamma $}
\tkzLabelPoint[left](D){$\varDelta  $}
\tkzLabelPoint[above right](E){$E$}

\tkzMarkRightAngle(C,A,B)
\tkzMarkRightAngle(D,E,C)
\tkzMarkAngles[size=1,mark=none](B,C,D D,C,A)
\tkzMarkSegments[mark=s|](A,D D,E)
\end{tikzpicture}
\end{center}

Να αποδείξετε ότι:

\begin{alist}
\item $A\varDelta  $=$\varDelta  E$, \monades{13}

\item $A\varDelta  $\textless $\varDelta   B$. \monades{12}
\end{alist}

\vspace{4mm}

% --- Άσκηση 34397 (34397.tex) ---
\Askhsh{34397}{2}
\begin{ylh}{1}
    \item 3.6. Κριτήρια ισότητας ορθογώνιων τριγώνων
    \item 4.6. Άθροισμα γωνιών τριγώνου.
\end{ylh}
% ===================================================================
%  Αρχείο: 34397.tex (Fragment)
%  Αυτόματη μετατροπή μέσω doc2latex.py
% ===================================================================



Δίνεται ορθογώνιο τρίγωνο $AB\varGamma $ ($\hat{A} = 90\degree$). Η διχοτόμος της γωνίας $\hat{B}$ τέμνει την πλευρά $A\varGamma $ στο σημείο $\varDelta  $. Φέρουμε τμήμα $\varDelta  E$ κάθετο στην πλευρά $B\varGamma $.

\begin{center}
\begin{tikzpicture}[scale=1]
\tkzDefPoint(0,0){A}
\tkzDefPoint(0,4.5){B}
\tkzDefPoint(6,0){C}
\draw[l3] (A)--(B)--(C)--cycle;

\tkzDefLine[bisector](A,B,C) \tkzGetPoint{bisB}
\tkzInterLL(B,bisB)(A,C) \tkzGetPoint{D}

\tkzDefPointBy[projection=onto B--C](D) \tkzGetPoint{E}

\draw[l3] (B)--(D);
\draw[l3] (D)--(E);

\tkzDrawPoints(A,B,C,D,E)
\tkzLabelPoint[below left](A){$A$}
\tkzLabelPoint[above left](B){$B$}
\tkzLabelPoint[below right](C){$\varGamma $}
\tkzLabelPoint[below](D){$\varDelta  $}
\tkzLabelPoint[above right](E){$E$}

\tkzMarkRightAngle(C,A,B)
\tkzMarkRightAngle(B,E,D)
\tkzMarkAngles[size=1,mark=none](A,B,D D,B,C)
\tkzMarkSegments[mark=s|](B,A B,E)
\end{tikzpicture}
\end{center}

Να αποδείξετε ότι:

\begin{alist}
\item $BA$=$BE$, \monades{12}

\item Αν επιπλέον $B\hat{\varDelta  }A = 55^\circ$
, να υπολογίσετε τις γωνίες του τριγώνου $\varGamma \varDelta  E$ \monades{13}
\end{alist}

\vspace{4mm}

% --- Άσκηση 34398 (34398.tex) ---
\Askhsh{34398}{2}
\begin{ylh}{1}
    \item 3.1. Είδη και στοιχεία τριγώνων
    \item 3.2. 1ο Κριτήριο ισότητας τριγώνων
    \item 4.2. Τέμνουσα δύο ευθειών - Ευκλείδειο αίτημα
    \item 5.6. Εφαρμογές στα τρίγωνα.
\end{ylh}
% ===================================================================
%  Αρχείο: 34398.tex (Fragment)
%  Αυτόματη μετατροπή μέσω doc2latex.py
% ===================================================================



Δίνεται ορθογώνιο και ισοσκελές τρίγωνο $AB\varGamma $ ($\hat{A} = 90\degree$) και η διχοτόμος $A\varDelta  $ της γωνίας του $\hat{A}$.

Από το σημείο $\varDelta  $ φέρουμε παράλληλη προς την $AB$ που τέμνει την πλευρά $A\varGamma $ στο σημείο $E$.

Να αποδείξετε ότι:

\begin{alist}
\item $\varDelta   E = \frac{A\varGamma }{2}$
, \monades{8}

\item το τρίγωνο $\varDelta  E\varGamma $ είναι ορθογώνιο και ισοσκελές, \monades{8}

\item $\varDelta   E = \frac{A\varGamma }{2}$

\begin{center}
\begin{tikzpicture}[scale=1]
\tkzDefPoint(0,5){B}
\tkzDefPoint(0,0){A}
\tkzDefPoint(5,0){C}
\draw[l3] (A)--(B)--(C)--cycle;

\tkzDefLine[bisector](B,A,C) \tkzGetPoint{bisA}
\tkzInterLL(A,bisA)(B,C) \tkzGetPoint{D}

\tkzDefLine[parallel=through D](A,B) \tkzGetPoint{par_dir}
\tkzInterLL(D,par_dir)(A,C) \tkzGetPoint{E}

\draw[l3] (A)--(D);
\draw[l3] (D)--(E);

\tkzDrawPoints(A,B,C,D,E)
\tkzLabelPoint[below left](A){$A$}
\tkzLabelPoint[above left](B){$B$}
\tkzLabelPoint[below right](C){$\varGamma $}
\tkzLabelPoint[above right](D){$\varDelta  $}
\tkzLabelPoint[below](E){$E$}

\tkzMarkRightAngle(C,A,B)
\tkzMarkSegments[mark=s|](A,E E,C) % wait, $AE$=$EG$=$DE$?
\tkzMarkAngles[size=0.6,mark=none](B,A,D D,A,E)
\end{tikzpicture}
\end{center}
. \monades{9}
\end{alist}

\vspace{4mm}

% --- Άσκηση 34399 (34399.tex) ---
\Askhsh{34399}{2}
\begin{ylh}{1}
    \item 3.1. Είδη και στοιχεία τριγώνων
    \item 3.2. 1ο Κριτήριο ισότητας τριγώνων
    \item 3.11. Ανισοτικές σχέσεις πλευρών και γωνιών
    \item 4.2. Τέμνουσα δύο ευθειών - Ευκλείδειο αίτημα.
\end{ylh}
% ===================================================================
%  Αρχείο: 34399.tex (Fragment)
%  Αυτόματη μετατροπή μέσω doc2latex.py
% ===================================================================



Σε ισοσκελές τρίγωνο $AB\varGamma $ ($AB=A\varGamma $) φέρουμε τη διχοτόμο $A\varDelta  $ και μια ευθεία (ε) παράλληλη προς την $B\varGamma $, που τέμνει τις πλευρές $AB$ και $A\varGamma $ στα σημεία $E$ και $Z$ αντίστοιχα.

Να αποδείξετε ότι:

\begin{alist}
\item το τρίγωνο $AEZ$ είναι ισοσκελές, \monades{12}

\item τα τρίγωνα Α$E\varDelta  $ και $AZ\varDelta  $ είναι ίσα. \monades{13}


\begin{center}
\begin{tikzpicture}[scale=1]
\tkzDefPoint(0,0){B}
\tkzDefPoint(4,0){C}
\tkzDefTriangle[two angles = 70 and 70](B,C) \tkzGetPoint{A}

\tkzDefLine[bisector](B,A,C) \tkzGetPoint{bisA}
\tkzInterLL(A,bisA)(B,C) \tkzGetPoint{D}

\tkzDefPointOnLine[pos=0.6](A,B) \tkzGetPoint{E}
\tkzDefLine[parallel=through E](B,C) \tkzGetPoint{par_dir}
\tkzInterLL(E,par_dir)(A,C) \tkzGetPoint{Z}

\draw[l3] (A)--(B)--(C)--cycle;
\draw[l3] (A)--(D);
\draw[l3] (E)--(Z);
\draw[l3, dashed] (E)--(D);
\draw[l3, dashed] (Z)--(D);

\tkzDrawPoints(A,B,C,D,E,Z)
\tkzLabelPoint[above](A){$A$}
\tkzLabelPoint[below left](B){$B$}
\tkzLabelPoint[below right](C){$\varGamma $}
\tkzLabelPoint[below](D){$\varDelta  $}
\tkzLabelPoint[left](E){$E$}
\tkzLabelPoint[right](Z){$Z$}

\tkzMarkSegments[mark=s|](A,B A,C)
\tkzMarkSegments[mark=s||](A,E A,Z)
\end{tikzpicture}
\end{center}
\end{alist}

\vspace{4mm}

% --- Άσκηση 34401 (34401.tex) ---
\Askhsh{34401}{2}
\begin{ylh}{1}
    \item 3.1. Είδη και στοιχεία τριγώνων
    \item 3.2. 1ο Κριτήριο ισότητας τριγώνων
    \item 3.6. Κριτήρια ισότητας ορθογώνιων τριγώνων.
\end{ylh}
% ===================================================================
%  Αρχείο: 34401.tex (Fragment)
%  Αυτόματη μετατροπή μέσω doc2latex.py
% ===================================================================



Δίνεται ισοσκελές τρίγωνο $AB\varGamma $ ($AB=A\varGamma $) και τα ύψη του $B\varDelta  $ και $\varGamma E$.

Να αποδείξετε ότι:

\begin{alist}
\item τα τρίγωνα $B\varDelta  $Γ και $\varGamma EB$ είναι ίσα, \monades{15}

\item $A\varDelta  $=$AE$. \monades{10}
\end{alist}

\vspace{4mm}

% --- Άσκηση 34404 (34404.tex) ---
\Askhsh{34404}{2}
\begin{ylh}{1}
    \item 3.1. Είδη και στοιχεία τριγώνων
    \item 3.2. 1ο Κριτήριο ισότητας τριγώνων
    \item 3.4. 3ο Κριτήριο ισότητας τριγώνων
    \item 3.6. Κριτήρια ισότητας ορθογώνιων τριγώνων.
\end{ylh}
% ===================================================================
%  Αρχείο: 34404.tex (Fragment)
%  Αυτόματη μετατροπή μέσω doc2latex.py
% ===================================================================



Θεωρούμε ισοσκελές τρίγωνο $AB\varGamma $ ($AB=A\varGamma $) και το μέσο $M$ της βάσης του $B\varGamma $. Φέρουμε τις αποστάσεις $MK$ και $M\varLambda $ του σημείου $M$ από τις ίσες πλευρές του τριγώνου $AB\varGamma $.

Να αποδείξετε ότι:

\begin{alist}
\item $MK$=$M\varLambda $, \monades{13}

\item η $AM$ είναι διχοτόμος της γωνίας Κ$\hat{Μ}$Λ. \monades{12}
\end{alist}

\vspace{4mm}

% --- Άσκηση 34405 (34405.tex) ---
\Askhsh{34405}{2}
\begin{ylh}{1}
    \item 3.1. Είδη και στοιχεία τριγώνων
    \item 3.2. 1ο Κριτήριο ισότητας τριγώνων
    \item 3.6. Κριτήρια ισότητας ορθογώνιων τριγώνων.
\end{ylh}
% ===================================================================
%  Αρχείο: 34405.tex (Fragment)
%  Αυτόματη μετατροπή μέσω doc2latex.py
% ===================================================================



Δίνεται ισοσκελές τρίγωνο $AB\varGamma $ με $AB=A\varGamma $. Από το μέσο $M$ της $B\varGamma $ φέρουμε τα κάθετα τμήματα $M\varDelta  $ και $ME$ στις πλευρές $AB$ και $A\varGamma $ αντίστοιχα.

Να αποδείξετε ότι:

\begin{alist}
\item $M\varDelta  $=$ME$ \monades{12}

\item το τρίγωνο $A\varDelta  E$ είναι ισοσκελές. \monades{13}
\end{alist}

\vspace{4mm}

% --- Άσκηση 34406 (34406.tex) ---
\Askhsh{34406}{2}
\begin{ylh}{1}
    \item 3.1. Είδη και στοιχεία τριγώνων
    \item 3.2. 1ο Κριτήριο ισότητας τριγώνων
    \item 3.11. Ανισοτικές σχέσεις πλευρών και γωνιών
    \item 4.6. Άθροισμα γωνιών τριγώνου
    \item 5.6. Εφαρμογές στα τρίγωνα
    \item 5.9. Μια ιδιότητα του ορθογώνιου τριγώνου.
\end{ylh}
% ===================================================================
%  Αρχείο: 34406.tex (Fragment)
%  Αυτόματη μετατροπή μέσω doc2latex.py
% ===================================================================



Δίνεται ορθογώνιο τρίγωνο $AB\varGamma $ ($\hat{A}=90\degree$) με $B\varGamma $ = 8 cm. Έστω $AM$ είναι διάμεσος του τριγώνου και 
$M\varDelta   \perp AC$.. Αν η γωνία 
$A\hat{M}\varGamma $ είναι ίση με 120\textsuperscript{ο}, τότε:

\begin{alist}
\item Να δείξετε ότι $AB$ = 4 cm. \monades{12}

\item Να βρείτε το μήκος της $M\varDelta  $
\item Να βρείτε το μήκος της $M\varDelta  $.

\begin{center}
\begin{tikzpicture}[scale=1]
\tkzDefPoint(0,0){A}
\tkzDefPoint(0,3.46){B} % $BC$=8, $AB$=4 => $AC$ = sqrt(64-16) = sqrt(48) = 6.92. Scale 1:2 => 3.46
\tkzDefPoint(6.92,0){C}
\draw[l3] (A)--(B)--(C)--cycle;

\tkzDefMidPoint(B,C) \tkzGetPoint{M}
\tkzDefPointBy[projection=onto A--C](M) \tkzGetPoint{D}

\draw[l3] (A)--(M);
\draw[l3] (M)--(D);

\tkzDrawPoints(A,B,C,M,D)
\tkzLabelPoint[below left](A){$A$}
\tkzLabelPoint[above left](B){$B$}
\tkzLabelPoint[below right](C){$\varGamma $}
\tkzLabelPoint[above right](M){$M$}
\tkzLabelPoint[below](D){$\varDelta  $}

\tkzMarkRightAngle(B,A,C)
\tkzMarkRightAngle(M,D,C)
\tkzMarkAngle[size=0.8,mark=none](C,M,A)
\tkzLabelAngle[pos=1.2](C,M,A){$120^\circ$}
\end{tikzpicture}
\end{center}
. \monades{13}
\end{alist}

\vspace{4mm}

% --- Άσκηση 34408 (34408.tex) ---
\Askhsh{34408}{2}
\begin{ylh}{1}
    \item 4.6. Άθροισμα γωνιών τριγώνου
    \item 5.3. Ορθογώνιο
    \item 5.9. Μια ιδιότητα του ορθογώνιου τριγώνου.
\end{ylh}
% ===================================================================
%  Αρχείο: 34408.tex (Fragment)
%  Αυτόματη μετατροπή μέσω doc2latex.py
% ===================================================================



Δίνεται τετράπλευρο $AB\varGamma \varDelta  $ με $\hat{A}$ = $\hat{\varDelta  }$ = 90\degree, $AB$ > $\varGamma \varDelta  $, $B\varGamma $ = 4 $\varDelta  \varGamma $ και $\hat{B}$ = 60\degree και $\varGamma H$ Ʇ $AB$.

Να αποδείξετε ότι:

\begin{alist}
\item $HB$ = 2ΔΓ, \monades{12}

\item το τετράπλευρο $AH\varGamma \varDelta  $ είναι ορθογώνιο με $AH$ = $\frac{1}{2}$ $HB$. \monades{13}


\begin{center}
\begin{tikzpicture}[scale=1]
\tkzDefPoint(0,4){D}
\tkzDefPoint(0,0){A}
\tkzDefPoint(2,4){C} % Just to have a length $CD$ = 2
\tkzDefPoint(6,0){B} % Just to have $AB$ > $CD$. Let's make it consistent with problem.
% $BC$ = 4 $CD$. If $CD$=x, $BC$=4x. Angle B = 60.
% $CD$ is top base, $AB$ is bottom. D, A are right angles.
% H is projection of C onto $AB$. $AH$ = $DC$. $HB$ = $AB$ - $DC$.
% In right triangle $CHB$, angle B= 60. $CH$ = $HB$ * tan 60. $BC$ = $HB$ / cos 60 = 2 $HB$.
% Problem says $BC$ = 4 $CD$. So 2 $HB$ = 4 $CD$ => $HB$ = 2 $CD$.
% So we need $HB$ = 2 $CD$.
\tkzDefPoint(0,0){A}
\tkzDefPoint(0,3.46){D} % height
\tkzDefPoint(1.5,3.46){C} % $CD$ = 1.5
\tkzDefPoint(1.5,0){H}
\tkzDefPoint(4.5,0){B} % $HB$ = 3 = 2 * $CD$.

\draw[l3] (A)--(B)--(C)--(D)--cycle;
\draw[l3] (C)--(H);

\tkzDrawPoints(A,B,C,D,H)
\tkzLabelPoint[below left](A){$A$}
\tkzLabelPoint[below right](B){$B$}
\tkzLabelPoint[above right](C){$\varGamma $}
\tkzLabelPoint[above left](D){$\varDelta  $}
\tkzLabelPoint[below](H){$H$}

\tkzMarkRightAngle(D,A,B)
\tkzMarkRightAngle(A,D,C)
\tkzMarkRightAngle(C,H,B)
\tkzMarkAngle[size=0.6,mark=none](C,B,H)
\tkzLabelAngle[pos=0.9](C,B,H){$60^\circ$}
\tkzMarkSegments[mark=s|](D,C A,H)
\tkzMarkSegments[mark=s||](H,B) % $HB$ = 2 $CD$. Maybe better mark?
\end{tikzpicture}
\end{center}
\end{alist}

\vspace{4mm}

% --- Άσκηση 34409 (34409.tex) ---
\Askhsh{34409}{2}
\begin{ylh}{1}
    \item 4.8. Άθροισμα γωνιών κυρτού ν-γώνου
    \item 5.10. Τραπέζιο
    \item 5.11. Ισοσκελές τραπέζιο.
\end{ylh}
% ===================================================================
%  Αρχείο: 34409.tex (Fragment)
%  Αυτόματη μετατροπή μέσω doc2latex.py
% ===================================================================



Δίνεται ισοσκελές τραπέζιο $AB\varGamma \varDelta  $ ($AB$ $\slash \slash$ $\varGamma \varDelta  $ και $A\varDelta  $ = $B\varGamma $), με $AB$ > $\varGamma \varDelta  $ και $MN$ η διάμεσός του.

\begin{alist}
\item Αν τα μήκη των βάσεων είναι $AB$ = 3x + 2, $\varGamma \varDelta  $ = x + 2 και το μήκος της διαμέσου του τραπεζίου είναι $MN$ = x + 4, τότε να δείξετε ότι x = 2. \monades{10}

\item Αν η γωνία $\hat{\varGamma }$ είναι διπλάσια της γωνίας $\hat{B}$, να υπολογίσετε τις γωνίες του τραπεζίου. \monades{15}


\begin{center}
\begin{tikzpicture}[scale=1]
\tkzDefPoint(1,2.5){D}
\tkzDefPoint(3,2.5){C}
\tkzDefPoint(0,0){A}
\tkzDefPoint(4,0){B}

\tkzDefMidPoint(A,D) \tkzGetPoint{M}
\tkzDefMidPoint(B,C) \tkzGetPoint{N}

\draw[l3] (A)--(B)--(C)--(D)--cycle;
\draw[l3] (M)--(N);

\tkzDrawPoints(A,B,C,D,M,N)
\tkzLabelPoint[below left](A){$A$}
\tkzLabelPoint[below right](B){$B$}
\tkzLabelPoint[above right](C){$\varGamma $}
\tkzLabelPoint[above left](D){$\varDelta  $}
\tkzLabelPoint[left](M){$M$}
\tkzLabelPoint[right](N){$N$}

\tkzMarkSegments[mark=s|](A,M M,D B,N N,C)
\tkzMarkSegments[mark=s||](A,D B,C)
\end{tikzpicture}
\end{center}
\end{alist}

\vspace{4mm}

% --- Άσκηση 34411 (34411.tex) ---
\Askhsh{34411}{2}
\begin{ylh}{1}
    \item 5.6. Εφαρμογές στα τρίγωνα
    \item 5.9. Μια ιδιότητα του ορθογώνιου τριγώνου.
\end{ylh}
% ===================================================================
%  Αρχείο: 34411.tex (Fragment)
%  Αυτόματη μετατροπή μέσω doc2latex.py
% ===================================================================



Δίνεται ισοσκελές τρίγωνο $AB\varGamma $ ($AB=A\varGamma $). Στην προέκταση της $BA$ (προς το μέρος της κορυφής Α) παίρνουμε σημείο $\varDelta  $ ώστε $AB$ = $A\varDelta  $ και στην προέκταση της $\varDelta  \varGamma $ (προς το μέρος της κορυφής Γ) παίρνουμε σημείο $E$ ώστε $\varDelta  \varGamma $ = $\varGamma E$.

Να αποδείξετε ότι:

\begin{alist}
\item το τρίγωνο $\varDelta  \varGamma B$ είναι ορθογώνιο, \monades{12}

\item $A\varGamma\parallel BE$ και $A\varGamma = \frac{BE}{2}$. \monades{13}

\begin{center}
\begin{tikzpicture}[scale=1]
\tkzDefPoint(0,0){B}
\tkzDefPoint(4,0){C}
\tkzDefTriangle[two angles = 30 and 30](B,C) \tkzGetPoint{A} % Just some angles
% Wait, A is such that $ABC$ is isosceles $AB$=$AC$.
% D is extension of $BA$ such that $AB$=$AD$.
\tkzDefPointBy[homothety=center A ratio -1](B) \tkzGetPoint{D}
% E is extension of $DC$ such that $DC$=$CE$.
\tkzDefPointBy[homothety=center C ratio -1](D) \tkzGetPoint{E}

\draw[l3] (A)--(B)--(C)--cycle;
\draw[l3] (A)--(D);
\draw[l3] (D)--(C);
\draw[l3] (C)--(E);
\draw[l3] (B)--(E);

\tkzDrawPoints(A,B,C,D,E)
\tkzLabelPoint[above](A){$A$}
\tkzLabelPoint[left](B){$B$}
\tkzLabelPoint[below](C){$\varGamma $}
\tkzLabelPoint[above](D){$\varDelta  $}
\tkzLabelPoint[right](E){$E$}

\tkzMarkSegments[mark=s|](A,B A,C A,D)
\tkzMarkSegments[mark=s||](D,C C,E)
\end{tikzpicture}
\end{center}


\end{alist}

\vspace{4mm}

% --- Άσκηση 34413 (34413.tex) ---
\Askhsh{34413}{2}
\begin{ylh}{1}
    \item 3.1. Είδη και στοιχεία τριγώνων
    \item 3.2. 1ο Κριτήριο ισότητας τριγώνων
    \item 4.2. Τέμνουσα δύο ευθειών - Ευκλείδειο αίτημα
    \item 4.6. Άθροισμα γωνιών τριγώνου.
\end{ylh}
% ===================================================================
%  Αρχείο: 34413.tex (Fragment)
%  Αυτόματη μετατροπή μέσω doc2latex.py
% ===================================================================



Ένας μαθητής της Α\textquotesingle{} λυκείου βρήκε έναν τρόπο να κατασκευάζει παράλληλες ευθείες. Στην αρχή σχεδιάζει μια τυχαία γωνία x$\hat{O}$y. Στη συνέχεια με κέντρο την κορυφή $O$ της γωνίας σχεδιάζει δυο ομόκεντρους διαφορετικούς κύκλους με τυχαίες ακτίνες. Ο μικρότερος κύκλος τέμνει τις πλευρές Οx και Οy της γωνίας στα σημεία $A$ και Β αντίστοιχα και ο με$\varGamma  A$λύτερος στα σημεία $\varGamma $ και $\varDelta$ αντίστοιχα. Ισχυρίζεται ότι οι ευθείες που ορίζονται από τις χορδές $AB$ και $\varGamma \varDelta  $ είναι παράλληλες. Μπορείτε να το δικαιολογήσετε; \monades{25}

\vspace{4mm}

% --- Άσκηση 34415 (34415.tex) ---
\Askhsh{34415}{2}
\begin{ylh}{1}
    \item 3.1. Είδη και στοιχεία τριγώνων
    \item 3.2. 1ο Κριτήριο ισότητας τριγώνων
    \item 3.11. Ανισοτικές σχέσεις πλευρών και γωνιών.
\end{ylh}
% ===================================================================
%  Αρχείο: 34415.tex (Fragment)
%  Αυτόματη μετατροπή μέσω doc2latex.py
% ===================================================================



Δίνεται ισοσκελές τρίγωνο $AB\varGamma $ ($AB=A\varGamma $). Οι διχοτόμοι των εξωτερικών γωνιών $\hat{B}$ και $\hat{\varGamma }$ τέμνονται στο σημείο $M$ και τα σημεία $K$ και Λ είναι αντίστοιχα τα μέσα των πλευρών $AB$ και $A\varGamma $.

Να αποδείξετε ότι:

\begin{alist}
\item το τρίγωνο $BM\varGamma $ είναι ισοσκελές με $MB$=$M\varGamma $, \monades{12}

\item $MK$=$M\varLambda $. \monades{13}


\begin{center}
\begin{tikzpicture}[scale=1]
\tkzDefPoint(0,0){B}
\tkzDefPoint(4,0){C}
\tkzDefTriangle[two angles = 70 and 70](B,C) \tkzGetPoint{A}

% Exterior bisectors of B and C
\tkzDefPointBy[homothety=center B ratio -0.5](C) \tkzGetPoint{X_B}
\tkzDefPointBy[homothety=center C ratio -0.5](B) \tkzGetPoint{X_C}

\tkzDefLine[bisector](A,B,X_B) \tkzGetPoint{bisBext}
\tkzDefLine[bisector](X_C,C,A) \tkzGetPoint{bisCext}

\tkzInterLL(B,bisBext)(C,bisCext) \tkzGetPoint{M}

\tkzDefMidPoint(A,B) \tkzGetPoint{K}
\tkzDefMidPoint(A,C) \tkzGetPoint{L}

\draw[l3] (A)--(B)--(C)--cycle;
\draw[l3] (B)--(X_B);
\draw[l3] (C)--(X_C);
\draw[l3] (B)--(M)--(C);
\draw[l3, dashed] (M)--(K);
\draw[l3, dashed] (M)--(L);

\tkzDrawPoints(A,B,C,M,K,L)
\tkzLabelPoint[above](A){$A$}
\tkzLabelPoint[below](B){$B$}
\tkzLabelPoint[below](C){$\varGamma $}
\tkzLabelPoint[below](M){$M$}
\tkzLabelPoint[left](K){$K$}
\tkzLabelPoint[right](L){$\Lambda$}

\tkzMarkSegments[mark=s|](A,K K,B A,L L,C)
\tkzMarkSegments[mark=s||](A,B A,C)
\end{tikzpicture}
\end{center}
\end{alist}

\vspace{4mm}

% --- Άσκηση 34418 (34418.tex) ---
\Askhsh{34418}{2}
\begin{ylh}{1}
    \item 3.1. Είδη και στοιχεία τριγώνων
    \item 3.2. 1ο Κριτήριο ισότητας τριγώνων
    \item 3.4. 3ο Κριτήριο ισότητας τριγώνων
    \item 3.11. Ανισοτικές σχέσεις πλευρών και γωνιών
    \item 4.6. Άθροισμα γωνιών τριγώνου.
\end{ylh}
% ===================================================================
%  Αρχείο: 34418.tex (Fragment)
%  Αυτόματη μετατροπή μέσω doc2latex.py
% ===================================================================



Δίνεται τρίγωνο $AB\varGamma $ στο οποίο η εξωτερική γωνία $A$ είναι διπλάσια της εσωτερικής γωνίας Α$\hat{B}$Γ.

\begin{alist}
\item Να αποδείξετε ότι το τρίγωνο $AB\varGamma $ είναι ισοσκελές με $AB=A\varGamma $. \monades{10}

\item Έστω ότι η μεσοκάθετος της πλευράς $AB$ τέμνει την πλευρά $A\varGamma $ σε εσωτερικό της σημείο $\varDelta  $.

Αν η γωνία Α$\hat{\varDelta  }$Β είναι ίση με 80˚, τότε να υπολογίσετε τις γωνίες του τριγώνου $AB\varGamma $.

\monades{15}


\begin{center}
\begin{tikzpicture}[scale=1]
\tkzDefPoint(0,0){B}
\tkzDefPoint(5,0){C}
\tkzDefTriangle[two angles = 40 and 40](B,C) \tkzGetPoint{A}

% Midpoint of $AB$
\tkzDefMidPoint(A,B) \tkzGetPoint{M_AB}
\tkzDefLine[orthogonal=through M_AB](A,B) \tkzGetPoint{perp_dir}
\tkzInterLL(M_AB,perp_dir)(A,C) \tkzGetPoint{D}

\tkzDefPointBy[homothety=center A ratio -0.5](C) \tkzGetPoint{X_A}

\draw[l3] (B)--(C)--(A)--cycle;
\draw[l3] (A)--(X_A);
\draw[l3] (D)--(B);
\draw[l3, dashed] (M_AB)--(D);

\tkzDrawPoints(A,B,C,D)
\tkzLabelPoint[above](A){$A$}
\tkzLabelPoint[below left](B){$B$}
\tkzLabelPoint[below right](C){$\varGamma $}
\tkzLabelPoint[right](D){$\varDelta  $}

\tkzMarkSegments[mark=s|](A,D D,B A,B)
\tkzMarkAngle[size=0.6,mark=none](A,D,B)
\tkzLabelAngle[pos=0.9](A,D,B){$80^\circ$}
\end{tikzpicture}
\end{center}
\end{alist}

\vspace{4mm}

% --- Άσκηση 34420 (34420.tex) ---
\Askhsh{34420}{2}
\begin{ylh}{1}
    \item 3.1. Είδη και στοιχεία τριγώνων
    \item 3.6. Κριτήρια ισότητας ορθογώνιων τριγώνων
    \item 5.9. Μια ιδιότητα του ορθογώνιου τριγώνου.
\end{ylh}
% ===================================================================
%  Αρχείο: 34420.tex (Fragment)
%  Αυτόματη μετατροπή μέσω doc2latex.py
% ===================================================================



Δίνεται ισοσκελές τρίγωνο $AB\varGamma $ ($AB=A\varGamma $). Φέρουμε, εκτός του τριγώνου, τις ημιευθείες Αx και Αy τέτοιες ώστε $Ax \perp AB$ και $Ay \perp AC$, όπως στο σχήμα που ακολουθεί. Στις Αx και Αy θεωρούμε τα σημεία $\varDelta  $ και $E$ αντίστοιχα, ώστε $A\varDelta  $=$AE$.

\begin{alist}
\item Να αποδείξετε ότι $B\varDelta  $=$\varGamma E$. \monades{12}

\item Αν $M$ και $N$ είναι τα μέσα των τμημάτων $B\varDelta  $ και $\varGamma E$ αντίστοιχα , να αποδείξετε ότι το τρίγωνο $AMN$ είναι ισοσκελές \monades{13}


\begin{center}
\begin{tikzpicture}[scale=1]
\tkzDefPoint(0,0){B}
\tkzDefPoint(4,0){C}
\tkzDefTriangle[two angles = 70 and 70](B,C) \tkzGetPoint{A}

% Ax perp $AB$
\tkzDefLine[orthogonal=through A](A,B) \tkzGetPoint{x_dir}
\tkzDefPointBy[homothety=center A ratio 0.8](x_dir) \tkzGetPoint{D}

% Ay perp $AC$
\tkzDefLine[orthogonal=through A](A,C) \tkzGetPoint{y_dir}
\tkzDefPointBy[homothety=center A ratio -0.8](y_dir) \tkzGetPoint{E}

\tkzDefMidPoint(B,D) \tkzGetPoint{M}
\tkzDefMidPoint(C,E) \tkzGetPoint{N}

\draw[l3] (A)--(B)--(C)--cycle;
\draw[l3] (A)--(D);
\draw[l3] (A)--(E);
\draw[l3] (B)--(D);
\draw[l3] (C)--(E);
\draw[l3, dashed] (A)--(M);
\draw[l3, dashed] (A)--(N);

\tkzDrawPoints(A,B,C,D,E,M,N)
\tkzLabelPoint[above](A){$A$}
\tkzLabelPoint[below](B){$B$}
\tkzLabelPoint[below](C){$\varGamma $}
\tkzLabelPoint[left](D){$\varDelta  $}
\tkzLabelPoint[right](E){$E$}
\tkzLabelPoint[left](M){$M$}
\tkzLabelPoint[right](N){$N$}

\tkzMarkRightAngle(B,A,D)
\tkzMarkRightAngle(E,A,C)
\tkzMarkSegments[mark=s|](A,D A,E)
\tkzMarkSegments[mark=s||](A,B A,C)
\end{tikzpicture}
\end{center}
\end{alist}

\vspace{4mm}

% --- Άσκηση 34422 (34422.tex) ---
\Askhsh{34422}{2}
\begin{ylh}{1}
    \item 3.1. Είδη και στοιχεία τριγώνων
    \item 3.2. 1ο Κριτήριο ισότητας τριγώνων
    \item 3.6. Κριτήρια ισότητας ορθογώνιων τριγώνων
    \item 3.11. Ανισοτικές σχέσεις πλευρών και γωνιών
    \item 4.6. Άθροισμα γωνιών τριγώνου
    \item 4.8. Άθροισμα γωνιών κυρτού ν-γώνου.
\end{ylh}
% ===================================================================
%  Αρχείο: 34422.tex (Fragment)
%  Αυτόματη μετατροπή μέσω doc2latex.py
% ===================================================================



Δίνεται το ισοσκελές τρίγωνο $AB\varGamma $ με $AB=A\varGamma $. Φέρουμε, εκτός του τριγώνου, τις ημιευθείες Αx και Αy τέτοιες ώστε $Ax \perp AB$ και $Ay \perp AC$, όπως στο σχήμα που ακολουθεί. Οι κάθετες στην πλευρά $B\varGamma $ στα σημεία $B$ και $\varGamma$ τέμνουν τις Αx και Αy στα σημεία $\varDelta  $ και $E$ αντίστοιχα.

\begin{alist}
\item Να αποδείξετε ότι $B\varDelta  $=$\varGamma E$. \monades{12}

\item Αν η γωνία Β$\hat{A}$Γ είναι ίση με 80\textsuperscript{ο}, να υπολογίσετε τις γωνίες του τριγώνου που έχει για κορυφές τα σημεία $A$, Ε και Δ. \monades{13}


\begin{center}
\begin{tikzpicture}[scale=1]
\tkzDefPoint(0,0){B}
\tkzDefPoint(4,0){C}
\tkzDefTriangle[two angles = 70 and 70](B,C) \tkzGetPoint{A}

% Ax perp $AB$
\tkzDefLine[orthogonal=through A](A,B) \tkzGetPoint{x_dir}
\tkzDefLine[orthogonal=through B](B,C) \tkzGetPoint{v_dir}
\tkzInterLL(A,x_dir)(B,v_dir) \tkzGetPoint{D}

% Ay perp $AC$
\tkzDefLine[orthogonal=through A](A,C) \tkzGetPoint{y_dir}
\tkzDefLine[orthogonal=through C](B,C) \tkzGetPoint{w_dir}
\tkzInterLL(A,y_dir)(C,w_dir) \tkzGetPoint{E}

\draw[l3] (A)--(B)--(C)--cycle;
\draw[l3] (A)--(D);
\draw[l3] (A)--(E);
\draw[l3] (B)--(D);
\draw[l3] (C)--(E);

\tkzDrawPoints(A,B,C,D,E)
\tkzLabelPoint[above](A){$A$}
\tkzLabelPoint[below](B){$B$}
\tkzLabelPoint[below](C){$\varGamma $}
\tkzLabelPoint[left](D){$\varDelta  $}
\tkzLabelPoint[right](E){$E$}

\tkzMarkRightAngle(B,A,D)
\tkzMarkRightAngle(E,A,C)
\tkzMarkRightAngle(C,B,D)
\tkzMarkRightAngle(B,C,E)
\end{tikzpicture}
\end{center}
\end{alist}

\vspace{4mm}

% --- Άσκηση 34423 (34423.tex) ---
\Askhsh{34423}{2}
\begin{ylh}{1}
    \item 3.1. Είδη και στοιχεία τριγώνων
    \item 3.2. 1ο Κριτήριο ισότητας τριγώνων
    \item 4.2. Τέμνουσα δύο ευθειών - Ευκλείδειο αίτημα
    \item 5.2. Παραλληλόγραμμα.
\end{ylh}
% ===================================================================
%  Αρχείο: 34423.tex (Fragment)
%  Αυτόματη μετατροπή μέσω doc2latex.py
% ===================================================================



Δίνεται παραλληλόγραμμο $AB\varGamma \varDelta  $ με $AB$=2$B\varGamma $ και $E$ το μέσο της πλευράς του $AB$.

Να αποδείξετε ότι:

\begin{alist}
\item το τρίγωνο $EA\varDelta  $ είναι ισοσκελές, \monades{10}

\item η $\varDelta  E$ είναι διχοτόμος της γωνίας $\hat{\varDelta  }$. \monades{15}

\begin{center}
\begin{tikzpicture}[scale=1]
\tkzDefPoint(0,0){A}
\tkzDefPoint(6,0){B}
\tkzInterCC[R](A,3)(B,10) \tkzGetPoints{D_alt}{D} % Pick D to make parallelogram
\tkzDefPoint(1.5,2.5){D} % Manual override for better look
\tkzDefPoint(6,0){B}
\tkzDefPointBy[translation=from A to B](D) \tkzGetPoint{C}

\tkzDefMidPoint(A,B) \tkzGetPoint{E}

\draw[l3] (A)--(B)--(C)--(D)--cycle;
\draw[l3] (D)--(E);

\tkzDrawPoints(A,B,C,D,E)
\tkzLabelPoint[below left](A){$A$}
\tkzLabelPoint[below](B){$B$}
\tkzLabelPoint[above right](C){$\varGamma $}
\tkzLabelPoint[above left](D){$\varDelta  $}
\tkzLabelPoint[below](E){$E$}

\tkzMarkSegments[mark=s|](A,D B,C A,E E,B)
\tkzMarkAngles[size=0.6,mark=none](A,D,E E,D,C)
\end{tikzpicture}
\end{center}
\end{alist}

\vspace{4mm}

% --- Άσκηση 34424 (34424.tex) ---
\Askhsh{34424}{2}
\begin{ylh}{1}
    \item 3.1. Είδη και στοιχεία τριγώνων
    \item 3.2. 1ο Κριτήριο ισότητας τριγώνων
    \item 3.4. 3ο Κριτήριο ισότητας τριγώνων
    \item 3.7. Κύκλος - Μεσοκάθετος - Διχοτόμος.
\end{ylh}
% ===================================================================
%  Αρχείο: 34424.tex (Fragment)
%  Αυτόματη μετατροπή μέσω doc2latex.py
% ===================================================================



Θεωρούμε ισοσκελές τρίγωνο $AB\varGamma $ ($AB=A\varGamma $) και Ι το σημείο τομής των διχοτόμων των γωνιών $\hat{B}$ και $\hat{\varGamma }$.

Να αποδείξετε ότι:

\begin{alist}
\item το τρίγωνο $BI\varGamma $ είναι ισοσκελές, \monades{8}

\item οι γωνίες $Α\hat{Ι}\varGamma $ και $A\hat{Ι}Β$ είναι ίσες, \monades{10}

\item η ευθεία $AI$ είναι μεσοκάθετος του τμήματος $B\varGamma $. \monades{7}
\end{alist}

\vspace{4mm}

% --- Άσκηση 34425 (34425.tex) ---
\Askhsh{34425}{2}
\begin{ylh}{1}
    \item 5.2. Παραλληλόγραμμα.
\end{ylh}
% ===================================================================
%  Αρχείο: 34425.tex (Fragment)
%  Αυτόματη μετατροπή μέσω doc2latex.py
% ===================================================================



Δίνεται τρίγωνο $AB\varGamma $ και η διάμεσός του $AM$. Στην προέκταση της διαμέσου $M\varDelta  $ του τριγώνου $AM\varGamma $ (προς το Δ) θεωρούμε σημείο $E$ ώστε $M\varDelta  $=$\varDelta  E$.

Να αποδείξετε ότι:

\begin{alist}
\item το τετράπλευρο $AM\varGamma E$ είναι παραλληλόγραμμο, \monades{12}

\item η $BE$ διέρχεται από το μέσο της διαμέσου $AM$. \monades{13}
\end{alist}

\vspace{4mm}

% --- Άσκηση 34426 (34426.tex) ---
\Askhsh{34426}{2}
\begin{ylh}{1}
    \item 3.4. 3ο Κριτήριο ισότητας τριγώνων
    \item 4.2. Τέμνουσα δύο ευθειών - Ευκλείδειο αίτημα
    \item 5.2. Παραλληλόγραμμα
    \item 5.3. Ορθογώνιο
    \item 5.6. Εφαρμογές στα τρίγωνα.
\end{ylh}
% ===================================================================
%  Αρχείο: 34426.tex (Fragment)
%  Αυτόματη μετατροπή μέσω doc2latex.py
% ===================================================================



Δίνεται ισοσκελές τρίγωνο $AB\varGamma $ ($AB=A\varGamma $) και η διάμεσός του $AM$. Στο τρίγωνο $AM\varGamma $ θεωρούμε

τη διάμεσο $M\varDelta  $ την οποία προεκτείνουμε προς το $\varDelta$ κατά τμήμα $\varDelta  E$=$\varDelta  M$. Φέρουμε από το

σημείο $\varDelta  $ τμήμα $\varDelta   Z$ κάθετο στην $AM$.

Να αποδείξετε ότι:

\begin{alist}
\item το τετράπλευρο $AM\varGamma E$ είναι ορθογώνιο, \monades{12}

\item $\varDelta   Z = \frac{B\varGamma }{4}$. \monades{13}

\begin{center}
\begin{tikzpicture}[scale=1]
\tkzDefPoint(0,0){B}
\tkzDefPoint(6,0){C}
\tkzDefTriangle[two angles = 70 and 70](B,C) \tkzGetPoint{A}
\tkzDefMidPoint(B,C) \tkzGetPoint{M}
\tkzDefMidPoint(M,C) \tkzGetPoint{D} % Problem says median of $AMG$?
% Wait: "Στο τρίγωνο $AM\varGamma $ θεωρούμε τη διάμεσο $M\varDelta  $".
% So D is midpoint of $AG$.
\tkzDefMidPoint(A,C) \tkzGetPoint{D}
% $DE$ = $DM$
\tkzDefPointBy[homothety=center D ratio -1](M) \tkzGetPoint{E}
% $DZ$ perp $AM$
\tkzDefPointBy[projection=onto A--M](D) \tkzGetPoint{Z}

\draw[l3] (A)--(B)--(C)--cycle;
\draw[l3] (A)--(M);
\draw[l3] (M)--(E);
\draw[l3] (A)--(E);
\draw[l3] (C)--(E);
\draw[l3] (D)--(Z);

\tkzDrawPoints(A,B,C,M,D,E,Z)
\tkzLabelPoint[above](A){$A$}
\tkzLabelPoint[below left](B){$B$}
\tkzLabelPoint[below right](C){$\varGamma $}
\tkzLabelPoint[below](M){$M$}
\tkzLabelPoint[right](D){$\varDelta  $}
\tkzLabelPoint[above right](E){$E$}
\tkzLabelPoint[left](Z){$Z$}

\tkzMarkRightAngle(A,M,C)
\tkzMarkRightAngle(D,Z,M)
\tkzMarkSegments[mark=s|](M,D D,E)
\tkzMarkSegments[mark=s||](A,D D,C)
\end{tikzpicture}
\end{center}


\end{alist}

\vspace{4mm}

% --- Άσκηση 34488 (34488.tex) ---
\Askhsh{34488}{2}
\begin{ylh}{1}
    \item 4.2. Τέμνουσα δύο ευθειών - Ευκλείδειο αίτημα
    \item 4.6. Άθροισμα γωνιών τριγώνου
    \item 5.3. Ορθογώνιο
    \item 5.9. Μια ιδιότητα του ορθογώνιου τριγώνου
    \item 5.11. Ισοσκελές τραπέζιο.
\end{ylh}
% ===================================================================
%  Αρχείο: 34488.tex (Fragment)
%  Αυτόματη μετατροπή μέσω doc2latex.py
% ===================================================================



Θεωρούμε ισοσκελές τραπέζιο $AB\varGamma \varDelta  $ ($AB$//$\varGamma \varDelta  $) με $\hat{\varGamma } = \hat{\varDelta  } = 60^{\circ}$
και τα κάθετα τμήματα $AE$, $BZ$ στη $\varDelta  \varGamma $.

Να αποδείξετε ότι:

\begin{alist}
\item $\varDelta  E$ = $\varGamma Z$, \monades{12}

\item το τετράπλευρο $AEZB$ είναι ορθογώνιο. \monades{13}


\begin{center}
\begin{tikzpicture}[scale=1]
\tkzDefPoint(0,0){D}
\tkzDefPoint(6,0){C}
\tkzDefPoint(1.5,2.59){A} % angle 60
\tkzDefPoint(4.5,2.59){B}

\tkzDefPointBy[projection=onto D--C](A) \tkzGetPoint{E}
\tkzDefPointBy[projection=onto D--C](B) \tkzGetPoint{Z}

\draw[l3] (A)--(B)--(C)--(D)--cycle;
\draw[l3, dashed] (A)--(E);
\draw[l3, dashed] (B)--(Z);

\tkzDrawPoints(A,B,C,D,E,Z)
\tkzLabelPoint[above left](A){$A$}
\tkzLabelPoint[above right](B){$B$}
\tkzLabelPoint[below right](C){$\varGamma $}
\tkzLabelPoint[below left](D){$\varDelta  $}
\tkzLabelPoint[below](E){$E$}
\tkzLabelPoint[below](Z){$Z$}

\tkzMarkRightAngle(A,E,C)
\tkzMarkRightAngle(B,Z,D)
\tkzMarkAngle[size=0.6,mark=none](A,D,E)
\tkzLabelAngle[pos=0.9](A,D,E){$60^\circ$}
\tkzMarkAngle[size=0.6,mark=none](B,C,Z)
\tkzLabelAngle[pos=0.9](B,C,Z){$60^\circ$}
\end{tikzpicture}
\end{center}
\end{alist}

\vspace{4mm}

% --- Άσκηση 34491 (34491.tex) ---
\Askhsh{34491}{2}
\begin{ylh}{1}
    \item 3.1. Είδη και στοιχεία τριγώνων
    \item 3.6. Κριτήρια ισότητας ορθογώνιων τριγώνων
    \item 4.2. Τέμνουσα δύο ευθειών - Ευκλείδειο αίτημα
    \item 5.3. Ορθογώνιο
    \item 5.10. Τραπέζιο
    \item 5.11. Ισοσκελές τραπέζιο.
\end{ylh}
% ===================================================================
%  Αρχείο: 34491.tex (Fragment)
%  Αυτόματη μετατροπή μέσω doc2latex.py
% ===================================================================



Θεωρούμε ισοσκελές τραπέζιο $AB\varGamma \varDelta  $ ($AB$//$\varGamma \varDelta  $). Από τα σημεία $A$ και Β φέρνουμε τα κάθετα τμήματα $AE$ και $BZ$ αντίστοιχα στη $\varDelta  \varGamma $.

Να αποδείξετε ότι:

\begin{alist}
\item $\varDelta  E$ = $\varGamma Z$, \monades{12}

\item $AZ$ = $BE$. \monades{13}


\begin{center}
\begin{tikzpicture}[scale=1]
\tkzDefPoint(0,0){D}
\tkzDefPoint(6,0){C}
\tkzDefPoint(1.5,2.5){A}
\tkzDefPoint(4.5,2.5){B}

\tkzDefPointBy[projection=onto D--C](A) \tkzGetPoint{E}
\tkzDefPointBy[projection=onto D--C](B) \tkzGetPoint{Z}

\draw[l3] (A)--(B)--(C)--(D)--cycle;
\draw[l3] (A)--(E);
\draw[l3] (B)--(Z);
\draw[l3, dashed] (A)--(Z);
\draw[l3, dashed] (B)--(E);

\tkzDrawPoints(A,B,C,D,E,Z)
\tkzLabelPoint[above left](A){$A$}
\tkzLabelPoint[above right](B){$B$}
\tkzLabelPoint[below right](C){$\varGamma $}
\tkzLabelPoint[below left](D){$\varDelta  $}
\tkzLabelPoint[below](E){$E$}
\tkzLabelPoint[below](Z){$Z$}

\tkzMarkRightAngle(A,E,C)
\tkzMarkRightAngle(B,Z,D)
\tkzMarkSegments[mark=s|](D,E C,Z)
\end{tikzpicture}
\end{center}
\end{alist}

\vspace{4mm}

% --- Άσκηση 34492 (34492.tex) ---
\Askhsh{34492}{2}
\begin{ylh}{1}
    \item 3.1. Είδη και στοιχεία τριγώνων
    \item 3.2. 1ο Κριτήριο ισότητας τριγώνων
    \item 3.6. Κριτήρια ισότητας ορθογώνιων τριγώνων
    \item 5.4. Ρόμβος
    \item 5.9. Μια ιδιότητα του ορθογώνιου τριγώνου.
\end{ylh}
% ===================================================================
%  Αρχείο: 34492.tex (Fragment)
%  Αυτόματη μετατροπή μέσω doc2latex.py
% ===================================================================



Θεωρούμε ισοσκελές τρίγωνο $AB\varGamma $ ($AB=A\varGamma $), το ύψος του $A\varDelta  $ και τα μέσα $E$ και $Z$ των πλευρών του $AB$ και $A\varGamma $ αντίστοιχα.

Να αποδείξετε ότι:

\begin{alist}
\item τα τρίγωνα $B\varDelta  $Ε και Γ$\varDelta   Z$ είναι ίσα, \monades{15}

\item το τετράπλευρο $AZ\varDelta  E$ είναι ρόμβος. \monades{10}


\begin{center}
\begin{tikzpicture}[scale=1]
\tkzDefPoint(0,0){B}
\tkzDefPoint(4,0){C}
\tkzDefTriangle[two angles = 70 and 70](B,C) \tkzGetPoint{A}
\tkzDefPointBy[projection=onto B--C](A) \tkzGetPoint{D}
\tkzDefMidPoint(A,B) \tkzGetPoint{E}
\tkzDefMidPoint(A,C) \tkzGetPoint{Z}

\draw[l3] (A)--(B)--(C)--cycle;
\draw[l3] (A)--(D);
\draw[l3] (E)--(D)--(Z);

\tkzDrawPoints(A,B,C,D,E,Z)
\tkzLabelPoint[above](A){$A$}
\tkzLabelPoint[below](B){$B$}
\tkzLabelPoint[below](C){$\varGamma $}
\tkzLabelPoint[below](D){$\varDelta  $}
\tkzLabelPoint[left](E){$E$}
\tkzLabelPoint[right](Z){$Z$}

\tkzMarkRightAngle(A,D,B)
\tkzMarkSegments[mark=s|](A,E E,B A,Z Z,C)
\tkzMarkSegments[mark=s||](E,D Z,D)
\end{tikzpicture}
\end{center}
\end{alist}

\vspace{4mm}

% --- Άσκηση 34493 (34493.tex) ---
\Askhsh{34493}{2}
\begin{ylh}{1}
    \item 3.1. Είδη και στοιχεία τριγώνων
    \item 3.2. 1ο Κριτήριο ισότητας τριγώνων
    \item 3.3. 2ο Κριτήριο ισότητας τριγώνων.
\end{ylh}
% ===================================================================
%  Αρχείο: 34493.tex (Fragment)
%  Αυτόματη μετατροπή μέσω doc2latex.py
% ===================================================================



Έστω δυο ισοσκελή τρίγωνα $AB\varGamma $ ($AB=A\varGamma $) και $A'B'\varGamma  ' (A'B'=A'\varGamma ')$.

\begin{alist}
\item Αν ισχύει $AB = A'B'$ και $\hat{A} = \hat{A'}$, να αποδείξετε ότι τα τρίγωνα $AB\varGamma $ και $A'B'\varGamma  '$ είναι ίσα. \monades{13}

\begin{center}
\begin{tikzpicture}[scale=0.8]
\tkzDefPoint(0,0){B}
\tkzDefPoint(3,0){C}
\tkzDefTriangle[two angles = 70 and 70](B,C) \tkzGetPoint{A}
\draw[l3] (A)--(B)--(C)--cycle;
\tkzDrawPoints(A,B,C)
\tkzLabelPoint[above](A){$A$}
\tkzLabelPoint[below](B){$B$}
\tkzLabelPoint[below](C){$\varGamma $}

\begin{scope}[xshift=5cm]
\tkzDefPoint(0,0){B'}
\tkzDefPoint(3,0){C'}
\tkzDefTriangle[two angles = 70 and 70](B',C') \tkzGetPoint{A'}
\draw[l3] (A')--(B')--(C')--cycle;
\tkzDrawPoints(A',B',C')
\tkzLabelPoint[above](A'){$A'$}
\tkzLabelPoint[below](B'){$B'$}
\tkzLabelPoint[below](C'){$\varGamma '$}
\end{scope}
\end{tikzpicture}
\end{center}

\item Αν ισχύει $A\varGamma  = A'\varGamma '$ και $\hat{B} = \hat{B'}$, να αποδείξετε ότι τα τρίγωνα $AB\varGamma $ και $A'B'\varGamma  '$ είναι ίσα. \monades{12}
\end{alist}

\vspace{4mm}

% --- Άσκηση 34494 (34494.tex) ---
\Askhsh{34494}{2}
\begin{ylh}{1}
    \item 5.2. Παραλληλόγραμμα
    \item 5.6. Εφαρμογές στα τρίγωνα.
\end{ylh}
% ===================================================================
%  Αρχείο: 34494.tex (Fragment)
%  Αυτόματη μετατροπή μέσω doc2latex.py
% ===================================================================



Θεωρούμε τρίγωνο $AB\varGamma $ και τα μέσα Δ, Ε και $Z$ των πλευρών του $AB$, $B\varGamma $ και $\varGamma  A$ αντίστοιχα.

Να αποδείξετε ότι:

\begin{alist}
\item το τετράπλευρο $\varDelta   BEZ$ είναι παραλληλόγραμμο, \monades{13}

\item η ευθεία $\varDelta   Z$ διχοτομεί το τμήμα $AE$. \monades{12}
\end{alist}

\vspace{4mm}

% --- Άσκηση 34495 (34495.tex) ---
\Askhsh{34495}{2}
\begin{ylh}{1}
    \item 3.1. Είδη και στοιχεία τριγώνων
    \item 3.6. Κριτήρια ισότητας ορθογώνιων τριγώνων
    \item 3.11. Ανισοτικές σχέσεις πλευρών και γωνιών
    \item 4.6. Άθροισμα γωνιών τριγώνου
    \item 5.9. Μια ιδιότητα του ορθογώνιου τριγώνου.
\end{ylh}
% ===================================================================
%  Αρχείο: 34495.tex (Fragment)
%  Αυτόματη μετατροπή μέσω doc2latex.py
% ===================================================================



Δίνεται οξυγώνιο τρίγωνο $AB\varGamma $ με $AB<A\varGamma $ και γωνία $\hat{\varGamma } = 30^{0}$. Θεωρούμε το ύψος του $A\varDelta  $ και το μέσο $Z$ της πλευράς $A\varGamma $.

\begin{alist}
\item Να αποδείξετε ότι $\varDelta   Z$ =$\frac{A\varGamma }{2}$. \monades{12}

\item Προεκτείνουμε το ύψος $A\varDelta  $ (προς το Δ) κατά ίσο τμήμα $\varDelta  E$. Να αποδείξετε ότι το τρίγωνο $A\varGamma E$ είναι ισόπλευρο. \monades{13}
\end{alist}

\vspace{4mm}

% --- Άσκηση 34496 (34496.tex) ---
\Askhsh{34496}{2}
\begin{ylh}{1}
    \item 3.1. Είδη και στοιχεία τριγώνων
    \item 3.2. 1ο Κριτήριο ισότητας τριγώνων
    \item 3.6. Κριτήρια ισότητας ορθογώνιων τριγώνων
    \item 3.11. Ανισοτικές σχέσεις πλευρών και γωνιών.
\end{ylh}
% ===================================================================
%  Αρχείο: 34496.tex (Fragment)
%  Αυτόματη μετατροπή μέσω doc2latex.py
% ===================================================================



Θεωρούμε τρίγωνο $AB\varGamma $ και τα ύψη του $B\varDelta  $ και $\varGamma E$ που αντιστοιχούν στις πλευρές του $A\varGamma $ και $AB$ αντίστοιχα.

Να αποδείξετε ότι:

\begin{alist}
\item Αν το τρίγωνο $AB\varGamma $ είναι ισοσκελές με $AB=A\varGamma $, τότε τα ύψη $B\varDelta  $ και $\varGamma E$ είναι ίσα.

\monades{12}

\item Αν τα ύψη $B\varDelta  $ και $\varGamma E$ είναι ίσα, τότε το τρίγωνο $AB\varGamma $ είναι ισοσκελές με $A\varGamma $=$AB$.

\monades{13}
\end{alist}

\vspace{4mm}

% --- Άσκηση 34497 (34497.tex) ---
\Askhsh{34497}{2}
\begin{ylh}{1}
    \item 3.1. Είδη και στοιχεία τριγώνων
    \item 3.2. 1ο Κριτήριο ισότητας τριγώνων
    \item 3.6. Κριτήρια ισότητας ορθογώνιων τριγώνων.
\end{ylh}
% ===================================================================
%  Αρχείο: 34497.tex (Fragment)
%  Αυτόματη μετατροπή μέσω doc2latex.py
% ===================================================================



Σε οξυγώνιο τρίγωνο $AB\varGamma $ προεκτείνουμε τη διάμεσο $AM$ (προς το Μ) κατά ίσο τμήμα $M\varDelta  $.

Να αποδείξετε ότι:

\begin{alist}
\item τα τρίγωνα $ABM$ και $M\varGamma \varDelta  $ είναι ίσα, \monades{12}

\item τα σημεία $A$ και $\varDelta$ ισαπέχουν από την πλευρά $B\varGamma $. \monades{13}
\end{alist}

\vspace{4mm}

% --- Άσκηση 34498 (34498.tex) ---
\Askhsh{34498}{2}
\begin{ylh}{1}
    \item 3.1. Είδη και στοιχεία τριγώνων
    \item 3.2. 1ο Κριτήριο ισότητας τριγώνων
    \item 5.4. Ρόμβος.
\end{ylh}
% ===================================================================
%  Αρχείο: 34498.tex (Fragment)
%  Αυτόματη μετατροπή μέσω doc2latex.py
% ===================================================================



Θεωρούμε οξυγώνιο τρίγωνο $AB\varGamma $ με $AB<A\varGamma $ και το ύψος του $A\varDelta  $. Προεκτείνουμε το $A\varDelta  $ (προς το Δ) κατά τμήμα $\varDelta  E$=$A\varDelta  $. Έστω Κ σημείο της $B\varGamma $ τέτοιο ώστε $B\varDelta  $ = $\varDelta  K$.

Να αποδείξετε ότι:

\begin{alist}
\item το τρίγωνο $ABK$ είναι ισοσκελές, \monades{12}

\item το τετράπλευρο Α$BE$Κ είναι ρόμβος. \monades{13}
\end{alist}

\vspace{4mm}

% --- Άσκηση 34499 (34499.tex) ---
\Askhsh{34499}{2}
\begin{ylh}{1}
    \item 3.1. Είδη και στοιχεία τριγώνων
    \item 3.6. Κριτήρια ισότητας ορθογώνιων τριγώνων.
\end{ylh}
% ===================================================================
%  Αρχείο: 34499.tex (Fragment)
%  Αυτόματη μετατροπή μέσω doc2latex.py
% ===================================================================



Δίνεται ορθογώνιο τρίγωνο $AB\varGamma $ ($\hat{A} = 90^{0}$) και $B\varDelta  $ η διχοτόμος της γωνίας $\hat{B}$. Από το $\varDelta$ φέρουμε ευθεία κάθετη στη $B\varGamma $ που την τέμνει σε σημείο $E$ και έστω $Z$ το σημείο στο οποίο η $E\varDelta  $ τέμνει την προέκταση της πλευράς $BA$ προς το A.

Να αποδείξετε ότι:

\begin{alist}
\item $AB$=$BE$, \monades{13}

\item τα τρίγωνα $AB\varGamma $ και $ZEB$ είναι ίσα. \monades{12}
\end{alist}

\vspace{4mm}

% --- Άσκηση 34500 (34500.tex) ---
\Askhsh{34500}{2}
\begin{ylh}{1}
    \item 3.1. Είδη και στοιχεία τριγώνων
    \item 3.2. 1ο Κριτήριο ισότητας τριγώνων
    \item 3.6. Κριτήρια ισότητας ορθογώνιων τριγώνων
    \item 4.6. Άθροισμα γωνιών τριγώνου.
\end{ylh}
% ===================================================================
%  Αρχείο: 34500.tex (Fragment)
%  Αυτόματη μετατροπή μέσω doc2latex.py
% ===================================================================



Θωρούμε ισοσκελές τρίγωνο $AB\varGamma $ ($AB=A\varGamma $) και σημεία $\varDelta  $ και $E$ στην ευθεία $B\varGamma $ τέτοια, ώστε $B\varDelta  $=$\varGamma E$. Έστω ότι $\varDelta   Z\bot ΑΒ$ και $EH\bot A\varGamma $.

\begin{alist}
\item Να αποδείξετε ότι:
\begin{rlist}
\item $BZ = \varGamma  H$, \monades{10}
\item το τρίγωνο $AZH$ είναι ισοσκελές. \monades{7}
\end{rlist}
\item Αν $\hat{A} = 50\degree$, να υπολογίσετε τις γωνίες του τριγώνου $AZH$. \monades{8}

\begin{center}
\begin{tikzpicture}[scale=1]
\tkzDefPoint(0,0){B}
\tkzDefPoint(4,0){C}
\tkzDefTriangle[two angles = 70 and 70](B,C) \tkzGetPoint{A}

\tkzDefPointBy[homothety=center B ratio -0.4](C) \tkzGetPoint{D}
\tkzDefPointBy[homothety=center C ratio -0.4](B) \tkzGetPoint{E}

\tkzDefPointBy[projection=onto A--B](D) \tkzGetPoint{Z}
\tkzDefPointBy[projection=onto A--C](E) \tkzGetPoint{H}

\draw[l3] (A)--(B)--(C)--cycle;
\draw[l3] (B)--(D);
\draw[l3] (C)--(E);
\draw[l3] (D)--(Z);
\draw[l3] (E)--(H);
\draw[l3, dashed] (Z)--(H);

\tkzDrawPoints(A,B,C,D,E,Z,H)
\tkzLabelPoint[above](A){$A$}
\tkzLabelPoint[below](B){$B$}
\tkzLabelPoint[below](C){$\varGamma $}
\tkzLabelPoint[below](D){$\varDelta  $}
\tkzLabelPoint[below](E){$E$}
\tkzLabelPoint[left](Z){$Z$}
\tkzLabelPoint[right](H){$H$}

\tkzMarkRightAngle(D,Z,A)
\tkzMarkRightAngle(E,H,A)
\tkzMarkSegments[mark=s|](B,D C,E)
\tkzMarkSegments[mark=s||](A,B A,C)
\end{tikzpicture}
\end{center}
\end{alist}

\vspace{4mm}

% --- Άσκηση 34502 (34502.tex) ---
\Askhsh{34502}{2}
\begin{ylh}{1}
    \item 3.1. Είδη και στοιχεία τριγώνων
    \item 3.2. 1ο Κριτήριο ισότητας τριγώνων
    \item 3.12. Τριγωνική ανισότητα.
\end{ylh}
% ===================================================================
%  Αρχείο: 34502.tex (Fragment)
%  Αυτόματη μετατροπή μέσω doc2latex.py
% ===================================================================



Στο ακόλουθο σχήμα, η $A\varDelta  $ είναι διάμεσος του τριγώνου $AB\varGamma $ και το $E$ είναι σημείο στην προέκταση της $A\varDelta  $, ώστε $\varDelta  E$=$A\varDelta  $.

Να αποδείξετε ότι:

\begin{alist}
\item $AB$ = $\varGamma E$, \monades{12}

\item $AE$ < $AB$ + $A\varGamma $. \monades{13}


\begin{center}
\begin{tikzpicture}[scale=1]
\tkzDefPoint(0,0){B}
\tkzDefPoint(5,0){C}
\tkzDefPoint(1,3){A}
\tkzDefMidPoint(B,C) \tkzGetPoint{D}
\tkzDefPointBy[homothety=center D ratio -1](A) \tkzGetPoint{E}

\draw[l3] (A)--(B)--(C)--cycle;
\draw[l3] (A)--(E);
\draw[l3] (B)--(E);
\draw[l3] (C)--(E);

\tkzDrawPoints(A,B,C,D,E)
\tkzLabelPoint[above](A){$A$}
\tkzLabelPoint[left](B){$B$}
\tkzLabelPoint[right](C){$\varGamma $}
\tkzLabelPoint[below right](D){$\varDelta  $}
\tkzLabelPoint[below](E){$E$}

\tkzMarkSegments[mark=s|](B,D D,C)
\tkzMarkSegments[mark=s||](A,D D,E)
\end{tikzpicture}
\end{center}
\end{alist}

\vspace{4mm}

% --- Άσκηση 34503 (34503.tex) ---
\Askhsh{34503}{2}
\begin{ylh}{1}
    \item 3.1. Είδη και στοιχεία τριγώνων
    \item 3.6. Κριτήρια ισότητας ορθογώνιων τριγώνων
    \item 3.7. Κύκλος - Μεσοκάθετος - Διχοτόμος.
\end{ylh}
% ===================================================================
%  Αρχείο: 34503.tex (Fragment)
%  Αυτόματη μετατροπή μέσω doc2latex.py
% ===================================================================



Δίνεται ορθογώνιο τρίγωνο $AB\varGamma $ ($\hat{A} = 90^{0}$) και η διχοτόμος της γωνίας του $\hat{\varGamma }$, η οποία τέμνει την πλευρά $AB$ στο Δ. Από το $\varDelta$ φέρουμε τμήμα $\varDelta  E$ κάθετο στην πλευρά $B\varGamma $.

Να αποδείξετε ότι:

\begin{alist}
\item τα τρίγωνα $A\varGamma \varDelta  $ και $\varDelta  \varGamma E$ είναι ίσα, \monades{13}

\item Το $\varGamma$ ισαπέχει από τα σημεία $A$ και $E$ και η ευθεία $\varGamma \varDelta  $ είναι μεσοκάθετος του τμήματος $AE$. \monades{12}
\end{alist}

\vspace{4mm}

% --- Άσκηση 34504 (34504.tex) ---
\Askhsh{34504}{2}
\begin{ylh}{1}
    \item 3.1. Είδη και στοιχεία τριγώνων
    \item 3.6. Κριτήρια ισότητας ορθογώνιων τριγώνων
    \item 5.2. Παραλληλόγραμμα
    \item 5.4. Ρόμβος.
\end{ylh}
% ===================================================================
%  Αρχείο: 34504.tex (Fragment)
%  Αυτόματη μετατροπή μέσω doc2latex.py
% ===================================================================



Το τετράπλευρο $AB\varGamma \varDelta  $ του σχήματος είναι παραλληλόγραμμο. Έστω ότι τα τμήματα $AZ$ και $AE$ είναι κάθετα στις πλευρές $\varDelta  \varGamma $ και $\varGamma B$ αντίστοιχα.

Να αποδείξετε ότι:

\begin{alist}
\item Αν το παραλληλόγραμμο $AB\varGamma \varDelta  $ είναι ρόμβος, τότε $AZ$=$AE$. \monades{12}

\item Αν για το παραλληλόγραμμο $AB\varGamma \varDelta  $ ισχύει $AZ$=$AE$, τότε αυτό είναι ρόμβος.

\monades{13}


\begin{center}
\begin{tikzpicture}[scale=1]
\tkzDefPoint(0,0){D}
\tkzDefPoint(5,0){C}
\tkzDefPoint(1.5,2.5){A}
\tkzDefPointBy[translation=from D to C](A) \tkzGetPoint{B}

\tkzDefPointBy[projection=onto D--C](A) \tkzGetPoint{Z}
\tkzDefPointBy[projection=onto B--C](A) \tkzGetPoint{E}

\draw[l3] (A)--(B)--(C)--(D)--cycle;
\draw[l3] (A)--(Z);
\draw[l3] (A)--(E);

\tkzDrawPoints(A,B,C,D,Z,E)
\tkzLabelPoint[above left](A){$A$}
\tkzLabelPoint[above right](B){$B$}
\tkzLabelPoint[right](C){$\varGamma $}
\tkzLabelPoint[left](D){$\varDelta  $}
\tkzLabelPoint[below](Z){$Z$}
\tkzLabelPoint[right](E){$E$}

\tkzMarkRightAngle(A,Z,C)
\tkzMarkRightAngle(A,E,B)
\end{tikzpicture}
\end{center}
\end{alist}

\vspace{4mm}

% --- Άσκηση 34505 (34505.tex) ---
\Askhsh{34505}{2}
\begin{ylh}{1}
    \item 3.1. Είδη και στοιχεία τριγώνων
    \item 3.6. Κριτήρια ισότητας ορθογώνιων τριγώνων
    \item 5.2. Παραλληλόγραμμα
    \item 5.4. Ρόμβος.
\end{ylh}
% ===================================================================
%  Αρχείο: 34505.tex (Fragment)
%  Αυτόματη μετατροπή μέσω doc2latex.py
% ===================================================================



Σε ημικύκλιο διαμέτρου $AB$ προεκτείνουμε την $AB$ προς το μέρος του Α και παίρνουμε ένα σημείο $\varGamma $. Θεωρούμε $E$ ένα σημείο του ημικυκλίου και έστω $\varDelta$ το σημείο τομής του τμήματος $\varGamma E$ με το ημικύκλιο. Αν το τμήμα $\varGamma \varDelta  $ είναι ίσο με το $OB$ και η γωνία $B\hat{\varGamma }E\ \varepsilon ί\nu\alpha\iota\ 15\degree$, τότε

\begin{alist}
\item να αποδείξετε ότι $O\hat{\varDelta  }E = 30\degree$. \monades{13}
\item να υπολογίσετε τη γωνία $E\hat{O}B = x$. \monades{12}

\begin{center}
\begin{tikzpicture}[scale=1]
\tkzDefPoint(0,0){O}
\tkzDefPoint(3,0){B}
\tkzDefPoint(-3,0){A}
\tkzDrawArc[l3](O,B)(A)
\draw[l3] (A)--(B);

\tkzDefPointBy[homothety=center O ratio 2](A) \tkzGetPoint{C}
\tkzDefPoint(45:3){E} % Just an example E
% Actually $ODE$=30, angle C=15.
% $CD$=$OB$=r.
% Let C = (-6,0). O=(0,0). E is such that angle $ECO$=15.
\tkzDefPoint(-6,0){C}
\tkzDefPointBy[rotation=center C angle 15](O) \tkzGetPoint{line_point}
\tkzInterLC(C,line_point)(O,B) \tkzGetPoints{E}{D} 

\draw[l3] (C)--(E);
\draw[l3] (O)--(D);
\draw[l3] (O)--(E);
\draw[l3] (A)--(C);

\tkzDrawPoints(O,A,B,C,D,E)
\tkzLabelPoint[below](O){$O$}
\tkzLabelPoint[below](A){$A$}
\tkzLabelPoint[below](B){$B$}
\tkzLabelPoint[below](C){$\varGamma $}
\tkzLabelPoint[above left](D){$\varDelta  $}
\tkzLabelPoint[above right](E){$E$}

\tkzMarkAngle[size=1.2,mark=none](O,C,E)
\tkzLabelAngle[pos=1.5](O,C,E){$15^\circ$}
\tkzMarkAngle[size=0.6,mark=none](B,O,E)
\tkzLabelAngle[pos=0.9](B,O,E){$x$}
\tkzMarkSegments[mark=s|](C,D O,B O,A O,D O,E)
\end{tikzpicture}
\end{center}
\end{alist}

\vspace{4mm}

% --- Άσκηση 34506 (34506.tex) ---
\Askhsh{34506}{2}
\begin{ylh}{1}
    \item 3.1. Είδη και στοιχεία τριγώνων
    \item 3.2. 1ο Κριτήριο ισότητας τριγώνων
    \item 4.2. Τέμνουσα δύο ευθειών - Ευκλείδειο αίτημα
    \item 4.6. Άθροισμα γωνιών τριγώνου.
\end{ylh}
% ===================================================================
%  Αρχείο: 34506.tex (Fragment)
%  Αυτόματη μετατροπή μέσω doc2latex.py
% ===================================================================



Δίνεται τετράπλευρο $AB\varGamma \varDelta  $ με $AB$ // $\varDelta  \varGamma $, στο οποίο η διαγώνιος $B\varDelta  $ είναι ίση με την πλευρά $A\varDelta  $. Αν είναι η γωνία $\hat{\varGamma } = 110\degree$ και η γωνία $\varDelta  \hat{B}\varGamma  = 30\degree$, να υπολογίσετε τη γωνία $A\hat{\varDelta  }B$.

\monades{25}

\begin{center}
\begin{tikzpicture}[scale=1]
\tkzDefPoint(0,0){D}
\tkzDefPoint(5,0){C}
\tkzDefTriangle[two angles = 40 and 110](D,C) \tkzGetPoint{B}
\tkzDefLine[parallel=through B](D,C) \tkzGetPoint{line_AB}
\tkzInterCC(D,B)(B,line_AB) \tkzGetPoints{A}{A_alt} 

\draw[l3] (A)--(B)--(C)--(D)--cycle;
\draw[l3] (B)--(D);

\tkzDrawPoints(A,B,C,D)
\tkzLabelPoint[above left](A){$A$}
\tkzLabelPoint[above right](B){$B$}
\tkzLabelPoint[right](C){$\varGamma $}
\tkzLabelPoint[left](D){$\varDelta  $}

\tkzMarkAngle[size=0.6,mark=none](B,C,D)
\tkzLabelAngle[pos=0.9](B,C,D){$110^\circ$}
\tkzMarkAngle[size=1.2,mark=none](C,B,D)
\tkzLabelAngle[pos=1.5](C,B,D){$30^\circ$}
\tkzMarkSegments[mark=s|](A,D B,D)
\end{tikzpicture}
\end{center}

\vspace{4mm}

% --- Άσκηση 34507 (34507.tex) ---
\Askhsh{34507}{2}
\begin{ylh}{1}
    \item 3.1. Είδη και στοιχεία τριγώνων
    \item 3.2. 1ο Κριτήριο ισότητας τριγώνων
    \item 3.6. Κριτήρια ισότητας ορθογώνιων τριγώνων
    \item 3.11. Ανισοτικές σχέσεις πλευρών και γωνιών.
\end{ylh}
% ===================================================================
%  Αρχείο: 34507.tex (Fragment)
%  Αυτόματη μετατροπή μέσω doc2latex.py
% ===================================================================



Θεωρούμε ισοσκελές τρίγωνο $AB\varGamma $ ($AB=A\varGamma $) και τα μέσα Δ, Ε των πλευρών του $AB$, $A\varGamma $ αντίστοιχα. Έστω ότι οι μεσοκάθετες ευθείες των πλευρών $AB$ και $A\varGamma $ τέμνονται στο $M$ και οι οποίες τέμνουν τον φορέα (ε) της βάσης $B\varGamma $ στα σημεία $H$ και Ζ.

\begin{alist}
\item Να συγκρίνετε τα τρίγωνα $\varDelta  BH$ και $EZ\varGamma $. \monades{15}

\item Να αποδείξετε ότι το τρίγωνο $MZH$ είναι ισοσκελές. \monades{10}


\begin{center}
\begin{tikzpicture}[scale=1]
\tkzDefPoint(0,0){B}
\tkzDefPoint(4,0){C}
\tkzDefTriangle[two angles = 70 and 70](B,C) \tkzGetPoint{A}

\tkzDefMidPoint(A,B) \tkzGetPoint{D}
\tkzDefMidPoint(A,C) \tkzGetPoint{E}

\tkzDefLine[orthogonal=through D](A,B) \tkzGetPoint{perpD}
\tkzDefLine[orthogonal=through E](A,C) \tkzGetPoint{perpE}

\tkzInterLL(D,perpD)(E,perpE) \tkzGetPoint{M}
\tkzInterLL(D,perpD)(B,C) \tkzGetPoint{H}
\tkzInterLL(E,perpE)(B,C) \tkzGetPoint{Z}

\draw[l3] (A)--(B)--(C)--cycle;
\draw[l3] (H)--(M)--(Z);
\draw[l3, dashed] (B)--(H);
\draw[l3, dashed] (C)--(Z);

\tkzDrawPoints(A,B,C,D,E,M,H,Z)
\tkzLabelPoint[above](A){$A$}
\tkzLabelPoint[below](B){$B$}
\tkzLabelPoint[below](C){$\varGamma $}
\tkzLabelPoint[left](D){$\varDelta  $}
\tkzLabelPoint[right](E){$E$}
\tkzLabelPoint[above](M){$M$}
\tkzLabelPoint[below](H){$H$}
\tkzLabelPoint[below](Z){$Z$}

\tkzMarkRightAngle(M,D,A)
\tkzMarkRightAngle(M,E,A)
\tkzMarkSegments[mark=s|](A,D D,B A,E E,C)
\tkzMarkSegments[mark=s||](A,B A,C)
\end{tikzpicture}
\end{center}
\end{alist}

\vspace{4mm}

% --- Άσκηση 34509 (34509.tex) ---
\Askhsh{34509}{2}
\begin{ylh}{1}
    \item 3.1. Είδη και στοιχεία τριγώνων
    \item 3.2. 1ο Κριτήριο ισότητας τριγώνων
    \item 3.11. Ανισοτικές σχέσεις πλευρών και γωνιών
    \item 5.11. Ισοσκελές τραπέζιο.
\end{ylh}
% ===================================================================
%  Αρχείο: 34509.tex (Fragment)
%  Αυτόματη μετατροπή μέσω doc2latex.py
% ===================================================================



Δίνεται ισοσκελές τραπέζιο $AB\varGamma \varDelta  $ με $AB$//$\varGamma \varDelta  $ και $AB$\textless $\varGamma \varDelta  $. Θεωρούμε τα σημεία $E$ και $Z$ πάνω στην $AB$ έτσι ώστε $AE$=$EZ$=$ZB$ και έστω Κ το σημείο τομής των $\varDelta   Z$ και $\varGamma E$.

Να αποδείξετε ότι:

\begin{alist}
\item $\varDelta   Z$ = $\varGamma E$, \monades{13}

\item τα τρίγωνα $EKZ$ και $\varDelta  K\varGamma $ είναι ισοσκελή. \monades{12}
\end{alist}

\vspace{4mm}

% --- Άσκηση 34511 (34511.tex) ---
\Askhsh{34511}{2}
\begin{ylh}{1}
    \item 3.2. 1ο Κριτήριο ισότητας τριγώνων
    \item 3.3. 2ο Κριτήριο ισότητας τριγώνων.
\end{ylh}
% ===================================================================
%  Αρχείο: 34511.tex (Fragment)
%  Αυτόματη μετατροπή μέσω doc2latex.py
% ===================================================================



Αν στο παρακάτω σχήμα είναι $\hat{\alpha} = \hat{\delta}$, $\hat{\beta} = \hat{\varGamma }$ και $AB=A\varGamma $, να αποδείξετε ότι:

\begin{alist}
\item τα τρίγωνα Α$B\varDelta  $ και $A\varGamma \varDelta  $ είναι ίσα, \monades{12}

\item οι γωνίες $\hat{\varepsilon}$ και $\hat{\zeta}$ είναι ίσες. \monades{13}


\begin{center}
\begin{tikzpicture}[scale=1]
\tkzDefPoint(0,0){B}
\tkzDefPoint(4,0){C}
\tkzDefTriangle[two angles = 70 and 70](B,C) \tkzGetPoint{A}
\tkzDefPoint(2, -1.5){D}

\draw[l3] (A)--(B)--(D)--(C)--cycle;
\draw[l3] (A)--(D);

\tkzDrawPoints(A,B,C,D)
\tkzLabelPoint[above](A){$A$}
\tkzLabelPoint[left](B){$B$}
\tkzLabelPoint[right](C){$\varGamma $}
\tkzLabelPoint[below](D){$\varDelta  $}

\tkzMarkAngle[size=0.6,mark=none](B,A,D)
\tkzLabelAngle[pos=0.8](B,A,D){$\alpha$}
\tkzMarkAngle[size=0.6,mark=none](D,A,C)
\tkzLabelAngle[pos=0.8](D,A,C){$\delta$}
\tkzMarkAngle[size=0.6,mark=none](A,D,B)
\tkzLabelAngle[pos=0.8](A,D,B){$\beta$}
\tkzMarkAngle[size=0.6,mark=none](C,D,A)
\tkzLabelAngle[pos=0.8](C,D,A){$\varGamma $}
\tkzMarkAngle[size=0.4,mark=none](A,B,D)
\tkzLabelAngle[pos=0.7](A,B,D){$\varepsilon$}
\tkzMarkAngle[size=0.4,mark=none](D,C,A)
\tkzLabelAngle[pos=0.7](D,C,A){$\zeta$}

\tkzMarkSegments[mark=s|](A,B A,C)
\end{tikzpicture}
\end{center}
\end{alist}

\vspace{4mm}

% --- Άσκηση 34512 (34512.tex) ---
\Askhsh{34512}{2}
\begin{ylh}{1}
    \item 5.2. Παραλληλόγραμμα
    \item 5.6. Εφαρμογές στα τρίγωνα.
\end{ylh}
% ===================================================================
%  Αρχείο: 34512.tex (Fragment)
%  Αυτόματη μετατροπή μέσω doc2latex.py
% ===================================================================



Δίνεται παραλληλόγραμμο $AB\varGamma \varDelta  $, του οποίου οι διαγώνιοι $A\varGamma $ και $\varDelta   B$ τέμνονται στο σημείο $O$. Έστω Ε, Ζ, Η και Θ είναι τα μέσα των $O\varDelta  $, $OA$, $OB$ και $O\varGamma $ αντίστοιχα.

\begin{alist}
\item Να αποδείξετε ότι τετράπλευρο $EZH\varTheta $ είναι παραλληλόγραμμο. \monades{15}

\item Αν η περίμετρος του παραλληλογράμμου $AB\varGamma \varDelta  $ είναι 40, να βρείτε την περίμετρο του $E\varTheta  HZ$.

\monades{10}


\begin{center}
\begin{tikzpicture}[scale=1]
\tkzDefPoint(0,0){A}
\tkzDefPoint(5,0){B}
\tkzDefPoint(6,2.5){C}
\tkzDefPoint(1,2.5){D}

\tkzInterLL(A,C)(B,D) \tkzGetPoint{O}
\tkzDefMidPoint(O,D) \tkzGetPoint{E}
\tkzDefMidPoint(O,A) \tkzGetPoint{Z}
\tkzDefMidPoint(O,B) \tkzGetPoint{H}
\tkzDefMidPoint(O,C) \tkzGetPoint{Theta}

\draw[l3] (A)--(B)--(C)--(D)--cycle;
\draw[l3] (A)--(C);
\draw[l3] (B)--(D);
\draw[l3, dashed] (Z)--(E)--(Theta)--(H)--cycle;

\tkzDrawPoints(A,B,C,D,O,E,Z,H,Theta)
\tkzLabelPoint[below left](A){$A$}
\tkzLabelPoint[below right](B){$B$}
\tkzLabelPoint[above right](C){$\varGamma $}
\tkzLabelPoint[above left](D){$\varDelta  $}
\tkzLabelPoint[below](O){$O$}
\tkzLabelPoint[left](Z){$Z$}
\tkzLabelPoint[right](H){$H$}
\tkzLabelPoint[above](E){$E$}
\tkzLabelPoint[below](Theta){$\varTheta $}

\tkzMarkSegments[mark=s|](O,Z Z,A O,H H,B O,E E,D O,Theta Theta,C)
\end{tikzpicture}
\end{center}
\end{alist}

\vspace{4mm}

% --- Άσκηση 34513 (34513.tex) ---
\Askhsh{34513}{2}
\begin{ylh}{1}
    \item 3.7. Κύκλος - Μεσοκάθετος - Διχοτόμος
    \item 5.4. Ρόμβος.
\end{ylh}
% ===================================================================
%  Αρχείο: 34513.tex (Fragment)
%  Αυτόματη μετατροπή μέσω doc2latex.py
% ===================================================================



Σε κύκλο κέντρου Ο, έστω $OA$ μία ακτίνα του. Φέρουμε τη μεσοκάθετη της $OA$ που τέμνει τον κύκλο στα σημεία $B$ και Γ.

Να αποδείξετε ότι:

\begin{alist}
\item το τρίγωνο $OBA$ είναι ισόπλευρο, \monades{13}

\item το τετράπλευρο $OBA\varGamma $ είναι ρόμβος. \monades{12}
\end{alist}

\vspace{4mm}

% --- Άσκηση 34514 (34514.tex) ---
\Askhsh{34514}{2}
\begin{ylh}{1}
    \item 3.1. Είδη και στοιχεία τριγώνων
    \item 3.2. 1ο Κριτήριο ισότητας τριγώνων
    \item 3.7. Κύκλος - Μεσοκάθετος - Διχοτόμος
    \item 3.11. Ανισοτικές σχέσεις πλευρών και γωνιών.
\end{ylh}
% ===================================================================
%  Αρχείο: 34514.tex (Fragment)
%  Αυτόματη μετατροπή μέσω doc2latex.py
% ===================================================================



Έστω κυρτό τετράπλευρο $AB\varGamma \varDelta  $ με $ΒΑ = B\varGamma $ και $\hat{A} = \hat{\varGamma }$.

Να αποδείξετε ότι:

\begin{alist}
\item B$\hat{A}$Γ = B$\hat{\varGamma }$Α, \monades{8}

\item το τρίγωνο $A\varDelta  \varGamma $ είναι ισοσκελές, \monades{10}

\item η ευθεία $B\varDelta  $ είναι μεσοκάθετος του τμήματος $A\varGamma $. \monades{7}


\begin{center}
\begin{tikzpicture}[scale=1]
\tkzDefPoint(0,0){A}
\tkzDefPoint(4,0){C}
\tkzDefPoint(2,1.5){B}
\tkzDefPoint(2,-2){D} % Kite shape

\draw[l3] (A)--(B)--(C)--(D)--cycle;
\draw[l3] (A)--(C);
\draw[l3] (B)--(D);

\tkzDrawPoints(A,B,C,D)
\tkzLabelPoint[left](A){$A$}
\tkzLabelPoint[right](C){$\varGamma $}
\tkzLabelPoint[above](B){$B$}
\tkzLabelPoint[below](D){$\varDelta  $}

\tkzMarkSegments[mark=s|](B,A B,C)
\tkzMarkSegments[mark=s||](D,A D,C)
\tkzMarkAngles[size=0.6](B,A,D B,C,D)
\end{tikzpicture}
\end{center}
\end{alist}

\vspace{4mm}

% --- Άσκηση 34515 (34515.tex) ---
\Askhsh{34515}{2}
\begin{ylh}{1}
    \item 3.1. Είδη και στοιχεία τριγώνων
    \item 3.2. 1ο Κριτήριο ισότητας τριγώνων
    \item 4.6. Άθροισμα γωνιών τριγώνου
    \item 5.9. Μια ιδιότητα του ορθογώνιου τριγώνου.
\end{ylh}
% ===================================================================
%  Αρχείο: 34515.tex (Fragment)
%  Αυτόματη μετατροπή μέσω doc2latex.py
% ===================================================================



Δίνεται γωνία $\overset{\land}{xOy}$ και σημείο $A$ στο εσωτερικό της. Από το Α φέρνουμε τις κάθετες $AB$, $A\varGamma $ προς τις πλευρές Οx , Oy της γωνίας αντίστοιχα, και ονομάζουμε $M$ το μέσο του $OA$.

Να αποδείξετε ότι:

\begin{alist}
\item το τρίγωνο $BMA$ είναι ισοσκελές, \monades{8}

\item το τρίγωνο $BM\varGamma $ είναι ισοσκελές, \monades{8}

\begin{center}
\begin{tikzpicture}[scale=1.2]
\tkzDefPoint(0,0){O}
\tkzDefPoint(5,0){x}
\tkzDefPoint(4,3){y}
\tkzDefPoint(4,1.5){A}

\tkzDefPointBy[projection=onto O--x](A) \tkzGetPoint{B}
\tkzDefPointBy[projection=onto O--y](A) \tkzGetPoint{C}
\tkzDefMidPoint(O,A) \tkzGetPoint{M}

\draw[l3] (O)--(x);
\draw[l3] (O)--(y);
\draw[l3] (B)--(A)--(C);
\draw[l3] (O)--(A);
\draw[l3] (B)--(M)--(C);

\tkzDrawPoints(O,A,B,C,M)
\tkzLabelPoint[below left](O){$O$}
\tkzLabelPoint[below](B){$B$}
\tkzLabelPoint[above left](C){$\varGamma $}
\tkzLabelPoint[right](A){$A$}
\tkzLabelPoint[above left](M){$M$}

\tkzMarkRightAngle(A,B,O)
\tkzMarkRightAngle(A,C,O)
\tkzMarkSegments[mark=s|](M,A M,O M,B M,C)
\end{tikzpicture}
\end{center}
\end{alist}

\vspace{4mm}

% --- Άσκηση 34516 (34516.tex) ---
\Askhsh{34516}{2}
\begin{ylh}{1}
    \item 3.1. Είδη και στοιχεία τριγώνων
    \item 3.2. 1ο Κριτήριο ισότητας τριγώνων
    \item 3.7. Κύκλος - Μεσοκάθετος - Διχοτόμος.
\end{ylh}
% ===================================================================
%  Αρχείο: 34516.tex (Fragment)
%  Αυτόματη μετατροπή μέσω doc2latex.py
% ===================================================================



Αν για το ισοσκελές τρίγωνο $AB\varGamma $ ($AB=A\varGamma $) του σχήματος ισχύουν $\hat{\alpha} = \hat{\beta}$ και $\hat{\varGamma } = \hat{\delta}$, να γράψετε μια απόδειξη για καθέναν από τους ακόλουθους ισχυρισμούς:

\begin{alist}
\item Τα τρίγωνα $AEB$ και $AE\varGamma $ είναι ίσα. \monades{8}

\item Το τρίγωνο $\varGamma EB$ είναι ισοσκελές. \monades{8}

\item Η ευθεία $A\varDelta  $ είναι μεσοκάθετος του τμήματος $B\varGamma $. \monades{9}


\begin{center}
\begin{tikzpicture}[scale=1]
\tkzDefPoint(0,0){B}
\tkzDefPoint(4,0){C}
\tkzDefTriangle[two angles = 70 and 70](B,C) \tkzGetPoint{A}
\tkzDefMidPoint(B,C) \tkzGetPoint{D}
\tkzDefPointBy[homothety=center D ratio 0.4](A) \tkzGetPoint{E}

\draw[l3] (A)--(B)--(C)--cycle;
\draw[l3] (A)--(D);
\draw[l3] (B)--(E)--(C);

\tkzDrawPoints(A,B,C,D,E)
\tkzLabelPoint[above](A){$A$}
\tkzLabelPoint[below](B){$B$}
\tkzLabelPoint[below](C){$\varGamma $}
\tkzLabelPoint[below](D){$\varDelta  $}
\tkzLabelPoint[above right](E){$E$}

\tkzMarkAngle[size=0.6,mark=none](B,A,D)
\tkzLabelAngle[pos=0.8](B,A,D){$\alpha$}
\tkzMarkAngle[size=0.6,mark=none](D,A,C)
\tkzLabelAngle[pos=0.8](D,A,C){$\beta$}
\tkzMarkAngle[size=0.4,mark=none](E,B,C)
\tkzLabelAngle[pos=0.7](E,B,C){$\varGamma $}
\tkzMarkAngle[size=0.4,mark=none](B,C,E)
\tkzLabelAngle[pos=0.7](B,C,E){$\delta$}

\tkzMarkSegments[mark=s|](A,B A,C)
\end{tikzpicture}
\end{center}
\end{alist}

\vspace{4mm}

% --- Άσκηση 34768 (34768.tex) ---
\Askhsh{34768}{2}
\begin{ylh}{1}
    \item 4.2. Τέμνουσα δύο ευθειών - Ευκλείδειο αίτημα
    \item 4.6. Άθροισμα γωνιών τριγώνου
    \item 5.6. Εφαρμογές στα τρίγωνα.
\end{ylh}
% ===================================================================
%  Αρχείο: 34768.tex (Fragment)
%  Αυτόματη μετατροπή μέσω doc2latex.py
% ===================================================================



Δίνεται τρίγωνο $AB\varGamma $ με $\hat{B} = 40^{0}$ και $\hat{\varGamma } = 60^{0}$. Επιπλέον, τα σημεία $\varDelta  $, Ε και $Z$ είναι τα μέσα των πλευρών του $AB$, $B\varGamma $ και $\varGamma  A$ αντίστοιχα

\begin{alist}
\item Να υπολογίσετε τη γωνία $\hat{A}$ του τριγώνου $AB\varGamma $. \monades{9}

\item Να αποδείξετε ότι $\varDelta  E$//$A\varGamma $ και $ZE$//$AB$. \monades{8}

\item Να υπολογίσετε τις γωνίες του τριγώνου $B\varDelta  $Ε. \monades{8}


\begin{center}
\begin{tikzpicture}[scale=1]
\tkzDefPoint(0,0){B}
\tkzDefPoint(5,0){C}
\tkzDefTriangle[two angles = 40 and 60](B,C) \tkzGetPoint{A}

\tkzDefMidPoint(A,B) \tkzGetPoint{D}
\tkzDefMidPoint(B,C) \tkzGetPoint{E}
\tkzDefMidPoint(C,A) \tkzGetPoint{Z}

\draw[l3] (A)--(B)--(C)--cycle;
\draw[l3] (D)--(E)--(Z)--cycle;

\tkzDrawPoints(A,B,C,D,E,Z)
\tkzLabelPoint[above](A){$A$}
\tkzLabelPoint[below](B){$B$}
\tkzLabelPoint[below](C){$\varGamma $}
\tkzLabelPoint[left](D){$\varDelta  $}
\tkzLabelPoint[below](E){$E$}
\tkzLabelPoint[right](Z){$Z$}

\tkzMarkAngle[size=0.6,mark=none](C,B,A)
\tkzLabelAngle[pos=0.9](C,B,A){$40^\circ$}
\tkzMarkAngle[size=0.6,mark=none](A,C,B)
\tkzLabelAngle[pos=0.9](A,C,B){$60^\circ$}
\tkzMarkSegments[mark=s|](A,D D,B)
\tkzMarkSegments[mark=s||](B,E E,C)
\tkzMarkSegments[mark=s|||](C,Z Z,A)
\end{tikzpicture}
\end{center}
\end{alist}

\vspace{4mm}

% --- Άσκηση 34770 (34770.tex) ---
\Askhsh{34770}{2}
\begin{ylh}{1}
    \item 4.2. Τέμνουσα δύο ευθειών - Ευκλείδειο αίτημα
    \item 4.6. Άθροισμα γωνιών τριγώνου.
\end{ylh}
% ===================================================================
%  Αρχείο: 34770.tex (Fragment)
%  Αυτόματη μετατροπή μέσω doc2latex.py
% ===================================================================



Δίνεται ευθεία ε του επιπέδου. Τα παράλληλα τμήματα $AB$ και $\varGamma \varDelta  $ καθώς και ένα τυχαίο σημείο $E$ βρίσκονται στο ίδιο ημιεπίπεδο της ε.

\begin{alist}
\item Αν το $E$ είναι εκτός των τμημάτων $AB$ και $\varGamma \varDelta  $, τότε να αποδείξετε ότι $\hat{\omega} = \hat{\varphi} + \hat{\theta}$. \monades{10}

\begin{center}
\begin{tikzpicture}[scale=1]
\tkzDefPoint(0,1.5){A_end}
\tkzDefPoint(4,1.5){A}
\tkzDefPoint(0,0){C_end}
\tkzDefPoint(4,0){C}
\tkzDefPoint(5.5, 0.75){E}

\draw[l3] (A_end)--(A);
\draw[l3] (C_end)--(C);
\draw[l3] (A)--(E)--(C);

\tkzDrawPoints(A,C,E)
\tkzLabelPoint[above](A){$B$}
\tkzLabelPoint[below](C){$\varDelta  $}
\tkzLabelPoint[right](E){$E$}

\tkzMarkAngle[size=0.6,mark=none](A_end,A,E)
\tkzLabelAngle[pos=0.8](A_end,A,E){$\varphi$}
\tkzMarkAngle[size=0.6,mark=none](E,C,C_end)
\tkzLabelAngle[pos=0.8](E,C,C_end){$\theta$}
\tkzMarkAngle[size=0.5,mark=none](C,E,A)
\tkzLabelAngle[pos=0.8](C,E,A){$\omega$}
\end{tikzpicture}
\end{center}

\item Αν το $E$ είναι ανάμεσα στα τμήματα $AB$ και $\varGamma \varDelta  $ και $EZ$//$AB$, τότε να αποδείξετε ότι $\hat{\theta} = \hat{\omega} + \hat{\varphi}$. \monades{15}

\begin{center}
\begin{tikzpicture}[scale=1]
\tkzDefPoint(0,2){A_start}
\tkzDefPoint(4,2){B}
\tkzDefPoint(0,0){C_start}
\tkzDefPoint(4,0){D}
\tkzDefPoint(2,1){E}
\tkzDefPoint(5,1){Z}

\draw[l3] (A_start)--(B);
\draw[l3] (C_start)--(D);
\draw[l3] (B)--(E)--(D);
\draw[l3, dashed] (E)--(Z);

\tkzDrawPoints(B,D,E,Z)
\tkzLabelPoint[above](B){$B$}
\tkzLabelPoint[below](D){$\varDelta  $}
\tkzLabelPoint[left](E){$E$}
\tkzLabelPoint[right](Z){$Z$}

\tkzMarkAngle[size=0.6,mark=none](A_start,B,E)
\tkzLabelAngle[pos=0.8](A_start,B,E){$\omega$}
\tkzMarkAngle[size=0.6,mark=none](E,D,C_start)
\tkzLabelAngle[pos=0.8](E,D,C_start){$\varphi$}
\tkzMarkAngle[size=0.5,mark=none](D,E,B)
\tkzLabelAngle[pos=0.8](D,E,B){$\theta$}
\end{tikzpicture}
\end{center}
\end{alist}

\vspace{4mm}

% --- Άσκηση 34773 (34773.tex) ---
\Askhsh{34773}{2}
\begin{ylh}{1}
    \item 3.1. Είδη και στοιχεία τριγώνων
    \item 3.2. 1ο Κριτήριο ισότητας τριγώνων
    \item 3.4. 3ο Κριτήριο ισότητας τριγώνων.
\end{ylh}
% ===================================================================
%  Αρχείο: 34773.tex (Fragment)
%  Αυτόματη μετατροπή μέσω doc2latex.py
% ===================================================================



Δίνεται ισοσκελές τρίγωνο $AB\varGamma $ ($AB=Α\varGamma $) και Κ εσωτερικό σημείο του τριγώνου τέτοιο ώστε $KB$ = $K\varGamma $.

\begin{alist}
\item Να αποδείξετε ότι:
\begin{rlist}
\item τα τρίγωνα $BAK$ και $KA\varGamma $ είναι ίσα, \monades{12}
\item η $AK$ είναι διχοτόμος της γωνίας Β$\hat{A}$Γ. \monades{6}
\end{rlist}

\item Αν η προέκταση της $AK$ τέμνει την $B\varGamma $ στο Ε, τότε να δείξετε ότι η $KE$ είναι διάμεσος του τριγώνου $BK\varGamma $. \monades{7}

\begin{center}
\begin{tikzpicture}[scale=1]
\tkzDefPoint(0,0){B}
\tkzDefPoint(4,0){C}
\tkzDefTriangle[two angles = 70 and 70](B,C) \tkzGetPoint{A}
\tkzDefMidPoint(B,C) \tkzGetPoint{E}
\tkzDefPointBy[homothety=center E ratio 0.4](A) \tkzGetPoint{K}

\draw[l3] (A)--(B)--(C)--cycle;
\draw[l3] (A)--(E);
\draw[l3] (B)--(K)--(C);

\tkzDrawPoints(A,B,C,E,K)
\tkzLabelPoint[above](A){$A$}
\tkzLabelPoint[below](B){$B$}
\tkzLabelPoint[below](C){$\varGamma $}
\tkzLabelPoint[below](E){$E$}
\tkzLabelPoint[left](K){$K$}

\tkzMarkSegments[mark=s|](A,B A,C)
\tkzMarkSegments[mark=s||](K,B K,C)
\end{tikzpicture}
\end{center}
\end{alist}

\vspace{4mm}

% --- Άσκηση 34774 (34774.tex) ---
\Askhsh{34774}{2}
\begin{ylh}{1}
    \item 3.1. Είδη και στοιχεία τριγώνων
    \item 3.2. 1ο Κριτήριο ισότητας τριγώνων.
\end{ylh}
% ===================================================================
%  Αρχείο: 34774.tex (Fragment)
%  Αυτόματη μετατροπή μέσω doc2latex.py
% ===================================================================



Δίνεται ισοσκελές τρίγωνο $AB\varGamma $ ($AB=A\varGamma $) και η διάμεσος του $AM$. Στην προέκταση της πλευράς $B\varGamma $ και προς τα δυο της άκρα, θεωρούμε σημεία $\varDelta  $ και $E$ αντίστοιχα έτσι ώστε $B\varDelta  $= $\varGamma E$.

Να αποδείξετε ότι:

\begin{alist}
\item Α$\hat{B}$Δ = Α$\hat{\varGamma }$Ε, \monades{6}

\item τα τρίγωνα Α$B\varDelta  $ και $A\varGamma E$ είναι ίσα, \monades{12}

\item η $AM$ είναι και διάμεσος του τριγώνου $A\varDelta  E$. \monades{7}


\begin{center}
\begin{tikzpicture}[scale=1]
\tkzDefPoint(0,0){B}
\tkzDefPoint(4,0){C}
\tkzDefTriangle[two angles = 70 and 70](B,C) \tkzGetPoint{A}
\tkzDefMidPoint(B,C) \tkzGetPoint{M}

\tkzDefPointBy[homothety=center B ratio -0.3](C) \tkzGetPoint{D}
\tkzDefPointBy[homothety=center C ratio -0.3](B) \tkzGetPoint{E}

\draw[l3] (A)--(D)--(E)--cycle;
\draw[l3] (A)--(B);
\draw[l3] (A)--(C);
\draw[l3] (A)--(M);

\tkzDrawPoints(A,B,C,D,E,M)
\tkzLabelPoint[above](A){$A$}
\tkzLabelPoint[below](B){$B$}
\tkzLabelPoint[below](C){$\varGamma $}
\tkzLabelPoint[below](D){$\varDelta  $}
\tkzLabelPoint[below](E){$E$}
\tkzLabelPoint[below](M){$M$}

\tkzMarkSegments[mark=s|](A,B A,C)
\tkzMarkSegments[mark=s||](B,D C,E)
\tkzMarkSegments[mark=s|||](B,M M,C)
\end{tikzpicture}
\end{center}
\end{alist}

\vspace{4mm}

% --- Άσκηση 34775 (34775.tex) ---
\Askhsh{34775}{2}
\begin{ylh}{1}
    \item 3.1. Είδη και στοιχεία τριγώνων
    \item 3.2. 1ο Κριτήριο ισότητας τριγώνων
    \item 3.4. 3ο Κριτήριο ισότητας τριγώνων
    \item 4.6. Άθροισμα γωνιών τριγώνου.
\end{ylh}
% ===================================================================
%  Αρχείο: 34775.tex (Fragment)
%  Αυτόματη μετατροπή μέσω doc2latex.py
% ===================================================================



Δίνεται ισοσκελές τρίγωνο $AB\varGamma $ με $AB=A\varGamma $ και $\hat{A} = 80^{0}$. Έστω Κ σημείο της διχοτόμου της γωνίας $\hat{A}$, τέτοιο ώστε $KB$=$KA$=$K\varGamma $.

\begin{alist}
\item Να αποδείξετε ότι τα τρίγωνα $BKA$ και $\varGamma KA$ είναι ίσα. \monades{10}

\item Να υπολογίσετε:
\begin{rlist}
\item τις γωνίες $A\hat{B}K$ και $A\hat{\varGamma }K$, \monades{8}
\item τη γωνία $B\hat{K}\varGamma $. \monades{7}
\end{rlist}

\begin{center}
\begin{tikzpicture}[scale=1]
\tkzDefPoint(0,0){B}
\tkzDefPoint(4,0){C}
\tkzDefTriangle[two angles = 50 and 50](B,C) \tkzGetPoint{A}
\tkzDefCircle[circum](A,B,C) \tkzGetPoint{K}

\draw[l3] (A)--(B)--(C)--cycle;
\draw[l3] (A)--(K);
\draw[l3] (B)--(K);
\draw[l3] (C)--(K);

\tkzDrawPoints(A,B,C,K)
\tkzLabelPoint[above](A){$A$}
\tkzLabelPoint[below](B){$B$}
\tkzLabelPoint[below](C){$\varGamma $}
\tkzLabelPoint[below](K){$K$}

\tkzMarkSegments[mark=s|](A,B A,C)
\tkzMarkSegments[mark=s||](K,A K,B K,C)
\end{tikzpicture}
\end{center}
\end{alist}

\vspace{4mm}

% --- Άσκηση 34776 (34776.tex) ---
\Askhsh{34776}{2}
\begin{ylh}{1}
    \item 3.1. Είδη και στοιχεία τριγώνων
    \item 3.11. Ανισοτικές σχέσεις πλευρών και γωνιών
    \item 4.2. Τέμνουσα δύο ευθειών - Ευκλείδειο αίτημα.
\end{ylh}
% ===================================================================
%  Αρχείο: 34776.tex (Fragment)
%  Αυτόματη μετατροπή μέσω doc2latex.py
% ===================================================================



Δίνεται ορθογώνιο τρίγωνο $AB\varGamma $ ($\hat{A}$ = 90\textsuperscript{ο}). Έστω $\varDelta$ σημείο της πλευράς $A\varGamma $ τέτοιο ώστε, η διχοτόμος $\varDelta  E$ της γωνίας Α$\hat{\varDelta  }$Β να είναι παράλληλη στην πλευρά $B\varGamma $.

\begin{alist}
\item Να αποδείξετε ότι:
\begin{rlist}
\item $E\hat{\varDelta  }B = \varDelta  \hat{B}\varGamma $ και $E\hat{\varDelta  }A = \hat{\varGamma }$, \monades{8}
\item το τρίγωνο $B\varDelta  \varGamma $ είναι ισοσκελές. \monades{8}
\end{rlist}

\item Αν είναι $A\hat{\varDelta  }B = 60\degree$, τότε να υπολογίσετε τη γωνία $\hat{\varGamma }$. \monades{9}

\begin{center}
\begin{tikzpicture}[scale=1]
\tkzDefPoint(0,0){A}
\tkzDefPoint(4,0){C}
\tkzDefPoint(0,3){B}
\tkzDefPoint(0,1.5){D} % dummy D

\tkzDefLine[parallel=through D](B,C) \tkzGetPoint{line_DE}
\tkzInterLL(D,line_DE)(A,B) \tkzGetPoint{E}

% We need $DE$ bisector of $ADB$? No, $ADB$ is fixed by D.
% If we want to $DRAW$ the condition "$DE$ bisector of $ADB$ $AND$ $DE$ // $BC$":
% This is only true for specific D.
% E is intersection of $DE$ and $AB$. 
% For diagram purposes, we'll pick D such that it looks right.
\tkzDefPoint(0,1.2){D}
\tkzDefLine[parallel=through D](B,C) \tkzGetPoint{DE_dir}
\tkzInterLL(D,DE_dir)(A,B) \tkzGetPoint{E}

\draw[l3] (A)--(B)--(C)--cycle;
\draw[l3] (B)--(D);
\draw[l3] (D)--(E);

\tkzDrawPoints(A,B,C,D,E)
\tkzLabelPoint[below left](A){$A$}
\tkzLabelPoint[above](B){$B$}
\tkzLabelPoint[below](C){$\varGamma $}
\tkzLabelPoint[left](D){$\varDelta  $}
\tkzLabelPoint[above left](E){$E$}

\tkzMarkRightAngle(B,A,C)
\tkzMarkSegments[mark=s|](D,B D,C)
\end{tikzpicture}
\end{center}
\end{alist}

\vspace{4mm}

% --- Άσκηση 34777 (34777.tex) ---
\Askhsh{34777}{2}
\begin{ylh}{1}
    \item 3.1. Είδη και στοιχεία τριγώνων
    \item 3.2. 1ο Κριτήριο ισότητας τριγώνων
    \item 4.2. Τέμνουσα δύο ευθειών - Ευκλείδειο αίτημα.
\end{ylh}
% ===================================================================
%  Αρχείο: 34777.tex (Fragment)
%  Αυτόματη μετατροπή μέσω doc2latex.py
% ===================================================================



Δίνεται ισοσκελές τρίγωνο $AB\varGamma $ ($AB=A\varGamma $) και το ύψος του $AM$. Φέρουμε ημιευθεία Γx κάθετη στη $B\varGamma $, προς το ημιεπίπεδο που δεν ανήκει το Α, και παίρνουμε σε αυτήν τμήμα $\varGamma \varDelta  $ = $AB$.

Να αποδείξετε ότι:

\begin{alist}
\item η γωνία Δ$\hat{A}$Γ είναι ίση με τη γωνία Γ$\hat{\varDelta  }$Α, \monades{12}

\item $\varGamma \varDelta  $ // $AM$, \monades{6}

\item η $A\varDelta  $ είναι διχοτόμος της γωνίας Μ$\hat{A}$Γ. \monades{7}


\begin{center}
\begin{tikzpicture}[scale=1]
\tkzDefPoint(0,0){B}
\tkzDefPoint(4,0){C}
\tkzDefTriangle[two angles = 70 and 70](B,C) \tkzGetPoint{A}
\tkzDefMidPoint(B,C) \tkzGetPoint{M}

\tkzDefPointBy[homothety=center C ratio 1](B) \tkzGetPoint{temp}
\tkzDefPointBy[rotation=center C angle -90](temp) \tkzGetPoint{D_dir}
\tkzCalcLength(A,B)\tkzGetLength{rAB}
\tkzInterLC[R](C,D_dir)(C,\rAB) \tkzGetPoints{D}{D_alt}

\draw[l3] (A)--(B)--(C)--cycle;
\draw[l3] (A)--(M);
\draw[l3] (C)--(D);
\draw[l3] (A)--(D);

\tkzDrawPoints(A,B,C,D,M)
\tkzLabelPoint[above](A){$A$}
\tkzLabelPoint[below](B){$B$}
\tkzLabelPoint[below](C){$\varGamma $}
\tkzLabelPoint[below](M){$M$}
\tkzLabelPoint[below](D){$\varDelta  $}

\tkzMarkRightAngle(A,M,C)
\tkzMarkRightAngle(B,C,D)
\tkzMarkSegments[mark=s|](A,B A,C C,D)
\end{tikzpicture}
\end{center}
\end{alist}

\vspace{4mm}

% --- Άσκηση 34778 (34778.tex) ---
\Askhsh{34778}{2}
\begin{ylh}{1}
    \item 3.11. Ανισοτικές σχέσεις πλευρών και γωνιών
    \item 4.2. Τέμνουσα δύο ευθειών - Ευκλείδειο αίτημα
    \item 4.6. Άθροισμα γωνιών τριγώνου.
\end{ylh}
% ===================================================================
%  Αρχείο: 34778.tex (Fragment)
%  Αυτόματη μετατροπή μέσω doc2latex.py
% ===================================================================



Δίνεται τρίγωνο $AB\varGamma $ με $AB$ < $A\varGamma $. Έστω Αx η διχοτόμος της εξωτερικής του γωνίας $\hat{A}$\textsubscript{εξ}, όπου

$\hat{A}$\textsubscript{εξ} = 120\textsuperscript{ο}. Από την κορυφή $B$ φέρνουμε ευθεία παράλληλη στην Αx, η οποία τέμνει την πλευρά $A\varGamma $ στο σημείο $\varDelta  $.

\begin{alist}
\item Να αποδείξετε ότι:
\begin{rlist}
\item $\varDelta  \hat{B}A = 60^{\circ}$, \monades{9}
\item το τρίγωνο $AB\varDelta  $ είναι ισόπλευρο, \monades{9}
\end{rlist}

\item Αν η γωνία $B\hat{\varDelta  }A$ είναι διπλάσια της γωνίας $\hat{\varGamma }$ του τριγώνου $AB\varGamma $, να υπολογίσετε τις γωνίες του τριγώνου $B\varDelta  \varGamma $. \monades{7}

\begin{center}
\begin{tikzpicture}[scale=1]
\tkzDefPoint(0,0){A}
\tkzDefPoint(4,0){C}
\tkzDefPoint(-2,0){A_ext}
\tkzDefTriangle[two angles = 60 and 30](A,C) \tkzGetPoint{B_dummy}
% We need angle A = 60. 
\tkzDefPoint(60:3){B}
\tkzDefPoint(120:3){x_pt} % Ax bisector of external angle (120 deg)

\tkzDefLine[parallel=through B](A,x_pt) \tkzGetPoint{BD_dir}
\tkzInterLL(B,BD_dir)(A,C) \tkzGetPoint{D}

\draw[l3] (A)--(B)--(C)--cycle;
\draw[l3] (B)--(D);
\draw[l3, dashed] (A)--(x_pt);
\draw[l3] (A)--(A_ext);

\tkzDrawPoints(A,B,C,D)
\tkzLabelPoint[below](A){$A$}
\tkzLabelPoint[above](B){$B$}
\tkzLabelPoint[below](C){$\varGamma $}
\tkzLabelPoint[below](D){$\varDelta  $}
\tkzLabelPoint[above left](x_pt){$x$}

\tkzMarkAngle[size=0.6,mark=none](x_pt,A,A_ext)
\tkzLabelAngle[pos=0.8](x_pt,A,A_ext){$60^\circ$}
\tkzMarkAngle[size=0.6,mark=none](B,A,x_pt)
\tkzLabelAngle[pos=0.8](B,A,x_pt){$60^\circ$}
\end{tikzpicture}
\end{center}
\end{alist}

\vspace{4mm}

% --- Άσκηση 34779 (34779.tex) ---
\Askhsh{34779}{2}
\begin{ylh}{1}
    \item 3.2. 1ο Κριτήριο ισότητας τριγώνων
    \item 4.2. Τέμνουσα δύο ευθειών - Ευκλείδειο αίτημα.
\end{ylh}
% ===================================================================
%  Αρχείο: 34779.tex (Fragment)
%  Αυτόματη μετατροπή μέσω doc2latex.py
% ===================================================================



Στις προεκτάσεις των πλευρών $BA$ (προς το Α) και $\varGamma  A$ (προς το Α) τριγώνου $AB\varGamma $ παίρνουμε τα τμήματα $A\varDelta  $=$AB$ και $AE$=$A\varGamma $.

Να αποδείξετε ότι:

\begin{alist}
\item τα τρίγωνα $AB\varGamma $ και $A\varDelta  E$ είναι ίσα, \monades{12}

\item $E\varDelta  $//$B\varGamma $. \monades{13}
\end{alist}

\vspace{4mm}

% --- Άσκηση 34780 (34780.tex) ---
\Askhsh{34780}{2}
\begin{ylh}{1}
    \item 3.2. 1ο Κριτήριο ισότητας τριγώνων
    \item 3.3. 2ο Κριτήριο ισότητας τριγώνων.
\end{ylh}
% ===================================================================
%  Αρχείο: 34780.tex (Fragment)
%  Αυτόματη μετατροπή μέσω doc2latex.py
% ===================================================================



Στις προεκτάσεις των πλευρών $BA$ και $\varGamma  A$ (προς το Α) τριγώνου $AB\varGamma $ παίρνουμε τα τμήματα

$A\varDelta  $ = $AB$ και $AE$ = $A\varGamma $.

\begin{alist}
\item Να αποδείξετε ότι $\varDelta  E$ = $B\varGamma $. \monades{12}

\item Αν $AM$ είναι η διάμεσος του τριγώνου $AB\varGamma $ και η προέκταση της $AM$ τέμνει την $E\varDelta  $ στο Ζ, να αποδείξετε ότι:


\end{alist}
\begin{alist}
\item
  τα τρίγωνα Α$\varDelta   Z$ και $ABM$ είναι ίσα, \monades{7}
\item
  $Z\varDelta  $ = $\frac{\varDelta   E}{2}$. \monades{6}
\end{alist}

\vspace{4mm}

% --- Άσκηση 34781 (34781.tex) ---
\Askhsh{34781}{2}
\begin{ylh}{1}
    \item 3.4. 3ο Κριτήριο ισότητας τριγώνων
    \item 3.6. Κριτήρια ισότητας ορθογώνιων τριγώνων
    \item 5.3. Ορθογώνιο.
\end{ylh}
% ===================================================================
%  Αρχείο: 34781.tex (Fragment)
%  Αυτόματη μετατροπή μέσω doc2latex.py
% ===================================================================



Σε ορθογώνιο $AB\varGamma \varDelta  $ θεωρούμε τα μέσα $M$ και $N$ των πλευρών του $AB$ και $\varGamma \varDelta  $ αντίστοιχα.

Να αποδείξετε ότι:

\begin{alist}
\item $M\varDelta  $=$M\varGamma $, \monades{12}

\item η ευθεία που ορίζουν τα σημεία $M$ και $N$ είναι μεσοκάθετος του τμήματος $\varGamma \varDelta  $.

\monades{13}
\end{alist}

\vspace{4mm}

% --- Άσκηση 34783 (34783.tex) ---
\Askhsh{34783}{2}
\begin{ylh}{1}
    \item 3.6. Κριτήρια ισότητας ορθογώνιων τριγώνων
    \item 4.2. Τέμνουσα δύο ευθειών - Ευκλείδειο αίτημα
    \item 5.2. Παραλληλόγραμμα.
\end{ylh}
% ===================================================================
%  Αρχείο: 34783.tex (Fragment)
%  Αυτόματη μετατροπή μέσω doc2latex.py
% ===================================================================



Θεωρούμε παραλληλόγραμμο $AB\varGamma \varDelta  $ και τμήματα $AA$΄, $\varGamma \varGamma $΄ κάθετα στη διαγώνιο $B\varDelta  $ από τις κορυφές Α, Γ αντίστοιχα. Αν τα σημεία $A$ ́ και Γ΄ δεν ταυτίζονται, να αποδείξετε ότι:

\begin{alist}
\item $AA$' // $\varGamma \varGamma $', \monades{8}

\item $AA$΄= $\varGamma \varGamma $΄, (Μονάδες10)

\item Το τετράπλευρο $A\varGamma $΄$\varGamma  A$ ́ είναι παραλληλόγραμμο. \monades{7}


\begin{center}
\begin{tikzpicture}[scale=1]
\tkzDefPoint(0,0){A}
\tkzDefPoint(4,0){B}
\tkzDefPoint(5,2){C}
\tkzDefPoint(1,2){D}

\draw[l3] (A)--(B)--(C)--(D)--cycle;
\draw[l3] (B)--(D);

\tkzDefPointBy[projection=onto B--D](A) \tkzGetPoint{Ap}
\tkzDefPointBy[projection=onto B--D](C) \tkzGetPoint{Cp}

\draw[l3] (A)--(Ap);
\draw[l3] (C)--(Cp);
\draw[l3, dashed] (A)--(Cp)--(C)--(Ap)--cycle;

\tkzDrawPoints(A,B,C,D,Ap,Cp)
\tkzLabelPoint[left](A){$A$}
\tkzLabelPoint[right](B){$B$}
\tkzLabelPoint[right](C){$\varGamma $}
\tkzLabelPoint[left](D){$\varDelta  $}
\tkzLabelPoint[above left](Ap){$A'$}
\tkzLabelPoint[below right](Cp){$\varGamma '$}

\tkzMarkRightAngle(A,Ap,D)
\tkzMarkRightAngle(C,Cp,B)
\end{tikzpicture}
\end{center}
\end{alist}

\vspace{4mm}

% --- Άσκηση 34784 (34784.tex) ---
\Askhsh{34784}{2}
\begin{ylh}{1}
    \item 3.1. Είδη και στοιχεία τριγώνων
    \item 3.4. 3ο Κριτήριο ισότητας τριγώνων.
\end{ylh}
% ===================================================================
%  Αρχείο: 34784.tex (Fragment)
%  Αυτόματη μετατροπή μέσω doc2latex.py
% ===================================================================



Θεωρούμε ισοσκελές τρίγωνο $AB\varGamma $ ($AB=Α\varGamma $) και σημείο $M$ εσωτερικό του τριγώνου τέτοιο ώστε $MB$ = $M\varGamma $.

Να αποδείξετε ότι:

\begin{alist}
\item τα τρίγωνα $AMB$ και $AM\varGamma $ είναι ίσα \monades{12}

\item η ευθεία $AM$ διχοτομεί τη γωνία Β$\hat{Μ}$Γ. \monades{13}
\end{alist}

\vspace{4mm}

% --- Άσκηση 34785 (34785.tex) ---
\Askhsh{34785}{2}
\begin{ylh}{1}
    \item 3.2. 1ο Κριτήριο ισότητας τριγώνων
    \item 4.6. Άθροισμα γωνιών τριγώνου.
\end{ylh}
% ===================================================================
%  Αρχείο: 34785.tex (Fragment)
%  Αυτόματη μετατροπή μέσω doc2latex.py
% ===================================================================



Δίνεται τρίγωνο ισοσκελές $AB\varGamma $ ($AB=A\varGamma $) με γωνία $\hat{A} = 50^{0}$ . Έστω $\varDelta$ είναι σημείο της πλευράς $A\varGamma $, τέτοιο ώστε $B\varDelta  $=$B\varGamma $.

\begin{alist}
\item Να υπολογίσετε τις γωνίες $\hat{B}$ και $\hat{\varGamma }$του τριγώνου $AB\varGamma $, \monades{12}

\item Να αποδείξετε ότι η γωνία $\varDelta  \hat{B}\varGamma $ είναι ίση με τη γωνία $\hat{A}$. \monades{13}


\begin{center}
\begin{tikzpicture}[scale=1]
\tkzDefPoint(0,0){B}
\tkzDefPoint(4,0){C}
\tkzDefTriangle[two angles = 65 and 65](B,C) \tkzGetPoint{A}

\tkzInterLC(A,C)(B,C) \tkzGetPoints{D_dummy}{D}

\draw[l3] (A)--(B)--(C)--cycle;
\draw[l3] (B)--(D);

\tkzDrawPoints(A,B,C,D)
\tkzLabelPoint[above](A){$A$}
\tkzLabelPoint[below](B){$B$}
\tkzLabelPoint[below](C){$\varGamma $}
\tkzLabelPoint[right](D){$\varDelta  $}

\tkzMarkSegments[mark=s|](A,B A,C)
\tkzMarkSegments[mark=s||](B,D B,C)
\tkzMarkAngle[size=0.6,mark=none](C,B,A)
\tkzLabelAngle[pos=0.9](C,B,A){$65^\circ$}
\tkzMarkAngle[size=0.6,mark=none](A,C,B)
\tkzLabelAngle[pos=0.9](A,C,B){$65^\circ$}
\end{tikzpicture}
\end{center}
\end{alist}

\vspace{4mm}

% --- Άσκηση 34786 (34786.tex) ---
\Askhsh{34786}{2}
\begin{ylh}{1}
    \item 4.6. Άθροισμα γωνιών τριγώνου.
\end{ylh}
% ===================================================================
%  Αρχείο: 34786.tex (Fragment)
%  Αυτόματη μετατροπή μέσω doc2latex.py
% ===================================================================



Θεωρούμε ορθογώνιο τρίγωνο $AB\varGamma $ ($\hat{A} = 90^{0}$) με $\hat{\varGamma } = 40^{0}$. Έστω $\varDelta$ τυχαίο σημείο της πλευράς $A\varGamma $ και $\varDelta   E\bot B\varGamma $.

Να υπολογίσετε:

\begin{alist}
\item τις γωνίες του τριγώνου $\varDelta  E\varGamma $, \monades{10}

\item τις γωνίες του τετραπλεύρου $A\varDelta  EB$. \monades{15}


\begin{center}
\begin{tikzpicture}[scale=1.2]
\tkzDefPoint(0,0){A}
\tkzDefPoint(4,0){C}
\tkzDefPoint(0,3.35){B} % Using tan(40) ~ 0.839. $AC$*tan(40) = 4 * 0.839 = 3.356

\tkzDefPoint(2,0){D}
\tkzDefPointBy[projection=onto B--C](D) \tkzGetPoint{E}

\draw[l3] (A)--(B)--(C)--cycle;
\draw[l3] (D)--(E);

\tkzDrawPoints(A,B,C,D,E)
\tkzLabelPoint[below left](A){$A$}
\tkzLabelPoint[above](B){$B$}
\tkzLabelPoint[below](C){$\varGamma $}
\tkzLabelPoint[below](D){$\varDelta  $}
\tkzLabelPoint[above right](E){$E$}

\tkzMarkRightAngle(B,A,C)
\tkzMarkRightAngle(D,E,C)
\tkzMarkAngle[size=0.6,mark=none](B,C,A)
\tkzLabelAngle[pos=0.9](B,C,A){$40^\circ$}
\end{tikzpicture}
\end{center}
\end{alist}

\vspace{4mm}

% --- Άσκηση 34787 (34787.tex) ---
\Askhsh{34787}{2}
\begin{ylh}{1}
    \item 3.2. 1ο Κριτήριο ισότητας τριγώνων
    \item 4.6. Άθροισμα γωνιών τριγώνου.
\end{ylh}
% ===================================================================
%  Αρχείο: 34787.tex (Fragment)
%  Αυτόματη μετατροπή μέσω doc2latex.py
% ===================================================================



Θεωρούμε ισοσκελές τρίγωνο $AB\varGamma $ ($AB=A\varGamma $) με γωνία κορυφής $\hat{A} = 40^{0}$. Στην προέκταση της $\varGamma B$ (προς το Β ) παίρνουμε τμήμα $B\varDelta  $ τέτοιο ώστε $B\varDelta   = ΑΒ$.

Να υπολογίσετε:

\begin{alist}
\item τις γωνίες του τριγώνου $AB\varGamma $, \monades{10}

\item τη γωνία $\varDelta  \hat{A}\varGamma $. \monades{15}
\end{alist}

\vspace{4mm}

% --- Άσκηση 36086 (36086.tex) ---
\Askhsh{36086}{2}
\begin{ylh}{1}
    \item 3.1. Είδη και στοιχεία τριγώνων
    \item 3.2. 1ο Κριτήριο ισότητας τριγώνων
    \item 4.6. Άθροισμα γωνιών τριγώνου
    \item 5.9. Μια ιδιότητα του ορθογώνιου τριγώνου.
\end{ylh}
% ===================================================================
%  Αρχείο: 36086.tex (Fragment)
%  Αυτόματη μετατροπή μέσω doc2latex.py
% ===================================================================



Θεωρούμε ορθογώνιο τρίγωνο $AB\varGamma $ ($\hat{A} = 90^{0}$) με γωνία $\hat{B} = 2\hat{\varGamma }$ και το μέσο $M$ της πλευράς του $B\varGamma $.

\begin{alist}
\item Να υπολογίσετε τις γωνίες $\hat{B}$ και $\hat{\varGamma }$ του τριγώνου $AB\varGamma $. \monades{8}

\item Να δείξετε ότι το τρίγωνο $AM\varGamma $ είναι ισοσκελές. \monades{8}

\item Να υπολογίσετε τη γωνία Α$\hat{Μ}$Γ. \monades{9}
\end{alist}

\vspace{4mm}

% --- Άσκηση 36087 (36087.tex) ---
\Askhsh{36087}{2}
\begin{ylh}{1}
    \item 3.1. Είδη και στοιχεία τριγώνων
    \item 3.2. 1ο Κριτήριο ισότητας τριγώνων
    \item 3.11. Ανισοτικές σχέσεις πλευρών και γωνιών
    \item 4.6. Άθροισμα γωνιών τριγώνου.
\end{ylh}
% ===================================================================
%  Αρχείο: 36087.tex (Fragment)
%  Αυτόματη μετατροπή μέσω doc2latex.py
% ===================================================================



Στο παρακάτω σχήμα ισχύουν $\varDelta   B$=$BA$=$A\varGamma $=$\varGamma E$ και $Β\hat{A}\varGamma  = 40^{0}$.

Να αποδείξετε ότι:

\begin{alist}
\item $Α\hat{B}\varDelta   = Α\hat{\varGamma }E = 110^{0}$, \monades{10}

\item τα τρίγωνα Α$B\varDelta  $ και $A\varGamma E$ είναι ίσα, \monades{10}

\item το τρίγωνο $\varDelta  AE$ είναι ισοσκελές. \monades{5}


\begin{center}
\begin{tikzpicture}[scale=0.8]
\tkzDefPoint(0,0){A}
\tkzDefPoint(-2,-4){B}
\tkzDefPoint(2,-4){C}

\tkzDefPointBy[homothety=center A ratio 2](B) \tkzGetPoint{D}
\tkzDefPointBy[homothety=center A ratio 2](C) \tkzGetPoint{E}

\draw[l3] (D)--(A)--(E);
\draw[l3] (B)--(C);
\draw[l3] (D)--(C); % Maybe D-C? No, usually D-A-E is the focus.
\draw[l3] (D)--(E); % and D-E

\tkzDrawPoints(A,B,C,D,E)
\tkzLabelPoint[above](A){$A$}
\tkzLabelPoint[left](B){$B$}
\tkzLabelPoint[right](C){$\varGamma $}
\tkzLabelPoint[below left](D){$\varDelta  $}
\tkzLabelPoint[below right](E){$E$}

\tkzMarkSegments[mark=s|](D,B B,A A,C C,E)
\tkzMarkAngle[size=0.6,mark=none](C,A,B)
\tkzLabelAngle[pos=0.8](C,A,B){$40^\circ$}
\end{tikzpicture}
\end{center}
\end{alist}

\vspace{4mm}

% --- Άσκηση 36088 (36088.tex) ---
\Askhsh{36088}{2}
\begin{ylh}{1}
    \item 3.1. Είδη και στοιχεία τριγώνων
    \item 3.11. Ανισοτικές σχέσεις πλευρών και γωνιών
    \item 4.6. Άθροισμα γωνιών τριγώνου
    \item 5.6. Εφαρμογές στα τρίγωνα.
\end{ylh}
% ===================================================================
%  Αρχείο: 36088.tex (Fragment)
%  Αυτόματη μετατροπή μέσω doc2latex.py
% ===================================================================



Δίνεται τρίγωνο $AB\varGamma $ με $\hat{A} = 40^{0}$ και $\hat{B} = 70^{0}$. Τα σημεία $\varDelta  $ και $E$ είναι τα μέσα των $AB$ και $A\varGamma $ με $\varDelta   E = 9$ και $E\varGamma  = 16$.

\begin{alist}
\item Να αποδείξετε ότι:
\begin{rlist}
\item το τρίγωνο $AB\varGamma $ είναι ισοσκελές και να βρείτε ποιες είναι οι ίσες πλευρές του, \monades{8}
\item $B\varGamma $ = 18. \monades{8}
\end{rlist}

\item Να υπολογίσετε την περίμετρο του τριγώνου $AB\varGamma $. \monades{9}

\begin{center}
\begin{tikzpicture}[scale=0.15] % Large dimensions (32, 18), scaled down
\tkzDefPoint(0,0){B}
\tkzDefPoint(18,0){C}
\tkzDefTriangle[two angles = 70 and 70](B,C) \tkzGetPoint{A}

\tkzDefMidPoint(A,B) \tkzGetPoint{D}
\tkzDefMidPoint(A,C) \tkzGetPoint{E}

\draw[l3] (A)--(B)--(C)--cycle;
\draw[l3] (D)--(E);

\tkzDrawPoints(A,B,C,D,E)
\tkzLabelPoint[above](A){$A$}
\tkzLabelPoint[below](B){$B$}
\tkzLabelPoint[below](C){$\varGamma $}
\tkzLabelPoint[left](D){$\varDelta  $}
\tkzLabelPoint[right](E){$E$}

\tkzMarkAngle[size=4,mark=none](C,B,A)
\tkzLabelAngle[pos=6](C,B,A){$70^\circ$}
\tkzMarkAngle[size=4,mark=none](A,C,B)
\tkzLabelAngle[pos=6](A,C,B){$70^\circ$}
\tkzMarkSegments[mark=s|](D,E)
\tkzLabelSegment[above](D,E){$9$}
\tkzMarkSegments[mark=s||](E,C)
\tkzLabelSegment[right](E,C){$16$}
\end{tikzpicture}
\end{center}
\end{alist}

\vspace{4mm}

% --- Άσκηση 36089 (36089.tex) ---
\Askhsh{36089}{2}
\begin{ylh}{1}
    \item 3.3. 2ο Κριτήριο ισότητας τριγώνων
    \item 5.2. Παραλληλόγραμμα.
\end{ylh}
% ===================================================================
%  Αρχείο: 36089.tex (Fragment)
%  Αυτόματη μετατροπή μέσω doc2latex.py
% ===================================================================



Θεωρούμε παραλληλόγραμμο $AB\varGamma \varDelta  $ . Αν οι διχοτόμοι των απέναντι γωνιών $\hat{\varDelta  }$ και $\hat{B}$ τέμνουν τις πλευρές $AB$ και $\varGamma \varDelta  $ στα σημεία $E$ και $Z$ αντίστοιχα, να αποδείξετε ότι:

\begin{alist}
\item τα τρίγωνα Α$E\varDelta  $ και $B\varGamma $Ζ είναι ίσα, \monades{12}

\item το τετράπλευρο $\varDelta   EBZ$ είναι παραλληλόγραμμο. \monades{13}
\end{alist}

\vspace{4mm}

% --- Άσκηση 36090 (36090.tex) ---
\Askhsh{36090}{2}
\begin{ylh}{1}
    \item 3.3. 2ο Κριτήριο ισότητας τριγώνων
    \item 4.2. Τέμνουσα δύο ευθειών - Ευκλείδειο αίτημα
    \item 5.2. Παραλληλόγραμμα.
\end{ylh}
% ===================================================================
%  Αρχείο: 36090.tex (Fragment)
%  Αυτόματη μετατροπή μέσω doc2latex.py
% ===================================================================



Στις πλευρές $A\varDelta  $ και $B\varGamma $ παραλληλογράμμου $AB\varGamma \varDelta  $ θεωρούμε σημεία $E$ και Z, τέτοια ώστε

$AE$ = $\varGamma Z$. Αν η ευθεία $ZE$ τέμνει τις προεκτάσεις των πλευρών $AB$ και $\varGamma \varDelta  $ στα σημεία $H$ και Θ, να αποδείξετε ότι:

\begin{alist}
\item $H\hat{B}Ζ = Ε\hat{\varDelta  }\varTheta $, \monades{8}

\item Β$\hat{Ζ}H = \varDelta  \hat{Ε}\varTheta $, \monades{8}

\item $BH$ = $\varTheta\varDelta  $. \monades{9}


\begin{center}
\begin{tikzpicture}[scale=0.8]
\tkzDefPoint(0,0){A}
\tkzDefPoint(4,0){B}
\tkzDefPoint(5,2.5){C}
\tkzDefPoint(1,2.5){D}

\tkzDefPointBy[homothety=center A ratio 0.4](D) \tkzGetPoint{E}
\tkzDefPointBy[homothety=center C ratio 0.4](B) \tkzGetPoint{Z}

\tkzInterLL(E,Z)(A,B) \tkzGetPoint{H}
\tkzInterLL(E,Z)(C,D) \tkzGetPoint{Theta}

\draw[l3] (A)--(B)--(C)--(D)--cycle;
\draw[l3] (H)--(Theta);
\draw[l3, dashed] (A)--(H);
\draw[l3, dashed] (D)--(Theta);

\tkzDrawPoints(A,B,C,D,E,Z,H,Theta)
\tkzLabelPoint[above left](A){$A$}
\tkzLabelPoint[above right](B){$B$}
\tkzLabelPoint[below right](C){$\varGamma $}
\tkzLabelPoint[below left](D){$\varDelta  $}
\tkzLabelPoint[left](E){$E$}
\tkzLabelPoint[right](Z){$Z$}
\tkzLabelPoint[above](H){$H$}
\tkzLabelPoint[below](Theta){$\varTheta $}

\tkzMarkSegments[mark=s|](A,E C,Z)
\end{tikzpicture}
\end{center}
\end{alist}

\vspace{4mm}

% --- Άσκηση 36091 (36091.tex) ---
\Askhsh{36091}{2}
\begin{ylh}{1}
    \item 4.2. Τέμνουσα δύο ευθειών - Ευκλείδειο αίτημα
    \item 4.6. Άθροισμα γωνιών τριγώνου
    \item 5.6. Εφαρμογές στα τρίγωνα.
\end{ylh}
% ===================================================================
%  Αρχείο: 36091.tex (Fragment)
%  Αυτόματη μετατροπή μέσω doc2latex.py
% ===================================================================



Δίνεται τρίγωνο $AB\varGamma $ με $\hat{B}$ = 50\degree. Έστω ότι τα σημεία $\varDelta  $ και $E$ είναι τα μέσα των πλευρών $B\varGamma $ και $A\varGamma $ αντίστοιχα, τέτοια ώστε $\varDelta  \hat{Ε}\varGamma $ = 70\degree
\begin{alist}
\item Να δικαιολογήσετε γιατί $\varDelta   E // AB$. \monades{8}

\item Να υπολογίσετε:
\begin{rlist}
\item τη γωνία $\hat{x}$, \monades{8}
\item τις γωνίες $\hat{A}$ και $\hat{\varGamma }$ του τριγώνου $AB\varGamma $. \monades{9}
\end{rlist}

\begin{center}
\begin{tikzpicture}[scale=1]
\tkzDefPoint(0,0){B}
\tkzDefPoint(5,0){C}
\tkzDefTriangle[two angles = 50 and 60](B,C) \tkzGetPoint{A}

\tkzDefMidPoint(B,C) \tkzGetPoint{D}
\tkzDefMidPoint(A,C) \tkzGetPoint{E}

\draw[l3] (A)--(B)--(C)--cycle;
\draw[l3] (D)--(E);

\tkzDrawPoints(A,B,C,D,E)
\tkzLabelPoint[above](A){$A$}
\tkzLabelPoint[below](B){$B$}
\tkzLabelPoint[below](C){$\varGamma $}
\tkzLabelPoint[below](D){$\varDelta  $}
\tkzLabelPoint[right](E){$E$}

\tkzMarkAngle[size=0.6,mark=none](C,B,A)
\tkzLabelAngle[pos=0.9, xshift=-2mm](C,B,A){$50^\circ$}
\tkzMarkAngle[size=0.6,mark=none](C,D,E)
\tkzLabelAngle[pos=0.9](C,D,E){$x$}
\tkzMarkAngle[size=0.5,mark=none](C,E,D)
\tkzLabelAngle[pos=0.8](C,E,D){$70^\circ$}
\end{tikzpicture}
\end{center}
\end{alist}

\vspace{4mm}

% --- Άσκηση 36092 (36092.tex) ---
\Askhsh{36092}{2}
\begin{ylh}{1}
    \item 5.6. Εφαρμογές στα τρίγωνα
    \item 5.10. Τραπέζιο.
\end{ylh}
% ===================================================================
%  Αρχείο: 36092.tex (Fragment)
%  Αυτόματη μετατροπή μέσω doc2latex.py
% ===================================================================



Έστω τρίγωνο $AB\varGamma $ με $\varDelta$ και $E$ τα μέσα των πλευρών $AB$ και $A\varGamma $ αντίστοιχα, $A\varDelta  $=9 , $E\varGamma $=10 και $B\varGamma $=30 .

\begin{alist}
\item Να υπολογίσετε:
\begin{rlist}
\item το μήκος $x$ του τμήματος $\varDelta   E$, \monades{8}
\item την περίμετρο του τριγώνου $AB\varGamma $. \monades{9}
\end{rlist}

\item Να αποδείξετε ότι το τετράπλευρο $\varDelta   E\varGamma  B$ είναι τραπέζιο. \monades{8}

\begin{center}
\begin{tikzpicture}[scale=0.15]
\tkzDefPoint(0,0){B}
\tkzDefPoint(30,0){C}
\tkzInterCC[R](B,18)(C,20) \tkzGetPoints{A}{A_dummy}

\tkzDefMidPoint(A,B) \tkzGetPoint{D}
\tkzDefMidPoint(A,C) \tkzGetPoint{E}

\draw[l3] (A)--(B)--(C)--cycle;
\draw[l3] (D)--(E);

\tkzDrawPoints(A,B,C,D,E)
\tkzLabelPoint[above](A){$A$}
\tkzLabelPoint[below](B){$B$}
\tkzLabelPoint[below](C){$\varGamma $}
\tkzLabelPoint[left](D){$\varDelta  $}
\tkzLabelPoint[right](E){$E$}

\tkzLabelSegment[left](A,D){$9$}
\tkzLabelSegment[right](A,E){$10$}
\tkzLabelSegment[below](B,C){$30$}
\tkzLabelSegment[above](D,E){$x$}
\end{tikzpicture}
\end{center}
\end{alist}

\vspace{4mm}

% --- Άσκηση 36093 (36093.tex) ---
\Askhsh{36093}{2}
\begin{ylh}{1}
    \item 5.6. Εφαρμογές στα τρίγωνα
    \item 5.9. Μια ιδιότητα του ορθογώνιου τριγώνου.
\end{ylh}
% ===================================================================
%  Αρχείο: 36093.tex (Fragment)
%  Αυτόματη μετατροπή μέσω doc2latex.py
% ===================================================================



Δίνεται τρίγωνο $AB\varGamma $. Τα σημεία $\varDelta  $ και $E$ είναι τα μέσα των πλευρών $AB$ και $A\varGamma $ αντίστοιχα. Επιπλέον ισχύουν $A\varDelta  $=$E\varDelta  $=$\varDelta   B$ με $AE$=8 και $\varDelta   B$=10.

\begin{alist}
\item Να αποδείξετε ότι:
\begin{rlist}
\item το τρίγωνο $AEB$ είναι ορθογώνιο, \monades{8}
\item $B\varGamma $ = 20. \monades{8}
\end{rlist}

\item Να υπολογίσετε την περίμετρο του τριγώνου $AB\varGamma $. \monades{9}

\begin{center}
\begin{tikzpicture}[scale=0.2]
\tkzDefPoint(0,0){B}
\tkzDefPoint(20,0){C}
\tkzInterCC[R](B,20)(C,16) \tkzGetPoints{A}{A_dummy}

\tkzDefMidPoint(A,B) \tkzGetPoint{D}
\tkzDefMidPoint(A,C) \tkzGetPoint{E}

\draw[l3] (A)--(B)--(C)--cycle;
\draw[l3] (D)--(E);
\draw[l3, dashed] (B)--(E);

\tkzDrawPoints(A,B,C,D,E)
\tkzLabelPoint[above](A){$A$}
\tkzLabelPoint[below](B){$B$}
\tkzLabelPoint[below](C){$\varGamma $}
\tkzLabelPoint[left](D){$\varDelta  $}
\tkzLabelPoint[right](E){$E$}

\tkzMarkSegments[mark=s|](A,D D,B D,E)
\tkzMarkRightAngle(A,E,B)
\tkzLabelSegment[left](A,D){$10$}
\tkzLabelSegment[right](A,E){$8$}
\end{tikzpicture}
\end{center}
\end{alist}

\vspace{4mm}

% --- Άσκηση 36095 (36095.tex) ---
\Askhsh{36095}{2}
\begin{ylh}{1}
    \item 3.2. 1ο Κριτήριο ισότητας τριγώνων
    \item 3.15. Εφαπτόμενα τμήματα.
\end{ylh}
% ===================================================================
%  Αρχείο: 36095.tex (Fragment)
%  Αυτόματη μετατροπή μέσω doc2latex.py
% ===================================================================



Από εξωτερικό σημείο $P$ ενός κύκλου (Ο,ρ) φέρνουμε τα εφαπτόμενα τμήματα $PA$ και $PB$. Αν $M$ είναι ένα τυχαίο εσωτερικό σημείο του ευθυγράμμου τμήματος $OP$, να αποδείξετε ότι:

\begin{alist}
\item τα τρίγωνα $PAM$ και $PMB$ είναι ίσα, \monades{12}

\item οι γωνίες $Μ\hat{A}Ο$ και $Μ\hat{B}Ο$ είναι ίσες. \monades{13}


\begin{center}
\begin{tikzpicture}[scale=1]
\tkzDefPoint(0,0){O}
\tkzDefPoint(4,0){P}
\tkzDefPoint(1.2,0){M}
\tkzDefPoint(1.5,0){O_rad}
\tkzDrawCircle(O,O_rad)
\tkzDefTangent[from=P](O,O_rad) \tkzGetPoints{A}{B}

\draw[l3] (P)--(A)--(O)--(B)--cycle;
\draw[l3] (P)--(B);
\draw[l3] (P)--(O);
\draw[l3] (A)--(M)--(B);

\tkzDrawPoints(O,P,A,B,M)
\tkzLabelPoint[left](O){$O$}
\tkzLabelPoint[right](P){$P$}
\tkzLabelPoint[above](A){$A$}
\tkzLabelPoint[below](B){$B$}
\tkzLabelPoint[below left](M){$M$}

\tkzMarkRightAngle(O,A,P)
\tkzMarkRightAngle(O,B,P)
\end{tikzpicture}
\end{center}
\end{alist}

\vspace{4mm}

% --- Άσκηση 36096 (36096.tex) ---
\Askhsh{36096}{2}
\begin{ylh}{1}
    \item 3.3. 2ο Κριτήριο ισότητας τριγώνων
    \item 4.2. Τέμνουσα δύο ευθειών - Ευκλείδειο αίτημα
    \item 5.2. Παραλληλόγραμμα.
\end{ylh}
% ===================================================================
%  Αρχείο: 36096.tex (Fragment)
%  Αυτόματη μετατροπή μέσω doc2latex.py
% ===================================================================



Θεωρούμε δύο παράλληλες ευθείες ε\textsubscript{1} και ε\textsubscript{2} και τα σημεία $A$, Β στην ε\textsubscript{1} και $\varDelta$ και $\varGamma$ στην ε\textsubscript{2} ώστε τα τμήματα $A\varGamma $ και $B\varDelta  $ να τέμνονται στο μέσο $O$ του $B\varDelta  $.

Να αποδείξετε ότι:

\begin{alist}
\item τα τρίγωνα $AOB$ και $\varGamma O\varDelta  $ είναι ίσα και να αναφέρετε τα ίσα κύρια στοιχεία τους αιτιολογώντας την απάντησή σας. \monades{12}

\item το $AB\varGamma \varDelta  $ είναι παραλληλόγραμμο. \monades{13}


\begin{center}
\begin{tikzpicture}[scale=0.8]
\tkzDefPoint(0,0){D}
\tkzDefPoint(4,0){C}
\tkzDefPoint(1,2.5){A}
\tkzDefPoint(5,2.5){B}
\tkzInterLL(A,C)(B,D) \tkzGetPoint{O}

\draw[l3] (-1,2.5)--(6,2.5) node[right]{$\epsilon_1$};
\draw[l3] (-1,0)--(6,0) node[right]{$\epsilon_2$};
\draw[l3] (A)--(C);
\draw[l3] (B)--(D);
\draw[l3, dashed] (A)--(B)--(C)--(D)--cycle;

\tkzDrawPoints(A,B,C,D,O)
\tkzLabelPoint[above](A){$A$}
\tkzLabelPoint[above](B){$B$}
\tkzLabelPoint[below](C){$\varGamma $}
\tkzLabelPoint[below](D){$\varDelta  $}
\tkzLabelPoint[right](O){$O$}

\tkzMarkSegments[mark=s|](D,O O,B)
\end{tikzpicture}
\end{center}
\end{alist}

\vspace{4mm}

% --- Άσκηση 36097 (36097.tex) ---
\Askhsh{36097}{2}
\begin{ylh}{1}
    \item 3.1. Είδη και στοιχεία τριγώνων
    \item 4.2. Τέμνουσα δύο ευθειών - Ευκλείδειο αίτημα
    \item 4.6. Άθροισμα γωνιών τριγώνου
    \item 5.9. Μια ιδιότητα του ορθογώνιου τριγώνου.
\end{ylh}
% ===================================================================
%  Αρχείο: 36097.tex (Fragment)
%  Αυτόματη μετατροπή μέσω doc2latex.py
% ===================================================================



Στο παρακάτω σχήμα οι ευθείες ε\textsubscript{1}, ε\textsubscript{2} είναι παράλληλες και $AB$=6.

\begin{alist}
\item Να υπολογίσετε τις γωνίες $\hat{\varphi}$ και $\hat{\omega}$. \monades{10}

\item Να προσδιορίσετε το είδος του τριγώνου $ABK$ ως προς τις γωνίες του. \monades{7}

\item Να υπολογίσετε το μήκος της $AK$, αιτιολογώντας την απάντησή σας. \monades{8}


\begin{center}
\begin{tikzpicture}[scale=0.8]
\tkzDefPoint(0,2){k1}
\tkzDefPoint(6,2){k2}
\tkzDefPoint(0,0){l1}
\tkzDefPoint(6,0){l2}

\tkzDefPoint(1,2){A}
\tkzDefPoint(4,2){K}
\tkzDefPoint(2.5,0){B}

\draw[l3] (k1)--(k2) node[right]{$\epsilon_1$};
\draw[l3] (l1)--(l2) node[right]{$\epsilon_2$};
\draw[l3] (A)--(B)--(K)--cycle;

\tkzDrawPoints(A,B,K)
\tkzLabelPoint[above](A){$A$}
\tkzLabelPoint[above](K){$K$}
\tkzLabelPoint[below](B){$B$}

\tkzMarkAngle[size=0.6,mark=none](B,A,K)
\tkzLabelAngle[pos=0.8](B,A,K){$\varphi$}
\tkzMarkAngle[size=0.6,mark=none](K,B,A) % Assuming omega here
\tkzMarkAngle[size=0.6,mark=none](A,B,l1)
\tkzLabelAngle[pos=0.8](A,B,l1){$\omega$}

\tkzLabelSegment[above](A,B){$6$}
\end{tikzpicture}
\end{center}
\end{alist}

\vspace{4mm}

% --- Άσκηση 36098 (36098.tex) ---
\Askhsh{36098}{2}
\begin{ylh}{1}
    \item 3.1. Είδη και στοιχεία τριγώνων
    \item 3.2. 1ο Κριτήριο ισότητας τριγώνων
    \item 3.15. Εφαπτόμενα τμήματα.
\end{ylh}
% ===================================================================
%  Αρχείο: 36098.tex (Fragment)
%  Αυτόματη μετατροπή μέσω doc2latex.py
% ===================================================================



Στο παρακάτω σχήμα δίνεται κύκλος (O,R) και τα εφαπτόμενα τμήματα $MA$ και $MB$. Προεκτείνουμε την $AM$ κατά τμήμα $M\varGamma $=$MA$ και την $OM$ κατά τμήμα $M\varDelta  $=$OM$.

\begin{alist}
\item Να αποδείξετε ότι $MB$ = $M\varGamma $. \monades{10}

\item Να αποδείξετε ότι τα τρίγωνα $OMB$ και $M\varGamma \varDelta  $ είναι ίσα. \monades{15}


\begin{center}
\begin{tikzpicture}[scale=0.8]
\tkzDefPoint(0,0){O}
\tkzDefPoint(1.5,0){O_rad}
\tkzDrawCircle(O,O_rad)
\tkzDefPoint(4,1){M} % M is generic external? "τα εφαπτόμενα τμήματα $MA$ και $MB$"
\tkzDefTangent[from=M](O,O_rad) \tkzGetPoints{A}{B}

\tkzDefPointBy[homothety=center M ratio -1](A) \tkzGetPoint{C}
\tkzDefPointBy[homothety=center M ratio -1](O) \tkzGetPoint{D}

\draw[l3] (O)--(A)--(M)--(B)--cycle;
\draw[l3] (M)--(C)--(D)--cycle;
\draw[l3] (M)--(D);
\draw[l3, dashed] (O)--(M);

\tkzDrawPoints(O,M,A,B,C,D)
\tkzLabelPoint[left](O){$O$}
\tkzLabelPoint[right](M){$M$}
\tkzLabelPoint[above](A){$A$}
\tkzLabelPoint[below](B){$B$}
\tkzLabelPoint[below](C){$\varGamma $}
\tkzLabelPoint[right](D){$\varDelta  $}

\tkzMarkSegments[mark=s|](M,A M,B M,C)
\tkzMarkSegments[mark=s||](O,M M,D)
\end{tikzpicture}
\end{center}
\end{alist}

\vspace{4mm}

% --- Άσκηση 36099 (36099.tex) ---
\Askhsh{36099}{2}
\begin{ylh}{1}
    \item 3.1. Είδη και στοιχεία τριγώνων
    \item 3.2. 1ο Κριτήριο ισότητας τριγώνων.
\end{ylh}
% ===================================================================
%  Αρχείο: 36099.tex (Fragment)
%  Αυτόματη μετατροπή μέσω doc2latex.py
% ===================================================================



Δίνεται ισοσκελές τρίγωνο $AB\varGamma $ ($AB=A\varGamma $) και στις ίσες πλευρές $AB$, $A\varGamma $ παίρνουμε αντίστοιχα τμήματα $A\varDelta   = \frac{1}{3} AB$ και $AE = \frac{1}{3} A\varGamma $. Αν $M$ είναι το μέσο της $B\varGamma $, να αποδείξετε ότι:

\begin{alist}
\item τα τμήματα $B\varDelta  $ και $\varGamma  E$ είναι ίσα, \monades{5}
\item τα τρίγωνα $B\varDelta   M$ και $M E\varGamma $ είναι ίσα, \monades{10}
\item το τρίγωνο $\varDelta   EM$ είναι ισοσκελές. \monades{10}

\begin{center}
\begin{tikzpicture}[scale=1.2]
\tkzDefPoint(0,0){B}
\tkzDefPoint(4,0){C}
\tkzDefTriangle[two angles = 70 and 70](B,C) \tkzGetPoint{A}

\tkzDefPointBy[homothety=center A ratio 0.333](B) \tkzGetPoint{D}
\tkzDefPointBy[homothety=center A ratio 0.333](C) \tkzGetPoint{E}
\tkzDefMidPoint(B,C) \tkzGetPoint{M}

\draw[l3] (A)--(B)--(C)--cycle;
\draw[l3] (D)--(M)--(E);

\tkzDrawPoints(A,B,C,D,E,M)
\tkzLabelPoint[above](A){$A$}
\tkzLabelPoint[below](B){$B$}
\tkzLabelPoint[below](C){$\varGamma $}
\tkzLabelPoint[left](D){$\varDelta  $}
\tkzLabelPoint[right](E){$E$}
\tkzLabelPoint[below](M){$M$}

\tkzMarkSegments[mark=s|](A,B A,C)
\tkzMarkSegments[mark=s||](B,M M,C)
\end{tikzpicture}
\end{center}
\end{alist}

\vspace{4mm}

% --- Άσκηση 36100 (36100.tex) ---
\Askhsh{36100}{2}
\begin{ylh}{1}
    \item 3.1. Είδη και στοιχεία τριγώνων
    \item 3.2. 1ο Κριτήριο ισότητας τριγώνων.
\end{ylh}
% ===================================================================
%  Αρχείο: 36100.tex (Fragment)
%  Αυτόματη μετατροπή μέσω doc2latex.py
% ===================================================================



Δίνεται ισοσκελές τρίγωνο $KAB$ ($KA$=$KB$) και $K\varGamma $ διχοτόμος της γωνίας $\hat{K}$. Στην προέκταση της $BA$ (προς το Α) παίρνουμε σημείο $\varLambda $ και στην προέκταση της $AB$ (προς το Β) παίρνουμε σημείο $M$, έτσι ώστε $A\varLambda $=$BM$. Να αποδείξετε ότι:

\begin{alist}
\item το τρίγωνο $K\varLambda M$ είναι ισοσκελές, \monades{12}

\item η $K\varGamma $ είναι διάμεσος του τριγώνου $K\varLambda M$. \monades{13}
\end{alist}

\vspace{4mm}

% --- Άσκηση 36101 (36101.tex) ---
\Askhsh{36101}{2}
\begin{ylh}{1}
    \item 4.4. Γωνίες με πλευρές παράλληλες
    \item 4.6. Άθροισμα γωνιών τριγώνου.
\end{ylh}
% ===================================================================
%  Αρχείο: 36101.tex (Fragment)
%  Αυτόματη μετατροπή μέσω doc2latex.py
% ===================================================================



Δίνεται τρίγωνο $AB\varGamma $ με $\hat{A} = 80\degree$ και $\hat{B} = 20\degree + \hat{\varGamma }$, και $A\varDelta  $ η διχοτόμος της γωνίας $\hat{A}$.

\begin{alist}
\item Να υπολογίσετε τις γωνίες $\hat{B}$ και $\hat{\varGamma }$. \monades{12}

\item Φέρουμε από το $\varDelta  $ ευθεία παράλληλη στην $AB$, που τέμνει την $A\varGamma $ στο $E$. Να υπολογίσετε τις γωνίες $A\hat{\varDelta  }E$ και $E\hat{\varDelta  }\varGamma $. \monades{13}

\begin{center}
\begin{tikzpicture}[scale=1]
\tkzDefPoint(0,0){B}
\tkzDefPoint(5,0){C}
\tkzDefTriangle[two angles = 60 and 40](B,C) \tkzGetPoint{A}

\tkzDefLine[bisector](B,A,C) \tkzGetPoint{AD_dir}
\tkzInterLL(A,AD_dir)(B,C) \tkzGetPoint{D}

\tkzDefLine[parallel=through D](A,B) \tkzGetPoint{DE_dir}
\tkzInterLL(D,DE_dir)(A,C) \tkzGetPoint{E}

\draw[l3] (A)--(B)--(C)--cycle;
\draw[l3] (A)--(D);
\draw[l3] (D)--(E);

\tkzDrawPoints(A,B,C,D,E)
\tkzLabelPoint[above](A){$A$}
\tkzLabelPoint[below](B){$B$}
\tkzLabelPoint[below](C){$\varGamma $}
\tkzLabelPoint[below](D){$\varDelta  $}
\tkzLabelPoint[right](E){$E$}

\tkzMarkAngle[size=0.6,mark=none](B,A,D)
\tkzMarkAngle[size=0.6,mark=none](D,A,C)
\tkzMarkAngle[size=0.5,mark=none](C,D,E)
\end{tikzpicture}
\end{center}
\end{alist}

\vspace{4mm}

% --- Άσκηση 36102 (36102.tex) ---
\Askhsh{36102}{2}
\begin{ylh}{1}
    \item 3.4. 3ο Κριτήριο ισότητας τριγώνων.
\end{ylh}
% ===================================================================
%  Αρχείο: 36102.tex (Fragment)
%  Αυτόματη μετατροπή μέσω doc2latex.py
% ===================================================================



Δίνεται τετράπλευρο $AB\varGamma \varDelta  $ με $BA$=$B\varGamma $ και $\varDelta  A$=$\varDelta  \varGamma $. Οι διαγώνιοι $A\varGamma $, $B\varDelta  $ του τετραπλεύρου είναι ίσες και τέμνονται κάθετα.

Να αποδείξετε ότι:

\begin{alist}
\item η $B\varDelta  $ είναι διχοτόμος των γωνιών Β και $\varDelta$ του τετραπλεύρου $AB\varGamma \varDelta  $, \monades{12}

\item η $B\varDelta  $ είναι μεσοκάθετος του τμήματος $A\varGamma $. \monades{13}


\begin{center}
\begin{tikzpicture}[scale=1]
\tkzDefPoint(0,0){A}
\tkzDefPoint(4,0){C}
\tkzDefMidPoint(A,C) \tkzGetPoint{O}
% $BD$ = $AC$ = 4. 
\tkzDefPoint(2,1){B}
\tkzDefPoint(2,-3){D}

\draw[l3] (A)--(B)--(C)--(D)--cycle;
\draw[l3] (A)--(C);
\draw[l3] (B)--(D);

\tkzDrawPoints(A,B,C,D,O)
\tkzLabelPoint[left](A){$A$}
\tkzLabelPoint[above](B){$B$}
\tkzLabelPoint[right](C){$\varGamma $}
\tkzLabelPoint[below](D){$\varDelta  $}

\tkzMarkRightAngle(B,O,C)
\tkzMarkSegments[mark=s|](A,B B,C)
\tkzMarkSegments[mark=s||](A,D D,C)
\end{tikzpicture}
\end{center}
\end{alist}

\vspace{4mm}

% --- Άσκηση 36103 (36103.tex) ---
\Askhsh{36103}{2}
\begin{ylh}{1}
    \item 4.6. Άθροισμα γωνιών τριγώνου
    \item 5.4. Ρόμβος
    \item 5.6. Εφαρμογές στα τρίγωνα
    \item 5.9. Μια ιδιότητα του ορθογώνιου τριγώνου.
\end{ylh}
% ===================================================================
%  Αρχείο: 36103.tex (Fragment)
%  Αυτόματη μετατροπή μέσω doc2latex.py
% ===================================================================



Δίνεται ισόπλευρο τρίγωνο $AB\varGamma $. Φέρουμε την εξωτερική διχοτόμο $Ax$ της γωνίας $\hat{A}$ και από το σημείο $\varGamma $ την κάθετο $\varGamma \varDelta  $ στην $Ax$. Τα σημεία $E$ και $Z$ είναι τα μέσα των πλευρών $AB$ και $A\varGamma $ αντίστοιχα.

\begin{alist}
\item $EZ = $AE$ = AZ$, \monades{9}
\item η γωνία $A\hat{\varGamma }\varDelta  $ είναι ίση με $30\degree$, \monades{10}
\item το τετράπλευρο $A\varDelta   ZE$ είναι ρόμβος. \monades{6}

\begin{center}
\begin{tikzpicture}[scale=1]
\tkzDefPoint(0,0){B}
\tkzDefPoint(4,0){C}
\tkzDefTriangle[equilateral](B,C) \tkzGetPoint{A}

\tkzDefPoint(-2,3.46){x_pt} % Ax parallel to $BC$ since external bisector of 60 deg is 60 deg from sides.
% External is 120. Half is 60. So Ax makes 60 with $AB$. 
% Internal A is 60. So Ax is parallel to $BC$.

\tkzDefPointBy[projection=onto A--x_pt](C) \tkzGetPoint{D}
\tkzDefMidPoint(A,B) \tkzGetPoint{E}
\tkzDefMidPoint(A,C) \tkzGetPoint{Z}

\draw[l3] (A)--(B)--(C)--cycle;
\draw[l3, dashed] (A)--(x_pt) node[left]{$x$};
\draw[l3] (C)--(D);
\draw[l3] (D)--(Z);
\draw[l3] (E)--(Z);

\tkzDrawPoints(A,B,C,D,E,Z)
\tkzLabelPoint[above](A){$A$}
\tkzLabelPoint[below](B){$B$}
\tkzLabelPoint[below](C){$\varGamma $}
\tkzLabelPoint[above](D){$\varDelta  $}
\tkzLabelPoint[left](E){$E$}
\tkzLabelPoint[right](Z){$Z$}

\tkzMarkRightAngle(C,D,A)
\end{tikzpicture}
\end{center}
\end{alist}

\vspace{4mm}

% --- Άσκηση 36104 (36104.tex) ---
\Askhsh{36104}{2}
\begin{ylh}{1}
    \item 3.1. Είδη και στοιχεία τριγώνων
    \item 3.2. 1ο Κριτήριο ισότητας τριγώνων.
\end{ylh}
% ===================================================================
%  Αρχείο: 36104.tex (Fragment)
%  Αυτόματη μετατροπή μέσω doc2latex.py
% ===================================================================



Δίνεται γωνία x$\hat{O}$y και η διχοτόμος της Οδ. Θεωρούμε σημείο $M$ της Οδ και σημεία $A$ και Β στις ημιευθείες Οx και Oy αντίστοιχα, τέτοια ώστε $OA$=$OB$.

Να αποδείξετε ότι:

\begin{alist}
\item $MA$=$MB$, \monades{15}

\item Η Οδ είναι διχοτόμος της γωνίας Α$\hat{Μ}$Β. \monades{10}
\end{alist}

\vspace{4mm}

% --- Άσκηση 36105 (36105.tex) ---
\Askhsh{36105}{2}
\begin{ylh}{1}
    \item 3.6. Κριτήρια ισότητας ορθογώνιων τριγώνων
    \item 4.2. Τέμνουσα δύο ευθειών - Ευκλείδειο αίτημα
    \item 5.2. Παραλληλόγραμμα.
\end{ylh}
% ===================================================================
%  Αρχείο: 36105.tex (Fragment)
%  Αυτόματη μετατροπή μέσω doc2latex.py
% ===================================================================



Σε παραλληλόγραμμο $AB\varGamma \varDelta  $ ($AB$//$\varGamma \varDelta  $) με $AB$ > $B\varGamma $ φέρουμε από τις κορυφές Α και $\varGamma$ καθέτους στη διαγώνιο $B\varDelta  $, οι οποίες την τέμνουν σε διαφορετικά σημεία $E$ και $Z$ αντίστοιχα.

Να αποδείξετε ότι:

\begin{alist}
\item $AE$=$\varGamma Z$, \monades{15}

\item το τετράπλευρο $AE\varGamma Z$ είναι παραλληλόγραμμο. \monades{10}


\begin{center}
\begin{tikzpicture}[scale=0.8]
\tkzDefPoint(0,0){A}
\tkzDefPoint(5,0){B}
\tkzDefPoint(6.5,3){C}
\tkzDefPoint(1.5,3){D}

\draw[l3] (A)--(B)--(C)--(D)--cycle;
\draw[l3] (B)--(D);

\tkzDefPointBy[projection=onto B--D](A) \tkzGetPoint{E}
\tkzDefPointBy[projection=onto B--D](C) \tkzGetPoint{Z}

\draw[l3] (A)--(E);
\draw[l3] (C)--(Z);
\draw[l3, dashed] (A)--(Z);
\draw[l3, dashed] (C)--(E);

\tkzDrawPoints(A,B,C,D,E,Z)
\tkzLabelPoint[left](A){$A$}
\tkzLabelPoint[right](B){$B$}
\tkzLabelPoint[right](C){$\varGamma $}
\tkzLabelPoint[left](D){$\varDelta  $}
\tkzLabelPoint[above left](E){$E$}
\tkzLabelPoint[below right](Z){$Z$}

\tkzMarkRightAngle(A,E,B)
\tkzMarkRightAngle(C,Z,D)
\end{tikzpicture}
\end{center}
\end{alist}

\vspace{4mm}

% --- Άσκηση 36106 (36106.tex) ---
\Askhsh{36106}{2}
\begin{ylh}{1}
    \item 3.11. Ανισοτικές σχέσεις πλευρών και γωνιών
    \item 4.2. Τέμνουσα δύο ευθειών - Ευκλείδειο αίτημα
    \item 4.6. Άθροισμα γωνιών τριγώνου
    \item 5.1 Εισαγωγή
    \item 5.11. Ισοσκελές τραπέζιο.
\end{ylh}
% ===================================================================
%  Αρχείο: 36106.tex (Fragment)
%  Αυτόματη μετατροπή μέσω doc2latex.py
% ===================================================================



Θεωρούμε ισοσκελές τραπέζιο $AB\varGamma \varDelta  $ ($AB$//$\varGamma \varDelta  $) με $\varGamma \varDelta  $\textgreater $AB$ και $\hat{B}$ = 135\textsuperscript{ο}. Από τις κορυφές Α και Β φέρουμε τα ύψη του $AE$ και $BZ$.

\begin{alist}
\item Να υπολογίσετε τις γωνίες του τραπεζίου. \monades{10}

\item Να αποδείξετε ότι $AE$=$E\varDelta  $=$BZ$=$\varGamma Z$. \monades{15}


\begin{center}
\begin{tikzpicture}[scale=1]
\tkzDefPoint(0,0){E}
\tkzDefPoint(0,2){A}
\tkzDefPoint(-2,0){D}
\tkzDefPoint(4,0){Z}
\tkzDefPoint(4,2){B}
\tkzDefPoint(6,0){C}

\draw[l3] (A)--(B)--(C)--(D)--cycle;
\draw[l3] (A)--(E);
\draw[l3] (B)--(Z);

\tkzDrawPoints(A,B,C,D,E,Z)
\tkzLabelPoint[above](A){$A$}
\tkzLabelPoint[above](B){$B$}
\tkzLabelPoint[below](C){$\varGamma $}
\tkzLabelPoint[below](D){$\varDelta  $}
\tkzLabelPoint[below](E){$E$}
\tkzLabelPoint[below](Z){$Z$}

\tkzMarkRightAngle(A,E,D)
\tkzMarkRightAngle(B,Z,C)
\tkzMarkAngle[size=0.6,mark=none](B,C,Z)
\tkzLabelAngle[pos=0.9](B,C,Z){$45^\circ$}
\tkzMarkAngle[size=0.6,mark=none](E,D,A)
\tkzLabelAngle[pos=0.9](E,D,A){$45^\circ$}
\end{tikzpicture}
\end{center}
\end{alist}

\vspace{4mm}

% --- Άσκηση 36107 (36107.tex) ---
\Askhsh{36107}{2}
\begin{ylh}{1}
    \item 3.4. 3ο Κριτήριο ισότητας τριγώνων
    \item 3.6. Κριτήρια ισότητας ορθογώνιων τριγώνων
    \item 5.4. Ρόμβος.
\end{ylh}
% ===================================================================
%  Αρχείο: 36107.tex (Fragment)
%  Αυτόματη μετατροπή μέσω doc2latex.py
% ===================================================================



Δίνεται οξυγώνιο τρίγωνο $A\varGamma B$. Φέρουμε από τη κορυφή $A$ ευθεία (ε) παράλληλη στη $B\varGamma $. Η κάθετη στο μέσο $M$ της πλευράς $AB$ τέμνει την (ε) στο $\varDelta$ και την $B\varGamma $ στο Ε.

\begin{alist}
\item Να αποδείξετε ότι $\varDelta  A$=$\varDelta   B$ και $EA$=$EB$. \monades{6}

\item Να συγκρίνετε τα τρίγωνα $AM\varDelta  $ και $EMB$. \monades{10}

\item Να αποδείξετε ότι το τετράπλευρο ΑΔ$BE$ είναι ρόμβος. \monades{9}


\begin{center}
\begin{tikzpicture}[scale=0.8]
\tkzDefPoint(1,4){A}
\tkzDefPoint(0,0){B}
\tkzDefPoint(5,0){C}
\tkzDefMidPoint(A,B) \tkzGetPoint{M}

\tkzDefLine[parallel=through A](B,C) \tkzGetPoint{e_pt}
\tkzDefLine[orthogonal=through M](A,B) \tkzGetPoint{perp_pt}
\tkzInterLL(M,perp_pt)(A,e_pt) \tkzGetPoint{D}
\tkzInterLL(M,perp_pt)(B,C) \tkzGetPoint{E}

\draw[l3] (A)--(B)--(C)--cycle;
\draw[l3] (A)--(e_pt) node[right]{$\epsilon$};
\draw[l3] (D)--(E);
\draw[l3] (A)--(D);
\draw[l3] (A)--(E);
\draw[l3] (B)--(D);

\tkzDrawPoints(A,B,C,M,D,E)
\tkzLabelPoint[above](A){$A$}
\tkzLabelPoint[left](B){$B$}
\tkzLabelPoint[right](C){$\varGamma $}
\tkzLabelPoint[above left](M){$M$}
\tkzLabelPoint[above right](D){$\varDelta  $}
\tkzLabelPoint[below](E){$E$}

\tkzMarkSegments[mark=s|](A,M M,B)
\tkzMarkRightAngle(D,M,A)
\end{tikzpicture}
\end{center}
\end{alist}

\vspace{4mm}

% --- Άσκηση 36109 (36109.tex) ---
\Askhsh{36109}{2}
\begin{ylh}{1}
    \item 4.6. Άθροισμα γωνιών τριγώνου
    \item 5.9. Μια ιδιότητα του ορθογώνιου τριγώνου.
\end{ylh}
% ===================================================================
%  Αρχείο: 36109.tex (Fragment)
%  Αυτόματη μετατροπή μέσω doc2latex.py
% ===================================================================



Σε τρίγωνο $AB\varGamma $ ισχύει $\hat{A}$ + $\hat{\varGamma }$ = 120\textsuperscript{ο} και $\hat{A} = 3\hat{\varGamma }$.

\begin{alist}
\item Να αποδείξετε ότι το τρίγωνο $AB\varGamma $ είναι ορθογώνιο και να υπολογίσετε τις γωνίες του. \monades{15}

\item Αν η πλευρά $B\varGamma $=2cm να βρείτε το μήκος της $AB$. \monades{10}
\end{alist}

\vspace{4mm}

% --- Άσκηση 36110 (36110.tex) ---
\Askhsh{36110}{2}
\begin{ylh}{1}
    \item 3.2. 1ο Κριτήριο ισότητας τριγώνων.
\end{ylh}
% ===================================================================
%  Αρχείο: 36110.tex (Fragment)
%  Αυτόματη μετατροπή μέσω doc2latex.py
% ===================================================================



Αν στο σχήμα που ακολουθεί είναι $A\hat{O}B = B\hat{O}\varGamma  = \varGamma \hat{O}\varDelta  $ και $OA=$OB$=O\varGamma =O\varDelta  $, τότε να αποδείξετε ότι:

\begin{alist}
\item $A\varGamma =B\varDelta  $, \monades{10}
\item το $M$ είναι μέσον της $B\varDelta  $, όπου $M$ το σημείο τομής των τμημάτων $O\varGamma $ και $B\varDelta  $. \monades{15}

\begin{center}
\begin{tikzpicture}[scale=1]
\tkzDefPoint(0,0){O}
\tkzDefPoint(3,0){X}
\tkzDefPointBy[rotation=center O angle 30](X) \tkzGetPoint{Delta}
\tkzDefPointBy[rotation=center O angle 60](X) \tkzGetPoint{Gamma}
\tkzDefPointBy[rotation=center O angle 90](X) \tkzGetPoint{B}
\tkzDefPointBy[rotation=center O angle 120](X) \tkzGetPoint{A}

\tkzInterLL(O,Gamma)(Delta,B) \tkzGetPoint{M}

\draw[l3] (O)--(A);
\draw[l3] (O)--(B);
\draw[l3] (O)--(Gamma);
\draw[l3] (O)--(Delta);
\draw[l3] (A)--(Gamma);
\draw[l3] (B)--(Delta);

\tkzDrawPoints(O,A,B,Gamma,Delta,M)
\tkzLabelPoint[below](O){$O$}
\tkzLabelPoint[above left](A){$A$}
\tkzLabelPoint[above](B){$B$}
\tkzLabelPoint[above right](Gamma){$\varGamma $}
\tkzLabelPoint[right](Delta){$\varDelta  $}
\tkzLabelPoint[above right](M){$M$}

\tkzMarkAngles[size=0.7,mark=none](B,O,A Gamma,O,B Delta,O,Gamma)
\tkzMarkSegments[mark=s|](O,A O,B O,Gamma O,Delta)
\end{tikzpicture}
\end{center}
\end{alist}

\vspace{4mm}

% --- Άσκηση 36111 (36111.tex) ---
\Askhsh{36111}{2}
\begin{ylh}{1}
    \item 3.2. 1ο Κριτήριο ισότητας τριγώνων
    \item 4.6. Άθροισμα γωνιών τριγώνου
    \item 5.9. Μια ιδιότητα του ορθογώνιου τριγώνου.
\end{ylh}
% ===================================================================
%  Αρχείο: 36111.tex (Fragment)
%  Αυτόματη μετατροπή μέσω doc2latex.py
% ===================================================================



Δίνεται ορθογώνιο τρίγωνο $AB\varGamma $ με $\hat{A} = 90^{\circ}$ και $\hat{\varGamma } = 25^{\circ}$. Δίνονται επίσης η διάμεσος $AM$, το ύψος $AH$ από την κορυφή $A$ και η διχοτόμος $A\varDelta  $ της γωνίας $\hat{A}$.

\begin{alist}
\item Να υπολογίσετε τις γωνίες $A\hat{M}B$, $H\hat{A}B$ και $A\hat{\varDelta  }B$. \monades{15}
\item Να αποδείξετε ότι $M\hat{A}\varDelta   = \varDelta  \hat{A}H = 20^{\circ}$. \monades{10}

\begin{center}
\begin{tikzpicture}[scale=1.2]
\tkzDefPoint(0,0){A}
\tkzDefPoint(0,3){B}
\tkzDefPoint(6.43,0){C} % tan(25) = 3/$AC$ -> $AC$ = 3 / 0.466 = 6.43
\tkzDefMidPoint(B,C) \tkzGetPoint{M}
\tkzDefPointBy[projection=onto B--C](A) \tkzGetPoint{H}
\tkzDefLine[bisector](B,A,C) \tkzGetPoint{AD_dir}
\tkzInterLL(A,AD_dir)(B,C) \tkzGetPoint{D}

\draw[l3] (A)--(B)--(C)--cycle;
\draw[l3] (A)--(M);
\draw[l3] (A)--(H);
\draw[l3] (A)--(D);

\tkzDrawPoints(A,B,C,M,H,D)
\tkzLabelPoint[below left](A){$A$}
\tkzLabelPoint[above](B){$B$}
\tkzLabelPoint[right](C){$\varGamma $}
\tkzLabelPoint[above right](M){$M$}
\tkzLabelPoint[above](H){$H$}
\tkzLabelPoint[below](D){$\varDelta  $}

\tkzMarkRightAngle(B,A,C)
\tkzMarkRightAngle(A,H,B)
\tkzMarkAngle[size=0.6,mark=none](A,B,H)
\tkzLabelAngle[pos=0.9](A,C,B){$25^\circ$}
\end{tikzpicture}
\end{center}
\end{alist}

\vspace{4mm}

% --- Άσκηση 36112 (36112.tex) ---
\Askhsh{36112}{2}
\begin{ylh}{1}
    \item 3.6. Κριτήρια ισότητας ορθογώνιων τριγώνων
    \item 4.2. Τέμνουσα δύο ευθειών - Ευκλείδειο αίτημα
    \item 5.3. Ορθογώνιο
    \item 5.9. Μια ιδιότητα του ορθογώνιου τριγώνου
    \item 5.10. Τραπέζιο
    \item 5.11. Ισοσκελές τραπέζιο.
\end{ylh}
% ===================================================================
%  Αρχείο: 36112.tex (Fragment)
%  Αυτόματη μετατροπή μέσω doc2latex.py
% ===================================================================



Δίνεται ισοσκελές τραπέζιο $AB\varGamma \varDelta  $ ($AB//\varGamma \varDelta  $), με $AB=6$, $B\varGamma =4$ και $\hat{\varGamma } = 60^{\circ}$. Δίνονται επίσης τα ύψη $AE$ και $BZ$ από τις κορυφές $A$ και $B$ αντίστοιχα.

\begin{alist}
\item Να υπολογίσετε τις υπόλοιπες γωνίες του τραπεζίου $AB\varGamma \varDelta  $. \monades{6}
\item Να αποδείξετε τα τρίγωνα $AE\varDelta  $, $BZ\varGamma $ είναι ίσα. \monades{10}
\item Να υπολογίσετε την περίμετρο του $AB\varGamma \varDelta  $. \monades{9}

\begin{center}
\begin{tikzpicture}[scale=0.6]
\tkzDefPoint(0,0){Z}
\tkzDefPoint(0,3.46){B}
\tkzDefPoint(2,0){C}
\tkzDefPoint(-6,3.46){A}
\tkzDefPoint(-6,0){E}
\tkzDefPoint(-8,0){Delta}

\draw[l3] (A)--(B)--(C)--(Delta)--cycle;
\draw[l3] (A)--(E);
\draw[l3] (B)--(Z);

\tkzDrawPoints(A,B,C,Delta,E,Z)
\tkzLabelPoint[above](A){$A$}
\tkzLabelPoint[above](B){$B$}
\tkzLabelPoint[below](C){$\varGamma $}
\tkzLabelPoint[below](Delta){$\varDelta  $}
\tkzLabelPoint[below](E){$E$}
\tkzLabelPoint[below](Z){$Z$}

\tkzMarkRightAngle(A,E,Delta)
\tkzMarkRightAngle(B,Z,C)
\tkzMarkAngle[size=0.8,mark=none](B,C,Delta)
\tkzLabelAngle[pos=1.2](B,C,Delta){$60^\circ$}
\end{tikzpicture}
\end{center}
\end{alist}

\vspace{4mm}

% --- Άσκηση 36113 (36113.tex) ---
\Askhsh{36113}{2}
\begin{ylh}{1}
    \item 4.2. Τέμνουσα δύο ευθειών - Ευκλείδειο αίτημα
    \item 4.6. Άθροισμα γωνιών τριγώνου
    \item 5.3. Ορθογώνιο
    \item 5.9. Μια ιδιότητα του ορθογώνιου τριγώνου
    \item 5.10. Τραπέζιο.
\end{ylh}
% ===================================================================
%  Αρχείο: 36113.tex (Fragment)
%  Αυτόματη μετατροπή μέσω doc2latex.py
% ===================================================================



Δίνεται τραπέζιο $AB\varGamma \varDelta  $ ($AB//\varGamma \varDelta  $), με $AB=B\varGamma =4$, $\hat{A} = 90^{\circ}$ και $\hat{\varGamma } = 60^{\circ}$. Φέρουμε το ύψος $BE$ από τη κορυφή $B$.

\begin{alist}
\item Να υπολογίσετε τις άλλες δυο γωνίες του τραπεζίου $AB\varGamma \varDelta  $. \monades{8}
\item Να αποδείξετε ότι $2E\varGamma =B\varGamma $. \monades{9}
\item Αν $M,N$ τα μέσα των πλευρών $A\varDelta  $, $B\varGamma $ αντίστοιχα να βρείτε το μήκος του ευθυγράμμου τμήματος $MN$. \monades{8}

\begin{center}
\begin{tikzpicture}[scale=0.8]
\tkzDefPoint(0,3.46){A}
\tkzDefPoint(4,3.46){B}
\tkzDefPoint(6,0){C}
\tkzDefPoint(0,0){Delta}
\tkzDefPoint(4,0){E}

\draw[l3] (A)--(B)--(C)--(Delta)--cycle;
\draw[l3] (B)--(E);

\tkzDefMidPoint(A,Delta) \tkzGetPoint{M}
\tkzDefMidPoint(B,C) \tkzGetPoint{N}
\draw[l3, dashed] (M)--(N);

\tkzDrawPoints(A,B,C,Delta,E,M,N)
\tkzLabelPoint[above](A){$A$}
\tkzLabelPoint[above](B){$B$}
\tkzLabelPoint[below](C){$\varGamma $}
\tkzLabelPoint[below](Delta){$\varDelta  $}
\tkzLabelPoint[below](E){$E$}
\tkzLabelPoint[left](M){$M$}
\tkzLabelPoint[right](N){$N$}

\tkzMarkRightAngle(B,E,C)
\tkzMarkAngle[size=0.6,mark=none](B,C,Delta)
\tkzLabelAngle[pos=0.9](B,C,Delta){$60^\circ$}
\tkzMarkSegments[mark=s|](A,B B,C)
\end{tikzpicture}
\end{center}
\end{alist}

\vspace{4mm}

% --- Άσκηση 36114 (36114.tex) ---
\Askhsh{36114}{2}
\begin{ylh}{1}
    \item 3.6. Κριτήρια ισότητας ορθογώνιων τριγώνων
    \item 3.15. Εφαπτόμενα τμήματα
    \item 4.8. Άθροισμα γωνιών κυρτού ν-γώνου.
\end{ylh}
% ===================================================================
%  Αρχείο: 36114.tex (Fragment)
%  Αυτόματη μετατροπή μέσω doc2latex.py
% ===================================================================



Δίνεται κύκλος κέντρου Ο και ένα εξωτερικό του σημείο $P$, από το οποίο φέρουμε τα εφαπτόμενα τμήματα του κύκλου $PA$ και $PB$. Έστω ότι το τμήμα $PO$ τέμνει τον κύκλο στο σημείο $M$ και η εφαπτομένη του κύκλου στο $M$ τέμνει τα $PA$ και $PB$ στα σημεία $\varDelta  $ και $\varGamma$ αντίστοιχα.

\begin{alist}
\item Να αποδείξετε ότι το τρίγωνο $P\varDelta  \varGamma $ είναι ισοσκελές. \monades{13}

\item Αν η γωνία Α$\hat{Ρ}$Β είναι 40\degree να υπολογίσετε τη γωνία Α$\hat{O}$Β. \monades{12}


\begin{center}
\begin{tikzpicture}[scale=1]
\tkzDefPoint(0,0){O}
\tkzDefPoint(1.5,0){O_rad}
\tkzDrawCircle(O,O_rad)
\tkzDefPoint(4,0){P}
\tkzDefTangent[from=P](O,O_rad) \tkzGetPoints{A}{B}
\tkzInterLC(P,O)(O,O_rad) \tkzGetPoints{M_dummy}{M}

\tkzDefLine[orthogonal=through M](O,M) \tkzGetPoint{tangent_M}
\tkzInterLL(M,tangent_M)(P,A) \tkzGetPoint{Delta}
\tkzInterLL(M,tangent_M)(P,B) \tkzGetPoint{Gamma}

\draw[l3] (P)--(A);
\draw[l3] (P)--(B);
\draw[l3] (O)--(A);
\draw[l3] (O)--(B);
\draw[l3] (P)--(O);
\draw[l3] (Delta)--(Gamma);

\tkzDrawPoints(O,P,A,B,M,Delta,Gamma)
\tkzLabelPoint[left](O){$O$}
\tkzLabelPoint[right](P){$P$}
\tkzLabelPoint[above](A){$A$}
\tkzLabelPoint[below](B){$B$}
\tkzLabelPoint[below left](M){$M$}
\tkzLabelPoint[above right](Delta){$\varDelta  $}
\tkzLabelPoint[below right](Gamma){$\varGamma $}

\tkzMarkRightAngle(O,A,P)
\tkzMarkRightAngle(O,B,P)
\tkzMarkRightAngle(O,M,Delta)
\end{tikzpicture}
\end{center}
\end{alist}

\vspace{4mm}

% --- Άσκηση 36115 (36115.tex) ---
\Askhsh{36115}{2}
\begin{ylh}{1}
    \item 3.2. 1ο Κριτήριο ισότητας τριγώνων
    \item 4.2. Τέμνουσα δύο ευθειών - Ευκλείδειο αίτημα
    \item 4.6. Άθροισμα γωνιών τριγώνου
    \item 5.2. Παραλληλόγραμμα.
\end{ylh}
% ===================================================================
%  Αρχείο: 36115.tex (Fragment)
%  Αυτόματη μετατροπή μέσω doc2latex.py
% ===================================================================



Δίνεται ισόπλευρο τρίγωνο $AB\varGamma $ . Στην προέκταση της $B\varGamma $ (προς το μέρος του Γ) θεωρούμε τμήμα $\varGamma \varDelta  $=$B\varGamma $. Φέρουμε τμήμα $\varDelta  E$ κάθετο στην $A\varDelta  $ στο σημείο της Δ , τέτοιο ώστε $\varDelta  E$=$B\varGamma $ (τα Α και $E$ είναι στο ίδιο ημιεπίπεδο ως προς την $B\varDelta  $).

\begin{alist}
\item Να βρείτε τις γωνίες του τριγώνου Α$B\varDelta  $. \monades{12}

\item Να αποδείξετε ότι το Α$B\varDelta  $Ε είναι παραλληλόγραμμο. \monades{13}


\begin{center}
\begin{tikzpicture}[scale=0.8]
\tkzDefPoint(0,0){B}
\tkzDefPoint(3,0){C}
\tkzDefTriangle[equilateral](B,C) \tkzGetPoint{A}
\tkzDefPoint(6,0){Delta}

\tkzDefLine[orthogonal=through Delta](A,Delta) \tkzGetPoint{perp_dir}
\tkzDefPointBy[homothety=center Delta ratio 1](perp_dir) \tkzGetPoint{E_trial}
\tkzDefPointBy[homothety=center Delta ratio 0.577](perp_dir) % Wait, $DE$=$BC$=3.
% perp_dir distance?
\tkzDefCircle[R](Delta,3) \tkzGetPoint{dummy}
\tkzInterLC(Delta,perp_dir)(Delta,dummy) \tkzGetPoints{E}{E_dummy}

\draw[l3] (A)--(B)--(Delta)--cycle;
\draw[l3] (A)--(C);
\draw[l3] (Delta)--(E);
\draw[l3] (E)--(A); % Prove it's a parallelogram $ABD$...E?
\draw[l3, dashed] (A)--(Delta);

\tkzDrawPoints(A,B,C,Delta,E)
\tkzLabelPoint[above](A){$A$}
\tkzLabelPoint[below](B){$B$}
\tkzLabelPoint[below](C){$\varGamma $}
\tkzLabelPoint[below](Delta){$\varDelta  $}
\tkzLabelPoint[above right](E){$E$}

\tkzMarkRightAngle(A,Delta,E)
\tkzMarkSegments[mark=s|](B,C C,Gamma A,B A,C Gamma,Delta Delta,E)
\end{tikzpicture}
\end{center}
\end{alist}

\vspace{4mm}

% --- Άσκηση 36116 (36116.tex) ---
\Askhsh{36116}{2}
\begin{ylh}{1}
    \item 3.2. 1ο Κριτήριο ισότητας τριγώνων
    \item 4.6. Άθροισμα γωνιών τριγώνου
    \item 5.9. Μια ιδιότητα του ορθογώνιου τριγώνου.
\end{ylh}
% ===================================================================
%  Αρχείο: 36116.tex (Fragment)
%  Αυτόματη μετατροπή μέσω doc2latex.py
% ===================================================================



Δίνεται ορθογώνιο τρίγωνο $AB\varGamma $ (με $\hat{A} = 90^{\circ}$) και η διχοτόμος της γωνίας $\hat{\varGamma }$ τέμνει την πλευρά $AB$ στο σημείο $\varDelta  $, τέτοιο ώστε $\varGamma \varDelta  =\varDelta   B=2\text{cm}$.

Να αποδείξετε ότι:

\begin{alist}
\item $\hat{B} = 30^{\circ}$, \monades{12}
\item $AB=3\text{cm}$. \monades{13}

\begin{center}
\begin{tikzpicture}[scale=1.5]
\tkzDefPoint(0,0){A}
\tkzDefPoint(1.732,0){B} % sqrt(3)
\tkzDefPoint(0,1){C}
\tkzDefLine[bisector](B,C,A) \tkzGetPoint{CD_dir}
\tkzInterLL(C,CD_dir)(A,B) \tkzGetPoint{Delta}

\draw[l3] (A)--(B)--(C)--cycle;
\draw[l3] (C)--(Delta);

\tkzDrawPoints(A,B,C,Delta)
\tkzLabelPoint[below left](A){$A$}
\tkzLabelPoint[below right](B){$B$}
\tkzLabelPoint[above left](C){$\varGamma $}
\tkzLabelPoint[below](Delta){$\varDelta  $}

\tkzMarkRightAngle(C,A,B)
\tkzMarkAngle[size=0.4,mark=none](B,C,Delta)
\tkzMarkAngle[size=0.4,mark=none](Delta,C,A)
\tkzMarkSegments[mark=s|](Delta,B Delta,C)
\end{tikzpicture}
\end{center}
\end{alist}

\vspace{4mm}

% --- Άσκηση 36117 (36117.tex) ---
\Askhsh{36117}{2}
\begin{ylh}{1}
    \item 3.11. Ανισοτικές σχέσεις πλευρών και γωνιών
    \item 4.6. Άθροισμα γωνιών τριγώνου
    \item 4.8. Άθροισμα γωνιών κυρτού ν-γώνου.
\end{ylh}
% ===================================================================
%  Αρχείο: 36117.tex (Fragment)
%  Αυτόματη μετατροπή μέσω doc2latex.py
% ===================================================================



Στα ορθογώνια τρίγωνα $AB\varGamma $ και $A\varDelta   E$ (γωνία $A$ ορθή) του παρακάτω σχήματος ισχύει $\hat{B} = \hat{\varDelta  } = 30^{\circ}$ και $Z$ το σημείο τομής των πλευρών τους $B\varGamma $ και $\varDelta   E$ αντίστοιχα.

\begin{alist}
\item Να υπολογίσετε τις γωνίες του τετραπλεύρου $AEZ\varGamma $. \monades{13}
\item Να αποδείξετε ότι τα τρίγωνα $\varGamma  Z\varDelta  $ και $EBZ$ είναι ισοσκελή. \monades{12}

\begin{center}
\begin{tikzpicture}[scale=1]
\tkzDefPoint(0,0){A}
\tkzDefPoint(4,0){B}
\tkzDefPoint(0,2.3){C}
\tkzDefPoint(0,4){E}
\tkzDefPoint(2.3,0){Delta}

\tkzInterLL(B,C)(Delta,E) \tkzGetPoint{Z}

\draw[l3] (A)--(B)--(C)--cycle;
\draw[l3] (A)--(Delta)--(E)--cycle;

\tkzDrawPoints(A,B,C,Delta,E,Z)
\tkzLabelPoint[below left](A){$A$}
\tkzLabelPoint[below](B){$B$}
\tkzLabelPoint[left](C){$\varGamma $}
\tkzLabelPoint[below](Delta){$\varDelta  $}
\tkzLabelPoint[above](E){$E$}
\tkzLabelPoint[above right](Z){$Z$}

\tkzMarkRightAngle(C,A,B)
\tkzMarkAngle[size=0.6,mark=none](A,B,C)
\tkzLabelAngle[pos=0.9](A,B,C){$30^\circ$}
\tkzMarkAngle[size=0.6,mark=none](E,Delta,A)
\tkzLabelAngle[pos=0.9](E,Delta,A){$30^\circ$}
\end{tikzpicture}
\end{center}
\end{alist}

\vspace{4mm}

% --- Άσκηση 36118 (36118.tex) ---
\Askhsh{36118}{2}
\begin{ylh}{1}
    \item 3.2. 1ο Κριτήριο ισότητας τριγώνων
    \item 4.2. Τέμνουσα δύο ευθειών - Ευκλείδειο αίτημα
    \item 4.6. Άθροισμα γωνιών τριγώνου.
\end{ylh}
% ===================================================================
%  Αρχείο: 36118.tex (Fragment)
%  Αυτόματη μετατροπή μέσω doc2latex.py
% ===================================================================



Στο παρακάτω σχήμα, οι $A\varDelta  $ και $BE$ είναι παράλληλες. Επιπλέον ισχύουν $A\varDelta  =A\varGamma $, $BE=B\varGamma $ και $\hat{A} = 70^{\circ}$.

\begin{alist}
\item Να υπολογίσετε τις γωνίες των τριγώνων $A\varDelta  \varGamma $ και $B\varGamma  E$. \monades{16}
\item Να αποδείξετε ότι $\varDelta  \hat{\varGamma }E = 90^{\circ}$. \monades{9}

\begin{center}
\begin{tikzpicture}[scale=1]
\tkzDefPoint(0,0){A}
\tkzDefPoint(4,0){B}
\tkzDefPointBy[rotation=center A angle 70](B) \tkzGetPoint{A_dir}
\tkzDefLine[parallel=through A](A,A_dir) \tkzGetPoint{AD_dir}
\tkzDefLine[parallel=through B](A,A_dir) \tkzGetPoint{BE_dir}

\tkzDefPoint(1.5,2.5){C} % Assuming C is generic? No, $AC$=$AD$.
% It's easier to just pick G (gamma) and then construct D and E on parallel lines.
\tkzDefPoint(1.2,1.8){Gamma}
\tkzDrawCircle(A,Gamma) % $AD$ = $AG$
\tkzInterLC(A,AD_dir)(A,Gamma) \tkzGetPoints{Delta}{Delta_alt}
\tkzDrawCircle(B,Gamma) % $BE$ = $BG$
\tkzInterLC(B,BE_dir)(B,Gamma) \tkzGetPoints{E}{E_alt}

\draw[l3] (A)--(Delta);
\draw[l3] (B)--(E);
\draw[l3] (Delta)--(Gamma)--(E);
\draw[l3] (A)--(Gamma)--(B);

\tkzDrawPoints(A,B,Gamma,Delta,E)
\tkzLabelPoint[below](A){$A$}
\tkzLabelPoint[below](B){$B$}
\tkzLabelPoint[right](Gamma){$\varGamma $}
\tkzLabelPoint[above left](Delta){$\varDelta  $}
\tkzLabelPoint[above left](E){$E$}

\tkzMarkSegments[mark=s|](A,Delta A,Gamma)
\tkzMarkSegments[mark=s||](B,E B,Gamma)
\end{tikzpicture}
\end{center}
\end{alist}

\vspace{4mm}

% --- Άσκηση 36163 (36163.tex) ---
\Askhsh{36163}{2}
\begin{ylh}{1}
    \item 3.1. Είδη και στοιχεία τριγώνων
    \item 3.6. Κριτήρια ισότητας ορθογώνιων τριγώνων
    \item 4.6. Άθροισμα γωνιών τριγώνου.
\end{ylh}
% ===================================================================
%  Αρχείο: 36163.tex (Fragment)
%  Αυτόματη μετατροπή μέσω doc2latex.py
% ===================================================================



Στο παρακάτω σχήμα οι γωνίες $\hat{A}$ , $\hat{B}$ είναι ορθές και επιπλέον $A\varDelta  $=$B\varGamma $ και $A\varGamma $=$BE$.

Να αποδείξετε ότι:

\begin{alist}
\item Τα τρίγωνα $A\varGamma \varDelta  $ και $B\varGamma  E$ είναι ίσα. \monades{13}
\item Αν η γωνία $E\hat{\varGamma }B = 40^{\circ}$, τότε το τρίγωνο $\varDelta  \varGamma  E$ είναι ορθογώνιο και ισοσκελές. \monades{12}

\begin{center}
\begin{tikzpicture}[scale=1]
\tkzDefPoint(0,3){A}
\tkzDefPoint(2,3){Delta}
\tkzDefPoint(0,0){B}
\tkzDefPoint(2,0){Gamma}
\tkzDefPoint(0,-3.6){E} % Assuming $AC$ = $BE$. $AC$ = sqrt(4+9) = 3.6.

\draw[l3] (Delta)--(A)--(Gamma)--cycle;
\draw[l3] (Gamma)--(B)--(E)--cycle;
\draw[l3] (Delta)--(Gamma)--(E);

\tkzDrawPoints(A,B,Gamma,Delta,E)
\tkzLabelPoint[left](A){$A$}
\tkzLabelPoint[left](B){$B$}
\tkzLabelPoint[right](Gamma){$\varGamma $}
\tkzLabelPoint[above](Delta){$\varDelta  $}
\tkzLabelPoint[below](E){$E$}

\tkzMarkRightAngle(Delta,A,Gamma)
\tkzMarkRightAngle(Gamma,B,E)
\tkzMarkAngle[size=0.6,mark=none](B,Gamma,E)
\tkzLabelAngle[pos=0.9](B,Gamma,E){$40^\circ$}
\tkzMarkSegments[mark=s|](A,Delta B,Gamma)
\tkzMarkSegments[mark=s||](A,Gamma B,E)
\end{tikzpicture}
\end{center}
\end{alist}

\vspace{4mm}

% --- Άσκηση 36164 (36164.tex) ---
\Askhsh{36164}{2}
\begin{ylh}{1}
    \item 5.2. Παραλληλόγραμμα.
\end{ylh}
% ===================================================================
%  Αρχείο: 36164.tex (Fragment)
%  Αυτόματη μετατροπή μέσω doc2latex.py
% ===================================================================



Δίνεται τρίγωνο $AB\varGamma $ στο οποίο ισχύει $B\varGamma $=2ΑΒ και έστω $M$ το μέσο της $B\varGamma $. Αν η $A\varDelta  $ είναι διάμεσος του τριγώνου $ABM$ και $E$ σημείο στην προέκτασή της ώστε $A\varDelta  $=$\varDelta  E$.

Να αποδείξετε ότι:

\begin{alist}
\item το τετράπλευρο Α$BE$Μ είναι παραλληλόγραμμο, \monades{12}

\item $ME$=$M\varGamma $. \monades{13}


\begin{center}
\begin{tikzpicture}[scale=1]
\tkzDefPoint(0,3){A}
\tkzDefPoint(0,0){B}
\tkzDefPoint(6,0){C}
\tkzDefMidPoint(B,C) \tkzGetPoint{M}
\tkzDefMidPoint(B,M) \tkzGetPoint{Delta}
\tkzDefPointBy[symmetry=center Delta](A) \tkzGetPoint{E}

\draw[l3] (A)--(B)--(C)--cycle;
\draw[l3] (A)--(M)--(E)--cycle;
\draw[l3] (B)--(E);

\tkzDrawPoints(A,B,Gamma,Delta,E,M)
\tkzLabelPoint[above](A){$A$}
\tkzLabelPoint[below](B){$B$}
\tkzLabelPoint[below](C){$\varGamma $}
\tkzLabelPoint[below](M){$M$}
\tkzLabelPoint[below](Delta){$\varDelta  $}
\tkzLabelPoint[below](E){$E$}

\tkzMarkSegments[mark=s|](A,B B,M M,C)
\tkzMarkSegments[mark=s||](B,Delta Delta,M)
\tkzMarkSegments[mark=s|||](A,Delta Delta,E)
\end{tikzpicture}
\end{center}
\end{alist}

\vspace{4mm}

% --- Άσκηση 36165 (36165.tex) ---
\Askhsh{36165}{2}
\begin{ylh}{1}
    \item 3.6. Κριτήρια ισότητας ορθογώνιων τριγώνων
    \item 4.2. Τέμνουσα δύο ευθειών - Ευκλείδειο αίτημα.
\end{ylh}
% ===================================================================
%  Αρχείο: 36165.tex (Fragment)
%  Αυτόματη μετατροπή μέσω doc2latex.py
% ===================================================================



Θεωρούμε τετράγωνο $AB\varGamma \varDelta  $ και σημεία $E$ και $Z$ στις προεκτάσεις των $AB$ (προς το Β) και $B\varGamma $ (προς το Γ) αντίστοιχα, ώστε $BE$=$\varGamma Z$.

Να αποδείξετε ότι:

\begin{alist}
\item τα τρίγωνα $ABZ$ και Α$E\varDelta  $ είναι ίσα, \monades{12}

\item οι γωνίες Ε$\hat{\varDelta  }$Γ και Α$\hat{Ζ}$Β είναι ίσες. \monades{13}


\begin{center}
\begin{tikzpicture}[scale=1]
\tkzDefPoint(0,3){D}
\tkzDefPoint(3,3){C}
\tkzDefPoint(3,0){B}
\tkzDefPoint(0,0){A}

\tkzDefPoint(4.5,0){E} % $BE$ = 1.5
\tkzDefPoint(3,4.5){Z} % $GZ$ = 1.5. Wait, G is Gamma (C). 
% So $CZ$ = 1.5. Extensions A-B-E and B-C-Z.
% C is (3,3). B is (3,0). A is (0,0).
% Extension of $AB$ (A->B) to E: B=(3,0), vector $AB$ is (3,0). E = (4.5, 0).
% Extension of $BC$ (B->C) to Z: C=(3,3), vector $BC$ is (0,3). Z = (3, 4.5).

\draw[l3] (A)--(B)--(C)--(D)--cycle;
\draw[l3] (B)--(E)--(D);
\draw[l3] (C)--(Z)--(A);

\tkzDrawPoints(A,B,C,Delta,E,Z)
\tkzLabelPoint[below left](A){$A$}
\tkzLabelPoint[below](B){$B$}
\tkzLabelPoint[below right](C){$\varGamma $}
\tkzLabelPoint[above left](D){$\varDelta  $}
\tkzLabelPoint[right](E){$E$}
\tkzLabelPoint[above](Z){$Z$}

\tkzMarkSegments[mark=s|](B,E C,Z)
\end{tikzpicture}
\end{center}
\end{alist}

\vspace{4mm}

% --- Άσκηση 36166 (36166.tex) ---
\Askhsh{36166}{2}
\begin{ylh}{1}
    \item 5.2. Παραλληλόγραμμα
    \item 5.10. Τραπέζιο.
\end{ylh}
% ===================================================================
%  Αρχείο: 36166.tex (Fragment)
%  Αυτόματη μετατροπή μέσω doc2latex.py
% ===================================================================



Δίνεται τραπέζιο $AB\varGamma \varDelta  $ ($AB$//$\varGamma \varDelta  $) με $AB$=3, $\varGamma \varDelta  $=4. Θεωρούμε σημείο $E$ στην $AB$ ώστε $AE$=1. Στο τραπέζιο Ε$B\varGamma $Δ θεωρούμε τα Κ και Λ, μέσα των $E\varDelta  $ και $B\varGamma $ αντίστοιχα

\begin{alist}
\item Να υπολογίσετε τη διάμεσο $K\varLambda $ του τραπεζίου Ε$B\varGamma $Δ. \monades{13}

\item Να αποδείξετε ότι το τετράπλευρο $AB\varLambda K$ είναι παραλληλόγραμμο. \monades{12}


\begin{center}
\begin{tikzpicture}[scale=1]
\tkzDefPoint(0,2){A}
\tkzDefPoint(3,2){B}
\tkzDefPoint(1,2){E}
\tkzDefPoint(-0.5,0){Delta}
\tkzDefPoint(3.5,0){C}

\tkzDefMidPoint(E,Delta) \tkzGetPoint{K}
\tkzDefMidPoint(B,C) \tkzGetPoint{L}

\draw[l3] (A)--(B)--(C)--(Delta)--cycle;
\draw[l3] (E)--(Delta);
\draw[l3] (K)--(L);

\tkzDrawPoints(A,B,C,Delta,E,K,L)
\tkzLabelPoint[above](A){$A$}
\tkzLabelPoint[above](B){$B$}
\tkzLabelPoint[below right](C){$\varGamma $}
\tkzLabelPoint[below left](Delta){$\varDelta  $}
\tkzLabelPoint[above](E){$E$}
\tkzLabelPoint[left](K){$K$}
\tkzLabelPoint[right](L){$\Lambda$}

\tkzMarkSegments[mark=s|](E,K K,Delta)
\tkzMarkSegments[mark=s||](B,L L,C)
\end{tikzpicture}
\end{center}
\end{alist}

\vspace{4mm}

% --- Άσκηση 36167 (36167.tex) ---
\Askhsh{36167}{2}
\begin{ylh}{1}
    \item 3.11. Ανισοτικές σχέσεις πλευρών και γωνιών
    \item 4.6. Άθροισμα γωνιών τριγώνου.
\end{ylh}
% ===================================================================
%  Αρχείο: 36167.tex (Fragment)
%  Αυτόματη μετατροπή μέσω doc2latex.py
% ===================================================================



Σε τρίγωνο $AB\varGamma $ ισχύουν $\hat{A} + \hat{\varGamma } = 2\hat{B}$ και $\hat{A} = 3\hat{\varGamma }$.

\begin{alist}
\item Να αποδείξετε ότι η γωνία $\hat{B}$ είναι $60\degree$. \monades{10}
\item Αν το ύψος $A\varDelta  $ και η διχοτόμος $BE$ τέμνονται στο σημείο $Z$, να αποδείξετε ότι το τρίγωνο $AZE$ είναι ισόπλευρο. \monades{15}

\begin{center}
\begin{tikzpicture}[scale=1.2]
\tkzDefPoint(0,0){A}
\tkzDefPoint(0,3){B}
\tkzDefPoint(5.2,0){C} % $AC$ = 3 * cot(30) = 5.196
\tkzDefPointBy[projection=onto B--C](A) \tkzGetPoint{Delta}
\tkzDefLine[bisector](A,B,C) \tkzGetPoint{BE_dir}
\tkzInterLL(B,BE_dir)(A,C) \tkzGetPoint{E}
\tkzInterLL(A,Delta)(B,E) \tkzGetPoint{Z}

\draw[l3] (A)--(B)--(C)--cycle;
\draw[l3] (A)--(Delta);
\draw[l3] (B)--(E);

\tkzDrawPoints(A,B,C,Delta,E,Z)
\tkzLabelPoint[below left](A){$A$}
\tkzLabelPoint[above](B){$B$}
\tkzLabelPoint[right](C){$\varGamma $}
\tkzLabelPoint[above right](Delta){$\varDelta  $}
\tkzLabelPoint[below](E){$E$}
\tkzLabelPoint[above right](Z){$Z$}

\tkzMarkRightAngle(B,A,C)
\tkzMarkRightAngle(A,Delta,B)
\tkzMarkAngle[size=0.6,mark=none](A,B,E)
\tkzMarkAngle[size=0.6,mark=none](E,B,C)
\tkzLabelAngle[pos=0.9](A,C,B){$30^\circ$}
\end{tikzpicture}
\end{center}
\end{alist}

\vspace{4mm}

% --- Άσκηση 36168 (36168.tex) ---
\Askhsh{36168}{2}
\begin{ylh}{1}
    \item 3.6. Κριτήρια ισότητας ορθογώνιων τριγώνων
    \item 3.11. Ανισοτικές σχέσεις πλευρών και γωνιών.
\end{ylh}
% ===================================================================
%  Αρχείο: 36168.tex (Fragment)
%  Αυτόματη μετατροπή μέσω doc2latex.py
% ===================================================================



Στο παρακάτω σχήμα το τρίγωνο $AB\varGamma $ είναι ορθογώνιο με ορθή τη γωνία $\hat{A}$ και η γωνία $\hat{\varGamma }$ είναι μικρότερη της γωνίας $\hat{B}$. Η $B\varDelta  $ είναι διχοτόμος της γωνίας $\hat{B}$ και η $\varDelta  E$ είναι κάθετη στην $B\varGamma $.

Να αποδείξετε ότι:

\begin{alist}
\item $A\varDelta  $=$\varDelta  E$, \monades{8}

\item $A\varDelta  $ < $\varDelta  \varGamma $, \monades{9}

\item $A\varGamma $\textgreater $AB$. \monades{8}


\begin{center}
\begin{tikzpicture}[scale=1.2]
\tkzDefPoint(0,0){A}
\tkzDefPoint(3,0){C}
\tkzDefPoint(0,1.732){B} % tan(60) = sqrt(3)
\tkzDefLine[bisector](A,B,C) \tkzGetPoint{BD_dir}
\tkzInterLL(B,BD_dir)(A,C) \tkzGetPoint{D}
\tkzDefPointBy[projection=onto B--C](D) \tkzGetPoint{E}

\draw[l3] (A)--(B)--(C)--cycle;
\draw[l3] (B)--(D);
\draw[l3] (D)--(E);

\tkzDrawPoints(A,B,C,Delta,E)
\tkzLabelPoint[below left](A){$A$}
\tkzLabelPoint[above](B){$B$}
\tkzLabelPoint[below](C){$\varGamma $}
\tkzLabelPoint[below](D){$\varDelta  $}
\tkzLabelPoint[above right](E){$E$}

\tkzMarkRightAngle(C,A,B)
\tkzMarkRightAngle(D,E,C)
\tkzMarkAngle[size=0.5,mark=none](A,B,D)
\tkzMarkAngle[size=0.5,mark=none](D,B,C)
\end{tikzpicture}
\end{center}
\end{alist}

\vspace{4mm}

% --- Άσκηση 36169 (36169.tex) ---
\Askhsh{36169}{2}
\begin{ylh}{1}
    \item 3.2. 1ο Κριτήριο ισότητας τριγώνων
    \item 4.6. Άθροισμα γωνιών τριγώνου
    \item 5.6. Εφαρμογές στα τρίγωνα.
\end{ylh}
% ===================================================================
%  Αρχείο: 36169.tex (Fragment)
%  Αυτόματη μετατροπή μέσω doc2latex.py
% ===================================================================



Δίνεται ορθογώνιο τρίγωνο $AB\varGamma $ με $\hat{A} = 90^{\circ}$, $\hat{B} = 35^{\circ}$ και $M$ το μέσο της $B\varGamma $.

\begin{alist}
\item Να υπολογίσετε τη γωνία $\hat{\varGamma }$. \monades{10}
\item Να υπολογίσετε τις γωνίες του τριγώνου $AMB$. \monades{15}

\begin{center}
\begin{tikzpicture}[scale=1]
\tkzDefPoint(0,0){A}
\tkzDefPoint(0,2.8){B} % tan(35) = 2/$AC$ -> $AC$ = 2/0.7 = 2.8
\tkzDefPoint(4,0){C}
\tkzDefMidPoint(B,C) \tkzGetPoint{M}

\draw[l3] (A)--(B)--(C)--cycle;
\draw[l3] (A)--(M);

\tkzDrawPoints(A,B,C,M)
\tkzLabelPoint[below left](A){$A$}
\tkzLabelPoint[above](B){$B$}
\tkzLabelPoint[right](C){$\varGamma $}
\tkzLabelPoint[above right](M){$M$}

\tkzMarkRightAngle(B,A,C)
\tkzMarkAngle[size=0.6,mark=none](A,B,C)
\tkzLabelAngle[pos=0.9](A,B,C){$35^\circ$}
\tkzMarkSegments[mark=s|](B,M M,C A,M)
\end{tikzpicture}
\end{center}
\end{alist}

\vspace{4mm}

% --- Άσκηση 36170 (36170.tex) ---
\Askhsh{36170}{2}
\begin{ylh}{1}
    \item 3.2. 1ο Κριτήριο ισότητας τριγώνων.
\end{ylh}
% ===================================================================
%  Αρχείο: 36170.tex (Fragment)
%  Αυτόματη μετατροπή μέσω doc2latex.py
% ===================================================================



Δίνεται ισοσκελές τρίγωνο $AB\varGamma $ με $AB=A\varGamma $. Στις προεκτάσεις των πλευρών $BA$ και $\varGamma  A$ (προς το Α) θεωρούμε τα σημεία $E$ και $\varDelta$ αντίστοιχα τέτοια ώστε $AE$ = $A\varDelta  $.

Να αποδείξετε ότι:

\begin{alist}
\item $BE$ = $\varGamma \varDelta  $, \monades{6}

\item $B\varDelta  $ = $\varGamma E$, \monades{10}

\item $\varDelta  \hat{B}\varGamma  = Ε\hat{\varGamma }B$. \monades{9}
\end{alist}

\vspace{4mm}

% --- Άσκηση 36171 (36171.tex) ---
\Askhsh{36171}{2}
\begin{ylh}{1}
    \item 4.6. Άθροισμα γωνιών τριγώνου
    \item 5.9. Μια ιδιότητα του ορθογώνιου τριγώνου.
\end{ylh}
% ===================================================================
%  Αρχείο: 36171.tex (Fragment)
%  Αυτόματη μετατροπή μέσω doc2latex.py
% ===================================================================



Δίνεται ορθογώνιο τρίγωνο $AB\varGamma $ με τη γωνία $A$ ορθή, $2\hat{\varGamma } = \hat{B}$ και το ύψος του $A\varDelta  $.

\begin{alist}
\item Να υπολογιστούν οι οξείες γωνίες του τριγώνου $AB\varGamma $. \monades{9}
\item Να υπολογιστεί η γωνία $B\hat{A}\varDelta  $. \monades{7}
\item Να αποδείξετε ότι $B\varDelta   = \frac{AB}{2}$. \monades{9}

\begin{center}
\begin{tikzpicture}[scale=1.2]
\tkzDefPoint(0,0){A}
\tkzDefPoint(0,2.6){B}
\tkzDefPoint(4.5,0){C} % tan(30) = 2.6/$AC$ -> $AC$ = 4.5
\tkzDefPointBy[projection=onto B--C](A) \tkzGetPoint{Delta}

\draw[l3] (A)--(B)--(C)--cycle;
\draw[l3] (A)--(Delta);

\tkzDrawPoints(A,B,C,Delta)
\tkzLabelPoint[below left](A){$A$}
\tkzLabelPoint[above](B){$B$}
\tkzLabelPoint[right](C){$\varGamma $}
\tkzLabelPoint[above right](Delta){$\varDelta  $}

\tkzMarkRightAngle(B,A,C)
\tkzMarkRightAngle(A,Delta,B)
\tkzMarkAngle[size=0.6,mark=none](A,B,C)
\tkzLabelAngle[pos=0.9](A,B,C){$60^\circ$}
\tkzMarkAngle[size=0.8,mark=none](A,C,B)
\tkzLabelAngle[pos=1.2](A,C,B){$30^\circ$}
\end{tikzpicture}
\end{center}
\end{alist}

\vspace{4mm}

% --- Άσκηση 36172 (36172.tex) ---
\Askhsh{36172}{2}
\begin{ylh}{1}
    \item 3.2. 1ο Κριτήριο ισότητας τριγώνων
    \item 3.11. Ανισοτικές σχέσεις πλευρών και γωνιών
    \item 4.2. Τέμνουσα δύο ευθειών - Ευκλείδειο αίτημα
    \item 4.6. Άθροισμα γωνιών τριγώνου.
\end{ylh}
% ===================================================================
%  Αρχείο: 36172.tex (Fragment)
%  Αυτόματη μετατροπή μέσω doc2latex.py
% ===================================================================



Δίνεται τραπέζιο $AB\varGamma \varDelta  $ με $AB$ // $\varGamma \varDelta  $ και $B\varDelta  $ = $B\varGamma $. Αν $\varDelta  \hat{B}\varGamma  = 110\degree$ και $A\hat{\varDelta  }B = 25\degree$ να υπολογίσετε:

\begin{alist}
\item τη γωνία $\hat{\varGamma }$, \monades{11}

\item τη γωνία $\hat{A}$. \monades{14}


\begin{center}
\begin{tikzpicture}[scale=0.8]
\tkzDefPoint(0,0){Delta}
\tkzDefPoint(5,0){C}
\tkzDefPointBy[rotation=center Delta angle 35](C) \tkzGetPoint{B_dir}
\tkzDefCircle[R](Delta,5) \tkzGetPoint{dummy}
\tkzInterLC(Delta,B_dir)(Delta,dummy) \tkzGetPoints{B_dummy}{B} % $BD$=$BC$? Wait.
% If $BD$=$BC$ and $CD$=5, then tri $BCD$ is isosceles.
% Base $CD$=5. Angles at C and D are 35.
% B is intersection of lines from C and D at 35.
\tkzDefPointBy[rotation=center Delta angle 35](C) \tkzGetPoint{BD_line}
\tkzDefPointBy[rotation=center C angle -35](Delta) \tkzGetPoint{BC_line}
\tkzInterLL(Delta,BD_line)(C,BC_line) \tkzGetPoint{B}

\tkzDefLine[parallel=through B](Delta,C) \tkzGetPoint{AB_line}
\tkzDefPointBy[rotation=center Delta angle 60](C) \tkzGetPoint{AD_line} % Total D = 35+25=60.
\tkzInterLL(B,AB_line)(Delta,AD_line) \tkzGetPoint{A}

\draw[l3] (A)--(B)--(C)--(Delta)--cycle;
\draw[l3] (B)--(Delta);

\tkzDrawPoints(A,B,C,Delta)
\tkzLabelPoint[left](A){$A$}
\tkzLabelPoint[above](B){$B$}
\tkzLabelPoint[right](C){$\varGamma $}
\tkzLabelPoint[below](Delta){$\varDelta  $}

\tkzMarkAngle[size=0.6,mark=none](C,B,Delta)
\tkzLabelAngle[pos=0.9](C,B,Delta){$110^\circ$}
\tkzMarkAngle[size=0.8,mark=none](A,Delta,B)
\tkzLabelAngle[pos=1.2](A,Delta,B){$25^\circ$}
\tkzMarkSegments[mark=s|](Delta,B B,C)
\end{tikzpicture}
\end{center}
\end{alist}

\vspace{4mm}

% --- Άσκηση 36173 (36173.tex) ---
\Askhsh{36173}{2}
\begin{ylh}{1}
    \item 3.1. Είδη και στοιχεία τριγώνων
    \item 3.2. 1ο Κριτήριο ισότητας τριγώνων
    \item 4.6. Άθροισμα γωνιών τριγώνου
    \item 5.5. Τετράγωνο.
\end{ylh}
% ===================================================================
%  Αρχείο: 36173.tex (Fragment)
%  Αυτόματη μετατροπή μέσω doc2latex.py
% ===================================================================



Δίνεται ισόπλευρο τρίγωνο $AB\varGamma $ και εκτός αυτού κατασκευάζουμε τετράγωνο $B\varGamma \varDelta   E$.

\begin{alist}
\item Να υπολογίσετε τις γωνίες:
\begin{alist}
\item $A\hat{B}E$, \monades{8}
\item $B\hat{E}A$. \monades{9}
\end{alist}
\item Να αποδείξετε ότι το τρίγωνο $AE\varDelta  $ είναι ισοσκελές. \monades{8}

\begin{center}
\begin{tikzpicture}[scale=1]
\tkzDefPoint(0,0){B}
\tkzDefPoint(3,0){C}
\tkzDefTriangle[equilateral](B,C) \tkzGetPoint{A}
\tkzDefSquare(B,C) \tkzGetPoints{Delta}{E}

\draw[l3] (A)--(B)--(C)--cycle;
\draw[l3] (B)--(E)--(Delta)--(C);
\draw[l3] (A)--(E)--(Delta)--cycle;

\tkzDrawPoints(A,B,C,Delta,E)
\tkzLabelPoint[above](A){$A$}
\tkzLabelPoint[left](B){$B$}
\tkzLabelPoint[right](C){$\varGamma $}
\tkzLabelPoint[below right](Delta){$\varDelta  $}
\tkzLabelPoint[below left](E){$E$}

\tkzMarkSegments[mark=s|](A,B B,C A,C B,E E,Delta Delta,C)
\end{tikzpicture}
\end{center}
\end{alist}

\vspace{4mm}

% --- Άσκηση 36174 (36174.tex) ---
\Askhsh{36174}{2}
\begin{ylh}{1}
    \item 3.1. Είδη και στοιχεία τριγώνων
    \item 3.2. 1ο Κριτήριο ισότητας τριγώνων
    \item 4.6. Άθροισμα γωνιών τριγώνου
    \item 5.5. Τετράγωνο.
\end{ylh}
% ===================================================================
%  Αρχείο: 36174.tex (Fragment)
%  Αυτόματη μετατροπή μέσω doc2latex.py
% ===================================================================



Δίνεται ισοσκελές τρίγωνο $AB\varGamma $ με $AB=A\varGamma $. Κατασκευάζουμε εξωτερικά του τριγώνου το τετράγωνο Α$B\varDelta  $Ε.

Να αποδείξετε ότι:

\begin{alist}
\item το τρίγωνο $A\varGamma  E$ είναι ισοσκελές, \monades{10}
\item $2 \cdot E\hat{\varGamma }A = 90^{\circ} - B\hat{A}\varGamma $. \monades{15}

\begin{center}
\begin{tikzpicture}[scale=0.8]
\tkzDefPoint(0,0){A}
\tkzDefPoint(3,0){B}
\tkzDefPointBy[rotation=center A angle 40](B) \tkzGetPoint{C_dir}
\tkzDefPointBy[homothety=center A ratio 1](C_dir) \tkzGetPoint{C}
\tkzDefSquare(A,B) \tkzGetPoints{Delta}{E}

\draw[l3] (A)--(B)--(C)--cycle;
\draw[l3] (A)--(E)--(Delta)--(B);
\draw[l3] (E)--(C);

\tkzDrawPoints(A,B,C,Delta,E)
\tkzLabelPoint[left](A){$A$}
\tkzLabelPoint[right](B){$B$}
\tkzLabelPoint[above](C){$\varGamma $}
\tkzLabelPoint[below](Delta){$\varDelta  $}
\tkzLabelPoint[left](E){$E$}

\tkzMarkSegments[mark=s|](A,B A,C A,E E,Delta Delta,B)
\end{tikzpicture}
\end{center}
\end{alist}

\vspace{4mm}

% --- Άσκηση 36175 (36175.tex) ---
\Askhsh{36175}{2}
\begin{ylh}{1}
    \item 3.1. Είδη και στοιχεία τριγώνων
    \item 4.2. Τέμνουσα δύο ευθειών - Ευκλείδειο αίτημα
    \item 5.2. Παραλληλόγραμμα
    \item 5.3. Ορθογώνιο.
\end{ylh}
% ===================================================================
%  Αρχείο: 36175.tex (Fragment)
%  Αυτόματη μετατροπή μέσω doc2latex.py
% ===================================================================



Στο παρακάτω σχήμα το τετράπλευρο $AB\varGamma \varDelta  $ είναι παραλληλόγραμμο και το $A\varGamma \varDelta  E$ είναι ορθογώνιο.

Να αποδείξετε ότι:

\begin{alist}
\item το σημείο $A$ είναι μέσο του $BE$, \monades{8}
\item το τρίγωνο $BE\varGamma $ είναι ισοσκελές, \monades{9}
\item $B\hat{\varGamma }A = A\hat{\varDelta  }E$. \monades{8}

\begin{center}
\begin{tikzpicture}[scale=0.8]
\tkzDefPoint(0,0){A}
\tkzDefPoint(3,0){B}
\tkzDefPoint(1,2.5){Gamma}
\tkzDefPointBy[translation=from B to A](Gamma) \tkzGetPoint{Delta}
\tkzDefLine[parallel=through A](Delta,Gamma) \tkzGetPoint{E_line}
\tkzDefPointBy[translation=from Gamma to Delta](A) \tkzGetPoint{E}

\draw[l3] (A)--(B)--(Gamma)--(Delta)--cycle;
\draw[l3] (A)--(E)--(Delta)--(Gamma)--cycle;
\draw[l3] (E)--(B)--(Gamma);

\tkzDrawPoints(A,B,Gamma,Delta,E)
\tkzLabelPoint[below](A){$A$}
\tkzLabelPoint[below right](B){$B$}
\tkzLabelPoint[above right](Gamma){$\varGamma $}
\tkzLabelPoint[above left](Delta){$\varDelta  $}
\tkzLabelPoint[below left](E){$E$}

\tkzMarkRightAngle(E,A,Gamma)
\tkzMarkSegments[mark=s|](A,B Delta,Gamma A,E)
\end{tikzpicture}
\end{center}
\end{alist}

\vspace{4mm}

% --- Άσκηση 36224 (36224.tex) ---
\Askhsh{36224}{2}
\begin{ylh}{1}
    \item 3.2. 1ο Κριτήριο ισότητας τριγώνων
    \item 4.6. Άθροισμα γωνιών τριγώνου
    \item 5.6. Εφαρμογές στα τρίγωνα
    \item 5.9. Μια ιδιότητα του ορθογώνιου τριγώνου.
\end{ylh}
% ===================================================================
%  Αρχείο: 36224.tex (Fragment)
%  Αυτόματη μετατροπή μέσω doc2latex.py
% ===================================================================



Έστω ισοσκελές τρίγωνο $AB\varGamma $ με $AB=Α\varGamma $ και γωνία $\hat{B}$ ίση με 30\textsuperscript{Ο}. Θεωρούμε $\varDelta$ και $E$ τα μέσα των $A\varGamma $ και $B\varGamma $ αντίστοιχα.

Να αποδείξετε ότι:

\begin{alist}
\item το τρίγωνο $E\varDelta  \varGamma $ είναι ισοσκελές και να υπολογίσετε τις γωνίες του. \monades{16}
\item το τρίγωνο $A\varDelta   E$ είναι ισόπλευρο. \monades{9}

\begin{center}
\begin{tikzpicture}[scale=0.8]
\tkzDefPoint(0,0){B}
\tkzDefPoint(6,0){C}
\tkzDefPointBy[rotation=center B angle 30](C) \tkzGetPoint{B_dir}
\tkzDefPointBy[rotation=center C angle -30](B) \tkzGetPoint{C_dir}
\tkzInterLL(B,B_dir)(C,C_dir) \tkzGetPoint{A}
\tkzDefMidPoint(A,C) \tkzGetPoint{Delta}
\tkzDefMidPoint(B,C) \tkzGetPoint{E}

\draw[l3] (A)--(B)--(C)--cycle;
\draw[l3] (Delta)--(E);
\draw[l3] (A)--(E);

\tkzDrawPoints(A,B,C,Delta,E)
\tkzLabelPoint[above](A){$A$}
\tkzLabelPoint[left](B){$B$}
\tkzLabelPoint[right](C){$\varGamma $}
\tkzLabelPoint[above right](Delta){$\varDelta  $}
\tkzLabelPoint[below](E){$E$}

\tkzMarkAngle[size=0.8,mark=none](C,B,A)
\tkzLabelAngle[pos=1.2](C,B,A){$30^\circ$}
\tkzMarkSegments[mark=s|](A,Delta Delta,C)
\tkzMarkSegments[mark=s||](B,E E,C)
\end{tikzpicture}
\end{center}
\end{alist}

\vspace{4mm}

% --- Άσκηση 36225 (36225.tex) ---
\Askhsh{36225}{2}
\begin{ylh}{1}
    \item 3.2. 1ο Κριτήριο ισότητας τριγώνων
    \item 5.2. Παραλληλόγραμμα.
\end{ylh}
% ===================================================================
%  Αρχείο: 36225.tex (Fragment)
%  Αυτόματη μετατροπή μέσω doc2latex.py
% ===================================================================



Έστω παραλληλόγραμμο $AB\varGamma \varDelta  $. Προεκτείνουμε την πλευρά $BA$ (προς το Α) και την πλευρά $\varDelta  \varGamma $ (προς το Γ) κατά τμήματα $AE$ = $AB$ και $\varGamma Z$ = $\varDelta  \varGamma $.

Να αποδείξετε ότι:

\begin{alist}
\item τα τρίγωνα $A\varDelta  E$ και $B\varGamma $Ζ είναι ίσα, \monades{13}

\item το τετράπλευρο $EBZ\varDelta  $ είναι παραλληλόγραμμο. \monades{12}


\begin{center}
\begin{tikzpicture}[scale=0.8]
\tkzDefPoint(0,0){A}
\tkzDefPoint(3,0){B}
\tkzDefPoint(4,2){C}
\tkzDefPointBy[translation=from B to A](C) \tkzGetPoint{D}
\tkzDefPointBy[symmetry=center A](B) \tkzGetPoint{E}
\tkzDefPointBy[symmetry=center C](D) \tkzGetPoint{Z}

\draw[l3] (A)--(B)--(C)--(D)--cycle;
\draw[l3, dashed] (A)--(E);
\draw[l3, dashed] (C)--(Z);
\draw[l3] (E)--(B)--(Z)--(D)--cycle;

\tkzDrawPoints(A,B,C,D,E,Z)
\tkzLabelPoint[below](A){$A$}
\tkzLabelPoint[below](B){$B$}
\tkzLabelPoint[above](C){$\varGamma $}
\tkzLabelPoint[above](D){$\varDelta  $}
\tkzLabelPoint[left](E){$E$}
\tkzLabelPoint[right](Z){$Z$}

\tkzMarkSegments[mark=s|](E,A A,B D,C C,Z)
\end{tikzpicture}
\end{center}
\end{alist}

\vspace{4mm}

% --- Άσκηση 36226 (36226.tex) ---
\Askhsh{36226}{2}
\begin{ylh}{1}
    \item 3.7. Κύκλος - Μεσοκάθετος - Διχοτόμος.
\end{ylh}
% ===================================================================
%  Αρχείο: 36226.tex (Fragment)
%  Αυτόματη μετατροπή μέσω doc2latex.py
% ===================================================================



Στο παρακάτω σχήμα έχουμε το χάρτη μίας περιοχής όπου είναι κρυμμένος ένας θησαυρός. Οι ημιευθείες Αx και Αy παριστάνουν δύο ποτάμια και στα σημεία $B$ και $\varGamma$ βρίσκονται δύο πλατάνια. Ο πλάτανος που βρίσκεται στο σημείο $B$ έχει μικρότερη απόσταση από το σημείο $A$, σε σχέση με την απόσταση που έχει από το σημείο $A$ ο πλάτανος που βρίσκεται στο σημείο $\varGamma $.

Να προσδιορίσετε γεωμετρικά τις δυνατές θέσεις του θησαυρού, αν είναι γνωστό ότι:

\begin{alist}
\item ο θησαυρός ισαπέχει από τα δύο πλατάνια. \monades{9}

\item ο θησαυρός ισαπέχει από τα δύο ποτάμια. \monades{9}

\item ο θησαυρός ισαπέχει από τα δύο πλατάνια και ισαπέχει και από τα δύο ποτάμια.

\monades{7}

Να αιτιολογήσετε την απάντησή σας σε κάθε περίπτωση.


\begin{center}
\begin{tikzpicture}[scale=1]
\tkzDefPoint(0,0){A}
\tkzDefPoint(5,0){x_axis}
\tkzDefPointBy[rotation=center A angle 50](x_axis) \tkzGetPoint{y_axis}
\tkzDefPoint(2,0){B}
\tkzDefPointBy[rotation=center A angle 50](x_axis) \tkzGetPoint{C_dir}
\tkzDefPointBy[homothety=center A ratio 0.8](C_dir) \tkzGetPoint{C}

\draw[l3, ->] (A)--(x_axis) node[right]{$x$};
\draw[l3, ->] (A)--(y_axis) node[above]{$y$};
\tkzDrawPoints(B,C)
\tkzLabelPoint[below](B){$B$}
\tkzLabelPoint[above left](C){$\varGamma $}

\node at (1,-1) {Ποτάμι $Ax$};
\node at (-1,2) {Ποτάμι $Ay$};
\end{tikzpicture}
\end{center}
\end{alist}

\vspace{4mm}

% --- Άσκηση 36228 (36228.tex) ---
\Askhsh{36228}{2}
\begin{ylh}{1}
    \item 3.2. 1ο Κριτήριο ισότητας τριγώνων
    \item 3.11. Ανισοτικές σχέσεις πλευρών και γωνιών
    \item 4.6. Άθροισμα γωνιών τριγώνου.
\end{ylh}
% ===================================================================
%  Αρχείο: 36228.tex (Fragment)
%  Αυτόματη μετατροπή μέσω doc2latex.py
% ===================================================================



Δίνεται ισόπλευρο τρίγωνο $AB\varGamma $. Θεωρούμε σημείο $E$ στην προέκταση της $BA$ (προς το Α) και σημείο $\varDelta  $ στο εσωτερικό της πλευράς $A\varGamma $, ώστε $AE$=$A\varDelta  $.

\begin{alist}
\item Να υπολογίσετε τις γωνίες του τριγώνου $A\varDelta  E$. \monades{10}

\item Αν $Z$ είναι το σημείο τομής της προέκτασης της $E\varDelta  $ (προς το Δ) με την $B\varGamma $, να αποδείξετε ότι η $EZ$ είναι κάθετη στην $B\varGamma $. \monades{15}


\begin{center}
\begin{tikzpicture}[scale=1]
\tkzDefPoint(0,0){B}
\tkzDefPoint(4,0){C}
\tkzDefTriangle[equilateral](B,C) \tkzGetPoint{A}
\tkzDefMidPoint(A,C) \tkzGetPoint{Delta}
\tkzDefPointBy[homothety=center A ratio -0.5](B) \tkzGetPoint{E} % $AE$ = 0.5 $AB$ = 2.

\tkzInterLL(E,Delta)(B,C) \tkzGetPoint{Z}

\draw[l3] (A)--(B)--(C)--cycle;
\draw[l3] (E)--(B);
\draw[l3] (E)--(Z);

\tkzDrawPoints(A,B,C,Delta,E,Z)
\tkzLabelPoint[right](A){$A$}
\tkzLabelPoint[below](B){$B$}
\tkzLabelPoint[below](C){$\varGamma $}
\tkzLabelPoint[right](Delta){$\varDelta  $}
\tkzLabelPoint[above left](E){$E$}
\tkzLabelPoint[below](Z){$Z$}

\tkzMarkSegments[mark=s|](A,Delta A,E)
\tkzMarkRightAngle(E,Z,B)
\end{tikzpicture}
\end{center}
\end{alist}

\vspace{4mm}

% --- Άσκηση 36327 (36327.tex) ---
\Askhsh{36327}{2}
\begin{ylh}{1}
    \item 3.1. Είδη και στοιχεία τριγώνων
    \item 3.4. 3ο Κριτήριο ισότητας τριγώνων
    \item 4.2. Τέμνουσα δύο ευθειών - Ευκλείδειο αίτημα
    \item 5.2. Παραλληλόγραμμα.
\end{ylh}
% ===================================================================
%  Αρχείο: 36327.tex (Fragment)
%  Αυτόματη μετατροπή μέσω doc2latex.py
% ===================================================================



Δίνονται τα παραλληλόγραμμα Α$B\varDelta  $Γ και $B\varDelta  $ΕΖ.

Να αποδείξετε ότι:

\begin{alist}
\item το τετράπλευρο $A\varGamma EZ$ είναι παραλληλόγραμμο, \monades{13}

\item $Α\overset{\land}{B}Ζ = \varGamma \overset{\land}{\varDelta  }Ε$. \monades{12}


\begin{center}
\begin{tikzpicture}[scale=0.8]
\tkzDefPoint(0,2){A}
\tkzDefPoint(3,2){B}
\tkzDefPoint(1,0){Gamma}
\tkzDefPoint(4,0){Delta}
\tkzDefPoint(7,0){E}
\tkzDefPoint(6,2){Z}

\draw[l3] (A)--(B)--(Delta)--(Gamma)--cycle;
\draw[l3] (B)--(Delta)--(E)--(Z)--cycle;
\draw[l3, dashed] (A)--(Gamma)--(E)--(Z)--cycle;

\tkzDrawPoints(A,B,Gamma,Delta,E,Z)
\tkzLabelPoint[above left](A){$A$}
\tkzLabelPoint[above](B){$B$}
\tkzLabelPoint[below left](Gamma){$\varGamma $}
\tkzLabelPoint[below](Delta){$\varDelta  $}
\tkzLabelPoint[below right](E){$E$}
\tkzLabelPoint[above right](Z){$Z$}
\end{tikzpicture}
\end{center}
\end{alist}

\vspace{4mm}

% --- Άσκηση 36328 (36328.tex) ---
\Askhsh{36328}{2}
\begin{ylh}{1}
    \item 3.2. 1ο Κριτήριο ισότητας τριγώνων
    \item 4.2. Τέμνουσα δύο ευθειών - Ευκλείδειο αίτημα
    \item 5.9. Μια ιδιότητα του ορθογώνιου τριγώνου.
\end{ylh}
% ===================================================================
%  Αρχείο: 36328.tex (Fragment)
%  Αυτόματη μετατροπή μέσω doc2latex.py
% ===================================================================



Δίνεται ορθογώνιο τρίγωνο $AB\varGamma $ με τη γωνία $A$ ορθή και $M$ το μέσο της $B\varGamma $. Φέρουμε ημιευθεία Αx παράλληλη στη $B\varGamma $ (στο ημιεπίπεδο που ορίζει η $AM$ με το σημείο $\varGamma $).

Να αποδείξετε ότι:

\begin{alist}
\item $M\hat{A}\varGamma  = M\hat{\varGamma }A$, \monades{12}
\item η $A\varGamma $ είναι διχοτόμος της γωνίας $M\hat{A}x$. \monades{13}

\begin{center}
\begin{tikzpicture}[scale=1]
\tkzDefPoint(0,0){A}
\tkzDefPoint(0,3){B}
\tkzDefPoint(4,0){C}
\tkzDefMidPoint(B,C) \tkzGetPoint{M}
\tkzDefLine[parallel=through A](B,C) \tkzGetPoint{x_dir}

\draw[l3] (A)--(B)--(C)--cycle;
\draw[l3] (A)--(M);
\draw[l3, ->] (A)--(x_dir) node[right]{$x$};

\tkzDrawPoints(A,B,C,M)
\tkzLabelPoint[below left](A){$A$}
\tkzLabelPoint[above](B){$B$}
\tkzLabelPoint[right](C){$\varGamma $}
\tkzLabelPoint[above right](M){$M$}

\tkzMarkRightAngle(B,A,C)
\tkzMarkSegments[mark=s|](B,M M,C A,M)
\end{tikzpicture}
\end{center}
\end{alist}

\vspace{4mm}

% --- Άσκηση 36329 (36329.tex) ---
\Askhsh{36329}{2}
\begin{ylh}{1}
    \item 3.1. Είδη και στοιχεία τριγώνων
    \item 3.6. Κριτήρια ισότητας ορθογώνιων τριγώνων
    \item 4.6. Άθροισμα γωνιών τριγώνου.
\end{ylh}
% ===================================================================
%  Αρχείο: 36329.tex (Fragment)
%  Αυτόματη μετατροπή μέσω doc2latex.py
% ===================================================================



Δίνεται τρίγωνο $AB\varGamma $ και $M\varDelta  $, $NE$ οι μεσοκάθετοι των πλευρών του $AB$, $A\varGamma $ αντίστοιχα.

Να αποδείξετε ότι:

\begin{alist}
\item Αν είναι $M\varDelta  $ = $NE$, τότε το τρίγωνο $AB\varGamma $ είναι ισοσκελές. \monades{12}

\item Αν είναι $AB=A\varGamma $, τότε $M\varDelta  $ = $NE$. \monades{13}


\begin{center}
\begin{tikzpicture}[scale=1]
\tkzDefPoint(0,0){B}
\tkzDefPoint(4,0){C}
\tkzDefPoint(1.5,3){A}
\tkzDefMidPoint(A,B) \tkzGetPoint{D}
\tkzDefMidPoint(A,C) \tkzGetPoint{E}

\tkzDefLine[orthogonal=through D](A,B) \tkzGetPoint{D_perp}
\tkzDefLine[orthogonal=through E](A,C) \tkzGetPoint{E_perp}

% Let's pick points M, N on these lines.
\tkzDefPointBy[homothety=center D ratio 1](D_perp) \tkzGetPoint{M}
\tkzDefPointBy[homothety=center E ratio 1](E_perp) \tkzGetPoint{N}

\draw[l3] (A)--(B)--(C)--cycle;
\draw[l3, dashed] (D)--(M);
\draw[l3, dashed] (E)--(N);

\tkzDrawPoints(A,B,C,D,E,M,N)
\tkzLabelPoint[above](A){$A$}
\tkzLabelPoint[below left](B){$B$}
\tkzLabelPoint[below right](C){$\varGamma $}
\tkzLabelPoint[left](D){$\varDelta  $}
\tkzLabelPoint[right](E){$E$}
\tkzLabelPoint[above left](M){$M$}
\tkzLabelPoint[above right](N){$N$}

\tkzMarkRightAngle(M,D,A)
\tkzMarkRightAngle(N,E,A)
\tkzMarkSegments[mark=s|](A,D D,B)
\tkzMarkSegments[mark=s||](A,E E,C)
\end{tikzpicture}
\end{center}
\end{alist}

\vspace{4mm}

% --- Άσκηση 36330 (36330.tex) ---
\Askhsh{36330}{2}
\begin{ylh}{1}
    \item 3.1. Είδη και στοιχεία τριγώνων
    \item 3.2. 1ο Κριτήριο ισότητας τριγώνων
    \item 3.6. Κριτήρια ισότητας ορθογώνιων τριγώνων.
\end{ylh}
% ===================================================================
%  Αρχείο: 36330.tex (Fragment)
%  Αυτόματη μετατροπή μέσω doc2latex.py
% ===================================================================



Δίνεται τρίγωνο $AB\varGamma $ και από σημείο $M$ της πλευράς $B\varGamma $ φέρουμε τα κάθετα τμήματα $M\varDelta  $ και $ME$ στις πλευρές $AB$ και $A\varGamma $ αντίστοιχα.

Να αποδείξετε ότι:

\begin{alist}
\item Αν είναι $M\varDelta  $ = $ME$, τότε τα τρίγωνα $AM\varDelta  $ και $AME$ είναι ίσα. \monades{13}

\item Αν είναι $AB=A\varGamma $ και $M$ μέσο του $B\varGamma $, τότε $M\varDelta  $ = $ME$. \monades{12}


\begin{center}
\begin{tikzpicture}[scale=1]
\tkzDefPoint(0,0){B}
\tkzDefPoint(5,0){C}
\tkzDefPoint(1,3){A}
\tkzDefPoint(2,0){M}
\tkzDefPointBy[projection=onto A--B](M) \tkzGetPoint{Delta}
\tkzDefPointBy[projection=onto A--C](M) \tkzGetPoint{E}

\draw[l3] (A)--(B)--(C)--cycle;
\draw[l3] (M)--(Delta);
\draw[l3] (M)--(E);
\draw[l3] (A)--(M);

\tkzDrawPoints(A,B,C,M,Delta,E)
\tkzLabelPoint[above](A){$A$}
\tkzLabelPoint[below left](B){$B$}
\tkzLabelPoint[below right](C){$\varGamma $}
\tkzLabelPoint[below](M){$M$}
\tkzLabelPoint[left](Delta){$\varDelta  $}
\tkzLabelPoint[right](E){$E$}

\tkzMarkRightAngle(M,Delta,A)
\tkzMarkRightAngle(M,E,A)
\end{tikzpicture}
\end{center}
\end{alist}

\vspace{4mm}

% --- Άσκηση 36331 (36331.tex) ---
\Askhsh{36331}{2}
\begin{ylh}{1}
    \item 3.2. 1ο Κριτήριο ισότητας τριγώνων
    \item 3.6. Κριτήρια ισότητας ορθογώνιων τριγώνων
    \item 3.11. Ανισοτικές σχέσεις πλευρών και γωνιών.
\end{ylh}
% ===================================================================
%  Αρχείο: 36331.tex (Fragment)
%  Αυτόματη μετατροπή μέσω doc2latex.py
% ===================================================================



Δίνεται ορθογώνιο τρίγωνο $AB\varGamma $ με τη γωνία $A$ ορθή και από το μέσο $M$ της πλευράς $B\varGamma $ φέρουμε τα κάθετα τμήματα $M\varDelta  $ και $ME$ στις πλευρές $AB$ και $A\varGamma $ αντίστοιχα.

Να αποδείξετε ότι:

\begin{alist}
\item Αν είναι $M\varDelta   = ME$, τότε:
\begin{alist}
\item τα τρίγωνα $B\varDelta   M$ και $\varGamma  EM$ είναι ίσα, \monades{8}
\item το τρίγωνο $AB\varGamma $ είναι ισοσκελές. \monades{9}
\end{alist}
\item Αν είναι $AB = A\varGamma $, τότε $M\varDelta   = ME$. \monades{8}

\begin{center}
\begin{tikzpicture}[scale=1]
\tkzDefPoint(0,0){A}
\tkzDefPoint(0,3){B}
\tkzDefPoint(3,0){C}
\tkzDefMidPoint(B,C) \tkzGetPoint{M}
\tkzDefPointBy[projection=onto A--B](M) \tkzGetPoint{Delta}
\tkzDefPointBy[projection=onto A--C](M) \tkzGetPoint{E}

\draw[l3] (A)--(B)--(C)--cycle;
\draw[l3] (M)--(Delta);
\draw[l3] (M)--(E);

\tkzDrawPoints(A,B,C,M,Delta,E)
\tkzLabelPoint[below left](A){$A$}
\tkzLabelPoint[left](B){$B$}
\tkzLabelPoint[right](C){$\varGamma $}
\tkzLabelPoint[above right](M){$M$}
\tkzLabelPoint[left](Delta){$\varDelta  $}
\tkzLabelPoint[below](E){$E$}

\tkzMarkRightAngle(B,A,C)
\tkzMarkRightAngle(M,Delta,B)
\tkzMarkRightAngle(M,E,C)
\tkzMarkSegments[mark=s|](B,M M,C)
\end{tikzpicture}
\end{center}
\end{alist}

\vspace{4mm}

% --- Άσκηση 36332 (36332.tex) ---
\Askhsh{36332}{2}
\begin{ylh}{1}
    \item 3.1. Είδη και στοιχεία τριγώνων
    \item 3.2. 1ο Κριτήριο ισότητας τριγώνων
    \item 3.6. Κριτήρια ισότητας ορθογώνιων τριγώνων.
\end{ylh}
% ===================================================================
%  Αρχείο: 36332.tex (Fragment)
%  Αυτόματη μετατροπή μέσω doc2latex.py
% ===================================================================



Δίνεται ισοσκελές τρίγωνο $AB\varGamma $ με $AB=A\varGamma $. Στην προέκταση της $B\varGamma $ (προς το Γ) θεωρούμε σημείο $\varDelta  $ και στην προέκταση της $\varGamma B$ (προς το Β) θεωρούμε σημείο $E$ έτσι ώστε $\varGamma \varDelta  $ = $BE$. Από το $\varDelta$ φέρουμε $\varDelta  H$ κάθετη στην ευθεία $A\varGamma $ και από το $E$ φέρουμε $EZ$ κάθετη στην ευθεία $AB$.

Να αποδείξετε ότι:

\begin{alist}
\item $A\varDelta  $ = $AE$ \monades{12}

\item $EZ$ = $\varDelta  H$ \monades{13}


\begin{center}
\begin{tikzpicture}[scale=0.8]
\tkzDefPoint(0,3){A}
\tkzDefPoint(-2,0){B}
\tkzDefPoint(2,0){C}
\tkzDefPoint(3.5,0){Delta} % $CD$ = 1.5
\tkzDefPoint(-3.5,0){E} % $BE$ = 1.5
\tkzDefPointBy[projection=onto A--C](Delta) \tkzGetPoint{H}
\tkzDefPointBy[projection=onto A--B](E) \tkzGetPoint{Z}

\draw[l3] (A)--(B)--(C)--cycle;
\draw[l3] (B)--(E);
\draw[l3] (C)--(Delta);
\draw[l3] (A)--(Delta);
\draw[l3] (A)--(E);
\draw[l3] (Delta)--(H);
\draw[l3] (E)--(Z);

\tkzDrawPoints(A,B,C,Delta,E,H,Z)
\tkzLabelPoint[above](A){$A$}
\tkzLabelPoint[below](B){$B$}
\tkzLabelPoint[below](C){$\varGamma $}
\tkzLabelPoint[below](Delta){$\varDelta  $}
\tkzLabelPoint[below](E){$E$}
\tkzLabelPoint[above right](H){$H$}
\tkzLabelPoint[above left](Z){$Z$}

\tkzMarkSegments[mark=s|](A,B A,C)
\tkzMarkSegments[mark=s||](B,E C,Delta)
\tkzMarkRightAngle(Delta,H,C)
\tkzMarkRightAngle(E,Z,B)
\end{tikzpicture}
\end{center}
\end{alist}

\vspace{4mm}

% --- Άσκηση 36333 (36333.tex) ---
\Askhsh{36333}{2}
\begin{ylh}{1}
    \item 3.1. Είδη και στοιχεία τριγώνων
    \item 3.2. 1ο Κριτήριο ισότητας τριγώνων.
\end{ylh}
% ===================================================================
%  Αρχείο: 36333.tex (Fragment)
%  Αυτόματη μετατροπή μέσω doc2latex.py
% ===================================================================



Δίνεται τρίγωνο $AB\varGamma $ και $E$ το μέσο της διαμέσου του $AM$. Αν είναι $B\varGamma $ = 2 $BE$, τότε να αποδείξετε ότι:

\begin{alist}
\item $Α\overset{\land}{Ε}B = Ε\overset{\land}{Μ}\varGamma $ \monades{12}

\item $AB$ = $E\varGamma $ \monades{13}


\begin{center}
\begin{tikzpicture}[scale=1]
\tkzDefPoint(0,0){B}
\tkzDefPoint(4,0){C}
\tkzDefMidPoint(B,C) \tkzGetPoint{M}
\tkzDefPoint(1.2,1.6){E} % $BE$ = sqrt(1.44 + 2.56) = 2. $BM$ = 2. 
\tkzDefPointBy[symmetry=center E](M) \tkzGetPoint{A}

\draw[l3] (A)--(B)--(C)--cycle;
\draw[l3] (A)--(M);
\draw[l3] (B)--(E)--(C);

\tkzDrawPoints(A,B,C,M,E)
\tkzLabelPoint[above](A){$A$}
\tkzLabelPoint[below](B){$B$}
\tkzLabelPoint[below](C){$\varGamma $}
\tkzLabelPoint[below](M){$M$}
\tkzLabelPoint[above left](E){$E$}

\tkzMarkSegments[mark=s|](A,E E,M)
\tkzMarkSegments[mark=s||](B,M M,C B,E)
\end{tikzpicture}
\end{center}
\end{alist}

\vspace{4mm}

% --- Άσκηση 36336 (36336.tex) ---
\Askhsh{36336}{2}
\begin{ylh}{1}
    \item 3.1. Είδη και στοιχεία τριγώνων
    \item 3.2. 1ο Κριτήριο ισότητας τριγώνων
    \item 3.4. 3ο Κριτήριο ισότητας τριγώνων
    \item 3.11. Ανισοτικές σχέσεις πλευρών και γωνιών
    \item 4.6. Άθροισμα γωνιών τριγώνου.
\end{ylh}
% ===================================================================
%  Αρχείο: 36336.tex (Fragment)
%  Αυτόματη μετατροπή μέσω doc2latex.py
% ===================================================================



Δίνεται τρίγωνο $AB\varGamma $ με $AB = A\varGamma $ και η διάμεσός του $A\varDelta  $ τέτοια ώστε $B\hat{A}\varDelta   = 30^{\circ}$. Θεωρούμε σημείο $E$ στην $A\varGamma $ τέτοιο ώστε $A\varDelta   = AE$.

\begin{alist}
\item Να αποδείξετε ότι το τρίγωνο $AB\varGamma $ είναι ισόπλευρο. \monades{8}
\item Να υπολογίσετε τις γωνίες του τριγώνου $A\varDelta   E$. \monades{9}
\item Να υπολογίσετε τη γωνία $E\hat{\varDelta  }\varGamma $. \monades{8}

\begin{center}
\begin{tikzpicture}[scale=1.2]
\tkzDefPoint(0,0){B}
\tkzDefPoint(4,0){C}
\tkzDefTriangle[equilateral](B,C) \tkzGetPoint{A}
\tkzDefMidPoint(B,C) \tkzGetPoint{Delta}
\tkzDefPointBy[homothety=center A ratio 0.866](C) \tkzGetPoint{E} % $AD$ = 2*sqrt(3) = 3.46. $AC$=4. 3.46/4 = 0.866.

\draw[l3] (A)--(B)--(C)--cycle;
\draw[l3] (A)--(Delta);
\draw[l3] (Delta)--(E);

\tkzDrawPoints(A,B,C,Delta,E)
\tkzLabelPoint[above](A){$A$}
\tkzLabelPoint[below](B){$B$}
\tkzLabelPoint[below](C){$\varGamma $}
\tkzLabelPoint[below](Delta){$\varDelta  $}
\tkzLabelPoint[right](E){$E$}

\tkzMarkRightAngle(A,Delta,B)
\tkzMarkAngle[size=0.6,mark=none](B,A,Delta)
\tkzLabelAngle[pos=0.9](B,A,Delta){$30^\circ$}
\tkzMarkSegments[mark=s|](A,Delta A,E)
\end{tikzpicture}
\end{center}
\end{alist}

\vspace{4mm}

% --- Άσκηση 36337 (36337.tex) ---
\Askhsh{36337}{2}
\begin{ylh}{1}
    \item 5.4. Ρόμβος
    \item 5.5. Τετράγωνο
    \item 5.6. Εφαρμογές στα τρίγωνα
    \item 5.11. Ισοσκελές τραπέζιο.
\end{ylh}
% ===================================================================
%  Αρχείο: 36337.tex (Fragment)
%  Αυτόματη μετατροπή μέσω doc2latex.py
% ===================================================================



Σε ορθογώνιο και ισοσκελές τρίγωνο $AB\varGamma $ ($\hat{A} = 90^{\circ}$) θεωρούμε τα μέσα $\varDelta  , E$ και $Z$ των πλευρών του $AB, A\varGamma $ και $B\varGamma $ αντίστοιχα.

\begin{alist}
\item το τετράπλευρο $AEZ\varDelta  $ είναι τετράγωνο, \monades{15}
\item το τετράπλευρο $E\varDelta   B\varGamma $ είναι ισοσκελές τραπέζιο. \monades{10}

\begin{center}
\begin{tikzpicture}[scale=1]
\tkzDefPoint(0,0){A}
\tkzDefPoint(0,3){B}
\tkzDefPoint(3,0){C}
\tkzDefMidPoint(A,B) \tkzGetPoint{Delta}
\tkzDefMidPoint(A,C) \tkzGetPoint{E}
\tkzDefMidPoint(B,C) \tkzGetPoint{Z}

\draw[l3] (A)--(B)--(C)--cycle;
\draw[l3] (Delta)--(Z)--(E);

\tkzDrawPoints(A,B,C,Delta,E,Z)
\tkzLabelPoint[below left](A){$A$}
\tkzLabelPoint[left](B){$B$}
\tkzLabelPoint[right](C){$\varGamma $}
\tkzLabelPoint[left](Delta){$\varDelta  $}
\tkzLabelPoint[below](E){$E$}
\tkzLabelPoint[above right](Z){$Z$}

\tkzMarkRightAngle(B,A,C)
\tkzMarkSegments[mark=s|](A,Delta Delta,B A,E E,C)
\end{tikzpicture}
\end{center}
\end{alist}

\vspace{4mm}

% --- Άσκηση 36338 (36338.tex) ---
\Askhsh{36338}{2}
\begin{ylh}{1}
    \item 3.1. Είδη και στοιχεία τριγώνων
    \item 3.6. Κριτήρια ισότητας ορθογώνιων τριγώνων
    \item 3.14. Σχετικές θέσεις ευθείας και κύκλου
    \item 3.15. Εφαπτόμενα τμήματα.
\end{ylh}
% ===================================================================
%  Αρχείο: 36338.tex (Fragment)
%  Αυτόματη μετατροπή μέσω doc2latex.py
% ===================================================================



Δίνονται δύο ομόκεντροι κύκλοι με κέντρο Ο και ακτίνες ρ και R (ρ\textless R). Οι χορδές $\varDelta  \varGamma $ και $ZE$ του κύκλου (Ο,R) εφάπτονται του κύκλου (Ο, ρ) στα σημεία $A$ και Β αντίστοιχα.

\begin{alist}
\item Να αποδείξετε ότι $\varDelta  \varGamma $=$ZE$. \monades{12}

\item Αν οι $\varDelta  \varGamma $ και $ZE$ προεκτεινόμενες τέμνονται στο σημείο $K$, να αποδείξετε ότι το τρίγωνο $KE\varGamma $ είναι ισοσκελές. \monades{13}


\begin{center}
\begin{tikzpicture}[scale=1]
\tkzDefPoint(0,0){O}
% Inner circle tangent point A on right side
\tkzDefPoint(1.5,0){A}
% Tangent point B on inner circle, rotated 70 degrees from A
\tkzDefPointBy[rotation=center O angle 70](A) \tkzGetPoint{B}

% Chord $\varDelta  \varGamma $ tangent at A (perpendicular to $OA$ at A)
\tkzDefLine[orthogonal=through A](O,A) \tkzGetPoint{chord1_dir}
% Find intersections of tangent line through A with outer circle (radius 2.5)
\tkzDefPoint(2.5,0){R_point}
\tkzInterLC(A,chord1_dir)(O,R_point) \tkzGetPoints{Gamma}{Delta}

% Chord $ZE$ tangent at B (perpendicular to $OB$ at B)
\tkzDefLine[orthogonal=through B](O,B) \tkzGetPoint{chord2_dir}
% Find intersections of tangent line through B with outer circle
\tkzInterLC(B,chord2_dir)(O,R_point) \tkzGetPoints{E}{Z}

% K is intersection of lines $\varDelta  \varGamma $ and $ZE$ (extended)
\tkzInterLL(Delta,Gamma)(Z,E) \tkzGetPoint{K}

% Draw circles
\tkzDrawCircle(O,A)
\tkzDrawCircle(O,R_point)

% Draw chords
\draw[l3] (Delta)--(Gamma);
\draw[l3] (Z)--(E);
\draw[l3, dashed] (Gamma)--(K)--(E);

% Draw radii to tangent points
\draw[l3, dashed] (O)--(A);
\draw[l3, dashed] (O)--(B);

\tkzDrawPoints(O,A,B,Delta,Gamma,Z,E,K)
\tkzLabelPoint[below](O){$ O$}
\tkzLabelPoint[right](A){$A$}
\tkzLabelPoint[above right](B){$B$}
\tkzLabelPoint[above](Delta){$\varDelta  $}
\tkzLabelPoint[below](Gamma){$\varGamma $}
\tkzLabelPoint[left](Z){$ Z$}
\tkzLabelPoint[above right](E){$ E$}
\tkzLabelPoint[right](K){$ K$}

\tkzMarkRightAngle(O,A,Delta)
\tkzMarkRightAngle(O,B,Z)
\end{tikzpicture}
\end{center}
\end{alist}

\vspace{4mm}

% --- Άσκηση 36339 (36339.tex) ---
\Askhsh{36339}{2}
\begin{ylh}{1}
    \item 3.6. Κριτήρια ισότητας ορθογώνιων τριγώνων
    \item 4.2. Τέμνουσα δύο ευθειών - Ευκλείδειο αίτημα
    \item 5.2. Παραλληλόγραμμα
    \item 5.3. Ορθογώνιο.
\end{ylh}
% ===================================================================
%  Αρχείο: 36339.tex (Fragment)
%  Αυτόματη μετατροπή μέσω doc2latex.py
% ===================================================================



Στο σχήμα που ακολουθεί, το $M$ είναι το μέσο της πλευράς $B\varGamma $ τριγώνου $AB\varGamma $, και τα τμήματα $B\varDelta  $ και $E\varGamma $ είναι κάθετα στη $B\varGamma $ στα σημεία $B$, Γ αντίστοιχα, τέτοια ώστε $M\varDelta  $ = $ME$. Να αποδείξετε ότι:

\begin{alist}
\item τα τμήματα $B\varDelta  $ και $\varGamma E$ είναι ίσα, \monades{13}

\item το τετράπλευρο $B\varDelta  $ΕΓ είναι ορθογώνιο παραλληλόγραμμο. \monades{12}


\begin{center}
\begin{tikzpicture}[scale=1]
\tkzDefPoint(0,0){B}
\tkzDefPoint(4,0){Gamma}
\tkzDefMidPoint(B,Gamma) \tkzGetPoint{M}
\tkzDefPoint(0,2.5){Delta}
\tkzDefPoint(4,2.5){E}
\tkzDefPoint(1.5,3.5){A}

\draw[l3] (A)--(B)--(Gamma)--cycle;
\draw[l3] (B)--(Delta)--(E)--(Gamma);
\draw[l3] (M)--(Delta);
\draw[l3] (M)--(E);

\tkzDrawPoints(A,B,Gamma,Delta,E,M)
\tkzLabelPoint[above](A){$A$}
\tkzLabelPoint[below](B){$B$}
\tkzLabelPoint[below](Gamma){$\varGamma $}
\tkzLabelPoint[left](Delta){$\varDelta  $}
\tkzLabelPoint[right](E){$E$}
\tkzLabelPoint[below](M){$M$}

\tkzMarkRightAngle(Delta,B,M)
\tkzMarkRightAngle(E,Gamma,M)
\tkzMarkSegments[mark=s|](B,M M,Gamma)
\tkzMarkSegments[mark=s||](M,Delta M,E)
\end{tikzpicture}
\end{center}
\end{alist}

\vspace{4mm}

% --- Άσκηση 36340 (36340.tex) ---
\Askhsh{36340}{2}
\begin{ylh}{1}
    \item 3.2. 1ο Κριτήριο ισότητας τριγώνων
    \item 5.11. Ισοσκελές τραπέζιο.
\end{ylh}
% ===================================================================
%  Αρχείο: 36340.tex (Fragment)
%  Αυτόματη μετατροπή μέσω doc2latex.py
% ===================================================================



Δίνεται ισοσκελές τραπέζιο $AB\varGamma \varDelta  $ ($AB$ // $\varDelta  \varGamma $ και $A\varDelta  $ = $B\varGamma $), το μέσο $M$ της πλευράς $\varDelta  \varGamma $ και τα μέσα Κ και Λ των μη παράλληλων πλευρών του $A\varDelta  $ και $B\varGamma $ αντίστοιχα.

Να αποδείξετε ότι:

\begin{alist}
\item τα τμήματα $KM$ και $\varLambda M$ είναι ίσα, \monades{12}

\item Τα τμήματα $AM$ και $BM$ είναι ίσα. \monades{13}


\begin{center}
\begin{tikzpicture}[scale=0.8]
\tkzDefPoint(1,2.5){A}
\tkzDefPoint(4,2.5){B}
\tkzDefPoint(6,0){C}
\tkzDefPoint(-1,0){Delta}
\tkzDefMidPoint(Delta,C) \tkzGetPoint{M}
\tkzDefMidPoint(A,Delta) \tkzGetPoint{K}
\tkzDefMidPoint(B,C) \tkzGetPoint{L}

\draw[l3] (A)--(B)--(C)--(Delta)--cycle;
\draw[l3] (K)--(M)--(L);
\draw[l3] (A)--(M)--(B);

\tkzDrawPoints(A,B,C,Delta,M,K,L)
\tkzLabelPoint[above](A){$A$}
\tkzLabelPoint[above](B){$B$}
\tkzLabelPoint[below right](C){$\varGamma $}
\tkzLabelPoint[below left](Delta){$\varDelta  $}
\tkzLabelPoint[below](M){$M$}
\tkzLabelPoint[left](K){$K$}
\tkzLabelPoint[right](L){$\Lambda$}

\tkzMarkSegments[mark=s|](Delta,M M,C)
\tkzMarkSegments[mark=s||](A,Delta B,C)
\end{tikzpicture}
\end{center}
\end{alist}

\vspace{4mm}

% --- Άσκηση 36341 (36341.tex) ---
\Askhsh{36341}{2}
\begin{ylh}{1}
    \item 3.1. Είδη και στοιχεία τριγώνων
    \item 3.2. 1ο Κριτήριο ισότητας τριγώνων
    \item 3.4. 3ο Κριτήριο ισότητας τριγώνων.
\end{ylh}
% ===================================================================
%  Αρχείο: 36341.tex (Fragment)
%  Αυτόματη μετατροπή μέσω doc2latex.py
% ===================================================================



Δίνεται γωνία xΑy και η διχοτόμος της Αδ. Από τυχαίο σημείο $B$ της Αx φέρνουμε κάθετη στη διχοτόμο Αδ, η οποία τέμνει την Αδ στο σημείο $\varDelta  $ και την Ay στο σημείο $\varGamma $.

\begin{alist}
\item Να αποδείξετε ότι τα τμήματα $AB$ και $A\varGamma $ είναι ίσα. (Μονάδες12)

\item Αν $E$ τυχαίο σημείο της Αδ, να αποδείξετε ότι το $E$ ισαπέχει από τα Β και Γ.

\monades{13}


\begin{center}
\begin{tikzpicture}[scale=1]
\tkzDefPoint(0,0){A}
\tkzDefPoint(5,0){x_axis}
\tkzDefPointBy[rotation=center A angle 40](x_axis) \tkzGetPoint{y_axis}
\tkzDefLine[bisector](x_axis,A,y_axis) \tkzGetPoint{Ad_dir}
\tkzDefPoint(3,0){B}
\tkzDefLine[orthogonal=through B](A,Ad_dir) \tkzGetPoint{perp_dir}
\tkzInterLL(B,perp_dir)(A,Ad_dir) \tkzGetPoint{Delta}
\tkzInterLL(B,perp_dir)(A,y_axis) \tkzGetPoint{Gamma}

\draw[l3, ->] (A)--(x_axis) node[right]{$x$};
\draw[l3, ->] (A)--(y_axis) node[above]{$y$};
\draw[l3, ->] (A)--(Ad_dir) node[right]{$\delta$};
\draw[l3] (B)--(Gamma);

\tkzDrawPoints(A,B,Gamma,Delta)
\tkzLabelPoint[below left](A){$A$}
\tkzLabelPoint[below](B){$B$}
\tkzLabelPoint[left](Gamma){$\varGamma $}
\tkzLabelPoint[above right](Delta){$\varDelta  $}

\tkzMarkAngle[size=0.6,mark=none](B,A,Delta)
\tkzMarkAngle[size=0.6,mark=none](Delta,A,Gamma)
\tkzMarkRightAngle(A,Delta,B)
\end{tikzpicture}
\end{center}
\end{alist}

\vspace{4mm}

% --- Άσκηση 36342 (36342.tex) ---
\Askhsh{36342}{2}
\begin{ylh}{1}
    \item 5.6. Εφαρμογές στα τρίγωνα.
\end{ylh}
% ===================================================================
%  Αρχείο: 36342.tex (Fragment)
%  Αυτόματη μετατροπή μέσω doc2latex.py
% ===================================================================



Έστω ορθογώνιο τρίγωνο $AB\varGamma $ με $\hat{A} = 90^{\circ}$ και $\hat{B} = 30^{\circ}$. Αν τα σημεία $E$ και $\varDelta  $ είναι τα μέσα των $AB$ και $B\varGamma $ αντίστοιχα με $E\varDelta  =1$, να υπολογίσετε τα τμήματα:

\begin{alist}
\item $A\varGamma $ \monades{8}
\item $B\varGamma $ \monades{9}
\item $A\varDelta  $ \monades{8}

\begin{center}
\begin{tikzpicture}[scale=1.5]
\tkzDefPoint(0,0){A}
\tkzDefPoint(0,1.732){B}
\tkzDefPoint(1,0){C}
\tkzDefMidPoint(A,B) \tkzGetPoint{E}
\tkzDefMidPoint(B,C) \tkzGetPoint{Delta}

\draw[l3] (A)--(B)--(C)--cycle;
\draw[l3] (E)--(Delta);
\draw[l3] (A)--(Delta);

\tkzDrawPoints(A,B,C,Delta,E)
\tkzLabelPoint[below left](A){$A$}
\tkzLabelPoint[above](B){$B$}
\tkzLabelPoint[right](C){$\varGamma $}
\tkzLabelPoint[left](E){$E$}
\tkzLabelPoint[above right](Delta){$\varDelta  $}

\tkzMarkRightAngle(B,A,C)
\tkzMarkAngle[size=0.4,mark=none](A,B,C)
\tkzLabelAngle[pos=0.6](A,B,C){$30^\circ$}
\tkzMarkSegments[mark=s|](A,E E,B)
\tkzMarkSegments[mark=s||](B,Delta Delta,C)
\end{tikzpicture}
\end{center}
\end{alist}

\vspace{4mm}

% --- Άσκηση 36343 (36343.tex) ---
\Askhsh{36343}{2}
\begin{ylh}{1}
    \item 3.6. Κριτήρια ισότητας ορθογώνιων τριγώνων
    \item 5.6. Εφαρμογές στα τρίγωνα.
\end{ylh}
% ===================================================================
%  Αρχείο: 36343.tex (Fragment)
%  Αυτόματη μετατροπή μέσω doc2latex.py
% ===================================================================



Έστω ισοσκελές τρίγωνο $AB\varGamma $ με $AB=A\varGamma $. Από τα μέσα Κ και Λ των πλευρών $A\varGamma $ και $AB$ αντίστοιχα, φέρουμε τα κάθετα τμήματα $KE$ και $\varLambda Z$ στην πλευρά $B\varGamma $.

Να αποδείξετε ότι:

\begin{alist}
\item τα τρίγωνα $KE\varGamma $ και $\varLambda ZB$ είναι ίσα. \monades{15}

\item $EH$=$Z\varTheta $, όπου Η , Θ τα μέσα των τμημάτων $K\varGamma $, $\varLambda B$ αντίστοιχα. \monades{10}


\begin{center}
\begin{tikzpicture}[scale=1]
\tkzDefPoint(0,4){A}
\tkzDefPoint(-2,0){B}
\tkzDefPoint(2,0){C}
\tkzDefMidPoint(A,C) \tkzGetPoint{K}
\tkzDefMidPoint(A,B) \tkzGetPoint{L}
\tkzDefPointBy[projection=onto B--C](K) \tkzGetPoint{E}
\tkzDefPointBy[projection=onto B--C](L) \tkzGetPoint{Z}
\tkzDefMidPoint(K,C) \tkzGetPoint{H}
\tkzDefMidPoint(L,B) \tkzGetPoint{Theta}

\draw[l3] (A)--(B)--(C)--cycle;
\draw[l3] (K)--(E);
\draw[l3] (L)--(Z);

\tkzDrawPoints(A,B,C,K,L,E,Z,H,Theta)
\tkzLabelPoint[above](A){$A$}
\tkzLabelPoint[below](B){$B$}
\tkzLabelPoint[below](C){$\varGamma $}
\tkzLabelPoint[above right](K){$K$}
\tkzLabelPoint[above left](L){$\Lambda$}
\tkzLabelPoint[below](E){$E$}
\tkzLabelPoint[below](Z){$Z$}
\tkzLabelPoint[right](H){$H$}
\tkzLabelPoint[left](Theta){$\varTheta $}

\tkzMarkSegments[mark=s|](A,L L,B A,K K,C)
\tkzMarkRightAngle(K,E,C)
\tkzMarkRightAngle(L,Z,B)
\end{tikzpicture}
\end{center}
\end{alist}

\vspace{4mm}

% --- Άσκηση 36344 (36344.tex) ---
\Askhsh{36344}{2}
\begin{ylh}{1}
    \item 3.1. Είδη και στοιχεία τριγώνων
    \item 3.4. 3ο Κριτήριο ισότητας τριγώνων
    \item 3.6. Κριτήρια ισότητας ορθογώνιων τριγώνων
    \item 3.14. Σχετικές θέσεις ευθείας και κύκλου.
\end{ylh}
% ===================================================================
%  Αρχείο: 36344.tex (Fragment)
%  Αυτόματη μετατροπή μέσω doc2latex.py
% ===================================================================



Έστω κύκλος με κέντρο Ο και ακτίνα ρ. Σε σημείο $N$ του κύκλου φέρουμε την εφαπτόμενή του, και εκατέρωθεν του $N$ θεωρούμε σημεία $A$ και Β, τέτοια ώστε $NA$=$NB$. Οι $OA$ και $OB$ τέμνουν τον κύκλο στα Κ και Λ αντίστοιχα.

Να αποδείξετε ότι:

\begin{alist}
\item το τρίγωνο $AOB$ είναι ισοσκελές, \monades{13}

\item το σημείο $N$ είναι μέσο του τόξου $K\varLambda $. \monades{12}


\begin{center}
\begin{tikzpicture}[scale=1]
\tkzDefPoint(0,0){O}
\tkzDefPoint(0,2.5){N}
\tkzDrawCircle(O,N)
\tkzDefLine[tangent at=N](O) \tkzGetPoint{tan_dir}
\tkzDefPoint(-2,2.5){A}
\tkzDefPoint(2,2.5){B}

\draw[l3] (A)--(B);
\draw[l3] (O)--(A);
\draw[l3] (O)--(B);
\tkzInterLC(O,A)(O,N) \tkzGetPoints{K_dummy}{K}
\tkzInterLC(O,B)(O,N) \tkzGetPoints{L}{L_dummy}

\tkzDrawPoints(O,N,A,B,K,L)
\tkzLabelPoint[below](O){$O$}
\tkzLabelPoint[above](N){$N$}
\tkzLabelPoint[left](A){$A$}
\tkzLabelPoint[right](B){$B$}
\tkzLabelPoint[above left](K){$K$}
\tkzLabelPoint[above right](L){$\Lambda$}

\tkzMarkRightAngle(O,N,B)
\tkzMarkSegments[mark=s|](N,A N,B)
\end{tikzpicture}
\end{center}
\end{alist}

\vspace{4mm}

% --- Άσκηση 36345 (36345.tex) ---
\Askhsh{36345}{2}
\begin{ylh}{1}
    \item 3.1. Είδη και στοιχεία τριγώνων
    \item 3.6. Κριτήρια ισότητας ορθογώνιων τριγώνων.
\end{ylh}
% ===================================================================
%  Αρχείο: 36345.tex (Fragment)
%  Αυτόματη μετατροπή μέσω doc2latex.py
% ===================================================================



Έστω κύκλος με κέντρο $O$ και ακτίνα $\rho$. Θεωρούμε διάμετρο $AB$ και τυχαίο σημείο $\varGamma $ του κύκλου. Αν τα $AE, \varGamma \varDelta  $ είναι κάθετα τμήματα στις $O\varGamma , OA$ αντίστοιχα και $Z$ το σημείο τομής τους, να αποδείξετε ότι:

\begin{alist}
\item το τρίγωνο $\varDelta   OE$ είναι ισοσκελές, \monades{13}
\item η $OZ$ διχοτομεί τη γωνία $A\hat{O}\varGamma $ και προεκτεινόμενη διέρχεται από το μέσο του τόξου $A\varGamma $. \monades{12}

\begin{center}
\begin{tikzpicture}[scale=1.5]
\tkzDefPoint(0,0){O}
\tkzDefPoint(-2,0){A}
\tkzDefPoint(2,0){B}
\tkzDrawCircle(O,A)
\tkzDefPointBy[rotation=center O angle 60](B) \tkzGetPoint{Gamma}
\tkzDefPointBy[projection=onto O--Gamma](A) \tkzGetPoint{E}
\tkzDefPointBy[projection=onto O--A](Gamma) \tkzGetPoint{Delta}
\tkzInterLL(A,E)(Gamma,Delta) \tkzGetPoint{Z}

\draw[l3] (A)--(B);
\draw[l3] (O)--(Gamma);
\draw[l3] (A)--(E);
\draw[l3] (Gamma)--(Delta);
\draw[l3] (O)--(Z);

\tkzDrawPoints(O,A,B,Gamma,Delta,E,Z)
\tkzLabelPoint[below](O){$O$}
\tkzLabelPoint[below left](A){$A$}
\tkzLabelPoint[below right](B){$B$}
\tkzLabelPoint[above right](Gamma){$\varGamma $}
\tkzLabelPoint[below](Delta){$\varDelta  $}
\tkzLabelPoint[above left](E){$E$}
\tkzLabelPoint[below left](Z){$Z$}

\tkzMarkRightAngle(O,Delta,Gamma)
\tkzMarkRightAngle(O,E,A)
\end{tikzpicture}
\end{center}
\end{alist}

\vspace{4mm}

% --- Άσκηση 36348 (36348.tex) ---
\Askhsh{36348}{2}
\begin{ylh}{1}
    \item 3.1. Είδη και στοιχεία τριγώνων
    \item 3.14. Σχετικές θέσεις ευθείας και κύκλου
    \item 4.2. Τέμνουσα δύο ευθειών - Ευκλείδειο αίτημα
    \item 5.2. Παραλληλόγραμμα.
\end{ylh}
% ===================================================================
%  Αρχείο: 36348.tex (Fragment)
%  Αυτόματη μετατροπή μέσω doc2latex.py
% ===================================================================



Έστω κύκλος με κέντρο Ο και ακτίνα ρ. Θεωρούμε κάθετες ακτίνες $OA$ , $O\varGamma $ και εφαπτόμενο στον κύκλο τμήμα $AB$ τέτοιο ώστε $AB$ = $O\varGamma $.

\begin{alist}
\item Να αποδείξετε ότι τα τμήματα $AO$ και $B\varGamma $ διχοτομούνται. \monades{10}

\item Να υπολογίσετε τις γωνίες του τετραπλεύρου $ABO\varGamma $. \monades{15}


\begin{center}
\begin{tikzpicture}[scale=1.5]
\tkzDefPoint(0,0){O}
\tkzDefPoint(2,0){A}
\tkzDefPoint(0,2){Gamma}
\tkzDrawArc(O,A)(Gamma)
\tkzDefPoint(2,2){B}

\draw[l3] (O)--(A)--(B)--(Gamma)--cycle;
\draw[l3] (A)--(Gamma);
\draw[l3] (O)--(B);

\tkzDrawPoints(O,A,B,Gamma)
\tkzLabelPoint[below left](O){$O$}
\tkzLabelPoint[below right](A){$A$}
\tkzLabelPoint[above right](B){$B$}
\tkzLabelPoint[above left](Gamma){$\varGamma $}

\tkzMarkRightAngle(A,O,Gamma)
\tkzMarkRightAngle(O,A,B)
\tkzMarkSegments[mark=s|](O,A A,B B,Gamma Gamma,O)
\end{tikzpicture}
\end{center}
\end{alist}

\vspace{4mm}

% --- Άσκηση 36349 (36349.tex) ---
\Askhsh{36349}{2}
\begin{ylh}{1}
    \item 3.1. Είδη και στοιχεία τριγώνων
    \item 3.6. Κριτήρια ισότητας ορθογώνιων τριγώνων
    \item 4.6. Άθροισμα γωνιών τριγώνου
    \item 5.4. Ρόμβος.
\end{ylh}
% ===================================================================
%  Αρχείο: 36349.tex (Fragment)
%  Αυτόματη μετατροπή μέσω doc2latex.py
% ===================================================================



Έστω κύκλος με κέντρο Ο και ακτίνα ρ. Θεωρούμε ακτίνα $OA$ και χορδή $B\varGamma $ κάθετη στο μέσο της Μ.

\begin{alist}
\item Να αποδείξετε ότι το τετράπλευρο $A\varGamma OB$ είναι ρόμβος. \monades{10}

\item Να υπολογίσετε τις γωνίες του τετραπλεύρου $A\varGamma OB$. \monades{15}


\begin{center}
\begin{tikzpicture}[scale=1.5]
\tkzDefPoint(0,0){O}
\tkzDefPoint(2,0){A}
\tkzDrawCircle(O,A)
\tkzDefMidPoint(O,A) \tkzGetPoint{M}
\tkzDefLine[orthogonal=through M](O,A) \tkzGetPoint{chord_dir}
\tkzInterLC(M,chord_dir)(O,A) \tkzGetPoints{Gamma}{B}

\draw[l3] (O)--(Gamma)--(A)--(B)--cycle;
\draw[l3] (O)--(A);
\draw[l3] (B)--(Gamma);

\tkzDrawPoints(O,A,B,Gamma,M)
\tkzLabelPoint[left](O){$O$}
\tkzLabelPoint[right](A){$A$}
\tkzLabelPoint[below](B){$B$}
\tkzLabelPoint[above](Gamma){$\varGamma $}
\tkzLabelPoint[below right](M){$M$}

\tkzMarkRightAngle(O,M,Gamma)
\tkzMarkSegments[mark=s|](O,M M,A)
\end{tikzpicture}
\end{center}
\end{alist}

\vspace{4mm}

% --- Άσκηση 36350 (36350.tex) ---
\Askhsh{36350}{2}
\begin{ylh}{1}
    \item 3.1. Είδη και στοιχεία τριγώνων
    \item 3.4. 3ο Κριτήριο ισότητας τριγώνων
    \item 3.6. Κριτήρια ισότητας ορθογώνιων τριγώνων
    \item 5.9. Μια ιδιότητα του ορθογώνιου τριγώνου.
\end{ylh}
% ===================================================================
%  Αρχείο: 36350.tex (Fragment)
%  Αυτόματη μετατροπή μέσω doc2latex.py
% ===================================================================



Έστω ισοσκελές τρίγωνο $ABC$ ($AB = AC$). Στις προεκτάσεις των πλευρών $AB$ και $AC$ προς το $A$ φέρουμε τμήματα $B\varDelta  $ και $\varGamma  E$ κάθετα στις $A\varGamma $ και $AB$ αντίστοιχα.

\begin{alist}
\item Να αποδείξετε ότι $B\varDelta   = \varGamma  E$. \monades{10}
\item Αν το σημείο $M$ είναι το μέσο της $B\varGamma $, τότε να αποδείξετε ότι:
\begin{alist}
\item $M\varDelta   = ME$, \monades{8}
\item η $MA$ διχοτομεί τη γωνία $\hat{\varDelta ME}$. \monades{7}
\end{alist}

\begin{center}
\begin{tikzpicture}[scale=1]
\tkzDefPoint(0,4){A}
\tkzDefPoint(-2,0){B}
\tkzDefPoint(2,0){C}
\tkzDefLine[orthogonal=through B](A,C) \tkzGetPoint{BD_dir}
\tkzInterLL(B,BD_dir)(A,C) \tkzGetPoint{Delta}
\tkzDefLine[orthogonal=through C](A,B) \tkzGetPoint{CE_dir}
\tkzInterLL(C,CE_dir)(A,B) \tkzGetPoint{E}
\tkzDefMidPoint(B,C) \tkzGetPoint{M}

\draw[l3] (A)--(B)--(C)--cycle;
\draw[l3] (B)--(Delta);
\draw[l3] (C)--(E);
\draw[l3] (M)--(Delta);
\draw[l3] (M)--(E);
\draw[l3] (M)--(A);

\tkzDrawPoints(A,B,C,Delta,E,M)
\tkzLabelPoint[above](A){$A$}
\tkzLabelPoint[left](B){$B$}
\tkzLabelPoint[right](C){$\varGamma $}
\tkzLabelPoint[right](Delta){$\varDelta  $}
\tkzLabelPoint[left](E){$E$}
\tkzLabelPoint[below](M){$M$}

\tkzMarkRightAngle(B,Delta,C)
\tkzMarkRightAngle(C,E,B)
\tkzMarkSegments[mark=s|](A,B A,C)
\tkzMarkSegments[mark=s||](B,M M,C)
\end{tikzpicture}
\end{center}
\end{alist}

\vspace{4mm}

% --- Άσκηση 36351 (36351.tex) ---
\Askhsh{36351}{2}
\begin{ylh}{1}
    \item 3.1. Είδη και στοιχεία τριγώνων
    \item 3.2. 1ο Κριτήριο ισότητας τριγώνων
    \item 3.6. Κριτήρια ισότητας ορθογώνιων τριγώνων
    \item 5.4. Ρόμβος.
\end{ylh}
% ===================================================================
%  Αρχείο: 36351.tex (Fragment)
%  Αυτόματη μετατροπή μέσω doc2latex.py
% ===================================================================



Δίνεται ρόμβος Α$B\varDelta  $Γ. Στην προέκταση της διαγωνίου του $A\varDelta  $ (προς το Δ) παίρνουμε τυχαίο σημείο $E$.

Να αποδείξετε ότι:

\begin{alist}
\item Το σημείο $E$ ισαπέχει από τις προεκτάσεις των πλευρών $AB$ και $A\varGamma $ (προς το μέρος των Β και $\varGamma$ αντίστοιχα). \monades{10}

\item Το σημείο $E$ ισαπέχει από τα σημεία $B$ και Γ. \monades{15}


\begin{center}
\begin{tikzpicture}[scale=0.8]
\tkzDefPoint(0,0){A}
\tkzDefPoint(4,1){B}
\tkzDefPoint(8,0){Delta}
\tkzDefPoint(4,-1){Gamma}
\tkzDefPoint(11,0){E}

\draw[l3] (A)--(B)--(Delta)--(Gamma)--cycle;
\draw[l3] (A)--(E);
\draw[l3] (B)--(E)--(Gamma);

\tkzDrawPoints(A,B,Delta,Gamma,E)
\tkzLabelPoint[left](A){$A$}
\tkzLabelPoint[above](B){$B$}
\tkzLabelPoint[below](Delta){$\varDelta  $}
\tkzLabelPoint[below](Gamma){$\varGamma $}
\tkzLabelPoint[right](E){$E$}

\tkzMarkSegments[mark=s|](A,B B,Delta Delta,Gamma Gamma,A)
\end{tikzpicture}
\end{center}
\end{alist}

\vspace{4mm}

% --- Άσκηση 36353 (36353.tex) ---
\Askhsh{36353}{2}
\begin{ylh}{1}
    \item 3.2. 1ο Κριτήριο ισότητας τριγώνων
    \item 5.3. Ορθογώνιο.
\end{ylh}
% ===================================================================
%  Αρχείο: 36353.tex (Fragment)
%  Αυτόματη μετατροπή μέσω doc2latex.py
% ===================================================================



Σε κύκλο κέντρου Ο και ακτίνας ρ φέρουμε δυο διαμέτρους του $AB$ και $\varGamma \varDelta  $.

Να αποδείξετε ότι:

\begin{alist}
\item οι χορδές $A\varGamma $ και $B\varDelta  $ του κύκλου είναι ίσες, \monades{13}

\item το τετράπλευρο με κορυφές τα σημεία $A$, Γ, Β και $\varDelta$ είναι ορθογώνιο. \monades{12}
\end{alist}

\vspace{4mm}

% --- Άσκηση 36354 (36354.tex) ---
\Askhsh{36354}{2}
\begin{ylh}{1}
    \item 3.4. 3ο Κριτήριο ισότητας τριγώνων
    \item 3.6. Κριτήρια ισότητας ορθογώνιων τριγώνων
    \item 3.14. Σχετικές θέσεις ευθείας και κύκλου
    \item 3.15. Εφαπτόμενα τμήματα.
\end{ylh}
% ===================================================================
%  Αρχείο: 36354.tex (Fragment)
%  Αυτόματη μετατροπή μέσω doc2latex.py
% ===================================================================



Έστω κύκλος (Ο, ρ) και ένα εξωτερικό του σημείο $A$. Από το Α φέρουμε τα εφαπτόμενα τμήματα $AB$ και $A\varGamma $ του κύκλου και έστω $E$ και $\varDelta$ τα αντιδιαμετρικά σημεία των Β και $\varGamma$ αντίστοιχα.

Να αποδείξετε ότι:

\begin{alist}
\item τα τρίγωνα Α$BE$ και $A\varGamma \varDelta  $ είναι ίσα, \monades{13}

\item τα τρίγωνα Α$B\varDelta  $ και $A\varGamma E$ είναι ίσα. \monades{12}


\begin{center}
\begin{tikzpicture}[scale=0.8]
\tkzDefPoint(0,0){O}
\tkzDefPoint(2,0){temp} % radius 2
\tkzDefPoint(5,0){A}
\tkzDefMidPoint(O,A) \tkzGetPoint{M}
\tkzInterCC(O,temp)(M,O) \tkzGetPoints{Gamma}{B}
\tkzDefPointBy[symmetry=center O](B) \tkzGetPoint{E}
\tkzDefPointBy[symmetry=center O](Gamma) \tkzGetPoint{Delta}

\draw[l3] (A)--(B);
\draw[l3] (A)--(Gamma);
\draw[l3] (B)--(E);
\draw[l3] (Gamma)--(Delta);
\draw[l3] (A)--(E);
\draw[l3] (A)--(Delta);
\draw[l3] (O)--(B);
\draw[l3] (O)--(Gamma);

\tkzDrawPoints(O,A,B,Gamma,Delta,E)
\tkzLabelPoint[below](O){$O$}
\tkzLabelPoint[right](A){$A$}
\tkzLabelPoint[above left](B){$B$}
\tkzLabelPoint[below left](Gamma){$\varGamma $}
\tkzLabelPoint[above right](Delta){$\varDelta  $}
\tkzLabelPoint[below right](E){$E$}

\tkzMarkRightAngle(O,B,A)
\tkzMarkRightAngle(O,Gamma,A)
\end{tikzpicture}
\end{center}
\end{alist}

\vspace{4mm}

% --- Άσκηση 36355 (36355.tex) ---
\Askhsh{36355}{2}
\begin{ylh}{1}
    \item 4.6. Άθροισμα γωνιών τριγώνου
    \item 5.9. Μια ιδιότητα του ορθογώνιου τριγώνου.
\end{ylh}
% ===================================================================
%  Αρχείο: 36355.tex (Fragment)
%  Αυτόματη μετατροπή μέσω doc2latex.py
% ===================================================================



Έστω ορθογώνιο τρίγωνο $AB\varGamma $ με $\hat{B} = 90\degree$ και $Z$ το μέσο του $A\varGamma $. Με υποτείνουσα το $A\varGamma $ κατασκευάζουμε ορθογώνιο ισοσκελές τρίγωνο $A\varDelta  \varGamma $ με $\hat{\varDelta  } = 90^{Ο}$ και $\varDelta  A$ = $\varDelta  \varGamma $.

\begin{alist}
\item Να αποδείξετε ότι $ΒΖ = \varDelta   Ζ$. \monades{13}

\item Αν είναι $A\hat{\varGamma }B = 30\degree$, τότε να υπολογίσετε τις γωνίες Β$\hat{A}$Δ και Β$\hat{\varGamma }$Δ. \monades{12}


\begin{center}
\begin{tikzpicture}[scale=1]
\tkzDefPoint(0,0){B}
\tkzDefPoint(0,3){A}
\tkzDefPoint(4,0){Gamma}
\tkzDefMidPoint(A,Gamma) \tkzGetPoint{Z}
\tkzDefPointBy[rotation=center Z angle 90](A) \tkzGetPoint{Delta}

\draw[l3] (A)--(B)--(Gamma)--cycle;
\draw[l3] (A)--(Delta)--(Gamma);
\draw[l3] (B)--(Z);
\draw[l3] (Delta)--(Z);

\tkzDrawPoints(A,B,Gamma,Delta,Z)
\tkzLabelPoint[left](A){$A$}
\tkzLabelPoint[below left](B){$B$}
\tkzLabelPoint[right](Gamma){$\varGamma $}
\tkzLabelPoint[above right](Delta){$\varDelta  $}
\tkzLabelPoint[above](Z){$Z$}

\tkzMarkRightAngle(A,B,Gamma)
\tkzMarkRightAngle(A,Delta,Gamma)
\tkzMarkSegments[mark=s|](A,Z Z,Gamma B,Z Delta,Z)
\end{tikzpicture}
\end{center}
\end{alist}

\vspace{4mm}

% --- Άσκηση 36356 (36356.tex) ---
\Askhsh{36356}{2}
\begin{ylh}{1}
    \item 3.1. Είδη και στοιχεία τριγώνων
    \item 3.6. Κριτήρια ισότητας ορθογώνιων τριγώνων
    \item 5.6. Εφαρμογές στα τρίγωνα.
\end{ylh}
% ===================================================================
%  Αρχείο: 36356.tex (Fragment)
%  Αυτόματη μετατροπή μέσω doc2latex.py
% ===================================================================



Θεωρούμε τρίγωνα $AB\varDelta  $ με $AB = B\varDelta   = 5$ και $A\varGamma  E$ με $A\varGamma  = \varGamma  E = 5$ έτσι ώστε τα σημεία $\varDelta  , B, \varGamma $ και $E$ να είναι συνευθειακά. Θεωρούμε τα ύψη τους $BK$ και $\varGamma  \Lambda$ αντίστοιχα.

\begin{alist}
\item Να αποδείξετε ότι:
\begin{alist}
\item τα τρίγωνα $AB\varDelta  $ και $A\varGamma  E$ είναι ισοσκελή, \monades{8}
\item τα σημεία $K$ και $\Lambda$ είναι τα μέσα των τμημάτων $A\varDelta  $ και $AE$ αντίστοιχα. \monades{8}
\end{alist}
\item Αν η περίμετρος του τριγώνου $AB\varGamma $ είναι 12, να υπολογίσετε το τμήμα $K\Lambda$. \monades{9}

\begin{center}
\begin{tikzpicture}[scale=0.8]
\tkzDefPoint(0,0){Delta}
\tkzDefPoint(5,0){B}
\tkzDefPoint(9,0){Gamma}
\tkzDefPoint(14,0){E}
\tkzInterCC[R](B,5)(Gamma,5) \tkzGetPoints{A}{A_dummy}
\tkzDefPointBy[projection=onto A--Delta](B) \tkzGetPoint{K}
\tkzDefPointBy[projection=onto A--E](Gamma) \tkzGetPoint{Lambda}

\draw[l3] (Delta)--(E);
\draw[l3] (A)--(Delta);
\draw[l3] (A)--(B);
\draw[l3] (A)--(Gamma);
\draw[l3] (A)--(E);
\draw[l3] (B)--(K);
\draw[l3] (Gamma)--(Lambda);

\tkzDrawPoints(A,B,Gamma,Delta,E,K,Lambda)
\tkzLabelPoint[above](A){$A$}
\tkzLabelPoint[below](B){$B$}
\tkzLabelPoint[below](Gamma){$\varGamma $}
\tkzLabelPoint[below](Delta){$\varDelta  $}
\tkzLabelPoint[below](E){$E$}
\tkzLabelPoint[left](K){$K$}
\tkzLabelPoint[right](Lambda){$\Lambda$}

\tkzMarkSegments[mark=s|](Delta,B A,B)
\tkzMarkSegments[mark=s||](E,Gamma A,Gamma)
\tkzMarkRightAngle(B,K,Delta)
\tkzMarkRightAngle(Gamma,Lambda,E)
\end{tikzpicture}
\end{center}
\end{alist}

\vspace{4mm}

% --- Άσκηση 37006 (37006.tex) ---
\Askhsh{37006}{2}
\begin{ylh}{1}
    \item 3.1. Είδη και στοιχεία τριγώνων
    \item 3.2. 1ο Κριτήριο ισότητας τριγώνων
    \item 4.6. Άθροισμα γωνιών τριγώνου
    \item 5.9. Μια ιδιότητα του ορθογώνιου τριγώνου.
\end{ylh}
% ===================================================================
%  Αρχείο: 37006.tex (Fragment)
%  Αυτόματη μετατροπή μέσω doc2latex.py
% ===================================================================



Σε ορθογώνιο τρίγωνο $ΑΒ\varGamma $ με $\hat{A} = 90^{\circ}$ και $\hat{B} > \hat{\varGamma }$ φέρουμε το ύψος $A\varDelta  $ και την διάμεσο $ΑΜ$ στην πλευρά $B\varGamma $.

Να αποδείξετε ότι:

\begin{alist}
\item οι γωνίες $\hat{B}$ και $\varGamma \hat{A}\varDelta  $ είναι ίσες, \monades{12}
\item $A\hat{M}\varDelta   = 2 \cdot \hat{\varGamma }$. \monades{13}

\begin{center}
\begin{tikzpicture}[scale=1.2]
\tkzDefPoint(0,0){A}
\tkzDefPoint(4,0){C}
\tkzDefPoint(0,2){B}
\tkzDefMidPoint(B,C) \tkzGetPoint{M}
\tkzDefPointBy[projection=onto B--C](A) \tkzGetPoint{Delta}

\draw[l3] (A)--(B)--(C)--cycle;
\draw[l3] (A)--(Delta);
\draw[l3] (A)--(M);

\tkzDrawPoints(A,B,C,Delta,M)
\tkzLabelPoint[below left](A){$A$}
\tkzLabelPoint[above left](B){$B$}
\tkzLabelPoint[below right](C){$\varGamma $}
\tkzLabelPoint[above right](Delta){$\varDelta  $}
\tkzLabelPoint[above right](M){$M$}

\tkzMarkRightAngle(B,A,C)
\tkzMarkRightAngle(A,Delta,C)
\tkzMarkSegments[mark=s|](B,M M,C A,M)
\end{tikzpicture}
\end{center}
\end{alist}

\vspace{4mm}

% --- Άσκηση 37007 (37007.tex) ---
\Askhsh{37007}{2}
\begin{ylh}{1}
    \item 3.6. Κριτήρια ισότητας ορθογώνιων τριγώνων
    \item 4.6. Άθροισμα γωνιών τριγώνου
    \item 5.2. Παραλληλόγραμμα
    \item 5.3. Ορθογώνιο
    \item 5.9. Μια ιδιότητα του ορθογώνιου τριγώνου.
\end{ylh}
% ===================================================================
%  Αρχείο: 37007.tex (Fragment)
%  Αυτόματη μετατροπή μέσω doc2latex.py
% ===================================================================



Δίνεται παραλληλόγραμμο $ΑΒ\varGamma \varDelta  $ με $\hat{B} = 60^{\circ}$. Φέρουμε τα ύψη $ΑΕ$ και $BZ$ του παραλληλογράμμου $ΑΒ\varGamma \varDelta  $ που αντιστοιχούν στην ευθεία $\varDelta  \varGamma $.

Να αποδείξετε ότι:

\begin{alist}
\item $\varGamma  Z = \frac{A\varDelta  }{2}$, \monades{8}
\item το τρίγωνο $A\varDelta   E$ είναι ίσο με το τρίγωνο $B\varGamma  Z$, \monades{9}
\item το τετράπλευρο $ABZE$ είναι ορθογώνιο. \monades{8}

\begin{center}
\begin{tikzpicture}[scale=1]
\tkzDefPoint(0,0){D}
\tkzDefPoint(4,0){C}
\tkzDefPoint(1.5,2.5){A}
\tkzDefPoint(5.5,2.5){B}

\tkzDefPointBy[projection=onto D--C](A) \tkzGetPoint{E}
\tkzDefPointBy[projection=onto D--C](B) \tkzGetPoint{Z}

\draw[l3] (A)--(B)--(C)--(D)--cycle;
\draw[l3] (A)--(E);
\draw[l3] (B)--(Z);
\draw[l3, dashed] (C)--(Z);

\tkzDrawPoints(A,B,C,D,E,Z)
\tkzLabelPoint[above](A){$A$}
\tkzLabelPoint[above](B){$B$}
\tkzLabelPoint[below](C){$\varGamma $}
\tkzLabelPoint[below](D){$\varDelta  $}
\tkzLabelPoint[below](E){$E$}
\tkzLabelPoint[below](Z){$Z$}

\tkzMarkRightAngle(A,E,C)
\tkzMarkRightAngle(B,Z,C)
\tkzMarkAngle[size=0.5](A,B,C)
\tkzLabelAngle[pos=0.8](A,B,C){$60\degree$}
\end{tikzpicture}
\end{center}
\end{alist}

\vspace{4mm}

% --- Άσκηση 37008 (37008.tex) ---
\Askhsh{37008}{2}
\begin{ylh}{1}
    \item 3.6. Κριτήρια ισότητας ορθογώνιων τριγώνων
    \item 5.2. Παραλληλόγραμμα.
\end{ylh}
% ===================================================================
%  Αρχείο: 37008.tex (Fragment)
%  Αυτόματη μετατροπή μέσω doc2latex.py
% ===================================================================



Έστω ορθογώνιο $AB\varGamma \varDelta  $ και τα σημεία $N$ και Κ των $AB$ και $\varDelta  \varGamma $ αντίστοιχα, τέτοια ώστε $AN$ = $K\varGamma $.

Να αποδείξετε ότι:

\begin{alist}
\item τα τρίγωνα $AN\varDelta  $ και $B\varGamma $Κ είναι ίσα , \monades{12}

\item το τετράπλευρο $NBK\varDelta  $ είναι παραλληλόγραμμο. \monades{13}
\end{alist}

\vspace{4mm}

% --- Άσκηση 37009 (37009.tex) ---
\Askhsh{37009}{2}
\begin{ylh}{1}
    \item 3.1. Είδη και στοιχεία τριγώνων
    \item 4.2. Τέμνουσα δύο ευθειών - Ευκλείδειο αίτημα
    \item 4.6. Άθροισμα γωνιών τριγώνου.
\end{ylh}
% ===================================================================
%  Αρχείο: 37009.tex (Fragment)
%  Αυτόματη μετατροπή μέσω doc2latex.py
% ===================================================================



Δίνεται ορθογώνιο τρίγωνο $AB\varGamma $ ($\hat{A} = 90^{0}$) και $A\varDelta  $ η διχοτόμος της γωνίας $\hat{A}$. Από το σημείο $\varDelta  $ φέρουμε την παράλληλη προς την $AB$ που τέμνει την $A\varGamma $ στο Ε.

\begin{alist}
\item Να αποδείξετε ότι το τρίγωνο $E\varDelta  $Γ είναι ορθογώνιο. \monades{9}

\item Να υπολογίσετε τη γωνία Α$\hat{\varDelta  }$Ε. \monades{9}

\item Αν η γωνία $\hat{B}$ είναι 20\textsuperscript{ο} με$\varGamma  A$λύτερη της γωνίας $\hat{\varGamma }$, να υπολογίσετε τη γωνία $Ε\hat{\varDelta  }\varGamma $.

\monades{7}
\end{alist}

\vspace{4mm}

% --- Άσκηση 37010 (37010.tex) ---
\Askhsh{37010}{2}
\begin{ylh}{1}
    \item 3.6. Κριτήρια ισότητας ορθογώνιων τριγώνων
    \item 5.10. Τραπέζιο
    \item 5.11. Ισοσκελές τραπέζιο.
\end{ylh}
% ===================================================================
%  Αρχείο: 37010.tex (Fragment)
%  Αυτόματη μετατροπή μέσω doc2latex.py
% ===================================================================



Δίνεται ισοσκελές τραπέζιο $AB\varGamma \varDelta  $ ($AB$//$\varGamma \varDelta  $ ) με $AB$=8 και $\varDelta  \varGamma $=12. Αν $AH$ και $B\varTheta $ είναι τα ύψη του τραπεζίου $AB\varGamma \varDelta  $,

\begin{alist}
\item να αποδείξετε ότι $\varDelta  H$ = $\varTheta \varGamma $. \monades{12}

\item να υπολογίσετε τη διάμεσο του τραπεζίου. \monades{13}
\end{alist}

\vspace{4mm}

% --- Άσκηση 37011 (37011.tex) ---
\Askhsh{37011}{2}
\begin{ylh}{1}
    \item 3.1. Είδη και στοιχεία τριγώνων
    \item 3.11. Ανισοτικές σχέσεις πλευρών και γωνιών
    \item 4.2. Τέμνουσα δύο ευθειών - Ευκλείδειο αίτημα
    \item 4.6. Άθροισμα γωνιών τριγώνου
    \item 5.2. Παραλληλόγραμμα
    \item 5.4. Ρόμβος.
\end{ylh}
% ===================================================================
%  Αρχείο: 37011.tex (Fragment)
%  Αυτόματη μετατροπή μέσω doc2latex.py
% ===================================================================



Στο τραπέζιο του παρακάτω σχήματος έχουμε $AB = A\varDelta   = \frac{\varGamma \varDelta  }{2}$, $\hat{\varDelta  } = 60^{\circ}$ και $M$ το μέσο της πλευράς $\varGamma \varDelta  $.

Να αποδείξετε ότι:

\begin{alist}
\item η διαγώνιος $\varDelta   B$ του τραπεζίου είναι διχοτόμος της γωνίας $\hat{\varDelta  }$, \monades{9}
\item η $BM$ χωρίζει το τραπέζιο σε ένα ρόμβο και ένα ισόπλευρο τρίγωνο. \monades{16}

\begin{center}
\begin{tikzpicture}[scale=1.2]
\tkzDefPoint(0,0){D}
\tkzDefPoint(4,0){C}
\tkzDefMidPoint(D,C) \tkzGetPoint{M}
\tkzDefPointBy[rotation=center D angle 60](M) \tkzGetPoint{A}
\tkzDefPointBy[translation=from D to M](A) \tkzGetPoint{B}

\draw[l3] (A)--(B)--(C)--(D)--cycle;
\draw[l3] (D)--(B);
\draw[l3] (B)--(M);

\tkzDrawPoints(A,B,C,D,M)
\tkzLabelPoint[above left](A){$A$}
\tkzLabelPoint[above right](B){$B$}
\tkzLabelPoint[below right](C){$\varGamma $}
\tkzLabelPoint[below left](D){$\varDelta  $}
\tkzLabelPoint[below](M){$M$}

\tkzMarkAngle[size=0.5](C,D,B)
\tkzMarkAngle[size=0.6](B,D,A)
\tkzMarkAngle[size=0.4](M,C,B)
\tkzLabelAngle[pos=0.7](M,C,B){$60^{\circ}$}
\tkzMarkSegments[mark=s|](A,B A,D D,M M,C B,C B,M)
\end{tikzpicture}
\end{center}
\end{alist}

\vspace{4mm}

% --- Άσκηση 37012 (37012.tex) ---
\Askhsh{37012}{2}
\begin{ylh}{1}
    \item 3.1. Είδη και στοιχεία τριγώνων
    \item 3.2. 1ο Κριτήριο ισότητας τριγώνων
    \item 3.6. Κριτήρια ισότητας ορθογώνιων τριγώνων.
\end{ylh}
% ===================================================================
%  Αρχείο: 37012.tex (Fragment)
%  Αυτόματη μετατροπή μέσω doc2latex.py
% ===================================================================



Δίνεται ισοσκελές τρίγωνο $ΑΒ\varGamma $ ($ΑΒ=A\varGamma $). Στα σημεία $B$ και $\varGamma $ της $B\varGamma $ φέρουμε προς το ίδιο μέρος της $B\varGamma $, τα τμήματα $B\varDelta   \perp B\varGamma $ και $\varGamma  E \perp B\varGamma $ τέτοια ώστε $B\varDelta   = \varGamma  E$. Αν $M$ είναι το μέσο της $B\varGamma $, να αποδείξετε ότι:

\begin{alist}
\item τα τρίγωνα $B\varDelta   M$ και $\varGamma  EM$ είναι ίσα, \monades{12}
\item $A\varDelta   = AE$. \monades{13}

\begin{center}
\begin{tikzpicture}[scale=1]
\tkzDefPoint(0,4){A}
\tkzDefPoint(-2,0){B}
\tkzDefPoint(2,0){C}
\tkzDefMidPoint(B,C) \tkzGetPoint{M}
\tkzDefShiftPoint[B](0,1.5){Delta}
\tkzDefShiftPoint[C](0,1.5){E}

\draw[l3] (A)--(B)--(C)--cycle;
\draw[l3] (B)--(Delta);
\draw[l3] (C)--(E);
\draw[l3] (Delta)--(M)--(E);
\draw[l3] (A)--(Delta);
\draw[l3] (A)--(E);

\tkzDrawPoints(A,B,C,Delta,E,M)
\tkzLabelPoint[above](A){$A$}
\tkzLabelPoint[left](B){$B$}
\tkzLabelPoint[right](C){$\varGamma $}
\tkzLabelPoint[left](Delta){$\varDelta  $}
\tkzLabelPoint[right](E){$E$}
\tkzLabelPoint[below](M){$M$}

\tkzMarkRightAngle(Delta,B,C)
\tkzMarkRightAngle(E,C,B)
\tkzMarkSegments[mark=s|](B,Delta C,E)
\tkzMarkSegments[mark=s||](B,M M,C)
\tkzMarkSegments[mark=s|||](A,B A,C)
\end{tikzpicture}
\end{center}
\end{alist}

\vspace{4mm}

% --- Άσκηση 37013 (37013.tex) ---
\Askhsh{37013}{2}
\begin{ylh}{1}
    \item 3.6. Κριτήρια ισότητας ορθογώνιων τριγώνων
    \item 4.6. Άθροισμα γωνιών τριγώνου.
\end{ylh}
% ===================================================================
%  Αρχείο: 37013.tex (Fragment)
%  Αυτόματη μετατροπή μέσω doc2latex.py
% ===================================================================



Έστω ισοσκελές τρίγωνο $AB\varGamma $ ($AB=A\varGamma $).

\begin{alist}
\item Να αποδείξετε ότι τα μέσα $\varDelta$ και $E$ των πλευρών $AB$ και $A\varGamma $ αντίστοιχα, ισαπέχουν από τη βάση $B\varGamma $. \monades{13}

\item Αν $\overset{\land}{A} = 75\degree + \overset{\land}{B}$, $\hat{\mathbf{A}}\mathbf{=}\mathbf{75}^{\mathbf{Ο}}\mathbf{+ \ }\hat{B}$να υπολογίσετε τις γωνίες του τριγώνου $AB\varGamma $. \monades{12}
\end{alist}

\vspace{4mm}

% --- Άσκηση 37014 (37014.tex) ---
\Askhsh{37014}{2}
\begin{ylh}{1}
    \item 3.11. Ανισοτικές σχέσεις πλευρών και γωνιών
    \item 4.6. Άθροισμα γωνιών τριγώνου.
\end{ylh}
% ===================================================================
%  Αρχείο: 37014.tex (Fragment)
%  Αυτόματη μετατροπή μέσω doc2latex.py
% ===================================================================



Στο παρακάτω σχήμα, να αποδείξετε ότι:

\begin{alist}
\item το τρίγωνο $ΑΒ\varGamma $ είναι ισοσκελές, \monades{12}
\item η γωνία $ΑE\varDelta  $ είναι ορθή. \monades{13}

\begin{center}
\begin{tikzpicture}[scale=0.8]
\tkzDefPoint(0,0){A}
\tkzDefPoint(3,0){B}
\tkzDefPoint(7,0){Delta}

% Construct isosceles triangle with base angles 70°
\tkzDefPointBy[rotation=center A angle 70](B) \tkzGetPoint{A_ray}
\tkzDefPointBy[rotation=center B angle -70](A) \tkzGetPoint{B_ray}
\tkzInterLL(A,A_ray)(B,B_ray) \tkzGetPoint{Gamma}

% E is on line $A\varGamma $ such that angle $E\varDelta  $Α = 20°
\tkzDefPointBy[rotation=center Delta angle 20](A) \tkzGetPoint{Delta_ray}
\tkzInterLL(A,Gamma)(Delta,Delta_ray) \tkzGetPoint{E}

\draw[l3] (A)--(Gamma)--(B)--cycle;
\draw[l3] (B)--(Delta)--(E);
\draw[l3] (A)--(E);

\tkzDrawPoints(A,B,Gamma,Delta,E)
\tkzLabelPoint[left](A){$A$}
\tkzLabelPoint[below](B){$B$}
\tkzLabelPoint[above](Gamma){$\varGamma $}
\tkzLabelPoint[right](Delta){$\varDelta  $}
\tkzLabelPoint[left](E){$ E$}

\tkzMarkAngle[size=0.6](B,A,Gamma)
\tkzLabelAngle[pos=0.9,font=\scriptsize](B,A,Gamma){$70\degree$}
\tkzMarkAngle[size=0.6](Gamma,B,A)
\tkzLabelAngle[pos=0.9,font=\scriptsize](Gamma,B,A){$70\degree$}
\tkzMarkAngle[size=0.8](A,Gamma,B)
\tkzLabelAngle[pos=1.2,font=\scriptsize](A,Gamma,B){$40\degree$}
\tkzMarkAngle[size=1.2](A,Delta,E)
\tkzLabelAngle[pos=1.5,font=\scriptsize](A,Delta,E){$20\degree$}
\tkzMarkAngle(Delta,E,A)
\end{tikzpicture}
\end{center}
\end{alist}

\vspace{4mm}

% --- Άσκηση 37015 (37015.tex) ---
\Askhsh{37015}{2}
\begin{ylh}{1}
    \item 5.2. Παραλληλόγραμμα.
\end{ylh}
% ===================================================================
%  Αρχείο: 37015.tex (Fragment)
%  Αυτόματη μετατροπή μέσω doc2latex.py
% ===================================================================



Δίνεται τρίγωνο $AB\varGamma $ με $AB<A\varGamma $ και $M$ το μέσο της $B\varGamma $. Προεκτείνουμε τη διάμεσο $AM$ κατά τμήμα $M\varDelta  $=$MA$. Από το Α φέρουμε παράλληλη προς τη $B\varGamma $ η οποία τέμνει την προέκταση της $\varDelta  \varGamma $ στο σημείο $E$.

Να αποδείξετε ότι:

\begin{alist}
\item το τετράπλευρο Α$B\varDelta  $Γ είναι παραλληλόγραμμο, \monades{12}

\item $ΒΜ =$ $\frac{AE}{2}$ . \monades{13}
\end{alist}

\vspace{4mm}

% --- Άσκηση 37016 (37016.tex) ---
\Askhsh{37016}{2}
\begin{ylh}{1}
    \item 3.2. 1ο Κριτήριο ισότητας τριγώνων
    \item 4.6. Άθροισμα γωνιών τριγώνου
    \item 5.9. Μια ιδιότητα του ορθογώνιου τριγώνου.
\end{ylh}
% ===================================================================
%  Αρχείο: 37016.tex (Fragment)
%  Αυτόματη μετατροπή μέσω doc2latex.py
% ===================================================================



Δίνεται τρίγωνο $ΑΒ\varGamma $ τέτοιο, ώστε $A\varGamma  < AB$. Στην πλευρά $AB$ θεωρούμε σημείο $\varDelta  $ τέτοιο ώστε $A\varDelta   = A\varGamma $ και στην προέκταση της $BA$ (προς το $A$) θεωρούμε σημείο $E$ τέτοιο ώστε $AE = A\varGamma $.

Να αποδείξετε ότι:

\begin{alist}
\item τα τμήματα $\varDelta  \varGamma $ και $E\varGamma $ είναι κάθετα μεταξύ τους, \monades{12}
\item η γωνία $E\hat{A}\varGamma $ είναι διπλάσια της γωνίας $A\hat{\varDelta  }\varGamma $. \monades{13}

\begin{center}
\begin{tikzpicture}[scale=1]
\tkzDefPoint(0,0){A}
\tkzDefPoint(5,0){B}
\tkzDefPoint(70:3){Gamma}
\tkzDefPoint(3,0){Delta}
\tkzDefPoint(-3,0){E}

\draw[l3] (A)--(B)--(Gamma)--cycle;
\draw[l3] (E)--(Gamma)--({Delta});
\draw[l3] (E)--(Delta);

\tkzDrawPoints(A,B,Gamma,Delta,E)
\tkzLabelPoint[below](A){$A$}
\tkzLabelPoint[below](B){$B$}
\tkzLabelPoint[above](Gamma){$\varGamma $}
\tkzLabelPoint[below](Delta){$\varDelta  $}
\tkzLabelPoint[below](E){$E$}

\tkzMarkSegments[mark=s|](A,Gamma A,Delta A,E)
\tkzMarkRightAngle(E,Gamma,Delta)
\tkzMarkAngle[size=0.6](A,Delta,Gamma)
\tkzMarkAngle[size=0.5](Gamma,E,A)
\tkzMarkAngle[size=0.4](E,A,Gamma)
\end{tikzpicture}
\end{center}
\end{alist}

\vspace{4mm}

% --- Άσκηση 37017 (37017.tex) ---
\Askhsh{37017}{2}
\begin{ylh}{1}
    \item 4.2. Τέμνουσα δύο ευθειών - Ευκλείδειο αίτημα
    \item 4.6. Άθροισμα γωνιών τριγώνου
    \item 5.2. Παραλληλόγραμμα
    \item 5.9. Μια ιδιότητα του ορθογώνιου τριγώνου.
\end{ylh}
% ===================================================================
%  Αρχείο: 37017.tex (Fragment)
%  Αυτόματη μετατροπή μέσω doc2latex.py
% ===================================================================



Σε παραλληλόγραμμο $ΑΒ\varGamma \varDelta  $ είναι $\hat{B} = 120^{\circ}$ και $\varDelta   Ε \perp B\varGamma $. Έστω $EZ$ η διάμεσος του τριγώνου $\varDelta   E\varGamma $.

\begin{alist}
\item Να υπολογίσετε τις γωνίες $\hat{A}$ και $\hat{\varGamma }$ του παραλληλογράμμου. \monades{8}
\item Αν $K$ είναι το μέσο της πλευράς $AB$, να αποδείξετε ότι $EZ = AK$. \monades{9}
\item Να υπολογίσετε τη γωνία $E\hat{Z}\varGamma $. \monades{8}

\begin{center}
\begin{tikzpicture}[scale=1]
\tkzDefPoint(0,0){D}
\tkzDefPoint(4,0){C}
\tkzDefPointBy[rotation=center D angle 120](C) \tkzGetPoint{A_ray}
\tkzDefPoint(2,0){temp}
\tkzDefPointBy[rotation=center D angle 120](temp) \tkzGetPoint{A}
\tkzDefPointBy[translation=from D to C](A) \tkzGetPoint{B}

\tkzDefLine[orthogonal=through D](B,C) \tkzGetPoint{DE_dir}
\tkzInterLL(D,DE_dir)(B,C) \tkzGetPoint{E}
\tkzDefMidPoint(D,C) \tkzGetPoint{Z}
\tkzDefMidPoint(A,B) \tkzGetPoint{K}

\draw[l3] (A)--(B)--(C)--(D)--cycle;
\draw[l3] (D)--(E);
\draw[l3] (E)--(Z);

\tkzDrawPoints(A,B,C,D,E,Z,K)
\tkzLabelPoint[above](A){$A$}
\tkzLabelPoint[above](B){$B$}
\tkzLabelPoint[below](C){$\varGamma $}
\tkzLabelPoint[below](D){$\varDelta  $}
\tkzLabelPoint[right](E){$E$}
\tkzLabelPoint[below](Z){$Z$}
\tkzLabelPoint[above](K){$K$}

\tkzMarkRightAngle(D,E,C)
\tkzMarkAngle[size=0.5](A,B,C)
\tkzLabelAngle[pos=0.8](A,B,C){$120\degree$}
\end{tikzpicture}
\end{center}
\end{alist}

\vspace{4mm}

% --- Άσκηση 37075 (37075.tex) ---
\Askhsh{37075}{4}
\begin{ylh}{1}
    \item 3.1. Είδη και στοιχεία τριγώνων
    \item 3.2. 1ο Κριτήριο ισότητας τριγώνων
    \item 3.4. 3ο Κριτήριο ισότητας τριγώνων
    \item 3.11. Ανισοτικές σχέσεις πλευρών και γωνιών
    \item 4.6. Άθροισμα γωνιών τριγώνου
    \item 5.9. Μια ιδιότητα του ορθογώνιου τριγώνου.
\end{ylh}
% ===================================================================
%  Αρχείο: 37075.tex (Fragment)
%  Αυτόματη μετατροπή μέσω doc2latex.py
% ===================================================================



Δίνεται ευθύγραμμο τμήμα $ΑΒ$ και στο εσωτερικό του θεωρούμε τα σημεία $\varGamma , \varDelta  $ ώστε να ισχύει $Α\varGamma  = \varGamma \varDelta   = \varDelta   B$. Επίσης θεωρούμε σημείο $Ο$ εκτός του ευθυγράμμμου τμήματος $ΑΒ$ έτσι ώστε να ισχύουν $Ο\varGamma  = A\varGamma $ και $Ο\varDelta   = \varDelta   B$.

\begin{alist}
\item Να αποδείξετε ότι:
\begin{alist}
\item η γωνία $\varGamma \hat{O}\varDelta  $ είναι $60^{\circ}$, \monades{9}
\item οι γωνίες $Ο\hat{A}\varGamma $ και $Ο\hat{B}\varDelta  $ είναι ίσες και κάθε μία ίση με $30\degree$. \monades{9}
\end{alist}
\item Αν $M$ το μέσον του ευθυγράμμου τμήματος $ΑΒ$, να αποδείξετε ότι $2ΟΜ = ΟΑ$. \monades{7}

\begin{center}
\begin{tikzpicture}[scale=1.2]
\tkzDefPoint(0,0){A}
\tkzDefPoint(6,0){B}
\tkzDefPoint(2,0){Gamma}
\tkzDefPoint(4,0){Delta}
\tkzDefEquilateral(Gamma,Delta) \tkzGetPoint{O}
\tkzDefMidPoint(A,B) \tkzGetPoint{M}

\draw[l3] (A)--(B);
\draw[l3] (A)--(O)--(B);
\draw[l3] (O)--(Gamma);
\draw[l3] (O)--(Delta);
\draw[l3] (O)--(M);

\tkzDrawPoints(A,B,Gamma,Delta,O,M)
\tkzLabelPoint[below](A){$A$}
\tkzLabelPoint[below](B){$B$}
\tkzLabelPoint[below](Gamma){$\varGamma $}
\tkzLabelPoint[below](Delta){$\varDelta  $}
\tkzLabelPoint[above](O){$O$}
\tkzLabelPoint[below](M){$M$}

\tkzMarkSegments[mark=s|](A,Gamma Gamma,Delta Delta,B O,Gamma O,Delta)
\tkzMarkAngle[size=0.6](Gamma,O,Delta)
\tkzLabelAngle[pos=0.8](Gamma,O,Delta){$60\degree$}
\tkzMarkAngle[size=0.8](Gamma,A,O)
\tkzLabelAngle[pos=1.1](Gamma,A,O){$30\degree$}
\tkzMarkAngle[size=0.8](O,B,Delta)
\tkzLabelAngle[pos=1.1](O,B,Delta){$30\degree$}
\end{tikzpicture}
\end{center}
\end{alist}

\vspace{4mm}

% --- Άσκηση 37080 (37080.tex) ---
\Askhsh{37080}{4}
\begin{ylh}{1}
    \item 5.2. Παραλληλόγραμμα
    \item 5.6. Εφαρμογές στα τρίγωνα
    \item 5.10. Τραπέζιο.
\end{ylh}
% ===================================================================
%  Αρχείο: 37080.tex (Fragment)
%  Αυτόματη μετατροπή μέσω doc2latex.py
% ===================================================================



Σε τραπέζιο $ΑΒ\varGamma \varDelta  $ ($ΑΒ//\varGamma \varDelta  $) είναι $\varGamma \varDelta   = 2ΑΒ$. Επίσης, τα σημεία $Ζ, H$ και $E$ είναι τα μέσα των $Α\varDelta  , B\varGamma $ και $\varDelta  \varGamma $ αντίστοιχα. Ακόμη η $ΖH$ τέμνει τις $AE, BE$ στα σημεία $\varTheta , Ι$ αντίστοιχα.

\begin{alist}
\item Να αποδείξετε ότι το τετράπλευρο $ΑΒ\varGamma  E$ είναι παραλληλόγραμμο. \monades{10}
\item Να δείξετε ότι:
\begin{alist}
\item τα σημεία $\varTheta , Ι$ είναι μέσα των $AE, BE$ αντίστοιχα. \monades{5}
\item $ΖH = \frac{3}{2}ΑΒ$. \monades{10}
\end{alist}

\begin{center}
\begin{tikzpicture}[scale=1.2]
\tkzDefPoint(0,0){D}
\tkzDefPoint(4,0){C}
\tkzDefMidPoint(D,C) \tkzGetPoint{E}
\tkzDefPoint(0,2){A}
\tkzDefPoint(2,2){B}

\tkzDefMidPoint(A,D) \tkzGetPoint{Z}
\tkzDefMidPoint(B,C) \tkzGetPoint{H}
\tkzInterLL(Z,H)(A,E) \tkzGetPoint{Theta}
\tkzInterLL(Z,H)(B,E) \tkzGetPoint{I}

\draw[l3] (A)--(B)--(C)--(D)--cycle;
\draw[l3] (A)--(E)--(B);
\draw[l3] (Z)--(H);

\tkzDrawPoints(A,B,C,D,E,Z,H,Theta,I)
\tkzLabelPoint[above left](A){$A$}
\tkzLabelPoint[above right](B){$B$}
\tkzLabelPoint[below right](C){$\varGamma $}
\tkzLabelPoint[below left](D){$\varDelta  $}
\tkzLabelPoint[below](E){$E$}
\tkzLabelPoint[left](Z){$Z$}
\tkzLabelPoint[right](H){$H$}
\tkzLabelPoint[above left](Theta){$\varTheta $}
\tkzLabelPoint[above left](I){$I$}

\tkzMarkSegments[mark=s|](A,Z Z,D B,H H,C D,E E,C)
\tkzMarkSegments[mark=s||](A,B D,E E,C)
\end{tikzpicture}
\end{center}
\end{alist}

\vspace{4mm}

% --- Άσκηση 37082 (37082.tex) ---
\Askhsh{37082}{4}
\begin{ylh}{1}
    \item 4.6. Άθροισμα γωνιών τριγώνου
    \item 5.6. Εφαρμογές στα τρίγωνα.
\end{ylh}
% ===================================================================
%  Αρχείο: 37082.tex (Fragment)
%  Αυτόματη μετατροπή μέσω doc2latex.py
% ===================================================================



Σε τρίγωνο $ΑΒ\varGamma $ ισχύει $\hat{A} + \hat{\varGamma } = 2\hat{B}$ και έστω $A\varDelta  $ ύψος και $BE$ διχοτόμος του τριγώνου που τέμνονται στο $Z$.

\begin{alist}
\item Να αποδείξετε ότι:
\begin{alist}
\item $\hat{B} = 60\degree$ και $AZ = BZ$, \monades{10}
\item $A\varDelta   = \frac{3}{2}BZ$ \monades{8}
\end{alist}
\item Αν είναι γνωστό ότι το τρίγωνο $AZE$ είναι ισόπλευρο, να υπολογίσετε τις άλλες γωνίες του τριγώνου $ΑΒ\varGamma $. \monades{7}

\begin{center}
\begin{tikzpicture}[scale=1.2]
\tkzDefPoint(0,0){B}
\tkzDefPoint(5,0){C}
\tkzDefPointBy[rotation=center B angle 60](C) \tkzGetPoint{B_ray}
\tkzDefPointBy[rotation=center C angle -50](B) \tkzGetPoint{C_ray}
\tkzInterLL(B,B_ray)(C,C_ray) \tkzGetPoint{A}

\tkzDefPointBy[projection=onto B--C](A) \tkzGetPoint{Delta}
\tkzDefLine[bisector](A,B,C) \tkzGetPoint{BE_dir}
\tkzInterLL(B,BE_dir)(A,C) \tkzGetPoint{E}
\tkzInterLL(A,Delta)(B,E) \tkzGetPoint{Z}

\draw[l3] (A)--(B)--(C)--cycle;
\draw[l3] (A)--(Delta);
\draw[l3] (B)--(E);

\tkzDrawPoints(A,B,C,Delta,E,Z)
\tkzLabelPoint[above](A){$A$}
\tkzLabelPoint[below](B){$B$}
\tkzLabelPoint[below](C){$\varGamma $}
\tkzLabelPoint[below](Delta){$\varDelta  $}
\tkzLabelPoint[right](E){$E$}
\tkzLabelPoint[below left](Z){$Z$}

\tkzMarkRightAngle(A,Delta,C)
\tkzMarkAngle[size=0.6,mark=s|](C,B,E)
\tkzMarkAngle[size=0.7,mark=s|](E,B,A)
\tkzMarkAngle[size=0.5](A,B,C)
\tkzLabelAngle[pos=0.8](A,B,C){$60\degree$}
\tkzMarkSegments[mark=s||](A,Z B,Z)
\end{tikzpicture}
\end{center}
\end{alist}

\vspace{4mm}

% --- Άσκηση 37083 (37083.tex) ---
\Askhsh{37083}{4}
\begin{ylh}{1}
    \item 3.6. Κριτήρια ισότητας ορθογώνιων τριγώνων
    \item 5.2. Παραλληλόγραμμα
    \item 5.3. Ορθογώνιο.
\end{ylh}
% ===================================================================
%  Αρχείο: 37083.tex (Fragment)
%  Αυτόματη μετατροπή μέσω doc2latex.py
% ===================================================================



Στην παρακάτω εικόνα φαίνεται μια κρεμάστρα τοίχου η οποία αποτελείται από έξι ίσα ευθύγραμμα κομμάτια ξύλου ($A\varDelta  , B\varGamma , \varGamma  Z, \varDelta   H, $ZK$, H\Lambda$) που είναι στερεωμένα με έντεκα καρφιά ($A, B, \varGamma , \varDelta  , \varTheta , Ε, Μ, H, Κ, \Lambda, Ζ$). Αν το σημείο $\varTheta $ είναι μέσο των τμημάτων $A\varDelta  $ και $B\varGamma $, ενώ το σημείο $Ε$ είναι μέσο των τμημάτων $\varGamma  Z$ και $\varDelta   H$, να αποδείξετε ότι:

\begin{alist}
\item Το τετράπλευρο $\varGamma  H Ζ \varDelta  $ είναι ορθογώνιο. \monades{10}
\item Τα σημεία $B, \varDelta  , Ζ$ είναι συνευθειακά. \monades{9}
\item Το τετράπλευρο $A\varGamma  Ζ \varDelta  $ είναι παραλληλόγραμμο. \monades{6}

\begin{center}
\begin{tikzpicture}[scale=0.8]
\tkzDefPoint(0,2){A}
\tkzDefPoint(0,0){B}
\tkzDefPoint(2,2){G}
\tkzDefPoint(2,0){D}
\tkzDefPoint(4,0){Z}
\tkzDefPoint(4,2){H}
\tkzDefPoint(6,2){K}
\tkzDefPoint(6,0){L}

\tkzDefMidPoint(A,D) \tkzGetPoint{Theta}
\tkzDefMidPoint(G,Z) \tkzGetPoint{E}
\tkzDefMidPoint(Z,K) \tkzGetPoint{M}

\draw[l3, line width=2pt, brown] (A)--(D);
\draw[l3, line width=2pt, brown] (B)--(G);
\draw[l3, line width=2pt, brown] (G)--(Z);
\draw[l3, line width=2pt, brown] (D)--(H);
\draw[l3, line width=2pt, brown] (Z)--(K);
\draw[l3, line width=2pt, brown] (H)--(L);

\tkzDrawPoints[size=5](A,B,G,D,Theta,E,M,H,K,L,Z)
\tkzLabelPoint[above](A){$A$}
\tkzLabelPoint[below](B){$B$}
\tkzLabelPoint[above](G){$\varGamma $}
\tkzLabelPoint[below](D){$\varDelta  $}
\tkzLabelPoint[below](Z){$Z$}
\tkzLabelPoint[above](H){$H$}
\tkzLabelPoint[above](K){$K$}
\tkzLabelPoint[below](L){$\Lambda$}
\tkzLabelPoint[above](Theta){$\varTheta $}
\tkzLabelPoint[below](E){$E$}
\tkzLabelPoint[above](M){$M$}
\end{tikzpicture}
\end{center}
\end{alist}

\vspace{4mm}

% --- Άσκηση 37084 (37084.tex) ---
\Askhsh{37084}{4}
\begin{ylh}{1}
    \item 3.6. Κριτήρια ισότητας ορθογώνιων τριγώνων
    \item 4.2. Τέμνουσα δύο ευθειών - Ευκλείδειο αίτημα
    \item 5.3. Ορθογώνιο
    \item 5.10. Τραπέζιο.
\end{ylh}
% ===================================================================
%  Αρχείο: 37084.tex (Fragment)
%  Αυτόματη μετατροπή μέσω doc2latex.py
% ===================================================================



Δίνεται ευθεία (ε) και δυο σημεία $A$, Β εκτός αυτής έτσι ώστε η ευθεία $AB$ να μην είναι κάθετη στην (ε). Φέρουμε $A\varDelta  $, $B\varGamma $ κάθετες στην (ε) και Μ, Ν μέσα των $AB$ και $\varGamma \varDelta  $ αντίστοιχα.

\begin{alist}
\item Αν τα Α, Β είναι στο ίδιο ημιεπίπεδο σε σχέση με την (ε),


\end{alist}
\begin{alist}
\item
  να εξετάσετε αν το τετράπλευρο $AB\varGamma \varDelta  $ είναι παραλληλόγραμμο, τραπέζιο ή ορθογώνιο σε καθεμία από τις παρακάτω περιπτώσεις, αιτιολογώντας την απάντησή σας:

\end{alist}
\begin{alist}
\item
    $A\varDelta  $\textless $B\varGamma $ \monades{5}
\item
    $A\varDelta  $=$B\varGamma $ \monades{5}
\item
  να εκφράσετε το τμήμα $MN$ σε σχέση με τα τμήματα $A\varDelta  $, $B\varGamma $ στις δυο προηγούμενες περιπτώσεις. \monades{6}

\item Αν η ευθεία (ε) τέμνει το τμήμα $AB$ στο μέσο του Μ, να βρείτε το είδος του τετραπλεύρου $AB\varGamma \varDelta  $ (παραλληλόγραμμο, τραπέζιο, ορθογώνιο), αιτιολογώντας την απάντησή σας.

\monades{9}
\end{alist}

\vspace{4mm}

% --- Άσκηση 37085 (37085.tex) ---
\Askhsh{37085}{4}
\begin{ylh}{1}
    \item 4.6. Άθροισμα γωνιών τριγώνου
    \item 5.8. Το ορθόκεντρο τριγώνου
    \item 5.9. Μια ιδιότητα του ορθογώνιου τριγώνου.
\end{ylh}
% ===================================================================
%  Αρχείο: 37085.tex (Fragment)
%  Αυτόματη μετατροπή μέσω doc2latex.py
% ===================================================================



Στο παρακάτω σχήμα δίνεται οξυγώνιο τρίγωνο $ΑΒ\varGamma $, τα ύψη του $Β\varDelta  $ και $\varGamma  E$ που τέμνονται στο σημείο $H$ και το μέσο $M$ της πλευράς $B\varGamma $.

\begin{alist}
\item Να αποδείξετε ότι:
\begin{alist}
\item $M\varDelta   = ME$, \monades{10}
\item η ευθεία $AH$ τέμνει κάθετα τη $B\varGamma $ και ότι $A\hat{H}\varDelta   = \hat{\varGamma }$, όπου $\hat{\varGamma }$ η γωνία του τριγώνου $ΑΒ\varGamma $. \monades{5}
\end{alist}
\item Να βρείτε το ορθόκεντρο του τριγώνου $ABH$. \monades{10}

\begin{center}
\begin{tikzpicture}[scale=1.2]
\tkzDefPoint(1.2,3.5){A}
\tkzDefPoint(0,0){B}
\tkzDefPoint(4,0){C}
\tkzDefPointBy[projection=onto A--C](B) \tkzGetPoint{Delta}
\tkzDefPointBy[projection=onto A--B](C) \tkzGetPoint{E}
\tkzInterLL(B,Delta)(C,E) \tkzGetPoint{H}
\tkzDefMidPoint(B,C) \tkzGetPoint{M}

\draw[l3] (A)--(B)--(C)--cycle;
\draw[l3] (B)--(Delta);
\draw[l3] (C)--(E);
\draw[l3] (M)--(Delta);
\tkzDefLine[orthogonal=through A](B,C) \tkzGetPoint{h_dir}
\draw[l3, dashed] (A)--(h_dir);

\tkzDrawPoints(A,B,C,Delta,E,H,M)
\tkzLabelPoint[above](A){$A$}
\tkzLabelPoint[below left](B){$B$}
\tkzLabelPoint[below right](C){$\varGamma $}
\tkzLabelPoint[right](Delta){$\varDelta  $}
\tkzLabelPoint[left](E){$E$}
\tkzLabelPoint[below right](H){$H$}
\tkzLabelPoint[below](M){$M$}

\tkzMarkRightAngle(B,Delta,C)
\tkzMarkRightAngle(C,E,A)
\tkzMarkSegments[mark=s|](B,M M,C)
\end{tikzpicture}
\end{center}
\end{alist}

\vspace{4mm}

% --- Άσκηση 37086 (37086.tex) ---
\Askhsh{37086}{4}
\begin{ylh}{1}
    \item 3.6. Κριτήρια ισότητας ορθογώνιων τριγώνων
    \item 5.2. Παραλληλόγραμμα
    \item 5.3. Ορθογώνιο
    \item 5.6. Εφαρμογές στα τρίγωνα
    \item 5.10. Τραπέζιο.
\end{ylh}
% ===================================================================
%  Αρχείο: 37086.tex (Fragment)
%  Αυτόματη μετατροπή μέσω doc2latex.py
% ===================================================================



Θεωρούμε παραλληλόγραμμο $ΑΒ\varGamma \varDelta  $ και τις προβολές $Α', B', \varGamma ', \varDelta  '$ και $Ο'$ των κορυφών του $Α, B, \varGamma , \varDelta  $ και του κέντρου του $Ο$ αντίστοιχα, σε μία ευθεία $(\varepsilon)$.

\begin{alist}
\item Αν η ευθεία $(\varepsilon)$ αφήνει τις κορυφές του παραλληλογράμμου στο ίδιο ημιεπίπεδο και είναι $AA'=3, $BB$'=2, \varGamma \varGamma '=5$, τότε:
\begin{alist}
\item Να αποδείξετε ότι η απόσταση $ΟΟ'$ του κέντρου $Ο$ του παραλληλογράμμου από την $(\varepsilon)$ είναι ίση με $4$. \monades{8}
\item Να βρείτε την απόσταση $\Delta\varDelta  '$. \monades{9}
\end{alist}

\begin{center}
\begin{tikzpicture}[scale=0.8]
\tkzDefPoint(1,3){A}
\tkzDefPoint(5,2){B}
\tkzDefPoint(7,5){C}
\tkzDefPoint(3,6){D}
\tkzInterLL(A,C)(B,D) \tkzGetPoint{O}
\draw[maincolor, thick] (-1,0) -- (9,0) node[right]{$(\epsilon)$};

\tkzDefPoint(1,0){A_p}
\tkzDefPoint(5,0){B_p}
\tkzDefPoint(7,0){C_p}
\tkzDefPoint(3,0){D_p}
\tkzDefPoint(4,0){O_p}

\draw[l3] (A)--(B)--(C)--(D)--cycle;
\draw[l3] (A)--(C);
\draw[l3] (B)--(D);

\draw[dashed] (A)--(A_p);
\draw[dashed] (B)--(B_p);
\draw[dashed] (C)--(C_p);
\draw[dashed] (D)--(D_p);
\draw[dashed] (O)--(O_p);

\tkzDrawPoints(A,B,C,D,O,A_p,B_p,C_p,D_p,O_p)
\tkzLabelPoint[above](A){$A$}
\tkzLabelPoint[below](B){$B$}
\tkzLabelPoint[above](C){$\varGamma $}
\tkzLabelPoint[above](D){$\varDelta  $}
\tkzLabelPoint[above right](O){$O$}

\tkzLabelPoint[below](A_p){$A'$}
\tkzLabelPoint[below](B_p){$B'$}
\tkzLabelPoint[below](C_p){$\varGamma '$}
\tkzLabelPoint[below](D_p){$\varDelta  '$}
\tkzLabelPoint[below](O_p){$O'$}

\tkzMarkRightAngle(A,A_p,B_p)
\tkzMarkRightAngle(B,B_p,C_p)
\tkzMarkRightAngle(C,C_p,D_p)
\tkzMarkRightAngle(D,D_p,A_p)
\end{tikzpicture}
\end{center}

\item Αν η ευθεία $(\varepsilon)$ διέρχεται από το κέντρο του παραλληλογράμμου και είναι παράλληλη προς δύο απέναντι πλευρές του, τι παρατηρείτε για τις αποστάσεις $AA', $BB$', \varGamma \varGamma ', \Delta\varDelta  '$.
Να αιτιολογήσετε την απάντησή σας. \monades{8}
\end{alist}

\vspace{4mm}

% --- Άσκηση 37087 (37087.tex) ---
\Askhsh{37087}{4}
\begin{ylh}{1}
    \item 3.6. Κριτήρια ισότητας ορθογώνιων τριγώνων
    \item 4.2. Τέμνουσα δύο ευθειών - Ευκλείδειο αίτημα
    \item 5.3. Ορθογώνιο
    \item 5.6. Εφαρμογές στα τρίγωνα
    \item 5.8. Το ορθόκεντρο τριγώνου.
\end{ylh}
% ===================================================================
%  Αρχείο: 37087.tex (Fragment)
%  Αυτόματη μετατροπή μέσω doc2latex.py
% ===================================================================



Δίνεται ισόπλευρο τρίγωνο $ΑΒ\varGamma $ και τα ύψη του $BK$ και $\varGamma \Lambda$, τα οποία τέμνονται στο $I$. Αν $M$ και $N$ είναι τα μέσα των $IB$ και $I\varGamma $ αντίστοιχα, να αποδείξετε ότι:

\begin{alist}
\item Το $AI$ προεκτεινόμενο διέρχεται από το μέσο της πλευράς $B\varGamma $. \monades{10}
\item Το τετράπλευρο $M\Lambda KN$ είναι ορθογώνιο παραλληλόγραμμο. \monades{15}

\begin{center}
\begin{tikzpicture}[scale=1.2]
\tkzDefPoint(0,0){B}
\tkzDefPoint(4,0){C}
\tkzDefEquilateral(B,C) \tkzGetPoint{A}
\tkzDefPointBy[projection=onto A--C](B) \tkzGetPoint{K}
\tkzDefPointBy[projection=onto A--B](C) \tkzGetPoint{L}
\tkzInterLL(B,K)(C,L) \tkzGetPoint{I}
\tkzDefMidPoint(I,B) \tkzGetPoint{M}
\tkzDefMidPoint(I,C) \tkzGetPoint{N}
\tkzDefMidPoint(B,C) \tkzGetPoint{midBC}

\draw[l3] (A)--(B)--(C)--cycle;
\draw[l3] (B)--(K);
\draw[l3] (C)--(L);
\draw[l3] (A)--(midBC);
\draw[l3] (M)--(L)--(K)--(N)--cycle;

\tkzDrawPoints(A,B,C,K,L,I,M,N)
\tkzLabelPoint[above](A){$A$}
\tkzLabelPoint[below left](B){$B$}
\tkzLabelPoint[below right](C){$\varGamma $}
\tkzLabelPoint[right](K){$K$}
\tkzLabelPoint[left](L){$\Lambda$}
\tkzLabelPoint[above right](I){$I$}
\tkzLabelPoint[below](M){$M$}
\tkzLabelPoint[below](N){$N$}

\tkzMarkRightAngle(B,K,C)
\tkzMarkRightAngle(C,L,B)
\end{tikzpicture}
\end{center}
\end{alist}

\vspace{4mm}

% --- Άσκηση 37088 (37088.tex) ---
\Askhsh{37088}{4}
\begin{ylh}{1}
    \item 3.14. Σχετικές θέσεις ευθείας και κύκλου
    \item 3.16. Σχετικές θέσεις δυο κύκλων
    \item 4.2. Τέμνουσα δύο ευθειών - Ευκλείδειο αίτημα
    \item 5.2. Παραλληλόγραμμα
    \item 5.3. Ορθογώνιο
    \item 5.9. Μια ιδιότητα του ορθογώνιου τριγώνου.
\end{ylh}
% ===================================================================
%  Αρχείο: 37088.tex (Fragment)
%  Αυτόματη μετατροπή μέσω doc2latex.py
% ===================================================================



Οι κύκλοι ($Κ, \rho$) και ($\Lambda, 3\rho$) εφάπτονται εξωτερικά στο σημείο $A$. Μία ευθεία $(\varepsilon)$ εφάπτεται εξωτερικά και στους δύο κύκλους στα σημεία $B$ και $\varGamma $ αντίστοιχα και τέμνει την προέκταση της διακέντρου $Κ\Lambda$ (προς το $Κ$) στο σημείο $Ε$. Φέρουμε από το σημείο $Κ$ παράλληλο τμήμα στην $(\varepsilon)$ που τέμνει το τμήμα $\Lambda\varGamma $ στο $\varDelta  $.

Να αποδείξετε ότι:

\begin{alist}
\item το τετράπλευρο $B\varGamma \varDelta   Κ$ είναι ορθογώνιο, \monades{9}
\item η γωνία $\varDelta  \hat{Κ}\Lambda$ είναι $30\degree$, \monades{8}
\item το τμήμα $E\Lambda = 6\rho$, όπου $\rho$ η ακτίνα του κύκλου ($Κ, \rho$). \monades{8}

\begin{center}
\begin{tikzpicture}[scale=1]
\tkzDefPoint(0,1){K}
\tkzDefPoint({2*sqrt(3)},3){L}
\tkzDefPoint(0,0){B}
\tkzDefPoint({2*sqrt(3)},0){C}
\tkzDefPoint({2*sqrt(3)},1){D}
\tkzDefPoint({-sqrt(3)},0){E}
\tkzDefPoint({0.5*sqrt(3)},1.5){A}

\tkzDrawCircle[l3](K,B)
\tkzDrawCircle[l3](L,C)
\draw[thick] (-3,0) -- (5,0) node[right]{$(\varepsilon)$};
\draw[l3] (E)--(L);
\draw[l3] (K)--(B);
\draw[l3] (L)--(C);
\draw[l3] (K)--(D);
\draw[l3] (K)--(L);

\tkzDrawPoints(K,L,A,B,C,D,E)
\tkzLabelPoint[above left](K){$K$}
\tkzLabelPoint[above](L){$\Lambda$}
\tkzLabelPoint[above right](A){$A$}
\tkzLabelPoint[below](B){$B$}
\tkzLabelPoint[below](C){$\varGamma $}
\tkzLabelPoint[right](D){$\varDelta  $}
\tkzLabelPoint[below](E){$E$}

\tkzMarkRightAngle(K,B,C)
\tkzMarkRightAngle(L,C,B)
\tkzMarkAngle[size=0.8](C,E,L)
\tkzLabelAngle[pos=1.1](C,E,L){$30\degree$}
\tkzMarkAngle[size=0.6](D,K,L)
\tkzLabelAngle[pos=0.9](D,K,L){$30\degree$}
\end{tikzpicture}
\end{center}
\end{alist}

\vspace{4mm}

% --- Άσκηση 37089 (37089.tex) ---
\Askhsh{37089}{4}
\begin{ylh}{1}
    \item 3.1. Είδη και στοιχεία τριγώνων
    \item 3.2. 1ο Κριτήριο ισότητας τριγώνων
    \item 3.4. 3ο Κριτήριο ισότητας τριγώνων
    \item 3.6. Κριτήρια ισότητας ορθογώνιων τριγώνων
    \item 4.2. Τέμνουσα δύο ευθειών - Ευκλείδειο αίτημα
    \item 5.10. Τραπέζιο
    \item 5.11. Ισοσκελές τραπέζιο.
\end{ylh}
% ===================================================================
%  Αρχείο: 37089.tex (Fragment)
%  Αυτόματη μετατροπή μέσω doc2latex.py
% ===================================================================



Δίνεται ορθογώνιο τρίγωνο $AB\varGamma $ ($\hat{A} = 90\degree$) και η διχοτόμος του $B\varDelta  $. Από το $\varDelta$ φέρουμε ΔΕ$\bot$$B\varGamma $ και ονομάζουμε $Z$ το σημείο στο οποίο η ευθεία $E\varDelta  $ τέμνει την προέκταση της $BA$.

Να αποδείξετε ότι:

\begin{alist}
\item το τρίγωνο Α$BE$ είναι ισοσκελές, \monades{6}

\item τα τρίγωνα $AB\varGamma $ και $BE$Ζ είναι ίσα, \monades{6}

\item η ευθεία $B\varDelta  $ είναι μεσοκάθετη των τμημάτων $AE$ και $Z\varGamma $, \monades{6}

\item το τετράπλευρο $AE\varGamma Z$ είναι ισοσκελές τραπέζιο. \monades{7}
\end{alist}

\vspace{4mm}

% --- Άσκηση 37094 (37094.tex) ---
\Askhsh{37094}{4}
\begin{ylh}{1}
    \item 3.2. 1ο Κριτήριο ισότητας τριγώνων
    \item 3.6. Κριτήρια ισότητας ορθογώνιων τριγώνων
    \item 3.11. Ανισοτικές σχέσεις πλευρών και γωνιών.
\end{ylh}
% ===================================================================
%  Αρχείο: 37094.tex (Fragment)
%  Αυτόματη μετατροπή μέσω doc2latex.py
% ===================================================================



Έστω $AB\varGamma $ τρίγωνο και τα ύψη του $BE$ και $\varGamma \varDelta  $ που αντιστοιχούν στις πλευρές $A\varGamma $ και $AB$ αντίστοιχα. Δίνεται η ακόλουθη πρόταση:

\textbf{Π:} Αν το τρίγωνο $AB\varGamma $ είναι ισοσκελές με $AB=A\varGamma $, τότε τα ύψη $BE$ και $\varGamma \varDelta  $ που αντιστοιχούν στις ίσες πλευρές του είναι ίσα.

\begin{alist}
\item Να εξετάσετε αν ισχύει η πρόταση \textbf{Π} αιτιολογώντας την απάντησή σας.

\monades{10}

\item Να διατυπώσετε την \textbf{αντίστροφη} πρόταση της \emph{\textbf{Π}} και να αποδείξετε ότι ισχύει.

\monades{10}

\item Να διατυπώσετε την πρόταση \textbf{Π} και την \textbf{αντίστροφή της} ως ενιαία πρόταση.

\monades{5}
\end{alist}

\vspace{4mm}

% --- Άσκηση 37095 (37095.tex) ---
\Askhsh{37095}{4}
\begin{ylh}{1}
    \item 3.1. Είδη και στοιχεία τριγώνων
    \item 3.2. 1ο Κριτήριο ισότητας τριγώνων
    \item 3.4. 3ο Κριτήριο ισότητας τριγώνων
    \item 3.11. Ανισοτικές σχέσεις πλευρών και γωνιών.
\end{ylh}
% ===================================================================
%  Αρχείο: 37095.tex (Fragment)
%  Αυτόματη μετατροπή μέσω doc2latex.py
% ===================================================================



Δίνεται οξεία γωνία $x\hat{O}y$ και δύο ομόκεντροι κύκλοι ($O, \rho_1$) και ($O, \rho_2$) με $\rho_1 < \rho_2$, που τέμνουν την $Ox$ στα σημεία $K, A$ και την $Oy$ στα $\Lambda, B$ αντίστοιχα.

Να αποδείξετε ότι:

\begin{alist}
\item $A\Lambda = BK$, \monades{8}
\item το τρίγωνο $APB$ είναι ισοσκελές, όπου $P$ το σημείο τομής των $A\Lambda$ και $BK$, \monades{8}
\item η $OP$ διχοτομεί την $x\hat{O}y$. \monades{9}

\begin{center}
\begin{tikzpicture}[scale=1]
\tkzDefPoint(0,0){O}
\tkzDefPoint(5,0){x_dir}
\tkzDefPoint(40:5){y_dir}
\tkzDefPoint(2,0){K}
\tkzInterLC(O,y_dir)(O,K) \tkzGetPoints{L}{L_prime}
\tkzDefPoint(4,0){A}
\tkzInterLC(O,y_dir)(O,A) \tkzGetPoints{B}{B_prime}
\tkzInterLL(A,L)(B,K) \tkzGetPoint{P}

\tkzDrawCircle[l3](O,K)
\tkzDrawCircle[l3](O,A)
\draw[l3] (O)--(x_dir) node[right]{$x$};
\draw[l3] (O)--(y_dir) node[above]{$y$};
\draw[l3] (A)--(L);
\draw[l3] (B)--(K);
\draw[l3] (O)--(P);

\tkzDrawPoints(O,K,A,L,B,P)
\tkzLabelPoint[left](O){$O$}
\tkzLabelPoint[below](K){$K$}
\tkzLabelPoint[below](A){$A$}
\tkzLabelPoint[above left](L){$\Lambda$}
\tkzLabelPoint[above left](B){$B$}
\tkzLabelPoint[right](P){$P$}
\end{tikzpicture}
\end{center}
\end{alist}

\vspace{4mm}

% --- Άσκηση 37096 (37096.tex) ---
\Askhsh{37096}{4}
\begin{ylh}{1}
    \item 3.1. Είδη και στοιχεία τριγώνων
    \item 4.2. Τέμνουσα δύο ευθειών - Ευκλείδειο αίτημα
    \item 5.6. Εφαρμογές στα τρίγωνα.
\end{ylh}
% ===================================================================
%  Αρχείο: 37096.tex (Fragment)
%  Αυτόματη μετατροπή μέσω doc2latex.py
% ===================================================================



\begin{alist}
\item Να αποδείξετε ότι το τρίγωνο με κορυφές τα μέσα πλευρών ισοσκελούς τριγώνου είναι ισοσκελές. (Μονάδες8)

\item Να διατυπώσετε και να αποδείξετε ανάλογη πρόταση για
\end{alist}
\begin{alist}
\item
  ισόπλευρο τρίγωνο. \monades{8}
\item
  ορθογώνιο και ισοσκελές τρίγωνο. \monades{9}
\end{alist}

\vspace{4mm}

% --- Άσκηση 37097 (37097.tex) ---
\Askhsh{37097}{4}
\begin{ylh}{1}
    \item 3.6. Κριτήρια ισότητας ορθογώνιων τριγώνων
    \item 4.2. Τέμνουσα δύο ευθειών - Ευκλείδειο αίτημα
    \item 4.6. Άθροισμα γωνιών τριγώνου.
\end{ylh}
% ===================================================================
%  Αρχείο: 37097.tex (Fragment)
%  Αυτόματη μετατροπή μέσω doc2latex.py
% ===================================================================



Δίνεται ισόπλευρο τρίγωνο $AB\varGamma $ και το ύψος του $\varGamma E$. Στην προέκταση της $\varGamma B$ (προς το μέρος του Β) θεωρούμε σημείο $\varDelta  $ τέτοιο ώστε $Β\varDelta   = \ \frac{Β\varGamma }{2}$. Έστω ότι, η ευθεία $\varDelta  E$ τέμνει την $A\varGamma $ στο $Z$ και $Z\varTheta $ // $B\varGamma $.

\begin{alist}
\item Να αποδείξετε ότι το τρίγωνο $B\varDelta  $Ε είναι ισοσκελές και το τρίγωνο $A\varTheta  Z$ είναι ισόπλευρο. \monades{10}

\item Να υπολογίσετε τις γωνίες του τριγώνου $\varTheta  EZ$. \monades{5}

\item Να αποδείξετε ότι $AE$ = 2$\cdot$ΘΖ. \monades{5}

\item Να αποδείξετε ότι 3$\cdot$ΑΒ = 4$\cdot$ΘΒ. \monades{5}

\begin{center}
\begin{tikzpicture}[scale=0.8]
% Equilateral triangle $AB\varGamma $
\tkzDefPoint(2,3.464){A}
\tkzDefPoint(0,0){B}
\tkzDefPoint(4,0){Gamma}

% Altitude $\varGamma E$
\tkzDefPointBy[projection=onto A--B](Gamma) \tkzGetPoint{E}

% Δ on extension of $\varGamma B$ past B: $B\varDelta  $ = $B\varGamma $/2
\tkzDefPoint(-2,0){Delta}

% $\varDelta  E$ line intersects $A\varGamma $ at Z
\tkzInterLL(Delta,E)(A,Gamma) \tkzGetPoint{Z}

% $Z\varTheta$ parallel to $B\varGamma $, intersecting $AB$ at Θ
\tkzDefLine[parallel=through Z](B,Gamma) \tkzGetPoint{par_Z}
\tkzInterLL(Z,par_Z)(A,B) \tkzGetPoint{Theta}

% Draw triangle and segments
\draw[l3] (A)--(B)--(Gamma)--cycle;
\draw[l3] (Gamma)--(Delta);
\draw[l3] (Gamma)--(E);
\draw[l3] (Delta)--(Z);
\draw[l3] (Z)--(Theta);

% Points and labels
\tkzDrawPoints(A,B,Gamma,Delta,E,Z,Theta)
\tkzLabelPoint[above](A){$A$}
\tkzLabelPoint[below](B){$B$}
\tkzLabelPoint[below](Gamma){$\varGamma $}
\tkzLabelPoint[below](Delta){$\varDelta  $}
\tkzLabelPoint[above left](E){$ E$}
\tkzLabelPoint[above right](Z){$ Z$}
\tkzLabelPoint[left](Theta){$\varTheta $}

% Right angle at E
\tkzMarkRightAngle[size=0.2](Gamma,E,A)
\end{tikzpicture}
\end{center}
\end{alist}

\vspace{4mm}

% --- Άσκηση 37098 (37098.tex) ---
\Askhsh{37098}{4}
\begin{ylh}{1}
    \item 4.2. Τέμνουσα δύο ευθειών - Ευκλείδειο αίτημα
    \item 4.6. Άθροισμα γωνιών τριγώνου
    \item 5.2. Παραλληλόγραμμα
    \item 5.3. Ορθογώνιο
    \item 5.6. Εφαρμογές στα τρίγωνα
    \item 5.10. Τραπέζιο.
\end{ylh}
% ===================================================================
%  Αρχείο: 37098.tex (Fragment)
%  Αυτόματη μετατροπή μέσω doc2latex.py
% ===================================================================



Σε μια ευθεία (ε) θεωρούμε διαδοχικά τα σημεία $A$, Β, Γ έτσι ώστε $AB$ = 2$\cdot$$B\varGamma $ και στο ίδιο ημιεπίπεδο θεωρούμε ισόπλευρα τρίγωνα Α$B\varDelta  $ και $B\varGamma $Ε. Αν Η είναι το μέσο του $A\varDelta  $ και η ευθεία $\varDelta  E$ τέμνει την ευθεία (ε) στο σημείο $Z$ να αποδείξετε ότι:

\begin{alist}
\item Το τετράπλευρο $BH\varDelta  E$ είναι ορθογώνιο. \monades{8}

\item Το τρίγωνο $\varGamma ZE$ είναι ισοσκελές. \monades{8}

\item Το τετράπλευρο ΗΕ$\varGamma  A$ είναι ισοσκελές τραπέζιο. \monades{9}

\begin{center}
\begin{tikzpicture}[scale=0.8]
% Line (ε) with A, B, Γ. $AB$ = 2*$B\varGamma $
\tkzDefPoint(0,0){A}
\tkzDefPoint(4,0){B}
\tkzDefPoint(6,0){Gamma}

% Equilateral triangle $AB\varDelta  $
\tkzDefTriangle[equilateral](A,B) \tkzGetPoint{Delta}

% Equilateral triangle $B\varGamma E$
\tkzDefTriangle[equilateral](B,Gamma) \tkzGetPoint{E}

% H is midpoint of $A\varDelta  $
\tkzDefMidPoint(A,Delta) \tkzGetPoint{H}

% $\varDelta  E$ intersects (ε) at Z. (ε) is A-Γ
\tkzInterLL(Delta,E)(A,Gamma) \tkzGetPoint{Z}

% Draw line (ε)
\draw[l3] ($(A)-(1,0)$)--($(Z)+(1,0)$) node[right]{$\varepsilon$};

% Draw triangles
\draw[l3] (A)--(B)--(Delta)--cycle;
\draw[l3] (B)--(Gamma)--(E)--cycle;

% Draw line $\varDelta  EZ$
\draw[l3, dashed] (Delta)--(Z);

% Draw H
\tkzDrawPoints(A,B,Gamma,Delta,E,H,Z)
\tkzLabelPoint[below](A){$A$}
\tkzLabelPoint[below](B){$B$}
\tkzLabelPoint[below](Gamma){$\varGamma $}
\tkzLabelPoint[above](Delta){$\varDelta  $}
\tkzLabelPoint[above](E){$ E$}
\tkzLabelPoint[left](H){$ H$}
\tkzLabelPoint[below](Z){$ Z$}
\end{tikzpicture}
\end{center}
\end{alist}

\vspace{4mm}

% --- Άσκηση 37099 (37099.tex) ---
\Askhsh{37099}{4}
\begin{ylh}{1}
    \item 4.4. Γωνίες με πλευρές παράλληλες
    \item 5.2. Παραλληλόγραμμα
    \item 5.9. Μια ιδιότητα του ορθογώνιου τριγώνου
    \item 5.10. Τραπέζιο.
\end{ylh}
% ===================================================================
%  Αρχείο: 37099.tex (Fragment)
%  Αυτόματη μετατροπή μέσω doc2latex.py
% ===================================================================



Δίνεται παραλληλόγραμμο $AB\varGamma \varDelta  $ και Κ το σημείο τομής των διαγωνίων του. Φέρουμε $AH$ κάθετη στην $B\varDelta  $ και στην προέκταση της $AH$ (προς το Η) θεωρούμε σημείο $E$ τέτοιο ώστε:

$AH$ = $HE$. Να αποδείξετε ότι:

\begin{alist}
\item Το τρίγωνο $AKE$ είναι ισοσκελές. \monades{7}

\item Το τρίγωνο $AE\varGamma $ είναι ορθογώνιο. \monades{9}

\item Το τετράπλευρο Δ$B\varGamma $Ε είναι ισοσκελές τραπέζιο. \monades{9}

\begin{center}
\begin{tikzpicture}[scale=0.9]
% Parallelogram $AB\varGamma \varDelta  $
\tkzDefPoint(0,0){A}
\tkzDefPoint(4,0){B}
\tkzDefPoint(5,3){Gamma}
\tkzDefPoint(1,3){Delta}

% Center K
\tkzInterLL(A,Gamma)(B,Delta) \tkzGetPoint{K}

% $AH$ perpendicular to $B\varDelta  $
\tkzDefPointBy[projection=onto B--Delta](A) \tkzGetPoint{H}

% E on extension of $AH$: $AH$ = $HE$ -> E is reflection of A over H
\tkzDefPointBy[symmetry=center H](A) \tkzGetPoint{E}

% Draw parallelogram
\draw[l3] (A)--(B)--(Gamma)--(Delta)--cycle;

% Draw diagonals
\draw[l3, dashed] (A)--(Gamma);
\draw[l3, dashed] (B)--(Delta);

% Draw $AE$
\draw[l3] (A)--(E);
\draw[l3, dashed] (E)--(Gamma); % Triangle $AE\varGamma $ mention
\draw[l3, dashed] (E)--(K);     % Triangle $AKE$ mention

% Right angle at H
\tkzMarkRightAngle[size=0.2](A,H,B)

\tkzDrawPoints(A,B,Gamma,Delta,K,H,E)
\tkzLabelPoint[below left](A){$A$}
\tkzLabelPoint[below right](B){$B$}
\tkzLabelPoint[above right](Gamma){$\varGamma $}
\tkzLabelPoint[above left](Delta){$\varDelta  $}
\tkzLabelPoint[above right](K){$ K$}
\tkzLabelPoint[below right](H){$ H$}
\tkzLabelPoint[below](E){$ E$}
\end{tikzpicture}
\end{center}
\end{alist}

\vspace{4mm}

% --- Άσκηση 37100 (37100.tex) ---
\Askhsh{37100}{4}
\begin{ylh}{1}
    \item 4.6. Άθροισμα γωνιών τριγώνου
    \item 5.9. Μια ιδιότητα του ορθογώνιου τριγώνου.
\end{ylh}
% ===================================================================
%  Αρχείο: 37100.tex (Fragment)
%  Αυτόματη μετατροπή μέσω doc2latex.py
% ===================================================================



Δίνεται ορθογώνιο τρίγωνο $AB\varGamma $ με τη γωνία $\hat{A}$ ορθή και $\hat{B} = 2\hat{\varGamma }$. Φέρουμε το ύψος του $A\varDelta  $ και σημείο $E$ στην προέκταση της $AB$ τέτοιο ώστε $BE$ = $B\varDelta  $.

\begin{alist}
\item Να υπολογίσετε τις γωνίες του τριγώνου $B\varDelta  $Ε. \monades{9}

\item Να αποδείξετε ότι:


\end{alist}
\begin{alist}
\item
  $ΒE = \ \frac{AB}{2}$. \monades{8}
\item
  $AE$ = $\varGamma \varDelta  $. \monades{8}

\begin{center}
\begin{tikzpicture}[scale=1.2]
% Right triangle $AB\varGamma $: Â=90°, B̂=60°, Γ̂=30°
\tkzDefPoint(0,0){A}
\tkzDefPoint(2,0){B}
\tkzDefPoint(0,3.464){Gamma}

% Altitude $A\varDelta  $
\tkzDefPoint(1.5,0.866){Delta}

% E on extension of $AB$: $BE$=$B\varDelta  $=1
\tkzDefPoint(3,0){E}

\draw[l3] (A)--(Gamma)--(B)--cycle;
\draw[l3] (A)--(Delta);
\draw[l3] (B)--(E);
\draw[l3, dashed] (Delta)--(E);

\tkzDrawPoints(A,B,Gamma,Delta,E)
\tkzLabelPoint[below left](A){$A$}
\tkzLabelPoint[below](B){$B$}
\tkzLabelPoint[above](Gamma){$\varGamma $}
\tkzLabelPoint[above right](Delta){$\varDelta  $}
\tkzLabelPoint[below](E){$ E$}

\tkzMarkRightAngle[size=0.2](Gamma,A,B)
\tkzMarkRightAngle[size=0.2](A,Delta,B)

\tkzMarkSegment[mark=||](B,Delta)
\tkzMarkSegment[mark=||](B,E)
\end{tikzpicture}
\end{center}
\end{alist}

\vspace{4mm}

% --- Άσκηση 37101 (37101.tex) ---
\Askhsh{37101}{4}
\begin{ylh}{1}
    \item 3.3. 2ο Κριτήριο ισότητας τριγώνων
    \item 4.2. Τέμνουσα δύο ευθειών - Ευκλείδειο αίτημα
    \item 5.2. Παραλληλόγραμμα
    \item 5.3. Ορθογώνιο
    \item 5.6. Εφαρμογές στα τρίγωνα.
\end{ylh}
% ===================================================================
%  Αρχείο: 37101.tex (Fragment)
%  Αυτόματη μετατροπή μέσω doc2latex.py
% ===================================================================



Δίνεται τρίγωνο $AB\varGamma $ με γωνίες $B$ και $\varGamma$ οξείες και Δ, Μ και $E$ τα μέσα των πλευρών του $AB$, $A\varGamma $ και $B\varGamma $ αντίστοιχα. Στις μεσοκάθετες των $AB$ και $B\varGamma $ και εκτός του τριγώνου $AB\varGamma $ θεωρούμε σημεία $Z$ και Η αντίστοιχα, τέτοια ώστε $\varDelta   Z = \ \frac{AB}{2}$ και $ΕH = \ \frac{Β\varGamma }{2}$.

\begin{alist}
\item Να αποδείξετε ότι:


\end{alist}
\begin{alist}
\item
  Το τετράπλευρο $B\varDelta  $ΜΕ είναι παραλληλόγραμμο. \monades{5}
\item
  Τα τρίγωνα $Z\varDelta  M$ και $EMH$ είναι ίσα. \monades{10}

\item Αν τα σημεία $Z$, Δ, Ε είναι συνευθειακά, να αποδείξετε ότι η γωνία $\hat{A} = \ 90\degree$. \monades{10}

\begin{center}
\begin{tikzpicture}[scale=0.8]
\tkzDefPoint(2,3){A}
\tkzDefPoint(0,0){B}
\tkzDefPoint(5,0){Gamma}

% Midpoints
\tkzDefMidPoint(A,B) \tkzGetPoint{Delta}
\tkzDefMidPoint(B,Gamma) \tkzGetPoint{E}
\tkzDefMidPoint(A,Gamma) \tkzGetPoint{M}

% Z on perp bisector of $AB$, outside triangle, dist = $AB$/2. Means $AZB$ is rt iso
\tkzDefSquare(A,B) \tkzGetPoints{sq1_1}{sq1_2}
\tkzInterLL(A,sq1_1)(B,sq1_2) \tkzGetPoint{Z}

% H on perp bisector of $B\varGamma $, outside, dist = $B\varGamma $/2. Means $B\varGamma H$ is rt iso
\tkzDefSquare(B,Gamma) \tkzGetPoints{sq2_1}{sq2_2}
\tkzInterLL(B,sq2_1)(Gamma,sq2_2) \tkzGetPoint{H}

\draw[l3] (A)--(B)--(Gamma)--cycle;
\draw[l3, dashed] (Delta)--(Z);
\draw[l3, dashed] (E)--(H);
\draw[l3] (A)--(Z)--(B);
\draw[l3] (B)--(H)--(Gamma);

% Segments connecting Δ, M, Z for triangle $Z\varDelta  M$ and $EMH$ mentioned
\draw[l3, dashed] (Z)--(M);
\draw[l3, dashed] (M)--(H);

\tkzDrawPoints(A,B,Gamma,Delta,M,E,Z,H)
\tkzLabelPoint[above](A){$A$}
\tkzLabelPoint[below left](B){$B$}
\tkzLabelPoint[below right](Gamma){$\varGamma $}
\tkzLabelPoint[above left](Delta){$\varDelta  $}
\tkzLabelPoint[above right](M){$\mathrm{M}$}
\tkzLabelPoint[below](E){$ E$}
\tkzLabelPoint[left](Z){$ Z$}
\tkzLabelPoint[below](H){$ H$}
\end{tikzpicture}
\end{center}
\end{alist}

\vspace{4mm}

% --- Άσκηση 37102 (37102.tex) ---
\Askhsh{37102}{4}
\begin{ylh}{1}
    \item 4.6. Άθροισμα γωνιών τριγώνου
    \item 5.2. Παραλληλόγραμμα.
\end{ylh}
% ===================================================================
%  Αρχείο: 37102.tex (Fragment)
%  Αυτόματη μετατροπή μέσω doc2latex.py
% ===================================================================



Δίνεται ορθογώνιο τρίγωνο $AB\varGamma $ ($\hat{A} = 90\degree$). Φέρουμε τη διάμεσο του $AM$ την οποία προεκτείνουμε (προς το μέρος του Μ) κατά τμήμα $M\varDelta  $ = $AM$. Θεωρούμε ευθεία $\varDelta  K$ κάθετη στη $B\varGamma $, η οποία τέμνει τη διχοτόμο της γωνίας $\hat{B}$ στο Ε. Να αποδείξετε ότι:

\begin{alist}
\item Το τετράπλευρο Α$B\varDelta  $Γ είναι ορθογώνιο \monades{8}

\item $Κ\hat{E}B = \ 90\degree - \ \frac{\hat{B}}{2}$ . \monades{8}

\item $\varDelta  E$ = $B\varDelta  $ \monades{9}

\begin{center}
\begin{tikzpicture}[scale=0.8]
% Right triangle A=90
\tkzDefPoint(0,0){A}
\tkzDefPoint(0,3){B}
\tkzDefPoint(4,0){Gamma}

% M midpoint $B\varGamma $
\tkzDefMidPoint(B,Gamma) \tkzGetPoint{M}

% Δ reflection of A over M ($M\varDelta  $=$AM$)
\tkzDefPointBy[symmetry=center M](A) \tkzGetPoint{Delta}

% $\varDelta  K$ perpendicular to $B\varGamma $
\tkzDefPointBy[projection=onto B--Gamma](Delta) \tkzGetPoint{K}

% Bisector of B: $BE$
\tkzDefLine[bisector](A,B,Gamma) \tkzGetPoint{bis_B}

% E is intersection of $\varDelta  K$ and bisector
\tkzInterLL(Delta,K)(B,bis_B) \tkzGetPoint{E}

% Draw rectangle $AB\varDelta  \varGamma $
\draw[l3] (A)--(Gamma)--(Delta)--(B)--cycle;

% Draw $A\varDelta  $
\draw[l3] (A)--(Delta);

% Draw bisector $BE$
\draw[l3] (B)--(E);

% Draw $\varDelta  E$ (extension of $\varDelta  K$)
\draw[l3] (Delta)--(E);

\tkzMarkRightAngle[size=0.2](B,A,Gamma)
\tkzMarkRightAngle[size=0.2](Delta,K,Gamma)

\tkzDrawPoints(A,B,Gamma,M,Delta,K,E)
\tkzLabelPoint[below left](A){$A$}
\tkzLabelPoint[above left](B){$B$}
\tkzLabelPoint[below right](Gamma){$\varGamma $}
\tkzLabelPoint[above left](M){$\mathrm{M}$}
\tkzLabelPoint[above right](Delta){$\varDelta  $}
\tkzLabelPoint[below right](K){$ K$}
\tkzLabelPoint[right](E){$ E$}
\end{tikzpicture}
\end{center}
\end{alist}

\vspace{4mm}

% --- Άσκηση 37103 (37103.tex) ---
\Askhsh{37103}{4}
\begin{ylh}{1}
    \item 3.2. 1ο Κριτήριο ισότητας τριγώνων
    \item 4.2. Τέμνουσα δύο ευθειών - Ευκλείδειο αίτημα
    \item 4.6. Άθροισμα γωνιών τριγώνου
    \item 5.2. Παραλληλόγραμμα
    \item 5.10. Τραπέζιο.
\end{ylh}
% ===================================================================
%  Αρχείο: 37103.tex (Fragment)
%  Αυτόματη μετατροπή μέσω doc2latex.py
% ===================================================================



Στο παρακάτω τετράπλευρο $AB\varGamma \varDelta  $ ισχύουν: $A\varDelta  $ = $B\varGamma $, $A\varGamma $ = $B\varDelta  $ και $AB$ < $\varGamma \varDelta  $.

\begin{alist}
\item Να αποδείξετε ότι τα τρίγωνα $AOB$ και $\varDelta  O\varGamma $ είναι ισοσκελή. \monades{9}

\item Να αποδείξετε ότι Δ$\hat{A}$B =$\ Α\hat{B}\varGamma $. \monades{8}

\item Να αποδείξετε ότι το τετράπλευρο $AB\varGamma \varDelta  $ είναι ισοσκελές τραπέζιο. \monades{8}

\begin{center}
\begin{tikzpicture}[scale=0.9]
% Isosceles trapezoid $AB\varGamma \varDelta  $
\tkzDefPoint(1,2.5){A}
\tkzDefPoint(4,2.5){B}
\tkzDefPoint(0,0){Delta}
\tkzDefPoint(5,0){Gamma}

% Diagonals intersection O
\tkzInterLL(A,Gamma)(B,Delta) \tkzGetPoint{O}

% Draw trapezoid
\draw[l3] (A)--(B)--(Gamma)--(Delta)--cycle;
\draw[l3] (A)--(Gamma);
\draw[l3] (B)--(Delta);

\tkzDrawPoints(A,B,Gamma,Delta,O)
\tkzLabelPoint[above left](A){$A$}
\tkzLabelPoint[above right](B){$B$}
\tkzLabelPoint[below right](Gamma){$\varGamma $}
\tkzLabelPoint[below left](Delta){$\varDelta  $}
\tkzLabelPoint[above](O){$ O$}
\end{tikzpicture}
\end{center}
\end{alist}

\vspace{4mm}

% --- Άσκηση 37104 (37104.tex) ---
\Askhsh{37104}{4}
\begin{ylh}{1}
    \item 4.6. Άθροισμα γωνιών τριγώνου
    \item 5.8. Το ορθόκεντρο τριγώνου
    \item 5.9. Μια ιδιότητα του ορθογώνιου τριγώνου.
\end{ylh}
% ===================================================================
%  Αρχείο: 37104.tex (Fragment)
%  Αυτόματη μετατροπή μέσω doc2latex.py
% ===================================================================



Δίνεται ορθογώνιο τρίγωνο $AB\varGamma $ ($\hat{A} = 90\degree$) και $\hat{\varGamma } = 30\degree$ με $M$ και $N$ τα μέσα των πλευρών $B\varGamma $ και $AB$ αντίστοιχα. Έστω ότι η μεσοκάθετος της πλευράς $B\varGamma $ τέμνει την $A\varGamma $ στο σημείο $E$.

\begin{alist}
\item Να αποδείξετε ότι:

i) η $BE$ είναι διχοτόμος της γωνίας $\hat{B}$. \monades{6}

ii) $ΑΕ = \ \frac{\varGamma  E}{2}$. \monades{6}

iii) η $BE$ είναι μεσοκάθετος της διαμέσου $AM$. \monades{7}

\item Αν $A\varDelta  $ είναι το ύψος του τριγώνου $AB\varGamma $ που τέμνει την $BE$ στο Η, να αποδείξετε ότι τα σημεία $M$, Η και $N$ είναι συνευθειακά. \monades{6}

\begin{center}
\begin{tikzpicture}[scale=1]
\tkzDefPoint(0,0){A}
\tkzDefPoint(0,2){B}
\tkzDefPoint(3.464,0){Gamma} % Γ=30°

\tkzDefMidPoint(B,Gamma) \tkzGetPoint{M}
\tkzDefMidPoint(A,B) \tkzGetPoint{N}

% Perp bisector of $B\varGamma $
\tkzDefLine[orthogonal=through M](B,Gamma) \tkzGetPoint{pM}
\tkzInterLL(M,pM)(A,Gamma) \tkzGetPoint{E}

% Altitude $A\varDelta  $
\tkzDefPointBy[projection=onto B--Gamma](A) \tkzGetPoint{Delta}

% $BE$ intersects $A\varDelta  $ at H
\tkzInterLL(A,Delta)(B,E) \tkzGetPoint{H}

\draw[l3] (A)--(B)--(Gamma)--cycle;
\draw[l3] (M)--(E);
\draw[l3] (B)--(E);
\draw[l3] (A)--(Delta);
\draw[l3] (A)--(M);

% dashed for collinear M, H, N (question ii)
\draw[l3, dashed] (M)--(N);

\tkzDrawPoints(A,B,Gamma,M,N,E,Delta,H)
\tkzLabelPoint[below left](A){$A$}
\tkzLabelPoint[above](B){$B$}
\tkzLabelPoint[below right](Gamma){$\varGamma $}
\tkzLabelPoint[above right](M){$\mathrm{M}$}
\tkzLabelPoint[left](N){$\mathrm{N}$}
\tkzLabelPoint[below](E){$ E$}
\tkzLabelPoint[below left](Delta){$\varDelta  $}
\tkzLabelPoint[above right](H){$ H$}

\tkzMarkRightAngle[size=0.2](B,A,Gamma)
\tkzMarkRightAngle[size=0.2](A,Delta,B)
\tkzMarkRightAngle[size=0.2](E,M,B)
\end{tikzpicture}
\end{center}
\end{alist}

\vspace{4mm}

% --- Άσκηση 37106 (37106.tex) ---
\Askhsh{37106}{4}
\begin{ylh}{1}
    \item 4.2. Τέμνουσα δύο ευθειών - Ευκλείδειο αίτημα
    \item 5.6. Εφαρμογές στα τρίγωνα.
\end{ylh}
% ===================================================================
%  Αρχείο: 37106.tex (Fragment)
%  Αυτόματη μετατροπή μέσω doc2latex.py
% ===================================================================



Δίνεται τρίγωνο $AB\varGamma $ με $AB$ < $B\varGamma $ και η διχοτόμος $BE$ της γωνίας $\hat{B}$. Αν $AZ$ $\bot$ $BE$, όπου $Z$ σημείο της $B\varGamma $ και $M$ το μέσον της $A\varGamma $, να αποδείξετε ότι :

\begin{alist}
\item Το τρίγωνο $ABZ$ είναι ισοσκελές. \monades{7}

\item $\varDelta  M$ // $B\varGamma $ και $\varDelta Μ = \ \frac{B\varGamma-AB}{2}$. \monades{10}

\item $E\hat{\varDelta  }Μ = \ \frac{\hat{B}}{2}$, όπου $\hat{B}$ η γωνία του τριγώνου $AB\varGamma $. \monades{8}

\begin{center}
\begin{tikzpicture}[scale=0.9]
\tkzDefPoint(1,3){A}
\tkzDefPoint(0,0){B}
\tkzDefPoint(6,0){Gamma}

\tkzDefLine[bisector](Gamma,B,A) \tkzGetPoint{$BE$}

% $AZ$ perp to Bx => Δ is intersection
\tkzDefPointBy[projection=onto B--$BE$](A) \tkzGetPoint{Delta}

% Ζ on $B\varGamma $
\tkzInterLL(A,Delta)(B,Gamma) \tkzGetPoint{Z}

% E is intersection of $BE$ with $A\varGamma $
\tkzInterLL(B,$BE$)(A,Gamma) \tkzGetPoint{E}

\tkzDefMidPoint(A,Gamma) \tkzGetPoint{M}

\draw[l3] (A)--(B)--(Gamma)--cycle;
\draw[l3] (B)--(E);
\draw[l3] (A)--(Z);
\draw[l3, dashed] (Delta)--(M);

\tkzDrawPoints(A,B,Gamma,Z,E,M,Delta)
\tkzLabelPoint[above](A){$A$}
\tkzLabelPoint[below left](B){$B$}
\tkzLabelPoint[below right](Gamma){$\varGamma $}
\tkzLabelPoint[below](Z){$ Z$}
\tkzLabelPoint[above right](E){$ E$}
\tkzLabelPoint[above right](M){$\mathrm{M}$}
\tkzLabelPoint[above left](Delta){$\varDelta  $}

\tkzMarkRightAngle[size=0.2](A,Delta,B)
\tkzMarkAngle[size=0.4,mark=s|](Gamma,B,E)
\tkzMarkAngle[size=0.45,mark=s|](E,B,A)
\end{tikzpicture}
\end{center}
\end{alist}

\vspace{4mm}

% --- Άσκηση 37107 (37107.tex) ---
\Askhsh{37107}{4}
\begin{ylh}{1}
    \item 4.4. Γωνίες με πλευρές παράλληλες
    \item 4.6. Άθροισμα γωνιών τριγώνου
    \item 5.3. Ορθογώνιο
    \item 5.6. Εφαρμογές στα τρίγωνα
    \item 5.11. Ισοσκελές τραπέζιο.
\end{ylh}
% ===================================================================
%  Αρχείο: 37107.tex (Fragment)
%  Αυτόματη μετατροπή μέσω doc2latex.py
% ===================================================================



Στο παρακάτω σχήμα δίνεται τρίγωνο $AB\varGamma $, η διχοτόμος Βx της γωνίας $\hat{B}$ του τριγώνου $AB\varGamma $ και η διχοτόμος Βy της εξωτερικής γωνίας $\hat{B}$. Αν $\varDelta$ και $E$ είναι οι προβολές της κορυφής Α του τριγώνου $AB\varGamma $ στην Βx και Βy αντίστοιχα, να αποδείξετε ότι:

\begin{alist}
\item Το τετράπλευρο ΑΔ$BE$ είναι ορθογώνιο. \monades{7}

\item Η ευθεία $E\varDelta  $ είναι παράλληλη προς τη $B\varGamma $ και διέρχεται από το μέσο $M$ της $A\varGamma $.

\monades{10}

\item Το τετράπλευρο $KM\varGamma B$ είναι τραπέζιο και η διάμεσός του είναι ίση με $\frac{3 \cdot \alpha}{4}$, όπου α=$B\varGamma $. \monades{8}

\begin{center}
\begin{tikzpicture}[scale=0.9]
\tkzDefPoint(1,3){A}
\tkzDefPoint(3,0){B}
\tkzDefPoint(7,0){Gamma}

% Ext angle at B. Let's define a point on Gamma-B extension
\tkzDefPoint(-1,0){ExtB}

\tkzDefLine[bisector](A,B,Gamma) \tkzGetPoint{Bx}
\tkzDefLine[bisector](ExtB,B,A) \tkzGetPoint{By}

% Δ projection of A to Bx
\tkzDefPointBy[projection=onto B--Bx](A) \tkzGetPoint{Delta}

% E projection of A to By
\tkzDefPointBy[projection=onto B--By](A) \tkzGetPoint{E}

\tkzDefMidPoint(A,Gamma) \tkzGetPoint{M}
\tkzInterLL(E,Delta)(A,B) \tkzGetPoint{K}

\draw[l3] (A)--(B)--(Gamma)--cycle;
\draw[l3] (B)--(ExtB); % Extension line
\draw[l3, dashed] (B)--(Delta); 
\draw[l3, dashed] (B)--(E);
\draw[l3, dashed] (A)--(Delta);
\draw[l3, dashed] (A)--(E);
\draw[l3] (E)--(M);

\tkzDrawPoints(A,B,Gamma,Delta,E,M,K)
\tkzLabelPoint[above](A){$A$}
\tkzLabelPoint[below](B){$B$}
\tkzLabelPoint[below right](Gamma){$\varGamma $}
\tkzLabelPoint[below](Delta){$\varDelta  $}
\tkzLabelPoint[left](E){$ E$}
\tkzLabelPoint[above right](M){$\mathrm{M}$}
\tkzLabelPoint[above left](K){$ K$}

\tkzMarkRightAngle[size=0.2](A,Delta,B)
\tkzMarkRightAngle[size=0.2](A,E,B)
\end{tikzpicture}
\end{center}
\end{alist}

\vspace{4mm}

% --- Άσκηση 37108 (37108.tex) ---
\Askhsh{37108}{4}
\begin{ylh}{1}
    \item 3.2. 1ο Κριτήριο ισότητας τριγώνων
    \item 5.2. Παραλληλόγραμμα.
\end{ylh}
% ===================================================================
%  Αρχείο: 37108.tex (Fragment)
%  Αυτόματη μετατροπή μέσω doc2latex.py
% ===================================================================



Σε παραλληλόγραμμο $AB\varGamma \varDelta  $ θεωρούμε σημεία $E$, Ζ, Η, Θ στις πλευρές $AB$, $B\varGamma $, $\varGamma \varDelta  $, $\varDelta  A$ αντίστοιχα, με $AE$ = $\varGamma H$ και $BZ$ = $\varDelta  \varTheta $. Να αποδείξετε ότι:

\begin{alist}
\item Το τετράπλευρο $AE\varGamma H$ είναι παραλληλόγραμμο. (6 μονάδες)

\item Το τετράπλευρο $EZH\varTheta $ είναι παραλληλόγραμμο. (10 μονάδες)

\item Τα τμήματα $A\varGamma $, $B\varDelta  $, $EH$ και $Z\varTheta $ διέρχονται από το ίδιο σημείο. (9 μονάδες)

\begin{center}
\begin{tikzpicture}[scale=0.9]
\tkzDefPoint(0,0){A}
\tkzDefPoint(4,0){B}
\tkzDefPoint(5,3){Gamma}
\tkzDefPoint(1,3){Delta}

\tkzDefPointOnLine[pos=0.3](A,B) \tkzGetPoint{E}
\tkzDefPointOnLine[pos=0.3](Gamma,Delta) \tkzGetPoint{H}

\tkzDefPointOnLine[pos=0.4](B,Gamma) \tkzGetPoint{Z}
\tkzDefPointOnLine[pos=0.4](Delta,A) \tkzGetPoint{Theta}

\draw[l3] (A)--(B)--(Gamma)--(Delta)--cycle;
\draw[l3, dashed] (A)--(Gamma);
\draw[l3, dashed] (B)--(Delta);
\draw[l3, dashed] (E)--(H);
\draw[l3, dashed] (Z)--(Theta);
\draw[l3] (E)--(Z)--(H)--(Theta)--cycle;

\tkzDrawPoints(A,B,Gamma,Delta,E,Z,H,Theta)
\tkzLabelPoint[below left](A){$A$}
\tkzLabelPoint[below right](B){$B$}
\tkzLabelPoint[above right](Gamma){$\varGamma $}
\tkzLabelPoint[above left](Delta){$\varDelta  $}
\tkzLabelPoint[below](E){$ E$}
\tkzLabelPoint[right](Z){$ Z$}
\tkzLabelPoint[above](H){$ H$}
\tkzLabelPoint[left](Theta){$\varTheta $}
\end{tikzpicture}
\end{center}
\end{alist}

\vspace{4mm}

% --- Άσκηση 37109 (37109.tex) ---
\Askhsh{37109}{4}
\begin{ylh}{1}
    \item 5.2. Παραλληλόγραμμα
    \item 5.3. Ορθογώνιο
    \item 5.4. Ρόμβος.
\end{ylh}
% ===================================================================
%  Αρχείο: 37109.tex (Fragment)
%  Αυτόματη μετατροπή μέσω doc2latex.py
% ===================================================================



Δίνεται παραλληλόγραμμο $AB\varGamma \varDelta  $ και σημεία $K$, Λ της διαγωνίου του $B\varDelta  $, τέτοια ώστε να ισχύει $BK$ = $K\varLambda $ = $\varLambda \varDelta  $ .

\begin{alist}
\item Να αποδείξετε ότι το τετράπλευρο $AK\varGamma \varLambda $ είναι παραλληλόγραμμο. \monades{10}

\item Να αποδείξετε ότι, αν το αρχικό παραλληλόγραμμο $AB\varGamma \varDelta  $ είναι ρόμβος, τότε και το $AK\varGamma \varLambda $ είναι ρόμβος. \monades{8}

\item Ποια πρέπει να είναι η σχέση των διαγωνίων του αρχικού παραλληλογράμμου $AN\varGamma \varDelta  $, ώστε το $AK\varGamma \varLambda $ να είναι ορθογώνιο. Να αιτιολογήσετε την απάντησή σας. \monades{7}
\end{alist}

\vspace{4mm}

% --- Άσκηση 37114 (37114.tex) ---
\Askhsh{37114}{4}
\begin{ylh}{1}
    \item 5.2. Παραλληλόγραμμα
    \item 5.6. Εφαρμογές στα τρίγωνα
    \item 5.10. Τραπέζιο.
\end{ylh}
% ===================================================================
%  Αρχείο: 37114.tex (Fragment)
%  Αυτόματη μετατροπή μέσω doc2latex.py
% ===================================================================



Δίνεται παραλληλόγραμμο $AB\varGamma \varDelta  $ με $AB$ < $A\varDelta  $ και έστω Ο το σημείο τομής των διαγωνίων του $A\varGamma $ και $B\varDelta  $. Φέρνουμε την $AE$ κάθετη στην διαγώνιο $B\varDelta  $. Αν το $Z$ είναι το συμμετρικό του Α ως προς την διαγώνιο $B\varDelta  $ και δεν συμπίπτει με το σημείο $\varGamma $, τότε να αποδείξετε ότι:

\begin{alist}
\item Το τρίγωνο Α$\varDelta   Z$ είναι ισοσκελές. \monades{7}

\item $Z\varGamma $ = 2ΟΕ. \monades{9}

\item Το Β$\varDelta   Z$Γ είναι ισοσκελές τραπέζιο. \monades{9}

\begin{center}
\begin{tikzpicture}[scale=0.9]
% Parallelogram $AB\varGamma \varDelta  $ with $AB$ < $A\varDelta  $
\tkzDefPoint(1,3){A}
\tkzDefPoint(0,0){B}
\tkzDefPoint(5,0){Gamma}
\tkzDefPoint(6,3){Delta}

% O intersection of diagonals
\tkzInterLL(A,Gamma)(B,Delta) \tkzGetPoint{O}

% $AE$ perpendicular to $B\varDelta  $
\tkzDefPointBy[projection=onto B--Delta](A) \tkzGetPoint{E}

% Z symmetric of A with respect to $B\varDelta  $
\tkzDefPointBy[symmetry=center E](A) \tkzGetPoint{Z}

\draw[l3] (A)--(B)--(Gamma)--(Delta)--cycle;
\draw[l3, dashed] (A)--(Gamma);
\draw[l3, dashed] (B)--(Delta);

% Draw $AZ$ (perpendicular to $BD$)
\draw[l3] (A)--(Z);
% Draw sides for question
\draw[l3, dashed] (A)--(Delta);
\draw[l3, dashed] (Delta)--(Z);
\draw[l3, dashed] (Z)--(Gamma);
\draw[l3, dashed] (B)--(Z); 

\tkzMarkRightAngle[size=0.2](A,E,Delta)

\tkzDrawPoints(A,B,Gamma,Delta,O,E,Z)
\tkzLabelPoint[above left](A){$A$}
\tkzLabelPoint[below left](B){$B$}
\tkzLabelPoint[below right](Gamma){$\varGamma $}
\tkzLabelPoint[above right](Delta){$\varDelta  $}
\tkzLabelPoint[above](O){$ O$}
\tkzLabelPoint[below right](E){$ E$}
\tkzLabelPoint[below](Z){$ Z$}
\end{tikzpicture}
\end{center}
\end{alist}

\vspace{4mm}

% --- Άσκηση 37115 (37115.tex) ---
\Askhsh{37115}{4}
\begin{ylh}{1}
    \item 5.2. Παραλληλόγραμμα
    \item 5.6. Εφαρμογές στα τρίγωνα
    \item 5.10. Τραπέζιο.
\end{ylh}
% ===================================================================
%  Αρχείο: 37115.tex (Fragment)
%  Αυτόματη μετατροπή μέσω doc2latex.py
% ===================================================================



Δίνεται παραλληλόγραμμο $AB\varGamma \varDelta  $. Στην προέκταση της πλευράς $AB$ παίρνουμε τμήμα $BE$ = $AB$ και στην προέκταση της πλευράς $A\varDelta  $ τμήμα $\varDelta   Z$ = $A\varDelta  $.

\begin{alist}
\item Να αποδείξετε ότι:


\end{alist}
\begin{alist}
\item
  Τα τετράπλευρα $B\varDelta  $ΓΕ και Β$\varDelta   Z$Γ είναι παραλληλόγραμμα. \monades{7}
\item
  Τα σημεία $E$, Γ και $Z$ είναι συνευθειακά. \monades{9}

\item Αν Κ και Λ είναι τα μέσα των $BE$ και $\varDelta   Z$ αντίστοιχα, τότε $K\varLambda $ // $\varDelta   B$ και $Κ\Lambda = \frac{3}{2}\varDelta   B$. \monades{9}
\end{alist}

\vspace{4mm}

% --- Άσκηση 37116 (37116.tex) ---
\Askhsh{37116}{4}
\begin{ylh}{1}
    \item 3.6. Κριτήρια ισότητας ορθογώνιων τριγώνων
    \item 5.6. Εφαρμογές στα τρίγωνα
    \item 5.8. Το ορθόκεντρο τριγώνου.
\end{ylh}
% ===================================================================
%  Αρχείο: 37116.tex (Fragment)
%  Αυτόματη μετατροπή μέσω doc2latex.py
% ===================================================================



Δίνεται ισοσκελές τρίγωνο $AB\varGamma $ με $AB=A\varGamma $ και το ύψος του $AM$. Φέρουμε $M\varDelta  $ κάθετη στην $A\varGamma $ και θεωρούμε Η το μέσο του τμήματος $M\varDelta  $. Από το Η φέρουμε παράλληλη στη $B\varGamma $ η οποία τέμνει τις $AM$ και $A\varGamma $ στα σημεία $K$ και $Z$ αντίστοιχα.

Να αποδείξετε ότι:

\begin{alist}
\item $HΖ = \ \frac{B\varGamma }{4}$ \monades{9}

\item $MZ$ // $B\varDelta  $ \monades{8}

\item Η ευθεία $AH$ είναι κάθετη στη $B\varDelta  $. \monades{8}

\begin{center}
\begin{tikzpicture}[scale=0.9]
% Isosceles triangle A=(0,4), B=(-2,0), Γ=(2,0)
\tkzDefPoint(0,4){A}
\tkzDefPoint(-2,0){B}
\tkzDefPoint(2,0){Gamma}

% M altitude point on $B\varGamma $ -> M=(0,0)
\tkzDefPointBy[projection=onto B--Gamma](A) \tkzGetPoint{M}

% $M\varDelta  $ perpendicular to $A\varGamma $
\tkzDefPointBy[projection=onto A--Gamma](M) \tkzGetPoint{Delta}

% H midpoint of $M\varDelta  $
\tkzDefMidPoint(M,Delta) \tkzGetPoint{H}

% Parallel to $B\varGamma $ from H
\tkzDefLine[parallel=through H](B,Gamma) \tkzGetPoint{H_dir}

% K is intersection with $AM$, Z intersection with $A\varGamma $
\tkzInterLL(H,H_dir)(A,M) \tkzGetPoint{K}
\tkzInterLL(H,H_dir)(A,Gamma) \tkzGetPoint{Z}

\draw[l3] (A)--(B)--(Gamma)--cycle;
\draw[l3] (A)--(M);
\draw[l3] (M)--(Delta);
\draw[l3] (K)--(Z); % This is parallel to $BC$, passing through H
\draw[l3, dashed] (A)--(H);
\draw[l3, dashed] (B)--(Delta);
\draw[l3, dashed] (M)--(Z);

\tkzMarkRightAngle[size=0.2](A,M,Gamma)
\tkzMarkRightAngle[size=0.2](M,Delta,A)

\tkzDrawPoints(A,B,Gamma,M,Delta,H,K,Z)
\tkzLabelPoint[above](A){$A$}
\tkzLabelPoint[below left](B){$B$}
\tkzLabelPoint[below right](Gamma){$\varGamma $}
\tkzLabelPoint[below](M){$\mathrm{M}$}
\tkzLabelPoint[right](Delta){$\varDelta  $}
\tkzLabelPoint[below right](H){$ H$}
\tkzLabelPoint[left](K){$ K$}
\tkzLabelPoint[above right](Z){$ Z$}
\end{tikzpicture}
\end{center}
\end{alist}

\vspace{4mm}

% --- Άσκηση 37117 (37117.tex) ---
\Askhsh{37117}{4}
\begin{ylh}{1}
    \item 4.4. Γωνίες με πλευρές παράλληλες
    \item 4.6. Άθροισμα γωνιών τριγώνου
    \item 5.2. Παραλληλόγραμμα.
\end{ylh}
% ===================================================================
%  Αρχείο: 37117.tex (Fragment)
%  Αυτόματη μετατροπή μέσω doc2latex.py
% ===================================================================



Έστω τρίγωνο $AB\varGamma $ και $A\varDelta  $ η διχοτόμος της γωνίας Α, για την οποία ισχύει $A\varDelta  $ = $\varDelta  \varGamma $. Η $\varDelta  E$ είναι διχοτόμος της γωνίας Α$\hat{\varDelta  }$Β και η $\varDelta   Z$ παράλληλη στην $AB$.

Να αποδείξετε ότι:

\begin{alist}
\item Τα τμήματα $E\varDelta  $ και $A\varGamma $ είναι παράλληλα. \monades{9}

\item Το τρίγωνο $EA\varDelta  $ είναι ισοσκελές. \monades{8}

\item Τα τμήματα $A\varDelta  $ και $EZ$ διχοτομούνται. \monades{8}

\begin{center}
\begin{tikzpicture}[scale=1]
\tkzDefPoint(1,3){A}
\tkzDefPoint(0,0){B}
\tkzDefPoint(6,0){Gamma}

\tkzDefLine[bisector](B,A,Gamma) \tkzGetPoint{bis_A}
\tkzInterLL(A,bis_A)(B,Gamma) \tkzGetPoint{Delta}

\tkzDefLine[bisector](A,Delta,B) \tkzGetPoint{bis_D}
\tkzInterLL(Delta,bis_D)(A,B) \tkzGetPoint{E}

\tkzDefLine[parallel=through Delta](A,B) \tkzGetPoint{par_AB}
\tkzInterLL(Delta,par_AB)(A,Gamma) \tkzGetPoint{Z}

\draw[l3] (A)--(B)--(Gamma)--cycle;
\draw[l3] (A)--(Delta);
\draw[l3] (Delta)--(E);
\draw[l3] (Delta)--(Z);
\draw[l3, dashed] (E)--(Z);

\tkzDrawPoints(A,B,Gamma,Delta,E,Z)
\tkzLabelPoint[above](A){$A$}
\tkzLabelPoint[below left](B){$B$}
\tkzLabelPoint[below right](Gamma){$\varGamma $}
\tkzLabelPoint[below](Delta){$\varDelta  $}
\tkzLabelPoint[left](E){$ E$}
\tkzLabelPoint[right](Z){$ Z$}

\tkzMarkAngle[size=0.4,mark=s|](B,A,Delta)
\tkzMarkAngle[size=0.45,mark=s|](Delta,A,Gamma)
\tkzMarkSegment[mark=||](A,Delta)
\tkzMarkSegment[mark=||](Delta,Gamma)
\end{tikzpicture}
\end{center}
\end{alist}

\vspace{4mm}

% --- Άσκηση 37118 (37118.tex) ---
\Askhsh{37118}{4}
\begin{ylh}{1}
    \item 4.4. Γωνίες με πλευρές παράλληλες
    \item 4.6. Άθροισμα γωνιών τριγώνου
    \item 5.10. Τραπέζιο.
\end{ylh}
% ===================================================================
%  Αρχείο: 37118.tex (Fragment)
%  Αυτόματη μετατροπή μέσω doc2latex.py
% ===================================================================



Δίνεται το τρίγωνο $AB\varGamma $ με γωνία $\hat{B}\ $= 60\textsuperscript{ο}. Φέρνουμε τα ύψη $A\varDelta  $ και $\varGamma E$ που τέμνονται στο Η. Φέρνουμε $KZ$ διχοτόμο της γωνίας Ε$\hat{H}$Α και $\varTheta  H$ κάθετο στο ύψος $A\varDelta$.

Να αποδείξετε ότι:

\begin{alist}
\item Για το τμήμα $ZE$ ισχύει $ZH$ = 2ΕΖ. \monades{9}

\item Το τρίγωνο $\varTheta ZH$ είναι ισόπλευρο. \monades{8}

\item Το τετράπλευρο $\varTheta HKB$ είναι ισοσκελές τραπέζιο. \monades{8}

\begin{center}
\begin{tikzpicture}[scale=1.2]
%\tkzDefPoint(0,0){B}
%\tkzDefPoint(4,0){Gamma}
%% B=60°. 
%\tkzDefPoint(2,3.464){A} % equilateral
%
%\tkzDefPointBy[projection=onto B--Gamma](A) \tkzGetPoint{Delta}
%\tkzDefPointBy[projection=onto A--B](Gamma) \tkzGetPoint{E}
%\tkzInterLL(A,Delta)(Gamma,E) \tkzGetPoint{H}
%
%% $\varThetaH$ perp to $A\varDelta  $. Since $A\varDelta  $ is vertical, $\varThetaH$ is horizontal. Θ on $AB$
%\tkzDefPointBy[rotation=center H angle 90](Delta) \tkzGetPoint{dirTheta}
%\tkzInterLL(H,dirTheta)(A,B) \tkzGetPoint{Theta}
%
%% $KZ$ bisector of $EHA$
%\tkzDefLine[bisector](E,H,A) \tkzGetPoint{dirKZ}
%\tkzInterLL(H,dirKZ)(A,B) \tkzGetPoint{Z}
%\tkzInterLL(H,dirKZ)(A,Gamma) \tkzGetPoint{K}
%
%\draw[l3] (A)--(B)--(Gamma)--cycle;
%\draw[l3] (A)--(Delta);
%\draw[l3] (Gamma)--(E);
%\draw[l3, dashed] (Theta)--(H);
%\draw[l3, dashed] (Z)--(K);
%\draw[l3, dashed] (Theta)--(Z);
%
%\tkzMarkRightAngle[size=0.15](A,Delta,Gamma)
%\tkzMarkRightAngle[size=0.15](Gamma,E,A)
%
%\tkzDrawPoints(A,B,Gamma,Delta,E,H,Theta,Z,K)
%\tkzLabelPoint[above](A){$A$}
%\tkzLabelPoint[below left](B){$B$}
%\tkzLabelPoint[below right](Gamma){$\varGamma $}
%\tkzLabelPoint[below](Delta){$\Delta$}
%\tkzLabelPoint[above left](E){$E$}
%\tkzLabelPoint[right](H){$\mathrm{H}$}
%\tkzLabelPoint[left](Theta){$\Theta$}
%\tkzLabelPoint[left](Z){$\mathrm{Z}$}
%\tkzLabelPoint[right](K){$\mathrm{K}$}
\end{tikzpicture}
\end{center}
\end{alist}

\vspace{4mm}

% --- Άσκηση 37124 (37124.tex) ---
\Askhsh{37124}{4}
\begin{ylh}{1}
    \item 3.2. 1ο Κριτήριο ισότητας τριγώνων.
\end{ylh}
% ===================================================================
%  Αρχείο: 37124.tex (Fragment)
%  Αυτόματη μετατροπή μέσω doc2latex.py
% ===================================================================



Δίνεται τρίγωνο $AB\varGamma $ με $AB$ < $A\varGamma $. Στην προέκταση της $AB$ (προς το Β) θεωρούμε σημείο $E$ έτσι ώστε $AE$ = $A\varGamma $. Στην πλευρά $A\varGamma $ θεωρούμε σημείο $\varDelta  $ έτσι ώστε $A\varDelta  $ = $AB$. Αν τα τμήματα $\varDelta  E$ και $B\varGamma $ τέμνονται στο Κ και η προέκταση της $AK$ τέμνει την $E\varGamma $ στο Μ, τότε να αποδείξετε ότι:

\begin{alist}
\item $B\varGamma $ = $\varDelta  E$ \monades{6}

\item $BK$ = $K\varDelta  $ \monades{7}

\item Η $AK$ είναι διχοτόμος της γωνίας Α. \monades{6}

\item Η $AM$ είναι μεσοκάθετος της $E\varGamma $. \monades{6}

\begin{center}
\begin{tikzpicture}[scale=0.9]
\tkzDefPoint(0,3){A}
\tkzDefPoint(-1.5,0){B}
\tkzDefPoint(2,0){Gamma}

% E on extension of $AB$: $AE$ = $A\varGamma $
\tkzCalcLength[cm](A,Gamma) \tkzGetLength{rAG}
\tkzInterLC[R](A,B)(A,\rAG cm) \tkzGetPoints{E_temp}{E} % E on ray $AB$ past B

% Δ on $A\varGamma $: $A\varDelta  $ = $AB$
\tkzCalcLength[cm](A,B) \tkzGetLength{rAB}
\tkzInterLC[R](A,Gamma)(A,\rAB cm) \tkzGetPoints{Delta_temp}{Delta}

% K is intersection of $\varDelta  E$ and $B\varGamma $
\tkzInterLL(Delta,E)(B,Gamma) \tkzGetPoint{K}

% M is intersection of $AK$ and $E\varGamma $
\tkzInterLL(A,K)(E,Gamma) \tkzGetPoint{M}

\draw[l3] (A)--(B)--(Gamma)--cycle;
\draw[l3] (B)--(E);
\draw[l3] (Delta)--(E);
\draw[l3] (E)--(Gamma);
\draw[l3] (A)--(M);

\tkzDrawPoints(A,B,Gamma,E,Delta,K,M)
\tkzLabelPoint[above](A){$A$}
\tkzLabelPoint[left](B){$B$}
\tkzLabelPoint[right](Gamma){$\varGamma $}
\tkzLabelPoint[below left](E){$ E$}
\tkzLabelPoint[above right](Delta){$\varDelta  $}
\tkzLabelPoint[below](K){$ K$}
\tkzLabelPoint[below](M){$\mathrm{M}$}
\end{tikzpicture}
\end{center}
\end{alist}

\vspace{4mm}

% --- Άσκηση 37125 (37125.tex) ---
\Askhsh{37125}{4}
\begin{ylh}{1}
    \item 3.2. 1ο Κριτήριο ισότητας τριγώνων
    \item 3.6. Κριτήρια ισότητας ορθογώνιων τριγώνων.
\end{ylh}
% ===================================================================
%  Αρχείο: 37125.tex (Fragment)
%  Αυτόματη μετατροπή μέσω doc2latex.py
% ===================================================================



Στις πλευρές Αx΄ και Αx γωνίας $x΄\hat{A}x$ θεωρούμε σημεία $B$ και $\varGamma$ ώστε $AB=A\varGamma $. Οι κάθετες στις Αx΄ και Αx στα σημεία $B$ και $\varGamma$ αντίστοιχα, τέμνονται στο Δ.

Αν οι ημιευθείες Αy και Αz χωρίζουν τη γωνία $x΄\hat{A}x$ σε τρεις ίσες γωνίες και τέμνουν τις $B\varDelta  $ και $\varDelta  \varGamma $ στα σημεία $E$ και $Z$ αντίστοιχα, να αποδείξετε ότι:

\begin{alist}
\item Το τρίγωνο $EAZ$ είναι ισοσκελές. \monades{8}

\item Το $\varDelta$ ανήκει στη διχοτόμο της γωνίας $x΄\hat{A}x$. \monades{8}

\item Οι γωνίες Γ$B\varDelta  $ και $\varGamma  A$Δ είναι ίσες. \monades{9}

\begin{center}
\begin{tikzpicture}[scale=1]
\tkzDefPoint(0,0){A}
\tkzDefPoint(2,1.5){B}
\tkzDefPoint(2,-1.5){Gamma}

% Perpendiculars at B and Γ
\tkzDefLine[orthogonal=through B](A,B) \tkzGetPoint{pB}
\tkzDefLine[orthogonal=through Gamma](A,Gamma) \tkzGetPoint{pG}
\tkzInterLL(B,pB)(Gamma,pG) \tkzGetPoint{Delta}

% Ay, Az trisect angle
\tkzFindAngle(Gamma,A,B) \tkzGetAngle{ang}
\pgfmathsetmacro{\thirdang}{\ang/3}
\tkzDefPointBy[rotation=center A angle \thirdang](Gamma) \tkzGetPoint{dirZ}
\tkzDefPointBy[rotation=center A angle 2*\thirdang](Gamma) \tkzGetPoint{dirE}

\tkzInterLL(A,dirE)(B,Delta) \tkzGetPoint{E}
\tkzInterLL(A,dirZ)(Gamma,Delta) \tkzGetPoint{Z}

\draw[l3, ->] (A)--($(A)!1.5!(B)$) node[above]{$x'$};
\draw[l3, ->] (A)--($(A)!1.5!(Gamma)$) node[below]{$x$};

\draw[l3] (A)--(B);
\draw[l3] (A)--(Gamma);
\draw[l3] (B)--(Delta);
\draw[l3] (Gamma)--(Delta);
\draw[l3, dashed] (A)--(Delta);
\draw[l3, ->] (A)--($(A)!1.2!(E)$) node[above left]{$y$};
\draw[l3, ->] (A)--($(A)!1.2!(Z)$) node[below left]{$z$};

\tkzDrawPoints(A,B,Gamma,Delta,E,Z)
\tkzLabelPoint[left](A){$A$}
\tkzLabelPoint[above left](B){$B$}
\tkzLabelPoint[below left](Gamma){$\varGamma $}
\tkzLabelPoint[right](Delta){$\varDelta  $}
\tkzLabelPoint[above right](E){$ E$}
\tkzLabelPoint[below right](Z){$ Z$}

\tkzMarkRightAngle[size=0.2](A,B,Delta)
\tkzMarkRightAngle[size=0.2](A,Gamma,Delta)
\end{tikzpicture}
\end{center}
\end{alist}

\vspace{4mm}

% --- Άσκηση 37126 (37126.tex) ---
\Askhsh{37126}{4}
\begin{ylh}{1}
    \item 3.2. 1ο Κριτήριο ισότητας τριγώνων
    \item 3.3. 2ο Κριτήριο ισότητας τριγώνων
    \item 4.6. Άθροισμα γωνιών τριγώνου
    \item 5.3. Ορθογώνιο
    \item 5.5. Τετράγωνο
    \item 5.9. Μια ιδιότητα του ορθογώνιου τριγώνου.
\end{ylh}
% ===================================================================
%  Αρχείο: 37126.tex (Fragment)
%  Αυτόματη μετατροπή μέσω doc2latex.py
% ===================================================================



Στο παρακάτω σχήμα το ορθογώνιο $E Z H \varTheta $ παριστάνει ένα τραπέζι του μπιλιάρδου. Ένας παίκτης τοποθετεί μια μπάλα στο σημείο $A$ το οποίο ανήκει στη μεσοκάθετη της $EZ$. Όταν ο παίκτης χτυπήσει τη μπάλα αυτή ακολουθεί τη διαδρομή $A \rightarrow B \rightarrow \varGamma  \rightarrow \varDelta   \rightarrow A$ χτυπώντας στους τοίχους του μπιλιάρδου $E\varTheta , \varTheta  H, ZH$ διαδοχικά. Για τη διαδρομή αυτή ισχύει ότι κάθε γωνία πρόσπτωσης σε τοίχο (π.χ. η γωνία $A\hat{B}E$) είναι ίση με κάθε γωνία ανάκλασης σε τοίχο (π.χ. η γωνία $\varTheta \hat{B}\varGamma $) και η κάθε μια απ' αυτές είναι $45^{\circ}$.

\begin{alist}
\item Να αποδείξετε ότι:
\begin{alist}
\item Τα τρίγωνα $EAB$ και $ZA\varDelta  $ είναι ίσα. \monades{9}
\item Η διαδρομή $AB\varGamma \varDelta  $ της μπάλας είναι τετράγωνο. \monades{8}
\end{alist}
\item Αν η $AZ$ είναι διπλάσια από την απόσταση του $A$ από τον τοίχο $EZ$, να υπολογίσετε τις γωνίες του τριγώνου $AEZ$. \monades{8}

\begin{center}
\begin{tikzpicture}[scale=0.8]
\tkzDefPoint(0,4){E}
\tkzDefPoint(6,4){Z}
\tkzDefPoint(6,0){H}
\tkzDefPoint(0,0){Theta}
\tkzDefPoint(3,3){A}
\tkzDefPoint(0,0){B} % simplified for diagram, assumes corner reflection or generic
\tkzDefPoint(0,1){B} % improved
\tkzDefPoint(1,0){Gamma}
\tkzDefPoint(6,5){Delta_dummy} % to intersect
\tkzDefPoint(6,3){Delta} % midpoint of $ZH$ for square path
\tkzDefPoint(4,0){Gamma_real}
\tkzDefPoint(0,1){B_real}
% Let's use a symmetric square path
\tkzDefPoint(3,4){A_top} % intersection of perp bisector and wall
\tkzDefPoint(3,3){A}
\tkzDefPoint(0,0.5){B}
\tkzDefPoint(2.5,0){G}
\tkzDefPoint(6,3.5){D}
% No, text says 45 deg reflections.
% A(3,3). wall at x=0. Reflection at B(0,0). Angle A-B-Theta is 45? (3-0)/(3-0) = 1. Yes.
% B(0,0). wall at y=0. Reflection at G on y=0. No, B $IS$ on junction ETheta/ThetaH if A(3,3).
% But text says "walls ETheta, ThetaH, $ZH$ sequentially".
% Let's use A(3,2), table height 4. dist(A, $EZ$) = 2.
% B on x=0. A(3,2) -> B(0, yB). Slope 1 \implies yB = 2 - 3 = -1? No.
% A(3,2) -> B(0, 5)? No. 
% Downward: A(3,2) -> B(0, -1)? No.
% Let's use A(2, 3) in a 6x4 table.
% B on ETheta (x=0). A(2,3) -> B(0, 1). (dx=2, dy=2). Correct.
% Gamma on ThetaH (y=0). B(0,1) -> Gamma(1,0). (dx=1, dy=1). Correct.
% Delta on $ZH$ (x=6). Gamma(1,0) -> Delta(6, 5). (Out of bounds).
% Table must be wide enough.
\tkzDefPoint(0,6){E} \tkzDefPoint(10,6){Z} \tkzDefPoint(10,0){H} \tkzDefPoint(0,0){Theta}
\tkzDefPoint(5,3){A}
\tkzDefPoint(0,-2){B_dummy}
\tkzDefPoint(0,1){B} % dx=5, dy=2. Not 45.
% $FOR$ 45 $DEG$: dx=dy.
% A(5, 4). B on x=0. B(0, 4-5) = (0, -1). 
% Let's just draw a illustrative path.
\tkzDefPoint(3,1.5){A}
\tkzDefPoint(0,4.5){B}
\tkzDefPoint(4.5,0){Gamma}
\tkzDefPoint(6,1.5){Delta}
\tkzDefPoint(0,6){E} \tkzDefPoint(6,6){Z} \tkzDefPoint(6,0){H} \tkzDefPoint(0,0){Theta}

\draw[l3, line width=2pt] (E)--(Z)--(H)--(Theta)--cycle;
\draw[l3, secondarycolor, thick, ->] (A)--(B);
\draw[l3, secondarycolor, thick, ->] (B)--(Gamma);
\draw[l3, secondarycolor, thick, ->] (Gamma)--(Delta);
\draw[l3, secondarycolor, thick, ->] (Delta)--(A);

\tkzDrawPoints(E,Z,H,Theta,A,B,Gamma,Delta)
\tkzLabelPoint[above left](E){$E$}
\tkzLabelPoint[above right](Z){$Z$}
\tkzLabelPoint[below right](H){$H$}
\tkzLabelPoint[below left](Theta){$\varTheta $}
\tkzLabelPoint[right](A){$A$}
\tkzLabelPoint[left](B){$B$}
\tkzLabelPoint[below](Gamma){$\varGamma $}
\tkzLabelPoint[right](Delta){$\varDelta  $}

\end{tikzpicture}
\end{center}
\end{alist}

\vspace{4mm}

% --- Άσκηση 37128 (37128.tex) ---
\Askhsh{37128}{4}
\begin{ylh}{1}
    \item 3.1. Είδη και στοιχεία τριγώνων
    \item 3.6. Κριτήρια ισότητας ορθογώνιων τριγώνων
    \item 3.11. Ανισοτικές σχέσεις πλευρών και γωνιών
    \item 4.2. Τέμνουσα δύο ευθειών - Ευκλείδειο αίτημα
    \item 4.6. Άθροισμα γωνιών τριγώνου
    \item 5.4. Ρόμβος
    \item 5.11. Ισοσκελές τραπέζιο.
\end{ylh}
% ===================================================================
%  Αρχείο: 37128.tex (Fragment)
%  Αυτόματη μετατροπή μέσω doc2latex.py
% ===================================================================



Έστω ισοσκελές τραπέζιο $AB\varGamma \varDelta  $ ($AB$ // $\varDelta  \varGamma $) με $\hat{B} = 2 \cdot \hat{\varGamma }$ και $ΑΒ = B\varGamma  = A\varDelta   = \ \frac{\varGamma \varDelta  }{2}$. Φέρουμε τη διχοτόμο της γωνίας $\hat{B}$, η οποία τέμνει το $\varDelta  \varGamma $ στο Κ και η κάθετη από το Κ προς το $B\varGamma $ το τέμνει στο Μ.

\begin{alist}
\item Να υπολογίσετε τις γωνίες του $AB\varGamma \varDelta  $. \monades{10}

\item Να αποδείξετε ότι:


\end{alist}
\begin{alist}
\item
  Το τετράπλευρο $ABK\varDelta  $ είναι ρόμβος. \monades{8}
\item
  Το σημείο $M$ είναι το μέσο του $B\varGamma $. \monades{7}

\begin{center}
\begin{tikzpicture}[scale=1]
\tkzDefPoint(0,0){Delta}
\tkzDefPoint(4,0){Gamma}
\tkzDefPoint(1,1.732){A}
\tkzDefPoint(3,1.732){B}

% K is intersection of B bisector with $\varDelta  \varGamma $
\tkzDefLine[bisector](A,B,Gamma) \tkzGetPoint{bis_B}
\tkzInterLL(B,bis_B)(Delta,Gamma) \tkzGetPoint{K}

% M is intersection of perp from K to $B\varGamma $
\tkzDefPointBy[projection=onto B--Gamma](K) \tkzGetPoint{M}

\draw[l3] (A)--(B)--(Gamma)--(Delta)--cycle;
\draw[l3] (B)--(K);
\draw[l3] (K)--(M);

% For question: $ABK\varDelta  $ rhombus
\draw[l3, dashed] (A)--(K);

\tkzDrawPoints(A,B,Gamma,Delta,K,M)
\tkzLabelPoint[above left](A){$A$}
\tkzLabelPoint[above right](B){$B$}
\tkzLabelPoint[below right](Gamma){$\varGamma $}
\tkzLabelPoint[below left](Delta){$\varDelta  $}
\tkzLabelPoint[below](K){$ K$}
\tkzLabelPoint[right](M){$\mathrm{M}$}

\tkzMarkRightAngle[size=0.2](K,M,Gamma)
\end{tikzpicture}
\end{center}
\end{alist}

\vspace{4mm}

% --- Άσκηση 37129 (37129.tex) ---
\Askhsh{37129}{4}
\begin{ylh}{1}
    \item 5.2. Παραλληλόγραμμα
    \item 5.9. Μια ιδιότητα του ορθογώνιου τριγώνου
    \item 5.10. Τραπέζιο
    \item 5.11. Ισοσκελές τραπέζιο.
\end{ylh}
% ===================================================================
%  Αρχείο: 37129.tex (Fragment)
%  Αυτόματη μετατροπή μέσω doc2latex.py
% ===================================================================



Έστω τετράγωνο $AB\varGamma \varDelta  $ και $M$ το μέσο της πλευράς $\varDelta  A$ . Προεκτείνουμε το τμήμα $\varDelta  A$ (προς την πλευρά του Α)κατά τμήμα $ΑΝ = \ \frac{A\varDelta  }{2}$. Φέρουμε τα τμήματα $\varGamma M$ και $BN$ και θεωρούμε τα μέσα τους Κ και Λ αντίστοιχα.

Να αποδείξετε ότι:

\begin{alist}
\item Το τετράπλευρο ΜΝ$B\varGamma $ είναι παραλληλόγραμμο. \monades{8}

\item Το τετράπλευρο $A\varDelta  K\varLambda $ είναι παραλληλόγραμμο. \monades{9}

\item Το τετράπλευρο $AMK\varLambda $ είναι ισοσκελές τραπέζιο. \monades{8}

\begin{center}
\begin{tikzpicture}[scale=1]
\tkzDefPoint(0,2){A}
\tkzDefPoint(2,2){B}
\tkzDefPoint(2,0){Gamma}
\tkzDefPoint(0,0){Delta}

\tkzDefMidPoint(A,Delta) \tkzGetPoint{M}

% N on extension of $\varDelta  A$ past A, $AN$ = $A\varDelta  $/2. length=1.
\tkzDefPoint(0,3){N}

\tkzDefMidPoint(Gamma,M) \tkzGetPoint{K}
\tkzDefMidPoint(B,N) \tkzGetPoint{L}

\draw[l3] (A)--(B)--(Gamma)--(Delta)--cycle;
\draw[l3] (A)--(N);
\draw[l3] (Gamma)--(M);
\draw[l3] (B)--(N);

% For questions (selected dashed lines)
\draw[l3, dashed] (M)--(N)--(B)--(Gamma);
\draw[l3, dashed] (A)--(Delta)--(K)--(L)--cycle;

\tkzDrawPoints(A,B,Gamma,Delta,M,N,K,L)
\tkzLabelPoint[left](A){$A$}
\tkzLabelPoint[right](B){$B$}
\tkzLabelPoint[right](Gamma){$\varGamma $}
\tkzLabelPoint[left](Delta){$\varDelta  $}
\tkzLabelPoint[left](M){$\mathrm{M}$}
\tkzLabelPoint[left](N){$\mathrm{N}$}
\tkzLabelPoint[below](K){$ K$}
\tkzLabelPoint[above right](L){$\Lambda$}
\end{tikzpicture}
\end{center}
\end{alist}

\vspace{4mm}

% --- Άσκηση 37130 (37130.tex) ---
\Askhsh{37130}{4}
\begin{ylh}{1}
    \item 5.2. Παραλληλόγραμμα
    \item 5.9. Μια ιδιότητα του ορθογώνιου τριγώνου
    \item 5.10. Τραπέζιο
    \item 5.11. Ισοσκελές τραπέζιο.
\end{ylh}

Σε παραλληλόγραμμο $AB\varGamma \varDelta  $ με $AB>B\varGamma $ και $\hat{B}<90\degree$ θεωρούμε σημείο $Z$ στην προέκταση της $B\varGamma $ (προς το $\varGamma$) τέτοιο ώστε $\varGamma Z=B\varGamma $. Αν $E$ είναι σημείο της $AB$, τέτοιο ώστε $E\varGamma =\varGamma B$, να αποδείξετε ότι:

\begin{alist}
\item η γωνία $BE$Ζ είναι ορθή, \monades{8}
\item το τετράπλευρο $AE\varGamma \varDelta  $ είναι ισοσκελές τραπέζιο, \monades{8}
\item το τετράπλευρο $A\varGamma Z\varDelta  $ είναι παραλληλόγραμμο. \monades{9}
\end{alist}
\vspace{4mm}

% --- Άσκηση 37131 (37131.tex) ---
\Askhsh{37131}{4}
\begin{ylh}{1}
    \item 3.3. 2ο Κριτήριο ισότητας τριγώνων
    \item 3.4. 3ο Κριτήριο ισότητας τριγώνων
    \item 3.11. Ανισοτικές σχέσεις πλευρών και γωνιών
    \item 4.2. Τέμνουσα δύο ευθειών - Ευκλείδειο αίτημα
    \item 5.2. Παραλληλόγραμμα.
\end{ylh}
% ===================================================================
%  Αρχείο: 37131.tex (Fragment)
%  Αυτόματη μετατροπή μέσω doc2latex.py
% ===================================================================



Δίνεται τρίγωνο $AB\varGamma $, με $AK$ διχοτόμο της γωνίας Α. Στην προέκταση της $AK$ θεωρούμε σημείο $\varDelta  $ ώστε $AK$ = $K\varDelta  $. Η παράλληλη από το $\varDelta$ προς την $AB$ τέμνει τις $A\varGamma $ και $B\varGamma $ στα $E$ και $Z$ αντίστοιχα.

Να αποδείξετε ότι:

\begin{alist}
\item Το τρίγωνο Α$E\varDelta  $ είναι ισοσκελές. \monades{6}

\item Η $EK$ είναι μεσοκάθετος της $A\varDelta  $ . \monades{6}

\item Τα τρίγωνα $AKB$ και Κ$\varDelta   Z$ είναι ίσα. \monades{7}

\item Το τετράπλευρο $AZ\varDelta  B$ είναι παραλληλόγραμμο. \monades{6}

\begin{center}
\begin{tikzpicture}[scale=0.9]
\tkzDefPoint(1,3){A}
\tkzDefPoint(0,0){B}
\tkzDefPoint(5,0){Gamma}

\tkzDefLine[bisector](B,A,Gamma) \tkzGetPoint{bis_A}
\tkzInterLL(A,bis_A)(B,Gamma) \tkzGetPoint{K}

% Δ on extension of $AK$, $AK$=$K\varDelta  $
\tkzDefPointBy[symmetry=center K](A) \tkzGetPoint{Delta}

% parallel to $AB$ through Δ, intersects $A\varGamma $ and $B\varGamma $
\tkzDefLine[parallel=through Delta](A,B) \tkzGetPoint{par_AB}
\tkzInterLL(Delta,par_AB)(A,Gamma) \tkzGetPoint{E}
\tkzInterLL(Delta,par_AB)(B,Gamma) \tkzGetPoint{Z}

\draw[l3] (A)--(B)--(Gamma)--cycle;
\draw[l3] (A)--(Delta);
\draw[l3] (Delta)--(E);

% For questions (A,Z, Δ, B parallelogram)
\draw[l3, dashed] (A)--(Z);
\draw[l3, dashed] (B)--(Delta);
\draw[l3, dashed] (E)--(K);

\tkzDrawPoints(A,B,Gamma,K,Delta,E,Z)
\tkzLabelPoint[above](A){$A$}
\tkzLabelPoint[left](B){$B$}
\tkzLabelPoint[right](Gamma){$\varGamma $}
\tkzLabelPoint[above left](K){$ K$}
\tkzLabelPoint[below](Delta){$\varDelta  $}
\tkzLabelPoint[above right](E){$ E$}
\tkzLabelPoint[below right](Z){$ Z$}
\end{tikzpicture}
\end{center}
\end{alist}

\vspace{4mm}

% --- Άσκηση 37132 (37132.tex) ---
\Askhsh{37132}{4}
\begin{ylh}{1}
    \item 3.6. Κριτήρια ισότητας ορθογώνιων τριγώνων
    \item 4.2. Τέμνουσα δύο ευθειών - Ευκλείδειο αίτημα
    \item 5.4. Ρόμβος
    \item 5.6. Εφαρμογές στα τρίγωνα
    \item 5.9. Μια ιδιότητα του ορθογώνιου τριγώνου.
\end{ylh}
% ===================================================================
%  Αρχείο: 37132.tex (Fragment)
%  Αυτόματη μετατροπή μέσω doc2latex.py
% ===================================================================



Δίνεται τρίγωνο $AB\varGamma $. Στην προέκταση του ύψους του $AK$ θεωρούμε σημείο $\varDelta  $ ώστε $AK$ = $K\varDelta  $. Έστω Λ, Μ, Ν τα μέσα των πλευρών $AB$, $A\varGamma $ και $B\varDelta  $ αντίστοιχα.

Να αποδείξετε ότι:

\begin{alist}
\item Το τρίγωνο Α$B\varDelta  $ είναι ισοσκελές. \monades{7}

\item Το τετράπλευρο $B\varLambda KN$ είναι ρόμβος . \monades{9}

\item $\Lambda Μ\ \bot\ \Lambda Ν$ \monades{9}

\begin{center}
\begin{tikzpicture}[scale=0.9]
\tkzDefPoint(1,3){A}
\tkzDefPoint(0,0){B}
\tkzDefPoint(5,0){Gamma}

\tkzDefPointBy[projection=onto B--Gamma](A) \tkzGetPoint{K}
\tkzDefPointBy[symmetry=center K](A) \tkzGetPoint{Delta}

\tkzDefMidPoint(A,B) \tkzGetPoint{Lambda}
\tkzDefMidPoint(A,Gamma) \tkzGetPoint{M}
\tkzDefMidPoint(B,Delta) \tkzGetPoint{N}

\draw[l3] (A)--(B)--(Gamma)--cycle;
\draw[l3] (A)--(Delta);
\draw[l3] (Delta)--(B);
\draw[l3] (Delta)--(Gamma);

% Rhombus $B\varLambda KN$
\draw[l3, dashed] (B)--(Lambda)--(K)--(N)--cycle;
\draw[l3, dashed] (Lambda)--(M);
\draw[l3, dashed] (Lambda)--(N);

\tkzDrawPoints(A,B,Gamma,K,Delta,Lambda,M,N)
\tkzLabelPoint[above](A){$A$}
\tkzLabelPoint[left](B){$B$}
\tkzLabelPoint[right](Gamma){$\varGamma $}
\tkzLabelPoint[below right](K){$ K$}
\tkzLabelPoint[below](Delta){$\varDelta  $}
\tkzLabelPoint[above left](Lambda){$\Lambda$}
\tkzLabelPoint[above right](M){$ M$}
\tkzLabelPoint[below left](N){$\mathrm{N}$}

\tkzMarkRightAngle[size=0.2](A,K,B)
\end{tikzpicture}
\end{center}
\end{alist}

\vspace{4mm}

% --- Άσκηση 37133 (37133.tex) ---
\Askhsh{37133}{4}
\begin{ylh}{1}
    \item 3.2. 1ο Κριτήριο ισότητας τριγώνων
    \item 5.2. Παραλληλόγραμμα
    \item 5.3. Ορθογώνιο
    \item 5.6. Εφαρμογές στα τρίγωνα.
\end{ylh}
% ===================================================================
%  Αρχείο: 37133.tex (Fragment)
%  Αυτόματη μετατροπή μέσω doc2latex.py
% ===================================================================



Θεωρούμε τρίγωνο $ΑΒ\varGamma $ και τις μεσοκαθέτους $\mu_1, \mu_2$ των πλευρών του $ΑΒ$ και $Α\varGamma $, οι οποίες τέμνονται στο μέσο $M$ της $Β\varGamma $.

\begin{alist}
\item Να αποδείξετε ότι:
\begin{alist}
\item Το τρίγωνο είναι ορθογώνιο με $\hat{A} = 90\degree$. \monades{5}
\item Το τετράπλευρο $A\Lambda M K$ είναι ορθογώνιο παραλληλόγραμμο (όπου $K, \Lambda$ τα μέσα των $ΑΒ, Α\varGamma $). \monades{7}
\item $\Lambda\varTheta  = \frac{Β\varGamma }{4}$, όπου $\varTheta $ το σημείο τομής των $ΑΜ$ και $K\Lambda$. \monades{6}
\end{alist}
\item Αν $I$ σημείο της $Β\varGamma $ τέτοιο ώστε $ΒI = \frac{Β\varGamma }{4}$, να αποδείξετε ότι το τετράπλευρο $K\varTheta  IB$ είναι παραλληλόγραμμο. \monades{7}

\begin{center}
\begin{tikzpicture}[scale=1.2]
\tkzDefPoint(0,0){A}
\tkzDefPoint(4,0){C}
\tkzDefPoint(0,3){B}
\tkzDefMidPoint(B,C) \tkzGetPoint{M}
\tkzDefMidPoint(A,B) \tkzGetPoint{K}
\tkzDefMidPoint(A,C) \tkzGetPoint{L}
\tkzInterLL(A,M)(K,L) \tkzGetPoint{Theta}
\tkzDefPointBy[homothety=center B ratio 0.25](C) \tkzGetPoint{I}

\draw[l3] (A)--(B)--(C)--cycle;
\draw[l3] (A)--(M);
\draw[l3] (K)--(L);
\draw[l3] (M)--(K);
\draw[l3] (M)--(L);
\draw[l3, dashed] (K)--(Theta)--(I)--(B)--cycle;

\tkzDrawPoints(A,B,C,M,K,L,Theta,I)
\tkzLabelPoint[below left](A){$A$}
\tkzLabelPoint[above left](B){$B$}
\tkzLabelPoint[below right](C){$\varGamma $}
\tkzLabelPoint[above right](M){$M$}
\tkzLabelPoint[left](K){$K$}
\tkzLabelPoint[below](L){$\Lambda$}
\tkzLabelPoint[above](Theta){$\varTheta $}
\tkzLabelPoint[below right](I){$I$}

\tkzMarkRightAngle(B,A,C)
\tkzMarkSegments[mark=s|](B,M M,C A,M)
\end{tikzpicture}
\end{center}
\end{alist}

\vspace{4mm}

% --- Άσκηση 37134 (37134.tex) ---
\Askhsh{37134}{4}
\begin{ylh}{1}
    \item 3.1. Είδη και στοιχεία τριγώνων
    \item 3.2. 1ο Κριτήριο ισότητας τριγώνων
    \item 3.11. Ανισοτικές σχέσεις πλευρών και γωνιών
    \item 4.2. Τέμνουσα δύο ευθειών - Ευκλείδειο αίτημα
    \item 4.6. Άθροισμα γωνιών τριγώνου
    \item 5.2. Παραλληλόγραμμα
    \item 5.9. Μια ιδιότητα του ορθογώνιου τριγώνου.
\end{ylh}
% ===================================================================
%  Αρχείο: 37134.tex (Fragment)
%  Αυτόματη μετατροπή μέσω doc2latex.py
% ===================================================================



Δίνεται τραπέζιο $ΑΒ\varGamma \varDelta  $ με $ΑΒ // \varGamma \varDelta  $, $\varDelta  \varGamma  = 4ΑΒ$ και $Β\varGamma  = 2ΑΒ$. Θεωρούμε σημείο $Z$ της $\varGamma \varDelta  $, ώστε $\varDelta   Z = ΑΒ$. Αν η γωνία $\hat{\varGamma }$ είναι $60\degree$ και $BE$ το ύψος του τραπεζίου, να αποδείξετε ότι:

\begin{alist}
\item Το τετράπλευρο $ΑΒ\varGamma  E$ είναι παραλληλόγραμμο. \monades{8}
\item Το τρίγωνο $ZAE$ είναι ισόπλευρο. \monades{8}
\item Τα τρίγωνα $\varDelta   AZ$ και $\varGamma  AE$ είναι ίσα. \monades{9}

\begin{center}
\begin{tikzpicture}[scale=1]
\tkzDefPoint(0,0){D}
\tkzDefPoint(1,0){Z}
\tkzDefPoint(2,2){h_top} % height
% $AB$=1. $CD$=4. $BC$=2. Gamma=60.
% C is at (4,0). B is at C + (2*cos120, 2*sin120) = (4-1, sqrt(3)) = (3, 1.73).
% A is at B + (-1, 0) = (2, 1.73).
\tkzDefPoint(4,0){C}
\tkzDefPoint(3,1.732){B}
\tkzDefPoint(2,1.732){A}
\tkzDefPoint(0,0){D}
\tkzDefPoint(1,0){Z}
\tkzDefPointBy[projection=onto D--C](B) \tkzGetPoint{E}

\draw[l3] (A)--(B)--(C)--(D)--cycle;
\draw[l3] (A)--(Z);
\draw[l3] (A)--(E);
\draw[l3] (B)--(E);

\tkzDrawPoints(A,B,C,D,Z,E)
\tkzLabelPoint[above left](A){$A$}
\tkzLabelPoint[above right](B){$B$}
\tkzLabelPoint[below right](C){$\varGamma $}
\tkzLabelPoint[below left](D){$\varDelta  $}
\tkzLabelPoint[below](Z){$Z$}
\tkzLabelPoint[below](E){$E$}

\tkzMarkRightAngle(B,E,C)
\tkzMarkAngle[size=0.5](B,C,E)
\tkzLabelAngle[pos=0.8](B,C,E){$60^{\circ}$}
\tkzMarkSegments[mark=||](A,B B,C_p) % illustrative
\end{tikzpicture}
\end{center}
\end{alist}

\vspace{4mm}

% --- Άσκηση 37135 (37135.tex) ---
\Askhsh{37135}{4}
\begin{ylh}{1}
    \item 3.2. 1ο Κριτήριο ισότητας τριγώνων
    \item 3.6. Κριτήρια ισότητας ορθογώνιων τριγώνων
    \item 4.2. Τέμνουσα δύο ευθειών - Ευκλείδειο αίτημα
    \item 5.10. Τραπέζιο
    \item 5.11. Ισοσκελές τραπέζιο.
\end{ylh}
% ===================================================================
%  Αρχείο: 37135.tex (Fragment)
%  Αυτόματη μετατροπή μέσω doc2latex.py
% ===================================================================



Δίνεται ισοσκελές τραπέζιο $AB\varGamma \varDelta  $ με $AB$ // $\varGamma \varDelta  $ και $A\varDelta  $ = $B\varGamma $ = $AB$. Φέρουμε τμήματα $AE$ και $BZ$ κάθετα στις διαγώνιες $B\varDelta  $ και $A\varGamma $ αντίστοιχα.

Να αποδείξετε ότι:

\begin{alist}
\item Τα σημεία $Z$ και $E$ είναι μέσα των διαγωνίων $A\varGamma $ και $B\varDelta  $ αντίστοιχα. \monades{5}

\item $AE$ = $BZ$ . \monades{8}

\item Το τετράπλευρο $AEZB$ είναι ισοσκελές τραπέζιο. \monades{7}

\item Η $B\varDelta  $ είναι διχοτόμος της γωνίας Δ. \monades{5}

\begin{center}
\begin{tikzpicture}[scale=1]
\tkzDefPoint(0,0){Delta}
\tkzDefPoint(4,0){Gamma}
\tkzDefPoint(1,1.732){A}
\tkzDefPoint(3,1.732){B}

\draw[l3] (A)--(B)--(Gamma)--(Delta)--cycle;
\draw[l3] (A)--(Gamma);
\draw[l3] (B)--(Delta);

\tkzDefPointBy[projection=onto B--Delta](A) \tkzGetPoint{E}
\tkzDefPointBy[projection=onto A--Gamma](B) \tkzGetPoint{Z}

\draw[l3] (A)--(E);
\draw[l3] (B)--(Z);
\draw[l3, dashed] (E)--(Z);

\tkzDrawPoints(A,B,Gamma,Delta,E,Z)
\tkzLabelPoint[above left](A){$A$}
\tkzLabelPoint[above right](B){$B$}
\tkzLabelPoint[below right](Gamma){$\varGamma $}
\tkzLabelPoint[below left](Delta){$\varDelta  $}
\tkzLabelPoint[below](E){$ E$}
\tkzLabelPoint[below](Z){$ Z$}

\tkzMarkRightAngle[size=0.2](A,E,Delta)
\tkzMarkRightAngle[size=0.2](B,Z,Gamma)
\end{tikzpicture}
\end{center}
\end{alist}

\vspace{4mm}

% --- Άσκηση 37136 (37136.tex) ---
\Askhsh{37136}{4}
\begin{ylh}{1}
    \item 3.6. Κριτήρια ισότητας ορθογώνιων τριγώνων
    \item 3.11. Ανισοτικές σχέσεις πλευρών και γωνιών
    \item 5.2. Παραλληλόγραμμα
    \item 5.9. Μια ιδιότητα του ορθογώνιου τριγώνου.
\end{ylh}
Δίνεται παραλληλόγραμμο $AB\varGamma \varDelta  $ τέτοιο ώστε αν φέρουμε την κάθετη στην $A\varGamma $ στο κέντρο του Ο, αυτή τέμνει την προέκταση της $A\varDelta  $ σε σημείο $E$ τέτοιο ώστε $\varDelta  E$=$A\varDelta  $. Να αποδείξετε ότι:
\begin{alist}
\item Το τρίγωνο $AE\varGamma $ είναι ισοσκελές. \monades{7}
\item Το τετράπλευρο $B\varGamma $$E\varDelta  $ είναι παραλληλόγραμμο. \monades{9}
\item Το τρίγωνο $BO\varGamma $ είναι ισοσκελές. \monades{9}
\begin{center}
\begin{tikzpicture}[scale=0.6]
\tkzDefPoint(0,0){A}
\tkzDefPoint(2,0){Delta}
\tkzDefPoint(4,-4){Gamma}
\tkzDefPoint(2,-4){B}
\tkzDefPoint(4,0){E}
\tkzDefMidPoint(A,Gamma) \tkzGetPoint{O}
\draw[l3] (A)--(B)--(Gamma)--(Delta)--cycle;
\draw[l3] (A)--(Gamma);
\draw[l3] (B)--(Delta);
\draw[l3] (Delta)--(E);
\draw[l3] (E)--(O);
\draw[l3, dashed] (E)--(Gamma);
\draw[l3, dashed] (E)--(B);
\tkzMarkRightAngle[size=0.3](E,O,Gamma)
\tkzDrawPoints(A,B,Gamma,Delta,E,O)
\tkzLabelPoint[left](A){$A$}
\tkzLabelPoint[below left](B){$B$}
\tkzLabelPoint[below right](Gamma){$\varGamma $}
\tkzLabelPoint[above](Delta){$\varDelta  $}
\tkzLabelPoint[right](E){$ E$}
\tkzLabelPoint[below](O){$ O$}


\end{tikzpicture}
\end{center}
\end{alist}

\vspace{4mm}

% --- Άσκηση 37137 (37137.tex) ---
\Askhsh{37137}{4}
\begin{ylh}{1}
    \item 3.6. Κριτήρια ισότητας ορθογώνιων τριγώνων
    \item 3.11. Ανισοτικές σχέσεις πλευρών και γωνιών
    \item 5.8. Το ορθόκεντρο τριγώνου.
\end{ylh}
% ===================================================================
%  Αρχείο: 37137.tex (Fragment)
%  Αυτόματη μετατροπή μέσω doc2latex.py
% ===================================================================



Δίνεται ορθογώνιο τρίγωνο $AB\varGamma $ ($\hat{A} = 90^{o}$) με $B\varDelta  $ διχοτόμο και $AK$ ύψος, που τέμνονται στο Ε. Η κάθετη από το $E$ στην $AB$ τέμνει τις $AB$ και $B\varGamma $ στα Η και $Z$ αντίστοιχα.

\begin{alist}
\item Να αποδείξετε ότι:


\end{alist}
\begin{alist}
\item
  τα τρίγωνα $EHA$ και $EKZ$ είναι ίσα. \monades{6}
\item
  το τρίγωνο $BKH$ είναι ισοσκελές. \monades{6}
\item
  η $B\varDelta  $ είναι κάθετη στην $AZ$. \monades{7}

\item Αν επιπλέον το ορθογώνιο τρίγωνο $AB\varGamma $ είναι και ισοσκελές, να αποδείξετε ότι η $\varGamma E$ είναι διχοτόμος της γωνίας Γ. \monades{6}

%\includesvg[width=5.98105in,height=3.39583in]{/home/spyros/Μαθηματικά/Trapeza_Thematwn/A_Lykeiou/Geometry/extracted/37137/media/image2.svg}
\end{alist}

\vspace{4mm}

% --- Άσκηση 37138 (37138.tex) ---
\Askhsh{37138}{4}
\begin{ylh}{1}
    \item 3.1. Είδη και στοιχεία τριγώνων
    \item 3.4. 3ο Κριτήριο ισότητας τριγώνων
    \item 3.7. Κύκλος - Μεσοκάθετος - Διχοτόμος
    \item 4.6. Άθροισμα γωνιών τριγώνου
    \item 5.9. Μια ιδιότητα του ορθογώνιου τριγώνου.
\end{ylh}
% ===================================================================
%  Αρχείο: 37138.tex (Fragment)
%  Αυτόματη μετατροπή μέσω doc2latex.py
% ===================================================================



Δίνεται ισόπλευρο τρίγωνο $AB\varGamma $. Με βάση την $AB$ κατασκευάζουμε ισοσκελές τρίγωνο $A\varDelta  B$, εκτός του τριγώνου $AB\varGamma $, με γωνία $\hat{\varDelta  } = 120\degree$. Θεωρούμε τα μέσα $Z$ και Η των πλευρών $A\varDelta  $ και $A\varGamma $ αντίστοιχα.

\begin{alist}
\item Να αποδείξετε ότι η $\varDelta  \varGamma $ είναι μεσοκάθετος του $AB$. \monades{8}

\item Αν η $\varDelta  \varGamma $ τέμνει την $AB$ στο Θ, να αποδείξετε ότι η γωνία $Z\hat{\varTheta }H$ είναι ορθή. \monades{9}

\item Αν η $ZK$ είναι η κάθετη στην $AB$ από το σημείο $Z$, να αποδείξετε ότι $ΖΚ = \frac{Α\varDelta  }{4}$. \monades{8}

%\includesvg[width=5.49197in,height=4.34375in]{/home/spyros/Μαθηματικά/Trapeza_Thematwn/A_Lykeiou/Geometry/extracted/37138/media/image2.svg}
\end{alist}

\vspace{4mm}

% --- Άσκηση 37139 (37139.tex) ---
\Askhsh{37139}{4}
\begin{ylh}{1}
    \item 3.4. 3ο Κριτήριο ισότητας τριγώνων
    \item 3.7. Κύκλος - Μεσοκάθετος - Διχοτόμος
    \item 3.11. Ανισοτικές σχέσεις πλευρών και γωνιών
    \item 4.2. Τέμνουσα δύο ευθειών - Ευκλείδειο αίτημα
    \item 5.6. Εφαρμογές στα τρίγωνα
    \item 5.11. Ισοσκελές τραπέζιο.
\end{ylh}
% ===================================================================
%  Αρχείο: 37139.tex (Fragment)
%  Αυτόματη μετατροπή μέσω doc2latex.py
% ===================================================================



Δίνεται ισοσκελές τραπέζιο $AB\varGamma \varDelta  $ ($AB$//$\varDelta  \varGamma $) και Ο το σημείο τομής των διαγωνίων του. Η $A\varGamma $ είναι κάθετη στην $A\varDelta  $ και η $B\varDelta  $ είναι κάθετη στην $B\varGamma $. Θεωρούμε τα μέσα Μ, Ε και $Z$ των $\varGamma \varDelta  $, $B\varDelta  $ και $A\varGamma $ αντίστοιχα.

Να αποδείξετε ότι:

\begin{alist}
\item $ME$=$MZ$. \monades{6}

\item Η $MZ$ είναι κάθετη στην $A\varGamma $. \monades{6}

\item Τα τρίγωνα $M\varDelta  E$ και $MZ\varGamma $ είναι ίσα. \monades{7}

\item Η $OM$ είναι μεσοκάθετος του $EZ$. \monades{6}

%\includesvg[width=5.5625in,height=2.73958in]{/home/spyros/Μαθηματικά/Trapeza_Thematwn/A_Lykeiou/Geometry/extracted/37139/media/image2.svg}
\end{alist}

\vspace{4mm}

% --- Άσκηση 37140 (37140.tex) ---
\Askhsh{37140}{4}
\begin{ylh}{1}
    \item 3.2. 1ο Κριτήριο ισότητας τριγώνων
    \item 5.2. Παραλληλόγραμμα
    \item 5.4. Ρόμβος
    \item 5.6. Εφαρμογές στα τρίγωνα
    \item 5.9. Μια ιδιότητα του ορθογώνιου τριγώνου.
\end{ylh}
% ===================================================================
%  Αρχείο: 37140.tex (Fragment)
%  Αυτόματη μετατροπή μέσω doc2latex.py
% ===================================================================



Δίνεται ισόπλευρο τρίγωνο $AB\varGamma $ και τα μέσα Δ, Ε και $M$ των $AB$, $A\varGamma $ και $B\varGamma $ αντίστοιχα. Στην προέκταση του $M\varDelta  $ (προς το Δ) θεωρούμε τμήμα $\varDelta   Z$=$\varDelta  M$.

Να αποδείξετε ότι:

\begin{alist}
\item Τα τρίγωνα $AZ\varDelta  $ $KAI$ $BM\varDelta  $ είναι ίσα. \monades{6}

\item Το τετράπλευρο $ZA\varGamma M$ είναι παραλληλόγραμμο. \monades{6}

\item Τα τμήματα $ZE$ και $A\varDelta  $ τέμνονται κάθετα και διχοτομούνται. \monades{7}

\item Η $BZ$ είναι κάθετη στη $ZA$. \monades{6}

%\includesvg[width=4.01042in,height=3.3027in]{/home/spyros/Μαθηματικά/Trapeza_Thematwn/A_Lykeiou/Geometry/extracted/37140/media/image2.svg}
\end{alist}

\vspace{4mm}

% --- Άσκηση 37141 (37141.tex) ---
\Askhsh{37141}{4}
\begin{ylh}{1}
    \item 3.2. 1ο Κριτήριο ισότητας τριγώνων
    \item 3.6. Κριτήρια ισότητας ορθογώνιων τριγώνων
    \item 5.4. Ρόμβος.
\end{ylh}
% ===================================================================
%  Αρχείο: 37141.tex (Fragment)
%  Αυτόματη μετατροπή μέσω doc2latex.py
% ===================================================================



Δίνεται τρίγωνο $AB\varGamma $ με $AB$ < $A\varGamma $. Φέρνουμε τμήμα $B\varDelta  $ κάθετο στην $AB$ με $B\varDelta  $ = $A\varGamma $ και τμήμα $\varGamma E$ κάθετο στην $A\varGamma $ με $\varGamma E$ = $AB$. Θεωρούμε τα μέσα $Z$ και Θ των $A\varDelta  $ και $AE$ καθώς και τη διχοτόμο Αδ της γωνίας $\varDelta  \hat{A}E$.

\begin{alist}
\item Να αποδείξετε ότι $A\varDelta  $ = $AE$. \monades{9}

\item Αν Κ τυχαίο σημείο της διχοτόμου Αδ, να αποδείξετε ότι ισαπέχει από τα μέσα $Z$ και Θ. \monades{9}

\item Αν το Κ είναι σημείο της διχοτόμου Αδ τέτοιο ώστε $KZ$ = $AZ$, να αποδείξετε ότι το τετράπλευρο $AZK\varTheta $ είναι ρόμβος. \monades{7}

%\includesvg[width=6.29861in,height=4.70833in]{/home/spyros/Μαθηματικά/Trapeza_Thematwn/A_Lykeiou/Geometry/extracted/37141/media/image2.svg}
\end{alist}

\vspace{4mm}

% --- Άσκηση 37142 (37142.tex) ---
\Askhsh{37142}{4}
\begin{ylh}{1}
    \item 3.1. Είδη και στοιχεία τριγώνων
    \item 3.2. 1ο Κριτήριο ισότητας τριγώνων
    \item 3.6. Κριτήρια ισότητας ορθογώνιων τριγώνων
    \item 3.7. Κύκλος - Μεσοκάθετος - Διχοτόμος
    \item 5.9. Μια ιδιότητα του ορθογώνιου τριγώνου.
\end{ylh}
% ===================================================================
%  Αρχείο: 37142.tex (Fragment)
%  Αυτόματη μετατροπή μέσω doc2latex.py
% ===================================================================



Δίνεται οξυγώνιο ισοσκελές τρίγωνο $AB\varGamma $ με $AB=A\varGamma $. Φέρνουμε τμήμα $A\varDelta  $ κάθετο στην $AB$ και τμήμα $AE$ κάθετο στην $A\varGamma $ με $A\varDelta  $ = $AE$. Θεωρούμε τα μέσα Ζ, Η και $M$ των $\varDelta   B$, $E\varGamma $ και $B\varGamma $ αντίστοιχα.

\begin{alist}
\item Να αποδείξετε ότι:


\end{alist}
\begin{alist}
\item
  Τα τρίγωνα $A\varDelta  B$ και $AE\varGamma $ είναι ίσα. \monades{7}
\item
  Το τρίγωνο $ZAH$ είναι ισοσκελές. \monades{6}
\item
  Η $AM$ είναι μεσοκάθετος του $ZH$. \monades{7}

\item Ένας μαθητής συγκρίνοντας τα τρίγωνα $A\varDelta  B$ και $AE\varGamma $ έγραψε τα εξής:

« 1. \emph{$A\varDelta  $ = $AE$ από υπόθεση}

\emph{2. $AB=Α\varGamma $ πλευρές ισοσκελούς τριγώνου}

\emph{3.} $\varDelta  \hat{A}Β = Ε\hat{A}\varGamma $ \emph{ως κατακορυφήν}

\emph{Άρα τα τρίγωνα είναι ίσα έχοντας δύο πλευρές ίσες μια προς μια και την περιεχόμενη γωνία ίση»}.

Ο καθηγητής είπε ότι αυτή η λύση περιέχει λάθος. Μπορείς να το εντοπίσεις; \monades{5}

%\includesvg[width=4.67735in,height=4.125in]{/home/spyros/Μαθηματικά/Trapeza_Thematwn/A_Lykeiou/Geometry/extracted/37142/media/image2.svg}
\end{alist}

\vspace{4mm}

% --- Άσκηση 37156 (37156.tex) ---
\Askhsh{37156}{4}
\begin{ylh}{1}
    \item 3.11. Ανισοτικές σχέσεις πλευρών και γωνιών
    \item 4.6. Άθροισμα γωνιών τριγώνου
    \item 5.6. Εφαρμογές στα τρίγωνα
    \item 5.9. Μια ιδιότητα του ορθογώνιου τριγώνου.
\end{ylh}
% ===================================================================
%  Αρχείο: 37156.tex (Fragment)
%  Αυτόματη μετατροπή μέσω doc2latex.py
% ===================================================================



Έστω ισοσκελές τρίγωνο $BA\varGamma $ με $\hat{A} = \ 120\degree$. Φέρουμε ημιευθεία Αx κάθετη στην $A\varGamma $ στο Α, η οποία τέμνει τη $B\varGamma $ στο Δ. Έστω Λ το μέσο του $AB$ και Κ το μέσο του $\varDelta  \varGamma $.

Να αποδείξετε ότι:

\begin{alist}
\item Το τρίγωνο $A\varDelta  B$ είναι ισοσκελές \monades{8}

\item $\varDelta  \varGamma $ = 2$\cdot$$B\varDelta  $ \monades{8}

\item $\varLambda \varDelta  $ // $AK$ \monades{5}

\item $AK$ = 2$\cdot$ΛΔ \monades{4}

%\includesvg[width=5.72413in,height=2.88542in]{/home/spyros/Μαθηματικά/Trapeza_Thematwn/A_Lykeiou/Geometry/extracted/37156/media/image2.svg}
\end{alist}

\vspace{4mm}

% --- Άσκηση 37157 (37157.tex) ---
\Askhsh{37157}{4}
\begin{ylh}{1}
    \item 3.6. Κριτήρια ισότητας ορθογώνιων τριγώνων
    \item 3.11. Ανισοτικές σχέσεις πλευρών και γωνιών
    \item 4.2. Τέμνουσα δύο ευθειών - Ευκλείδειο αίτημα
    \item 5.2. Παραλληλόγραμμα
    \item 5.4. Ρόμβος
    \item 5.6. Εφαρμογές στα τρίγωνα.
\end{ylh}
% ===================================================================
%  Αρχείο: 37157.tex (Fragment)
%  Αυτόματη μετατροπή μέσω doc2latex.py
% ===================================================================



Έστω τρίγωνο $AB\varGamma $ με διάμεσο $AM$ τέτοια ώστε $AM$=$AB$. Φέρουμε το ύψος του $AK$ και το προεκτείνουμε (προς το Κ) κατά τμήμα $K\varDelta  $=$AK$. Προεκτείνουμε τη διάμεσο $AM$ (προς το Μ) κατά τμήμα $ME$=$AM$.

Να αποδείξετε ότι:

\begin{alist}
\item ΔΕ$\bot$ΑΔ και $\varDelta  E$=2ΚΜ. \monades{7}

\item Το τετράπλευρο Α$BE$Γ είναι παραλληλόγραμμο. \monades{6}

\item Το τετράπλευρο Α$B\varDelta  $Μ είναι ρόμβος. \monades{6}

\item Η προέκταση της $\varDelta  M$ τέμνει το $A\varGamma $ στο μέσον του Ζ. \monades{6}

%\includesvg[width=5.78125in,height=4.59593in]{/home/spyros/Μαθηματικά/Trapeza_Thematwn/A_Lykeiou/Geometry/extracted/37157/media/image2.svg}
\end{alist}

\vspace{4mm}

% --- Άσκηση 37158 (37158.tex) ---
\Askhsh{37158}{4}
\begin{ylh}{1}
    \item 3.2. 1ο Κριτήριο ισότητας τριγώνων
    \item 3.11. Ανισοτικές σχέσεις πλευρών και γωνιών
    \item 4.2. Τέμνουσα δύο ευθειών - Ευκλείδειο αίτημα
    \item 4.6. Άθροισμα γωνιών τριγώνου
    \item 5.2. Παραλληλόγραμμα
    \item 5.6. Εφαρμογές στα τρίγωνα.
\end{ylh}
% ===================================================================
%  Αρχείο: 37158.tex (Fragment)
%  Αυτόματη μετατροπή μέσω doc2latex.py
% ===================================================================



Έστω ορθογώνιο τρίγωνο $AB\varGamma $ με $\hat{A} = \ 90\degree$ και $\varDelta$ και $E$ τα μέσα των $AB$ και $A\varGamma $ αντίστοιχα. Στο τμήμα $B\varGamma $ θεωρούμε σημεία $K$ και Λ ώστε $\varDelta  K$ = $KB$ και $E\varLambda $ = $\varLambda \varGamma $.

Να αποδείξετε ότι:

\begin{alist}
\item $\varDelta  \hat{Κ}\Lambda = 2 \cdot \ \hat{B}$ και $E\hat{\Lambda}Κ = 2 \cdot \ \hat{\varGamma }$ . \monades{10}

\item Το τετράπλευρο $\varDelta  E\varLambda K$ είναι παραλληλόγραμμο. \monades{9}

\item $\varDelta   E = 2 \cdot \varDelta   Κ$. \monades{6}

%\includesvg[width=6.29861in,height=3.44514in]{/home/spyros/Μαθηματικά/Trapeza_Thematwn/A_Lykeiou/Geometry/extracted/37158/media/image2.svg}
\end{alist}

\vspace{4mm}

% --- Άσκηση 37159 (37159.tex) ---
\Askhsh{37159}{4}
\begin{ylh}{1}
    \item 3.2. 1ο Κριτήριο ισότητας τριγώνων
    \item 3.11. Ανισοτικές σχέσεις πλευρών και γωνιών
    \item 4.6. Άθροισμα γωνιών τριγώνου
    \item 5.6. Εφαρμογές στα τρίγωνα.
\end{ylh}
% ===================================================================
%  Αρχείο: 37159.tex (Fragment)
%  Αυτόματη μετατροπή μέσω doc2latex.py
% ===================================================================



Έστω τρίγωνο $AB\varGamma $ ($AB$\textgreater $A\varGamma $), $A\varDelta  $ το ύψος του και $M$ το μέσο του $AB$. Η προέκταση της $M\varDelta  $ τέμνει την προέκταση της $A\varGamma $ στο σημείο $E$ ώστε $\varGamma \varDelta  =\varGamma E$.

Να αποδείξετε ότι:

\begin{alist}
\item $\hat{B} = \ \hat{E}$. \monades{8}

\item $\ \hat{\varGamma } = \ 2 \cdot \hat{B} = \ A\hat{M}\varDelta  $. \monades{10}

\item $\varGamma E$ < $A\varGamma $. \monades{7}

%\includesvg[width=6.673in,height=3.15625in]{/home/spyros/Μαθηματικά/Trapeza_Thematwn/A_Lykeiou/Geometry/extracted/37159/media/image2.svg}
\end{alist}

\vspace{4mm}

% --- Άσκηση 37160 (37160.tex) ---
\Askhsh{37160}{4}
\begin{ylh}{1}
    \item 3.4. 3ο Κριτήριο ισότητας τριγώνων
    \item 3.11. Ανισοτικές σχέσεις πλευρών και γωνιών
    \item 4.2. Τέμνουσα δύο ευθειών - Ευκλείδειο αίτημα
    \item 5.2. Παραλληλόγραμμα.
\end{ylh}
% ===================================================================
%  Αρχείο: 37160.tex (Fragment)
%  Αυτόματη μετατροπή μέσω doc2latex.py
% ===================================================================



Έστω τρίγωνο $AB\varGamma $, $A\varDelta  $ η διχοτόμος της γωνίας $\hat{A}$ και $M$ το μέσον της $AB$. Η κάθετη από το $M$ στην $A\varDelta  $ τέμνει το $A\varGamma $ στο Ε. Η παράλληλη από το Β στο $A\varGamma $ τέμνει την προέκταση της $A\varDelta  $ στο Κ και την προέκταση της $EM$ στο Λ.

Να αποδείξετε ότι:

\begin{alist}
\item Τα τρίγωνα $AEM$, $MB\varLambda $ και $ABK$ είναι ισοσκελή. \monades{15}

\item Το τετράπλευρο ΑΛ$BE$ είναι παραλληλόγραμμο. \monades{10}

%\includesvg[width=5.92708in,height=4.73958in]{/home/spyros/Μαθηματικά/Trapeza_Thematwn/A_Lykeiou/Geometry/extracted/37160/media/image2.svg}
\end{alist}

\vspace{4mm}

% --- Άσκηση 37161 (37161.tex) ---
\Askhsh{37161}{4}
\begin{ylh}{1}
    \item 3.2. 1ο Κριτήριο ισότητας τριγώνων
    \item 3.4. 3ο Κριτήριο ισότητας τριγώνων
    \item 3.6. Κριτήρια ισότητας ορθογώνιων τριγώνων
    \item 3.11. Ανισοτικές σχέσεις πλευρών και γωνιών
    \item 4.2. Τέμνουσα δύο ευθειών - Ευκλείδειο αίτημα
    \item 4.6. Άθροισμα γωνιών τριγώνου
    \item 5.10. Τραπέζιο
    \item 5.11. Ισοσκελές τραπέζιο.
\end{ylh}
% ===================================================================
%  Αρχείο: 37161.tex (Fragment)
%  Αυτόματη μετατροπή μέσω doc2latex.py
% ===================================================================



Έστω ισοσκελές τρίγωνο $AB\varGamma $ ($AB=A\varGamma $) και $A\varDelta  $ διάμεσος. Στο τμήμα $A\varDelta  $ θεωρούμε τυχαίο σημείο $K$ από το οποίο φέρνουμε τα τμήματα $KZ$ και $KE$ κάθετα στις $AB$ και $A\varGamma $ αντίστοιχα.

\begin{alist}
\item Να αποδείξετε ότι τα τρίγωνα Κ$B\varGamma $ και $KZE$ είναι ισοσκελή. \monades{8}

\item Να αποδείξετε ότι το τετράπλευρο $ZE\varGamma B$ είναι ισοσκελές τραπέζιο. \monades{10}

\item Ένας μαθητής στην πορεία της λύσης του έδωσε το εξής επιχείρημα:

«\emph{Το τμήμα $A\varDelta  $ είναι διάμεσος στη βάση ισοσκελούς άρα ύψος και διχοτόμος του τριγώνου $AB\varGamma $ και μεσοκάθετος του $B\varGamma $. Οπότε και το τρίγωνο $BK\varGamma $ είναι ισοσκελές.}

\emph{Τα τρίγωνα $ABK$, $A\varGamma K$ έχουν:}


\end{alist}
\begin{alist}
\item
  \emph{$BK$ = $K\varGamma $}
\item
  $Β\hat{A}Κ = \ \varGamma \hat{A}Κ$ \emph{επειδή $AK$ διχοτόμος της} $\hat{A}$\emph{.}
\item
  $Α\hat{B}Κ = \ Α\hat{\varGamma }Κ$ \emph{ως διαφορές ίσων γωνιών ισοσκελών τριγώνων.}

\emph{Άρα τα τρίγωνα είναι ίσα βάση του κριτηρίου Γωνία Πλευρά Γωνία}.»

Ο καθηγητής είπε ότι η απάντησή του είναι ελλιπής. Να συμπληρώσετε την απάντηση του μαθητή ώστε να ικανοποιεί το κριτήριο Γωνία --Πλευρά- Γωνία διατηρώντας τις πλευρές $BK$ και $K\varGamma $. \monades{7}

%\includesvg[width=3.18513in,height=3.39583in]{/home/spyros/Μαθηματικά/Trapeza_Thematwn/A_Lykeiou/Geometry/extracted/37161/media/image2.svg}
\end{alist}

\vspace{4mm}

% --- Άσκηση 37162 (37162.tex) ---
\Askhsh{37162}{4}
\begin{ylh}{1}
    \item 3.2. 1ο Κριτήριο ισότητας τριγώνων
    \item 4.2. Τέμνουσα δύο ευθειών - Ευκλείδειο αίτημα
    \item 5.6. Εφαρμογές στα τρίγωνα
    \item 5.9. Μια ιδιότητα του ορθογώνιου τριγώνου
    \item 5.10. Τραπέζιο
    \item 5.11. Ισοσκελές τραπέζιο.
\end{ylh}
% ===================================================================
%  Αρχείο: 37162.tex (Fragment)
%  Αυτόματη μετατροπή μέσω doc2latex.py
% ===================================================================



Δίνεται τρίγωνο $AB\varGamma $ με $AB$ < $A\varGamma $ και το ύψος του $AH$. Αν Δ, Ε και $Z$ είναι τα μέσα των $AB$,$A\varGamma $ και $B\varGamma $ αντίστοιχα, να αποδείξετε ότι :

\begin{alist}
\item το τετράπλευρο $\varDelta  EZH$ είναι ισοσκελές τραπέζιο. \monades{8}

\item οι γωνίες $H\hat{\varDelta  }Ζ$ και $H\hat{Ε}Ζ$ είναι ίσες . \monades{8}

\item οι γωνίες $Ε\hat{\varDelta  }Ζ$ και $Ε\hat{H}Ζ$ είναι ίσες. \monades{9}
\end{alist}

\vspace{4mm}

% --- Άσκηση 37163 (37163.tex) ---
\Askhsh{37163}{4}
\begin{ylh}{1}
    \item 3.2. 1ο Κριτήριο ισότητας τριγώνων
    \item 3.7. Κύκλος - Μεσοκάθετος - Διχοτόμος
    \item 5.8. Το ορθόκεντρο τριγώνου.
\end{ylh}
% ===================================================================
%  Αρχείο: 37163.tex (Fragment)
%  Αυτόματη μετατροπή μέσω doc2latex.py
% ===================================================================



Δίνεται τρίγωνο $AB\varGamma $ με $AB<A\varGamma $ και η διχοτόμος του $A\varDelta  $. Στην πλευρά $A\varGamma $ θεωρούμε σημείο $E$ τέτοιο ώστε $AE$ = $AB$.

Να αποδείξετε ότι :

\begin{alist}
\item τα τρίγωνα Α$B\varDelta  $ και $A\varDelta  E$ είναι ίσα. \monades{7}

\item η ευθεία $A\varDelta  $ είναι μεσοκάθετος του τμήματος $BE$. \monades{9}

\item αν το ύψος από την κορυφή $B$ του τριγώνου $AB\varGamma $ τέμνει την $A\varDelta  $ στο Η τότε η ευθεία $EH$ είναι κάθετη στην $AB$. \monades{9}

%\includesvg[width=4.05208in,height=3.14008in]{/home/spyros/Μαθηματικά/Trapeza_Thematwn/A_Lykeiou/Geometry/extracted/37163/media/image2.svg}
\end{alist}

\vspace{4mm}

% --- Άσκηση 37164 (37164.tex) ---
\Askhsh{37164}{4}
\begin{ylh}{1}
    \item 3.4. 3ο Κριτήριο ισότητας τριγώνων
    \item 3.11. Ανισοτικές σχέσεις πλευρών και γωνιών
    \item 3.12. Τριγωνική ανισότητα
    \item 4.2. Τέμνουσα δύο ευθειών - Ευκλείδειο αίτημα
    \item 4.4. Γωνίες με πλευρές παράλληλες
    \item 4.6. Άθροισμα γωνιών τριγώνου.
\end{ylh}
% ===================================================================
%  Αρχείο: 37164.tex (Fragment)
%  Αυτόματη μετατροπή μέσω doc2latex.py
% ===================================================================



Έστω ισοσκελές τρίγωνο $AB\varGamma $ ($AB=A\varGamma $) και $M$ το μέσο της $B\varGamma $. Φέρουμε ΓΔ$\bot$$B\varGamma $ με $\varGamma \varDelta  $=$AB$ (Α,Δ εκατέρωθεν της $B\varGamma $).

Nα αποδείξετε ότι:

\begin{alist}
\item $AM$ // $\varGamma \varDelta  $. \monades{6}

\item η $A\varDelta  $ είναι διχοτόμος της γωνίας $Μ\hat{A}\varGamma $. \monades{7}

\item $\varDelta  \hat{A}\varGamma  = 45\degree - \ \frac{\hat{B}}{2}$. \monades{7}

\item $A\varDelta  $ < 2 $AB$. \monades{5}

%\includesvg[width=3.76042in,height=5.05208in]{/home/spyros/Μαθηματικά/Trapeza_Thematwn/A_Lykeiou/Geometry/extracted/37164/media/image2.svg}
\end{alist}

\vspace{4mm}

% --- Άσκηση 37165 (37165.tex) ---
\Askhsh{37165}{4}
\begin{ylh}{1}
    \item 3.4. 3ο Κριτήριο ισότητας τριγώνων
    \item 4.2. Τέμνουσα δύο ευθειών - Ευκλείδειο αίτημα
    \item 5.3. Ορθογώνιο
    \item 5.4. Ρόμβος
    \item 5.6. Εφαρμογές στα τρίγωνα.
\end{ylh}
% ===================================================================
%  Αρχείο: 37165.tex (Fragment)
%  Αυτόματη μετατροπή μέσω doc2latex.py
% ===================================================================



Δίνεται οξυγώνιο τρίγωνο $AB\varGamma $ με $AB<A\varGamma $. Από το Β φέρουμε κάθετη στην διχοτόμο $AM$ της γωνίας Α, η οποία τέμνει την $AM$ στο Η και την $A\varGamma $ στο Δ. Στην προέκταση της $AH$ (προς το Η) θεωρούμε σημείο $Z$ τέτοιο ώστε $AH$ = $HZ$ και έστω Θ το μέσο της πλευράς $B\varGamma $.

Να αποδείξετε ότι

\begin{alist}
\item το τετράπλευρο $ABZ\varDelta  $ είναι ρόμβος. \monades{9}

\item $H\varTheta $ // $BZ$. \monades{9}

\item $H\varTheta =\frac{A\varGamma  - AB}{2}$ \monades{7}

%\includesvg[width=3.97917in,height=3.79167in]{/home/spyros/Μαθηματικά/Trapeza_Thematwn/A_Lykeiou/Geometry/extracted/37165/media/image2.svg}
\end{alist}

\vspace{4mm}

% --- Άσκηση 37166 (37166.tex) ---
\Askhsh{37166}{4}
\begin{ylh}{1}
    \item 3.2. 1ο Κριτήριο ισότητας τριγώνων
    \item 3.6. Κριτήρια ισότητας ορθογώνιων τριγώνων
    \item 5.2. Παραλληλόγραμμα.
\end{ylh}
% ===================================================================
%  Αρχείο: 37166.tex (Fragment)
%  Αυτόματη μετατροπή μέσω doc2latex.py
% ===================================================================



Στο παρακάτω σχήμα φαίνονται οι θέσεις στο χάρτη πέντε χωριών Α, Β, Γ, Δ και $E$ και οι δρόμοι που τα συνδέουν. Το χωριό $E$ ισαπέχει από τα χωριά Β, Γ και επίσης από τα χωριά Α και Δ.
\begin{alist}
\item Να αποδείξετε ότι:


\end{alist}
\begin{alist}
\item
  η απόσταση των χωριών Α και Β είναι ίση με την απόσταση των χωριών $\varGamma$ και Δ. \monades{5}
\item
  αν οι δρόμοι $AB$ και $\varGamma \varDelta  $ έχουν δυνατότητα να προεκταθούν, να αποδείξετε ότι αποκλείεται να συναντηθούν. \monades{5}
\item
  τα χωριά Β και $\varGamma$ ισαπέχουν από τον δρόμο $A\varDelta  $. \monades{8}

\item Να προσδιορίσετε γεωμετρικά το σημείο του δρόμου $A\varGamma $ που ισαπέχει από τα χωριά Α και Δ. \monades{7}
\end{alist}

\vspace{4mm}

% --- Άσκηση 37167 (37167.tex) ---
\Askhsh{37167}{4}
\begin{ylh}{1}
    \item 3.2. 1ο Κριτήριο ισότητας τριγώνων
    \item 3.3. 2ο Κριτήριο ισότητας τριγώνων
    \item 3.11. Ανισοτικές σχέσεις πλευρών και γωνιών
    \item 4.2. Τέμνουσα δύο ευθειών - Ευκλείδειο αίτημα
    \item 5.2. Παραλληλόγραμμα
    \item 5.3. Ορθογώνιο.
\end{ylh}
% ===================================================================
%  Αρχείο: 37167.tex (Fragment)
%  Αυτόματη μετατροπή μέσω doc2latex.py
% ===================================================================



Δίνεται παραλληλόγραμμο $AB\varGamma \varDelta  $ με $AB$ > $A\varDelta  $ και οι διχοτόμοι των γωνιών του $AP$, $BE$, $\varGamma \varSigma$ και $\varDelta  T$ (όπου Ρ, Ε στην $\varDelta  \varGamma $ και Σ, Τ στην $AB$) τέμνονται στα σημεία $K$, Λ, M και N όπως φαίνεται στο παρακάτω σχήμα.

Να αποδείξετε ότι:

\begin{alist}
\item το τετράπλευρο $\varDelta  EBT$ είναι παραλληλόγραμμο. \monades{7}

\item το τετράπλευρο $K\varLambda MN$ είναι ορθογώνιο. \monades{8}

\item $\varLambda N$ // $AB$ \monades{5}

\item $\varLambda N$ = $AB$ -- $A\varDelta  $ \monades{5}

%\includesvg[width=5.52921in,height=2.63542in]{/home/spyros/Μαθηματικά/Trapeza_Thematwn/A_Lykeiou/Geometry/extracted/37167/media/image2.svg}
\end{alist}

\vspace{4mm}

% --- Άσκηση 37823 (37823.tex) ---
\Askhsh{37823}{4}
\begin{ylh}{1}
    \item 3.2. 1ο Κριτήριο ισότητας τριγώνων
    \item 3.4. 3ο Κριτήριο ισότητας τριγώνων
    \item 3.11. Ανισοτικές σχέσεις πλευρών και γωνιών.
\end{ylh}
% ===================================================================
%  Αρχείο: 37823.tex (Fragment)
%  Αυτόματη μετατροπή μέσω doc2latex.py
% ===================================================================



Δίνεται οξυγώνιο ισοσκελές τρίγωνο $ΑΒ\varGamma $ ($A\varGamma  = B\varGamma $). Η μεσοκάθετη $\varepsilon$ της $A\varGamma $ τέμνει την προέκταση της $AB$ (προς το μέρος του $B$) στο σημείο $M$ και την $A\varGamma $ στο $Z$. Στην προέκταση της $M\varGamma $ (προς το μέρος του $\varGamma $) παίρνουμε σημείο $E$ τέτοιο ώστε $\varGamma  E = BM$.

\begin{alist}
\item Να δείξετε ότι το τρίγωνο $AM\varGamma $ είναι ισοσκελές. \monades{8}
\item Να δειχτεί ότι τα τρίγωνα $A\varGamma  E$ και $\varGamma  BM$ είναι ίσα. \monades{10}
\item Να δειχτεί ότι το τρίγωνο $AME$ είναι ισοσκελές. \monades{7}

\begin{center}
\begin{tikzpicture}[scale=1]
\tkzDefPoint(0,0){A}
\tkzDefPoint(3,0){B}
% $AC$=$BC$. Let C(1.5, y).
\tkzDefPoint(1.5,4){C}
\tkzDefMidPoint(A,C) \tkzGetPoint{Z}
\tkzDefLine[orthogonal=through Z](A,C) \tkzGetPoint{eps_dir}
\tkzInterLL(A,B)(Z,eps_dir) \tkzGetPoint{M}
\tkzDefPointBy[homothety=center G ratio -1.4](M) \tkzGetPoint{E_dummy}
\tkzCalcLength(B,M)\tkzGetLength{rBM}
\tkzDefCircle[R](C,\rBM) \tkzGetPoint{circle_r}
\tkzInterLC(M,C)(C,circle_r) \tkzGetPoints{E}{E_prime}

\draw[l3] (A)--(B)--(C)--cycle;
\draw[l3] (B)--(M);
\draw[l3] (A)--(M);
\draw[l3] (M)--(E);
\draw[l3] (A)--(E);
\draw[l3, dashed] (Z)--(M);

\tkzDrawPoints(A,B,C,Z,M,E)
\tkzLabelPoint[below left](A){$A$}
\tkzLabelPoint[below](B){$B$}
\tkzLabelPoint[above](C){$\varGamma $}
\tkzLabelPoint[above left](Z){$Z$}
\tkzLabelPoint[below right](M){$M$}
\tkzLabelPoint[above right](E){$E$}

\tkzMarkRightAngle(M,Z,A)
\tkzMarkSegments[mark=s|](A,C B,C)
\end{tikzpicture}
\end{center}
\end{alist}

\vspace{4mm}

\end{document}
